% Generated mechanically from frozen input inputs/p06/ja/chow.tex.
% Do not edit this generated file; edit the bound locale source instead.

% OK, start here.
%

\setcounter{chapter}{41}
\chapter{Chow ホモロジーと Chern 類}


\phantomsection
\label{chow-section-phantom}




\section{序}
\label{chow-section-introduction}

\noindent
本章では Chow ホモロジー群、およびベクトル束の Chern 類を
作用素的 Chow コホモロジー群の元として構成する方法を論じる
（係数はすべて $\mathbf{Z}$ とする）。

\medskip\noindent
本章は、鍵となる補題 \ref{chow-lemma-milnor-gersten-low-degree} の
可能な限り短い代数的証明から始める。まず Herbrand 商を定義し
（第 \ref{chow-section-periodic-complexes} 節）、いくつかの場合にそれを計算する
（第 \ref{chow-section-calculation} 節）。次に、修正を施した後に適切な
分解が存在することに関する簡単な代数的補題を証明する
（第 \ref{chow-section-preparation-tame-symbol} 節）。これらを用いて
第 \ref{chow-section-tame-symbol} 節で tame 記号を構成し、定義する。
鍵となる補題の証明に必要なのは tame 記号の最も基本的な性質だけであり、
その証明は第 \ref{chow-section-key-lemma} 節で与える。

\medskip\noindent
続いて第 \ref{chow-section-setup} 節で、本章の残りを通じて用いる基本設定を導入する。
内容に多少の奥行きをもたせるため、通常よりやや一般的な場合を扱う。
すなわち、スキーム $X$ は、普遍的鎖状で次元関数を備えた固定された
局所 Noether 基底スキーム上で局所有限型であると仮定する。
これらの仮定のもとで Chow ホモロジー群 $\CH_*(X)$ と、
それに対する Chern 類とのキャップ積の作用を定義できる。
これは、このような基底上で局所有限型な代数スタックに対しても
同じ対象を定義できるはずであることを示唆している。

\medskip\noindent
次に \cite{F} の冒頭数章に従って、サイクル、平坦引き戻し、
固有押し出し、および有理同値を定義する。ただし、現れるサイクルの
台については、そこでの記述ほど厳密には扱わない。

\medskip\noindent
ここでは \cite{F} の提示とは異なり、上記の鍵となる補題を用いて、
第 \ref{chow-section-key} 節で基本的な可換関係を証明する。これにより、
可逆層との交差演算が有理同値を通して定まり、かつ可換であることを示す
（第 \ref{chow-section-commutativity} 節）。さらに鍵となる補題をもう一度適用すると、
有効 Cartier 因子の Gysin 写像が有理同値を通して定まることが分かる
（第 \ref{chow-section-gysin} 節）。ここまで証明すれば、ベクトル束の Chern 類を定義し、
加法性と分裂原理を証明し、Chern 指標と Todd 類を導入して、
Grothendieck--Riemann--Roch 定理を述べることは直ちにできる。

\medskip\noindent
付録は二つある。付録 A（第 \ref{chow-section-appendix-A} 節）では、
tame 記号の別の、より長い構成と、それに対応する鍵となる補題の証明を論じる。
最後に付録 B（第 \ref{chow-section-appendix-chow} 節）では、
連接層の $K$ 理論との関係を簡潔に述べ、いくつかのブローアップ補題を扱う。
詳しくは各付録の序を参照されたい。

\medskip\noindent
次章では、代数閉体上の滑らかな射影多様体の Chow 群 $\CH_*(X)$ に立ち返る。
\cite{Samuel}、\cite{ChevalleyI}、\cite{ChevalleyII} にあるような移動補題と
Serre の Tor 公式（\cite{Serre_local_algebra} または
\cite{Serre_algebre_locale} を参照）を用いて、$\CH_*(X)$ 上に環構造を定義する。
「相交理論」第 \ref{intersection-section-introduction} 節以降を参照されたい。








\section{周期複体と Herbrand 商}
\label{chow-section-periodic-complexes}

\noindent
周期複体には、もちろん非常に一般的な概念がある。
写像の周期性を課すことも、対象の周期性を課すこともできる。
必要に応じて後にそれらを導入するが、当面必要なのは次の場合だけである。

\begin{definition}
\label{chow-definition-periodic-complex}
$R$ を環とする。
\begin{enumerate}
\item $R$ 上の {\it $2$-周期複体}とは、$R$-加群 $M$, $N$ と
$R$-加群準同型 $\varphi : M \to N$, $\psi : N \to M$ からなる四つ組
$(M, N, \varphi, \psi)$ であって、
$$
\xymatrix{
\ldots \ar[r] &
M \ar[r]^\varphi &
N \ar[r]^\psi &
M \ar[r]^\varphi &
N \ar[r] & \ldots
}
$$
が複体となるものをいう。このとき、この複体の {\it コホモロジー加群}を
次の $R$-加群として定義する：
$$
H^0(M, N, \varphi, \psi) = \Ker(\varphi)/\Im(\psi)
\quad\text{and}\quad
H^1(M, N, \varphi, \psi) = \Ker(\psi)/\Im(\varphi).
$$
これらのコホモロジー加群が零であるとき、$2$-周期複体は {\it 完全}であるという。
\item $R$ 上の {\it $(2, 1)$-周期複体}とは、$R$-加群 $M$ と
$R$-加群準同型 $\varphi : M \to M$, $\psi : M \to M$ からなる三つ組
$(M, \varphi, \psi)$ であって、
$$
\xymatrix{
\ldots \ar[r] &
M \ar[r]^\varphi &
M \ar[r]^\psi &
M \ar[r]^\varphi &
M \ar[r] & \ldots
}
$$
が複体となるものをいう。これは $2$-周期複体の特別な場合なので、
{\it コホモロジー加群} $H^0(M, \varphi, \psi)$, $H^1(M, \varphi, \psi)$ と
完全性の概念をもつ。
\end{enumerate}
\end{definition}

\noindent
以下では、$2$-周期複体について証明した任意の結果を、断りなく
$(2, 1)$-周期複体にも用いる。$2$-周期複体全体が圏をなすことは明らかである。
射 $(f, g) : (M, N, \varphi, \psi) \to (M', N', \varphi', \psi')$ は、
$\varphi' \circ f = g \circ \varphi$ および $\psi' \circ g = f \circ \psi$ を満たす
二つの射 $f : M \to M'$ と $g : N \to N'$ の組である。
核と余核を「ホモロジー」補題 \ref{homology-lemma-cat-chain-abelian} と同様にとれば、
これは Abel 圏となる。

\begin{definition}
\label{chow-definition-periodic-length}
$R$ を環とし、$(M, N, \varphi, \psi)$ を、コホモロジー加群が有限長である
$R$ 上の $2$-周期複体とする。このとき $(M, N, \varphi, \psi)$ の
{\it 重複度}を次の整数として定義する：
$$
e_R(M, N, \varphi, \psi) =
\text{length}_R(H^0(M, N, \varphi, \psi))
-
\text{length}_R(H^1(M, N, \varphi, \psi))
$$
$(2, 1)$-周期複体 $(M, \varphi, \psi)$ の場合には、これを
$e_R(M, \varphi, \psi)$ と書き、{\it （加法的）Herbrand 商}と呼ぶこともある。
\end{definition}

\noindent
$(M, \varphi, \psi)$ のコホモロジー群が有限 Abel 群である場合、通常は
$$
q(M, \varphi, \psi) =
\frac{\# H^0(M, \varphi, \psi)}{\# H^1(M, \varphi, \psi)}
$$
を {\it （乗法的）Herbrand 商}と呼ぶ。すなわち、乗法的 Herbrand 商は
$H^0$ の元の個数を $H^1$ の元の個数で割ったものである。
$R$ が局所環で、その剰余体が $q$ 個の元をもつ有限体ならば、
$$
q(M, \varphi, \psi) = q^{e_R(M, \varphi, \psi)}
$$
となる。

\medskip\noindent
環 $R$ 上の $(2, 1)$-周期複体の例として、$M$ を $R$-加群、
$\psi$ を $R$-線形写像とする任意の三つ組 $(M, 0, \psi)$ がある。
$\psi$ の核と余核が有限長ならば、
\begin{equation}
\label{chow-equation-multiplicity-coker-ker}
e_R(M, 0, \psi) = \text{length}_R(\Coker(\psi)) - \text{length}_R(\Ker(\psi))
\end{equation}
を得る。以下、この記法に関する基本的な補題を述べて証明する。

\begin{lemma}
\label{chow-lemma-additivity-periodic-length}
$R$ を環とし、$2$-周期複体の短完全列
$$
0 \to (M_1, N_1, \varphi_1, \psi_1)
\to (M_2, N_2, \varphi_2, \psi_2)
\to (M_3, N_3, \varphi_3, \psi_3)
\to 0
$$
が与えられているとする。三つのうち二つのコホモロジー加群が有限長ならば、
残る一つも有限長であり、
$$
e_R(M_2, N_2, \varphi_2, \psi_2) =
e_R(M_1, N_1, \varphi_1, \psi_1) +
e_R(M_3, N_3, \varphi_3, \psi_3).
$$
が成り立つ。
\end{lemma}

\begin{proof}
$A = (M_1, N_1, \varphi_1, \psi_1)$,
$B = (M_2, N_2, \varphi_2, \psi_2)$, $C = (M_3, N_3, \varphi_3, \psi_3)$ と略記する。
コホモロジー長完全列
$$
\ldots
\to H^1(C)
\to H^0(A)
\to H^0(B)
\to H^0(C)
\to H^1(A)
\to H^1(B)
\to H^1(C)
\to \ldots
$$
がある。これから有限完全列
$$
0 \to I
\to H^0(A)
\to H^0(B)
\to H^0(C)
\to H^1(A)
\to H^1(B)
\to K \to 0
$$
を得て、さらに $0 \to K \to H^1(C) \to I \to 0$ がフィルトレーションを与える。
長さ関数の加法性（「代数」補題 \ref{algebra-lemma-length-additive}）から結論が従う。
\end{proof}

\begin{lemma}
\label{chow-lemma-finite-periodic-length}
$R$ を環とする。$(M, N, \varphi, \psi)$ を $2$-周期複体とし、
$M$, $N$ が有限長ならば、
$e_R(M, N, \varphi, \psi) = \text{length}_R(M) - \text{length}_R(N)$
である。特に、$(M, \varphi, \psi)$ が $(2, 1)$-周期複体で、
$M$ が有限長ならば $e_R(M, \varphi, \psi) = 0$ である。
\end{lemma}

\begin{proof}
$M$, $N$ が有限長である $2$-周期複体の圏では、量
``$\text{length}_R(M) - \text{length}_R(N)$'' が短完全列に関して加法的である
（厳密な主張は読者に委ねる）。短完全列
$$
0 \to (M, \Im(\varphi), \varphi, 0) \to
(M, N, \varphi, \psi) \to (0, N/\Im(\varphi), 0, 0) \to 0
$$
を考える。最初の注意と補題 \ref{chow-lemma-additivity-periodic-length} の加法性を併せると、
(a) $M = 0$ の場合と (b) $\varphi$ が全射の場合に帰着される。
これらの場合は読者に委ねる。
\end{proof}

\begin{lemma}
\label{chow-lemma-compare-periodic-lengths}
$R$ を環とする。$f : (M, \varphi, \psi) \to (M', \varphi', \psi')$ を、
コホモロジー加群が有限長である $(2, 1)$-周期複体の射とする。
$\Ker(f)$ と $\Coker(f)$ が有限長ならば、
$e_R(M, \varphi, \psi) = e_R(M', \varphi', \psi')$ である。
\end{lemma}

\begin{proof}
補題 \ref{chow-lemma-additivity-periodic-length} の加法性を適用する。
補題 \ref{chow-lemma-finite-periodic-length} により、
$(\Ker(f), \varphi, \psi)$ と $(\Coker(f), \varphi', \psi')$ の重複度は零である。
\end{proof}




\section{いくつかの重複度の計算}
\label{chow-section-calculation}

\noindent
後で特定のサイクルの等しさを証明するために、いくつかの重複度を計算する必要がある。
前節で与えた重複度に関する初等的な補題に加え、主な道具となるのは
「代数」補題 \ref{algebra-lemma-order-vanishing-determinant} である。

\begin{lemma}
\label{chow-lemma-length-multiplication}
$R$ を Noether 局所環、$M$ を有限 $R$-加群、$x \in R$ とする。
次を仮定する。
\begin{enumerate}
\item $\dim(\text{Supp}(M)) \leq 1$, かつ
\item $\dim(\text{Supp}(M/xM)) \leq 0$.
\end{enumerate}
$\text{Supp}(M) = \{\mathfrak m, \mathfrak q_1, \ldots, \mathfrak q_t\}$ と書く。このとき
$$
e_R(M, 0, x) =
\sum\nolimits_{i = 1, \ldots, t}
\text{ord}_{R/\mathfrak q_i}(x)
\text{length}_{R_{\mathfrak q_i}}(M_{\mathfrak q_i}).
$$
\end{lemma}

\begin{proof}
まず準備的な注意を述べる。$M$ が有限長、すなわち $t = 0$ ならば、
左辺も右辺も零なので補題の結論は成り立つ
（補題 \ref{chow-lemma-finite-periodic-length} を参照）。
また、(1), (2) を満たす加群の短完全列 $0 \to M \to M' \to M'' \to 0$ があるとき、
三つのうち二つについて補題が成り立てば、長さの加法性
（「代数」補題 \ref{algebra-lemma-length-additive}）と重複度の加法性
（補題 \ref{chow-lemma-additivity-periodic-length}）により残る一つについても成り立つ。

\medskip\noindent
$M_i$ を $M$ の $M_{\mathfrak q_i}$ における像とする。このとき
$\text{Supp}(M_i) = \{\mathfrak m, \mathfrak q_i\}$ である。
写像 $M \to \bigoplus M_i$ の核と余核は台 $\{\mathfrak m\}$ をもち、したがって有限長である。
準備的な注意から、各 $M_i$ について補題を証明すれば十分である。
ゆえに $\text{Supp}(M) = \{\mathfrak m, \mathfrak q\}$ と仮定してよい。
この場合、「代数」補題 \ref{algebra-lemma-Noetherian-power-ideal-kills-module} により
有限フィルトレーション
$M \supset \mathfrak qM \supset \mathfrak q^2M \supset \ldots \supset
\mathfrak q^nM = 0$ がある。再び加法性により、$M$ が $\mathfrak q$ で消される場合に
補題を証明すれば十分である。このとき $M$ を $R/\mathfrak q$-加群とみなせるので、
$R$ は分数体 $K$ をもつ次元 $1$ の Noether 局所整域であると仮定してよい。
torsion 部分加群、すなわち $M \to M \otimes_R K = V$ の核で割る。
この torsion は有限長であり準備的な注意で処理できるから、
$M \subset V$ は格子であると仮定してよい
（「代数」定義 \ref{algebra-definition-lattice}）。
すると $x : M \to M$ は単射であり、
$\text{length}_R(M/xM) = d(M, xM)$
である（「代数」定義 \ref{algebra-definition-distance}）。
$\text{length}_K(V) = \dim_K(V)$ であるから、
$\det(x : V \to V) = x^{\dim_K(V)}$ および
$\text{ord}_R(\det(x : V \to V)) = \dim_K(V) \text{ord}_R(x)$ を得る。
したがって、この場合の所望の等式は
「代数」補題 \ref{algebra-lemma-order-vanishing-determinant} から従う。
\end{proof}

\begin{lemma}
\label{chow-lemma-additivity-divisors-restricted}
$R$ を Noether 局所環、$x \in R$ とする。$M$ が有限 Cohen--Macaulay $R$-加群で、
$\dim(\text{Supp}(M)) = 1$ かつ $\dim(\text{Supp}(M/xM)) = 0$ ならば、
$$
\text{length}_R(M/xM)
=
\sum\nolimits_i \text{length}_R(R/(x, \mathfrak q_i))
\text{length}_{R_{\mathfrak q_i}}(M_{\mathfrak q_i}).
$$
ここで $\mathfrak q_1, \ldots, \mathfrak q_t$ は $M$ の台の極小素イデアルである。
また、$I \subset R$ をイデアルとし、$x$ が $R/I$ 上の非零因子で
$\dim(R/I) = 1$ ならば、
$$
\text{length}_R(R/(x, I))
=
\sum\nolimits_i \text{length}_R(R/(x, \mathfrak q_i))
\text{length}_{R_{\mathfrak q_i}}((R/I)_{\mathfrak q_i})
$$
である。ここで $\mathfrak q_1, \ldots, \mathfrak q_n$ は $I$ 上の極小素イデアルである。
\end{lemma}

\begin{proof}
いずれも補題 \ref{chow-lemma-length-multiplication} の特別な場合である。
\end{proof}

\noindent
重複度の値を決定できる別の場合を述べる。

\begin{lemma}
\label{chow-lemma-powers-period-length-zero}
$R$ を環、$M$ を $R$-加群とする。$\varphi : M \to M$ を自己準同型、$n > 0$ とし、
$\varphi^n = 0$ で、$\Ker(\varphi)/\Im(\varphi^{n - 1})$ が $R$-加群として有限長であるとする。
このとき
$$
e_R(M, \varphi^i, \varphi^{n - i}) = 0
$$
が $i = 0, \ldots, n$ に対して成り立つ。
\end{lemma}

\begin{proof}
慣例により $\varphi^0 = \text{id}_M$ なので、$i = 0, n$ の場合は自明である。
$t$ の乗法を $\varphi$ で定めて、$M$ を $R[t]$-加群とみなす。
$K_i = \Ker(t^i : M \to M)$ と書き、
$$
a_i = \text{length}_R(K_i/t^{n - i}M),\quad
b_i = \text{length}_R(K_i/tK_{i + 1}),\quad
c_i = \text{length}_R(K_1/t^iK_{i + 1})
$$
とおく。境界値は $a_0 = a_n = b_0 = c_0 = 0$ である。
$K_1/t^iK_{i + 1}$ は、仮定により有限長である $K_1/t^{n - 1}M$ の商なので、
$i < n$ に対して $c_i$ は整数である。
以下、$K_i \cap t^jM = t^jK_{i + j}$ を繰り返し用いる。
$0 < i < n - 1$ に対し、完全列
$$
0 \to
K_1/t^{n - i - 1}K_{n - i} \to
K_{i + 1}/t^{n - i - 1}M \xrightarrow{t} K_i/t^{n - i}M
\to K_i/tK_{i + 1} \to 0
$$
がある。$i$ に関する帰納法により、$i < n$ に対して $a_i$, $b_i$ は整数であり、
$$
c_{n - i - 1} - a_{i + 1} + a_i - b_i = 0
$$
が成り立つ。$0 < i < n - 1$ に対しては短完全列
$$
0 \to
K_i/tK_{i + 1} \to
K_{i + 1}/tK_{i + 2} \xrightarrow{t^i}
K_1/t^{i + 1}K_{i + 2} \to
K_1/t^iK_{i + 1} \to 0
$$
もあり、これから
$$
b_i - b_{i + 1} + c_{i + 1} - c_i = 0
$$
を得る。$b_0 = c_0$ なので $i < n$ に対して $b_i = c_i$ である。
したがって
$$
a_2 = a_1 + b_{n - 2} - b_1,\quad
a_3 = a_2 + b_{n - 3} - b_2,\quad \ldots
$$
を得る。これが所望の $a_i = a_{n - i}$ を含意することは直ちに分かる。
\end{proof}

\begin{lemma}
\label{chow-lemma-multiply-period-length}
$(R, \mathfrak m)$ を Noether 局所環とする。$(M, \varphi, \psi)$ を、
$M$ が有限でコホモロジー群が $R$ 上有限長である $R$ 上の $(2, 1)$-周期複体とする。
$x \in R$ が $\dim(\text{Supp}(M/xM)) \leq 0$ を満たすとする。このとき
$$
e_R(M, x\varphi, \psi) = e_R(M, \varphi, \psi) - e_R(\Im(\varphi), 0, x)
$$
および
$$
e_R(M, \varphi, x\psi) = e_R(M, \varphi, \psi) + e_R(\Im(\psi), 0, x)
$$
が成り立つ。
\end{lemma}

\begin{proof}
第二の公式もまったく同様に証明されるので、第一の公式だけを証明する。
$M' = M[x^\infty]$ を $M$ の $x$-冪 torsion 部分加群とする。
短完全列 $0 \to M' \to M \to M'' \to 0$ を考える。
すると $M''$ は $x$-冪 torsion をもたない
（「代数についてさらに」補題 \ref{more-algebra-lemma-divide-by-torsion}）。
$\varphi$, $\psi$ は $M'$ を $M'$ に写すので、$(2, 1)$-周期複体の短完全列
$$
0 \to (M', \varphi', \psi') \to (M, \varphi, \psi) \to
(M'', \varphi'', \psi'') \to 0
$$
を得る。また、短完全列
$0 \to M' \cap \Im(\varphi) \to \Im(\varphi) \to \Im(\varphi'') \to 0$ も得る。
補題 \ref{chow-lemma-compare-periodic-lengths} により
$e_R(M', \varphi, \psi) = e_R(M', x\varphi, \psi) =
e_R(M' \cap \Im(\varphi), 0, x) = 0$ である。
加法性（補題 \ref{chow-lemma-additivity-periodic-length}）から、
$(M'', \varphi'', \psi'')$ について補題を証明すれば十分である。
これは次の段落で扱う場合に帰着する。

\medskip\noindent
$x : M \to M$ が単射であると仮定する。このとき
$\Ker(x\varphi) = \Ker(\varphi)$ である。一方、短完全列
$$
0 \to \Im(\varphi)/x\Im(\varphi) \to
\Ker(\psi)/\Im(x\varphi) \to \Ker(\psi)/\Im(\varphi) \to 0
$$
がある。これと (\ref{chow-equation-multiplicity-coker-ker}) から公式が従う。
\end{proof}







\section{tame 記号のための準備}
\label{chow-section-preparation-tame-symbol}

\noindent
本節では、次節で tame 記号を定義する際に用いる補題をいくつか与える。

\begin{lemma}
\label{chow-lemma-glue-at-max}
$A$ を Noether 環とし、$\mathfrak m_1, \ldots, \mathfrak m_r$ を
$A$ の相異なる極大イデアルとする。$i = 1, \ldots, r$ に対し、
$\varphi_i : A_{\mathfrak m_i} \to B_i$ を、その核と余核が
$\mathfrak m_i$ のある冪で消される環準同型とする。このとき、環準同型
$\varphi : A \to B$ で次を満たすものが存在する。
\begin{enumerate}
\item $\mathfrak m_i$ における $\varphi$ の局所化は $\varphi_i$ と同型であり、
\item $\Ker(\varphi)$ と $\Coker(\varphi)$ は
$\mathfrak m_1 \cap \ldots \cap \mathfrak m_r$ のある冪で消される。
\end{enumerate}
さらに、各 $\varphi_i$ が有限、単射、または全射ならば、$\varphi$ も同じ性質をもつ。
\end{lemma}

\begin{proof}
$I = \mathfrak m_1 \cap \ldots \cap \mathfrak m_r$ とおき、
$A_i = A_{\mathfrak m_i}$ および $A' = \prod A_i$ とおく。
すると $IA' = \prod \mathfrak m_i A_i$ であり、$A \to A'$ は
$A/I \cong A'/IA'$ を満たす平坦環準同型である。
したがって「代数についてさらに」補題
\ref{more-algebra-lemma-application-formal-glueing} により、$A$-加群準同型
$\varphi : A \to B$ が存在し、その $\mathfrak m_i$ における局所化 $\varphi$ は
$\varphi_i$ と同型である。次に「代数についてさらに」注意
\ref{more-algebra-remark-formal-glueing-algebras} の議論を用い、
$B_i$ 上の所与の $A$-代数構造と整合する $A$-代数構造を $B$ に入れる。
補題の最後の主張は、
$\Ker(\varphi)_{\mathfrak m_i} \cong \Ker(\varphi_i)$ および
$\Coker(\varphi)_{\mathfrak m_i} \cong \Coker(\varphi_i)$ から直ちに従う。
\end{proof}

\noindent
次の補題は「代数」補題 \ref{algebra-lemma-nonregular-dimension-one} と非常によく似ている。

\begin{lemma}
\label{chow-lemma-Noetherian-domain-dim-1-two-elements}
$(R, \mathfrak m)$ を次元 $1$ の Noether 局所環とし、
$a, b \in R$ を非零因子とする。有限環拡大 $R \subset R'$ で、
$R'/R$ が $\mathfrak m$ のある冪で消され、さらに非零因子
$t, a', b' \in R'$ が存在して、$a = ta'$, $b = tb'$ および
$R' = a'R' + b'R'$ を満たすものがある。
\end{lemma}

\begin{proof}
$a$ または $b$ が単元ならば、$R = R'$ として結論が成り立つ。
したがって $a, b \in \mathfrak m$ と仮定してよい。$I = (a, b)$ とおく。
着想は $I$ に沿って $R$ をブローアップすることである。
代数的議論の代わりに幾何学的に論じる。
$X = \text{Proj}(\bigoplus_{d \geq 0} I^d)$ とおく。
「因子」補題 \ref{divisors-lemma-blowing-up-gives-effective-Cartier-divisor} により、
射 $X \to \Spec(R)$ は穴あきスペクトル
$U = \Spec(R) \setminus \{\mathfrak m\}$ 上で同型である。
そこで $U$ を $X$ の開部分スキームとみなす。
射 $X \to \Spec(R)$ は「因子」補題 \ref{divisors-lemma-blowing-up-projective} により射影的である。
また、例えば「因子」補題 \ref{divisors-lemma-blow-up-and-irreducible-components} により、
$X$ のすべての生成点は $U$ に属する。
「多様体」補題 \ref{varieties-lemma-finite-in-codim-1} から
$X \to \Spec(R)$ は有限である。ゆえに $X = \Spec(R')$ は affine で、
$R \to R'$ は有限である。$U = D(a)$ なので $R_a \cong R'_a$ である。
したがって $a$ のある冪が有限 $R$-加群 $R'/R$ を消す。
$\mathfrak m = \sqrt{(a)}$ であるから、$R'/R$ は $\mathfrak m$ のある冪で消される。
「因子」補題 \ref{divisors-lemma-blowing-up-gives-effective-Cartier-divisor} により、
$IR'$ は局所単項イデアルである。$R'$ は半局所なので $IR'$ は単項である
（「代数」補題 \ref{algebra-lemma-locally-free-semi-local-free} を参照）。
$IR' = (t)$ と書けば、$a = a't$, $b = b't$ となり、結論は明らかである。
\end{proof}

\begin{lemma}
\label{chow-lemma-not-infinitely-divisible}
$(R, \mathfrak m)$ を次元 $1$ の Noether 局所環とする。
$a, b \in R$ を $a \in \mathfrak m$ を満たす非零因子とする。
整数 $n = n(R, a, b)$ で、任意の有限環拡大 $R \subset R'$ と
$c \in R'$ に対し、$b = a^m c$ ならば $m \leq n$ となるものが存在する。
\end{lemma}

\begin{proof}
極小素イデアル $\mathfrak q \subset R$ を選ぶ。
$\dim(R/\mathfrak q) = 1$ であり、特に $R/\mathfrak q$ は体でない。
「代数」補題 \ref{algebra-lemma-dominate-by-dimension-1} により、
$R/\mathfrak q$ と同じ分数体をもち、それを支配する離散付値環 $A$ を選べる。
$R$ の非零因子は $\mathfrak q$ に含まれないので、$a$, $b$ の $A$ における像は非零である。
$v$ を $A$ 上の離散付値とする。$a \in \mathfrak m$ なので $v(a) > 0$ である。
$n = v(b)/v(a)$ が所望の整数であると主張する。

\medskip\noindent
$R \subset R'$ が与えられたとし、$A' = A \otimes_R R'$ とおく。
$\Spec(R') \to \Spec(R)$ は全射なので
（「代数」補題 \ref{algebra-lemma-integral-overring-surjective}）、
$\Spec(A') \to \Spec(A)$ も全射である
（「代数」補題 \ref{algebra-lemma-surjective-spec-radical-ideal}）。
$(0) \subset A$ の上にある素イデアル $\mathfrak q' \subset A'$ を選ぶ。
すると $A \subset A'' = A'/\mathfrak q'$ は有限環拡大であり、
再びスペクトル上の全射を誘導する。
$A$ の極大イデアルの上にある極大イデアル $\mathfrak m'' \subset A''$ と、
$A''_{\mathfrak m''}$ を支配する離散付値環 $A'''$ を選ぶ（上で引用した補題を参照）。
このとき $A \to A'''$ は離散付値環の拡大で、$A'''$ において
$b = a^m c$ である。ゆえに $v'''(b) \geq mv'''(a)$ である。
$A'''/A$ の分岐指数を $e$ とすれば $v''' = ev$ なので、所望の $m \leq n$ を得る。
\end{proof}

\begin{lemma}
\label{chow-lemma-prepare-tame-symbol}
$(A, \mathfrak m)$ を次元 $1$ の Noether 局所環とする。
$r \geq 2$ とし、$a_1, \ldots, a_r \in A$ を、すべてが単元ではない非零因子とする。
このとき次が存在する。
\begin{enumerate}
\item $B/A$ が $\mathfrak m$ のある冪で消される有限環拡大 $A \subset B$,
\item 各極大イデアル $\mathfrak m_j \subset B$ に対する非零因子
$\pi_j \in B_j = B_{\mathfrak m_j}$, および
\item $B_j$ における分解 $a_i = u_{i, j} \pi_j^{e_{i, j}}$ で、
$u_{i, j} \in B_j$ は単元、$e_{i, j} \geq 0$ であるもの。
\end{enumerate}
\end{lemma}

\begin{proof}
少なくとも一つの $a_i$ は単元でないので、$\mathfrak m$ は $A$ の随伴素イデアルでない。
さらに、主張にある任意の $A \subset B$ に対して、$\mathfrak m$ は $B$ の随伴素イデアルでなく、
$\mathfrak m_j$ は $B_j$ の随伴素イデアルでない。この点を念頭に置けば、
以下の議論を確認しやすい。

\medskip\noindent
まず $r = 2$ の場合を証明すれば十分であると主張する。
$r$ に関する帰納法で示すが、この証明は読み飛ばしてもよい。
$A \subset B$、$B_j = B_{\mathfrak m_j}$ 内の $\pi_j$、および
$i = 1, \ldots, r - 1$ に対する $B_j$ 内の分解
$a_i = u_{i, j} \pi_j^{e_{i, j}}$ が与えられ、
$u_{i, j} \in B_j$ は単元、$e_{i, j} \geq 0$ とする。
$B_j$ における $\pi_j$, $a_r$ に $r = 2$ の場合を適用すると、拡大
$B_j \subset C_j$ が得られ、$C_j$ の各極大イデアル $\mathfrak m_{j, k}$ に対して、
非零因子 $\pi_{j, k} \in C_{j, k} = (C_j)_{\mathfrak m_{j, k}}$ と分解
$$
\pi_j = v_{j, k} \pi_{j, k}^{f_{j, k}}
\quad\text{and}\quad
a_r = w_{j, k} \pi_{j, k}^{g_{j, k}}
$$
が補題のとおり得られる。補題 \ref{chow-lemma-glue-at-max} により、
$C/B$ が $\mathfrak m$ のある冪で消され、すべての $j$ に対して
$C_j \cong C_{\mathfrak m_j}$ となる一意な有限拡大 $B \subset C$ が存在する。
$C$ の極大イデアルは各局所化の極大イデアル $\mathfrak m_{j, k}$ と
$1$ 対 $1$ に対応し、
これらの局所化では
$$
a_i = u_{i, j} \pi_j^{e_{i, j}} =
u_{i, j} v_{j, k}^{e_{i, j}} \pi_{j, k}^{e_{i, j}f_{j, k}}
$$
が $i \leq r - 1$ に対して成り立つ。$a_r$ も正しく分解されるので、帰納段階が完了する。

\medskip\noindent
$r = 2$ の場合を証明する。次の量に関する帰納法を用いる：
$$
\ell = \min(\text{length}_A(A/a_1A),\ \text{length}_A(A/a_2A)).
$$
$\ell = 0$ ならば $a_1$ または $a_2$ が単元であり、$A = B$ として補題が成り立つ。
したがって $\ell > 0$ と仮定してよい。

\medskip\noindent
有限環拡大 $A \subset A'$ があり、$A'/A$ は $\mathfrak m$ のある冪で消され、
$\mathfrak m$ は $A'$ の随伴素イデアルでないとする。
$\mathfrak m_1, \ldots, \mathfrak m_r \subset A'$ を極大イデアルとし、
$A'_i = A'_{\mathfrak m_i}$ とおく。各 $A'_i$ で $a_1, a_2$ に対する問題を解ければ、
補題 \ref{chow-lemma-glue-at-max} を適用して $A$ における $a_1, a_2$ の解を作れる。
$\ell = \text{length}_A(A/xA)$ となる $x \in \{a_1, a_2\}$ を選ぶ。
補題 \ref{chow-lemma-compare-periodic-lengths} と
(\ref{chow-equation-multiplicity-coker-ker}) により、
$\text{length}_A(A/xA) = \text{length}_A(A'/xA')$ である。一方、
「代数」補題 \ref{algebra-lemma-pushdown-module} により
$$
\text{length}_A(A'/xA') =
\sum [\kappa(\mathfrak m_i) : \kappa(\mathfrak m)]
\text{length}_{A'_i}(A'_i/xA'_i)
$$
である。$x \in \mathfrak m$ なので右辺の各項は正である。
$r > 1$ の場合、または $r = 1$ かつ
$[\kappa(\mathfrak m_1) : \kappa(\mathfrak m)] > 1$ の場合には、
各 $A'_i$ 内の $a_1, a_2$ に帰納法の仮定を適用できる。
したがって、上のような各 $A'$ は $A$ と同じ剰余体をもつ局所環であると仮定してよい。

\medskip\noindent
前段落の議論を適用し、補題 \ref{chow-lemma-Noetherian-domain-dim-1-two-elements} で
$a_1, a_2 \in A$ に対して構成した環で $A$ を置き換えてよい。
$A$ は局所環なので、必要なら $a_1$, $a_2$ を入れ替えることで $a_2 \in (a_1)$ となる。
$a_2 = a_1^m c$ と書き、$m > 0$ を最大にとる。実際、補題
\ref{chow-lemma-not-infinitely-divisible} により、前段落のような任意の有限拡大
$A \subset A'$ で $A$ を置き換えた後にも $m$ が最大であると仮定できる。
$c$ が単元ならば終了する。そうでなければ、補題
\ref{chow-lemma-Noetherian-domain-dim-1-two-elements} で $a_1, c \in A$ に対して
構成した環で $A$ を置き換える。このとき (1) $c = a_1 c'$ または
(2) $a_1 = c a'_1$ である。第一の場合は $a_2 = a_1^{m + 1} c'$ となり、
$m$ の最大性に反するので起こらない。第二の場合には
$a_1 = c a'_1$ および $a_2 = c^{m + 1} (a'_1)^m$ を得る。
したがって $A$ と $c, a'_1$ に対して補題を証明すれば十分である。
$a'_1$ が単元ならば終了し、そうでなければ $cA$ は $a_1A$ より真に大きいので
$\text{length}_A(A/cA) < \ell$ である。ゆえに帰納法の仮定から結論が従う。
\end{proof}






\section{tame 記号}
\label{chow-section-tame-symbol}

\noindent
$(A, \mathfrak m)$ を次元 $1$ の Noether 局所環とする。
$Q(A)$ で $A$ の全商環を表す
（「代数」例 \ref{algebra-example-localize-at-prime} を参照）。
{\it tame 記号}とは、次の性質を満たす写像
$$
\partial_A(-, -) : Q(A)^* \times Q(A)^* \longrightarrow \kappa(\mathfrak m)^*
$$
である。
\begin{enumerate}
\item $\partial_A(f, gh) = \partial_A(f, g) \partial_A(f, h)$
\label{chow-item-bilinear}
が $f, g, h \in Q(A)^*$ に対して成り立つ。
\item $\partial_A(f, g) \partial_A(g, f) = 1$
\label{chow-item-skew}
が $f, g \in Q(A)^*$ に対して成り立つ。
\item $\partial_A(f, 1 - f) = 1$
\label{chow-item-1-x}
が、$1 - f \in Q(A)^*$ を満たす $f \in Q(A)^*$ に対して成り立つ。
\item $\partial_A(aa', b) = \partial_A(a, b)\partial_A(a', b)$
\label{chow-item-bilinear-better}
および $\partial_A(a, bb') = \partial_A(a, b)\partial_A(a, b')$ が、
非零因子 $a, a', b, b' \in A$ に対して成り立つ。
\item $\partial_A(b, b) = (-1)^m$
\label{chow-item-skew-better}
が、非零因子 $b \in A$ と $m = \text{length}_A(A/bA)$ に対して成り立つ。
\item $\partial_A(u, b) = u^m \bmod \mathfrak m$
\label{chow-item-normalization}
が、単元 $u \in A$、非零因子 $b \in A$、および
$m = \text{length}_A(A/bA)$ に対して成り立つ。
\item
\label{chow-item-1-x-better}
$\partial_A(a, b - a)\partial_A(b, b) = \partial_A(b, b - a)\partial_A(a, b)$
が、$a, b, b - a$ が非零因子となる $a, b \in A$ に対して成り立つ。
\end{enumerate}
$A$ の元を用いる方が扱いやすいため、$\partial_A$ をしばしば、
(\ref{chow-item-bilinear-better}), (\ref{chow-item-skew-better}),
(\ref{chow-item-normalization}), (\ref{chow-item-1-x-better}) を満たす
$A$ の非零因子の組上の写像とみなす。
$$
\partial_A(\frac{a}{b}, \frac{c}{d}) =
\partial_A(a, c) \partial_A(a, d)^{-1} \partial_A(b, c)^{-1} \partial_A(b, d)
$$
とおけば、他の性質も含め、(\ref{chow-item-bilinear}), (\ref{chow-item-skew}),
(\ref{chow-item-1-x}) を満たす well-defined な写像
$Q(A)^* \times Q(A)^* \to \kappa(\mathfrak m)^*$ が得られることは演習である。

\medskip\noindent
これらの性質をもつ写像が一意であるとは主張しない。
代わりに、そのような写像の構成法を与える。
非零因子 $a_1, a_2 \in A$ が与えられたとき、補題
\ref{chow-lemma-prepare-tame-symbol} のように環拡大 $A \subset B$ と局所分解を選ぶ。
そして次で定義する。
\begin{equation}
\label{chow-equation-tame-symbol}
\partial_A(a_1, a_2) = \prod\nolimits_j
\text{Norm}_{\kappa(\mathfrak m_j)/\kappa(\mathfrak m)}
((-1)^{e_{1, j}e_{2, j}}u_{1, j}^{e_{2, j}}u_{2, j}^{-e_{1, j}}
\bmod \mathfrak m_j)^{m_j}
\end{equation}
ここで $m_j = \text{length}_{B_j}(B_j/\pi_j B_j)$ であり、積は $B$ の極大イデアル
$\mathfrak m_1, \ldots, \mathfrak m_r$ 全体にわたる。

\begin{lemma}
\label{chow-lemma-well-defined-tame-symbol}
公式 (\ref{chow-equation-tame-symbol}) は well-defined な
$\kappa(\mathfrak m)^*$ の元を定める。すなわち、右辺は局所分解の選び方にも
$B$ の選び方にも依存しない。
\end{lemma}

\begin{proof}
まず分解の選び方からの独立性を示す。$B$ を Noether $1$-次元局所環、
$a_1, a_2 \in B$ を元とし、非零因子 $\pi, \theta$ によって
$$
a_1 = u_1 \pi^{e_1} = v_1 \theta^{f_1},\quad
a_2 = u_2 \pi^{e_2} = v_2 \theta^{f_2}
$$
と書けるとする。ここで $e_i, f_i \geq 0$ は整数、$u_i, v_i$ は $B$ の単元である。
すると
$$
a_1^{e_2} = u_1^{e_2}u_2^{-e_1}a_2^{e_1},\quad
a_1^{f_2} = v_1^{f_2}v_2^{-f_1}a_2^{f_1}
$$
である。一方、$m = \text{length}_B(B/\pi B)$,
$k = \text{length}_B(B/\theta B)$ とおけば、
$e_2 m = \text{length}_B(B/a_2 B) = f_2 k$ である。
上の式を用いて $a_1^{e_2m} = a_1^{f_2 k}$ を展開すると
$$
(u_1^{e_2}u_2^{-e_1})^m =  (v_1^{f_2}v_2^{-f_1})^k
$$
を得る。これで符号を除く所望の等式が示された。
符号については、$me_1e_2$ が偶数であることと $kf_1f_2$ が偶数であることが
同値だと示せばよい。これは $me_2 = kf_2$ および $me_1 = kf_1$ から従う
（上と同じ議論）。

\medskip\noindent
次に $B$ の選び方からの独立性を示す。補題 \ref{chow-lemma-prepare-tame-symbol} のような
二つの拡大 $A \subset B$, $A \subset B'$ が与えられたとする。このとき
$$
C = (B \otimes_A B')/(\mathfrak m\text{-power torsion})
$$
は第三の拡大となる。したがって $A \subset B \subset C$ と仮定してよく、
$B$ の局所環上での分解を $C$ の局所環上でも用いたとき、同じ
$\kappa(\mathfrak m)$ の元を与えることを示せばよい。
ノルムの推移性（「体」補題 \ref{fields-lemma-trace-and-norm-tower}）により、
問題は次に帰着する。$B$ が次元 $1$ の Noether 局所環で、
$\pi \in B$ が非零因子ならば、
$$
\lambda^m = \prod \text{Norm}_{\kappa_k/\kappa}(\lambda)^{m_k}
$$
である。ここで記法は次のとおりである：
(1) $\kappa$ は $B$ の剰余体、
(2) $\lambda$ は $\kappa$ の元、
(3) $\mathfrak m_k \subset C$ は $C$ の極大イデアル、
(4) $\kappa_k = \kappa(\mathfrak m_k)$ は $C_k = C_{\mathfrak m_k}$ の剰余体、
(5) $m = \text{length}_B(B/\pi B)$, および
(6) $m_k = \text{length}_{C_k}(C_k/\pi C_k)$.
表示された等式は、$\lambda \in \kappa$ に対して
$\text{Norm}_{\kappa_k/\kappa}(\lambda) = \lambda^{[\kappa_k : \kappa]}$ であり、
$m = \sum m_k[\kappa_k:\kappa]$ であることから従う。
まず補題 \ref{chow-lemma-compare-periodic-lengths} と
(\ref{chow-equation-multiplicity-coker-ker}) により
$m = \text{length}_B(B/xB) = \text{length}_B(C/\pi C)$ である。
最後に「代数」補題 \ref{algebra-lemma-pushdown-module} により
$\text{length}_B(C/\pi C) = \sum m_k[\kappa_k:\kappa]$ である。
\end{proof}

\begin{lemma}
\label{chow-lemma-tame-symbol}
tame 記号 (\ref{chow-equation-tame-symbol}) は
(\ref{chow-item-bilinear-better}), (\ref{chow-item-skew-better}),
(\ref{chow-item-normalization}), (\ref{chow-item-1-x-better}) を満たす。
したがって (\ref{chow-item-bilinear}), (\ref{chow-item-skew}), (\ref{chow-item-1-x}) を満たす写像
$\partial_A : Q(A)^* \times Q(A)^* \to \kappa(\mathfrak m)^*$ を与える。
\end{lemma}

\begin{proof}
まず (\ref{chow-item-bilinear-better}) を証明する。
$a_1, a_2, a_3 \in A$ を非零因子とし、補題
\ref{chow-lemma-prepare-tame-symbol} のような $A \subset B$ を
$a_1, a_2, a_3$ に対して選ぶ。このとき等式
$$
\partial_A(a_1a_2, a_3) = \partial_A(a_1, a_3) \partial_A(a_2, a_3)
$$
は、$B_j$ における等式
$$
(-1)^{(e_{1, j} + e_{2, j})e_{3, j}}
(u_{1, j}u_{2, j})^{e_{3, j}}u_{3, j}^{-e_{1, j} - e_{2, j}} =
(-1)^{e_{1, j}e_{3, j}}
u_{1, j}^{e_{3, j}}u_{3, j}^{-e_{1, j}}
(-1)^{e_{2, j}e_{3, j}}
u_{2, j}^{e_{3, j}}u_{3, j}^{-e_{2, j}}
$$
から従う。(\ref{chow-item-skew-better}) と (\ref{chow-item-normalization}) も同様に直ちに従う。

\medskip\noindent
次に (\ref{chow-item-1-x-better}) を証明する。
$a_1, a_2, a_1 - a_2 \in A$ を非零因子とし、$a_3 = a_1 - a_2$ とおく。
補題 \ref{chow-lemma-prepare-tame-symbol} のような $A \subset B$ を
$a_1, a_2, a_3$ に対して選ぶ。すると次を示せば十分である：
$$
(-1)^{e_{1, j}e_{2, j} + e_{1, j}e_{3, j} + e_{2, j}e_{3, j} + e_{2, j}}
u_{1, j}^{e_{2, j} - e_{3, j}}
u_{2, j}^{e_{3, j} - e_{1, j}}
u_{3, j}^{e_{1, j} - e_{2, j}} \bmod \mathfrak m_j = 1
$$
$e_{1, j} = e_{2, j} = e_{3, j}$ ならば明らかである。
$e_{1, j} > e_{2, j}$ とする。$a_3 = a_1 - a_2$ なので
$e_{3, j} = e_{2, j}$ であり、$u_{3, j}$ と $-u_{2, j}$ は同じ剰余類をもつ。
したがって公式は成り立ち、符号も整合する。この検証こそが
(\ref{chow-equation-tame-symbol}) での符号の選び方の理由である。
他の場合もまったく同様に扱える。
\end{proof}

\begin{lemma}
\label{chow-lemma-norm-down-tame-symbol}
$(A, \mathfrak m)$ を次元 $1$ の Noether 局所環とする。
$A \subset B$ を有限環拡大とし、$B/A$ は $\mathfrak m$ のある冪で消され、
$\mathfrak m$ は $B$ の随伴素イデアルでないとする。
非零因子 $a, b \in A$ に対して
$$
\partial_A(a, b) = \prod
\text{Norm}_{\kappa(\mathfrak m_j)/\kappa(\mathfrak m)}(\partial_{B_j}(a, b))
$$
が成り立つ。ここで積は $B$ の極大イデアル $\mathfrak m_j$ 全体にわたり、
$B_j = B_{\mathfrak m_j}$ である。
\end{lemma}

\begin{proof}
$a, b$ に対し、補題 \ref{chow-lemma-prepare-tame-symbol} のような
$B_j \subset C_j$ を選ぶ。補題 \ref{chow-lemma-glue-at-max} により、
すべての $j$ に対して $C_j \cong C_{\mathfrak m_j}$ となる有限環拡大
$B \subset C$ を選べる。$\mathfrak m_j$ の上にある $C$ の極大イデアルを
$\mathfrak m_{j, k} \subset C$ とする。$C_j \cong C_{\mathfrak m_j}$ の選び方によって存在する
局所分解を
$$
a = u_{j, k}\pi_{j, k}^{f_{j, k}},\quad
b = v_{j, k}\pi_{j, k}^{g_{j, k}}
$$
とする。定義により
$$
\partial_A(a, b) = 
\prod\nolimits_{j, k}
\text{Norm}_{\kappa(\mathfrak m_{j, k})/\kappa(\mathfrak m)}
((-1)^{f_{j, k}g_{j, k}}u_{j, k}^{g_{j, k}}v_{j, k}^{-f_{j, k}}
\bmod \mathfrak m_{j, k})^{m_{j, k}}
$$
および
$$
\partial_{B_j}(a, b) = 
\prod\nolimits_k
\text{Norm}_{\kappa(\mathfrak m_{j, k})/\kappa(\mathfrak m_j)}
((-1)^{f_{j, k}g_{j, k}}u_{j, k}^{g_{j, k}}v_{j, k}^{-f_{j, k}}
\bmod \mathfrak m_{j, k})^{m_{j, k}}
$$
である。結果は拡大 $\kappa(\mathfrak m_{j, k})/\kappa(\mathfrak m_j)/\kappa(\mathfrak m)$
に対するノルムの推移性から従う
（「体」補題 \ref{fields-lemma-trace-and-norm-tower}）。
\end{proof}

\begin{lemma}
\label{chow-lemma-tame-symbol-formally-smooth}
$(A, \mathfrak m, \kappa) \to (A', \mathfrak m', \kappa')$ を
Noether 局所環の局所準同型とする。$A \to A'$ は平坦で、
$\dim(A) = \dim(A') = 1$ と仮定する。
$m = \text{length}_{A'}(A'/\mathfrak mA')$ とおく。
非零因子 $a_1, a_2 \in A$ に対し、$\partial_A(a_1, a_2)^m$ は
$\kappa \to \kappa'$ によって $\partial_{A'}(a_1, a_2)$ に写る。
\end{lemma}

\begin{proof}
$a_1, a_2$ がともに単元ならば、$\partial_A(a_1, a_2) = 1$ および
$\partial_{A'}(a_1, a_2) = 1$ なので結論は成り立つ。
そうでなければ、補題 \ref{chow-lemma-prepare-tame-symbol} のような
環拡大 $A \subset B$ と局所分解を選べる。
$\mathfrak m_1, \ldots, \mathfrak m_m$ を $B$ の極大イデアルと書く。
$\mathfrak m_1, \ldots, \mathfrak m_m$ を、剰余体
$\kappa_1, \ldots, \kappa_m$ をもつ $B$ の極大イデアルとする。
各 $j \in \{1, \ldots, m\}$ に対し、補題のような分解
$a_i = u_{i, j}\pi_j^{e_{i, j}}$ を与える非零因子
$\pi_j \in B_j = B_{\mathfrak m_j}$ をとる。定義により
$$
\partial_A(a_1, a_2) = \prod\nolimits_j
\text{Norm}_{\kappa_j/\kappa}
((-1)^{e_{1, j}e_{2, j}}u_{1, j}^{e_{2, j}}u_{2, j}^{-e_{1, j}}
\bmod \mathfrak m_j)^{m_j}
$$
である。ここで $m_j = \text{length}_{B_j}(B_j/\pi_j B_j)$ である。

\medskip\noindent
$B' = A' \otimes_A B$ とおく。$A'$ は $A$ 上平坦なので、
$A' \subset B'$ は環拡大で、$B'/A'$ は $\mathfrak m'$ のある冪で消される。
$$
\mathfrak m'_{j, l},\quad l = 1, \ldots, n_j
$$
を、$\mathfrak m_j$ の上にある $B'$ の極大イデアルとする。
$\kappa'_{j, l}$ を $\mathfrak m'_{j, l}$ の剰余体、$B'_{j, l}$ を
$\mathfrak m'_{j, l}$ における $B'$ の局所化とする。
$B'_{j, l}$ における $a_1$, $a_2$ の分解には、$B_j$ で与えられた分解
$a_i = u_{i, j} \pi_j^{e_{i, j}}$ の像を用いる。定義により
$$
\partial_{A'}(a_1, a_2) = \prod\nolimits_{j, l}
\text{Norm}_{\kappa'_{j, l}/\kappa'}
((-1)^{e_{1, j}e_{2, j}}u_{1, j}^{e_{2, j}}u_{2, j}^{-e_{1, j}}
\bmod \mathfrak m'_{j, l})^{m'_{j, l}}
$$
である。ここで $m'_{j, l} = \text{length}_{B'_{j, l}}(B'_{j, l}/\pi_j B'_{j, l})$ である。

\medskip\noindent
上の公式を比較すると、各 $j$ と任意の単元 $u \in B_j$ に対して
\begin{equation}
\label{chow-equation-to-prove}
\left(\text{Norm}_{\kappa_j/\kappa}(u \bmod \mathfrak m_j)^{m_j}\right)^m
=
\prod\nolimits_l
\text{Norm}_{\kappa'_{j, l}/\kappa'}(u \bmod \mathfrak m'_{j, l})^{m'_{j, l}}
\end{equation}
を $\kappa'$ で示せば十分である。これには「代数についてさらに」第
\ref{more-algebra-section-determinants-finite-length} 節の、有限長加群の自己準同型の
行列式の構成を用いる。$M = B_j/\pi_j B_j$ とおく。
「代数についてさらに」補題 \ref{more-algebra-lemma-multiplication} により
$$
\text{Norm}_{\kappa_j/\kappa}(u \bmod \mathfrak m_j)^{m_j} =
\det\nolimits_\kappa(u : M \to M)
$$
である。したがって「代数についてさらに」補題
\ref{more-algebra-lemma-flat-base-change-det} により、
(\ref{chow-equation-to-prove}) の左辺は
$\det_{\kappa'}(u : M \otimes_A A' \to M \otimes_A A')$ に等しい。
$A'$-加群の同型
$$
M \otimes_A A' = (B_j/\pi_j B_j) \otimes_A A' =
\bigoplus\nolimits_l B'_{j, l}/\pi_j B'_{j, l}
$$
がある。$M'_l = B'_{j, l}/\pi_j B'_{j, l}$ とおけば、再び
「代数についてさらに」補題 \ref{more-algebra-lemma-multiplication} により
$\text{Norm}_{\kappa'_{j, l}/\kappa'}(u \bmod \mathfrak m'_{j, l})^{m'_{j, l}}
= \det_{\kappa'}(u_j : M'_l \to M'_l)$ である。
ゆえに行列式構成の乗法性から (\ref{chow-equation-to-prove}) が成り立つ
（「代数についてさらに」補題 \ref{more-algebra-lemma-ses}）。
\end{proof}






\section{鍵となる補題}
\label{chow-section-key-lemma}

\noindent
本節では上の結果を適用して補題 \ref{chow-lemma-milnor-gersten-low-degree} を証明する。
この補題は、Milnor K 理論に対して Quillen の K 理論における
Gersten--Quillen 複体と類似の複体が存在するという主張の低次数の場合である。
注意 \ref{chow-remark-gersten-complex-milnor} を参照されたい。

\begin{lemma}
\label{chow-lemma-perpare-key}
$(A, \mathfrak m)$ を $2$-次元 Noether 局所環とする。
$t \in \mathfrak m$ を非零因子とし、
$V(t) = \{\mathfrak m, \mathfrak q_1, \ldots, \mathfrak q_r\}$ と書く。
$A_{\mathfrak q_i} \subset B_i$ を有限環拡大とし、
$B_i/A_{\mathfrak q_i}$ は $t$ のある冪で消されるとする。
このとき剰余体を同一視する局所環の有限拡大 $A \subset B$ で、
$B_i \cong B_{\mathfrak q_i}$ かつ $B/A$ が $t$ のある冪で消されるものが存在する。
\end{lemma}

\begin{proof}
$B_i \subset t^{-n}A_{\mathfrak q_i}$ となる $n > 0$ を選ぶ。
$M \subset t^{-n}A$ と、\ $M' \subset t^{-2n}A$ を、それぞれ
$t^{-n}A_{\mathfrak q_i}$ と、\ $t^{-2n}A_{\mathfrak q_i}$ で
$B_i$ に写る元からなる $A$-部分加群とする。
$A$ は Noether 環で、局所化は完全なので、$M \subset M'$ は有限 $A$-加群であり、
$M_{\mathfrak q_i} = M'_{\mathfrak q_i} = B_i$ である。
したがって、ある $c > 0$ に対して $M'/M$ は $\mathfrak m^c$ で消される。
乗法 $t^{-n}A \times t^{-n}A \to t^{-2n}A$ のもとで
$M \cdot M \subset M'$ である。ゆえに
$B = A + \mathfrak m^{c + 1}M$ は、所望の局所化をもつ有限 $A$-代数である。
$B$ が極大イデアル $\mathfrak m + \mathfrak m^{c + 1}M$ をもつ局所環であることの
確認は省略する。
\end{proof}

\begin{lemma}
\label{chow-lemma-key-nonzerodivisors}
$(A, \mathfrak m)$ を $2$-次元 Noether 局所環とし、
$a, b \in A$ を非零因子とする。このとき
$$
\sum
\text{ord}_{A/\mathfrak q}(\partial_{A_{\mathfrak q}}(a, b))
=
0
$$
である。ここで和は $A$ の高さ $1$ の素イデアル $\mathfrak q$ 全体にわたる。
\end{lemma}

\begin{proof}
$\mathfrak q$ を $A$ の高さ $1$ の素イデアルとし、$a, b$ の像が
$A_\mathfrak q$ の単元であるならば、$\partial_{A_\mathfrak q}(a, b) = 1$ である。
したがって和は有限である。実際、
$V(ab) = \{\mathfrak m, \mathfrak q_1, \ldots, \mathfrak q_r\}$ ならば、
和は $i = 1, \ldots, r$ にわたる。各 $i$ に対し、$a, b$ について
補題 \ref{chow-lemma-prepare-tame-symbol} のような拡大
$A_{\mathfrak q_i} \subset B_i$ を選ぶ。$t = ab$ と所与の素イデアル列に
補題 \ref{chow-lemma-perpare-key} を適用すると、有限局所拡大 $A \subset B$ で、
$B/A$ が $ab$ のある冪で消され、各 $i$ に対して
$B_{\mathfrak q_i} \cong B_i$ となるものがあると仮定してよい。
$\mathfrak q_{i, j}$ を $\mathfrak q_i$ の上にある $B$ の素イデアルとすると、
補題 \ref{chow-lemma-norm-down-tame-symbol} と
「代数」補題 \ref{algebra-lemma-finite-extension-dim-1} により
$$
\text{ord}_{A/\mathfrak q_i}(\partial_{A_{\mathfrak q_i}}(a, b))
=
\sum\nolimits_j
\text{ord}_{B/\mathfrak q_{i, j}}(\partial_{B_{\mathfrak q_{i, j}}}(a, b))
$$
である。したがって $A$ を $B$ で置き換え、次段落の場合に帰着してよい。

\medskip\noindent
各 $i$ に対し、非零因子 $\pi_i \in A_{\mathfrak q_i}$ と単元
$u_i, v_i \in A_{\mathfrak q_i}$ があり、ある整数 $e_i, f_i \geq 0$ に対して
$$
a = u_i \pi_i^{e_i},\quad b = v_i \pi_i^{f_i}
$$
が $A_{\mathfrak q_i}$ で成り立つと仮定する。
$m_i = \text{length}_{A_{\mathfrak q_i}}(A_{\mathfrak q_i}/\pi_i)$ とおけば、定義により
$\partial_{A_{\mathfrak q_i}}(a, b) =
((-1)^{e_if_i}u_i^{f_i}v_i^{-e_i})^{m_i}$ である。
$a, b$ は非零因子なので、$(2, 1)$-周期複体 $(A/(ab), a, b)$ の
コホモロジーは消える。$M_i$ を $A/(ab)$ の $A_{\mathfrak q_i}/(ab)$ における像とする。
すると写像
$$
A/(ab) \longrightarrow \bigoplus M_i
$$
の核と余核は $\{\mathfrak m\}$ に台をもち、したがって有限長である。
補題 \ref{chow-lemma-compare-periodic-lengths} により
$$
\sum e_A(M_i, a, b) = 0
$$
である。ゆえに
$e_A(M_i, a, b) =
- \text{ord}_{A/\mathfrak q_i}(\partial_{A_{\mathfrak q_i}}(a, b))$
を示せば十分である。

\medskip\noindent
まず $\pi_i, u_i, v_i$ が元 $\pi_i, u_i, v_i \in A$ の像である場合を証明する
（同じ記号を用いても混乱はない）。この場合、
\begin{align*}
e_A(M_i, a, b) & =
e_A(M_i, u_i\pi_i^{e_i}, v_i\pi_i^{f_i}) \\
& =
e_A(M_i, \pi_i^{e_i}, \pi_i^{f_i}) -
e_A(\pi_i^{e_i}M_i, 0, u_i) +
e_A(\pi_i^{f_i}M_i, 0, v_i) \\
& =
0 -
f_im_i\text{ord}_{A/\mathfrak q_i}(u_i) +
e_im_i\text{ord}_{A/\mathfrak q_i}(v_i) \\
& =
-m_i\text{ord}_{A/\mathfrak q_i}(u_i^{f_i}v_i^{-e_i}) =
-\text{ord}_{A/\mathfrak q_i}(\partial_{A_{\mathfrak q_i}}(a, b))
\end{align*}
を得る。第二の等式は補題 \ref{chow-lemma-multiply-period-length} による。
さらに
$M_i \subset (M_i)_{\mathfrak q_i} = A_{\mathfrak q_i}/(\pi_i^{e_i + f_i})$,
$(\pi_i^{e_i}M_i)_{\mathfrak q_i} \cong A_{\mathfrak q_i}/\pi_i^{f_i}$,
$(\pi_i^{f_i}M_i)_{\mathfrak q_i} \cong A_{\mathfrak q_i}/\pi_i^{e_i}$ である。
第三の等式の $0$ は補題 \ref{chow-lemma-powers-period-length-zero} から、
他の二項は補題 \ref{chow-lemma-length-multiplication} から生じる。
最後の二等式は位数関数の乗法性と tame 記号の定義から従う。

\medskip\noindent
一般の場合には、まず $c\pi_i \in A$ となる
$c \in A$, $c \not \in \mathfrak q_i$ を選ぶ。
$\pi_i$ を $c\pi_i$ で、$u_i$ を $c^{-e_i}u_i$ で、$v_i$ を
$c^{-f_i}v_i$ で置き換え、$\pi_i$ は $A$ に属すると仮定してよい。
次に $cu_i, cv_i \in A$ となる $c \in A$, $c \not \in \mathfrak q_i$ を選ぶ。
補題 \ref{chow-lemma-length-multiplication} により
$$
e_A(M_i, ca, cb) = e_A(M_i, a, b) - e_A(aM_i, 0, c) + e_A(bM_i, 0, c)
$$
である。一方、$c$ は $A_{\mathfrak q_i}$ の単元なので、
$\kappa(\mathfrak q_i)^*$ において
$$
\partial_{A_{\mathfrak q_i}}(ca, cb) =
c^{m_i(f_i - e_i)}\partial_{A_{\mathfrak q_i}}(a, b)
$$
である。前段落の議論から
$e_A(M_i, ca, cb) = -
\text{ord}_{A/\mathfrak q_i}(\partial_{A_{\mathfrak q_i}}(ca, cb))$ である。
したがって
$$
e_A(aM_i, 0, c) = \text{ord}_{A/\mathfrak q_i}(c^{m_if_i})
\quad\text{and}\quad
e_A(bM_i, 0, c) = \text{ord}_{A/\mathfrak q_i}(c^{m_ie_i})
$$
を証明すれば十分である。これは、$\mathfrak q_i$ で局所化したときの上記の記述と
補題 \ref{chow-lemma-length-multiplication} から従う。
\end{proof}

\begin{lemma}[鍵となる補題]
\label{chow-lemma-milnor-gersten-low-degree}
\begin{reference}
$A$ が excellent 環の場合、これは \cite[Proposition 1]{Kato-Milnor-K} である。
\end{reference}
$A$ を分数体 $K$ をもつ $2$-次元 Noether 局所整域とし、$f, g \in K^*$ とする。
$\mathfrak q_1, \ldots, \mathfrak q_t$ を、$f$ または $g$ の少なくとも一方が
$A^*_{\mathfrak q}$ の元でないような $A$ の高さ $1$ の素イデアル
$\mathfrak q$ とする。このとき
$$
\sum\nolimits_{i = 1, \ldots, t}
\text{ord}_{A/\mathfrak q_i}(\partial_{A_{\mathfrak q_i}}(f, g))
=
0
$$
である。これは次のようにも書ける：
$$
\sum\nolimits_{\text{height}(\mathfrak q) = 1}
\text{ord}_{A/\mathfrak q}(\partial_{A_{\mathfrak q}}(f, g))
=
0
$$
実際、$f, g \in A^*_{\mathfrak q}$ となる $A$ の任意の高さ $1$ の素イデアル
$\mathfrak q$ に対して $\partial_{A_{\mathfrak q}}(f, g) = 1$ である。
\end{lemma}

\begin{proof}
tame 記号 $\partial_{A_{\mathfrak q}}(f, g)$ は双線形で、位数関数
$\text{ord}_{A/\mathfrak q}$ は加法的なので、$f$, $g$ が $A$ の元である場合に
公式を証明すれば十分である。この場合は補題
\ref{chow-lemma-key-nonzerodivisors} で証明されている。
\end{proof}

\begin{remark}[Milnor K 理論]
\label{chow-remark-gersten-complex-milnor}
体 $k$ に対し、$K^M_*(k)$ で、$k^*$ 上のテンソル代数を
$x \in k \setminus \{0, 1\}$ に対する元 $x \otimes 1 - x$ が生成する
両側イデアルで割った商を表す。したがって
$K^M_0(k) = \mathbf{Z}$, $K_1^M(k) = k^*$ であり、
$$
K^M_2(k) = k^* \otimes_\mathbf{Z} k^* / \langle x \otimes 1 - x \rangle
$$
である。$A$ が分数体 $F = \text{Frac}(A)$ と剰余体 $\kappa$ をもつ離散付値環ならば、
第 \ref{chow-section-tame-symbol} 節のように定義される tame 記号
$$
\partial_A : K_{i + 1}^M(F) \to K_i^M(\kappa)
$$
がある（\cite{Kato-Milnor-K} を参照）。より一般に、正規化と $K_i^M$ 上の
ノルム写像を用いると、この写像は $A$ が次元 $1$ の excellent 局所整域である場合に
拡張できる（\cite{Kato-Milnor-K} を参照）。おそらく第
\ref{chow-section-tame-symbol} 節の方法を用いれば、tame 記号 $\partial_A$ の構成を
任意の次元 $1$ の Noether 局所整域 $A$ に拡張できる。
次に、$X$ を次元関数 $\delta$ をもつ Noether スキームとする。
これらの tame 記号を用いて次の矢印を得る：
$$
\bigoplus\nolimits_{\delta(x) = j + 1} K^M_{i + 1}(\kappa(x))
\longrightarrow
\bigoplus\nolimits_{\delta(x) = j} K^M_i(\kappa(x))
\longrightarrow
\bigoplus\nolimits_{\delta(x) = j - 1} K^M_{i - 1}(\kappa(x))
$$
しかし合成が零であること、すなわち Abel 群の複体が得られることは自明でない。
$X$ が excellent の場合、これは \cite{Kato-Milnor-K} で示されている。
$i = 1$ で $j$ が任意の場合には、補題
\ref{chow-lemma-milnor-gersten-low-degree} から従う。
\end{remark}





















\section{設定}
\label{chow-section-setup}

\noindent
以下では一貫して、次元関数 $\delta$ を備えた、局所 Noether かつ普遍的鎖状な
基底 $S$ 上で議論する。本章の議論の一部はさらに一般化できると思われるが、
これは良い第一近似である。これはまさに \cite{Thorup} で扱われる一般性である。
通常、スキームが分離的または準コンパクトであるとは仮定しない。
興味深い代数スタックの多くは非分離的または非準コンパクトであり、
それらに対しても妥当な理論をどのように展開するかを見るための良い事例となる。
これらの仮定を参照できるよう番号を付ける。

\begin{situation}
\label{chow-situation-setup}
ここで $S$ は局所 Noether かつ普遍的鎖状なスキームである。
さらに、$S$ は次元関数 $\delta : S \longrightarrow \mathbf{Z}$ を備えると仮定する。
\end{situation}

\noindent
普遍的鎖状スキームの概念については「射」定義
\ref{morphisms-definition-universally-catenary} を、次元関数の概念については
「位相」定義 \ref{topology-definition-dimension-function} を参照されたい。
任意の局所 Noether 鎖状スキームは局所的に次元関数をもつ
（「性質」補題 \ref{properties-lemma-catenary-dimension-function}）。
また、普遍的鎖状なスキームは多数存在する
（「射」補題 \ref{morphisms-lemma-ubiquity-uc}）。

\medskip\noindent
$(S, \delta)$ を状況 \ref{chow-situation-setup} のとおりとする。
$S$ 上局所有限型な任意のスキーム $X$ は局所 Noether かつ鎖状である。
実際、$X$ は $(f : X \to S, \delta)$ に付随する標準的な次元関数
$$
\delta = \delta_{X/S} : X \longrightarrow \mathbf{Z}
$$
をもち、これは
$\delta_{X/S}(x) = \delta(f(x)) + \text{trdeg}_{\kappa(f(x))}\kappa(x)$
で与えられる。「射」補題 \ref{morphisms-lemma-dimension-function-propagates} を参照されたい。
さらに、$h : X \to Y$ が $S$ 上局所有限型なスキームの射で、
$x \in X$, $y = h(x)$ ならば、明らかに
$\delta_{X/S}(x) = \delta_{Y/S}(y) + \text{trdeg}_{\kappa(y)}\kappa(x)$ である。
以下、この関数とその性質を自由に用いる。

\medskip\noindent
上の設定の基本的な例を挙げる。実際、主に興味があるのは、基底が体のスペクトルである場合、
または基底が Dedekind 環のスペクトルである場合
（たとえば\ $\mathbf{Z}$ または離散付値環）である。

\begin{example}
\label{chow-example-field}
ここで $S = \Spec(k)$ であり、$k$ は体である。
$pt$ を $S$ の唯一の点として $\delta(pt) = 0$ とおく。
組 $(S, \delta)$ は状況 \ref{chow-situation-setup} の例である
（「射」補題 \ref{morphisms-lemma-ubiquity-uc} による）。
\end{example}

\begin{example}
\label{chow-example-domain-dimension-1}
ここで $S = \Spec(A)$ であり、$A$ は次元 $1$ の Noether 整域である。
例えば $A = \mathbf{Z}$ としてよい。$\mathfrak p$ が極大イデアルならば
$\delta(\mathfrak p) = 0$ とおき、$\mathfrak p = (0)$ が生成点に対応するときは
$\delta(\mathfrak p) = 1$ とおく。これは状況
\ref{chow-situation-setup} の例である（「射」補題
\ref{morphisms-lemma-ubiquity-uc} による）。
\end{example}

\begin{example}
\label{chow-example-CM-irreducible}
ここで $S$ は Cohen--Macaulay スキームである。
「射」補題 \ref{morphisms-lemma-ubiquity-uc} により $S$ は普遍的鎖状である。
$\delta(s) = -\dim(\mathcal{O}_{S, s})$ とおく。
$s' \leadsto s$ が $S$ の点の非自明な特殊化ならば、
$\mathcal{O}_{S, s'}$ は $\mathcal{O}_{S, s}$ の非極大素イデアル
$\mathfrak p \subset \mathcal{O}_{S, s}$ における局所化である
（「スキーム」補題 \ref{schemes-lemma-specialize-points}）。
したがって「代数」補題 \ref{algebra-lemma-CM-dim-formula} により
$\dim(\mathcal{O}_{S, s}) =
\dim(\mathcal{O}_{S, s'}) + \dim(\mathcal{O}_{S, s}/\mathfrak p) >
\dim(\mathcal{O}_{S, s'})$ である。ゆえに $\delta(s') > \delta(s)$ である。
$s' \leadsto s$ が直近の特殊化ならば、$\mathfrak p$ と $\mathfrak m_s$ の間に
真に挟まる素イデアルはなく、$\delta(s') = \delta(s) + 1$ を得る。
したがって $\delta$ は次元関数であり、組 $(S, \delta)$ は状況
\ref{chow-situation-setup} の例である。
\end{example}

\noindent
$S$ が Jacobson で、$\delta$ が閉点を零に写すならば、$\delta$ は各点を
その閉包の次元に写す関数である。

\begin{lemma}
\label{chow-lemma-delta-is-dimension}
$(S, \delta)$ を状況 \ref{chow-situation-setup} のとおりとする。
さらに $S$ は Jacobson スキームで、$S$ のすべての閉点 $s$ に対し
$\delta(s) = 0$ と仮定する。$X$ を $S$ 上局所有限型とする。
$Z \subset X$ を整閉部分スキームとし、$\xi \in Z$ をその生成点とする。
次の整数は等しい。
\begin{enumerate}
\item $\delta_{X/S}(\xi)$,
\item $\dim(Z)$, および
\item $z$ を $Z$ の閉点としたときの $\dim(\mathcal{O}_{Z, z})$.
\end{enumerate}
\end{lemma}

\begin{proof}
$X \to S$, $\xi \in Z \subset X$ を補題のとおりとする。
$X$ は $S$ 上局所有限型なので、$X$ は Jacobson である
（「射」補題 \ref{morphisms-lemma-Jacobson-universally-Jacobson}）。
したがって $X$ の閉点は $Z$ の各閉部分集合で稠密であり、$S$ の閉点に写る。
ゆえに $Z$ の既約閉部分集合の任意の鎖を $Z$ の閉点で終わらせられる。
したがって
$\dim(Z) = \sup_z(\dim(\mathcal{O}_{Z, z})$
である（「性質」補題 \ref{properties-lemma-codimension-local-ring}）。
ここで $z \in Z$ は $Z$ の閉点全体を走る。
次元関数の性質から
$\dim(\mathcal{O}_{Z, z}) = \delta(\xi) - \delta(z)$ である。
各閉点 $z \in Z$ に対して体拡大 $\kappa(z)/\kappa(f(z))$ は有限である
（「射」補題 \ref{morphisms-lemma-jacobson-finite-type-points}）。
したがって閉点 $z \in Z$ に対して
$\delta_{X/S}(z) = \delta(f(z)) = 0$ である。
ゆえに三つの整数はすべて等しい。
\end{proof}

\noindent
上の補題の状況では、閉既約部分集合の生成点における $\delta$ の値は
その既約閉部分集合の次元である。しかし一般には、この等式は期待できない。
例えば $S = \Spec(\mathbf{C}[[t]])$ かつ
$X = \Spec(\mathbf{C}((t)))$ ならば、$X$ の唯一の点について
$\delta(x) = 1$ となるが、$\dim(X) = 0$ である。
それでも $\delta_{X/S}$ は既約閉部分スキームの次元を与えるものと考えたい。
そこで次の用語を導入する。

\begin{definition}
\label{chow-definition-delta-dimension}
$(S, \delta)$ を状況 \ref{chow-situation-setup} のとおりとする。
$X$ を $S$ 上局所有限型な任意のスキーム、$Z \subset X$ を任意の既約閉部分集合とし、
$$
\dim_\delta(Z) = \delta(\xi)
$$
と定義する。ここで $\xi \in Z$ は $Z$ の生成点である。
これを {\it $Z$ の $\delta$-次元}と呼ぶ。
$Z$ が $X$ の閉部分スキームならば、$\dim_\delta(Z)$ を、その既約成分の
$\delta$-次元の上限として定義する。
\end{definition}







\section{サイクル}
\label{chow-section-cycles}

\noindent
ここではスキームが準コンパクトであるとは仮定しないので、サイクルの定義には
少し注意が必要である。例えば有理関数は無限個の極をもちうるため、無限和を
許さなければならない。いずれにせよ、$X$ が準コンパクトならば、サイクルは
通常どおり有限和である。

\begin{definition}
\label{chow-definition-cycles}
$(S, \delta)$ を状況 \ref{chow-situation-setup} のとおりとする。
$X$ を $S$ 上局所有限型とし、$k \in \mathbf{Z}$ とする。
\begin{enumerate}
\item {\it $X$ 上のサイクル}とは形式和
$$
\alpha = \sum n_Z [Z]
$$
であって、和は整閉部分スキーム $Z \subset X$ 全体にわたり、各
$n_Z \in \mathbf{Z}$ であり、族 $\{Z; n_Z \not = 0\}$ が局所有限であるものをいう
（「位相」定義 \ref{topology-definition-locally-finite}）。
\item $X$ 上の{\it $k$-サイクル}とは、
$$
\alpha = \sum n_Z [Z]
$$
と表され、$n_Z \not = 0 \Rightarrow \dim_\delta(Z) = k$ を満たすサイクルをいう。
\item $X$ 上のすべての $k$-サイクルからなる Abel 群を $Z_k(X)$ と表す。
\end{enumerate}
\end{definition}

\noindent
言い換えれば、$X$ 上の $k$-サイクルは、$\delta$-次元 $k$ の整閉部分スキームの
局所有限な形式的 $\mathbf{Z}$-線形結合である。$k$-サイクル
$\alpha = \sum n_Z[Z]$ と $\beta = \sum m_Z[Z]$ の和は
$$
\alpha + \beta = \sum (n_Z + m_Z)[Z],
$$
すなわち係数を加えることによって定める。

\begin{remark}
\label{chow-remark-cycles-pointwise}
$(S, \delta)$ を状況 \ref{chow-situation-setup} のとおりとする。
$X$ を $S$ 上局所有限型とし、$k \in \mathbf{Z}$ とする。このとき
$$
Z_k(X) = \bigoplus\nolimits_{\delta(x) = k}' K_0^M(\kappa(x))
\quad\subset\quad
\bigoplus\nolimits_{\delta(x) = k} K_0^M(\kappa(x))
$$
と書ける。ここで記法と規約は次のとおりである。
\begin{enumerate}
\item $K_0^M(\kappa(x)) = \mathbf{Z}$ は、点 $x \in X$ の剰余体
$\kappa(x)$ の Milnor K 理論の次数 $0$ の部分である
（注意 \ref{chow-remark-gersten-complex-milnor} を参照）。
\item 右辺の直和は、$\delta(x) = k$ を満たすすべての点 $x \in X$ にわたる。
\item 記号 $\bigoplus'_x$ は、局所有限な元からなる部分群を考えることを表す。
すなわち、任意の準コンパクト開集合 $U \subset X$ に対して、
$n_x \not = 0$ となる $x \in U$ の集合が有限であるような元 $\sum_x n_x$ をいう。
\end{enumerate}
\end{remark}

\begin{definition}
\label{chow-definition-support-cycle}
$(S, \delta)$ を状況 \ref{chow-situation-setup} のとおりとする。
$X$ を $S$ 上局所有限型とする。$X$ 上のサイクル
$\alpha = \sum n_Z [Z]$ の{\it 台}を
$$
\text{Supp}(\alpha) = \bigcup\nolimits_{n_Z \not = 0} Z \subset X
$$
とする。
\end{definition}

\noindent
族 $\{Z; n_Z \not = 0\}$ は局所有限なので、$\text{Supp}(\alpha)$ は
$X$ の閉部分集合である。$\alpha$ が $k$-サイクルならば、
$\text{Supp}(\alpha)$ の各既約成分 $Z$ の $\delta$-次元は $k$ である。

\begin{definition}
\label{chow-definition-effective-cycle}
$(S, \delta)$ を状況 \ref{chow-situation-setup} のとおりとする。
$X$ を $S$ 上局所有限型とする。$X$ 上のサイクル $\alpha$ が
すべての $Z$ に対して $n_Z \geq 0$ となる形
$\alpha =\sum n_Z [Z]$ に書けるとき、{\it 有効}であるという。
\end{definition}

\noindent
二つの有効サイクルの和は有効なので、有効サイクル全体の集合はモノイドであるが、
群ではない（$X = \emptyset$ の場合を除く）。




\section{閉部分スキームに付随するサイクル}
\label{chow-section-cycle-of-closed-subscheme}

\begin{lemma}
\label{chow-lemma-multiplicity-finite}
$(S, \delta)$ を状況 \ref{chow-situation-setup} のとおりとする。
$X$ を $S$ 上局所有限型とし、$Z \subset X$ を閉部分スキームとする。
\begin{enumerate}
\item $Z' \subset Z$ を既約成分とし、$\xi \in Z'$ をその生成点とする。このとき
$$
\text{length}_{\mathcal{O}_{X, \xi}} \mathcal{O}_{Z, \xi} < \infty
$$
である。
\item $\dim_\delta(Z) \leq k$ であり、$\xi \in Z$ が
$\delta(\xi) = k$ を満たすならば、$\xi$ は $Z$ のある既約成分の生成点である。
\end{enumerate}
\end{lemma}

\begin{proof}
$Z' \subset Z$, $\xi \in Z'$ を (1) のとおりとする。このとき
$\dim(\mathcal{O}_{Z, \xi}) = 0$ である（例えば「性質」補題
\ref{properties-lemma-codimension-local-ring} による）。したがって
$\mathcal{O}_{Z, \xi}$ は次元零の Noether 局所環であり、それ自身の上で有限長をもつ
（「代数」命題 \ref{algebra-proposition-dimension-zero-ring} を参照）。
ゆえに $\mathcal{O}_{X, \xi}$ 上でも有限長をもつ
（「代数」補題 \ref{algebra-lemma-length-independent}）。

\medskip\noindent
$\xi \in Z$ かつ $\delta(\xi) = k$ と仮定する。
閉包 $Z' = \overline{\{\xi\}}$ を考える。定義により、これは
$\dim_\delta(Z') = k$ を満たす既約閉部分スキームである。
$\dim_\delta(Z) = k$ なので、これは $Z$ の既約成分でなければならない。
したがって (2) が成り立つ。
\end{proof}

\begin{definition}
\label{chow-definition-cycle-associated-to-closed-subscheme}
$(S, \delta)$ を状況 \ref{chow-situation-setup} のとおりとする。
$X$ を $S$ 上局所有限型とし、$Z \subset X$ を閉部分スキームとする。
\begin{enumerate}
\item 生成点 $\xi$ をもつ任意の既約成分 $Z' \subset Z$ に対し、整数
$m_{Z', Z} = \text{length}_{\mathcal{O}_{X, \xi}} \mathcal{O}_{Z, \xi}$
（補題 \ref{chow-lemma-multiplicity-finite}）を、{\it $Z$ における $Z'$ の重複度}と呼ぶ。
\item $\dim_\delta(Z) \leq k$ と仮定する。{\it $Z$ に付随する $k$-サイクル}を
$$
[Z]_k
=
\sum m_{Z', Z}[Z']
$$
と定める。ここで和は、$\delta$-次元 $k$ の $Z$ の既約成分全体にわたる
（「因子」補題 \ref{divisors-lemma-components-locally-finite} により、これは
$k$-サイクルである）。
\end{enumerate}
\end{definition}

\noindent
$[Z]_k$ は、$Z$ の $\delta$-次元が $k$ 以下の場合に限って定義することに注意が必要である。
言い換えれば、規約により $[Z]_k$ と書いたならば、それは
$\dim_\delta(Z) \leq k$ を含意する。



\section{連接層に付随するサイクル}
\label{chow-section-cycle-of-coherent-sheaf}



\begin{lemma}
\label{chow-lemma-length-finite}
$(S, \delta)$ を状況 \ref{chow-situation-setup} のとおりとする。
$X$ を $S$ 上局所有限型とし、$\mathcal{F}$ を連接 $\mathcal{O}_X$-加群とする。
\begin{enumerate}
\item $\mathcal{F}$ の台の既約成分の族は局所有限である。
\item $Z' \subset \text{Supp}(\mathcal{F})$ を既約成分とし、
$\xi \in Z'$ をその生成点とする。このとき
$$
\text{length}_{\mathcal{O}_{X, \xi}} \mathcal{F}_\xi < \infty
$$
である。
\item $\dim_\delta(\text{Supp}(\mathcal{F})) \leq k$ であり、
$\xi \in \text{Supp}(\mathcal{F})$ が $\delta(\xi) = k$ を満たすならば、
$\xi$ は $\text{Supp}(\mathcal{F})$ のある既約成分の生成点である。
\end{enumerate}
\end{lemma}

\begin{proof}
「スキームのコホモロジー」補題 \ref{coherent-lemma-coherent-support-closed} により、
$\mathcal{F}$ の台 $Z$ は $X$ の閉部分集合である。$Z$ を $X$ の被約閉部分スキームと
みなしてよい（「スキーム」補題 \ref{schemes-lemma-reduced-closed-subscheme}）。
したがって (1) は $Z$ に「因子」補題
\ref{divisors-lemma-components-locally-finite} を適用すれば従い、(3) は
$Z$ に補題 \ref{chow-lemma-multiplicity-finite} を適用すれば従う。

\medskip\noindent
$\xi \in Z'$ を (2) のとおりとする。この場合、$X$ における任意の特殊化
$\xi' \leadsto \xi$ に対して $\mathcal{F}_{\xi'} = 0$ である。
$\mathcal{O}_{X, \xi}$ の非極大素イデアルは、$\xi$ に特殊化する $X$ の点に対応することを
思い出そう（「スキーム」補題 \ref{schemes-lemma-specialize-points}）。
したがって $\mathcal{F}_\xi$ は、台が $\{\mathfrak m_\xi\}$ である有限
$\mathcal{O}_{X, \xi}$-加群である。ゆえに「代数」補題
\ref{algebra-lemma-support-point} により有限長をもつ。
\end{proof}

\begin{definition}
\label{chow-definition-cycle-associated-to-coherent-sheaf}
$(S, \delta)$ を状況 \ref{chow-situation-setup} のとおりとする。
$X$ を $S$ 上局所有限型とし、$\mathcal{F}$ を連接 $\mathcal{O}_X$-加群とする。
\begin{enumerate}
\item 生成点 $\xi$ をもつ任意の既約成分
$Z' \subset \text{Supp}(\mathcal{F})$ に対し、整数
$m_{Z', \mathcal{F}} = \text{length}_{\mathcal{O}_{X, \xi}} \mathcal{F}_\xi$
（補題 \ref{chow-lemma-length-finite}）を、{\it $\mathcal{F}$ における $Z'$ の重複度}と呼ぶ。
\item $\dim_\delta(\text{Supp}(\mathcal{F})) \leq k$ と仮定する。
{\it $\mathcal{F}$ に付随する $k$-サイクル}を
$$
[\mathcal{F}]_k
=
\sum m_{Z', \mathcal{F}}[Z']
$$
と定める。ここで和は、$\text{Supp}(\mathcal{F})$ の $\delta$-次元 $k$ の
既約成分全体にわたる（補題 \ref{chow-lemma-length-finite} により、これは
$k$-サイクルである）。
\end{enumerate}
\end{definition}

\noindent
$[\mathcal{F}]_k$ は、$\mathcal{F}$ が連接であり、かつ
$\text{Supp}(\mathcal{F})$ の $\delta$-次元が $k$ 以下の場合に限って定義することに
注意が必要である。言い換えれば、規約により $[\mathcal{F}]_k$ と書いたならば、
$\mathcal{F}$ は $X$ 上連接であり、
$\dim_\delta(\text{Supp}(\mathcal{F})) \leq k$ であることを含意する。

\begin{lemma}
\label{chow-lemma-cycle-closed-coherent}
$(S, \delta)$ を状況 \ref{chow-situation-setup} のとおりとする。
$X$ を $S$ 上局所有限型とし、$Z \subset X$ を閉部分スキームとする。
$\dim_\delta(Z) \leq k$ ならば $[Z]_k = [{\mathcal O}_Z]_k$ である。
\end{lemma}

\begin{proof}
この場合、重複度 $m_{Z', Z}$ と $m_{Z', \mathcal{O}_Z}$ は定義により一致するからである。
\end{proof}

\begin{lemma}
\label{chow-lemma-additivity-sheaf-cycle}
$(S, \delta)$ を状況 \ref{chow-situation-setup} のとおりとする。
$X$ を $S$ 上局所有限型とする。$X$ 上の連接層の短完全列
$0 \to \mathcal{F} \to \mathcal{G} \to \mathcal{H} \to 0$
をとる。$\mathcal{F}$, $\mathcal{G}$, $\mathcal{H}$ の台の $\delta$-次元が
$\leq k$ であると仮定する。このとき
$[\mathcal{G}]_k = [\mathcal{F}]_k + [\mathcal{H}]_k$ である。
\end{lemma}

\begin{proof}
長さの加法性から直ちに従う。「代数」補題
\ref{algebra-lemma-length-additive} を参照されたい。
\end{proof}









\section{固有押し出しのための準備}
\label{chow-section-preparation-pushforward}

\begin{lemma}
\label{chow-lemma-equal-dimension}
$(S, \delta)$ を状況 \ref{chow-situation-setup} のとおりとする。
$X$, $Y$ を $S$ 上局所有限型とし、$f : X \to Y$ を射とする。
$X$, $Y$ は整で、$\dim_\delta(X) = \dim_\delta(Y)$ と仮定する。
このとき、$f(X)$ が $Y$ の真の閉部分スキームに含まれるか、または
$f$ が支配的で関数体の拡大 $R(X)/R(Y)$ が有限であるかのいずれかである。
\end{lemma}

\begin{proof}
$X$ は既約なので、閉包 $\overline{f(X)} \subset Y$ は既約である
（「位相」補題 \ref{topology-lemma-image-irreducible-space} および
\ref{topology-lemma-irreducible}）。$\overline{f(X)} \not = Y$ ならば結論を得る。
$\overline{f(X)} = Y$ ならば $f$ は支配的であり、「射」補題
\ref{morphisms-lemma-dominant-finite-number-irreducible-components} により、
$Y$ の生成点 $\eta_Y$ は $f$ の像に属する。したがってもちろん、
$X$ の生成点を $\eta_X \in X$ とすれば $f(\eta_X) = \eta_Y$ である。
$\delta(\eta_X) = \delta(\eta_Y)$ なので、
$R(Y) = \kappa(\eta_Y) \subset \kappa(\eta_X) = R(X)$ は超越次数 $0$ の拡大である。
ゆえに「射」補題 \ref{morphisms-lemma-finite-degree} により、
$R(Y) \subset R(X)$ は有限拡大である（この補題は「射」補題
\ref{morphisms-lemma-permanence-finite-type} により適用できる）。
\end{proof}

\begin{lemma}
\label{chow-lemma-quasi-compact-locally-finite}
$(S, \delta)$ を状況 \ref{chow-situation-setup} のとおりとする。
$X$, $Y$ を $S$ 上局所有限型とし、$f : X \to Y$ を射とする。
$f$ は準コンパクトであり、$\{Z_i\}_{i \in I}$ は $X$ の閉部分集合の
局所有限族であると仮定する。このとき
$\{\overline{f(Z_i)}\}_{i \in I}$ は $Y$ の閉部分集合の局所有限族である。
\end{lemma}

\begin{proof}
$V \subset Y$ を準コンパクト開部分集合とする。$f$ は準コンパクトなので、
開集合 $f^{-1}(V)$ は準コンパクトである。したがって集合
$\{i \in I \mid Z_i \cap f^{-1}(V) \not = \emptyset \}$
は、ここでは省略する簡単な位相的議論により有限である。これは集合
$$
\{i \in I \mid f(Z_i) \cap V \not = \emptyset \} =
\{i \in I \mid \overline{f(Z_i)} \cap V \not = \emptyset \}
$$
と同じなので、補題が従う。
\end{proof}









\section{固有押し出し}
\label{chow-section-proper-pushforward}

\begin{definition}
\label{chow-definition-proper-pushforward}
$(S, \delta)$ を状況 \ref{chow-situation-setup} のとおりとする。
$X$, $Y$ を $S$ 上局所有限型とし、$f : X \to Y$ を射とする。
$f$ は固有であると仮定する。
\begin{enumerate}
\item $Z \subset X$ を $\dim_\delta(Z) = k$ を満たす整閉部分スキームとする。
次のように定義する。
$$
f_*[Z] =
\left\{
\begin{matrix}
0 & \text{if} & \dim_\delta(f(Z))< k, \\
\deg(Z/f(Z)) [f(Z)] & \text{if} & \dim_\delta(f(Z)) = k.
\end{matrix}
\right.
$$
ここで $f(Z) \subset Y$ は整閉部分スキームとみなす。
$\dim_\delta(f(Z)) = \dim_\delta(Z)$ ならば、補題
\ref{chow-lemma-equal-dimension} により、$f(Z)$ 上の $Z$ の次数は有限である。
\item $\alpha = \sum n_Z [Z]$ を $X$ 上の $k$-サイクルとする。
$\alpha$ の{\it 押し出し}を和
$$
f_* \alpha = \sum n_Z f_*[Z]
$$
として定義する。ここで各 $f_*[Z]$ は上のとおり定義する。
補題 \ref{chow-lemma-quasi-compact-locally-finite} により、この和は局所有限である。
\end{enumerate}
\end{definition}

\noindent
定義により、サイクルの固有押し出し
$$
f_* : Z_k(X) \longrightarrow Z_k(Y)
$$
は Abel 群の準同型である。これは $X \mapsto Z_k(X)$ を、$S$ 上局所有限型な
スキームを対象、固有射を射とする圏上の共変関手にする。

\begin{lemma}
\label{chow-lemma-compose-pushforward}
$(S, \delta)$ を状況 \ref{chow-situation-setup} のとおりとする。
$X$, $Y$, $Z$ を $S$ 上局所有限型とする。
$f : X \to Y$ と $g : Y \to Z$ を固有射とする。このとき写像
$Z_k(X) \to Z_k(Z)$ として $g_* \circ f_* = (g \circ f)_*$ である。
\end{lemma}

\begin{proof}
$W \subset X$ を次元 $k$ の整閉部分スキームとする。
$W' = f(W) \subset Y$ と $W'' = g(f(W)) \subset Z$ を考える。
$f$, $g$ は固有なので、$W'$（それぞれ\ $W''$）は $Y$（それぞれ\ $Z$）の
整閉部分スキームである。$g_*(f_*[W]) = (g \circ f)_*[W]$ を示せばよい。
$\dim_\delta(W'') < k$ ならば両辺は零である。
$\dim_\delta(W'') = k$ ならば、誘導される射
$$
W \longrightarrow
W' \longrightarrow
W''
$$
はいずれも補題 \ref{chow-lemma-equal-dimension} の仮定を満たす。したがって
$$
g_*(f_*[W]) = \deg(W/W')\deg(W'/W'')[W''],
\quad
(g \circ f)_*[W] = \deg(W/W'')[W''].
$$
そこで「射」補題 \ref{morphisms-lemma-degree-composition} を適用すれば結論を得る。
\end{proof}

\noindent
閉埋め込みは固有である。$i : Z \to X$ が閉埋め込みならば、すべての写像
$$
i_* : Z_k(Z) \longrightarrow Z_k(X)
$$
は{\it 単射}である。

\begin{lemma}
\label{chow-lemma-exact-sequence-closed}
$(S, \delta)$ を状況 \ref{chow-situation-setup} のとおりとする。
$X$ を $S$ 上局所有限型とする。$X_1, X_2 \subset X$ を、集合論的に
$X = X_1 \cup X_2$ となる閉部分スキームとする。任意の
$k \in \mathbf{Z}$ に対し、Abel 群の列
$$
\xymatrix{
Z_k(X_1 \cap X_2) \ar[r] &
Z_k(X_1) \oplus Z_k(X_2) \ar[r] &
Z_k(X) \ar[r] &
0
}
$$
は完全である。ここで $X_1 \cap X_2$ はスキーム論的交叉であり、写像は
押し出し写像で、その一方には $-1$ を掛ける。
\end{lemma}

\begin{proof}
まず $X$ が準コンパクトであると仮定する。このとき $Z_k(X)$ は、
$Z \subset X$ が $\delta$-次元 $k$ の整閉部分スキームであるような元 $[Z]$ を
基底とする自由 $\mathbf{Z}$-加群である。群 $Z_k(X_1)$, $Z_k(X_2)$,
$Z_k(X_1 \cap X_2)$ は、それぞれ $Z \subset X_1$, $Z \subset X_2$,
$Z \subset X_1 \cap X_2$ を満たすこれらの $Z$ の部分集合上で自由である。
これにより、この場合の補題は直ちに従う。一般の場合も同様であり、証明を省略する。
\end{proof}

\begin{lemma}
\label{chow-lemma-cycle-push-sheaf}
$(S, \delta)$ を状況 \ref{chow-situation-setup} のとおりとする。
$f : X \to Y$ を、$S$ 上局所有限型なスキームの固有射とする。
\begin{enumerate}
\item $Z \subset X$ を $\dim_\delta(Z) \leq k$ を満たす閉部分スキームとする。このとき
$$
f_*[Z]_k = [f_*{\mathcal O}_Z]_k.
$$
\item $\mathcal{F}$ を、$\dim_\delta(\text{Supp}(\mathcal{F})) \leq k$ を満たす
$X$ 上の連接層とする。このとき
$$
f_*[\mathcal{F}]_k = [f_*{\mathcal F}]_k.
$$
\end{enumerate}
「スキームのコホモロジー」命題
\ref{coherent-proposition-proper-pushforward-coherent} により、
$f_*\mathcal{F}$ と $f_*\mathcal{O}_Z$ は連接 $\mathcal{O}_Y$-加群なので、
この主張は意味をもつことに注意する。
\end{lemma}

\begin{proof}
(1) は (2) と補題 \ref{chow-lemma-cycle-closed-coherent} から従う。
$\mathcal{F}$ を $X$ 上の連接層とし、
$\dim_\delta(\text{Supp}(\mathcal{F})) \leq k$ と仮定する。
「スキームのコホモロジー」補題 \ref{coherent-lemma-coherent-support-closed} により、
閉部分スキーム $i : Z \to X$ と連接 $\mathcal{O}_Z$-加群 $\mathcal{G}$ であって、
$i_*\mathcal{G} \cong \mathcal{F}$ かつ $\mathcal{F}$ の台が $Z$ となるものが存在する。
$Z' \subset Y$ を $f|_Z : Z \to Y$ のスキーム論的像とする。スキームの可換図式
$$
\xymatrix{
Z \ar[r]_i \ar[d]_{f|_Z} &
X \ar[d]^f \\
Z' \ar[r]^{i'} & Y
}
$$
を考える。図式を二通りに回ることにより
$f_*\mathcal{F} = f_*i_*\mathcal{G} = i'_*(f|_Z)_*\mathcal{G}$ である。
閉埋め込みと $f|_Z$ に対して結論が成り立つと仮定すると、
$$
f_*[\mathcal{F}]_k = f_*i_*[\mathcal{G}]_k
= (i')_*(f|_Z)_*[\mathcal{G}]_k =
(i')_*[(f|_Z)_*\mathcal{G}]_k =
[(i')_*(f|_Z)_*\mathcal{G}]_k = [f_*\mathcal{F}]_k
$$
を得て、望む結論となる。閉埋め込みの場合は直ちに分かるので省略する。
$f|_Z : Z \to Z'$ は支配的な射であることに注意する
（「射」補題 \ref{morphisms-lemma-quasi-compact-scheme-theoretic-image}）。
したがって $\dim_\delta(X) \leq k$ で、$f : X \to Y$ が固有かつ支配的である場合に帰着した。

\medskip\noindent
$\dim_\delta(X) \leq k$ で、$f : X \to Y$ は固有かつ支配的であると仮定する。
$f$ は支配的なので、生成点 $\eta$ をもつ任意の既約成分 $Z \subset Y$ に対し、
$f(\xi) = \eta$ となる点 $\xi \in X$ が存在する。したがって
$\delta(\eta) \leq \delta(\xi) \leq k$ である。ゆえに表示
$$
f_*[\mathcal{F}]_k = \sum n_Z[Z],
\quad
\text{and}
\quad
[f_*\mathcal{F}]_k = \sum m_Z[Z].
$$
において、$n_Z \not = 0$ または $m_Z \not = 0$ ならば、整閉部分スキーム
$Z$ は実際に $Y$ の $\delta$-次元 $k$ の既約成分である。このような整閉部分スキーム
$Z \subset Y$ を一つとり、その生成点を $\eta$ と書く。$f(\xi) = \eta$ を満たす
任意の $\xi \in X$ に対して $\delta(\xi) \geq k$ であり、したがって $\xi$ も
$X$ の $\delta$-次元 $k$ のある既約成分の生成点である
（補題 \ref{chow-lemma-multiplicity-finite}）。$f$ は準コンパクトで $X$ は局所 Noether なので、
このような点は有限個しかなく、ゆえに $f^{-1}(\{\eta\})$ は有限である。
「射」補題 \ref{morphisms-lemma-generically-finite} により、
$f^{-1}(V) \to V$ が有限となる開近傍 $\eta \in V \subset Y$ が存在する。
$Y$ を $V$ で、$X$ を $f^{-1}(V)$ で置き換えると、$Y$ がアフィンで $f$ が有限な場合に帰着する。

\medskip\noindent
$Y = \Spec(R)$, $X = \Spec(A)$ と書く（有限射はアフィンなので可能である）。
このとき $R$ と $A$ は Noether 環で、$A$ は $R$ 上有限である。また、ある有限
$A$-加群 $M$ に対して $\mathcal{F} = \widetilde{M}$ である。
$f_*\mathcal{F}$ は、$R$-加群とみなした $M$ に対応することに注意する。
$\mathfrak p \subset R$ を、$\eta \in Y$ に対応する極小素イデアルとする。
$[f_*\mathcal{F}]_k$ における $Z$ の係数は明らかに
$\text{length}_{R_{\mathfrak p}}(M_{\mathfrak p})$ である。
$\mathfrak q_i$, $i = 1, \ldots, t$ を、$\mathfrak p$ の上にある $A$ の素イデアルとする。
$A_{\mathfrak p}$ は次元零の Noether 局所環 $R_{\mathfrak p}$ 上有限な Artin 環なので、
$A_{\mathfrak p} = \prod A_{\mathfrak q_i}$ である。
$f_*[\mathcal{F}]_k$ における $Z$ の係数は明らかに
$$
\sum\nolimits_{i = 1, \ldots, t}
[\kappa(\mathfrak q_i) : \kappa(\mathfrak p)]
\text{length}_{A_{\mathfrak q_i}}(M_{\mathfrak q_i})
$$
である。したがって「代数」補題 \ref{algebra-lemma-pushdown-module} から望む等式が従う。
\end{proof}




















\section{平坦引き戻しのための準備}
\label{chow-section-preparation-flat-pullback}


\noindent
局所有限型な射 $f : X \to Y$ は、空でない各ファイバーが次元 $r$ の等次元であるとき、
相対次元 $r$ をもつということを思い出そう。「射」定義
\ref{morphisms-definition-relative-dimension-d} を参照されたい。

\begin{lemma}
\label{chow-lemma-flat-inverse-image-dimension}
$(S, \delta)$ を状況 \ref{chow-situation-setup} のとおりとする。
$X$, $Y$ を $S$ 上局所有限型とし、$f : X \to Y$ を射とする。
$f$ は相対次元 $r$ の平坦射であると仮定する。任意の閉部分集合
$Z \subset Y$ に対し、$f^{-1}(Z)$ が空でなければ
$$
\dim_\delta(f^{-1}(Z)) = \dim_\delta(Z) + r.
$$
である。$Z$ が既約で、$Z' \subset f^{-1}(Z)$ が既約成分ならば、
$Z'$ は $Z$ を支配し、$\dim_\delta(Z') = \dim_\delta(Z) + r$ である。
\end{lemma}

\begin{proof}
最後の主張を示せば十分である。$Y$ を整閉部分スキーム $Z$ で、$X$ を
スキーム論的逆像 $f^{-1}(Z) = Z \times_Y X$ で置き換えてよい。
したがって $Z = Y$ は整で、$f$ は相対次元 $r$ の平坦射であると仮定してよい。
$Y$ は局所 Noether なので、局所有限型な射 $f$ は実際に局所有限表示である。
したがって「射」補題 \ref{morphisms-lemma-fppf-open} が適用でき、$f$ は開である。
$\xi \in X$ を $X$ のある既約成分の生成点とする。$f$ は開なので、
$f(\xi)$ は $Z = Y$ の生成点 $\eta$ である。$f$ が相対次元 $r$ をもつという仮定から
$\dim_\xi(X_\eta) = r$ である。一方、$\xi$ は $X$ の生成点なので、
$\mathcal{O}_{X, \xi} = \mathcal{O}_{X_\eta, \xi}$ はただ一つの素イデアルをもち、
したがって次元 $0$ である。ゆえに「射」補題
\ref{morphisms-lemma-dimension-fibre-at-a-point} により、
$\kappa(\eta)$ 上の $\kappa(\xi)$ の超越次数は $r$ である。
言い換えれば $\delta(\xi) = \delta(\eta) + r$ となり、望む結論を得る。
\end{proof}

\noindent
閉部分スキームの局所有限族の平坦引き戻しが局所有限であることを示すため、
次の補題を用いる。

\begin{lemma}
\label{chow-lemma-inverse-image-locally-finite}
$(S, \delta)$ を状況 \ref{chow-situation-setup} のとおりとする。
$X$, $Y$ を $S$ 上局所有限型とし、$f : X \to Y$ を射とする。
$\{Z_i\}_{i \in I}$ は $Y$ の閉部分集合の局所有限族であると仮定する。
このとき $\{f^{-1}(Z_i)\}_{i \in I}$ は $X$ の閉部分集合の局所有限族である。
\end{lemma}

\begin{proof}
$U \subset X$ を準コンパクト開部分集合とする。像 $f(U) \subset Y$ は準コンパクトなので、
$f(U) \subset V$ を満たす準コンパクト開集合 $V \subset Y$ が存在する。ここで
$$
\{i \in I \mid f^{-1}(Z_i) \cap U \not = \emptyset \}
\subset
\{i \in I \mid Z_i \cap V \not = \emptyset \}.
$$
右辺は仮定により有限なので、結論を得る。
\end{proof}



\section{平坦引き戻し}
\label{chow-section-flat-pullback}

\noindent
以下では、スキームの射 $f : X \to Y$ と閉部分スキーム $Z \subset Y$ に対し、
$f^{-1}(Z)$ でその{\it スキーム論的逆像}を表す。スキーム論的逆像はファイバー積
$$
\xymatrix{
f^{-1}(Z) \ar[r] \ar[d] & X \ar[d] \\
Z \ar[r] & Y
}
$$
であることを思い出そう。また、$\mathcal{I} \subset \mathcal{O}_Y$ が $Y$ における
$Z$ に対応する準連接イデアル層ならば、これは準連接イデアル層
$f^{-1}(\mathcal{I})\mathcal{O}_X$ によって切り出される $X$ の閉部分スキームでもある。
（これは「スキーム」節 \ref{schemes-section-closed-immersion}、補題
\ref{schemes-lemma-fibre-product-immersion}、および定義
\ref{schemes-definition-inverse-image-closed-subscheme} で論じられている。）

\begin{definition}
\label{chow-definition-flat-pullback}
$(S, \delta)$ を状況 \ref{chow-situation-setup} のとおりとする。
$X$, $Y$ を $S$ 上局所有限型とし、$f : X \to Y$ を射とする。
$f$ は相対次元 $r$ の平坦射であると仮定する。
\begin{enumerate}
\item $Z \subset Y$ を $\delta$-次元 $k$ の整閉部分スキームとする。
スキーム論的逆像に付随する $X$ 上の $(k+r)$-サイクルとして $f^*[Z]$ を
$$
f^*[Z] = [f^{-1}(Z)]_{k+r}.
$$
と定義する。補題 \ref{chow-lemma-flat-inverse-image-dimension} により
$\dim_\delta(f^{-1}(Z)) = k + r$ なので、これは意味をもつ。
\item $\alpha = \sum n_i [Z_i]$ を $Y$ 上の $k$-サイクルとする。
{\it $f$ による $\alpha$ の平坦引き戻し}を和
$$
f^* \alpha = \sum n_i f^*[Z_i]
$$
として定義する。ここで各 $f^*[Z_i]$ は上のとおり定義する。
補題 \ref{chow-lemma-inverse-image-locally-finite} により、この和は局所有限である。
\item このように得られる Abel 群の写像を $f^* : Z_k(Y) \to Z_{k + r}(X)$ と表す。
\end{enumerate}
\end{definition}

\noindent
開埋め込みは平坦である。これは自明ではあるが、平坦射の重要な特別の場合である。
$U \subset X$ が開ならば、サイクルの $j : U \to X$ による引き戻しを、しばしば
$U$ へのサイクルの{\it 制限}と呼ぶ。この場合、すべての写像
$$
j^* : Z_k(X) \longrightarrow Z_k(U)
$$
は{\it 全射}であることに注意する。その理由は、任意の整閉部分スキーム
$Z' \subset U$ に対し、$Z'$ の $X$ における閉包を $Z$ として、それを $X$ の
被約閉部分スキームとみなせるからである（「スキーム」補題
\ref{schemes-lemma-reduced-closed-subscheme}）。明らかに $Z \cap U = Z'$、すなわち
$j^*[Z] = [Z']$ なので全射性が従う。実際、もう少し強いことが成り立つ。

\begin{lemma}
\label{chow-lemma-exact-sequence-open}
$(S, \delta)$ を状況 \ref{chow-situation-setup} のとおりとする。
$X$ を $S$ 上局所有限型とする。$U \subset X$ を開部分スキームとし、
$i : Y = X \setminus U \to X$ を $X$ の被約閉部分スキームとして表す。
任意の $k \in \mathbf{Z}$ に対し、列
$$
\xymatrix{
Z_k(Y) \ar[r]^{i_*} & Z_k(X) \ar[r]^{j^*} & Z_k(U) \ar[r] & 0
}
$$
は Abel 群の完全複体である。
\end{lemma}

\begin{proof}
まず $X$ が準コンパクトであると仮定する。このとき $Z_k(X)$ は、
$Z \subset X$ が $\delta$-次元 $k$ の整閉部分スキームであるような元 $[Z]$ を
基底とする自由 $\mathbf{Z}$-加群である。このような基底元は、基底元
$[Z \cap U]$ に写るか、$Z \subset Y$ ならば零に写る。したがってこの場合の補題は
明らかである。一般の場合も同様であり、証明を省略する。
\end{proof}

\begin{lemma}
\label{chow-lemma-compose-flat-pullback}
$(S, \delta)$ を状況 \ref{chow-situation-setup} のとおりとする。
$X, Y, Z$ を $S$ 上局所有限型とする。$f : X \to Y$ と $g : Y \to Z$ を、
それぞれ相対次元 $r$ と $s$ の平坦射とする。このとき $g \circ f$ は
相対次元 $r + s$ の平坦射であり、
$$
f^* \circ g^* = (g \circ f)^*
$$
が写像 $Z_k(Z) \to Z_{k + r + s}(X)$ として成り立つ。
\end{lemma}

\begin{proof}
「射」補題 \ref{morphisms-lemma-composition-relative-dimension-d} により、合成は
相対次元 $r + s$ の平坦射である。次を仮定する。
\begin{enumerate}
\item $W \subset Z$ は $\delta$-次元 $k$ の整閉部分スキームである。
\item $W' \subset Y$ は $\delta$-次元 $k + s$ の整閉部分スキームで、
$W' \subset g^{-1}(W)$ を満たす。
\item $W'' \subset Y$ は $\delta$-次元 $k + s + r$ の整閉部分スキームで、
$W'' \subset f^{-1}(W')$ を満たす。
\end{enumerate}
$(g \circ f)^*[W]$ における $[W'']$ の係数 $n$ と、
$f^*(g^*[W])$ における $[W'']$ の係数 $m$ が一致することを示せばよい。
これらの場合を確認すれば補題に十分であることは、補題
\ref{chow-lemma-flat-inverse-image-dimension} から従う。
$\xi'' \in W''$, $\xi' \in W'$, $\xi \in W$ をそれぞれ生成点とする。
局所環 $A = \mathcal{O}_{Z, \xi}$, $B = \mathcal{O}_{Y, \xi'}$,
$C = \mathcal{O}_{X, \xi''}$ を考える。このとき局所平坦環準同型
$A \to B$, $B \to C$ があり、さらに
$$
n = \text{length}_C(C/\mathfrak m_AC),
\quad
\text{and}
\quad
m = \text{length}_C(C/\mathfrak m_BC) \text{length}_B(B/\mathfrak m_AB)
$$
である。したがって「代数」補題 \ref{algebra-lemma-pullback-transitive} から等式が従う。
\end{proof}

\begin{lemma}
\label{chow-lemma-pullback-coherent}
$(S, \delta)$ を状況 \ref{chow-situation-setup} のとおりとする。
$X, Y$ を $S$ 上局所有限型とし、$f : X \to Y$ を相対次元 $r$ の平坦射とする。
\begin{enumerate}
\item $Z \subset Y$ を $\dim_\delta(Z) \leq k$ を満たす閉部分スキームとする。このとき
$\dim_\delta(f^{-1}(Z)) \leq k + r$ であり、$Z_{k + r}(X)$ において
$[f^{-1}(Z)]_{k + r} = f^*[Z]_k$ である。
\item $\mathcal{F}$ を、$\dim_\delta(\text{Supp}(\mathcal{F})) \leq k$ を満たす
$Y$ 上の連接層とする。このとき
$\dim_\delta(\text{Supp}(f^*\mathcal{F})) \leq k + r$ であり、
$$
f^*[{\mathcal F}]_k = [f^*{\mathcal F}]_{k+r}
$$
が $Z_{k + r}(X)$ において成り立つ。
\end{enumerate}
\end{lemma}

\begin{proof}
次元に関する主張は補題 \ref{chow-lemma-flat-inverse-image-dimension} から直ちに従う。
(1) は (2)、補題 \ref{chow-lemma-cycle-closed-coherent}、および
$f^*\mathcal{O}_Z = \mathcal{O}_{f^{-1}(Z)}$ であることから従う。

\medskip\noindent
(2) を証明する。$X$, $Y$ は局所 Noether なので、「スキームのコホモロジー」補題
\ref{coherent-lemma-coherent-Noetherian} を適用すると、$\mathcal{F}$ は有限型であり、
したがって $f^*\mathcal{F}$ は有限型であり（「加群」補題
\ref{modules-lemma-pullback-finite-type}）、ゆえに $f^*\mathcal{F}$ は連接である
（再び「スキームのコホモロジー」補題 \ref{coherent-lemma-coherent-Noetherian}）。
したがって補題は意味をもつ。$W \subset Y$ を $\delta$-次元 $k$ の整閉部分スキームとし、
$W' \subset X$ を $f$ により $W$ に写る次元 $k + r$ の整閉部分スキームとする。
$f^*[{\mathcal F}]_k$ における $[W']$ の係数 $n$ と、
$[f^*{\mathcal F}]_{k+r}$ における $[W']$ の係数 $m$ が一致することを示せばよい。
$\xi \in W$ と $\xi' \in W'$ を生成点とする。
$A = \mathcal{O}_{Y, \xi}$, $B = \mathcal{O}_{X, \xi'}$ とおき、
$M = \mathcal{F}_\xi$ を $A$-加群とみなす。（次元の仮定により $M$ は有限長をもつが、
実際にはこれを確認する必要はない。補題 \ref{chow-lemma-length-finite} を参照。）
$f^*\mathcal{F}_{\xi'} = B \otimes_A M$ である。したがって
$$
n = \text{length}_B(B \otimes_A M)
\quad
\text{and}
\quad
m = \text{length}_A(M) \text{length}_B(B/\mathfrak m_AB)
$$
であり、「代数」補題 \ref{algebra-lemma-pullback-module} から等式が従う。
\end{proof}



\section{押し出しと引き戻し}
\label{chow-section-push-pull}

\noindent
この節では、意味をもつ場合に固有押し出しと平坦引き戻しが両立することを確かめる。
上で行った準備により、これはコホモロジーと基底変換の帰結である。

\begin{lemma}
\label{chow-lemma-flat-pullback-proper-pushforward}
$(S, \delta)$ を状況 \ref{chow-situation-setup} のとおりとする。
$$
\xymatrix{
X' \ar[r]_{g'} \ar[d]_{f'} & X \ar[d]^f \\
Y' \ar[r]^g & Y
}
$$
を、$S$ 上局所有限型なスキームのファイバー積図式とする。
$f : X \to Y$ は固有で、$g : Y' \to Y$ は相対次元 $r$ の平坦射であると仮定する。
このとき $f'$ も固有であり、$g'$ も相対次元 $r$ の平坦射である。
$X$ 上の任意の $k$-サイクル $\alpha$ に対し、$Z_{k + r}(Y')$ において
$$
g^*f_*\alpha = f'_*(g')^*\alpha
$$
である。
\end{lemma}

\begin{proof}
$f'$ が固有であることは「射」補題 \ref{morphisms-lemma-base-change-proper} から従う。
$g'$ が相対次元 $r$ の平坦射であることは、「射」補題
\ref{morphisms-lemma-base-change-relative-dimension-d} および
\ref{morphisms-lemma-base-change-flat} から従う。
$W \subset X$ が $\delta$-次元 $k$ の整閉部分スキームで、
$\alpha = [W]$ である場合にサイクルの等式を示せば十分である。
この場合、補題 \ref{chow-lemma-cycle-closed-coherent} により
$\alpha = [\mathcal{O}_W]_k$ である。したがって補題
\ref{chow-lemma-cycle-push-sheaf} と \ref{chow-lemma-pullback-coherent} により、
$f'_*(g')^*\mathcal{O}_W$ が $g^*f_*\mathcal{O}_W$ と同型であることを示せば十分である。
これはコホモロジーと基底変換から従う。「スキームのコホモロジー」補題
\ref{coherent-lemma-flat-base-change-cohomology} を参照されたい。
\end{proof}

\begin{lemma}
\label{chow-lemma-finite-flat}
$(S, \delta)$ を状況 \ref{chow-situation-setup} のとおりとする。
$X$, $Y$ を $S$ 上局所有限型とする。$f : X \to Y$ を次数 $d$ の有限局所自由射とする
（「射」定義 \ref{morphisms-definition-finite-locally-free} を参照）。
このとき $f$ は固有かつ相対次元 $0$ の平坦射であり、任意の
$\alpha \in Z_k(Y)$ に対して
$$
f_*f^*\alpha = d\alpha
$$
である。
\end{lemma}

\begin{proof}
「射」補題 \ref{morphisms-lemma-finite-flat} により有限局所自由射は平坦かつ有限であり、
「射」補題 \ref{morphisms-lemma-finite-proper} により有限射は固有である。
有限射が相対次元 $0$ をもつことの証明は省略する。したがって公式は意味をもつ。
$Z \subset Y$ を $\delta$-次元 $k$ の整閉部分スキームとする。
$\alpha = [Z]$ に対して公式を示せば十分である。有限局所自由射の基底変換は有限局所自由なので
（「射」補題 \ref{morphisms-lemma-base-change-finite-locally-free}）、
$f_*f^*\mathcal{O}_Z$ は $Z$ 上階数 $d$ の有限局所自由層である。したがって
$$
f_*f^*[Z] = f_*f^*[\mathcal{O}_Z]_k =
[f_*f^*\mathcal{O}_Z]_k = d[Z]
$$
である。ここで補題 \ref{chow-lemma-pullback-coherent} と
\ref{chow-lemma-cycle-push-sheaf} を用いた。
\end{proof}








\section{主因子のための準備}
\label{chow-section-preparation-principal-divisors}

\noindent
この節の内容の一部は、「因子」節 \ref{divisors-section-Weil-divisors} の議論と重なる。

\begin{lemma}
\label{chow-lemma-divisor-delta-dimension}
$(S, \delta)$ を状況 \ref{chow-situation-setup} のとおりとする。
$X$ を $S$ 上局所有限型とし、$X$ は整であると仮定する。
\begin{enumerate}
\item $Z \subset X$ が整閉部分スキームならば、次は同値である。
\begin{enumerate}
\item $Z$ は素因子である。
\item $Z$ は $X$ において余次元 $1$ をもつ。
\item $\dim_\delta(Z) = \dim_\delta(X) - 1$ である。
\end{enumerate}
\item $Z$ が $X$ 上の有効 Cartier 因子の既約成分ならば、
$\dim_\delta(Z) = \dim_\delta(X) - 1$ である。
\end{enumerate}
\end{lemma}

\begin{proof}
(1) は素因子の定義（「因子」定義 \ref{divisors-definition-Weil-divisor}）と
次元関数の定義（「位相」定義 \ref{topology-definition-dimension-function}）から従う。
$\xi \in Z$ を、有効 Cartier 因子 $D \subset X$ の既約成分 $Z$ の生成点とする。
このとき $\dim(\mathcal{O}_{D, \xi}) = 0$ であり、ある非零因子
$f \in \mathcal{O}_{X, \xi}$ に対し
$\mathcal{O}_{D, \xi} = \mathcal{O}_{X, \xi}/(f)$ である
（「因子」補題 \ref{divisors-lemma-effective-Cartier-in-points}）。
「代数」補題 \ref{algebra-lemma-one-equation} により
$\dim(\mathcal{O}_{X, \xi}) = 1$ である。したがって「性質」補題
\ref{properties-lemma-codimension-local-ring} により $Z$ は (1) のとおりであり、証明が完了する。
\end{proof}

\begin{lemma}
\label{chow-lemma-finite-in-codimension-one}
$f : X \to Y$ をスキームの射とし、$\xi \in Y$ を点とする。次を仮定する。
\begin{enumerate}
\item $X$, $Y$ は整である。
\item $Y$ は局所 Noether である。
\item $f$ は固有かつ支配的で、$R(Y) \subset R(X)$ は有限である。
\item $\dim(\mathcal{O}_{Y, \xi}) = 1$ である。
\end{enumerate}
このとき、$f|_{f^{-1}(V)} : f^{-1}(V) \to V$ が有限となるような
$\xi$ の開近傍 $V \subset Y$ が存在する。
\end{lemma}

\begin{proof}
この補題は「多様体」補題 \ref{varieties-lemma-finite-in-codim-1} の特別な場合である。
ここではこの場合の直接的な議論を与える。「スキームのコホモロジー」補題
\ref{coherent-lemma-proper-finite-fibre-finite-in-neighbourhood} により、
$f^{-1}(\{\xi\})$ が有限であることを示せば十分である。
$Y$ をアフィン開集合で置き換え、$Y = \Spec(R)$ とする。
$Y$ は局所 Noether と仮定したので、$R$ は Noether である。
$f$ は固有なので準コンパクトである。したがって、各
$U_i = \Spec(A_i)$ がアフィンとなる有限アフィン開被覆
$X = U_1 \cup \ldots \cup U_n$ をとれる。$R \to A_i$ は整域の有限型単射準同型で、
誘導される分数体の拡大は有限である。ゆえに「代数」補題
\ref{algebra-lemma-finite-in-codim-1} から結論が従う。
\end{proof}


\section{主因子}
\label{chow-section-principal-divisors}

\noindent
次の定義は、現在の設定における「因子」定義
\ref{divisors-definition-principal-divisor} の類似である。

\begin{definition}
\label{chow-definition-principal-divisor}
$(S, \delta)$ を状況 \ref{chow-situation-setup} のとおりとする。
$X$ を $S$ 上局所有限型とし、$X$ は整で
$\dim_\delta(X) = n$ と仮定する。$f \in R(X)^*$ とする。
{\it $f$ に付随する主因子}とは、「因子」定義
\ref{divisors-definition-principal-divisor} で定義される $(n - 1)$-サイクル
$$
\text{div}(f) = \text{div}_X(f) = \sum \text{ord}_Z(f) [Z]
$$
である。補題 \ref{chow-lemma-divisor-delta-dimension} により素因子の $\delta$-次元は
$n - 1$ なので、これは意味をもつ。
\end{definition}

\noindent
定義の状況で $f, g \in R(X)^*$ に対し、$Z_{n - 1}(X)$ において
$$
\text{div}_X(fg) = \text{div}_X(f) + \text{div}_X(g)
$$
である。「因子」補題 \ref{divisors-lemma-div-additive} を参照されたい。
次の補題は、より一般的な補題 \ref{chow-lemma-flat-pullback-rational-equivalence} によって後に置き換えられる。

\begin{lemma}
\label{chow-lemma-flat-pullback-principal-divisor}
$(S, \delta)$ を状況 \ref{chow-situation-setup} のとおりとする。
$X$, $Y$ を $S$ 上局所有限型とする。$X$, $Y$ は整で、
$n = \dim_\delta(Y)$ と仮定する。$f : X \to Y$ を相対次元 $r$ の平坦射とし、
$g \in R(Y)^*$ とする。このとき $Z_{n + r - 1}(X)$ において
$$
f^*(\text{div}_Y(g)) = \text{div}_X(g)
$$
である。
\end{lemma}

\begin{proof}
$f$ は平坦なので支配的であり、したがって $f$ は埋め込み $R(Y) \subset R(X)$ を誘導する。
ゆえに $g$ を $R(X)^*$ の元とみなしてよい。
$Z \subset X$ を $\delta$-次元 $n + r - 1$ の整閉部分スキームとし、
$\xi \in Z$ をその生成点とする。$\dim_\delta(f(Z)) > n - 1$ ならば、
等式の左辺と右辺における $[Z]$ の係数はいずれも零である。
したがって $Z' = \overline{f(Z)}$ が $Y$ の $\delta$-次元 $n - 1$ の整閉部分スキームであると
仮定してよい。$\xi' = f(\xi)$ とおく。これは $Z'$ の生成点である。
$A = \mathcal{O}_{Y, \xi'}$, $B = \mathcal{O}_{X, \xi}$ とおく。
環準同型 $A \to B$ は、次元 $1$ の Noether 局所整域の平坦局所準同型である。
$g$ は $A$ の分数体に属する。示すべきことは
$$
\text{ord}_A(g) \text{length}_B(B/\mathfrak m_AB)
=
\text{ord}_B(g).
$$
である。これは「代数」補題 \ref{algebra-lemma-pullback-module} から従う
（詳細は省略する）。
\end{proof}











\section{主因子と押し出し}
\label{chow-section-two-fun}

\noindent
最初の補題から、生成的有限射に沿う主因子の押し出しは主因子であることが従う。

\begin{lemma}
\label{chow-lemma-proper-pushforward-alteration}
$(S, \delta)$ を状況 \ref{chow-situation-setup} のとおりとする。
$X$, $Y$ を $S$ 上局所有限型とする。$X$, $Y$ は整で、
$n = \dim_\delta(X) = \dim_\delta(Y)$ と仮定する。
$p : X \to Y$ を支配的固有射とし、$f \in R(X)^*$ とする。
$$
g = \text{Nm}_{R(X)/R(Y)}(f).
$$
とおく。このとき $p_*\text{div}(f) = \text{div}(g)$ である。
\end{lemma}

\begin{proof}
$Z \subset Y$ を $\delta$-次元 $n - 1$ の整閉部分スキームとする。
$p_*\text{div}(f)$ と $\text{div}(g)$ における $[Z]$ の係数が等しいことを示したい。
射 $p : X \to Y$ と生成点 $\xi \in Z$ に補題
\ref{chow-lemma-finite-in-codimension-one} を適用できる。したがって $Y$ を
$\xi$ のアフィン開近傍で置き換え、$p : X \to Y$ が有限であると仮定してよい。
$Y = \Spec(R)$, $X = \Spec(A)$ と書く。ここで $p$ は Noether 整域の有限準同型
$R \to A$ によって誘導され、これは分数体の有限体拡大 $L/K$ を誘導する。
いま $f \in L$, $g = \text{Nm}(f) \in K$ であり、
$\dim(R_{\mathfrak p}) = 1$ を満たす素イデアル $\mathfrak p \subset R$ がある。
$\text{div}_Y(g)$ における $[Z]$ の係数は $\text{ord}_{R_\mathfrak p}(g)$ である。
$p_*\text{div}_X(f)$ における $[Z]$ の係数は
$$
\sum\nolimits_{\mathfrak q\text{ lying over }\mathfrak p}
[\kappa(\mathfrak q) : \kappa(\mathfrak p)]
\text{ord}_{A_{\mathfrak q}}(f)
$$
である。したがって「代数」補題 \ref{algebra-lemma-finite-extension-dim-1} から
望む等式が従う。
\end{proof}

\noindent
主因子の議論では、$\mathbf{P}^1_S$ 上の「普遍的」主因子
$[0] - [\infty]$ が重要な役割を果たす。これを正確に述べるため、
\begin{equation}
\label{chow-equation-zero-infty}
D_0, D_\infty \subset
\mathbf{P}^1_S = \underline{\text{Proj}}_S(\mathcal{O}_S[T_0, T_1])
\end{equation}
で $\mathcal{O}(1)$ の切断 $T_1$ と\ $T_0$ によってそれぞれ切り出される閉部分スキームを表す。
これらは有効 Cartier 因子である。「因子」定義
\ref{divisors-definition-effective-Cartier-divisor} および補題
\ref{divisors-lemma-characterize-OD} を参照されたい。
次の補題は、大まかに言えば
``$\text{div}(T_1/T_0) = [D_0] - [D_1]$'' であり、これが普遍的主因子であることを述べる。

\begin{lemma}
\label{chow-lemma-rational-function}
$(S, \delta)$ を状況 \ref{chow-situation-setup} のとおりとする。
$X$ を $S$ 上局所有限型とする。$X$ は整で、$n = \dim_\delta(X)$ と仮定する。
$f \in R(X)^*$ とする。$U \subset X$ を空でない開集合で、$f$ が切断
$f \in \Gamma(U, \mathcal{O}_X^*)$ に対応するものとする。
$Y \subset X \times_S \mathbf{P}^1_S$ を $f : U \to \mathbf{P}^1_S$ のグラフの閉包とする。
このとき次が成り立つ。
\begin{enumerate}
\item 射影 $p : Y \to X$ は固有である。
\item $p|_{p^{-1}(U)} : p^{-1}(U) \to U$ は同型である。
\item 射 $q : Y \to \mathbf{P}^1_S$ による引き戻し
$Y_0 = q^{-1}D_0$ と $Y_\infty = q^{-1}D_\infty$ が定義される
（「因子」定義 \ref{divisors-definition-pullback-effective-Cartier-divisor}）。
\item
$$
\text{div}_Y(f) = [Y_0]_{n - 1} - [Y_\infty]_{n - 1}
$$
である。
\item
$$
\text{div}_X(f) = p_*\text{div}_Y(f)
$$
である。
\item 射 $p$ により $Y_0$ と $Y_\infty$ を $X$ の閉部分スキームとみなすと、
$$
\text{div}_X(f) = [Y_0]_{n - 1} - [Y_\infty]_{n - 1}
$$
である。
\end{enumerate}
\end{lemma}

\begin{proof}
$X$ は整なので $U$ も整である。したがって $Y$ は整で、
$(1, f)(U) \subset Y$ は開稠密部分スキームである。また、閉部分スキーム
$Y \subset X \times_S \mathbf{P}^1_S$ は開集合 $U$ の選び方に依存しないことに注意する。
実際、これは生成点を $\eta \in X$ としたときの一点集合
$\{\eta'\} = \{(1, f)(\eta)\}$ の閉包である。以上を踏まえて各主張を証明する。

\medskip\noindent
(1) について、$p$ は閉埋め込み $Y \to X \times_S \mathbf{P}^1_S = \mathbf{P}^1_X$ と
固有射 $\mathbf{P}^1_X \to X$ の合成である。固有射の合成は固有なので
（「射」補題 \ref{morphisms-lemma-composition-proper}）、結論を得る。

\medskip\noindent
$Y \cap U \times_S \mathbf{P}^1_S = (1, f)(U)$ は明らかである。
したがって (2) が従い、さらに $\dim_\delta(Y) = n$ も従う。

\medskip\noindent
$f \in R(X)^*$ なので、$q(\eta') = f(\eta)$ は $D_0$ にも $D_\infty$ にも含まれない。
したがって「因子」補題 \ref{divisors-lemma-pullback-effective-Cartier-defined} により (3) が従う。
補題 \ref{chow-lemma-divisor-delta-dimension} から
$\dim_\delta(Y_0) = n - 1$ および $\dim_\delta(Y_\infty) = n - 1$ を得る。

\medskip\noindent
有効 Cartier 因子 $Y_0$ を考える。各点 $\xi \in Y_0$ において
$f \in \mathcal{O}_{Y, \xi}$ であり、$Y_0$ の局所方程式は $f$ で与えられる。
特に $\delta(\xi) = n - 1$ で、$\xi$ が $\delta$-次元 $n - 1$ の整閉部分スキーム
$Z$ の生成点ならば、$\text{div}_Y(f)$ における $[Z]$ の係数は
$$
\text{ord}_Z(f) =
\text{length}_{\mathcal{O}_{Y, \xi}}
(\mathcal{O}_{Y, \xi}/f\mathcal{O}_{Y, \xi}) =
\text{length}_{\mathcal{O}_{Y, \xi}}
(\mathcal{O}_{Y_0, \xi})
$$
であり、これは $[Y_0]_{n - 1}$ における $[Z]$ の係数である。
有理関数 $1/f$ を用いる同様の議論から、$-[Y_\infty]$ は
$\text{div}_Y(f)$ の表示における負の係数をもつ項と一致する。したがって (4) が従う。

\medskip\noindent
$D_0 \to S$ は同型であることに注意する。したがって
$X \times_S D_0 \to X$ も同型である。$X \times_S \mathbf{P}^1_S$ の中で
$Y_0 = Y \cap X \times_S D_0$（スキーム論的交叉）であることは明らかである。
したがって $Y_0 \to X$ は実際に閉埋め込みであり、
$$
p_*\mathcal{O}_{Y_0} = \mathcal{O}_{Y'_0}
$$
を得る。ここで $Y'_0 \subset X$ は $Y_0 \to X$ の像である。
補題 \ref{chow-lemma-cycle-push-sheaf} により
$p_*[Y_0]_{n - 1} = [Y'_0]_{n - 1}$ である。
$D_\infty$ と $Y_\infty$ についても同じことが成り立つ。したがって (6) は (5) の帰結である。
最後に、(5) は補題 \ref{chow-lemma-proper-pushforward-alteration} から直ちに従う。
\end{proof}

\noindent
次の補題は、固有曲線上の主因子の次数が零であることを述べる。

\begin{lemma}
\label{chow-lemma-curve-principal-divisor}
$K$ を任意の体とする。$X$ を、固有射 $c : X \to \Spec(K)$ を備えた
$1$-次元整スキームとする。$f \in K(X)^*$ を可逆有理関数とする。このとき
$$
\sum\nolimits_{x \in X \text{ closed}}
[\kappa(x) : K] \text{ord}_{\mathcal{O}_{X, x}}(f)
=
0
$$
である。ここで $\text{ord}$ は「代数」定義 \ref{algebra-definition-ord} のとおりである。
言い換えれば $c_*\text{div}(f) = 0$ である。
\end{lemma}

\begin{proof}
$S = \Spec(K)$ 上の $X$ と有理関数 $f$ から補題
\ref{chow-lemma-rational-function} で構成した図式
$$
\xymatrix{
Y \ar[r]_p \ar[d]_q & X \ar[d]^c \\
\mathbf{P}^1_K \ar[r]^-{c'} & \Spec(K)
}
$$
を考える。この補題のすべての結論を以下では断りなく用いる。
$c_*\text{div}_X(f) = c_*p_*\text{div}_Y(f) = 0$ を示せばよい。
これは $c'_*q_*\text{div}_Y(f) = 0$ を示すことと同じである。
$q(Y)$ が $\mathbf{P}^1_K$ の閉点ならば $\text{div}_X(f) = 0$ であり、補題が成り立つ。
したがって $q$ は支配的であると仮定してよい。
$q : Y \to \mathbf{P}^1_K$ が次数 $d$ の有限局所自由射であることを示せると仮定する
（「射」定義 \ref{morphisms-definition-finite-locally-free} を参照）。
$\text{div}_Y(f) = [q^{-1}D_0]_0 - [q^{-1}D_\infty]_0$ なので、
平坦引き戻しの定義により
$\text{div}_Y(f) = q^*([D_0]_0 - [D_\infty]_0)$ である。
すると補題 \ref{chow-lemma-finite-flat} により
$q_*\text{div}_Y(f) = d([D_0]_0 - [D_\infty]_0)$ を得る。
$c'_*[D_0]_0 = c'_*[D_\infty]_0$ は明らかなので、結論が従う。

\medskip\noindent
$q$ が有限局所自由であることを示すのが残っている。
（$\mathbf{P}^1_K$ は連結なので、自動的にある一定の次数をもつ。）
$\dim(\mathbf{P}^1_K) = 1$ なので、例えば補題
\ref{chow-lemma-finite-in-codimension-one} により $q$ は有限である。
$\mathbf{P}^1_K$ の閉点におけるすべての局所環は、$K[T]$ の局所化なので
次元 $1$ の正則局所環、言い換えれば離散付値環である
（「代数」補題 \ref{algebra-lemma-dim-affine-space} を参照）。
したがって、閉点 $y\in Y$ に対し、局所環 $\mathcal{O}_{Y, y}$ が
捩れなしであれば、$\mathcal{O}_{\mathbf{P}^1_K, q(y)}$ 上平坦である
（「代数詳論」補題 \ref{more-algebra-lemma-dedekind-torsion-free-flat}）。
$\mathcal{O}_{Y, y}$ は整域で $q$ は支配的なので、これは明らかに成り立つ。
したがって $q$ は平坦である。ゆえに「射」補題
\ref{morphisms-lemma-finite-flat} により $q$ は有限局所自由である。
\end{proof}





\section{有理同値}
\label{chow-section-rational-equivalence}

\noindent
この節では $k$-サイクル上の{\it 有理同値}を定義する。
主因子の像（閉埋め込みによるもの）の局所有限和を許すことにする。
これはかなり奇妙な現象を引き起こす。例 \ref{chow-example-weird} を参照されたい。
しかしこれを許さなければ、直線束の Chern 類とのキャップ積が有理同値を経由することを
どのように証明すればよいか分からない。

\begin{definition}
\label{chow-definition-rational-equivalence}
$(S, \delta)$ を状況 \ref{chow-situation-setup} のとおりとする。
$X$ を $S$ 上局所有限型なスキームとし、$k \in \mathbf{Z}$ とする。
\begin{enumerate}
\item $\dim_\delta(W_j) = k + 1$ を満たす整閉部分スキームの任意の局所有限族
$\{W_j \subset X\}$ と任意の $f_j \in R(W_j)^*$ に対し、
$$
\sum (i_j)_*\text{div}(f_j) \in Z_k(X)
$$
を考えられる。ここで $i_j : W_j \to X$ は包含射である。
射 $\coprod i_j : \coprod W_j \to X$ は固有なので、これは意味をもつ。
\item $\alpha \in Z_k(X)$ が上に表示した形のサイクルであるとき、
$\alpha$ は{\it 零と有理同値}であるという。
\item $\alpha - \beta$ が零と有理同値であるとき、
$\alpha, \beta \in Z_k(X)$ は{\it 有理同値}であるといい、
$\alpha \sim_{rat} \beta$ と書く。
\item
$$
\CH_k(X) = Z_k(X) / \sim_{rat}
$$
を{\it $X$ 上の $k$-サイクルの Chow 群}と定義する。これは
{\it $X$ 上の有理同値を法とする $k$-サイクルの Chow 群}とも呼ばれる。
\end{enumerate}
\end{definition}

\noindent
ほかにも興味深い（adequate）同値関係が多数ある。有理同値はそれらの中で最も粗い。

\begin{remark}
\label{chow-remark-chow-group-pointwise}
$(S, \delta)$ を状況 \ref{chow-situation-setup} のとおりとする。
$X$ を $S$ 上局所有限型とし、$k \in \mathbf{Z}$ とする。表示
$$
\bigoplus\nolimits_{\delta(x) = k + 1}' K_1^M(\kappa(x))
\xrightarrow{\partial}
\bigoplus\nolimits_{\delta(x) = k}' K_0^M(\kappa(x)) \to
\CH_k(X) \to 0
$$
をもつことを示そう。ここでは注意 \ref{chow-remark-cycles-pointwise} で導入した記法と規約を用い、
さらに次のように定める。
\begin{enumerate}
\item $K_1^M(\kappa(x)) = \kappa(x)^*$ は、点 $x \in X$ の剰余体
$\kappa(x)$ の Milnor K 理論の次数 $1$ の部分である
（注意 \ref{chow-remark-gersten-complex-milnor} を参照）。
\item 微分 $\partial$ を次のように定義する。元 $\xi = \sum_x f_x$ に対し、
$W_x = \overline{x}$ で、生成点 $x$ をもつ $X$ の整閉部分スキームを表し、
$$
\partial(\xi) = \sum (W_x \to X)_*\text{div}(f_x)
$$
とおく。注意 \ref{chow-remark-cycles-pointwise} により、この複体の第二項は
$Z_k(X)$ と等しいことを既に見たので、これは $Z_k(X)$ に属して意味をもつ。
\end{enumerate}
これが $\CH_k(X)$ の表示を与えることは、定義
\ref{chow-definition-rational-equivalence} と比較すれば直ちに従う。
\end{remark}

\noindent
次の補題は非常に簡単だが重要である。

\begin{lemma}
\label{chow-lemma-restrict-to-open}
$(S, \delta)$ を状況 \ref{chow-situation-setup} のとおりとする。
$X$ を $S$ 上局所有限型なスキームとする。$U \subset X$ を開部分スキームとし、
$i : Y = X \setminus U \to X$ で $X$ の被約閉部分スキームを表す。
$k \in \mathbf{Z}$ とし、$\alpha, \beta \in Z_k(X)$ と仮定する。
$\alpha|_U \sim_{rat} \beta|_U$ ならば、あるサイクル $\gamma \in Z_k(Y)$ が存在して
$$
\alpha \sim_{rat} \beta + i_*\gamma.
$$
となる。言い換えれば、列
$$
\xymatrix{
\CH_k(Y) \ar[r]^{i_*} & \CH_k(X) \ar[r]^{j^*} & \CH_k(U) \ar[r] & 0
}
$$
は Abel 群の完全複体である。
\end{lemma}

\begin{proof}
$\{W_j\}_{j \in J}$ を、$U$ の $\delta$-次元 $k + 1$ の整閉部分スキームの
局所有限族とし、$f_j \in R(W_j)^*$ を、定義のとおり
$(\alpha - \beta)|_U = \sum (i_j)_*\text{div}(f_j)$ となる元とする。
$W_j' \subset X$ を $W_j$ の閉包とする。$V \subset X$ を準コンパクト開集合とする。
$V$ は Noether なので $V \cap U$ も $U$ の準コンパクト開集合である。したがって
$\{W_j\}$ の局所有限性により、集合
$\{j \in J \mid W_j \cap V \not = \emptyset\}
= \{j \in J \mid W'_j \cap V \not = \emptyset\}$
は有限である。言い換えれば $\{W'_j\}$ も局所有限である。
$R(W_j) = R(W'_j)$ なので
$$
\alpha - \beta - \sum (i'_j)_*\text{div}(f_j)
$$
は $Y$ 上に台をもつサイクルであり、補題が従う
（補題 \ref{chow-lemma-exact-sequence-open} を参照）。
\end{proof}

\begin{lemma}
\label{chow-lemma-exact-sequence-closed-chow}
$(S, \delta)$ を状況 \ref{chow-situation-setup} のとおりとする。
$X$ を $S$ 上局所有限型とする。$X_1, X_2 \subset X$ を、集合論的に
$X = X_1 \cup X_2$ となる閉部分スキームとする。
任意の $k \in \mathbf{Z}$ に対し、Abel 群の列
$$
\xymatrix{
\CH_k(X_1 \cap X_2) \ar[r] &
\CH_k(X_1) \oplus \CH_k(X_2) \ar[r] &
\CH_k(X) \ar[r] &
0
}
$$
は完全である。ここで $X_1 \cap X_2$ はスキーム論的交叉であり、写像は押し出し写像で、
その一方には $-1$ を掛ける。
\end{lemma}

\begin{proof}
補題 \ref{chow-lemma-exact-sequence-closed} により、矢印
$\CH_k(X_1) \oplus \CH_k(X_2) \to \CH_k(X)$ は全射である。
$(\alpha_1, \alpha_2)$ がこの写像で零に写ると仮定する。
$\alpha_1 = \sum n_{1, i}[W_{1, i}]$ および
$\alpha_2 = \sum n_{2, i}[W_{2, i}]$ と書く。このとき、$X$ の $\delta$-次元
$k + 1$ の整閉部分スキームの局所有限族 $\{W_j\}_{j \in J}$ と
$f_j \in R(W_j)^*$ であって、$X$ 上のサイクルとして
$$
\sum n_{1, i}[W_{1, i}] + \sum n_{2, i}[W_{2, i}] = \sum (i_j)_*\text{div}(f_j)
$$
となるものを得る。ここで $i_j : W_j \to X$ は包含射である。
$j \in J_1$ ならば $W_j \subset X_1$、$j \in J_2$ ならば
$W_j \subset X_2$ となるような非交和分解 $J = J_1 \amalg J_2$ を選ぶ。
（$W_j$ は整なので可能である。）すると上の等式を
$$
\sum n_{1, i}[W_{1, i}] - \sum\nolimits_{j \in J_1} (i_j)_*\text{div}(f_j) =
- \sum n_{2, i}[W_{2, i}] + \sum\nolimits_{j \in J_2} (i_j)_*\text{div}(f_j)
$$
と書ける。したがってこの式は $X_1 \cap X_2$ 上のサイクル（！）である。
言い換えれば、元 $(\alpha_1, \alpha_2)$ は最初の矢印の像に属し、証明が完了する。
\end{proof}

\begin{example}
\label{chow-example-weird}
「奇妙な」例を挙げる。$S$ を体 $k$ のスペクトルとし、$\delta$ は
例 \ref{chow-example-field} のとおりとする。$X = C_1 \cup C_2 \cup \ldots$ を、
曲線 $C_j \cong \mathbf{P}^1_k$ の無限和を次のように貼り合わせたものとする。
各 $j = 1, 2, 3, \ldots$ に対し、点 $\infty \in C_j$ を点
$0 \in C_{j + 1}$ と横断的に貼り合わせる。点 $0 \in C_1$ をとると、これは零次サイクル
$[0] \in Z_0(X)$ を与える。この状況の「奇妙さ」は、実際に
$[0] \sim_{rat} 0$ となることである！ すなわち、有理関数
$f_j \in R(C_j)$ を、$0$ に単純零点、$\infty$ に単純極をもち、ほかには零点も極も
もたない関数として選べる。このとき和 $\sum (i_j)_*\text{div}(f_j)$ はちょうど
$0$-サイクル $[0]$ である。実際、この例では $\CH_0(X) = 0$ となる。
これがあまりに奇妙だと感じるなら、空間を常に準コンパクトにしておけばよい
（その場合、この $X$ はそもそも存在しない）。
\end{example}

\begin{remark}
\label{chow-remark-infinite-sums-rational-equivalences}
$(S, \delta)$ を状況 \ref{chow-situation-setup} のとおりとする。
$X$ を $S$ 上局所有限型なスキームとする。$X$ 上の $k$-サイクルの無限族
$\alpha_i, \beta_i \in Z_k(X)$, $i \in I$ があると仮定する。
$\alpha_i$ と $\beta_i$ の台が $X$ の閉部分集合の局所有限族をなし、
$\sum \alpha_i$ と $\sum \beta_i$ がサイクルとして定義されるとする。
さらに各 $i$ に対して $\alpha_i \sim_{rat} \beta_i$ と仮定する。
このとき $\sum \alpha_i \sim_{rat} \sum \beta_i$ であるとは限らない。
問題は、有理同値が局所有限族
$\{W_{i, j}, f_{i, j} \in R(W_{i, j})^*\}_{j \in J_i}$ によって与えられていても、
その合併 $\{W_{i, j}\}_{i \in I, j\in J_i}$ が局所有限とは限らないことである。

\medskip\noindent
実際の多くの場合には、閉部分集合の局所有限族 $\{T_i\}_{i \in I}$ があり、
$\alpha_i, \beta_i$ は $T_i$ 上に台をもち、$\CH_k(T_i)$ において
$\alpha_i = \beta_i$ となる。言い換えれば、族
$\{W_{i, j}, f_{i, j} \in R(W_{i, j})^*\}_{j \in J_i}$ は部分スキーム
$W_{i, j} \subset T_i$ からなる。この場合には、族
$\{W_{i, j}\}_{i \in I, j\in J_i}$ が自動的に局所有限なので、$X$ 上で
$\sum \alpha_i \sim_{rat} \sum \beta_i$ が実際に成り立つ。
\end{remark}






\section{有理同値と押し出し・引き戻し}
\label{chow-section-properties-rational-equivalence}

\noindent
この節では、平坦引き戻しと固有押し出しが有理同値と可換であることを示す。

\begin{lemma}
\label{chow-lemma-prepare-flat-pullback-rational-equivalence}
$(S, \delta)$ を状況 \ref{chow-situation-setup} のとおりとする。
$X$, $Y$ を $S$ 上局所有限型なスキームとする。
$Y$ は整で $\dim_\delta(Y) = k$ と仮定する。
$f : X \to Y$ を相対次元 $r$ の平坦射とする。このとき $g \in R(Y)^*$ に対し、
$$
f^*\text{div}_Y(g) =
\sum n_j i_{j, *}\text{div}_{X_j}(g \circ f|_{X_j})
$$
が $X$ 上の $(k + r - 1)$-サイクルとして成り立つ。ここで和は $X$ の既約成分
$X_j$ 全体にわたり、$n_j$ は $X$ における $X_j$ の重複度である。
\end{lemma}

\begin{proof}
$Z \subset X$ を $\delta$-次元 $k + r - 1$ の整閉部分スキームとする。
$f^*\text{div}(g)$ における $[Z]$ の係数 $n$ と、
$\sum i_{j, *} \text{div}(g \circ f|_{X_j})$ における $[Z]$ の係数 $m$ が
等しいことを示せばよい。$Z'$ を $f(Z)$ の閉包とすると、これは $Y$ の整閉部分スキームである。
補題 \ref{chow-lemma-flat-inverse-image-dimension} により
$\dim_\delta(Z') \geq k - 1$ である。したがって $Z' = Y$ であるか、
$Z'$ は $Y$ 上の素因子である。$Z' = Y$ ならば係数 $n$ と $m$ はともに零である。
$n$ については $f^*$ の定義から明らかであり、$m$ については
$g \circ f|_{X_j}$ が $Y$ の生成点に写る $X_j$ の任意の点で単元だからである。
したがって $Z' \subset Y$ は素因子であると仮定してよい。

\medskip\noindent
$n$ と $m$ の等式を代数に翻訳する。$\xi' \in Z'$ と $\xi \in Z$ を生成点とし、
$A = \mathcal{O}_{Y, \xi'}$, $B = \mathcal{O}_{X, \xi}$ とおく。
$A$, $B$ は Noether で、$A \to B$ は平坦局所準同型、$A$ は整域であり、
$\mathfrak m_AB$ は局所環 $B$ の定義イデアルである。有理関数 $g$ は
$A$ の分数体 $Q(A)$ の元である。構成により、$\xi$ と交わる閉部分スキーム $X_j$ は
極小素イデアルと $1$ 対 $1$ に対応し、それらを
$$
\mathfrak q_1, \ldots, \mathfrak q_s \subset B
$$
と書く。整数 $n_j$ は対応する長さ
$$
n_i = \text{length}_{B_{\mathfrak q_i}}(B_{\mathfrak q_i})
$$
である。有理関数 $g \circ f|_{X_j}$ は、$g \in Q(A)$ の像
$g_i \in \kappa(\mathfrak q_i)^*$ に対応する。以上をまとめると
$$
n = \text{ord}_A(g) \text{length}_B(B/\mathfrak m_AB)
$$
および
$$
m = \sum \text{ord}_{B/\mathfrak q_i}(g_i)
\text{length}_{B_{\mathfrak q_i}}(B_{\mathfrak q_i})
$$
を得る。非零な $x, y \in A$ によって $g = x/y$ と書くと、
$$
\text{length}_A(A/(x)) \text{length}_B(B/\mathfrak m_AB) =
\text{length}_B(B/xB)
$$
（この等式では「代数」補題 \ref{algebra-lemma-pullback-module} を用いる）が
$$
\sum\nolimits_{i = 1, \ldots, s}
\text{length}_{B/\mathfrak q_i}(B/(x, \mathfrak q_i))
\text{length}_{B_{\mathfrak q_i}}(B_{\mathfrak q_i})
$$
に等しいことを示せば十分であり、$y$ についても同様である。
$A \to B$ は平坦なので $x$ は $B$ の非零因子である。
したがって補題 \ref{chow-lemma-additivity-divisors-restricted} から望む等式が従う。
\end{proof}

\begin{lemma}
\label{chow-lemma-flat-pullback-rational-equivalence}
$(S, \delta)$ を状況 \ref{chow-situation-setup} のとおりとする。
$X$, $Y$ を $S$ 上局所有限型なスキームとし、$f : X \to Y$ を相対次元 $r$ の平坦射とする。
$\alpha \sim_{rat} \beta$ を $Y$ 上の有理同値な $k$-サイクルとする。
このとき $X$ 上の $(k + r)$-サイクルとして
$f^*\alpha \sim_{rat} f^*\beta$ である。
\end{lemma}

\begin{proof}
何を示すべきかを述べる。各 $W_j$ が $\delta$-次元 $k + 1$ の整スキームである閉埋め込みの族
$$
i_j : W_j \longrightarrow Y
$$
と有理関数 $g_j \in R(W_j)^*$ が与えられ、さらに族
$\{i_j(W_j)\}_{j \in J}$ は $Y$ 上局所有限であると仮定する。このとき
$$
f^*(\sum i_{j, *}\text{div}(g_j)) = \sum f^*i_{j, *}\text{div}(g_j)
$$
が $X$ 上で零と有理同値であることを示せばよい。補題
\ref{chow-lemma-inverse-image-locally-finite} により $\{W_j\}$ は $X$ で局所有限なので、
右辺の和は意味をもつ。

\medskip\noindent
ファイバー積
$$
i'_j : W'_j = W_j \times_Y X \longrightarrow X.
$$
を考え、$f_j : W'_j \to W_j$ で第一射影を表す。
補題 \ref{chow-lemma-flat-pullback-proper-pushforward} により、上の和は
$$
\sum i'_{j, *}(f_j^*\text{div}(g_j))
$$
と書ける。補題 \ref{chow-lemma-prepare-flat-pullback-rational-equivalence} により、
各 $f_j^*\text{div}(g_j)$ は $W'_j$ 上で零と有理同値である。
したがって各 $i'_{j, *}(f_j^*\text{div}(g_j))$ は零と有理同値である。
すると注意 \ref{chow-remark-infinite-sums-rational-equivalences} の議論により、
表示された和についても同じことが成り立つ。
\end{proof}

\begin{lemma}
\label{chow-lemma-proper-pushforward-rational-equivalence}
$(S, \delta)$ を状況 \ref{chow-situation-setup} のとおりとする。
$X$, $Y$ を $S$ 上局所有限型なスキームとし、$p : X \to Y$ を固有射とする。
$\alpha, \beta \in Z_k(X)$ が有理同値であると仮定する。
このとき $p_*\alpha$ は $p_*\beta$ と有理同値である。
\end{lemma}

\begin{proof}
何を示すべきかを述べる。各 $W_j$ が $\delta$-次元 $k + 1$ の整スキームである閉埋め込みの族
$$
i_j : W_j \longrightarrow X
$$
と有理関数 $f_j \in R(W_j)^*$ が与えられ、さらに族
$\{i_j(W_j)\}_{j \in J}$ は $X$ 上局所有限であると仮定する。このとき
$$
p_*\left(\sum i_{j, *}\text{div}(f_j)\right)
$$
が $X$ 上で零と有理同値であることを示せばよい。

\medskip\noindent
この和は
$$
\sum p_*i_{j, *}\text{div}(f_j).
$$
に等しい。$W'_j \subset Y$ を $p \circ i_j$ の像である整閉部分スキームとする。
補題 \ref{chow-lemma-quasi-compact-locally-finite} により族 $\{W'_j\}$ は $Y$ で局所有限である。
したがって各 $j$ について、$p_*i_{j, *}\text{div}(f_j) = 0$ であるか、またはある
$g_j \in R(W'_j)^*$ に対して $i'_{j, *}\text{div}(g_j)$ に等しいことを示せば十分である。

\medskip\noindent
以上の議論により、$\delta$-次元 $k + 1$ の一つの整閉部分スキーム
$W \subset X$ の場合に帰着する。$f \in R(W)^*$ とし、上のとおり $W' = p(W)$ とおく。
射の可換図式
$$
\xymatrix{
W \ar[r]_i \ar[d]_{p'} & X \ar[d]^p \\
W' \ar[r]^{i'} & Y
}
$$
を得る。補題 \ref{chow-lemma-compose-pushforward} により
$p_*i_*\text{div}(f) = i'_*(p')_*\text{div}(f)$ である。
上で説明したように、$(p')_*\text{div}(f)$ が $W'$ 上の有理関数の因子または零であることを
示せばよい。三つの場合に分ける。

\medskip\noindent
$\dim_\delta(W') < k$ の場合。このとき自動的に
$(p')_*\text{div}(f) = 0$ であり、示すことはない。

\medskip\noindent
$\dim_\delta(W') = k$ の場合。この場合
$(p')_*\text{div}(f) = 0$ であることを示そう。$\eta \in W'$ を生成点とする。
$c : W_\eta \to \Spec(K)$ は $K = \kappa(\eta)$ 上の固有整曲線であり、
その関数体 $K(W_\eta)$ は $R(W)$ と同一視される。図式
$$
\xymatrix{
W_\eta \ar[r] \ar[d]_c & W \ar[d]^{p'} \\
\Spec(K) \ar[r] & W'
}
$$
がある。$f_\eta \in K(W_\eta)^*$ で $f \in R(W)^*$ に対応する有理関数を表す。
さらに、$W_\eta$ の閉点 $\xi$ は、$p'(Z) = W'$ を満たす
$\delta$-次元 $k$ の整閉部分スキーム $Z = Z_\xi \subset W$ と $1 - 1$ に対応する。
局所環 $\mathcal{O}_{W_\eta, \xi}$ と $\mathcal{O}_{W, \xi}$ は分数体の部分環として
同一視されるので、$\text{div}(f)$ における $Z_\xi$ の重複度は
$\text{ord}_{\mathcal{O}_{W_\eta, \xi}}(f_\eta)$ に等しい。
したがって $(p')_*\text{div}(f)$ における $[W']$ の重複度は、
$c_*\text{div}(f_\eta)$ における $[\Spec(K)]$ の重複度に等しい。
補題 \ref{chow-lemma-curve-principal-divisor} によりこれは零である。

\medskip\noindent
$\dim_\delta(W') = k + 1$ の場合。この場合には補題
\ref{chow-lemma-proper-pushforward-alteration} が適用でき、望むとおり、ある
$g \in R(W')^*$ に対して $p'_*\text{div}(f) = \text{div}(g)$ である。
\end{proof}















\section{有理同値と射影直線}
\label{chow-section-different-rational-equivalence}

\noindent
$(S, \delta)$ を状況 \ref{chow-situation-setup} のとおりとし、
$X$ を $S$ 上局所有限型なスキームとする。任意の閉部分スキーム
$Z \subset X \times_S \mathbf{P}^1_S = X \times \mathbf{P}^1$ に対し、
$Z_0$ と\ $Z_\infty$ で、スキーム論的閉部分スキーム
$Z_0 = \text{pr}_2^{-1}(D_0)$ および\ $Z_\infty = \text{pr}_2^{-1}(D_\infty)$ を表す。
ここで $D_0$, $D_\infty$ は (\ref{chow-equation-zero-infty}) のとおりである。

\begin{lemma}
\label{chow-lemma-rational-equivalence-family}
$(S, \delta)$ を状況 \ref{chow-situation-setup} のとおりとする。
$X$ を $S$ 上局所有限型なスキームとする。
$W \subset X \times_S \mathbf{P}^1_S$ を $\delta$-次元 $k + 1$ の整閉部分スキームとする。
$W \not = W_0$ かつ $W \not = W_\infty$ と仮定する。このとき次が成り立つ。
\begin{enumerate}
\item $W_0$, $W_\infty$ は $W$ の有効 Cartier 因子である。
\item $W_0$, $W_\infty$ は $X$ の閉部分スキームとみなせて、
$$
[W_0]_k \sim_{rat} [W_\infty]_k,
$$
である。
\item $\delta$-次元 $k + 1$ の整閉部分スキームの任意の局所有限族
$W_i \subset X \times_S \mathbf{P}^1_S$ で、$W_i \not = (W_i)_0$ かつ
$W_i \not = (W_i)_\infty$ を満たすものに対し、$X$ 上で
$\sum ([(W_i)_0]_k - [(W_i)_\infty]_k) \sim_{rat} 0$ である。
\item $\alpha \in Z_k(X)$ が $\alpha \sim_{rat} 0$ を満たすならば、上のような
整閉部分スキームの局所有限族 $W_i \subset X \times_S \mathbf{P}^1_S$ が存在して
$\alpha = \sum ([(W_i)_0]_k - [(W_i)_\infty]_k)$ となる。
\end{enumerate}
\end{lemma}

\begin{proof}
仮定により、$W$ の生成点は射影
$X \times_S \mathbf{P}^1_S \to \mathbf{P}^1_S$ の下で $D_0$ にも $D_\infty$ にも
写らない。したがって「因子」補題
\ref{divisors-lemma-pullback-effective-Cartier-defined} から (1) が従う。

\medskip\noindent
$X \times_S D_0 \to X$ は閉埋め込みなので、$W_0$ は $X$ の閉部分スキームと同型である。
$W_\infty$ についても同様である。射 $p : W \to X$ は、閉埋め込み
$W \to X \times_S \mathbf{P}^1_S$ と固有射 $X \times_S \mathbf{P}^1_S \to X$ の合成なので
固有である。補題 \ref{chow-lemma-rational-function} により、$W$ 上のサイクルとして
$[W_0]_k \sim_{rat} [W_\infty]_k$ である。明らかに
$p_*[W_0]_k = [W_0]_k$ であり、$W_\infty$ についても同様なので、補題
\ref{chow-lemma-proper-pushforward-rational-equivalence} から (2) が従う。

\medskip\noindent
(1) と (2) を踏まえると、(3) の内容は族 $\{(W_i)_0, (W_i)_\infty\}$ が
$X$ の閉部分スキームの局所有限族であることだけであり、これは明らかである。

\medskip\noindent
$\alpha \sim_{rat} 0$ と仮定する。定義により、$\delta$-次元 $k + 1$ の整閉部分スキーム
$V_i \subset X$ と有理関数 $f_i \in R(V_i)^*$ が存在して、族
$\{V_i\}_{i \in I}$ は $X$ で局所有限であり、
$\alpha = \sum (V_i \to X)_*\text{div}(f_i)$ となる。
$$
W_i \subset V_i \times_S \mathbf{P}^1_S \subset X \times_S \mathbf{P}^1_S
$$
を、補題 \ref{chow-lemma-rational-function} のとおり有理写像 $f_i$ のグラフの閉包とする。
同じ補題により $(V_i \to X)_*\text{div}(f_i)$ は
$[(W_i)_0]_k - [(W_i)_\infty]_k$ に等しい。したがって結論は明らかである。
\end{proof}

\begin{lemma}
\label{chow-lemma-closed-subscheme-cross-p1}
$(S, \delta)$ を状況 \ref{chow-situation-setup} のとおりとする。
$X$ を $S$ 上局所有限型なスキームとし、$Z$ を
$X \times \mathbf{P}^1$ の閉部分スキームとする。次を仮定する。
\begin{enumerate}
\item $\dim_\delta(Z) \leq k + 1$ である。
\item $\dim_\delta(Z_0) \leq k$, $\dim_\delta(Z_\infty) \leq k$ である。
\item $Z$ の任意の埋め込まれた点 $\xi$（「因子」定義
\ref{divisors-definition-embedded}）に対し、
$\xi \not \in Z_0 \cup Z_\infty$ または $\delta(\xi) < k$ である。
\end{enumerate}
このとき $X$ 上の $k$-サイクルとして $[Z_0]_k \sim_{rat} [Z_\infty]_k$ である。
\end{lemma}

\begin{proof}
$\{W_i\}_{i \in I}$ を、$Z$ の $\delta$-次元 $k + 1$ の既約成分全体とする。
定義のとおり $n_i > 0$ として
$$
[Z]_{k + 1} = \sum n_i[W_i]
$$
と書く。「因子」補題 \ref{divisors-lemma-components-locally-finite} により、
$\{W_i\}$ は $X \times_S \mathbf{P}^1_S$ の閉部分集合の局所有限族である。
$$
[Z_0]_k = \sum n_i[(W_i)_0]_k
$$
であると主張する。$[Z_\infty]_k$ についても同様である。
これを示せば、補題 \ref{chow-lemma-rational-equivalence-family} から結論が従う。

\medskip\noindent
$Z' \subset X$ を $\delta$-次元 $k$ の整閉部分スキームとする。
上の等式を示すには、$[Z_0]_k$ における $[Z']$ の係数 $n$ と、
$\sum n_i[(W_i)_0]_k$ における $[Z']$ の係数 $m$ が等しいことを示せば十分である。
$\xi' \in Z'$ を生成点とし、$\xi = (\xi', 0) \in  X \times_S \mathbf{P}^1_S$ とおく。
局所環 $A = \mathcal{O}_{X \times_S \mathbf{P}^1_S, \xi}$ を考える。
$I \subset A$ を $Z$ を切り出すイデアル、すなわち
$A/I = \mathcal{O}_{Z, \xi}$ となるものとする。$t \in A$ を
$X \times_S D_0$ を切り出す元、すなわち $\mathbf{P}^1$ の零における座標の引き戻しとする。
$\xi' \in Z'$ の選び方から $\delta(\xi) = k$ であり、したがって
$\dim(A/I) = 1$ である。仮定 (3) により $\xi$ は埋め込まれた点ではないので、
$A/I$ は Cohen--Macaulay である。$\dim_\delta(Z_0)
= k$ なので $\dim(A/(t, I)) = 0$ であり、これは $t$ が $A/I$ 上の非零因子であることを
含意する。最後に、$\xi$ を通る既約閉部分スキーム $W_i$ は、$I$ の上にある極小素イデアル
$I \subset \mathfrak q_i$ に対応する。重複度 $n_i$ は長さ
$\text{length}_{A_{\mathfrak q_i}}(A/I)_{\mathfrak q_i}$ に対応する。したがって
$$
n = \text{length}_A(A/(t, I))
$$
および
$$
m = \sum
\text{length}_A(A/(t, \mathfrak q_i))
\text{length}_{A_{\mathfrak q_i}}(A/I)_{\mathfrak q_i}
$$
を得る。ゆえに補題 \ref{chow-lemma-additivity-divisors-restricted} から結論が従う。
\end{proof}

\begin{lemma}
\label{chow-lemma-coherent-sheaf-cross-p1}
$(S, \delta)$ を状況 \ref{chow-situation-setup} のとおりとする。
$X$ を $S$ 上局所有限型なスキームとし、$\mathcal{F}$ を
$X \times \mathbf{P}^1$ 上の連接層とする。$i_0, i_\infty : X \to X \times \mathbf{P}^1$ を
$i_t(x) = (x, t)$ を満たす閉埋め込みとする。
$\mathcal{F}_0 = i_0^*\mathcal{F}$ および
$\mathcal{F}_\infty = i_\infty^*\mathcal{F}$ と表す。次を仮定する。
\begin{enumerate}
\item $\dim_\delta(\text{Supp}(\mathcal{F})) \leq k + 1$ である。
\item $\dim_\delta(\text{Supp}(\mathcal{F}_0)) \leq k$ および
$\dim_\delta(\text{Supp}(\mathcal{F}_\infty)) \leq k$ である。
\item $\mathcal{F}$ の任意の埋め込まれた随伴点 $\xi$ に対し、
$\xi \not \in (X \times \mathbf{P}^1)_0 \cup (X \times \mathbf{P}^1)_\infty$
または $\delta(\xi) < k$ である。
\end{enumerate}
このとき $X$ 上の $k$-サイクルとして
$[\mathcal{F}_0]_k \sim_{rat} [\mathcal{F}_\infty]_k$ である。
\end{lemma}

\begin{proof}
$\{W_i\}_{i \in I}$ を、$\text{Supp}(\mathcal{F})$ の $\delta$-次元
$k + 1$ の既約成分全体とする。定義のとおり $n_i > 0$ として
$$
[\mathcal{F}]_{k + 1} = \sum n_i[W_i]
$$
と書く。補題 \ref{chow-lemma-length-finite} により、$\{W_i\}$ は
$X \times_S \mathbf{P}^1_S$ の閉部分集合の局所有限族である。
$$
[\mathcal{F}_0]_k = \sum n_i[(W_i)_0]_k
$$
であると主張する。$[\mathcal{F}_\infty]_k$ についても同様である。
これを示せば、補題 \ref{chow-lemma-rational-equivalence-family} から結論が従う。

\medskip\noindent
$Z' \subset X$ を $\delta$-次元 $k$ の整閉部分スキームとする。
上の等式を示すには、$[\mathcal{F}_0]_k$ における $[Z']$ の係数 $n$ と、
$\sum n_i[(W_i)_0]_k$ における $[Z']$ の係数 $m$ が等しいことを示せば十分である。
$\xi' \in Z'$ を生成点とし、$\xi = (\xi', 0) \in  X \times_S \mathbf{P}^1_S$ とおく。
局所環 $A = \mathcal{O}_{X \times_S \mathbf{P}^1_S, \xi}$ を考え、
$M = \mathcal{F}_\xi$ を $A$-加群とみなす。$t \in A$ を
$X \times_S D_0$ を切り出す元、すなわち $\mathbf{P}^1$ の零における座標の引き戻しとする。
$\xi' \in Z'$ の選び方から $\delta(\xi) = k$ であり、したがって
$\dim(\text{Supp}(M)) = 1$ である。仮定 (3) により $\xi$ は $\mathcal{F}$ の随伴点では
ないので、$M$ は Cohen--Macaulay 加群である。
$\dim_\delta(\text{Supp}(\mathcal{F}_0)) = k$ なので
$\dim(\text{Supp}(M/tM)) = 0$ であり、これは $t$ が $M$ 上の非零因子であることを
含意する。最後に、$\xi$ を通る既約閉部分スキーム $W_i$ は
$\text{Ass}(M)$ の極小素イデアル $\mathfrak q_i$ に対応する。
重複度 $n_i$ は長さ $\text{length}_{A_{\mathfrak q_i}}M_{\mathfrak q_i}$ に対応する。
したがって
$$
n = \text{length}_A(M/tM)
$$
および
$$
m = \sum
\text{length}_A(A/(t, \mathfrak q_i)A)
\text{length}_{A_{\mathfrak q_i}}M_{\mathfrak q_i}
$$
を得る。ゆえに補題 \ref{chow-lemma-additivity-divisors-restricted} から結論が従う。
\end{proof}







\section{Chow 群と包絡射}
\label{chow-section-envelopes}

\noindent
定義は次のとおりである。

\begin{definition}
\label{chow-definition-envelope}
\begin{reference}
\cite[Definition 18.3]{F}
\end{reference}
$X$ をスキームとする。{\it 包絡射} とは、完全分解的な固有射 $f : Y \to X$
のことである
（『射の続論』の定義 \ref{more-morphisms-definition-cd-morphism}）。
\end{definition}

\noindent
補題 \ref{chow-lemma-envelope} の完全列が、この定義を導入する主な動機である。

\begin{lemma}
\label{chow-lemma-composition-envelope}
$(S, \delta)$ を状況 \ref{chow-situation-setup} のとおりとする。
$X$ を $S$ 上局所有限型なスキームとする。
$f : Y \to X$ と $g : Z \to Y$ が包絡射ならば、
$f \circ g$ も包絡射である。
\end{lemma}

\begin{proof}
『スキームの射』の補題 \ref{morphisms-lemma-composition-proper} および
『射の続論』の補題 \ref{more-morphisms-lemma-composition-cd} から従う。
\end{proof}

\begin{lemma}
\label{chow-lemma-base-change-envelope}
$(S, \delta)$ を状況 \ref{chow-situation-setup} のとおりとする。
$X' \to X$ を $S$ 上局所有限型なスキームの射とする。
$f : Y \to X$ が包絡射ならば、$f$ の基底変換 $f' : Y' \to X'$ も包絡射である。
\end{lemma}

\begin{proof}
『スキームの射』の補題 \ref{morphisms-lemma-base-change-proper} および
『射の続論』の補題 \ref{more-morphisms-lemma-base-change-cd} から従う。
\end{proof}

\begin{lemma}
\label{chow-lemma-envelope}
$(S, \delta)$ を状況 \ref{chow-situation-setup} のとおりとする。
$X$ を $S$ 上局所有限型なスキームとし、$f : Y \to X$ を包絡射とする。このとき、
任意の $k \in \mathbf{Z}$ に対して完全列
$$
\CH_k(Y \times_X Y) \xrightarrow{p_* - q_*}
\CH_k(Y) \xrightarrow{f_*}
\CH_k(X) \to 0
$$
が存在する。ここで $p, q : Y \times_X Y \to Y$ は射影である。
\end{lemma}

\begin{proof}
$f$ は包絡射なので、$f$ は固有であり、したがってサイクルおよびサイクル類の押し出しが
定義される。第 \ref{chow-section-proper-pushforward} 節および第 \ref{chow-section-push-pull} 節を参照せよ。
同様に、射 $p$ と $q$ は $f$ の基底変換なので固有である。さらに
$f_* \circ p_* = (p \circ f)_* = (q \circ f)_* = f_* \circ q_*$ であるから、
矢印の合成は零である。補題 \ref{chow-lemma-compose-pushforward} を参照せよ。

\medskip\noindent
$f_* : Z_k(Y) \to Z_k(X)$ が全射であることを示そう。
$Z_i \subset X$ が整閉部分スキームの局所有限族であり、
$\alpha = \sum n_i[Z_i] \in Z_k(X)$ とする。$x_i \in Z_i$ を生成点とする。
$f$ は包絡射、したがって完全分解的なので、$f(y_i) = x_i$ かつ
$\kappa(y_i)/\kappa(x_i)$ が自明となる点 $y_i \in Y$ が存在する。
$W_i \subset Y$ を生成点 $y_i$ をもつ整閉部分スキームとする。
$f$ は閉射なので $f(W_i) = Z_i$ である。したがって閉部分スキームの族 $W_i$ は
$Y$ 上局所有限である。$\kappa(y_i)/\kappa(x_i)$ は自明なので
$\dim_\delta(W_i) = \dim_\delta(Z_i) = k$ である。ゆえに
$\beta = \sum n_i[W_i]$ は $Z_k(Y)$ に属する。最後に、
$\kappa(y_i)/\kappa(x_i)$ が自明であることから、支配射
$f|_{W_i} : W_i \to Z_i$ の次数は $1$ であり、$f_*\beta = \alpha$ を得る。

\medskip\noindent
$f_* : Z_k(Y) \to Z_k(X)$ は全射なので、なおさら
$f_* : \CH_k(Y) \to \CH_k(X)$ も全射である。

\medskip\noindent
$\beta \in Z_k(Y)$ を、$f_*\beta$ が $\CH_k(X)$ において零となる元とする。
これは、$\dim_\delta(Z_j) = k + 1$ を満たす整閉部分スキーム
$Z_j \subset X$ の局所有限族と、次を満たす $f_j \in R(Z_j)^*$ が存在することを意味する。
$$
f_*\beta = \sum (Z_j \to X)_*\text{div}(f_j)
$$
ここでこれはサイクルとしての等式であり、$i_j : Z_j \to X$ は与えられた閉埋め込みである。
上とまったく同様にして、$f(W_j) = Z_j$ であり、$W_j \to Z_j$ が双有理、すなわち
同型 $R(Z_j) = R(W_j)$ を誘導する整閉部分スキーム $W_j \subset Y$ の
局所有限族を取ることができる。$f_j$ に対応する元を $g_j \in R(W_j)^*$ と表す。
$W_j \to Z_j$ は固有であり、$Z_j$ 上のサイクルとして
$(W_j \to Z_j)_*\text{div}(g_j) = \text{div}(f_j)$ であることに注意する。
したがって $\beta$ を有理同値なサイクル
$$
\beta' = \beta - \sum (W_j \to Y)_*\text{div}(g_j)
$$
で置き換えると、$f_*\beta' = 0$ となる。
（ここでは補題 \ref{chow-lemma-compose-pushforward} を用いた。）
ゆえに補題の証明を終えるには、次の段落の主張を示せば十分である。

\medskip\noindent
主張：$\beta \in Z_k(Y)$ かつ $f_*\beta = 0$ ならば、ある
$\gamma \in Z_k(Y \times_X Y)$ に対して $Z_k(Y)$ において
$\beta = \delta + p_*\gamma - q_*\gamma$ である。実際、
$\{W_j\}_{j \in J}$ を $\dim_\delta(W_j) = k$ を満たす $Y$ の整閉部分スキームの
局所有限族として、$\beta = \sum_{j \in J} n_j[W_j]$ と書く。
整閉部分スキーム $Z \subset X$ を一つ固定し、部分集合
$J_Z = \{j \in J : f(W_j) = Z\}$ を考える。これは有限集合である。次の三つの場合がある。
\begin{enumerate}
\item $J_Z = \emptyset$ の場合。このとき $\gamma_Z = 0$ とおく。
\item $J_Z \not = \emptyset$ かつ $\dim_\delta(Z) = k$ の場合。
条件 $f_*\beta = 0$ から、$Z$ の係数を比較して
$\sum_{j \in J_Z} n_j\deg(W_j/Z) = 0$ を得る。この場合、$Z$ に双有理に写る
整閉部分スキーム $W \subset Y$ を選ぶ（上を参照）。生成点を考えると、
$W_j \times_Z W$ は $W_j$ に双有理に写る唯一の既約成分
$W'_j \subset W_j \times_Z W \subset Y \times_X Y$ をもつ。
このとき $W'_j \to W$ は支配的で、$\deg(W'_j/W) = \deg(W_j/W)$ である。そこで
$\gamma_Z = \sum_{j \in J_Z} n_j[W'_j]$ とおけば、
$p_*\gamma_Z = \sum_{j \in J_Z} n_j[W_j]$ および
$q_*\gamma_Z = \sum_{j \in J_Z} n_j\deg(W'_j/W)[W] = 0$ を得る。
\item $J_Z \not = \emptyset$ かつ $\dim_\delta(Z) < k$ の場合。
この場合も $Z$ に双有理に写る整閉部分スキーム $W \subset Y$ を選ぶ（上を参照）。
生成点を考えると、$W_j \times_Z W$ は $W_j$ に双有理に写る唯一の既約成分
$W'_j \subset W_j \times_Z W \subset Y \times_X Y$ をもつ。
このとき $W'_j \to W$ は支配的で、
$k = \dim_\delta(W'_j) > \dim_\delta(W) = \dim_\delta(Z)$ である。
そこで $\gamma_Z = \sum_{j \in J_Z} n_j[W'_j]$ とおけば、
$p_*\gamma_Z = \sum_{j \in J_Z} n_j[W_j]$ および $q_*\gamma_Z = 0$ を得る。
\end{enumerate}
整閉部分スキームの族 $\{f(W_j)\}$ は $X$ 上局所有限なので
（補題 \ref{chow-lemma-quasi-compact-locally-finite}）、$Y \times_X Y$ 上の $k$-サイクル
$$
\gamma = \sum\nolimits_{Z \subset X\text{ integral closed}} \gamma_Z
$$
はよく定義される。上の計算から $p_*\gamma_Z = \beta$ および
$q_*\gamma_Z = 0$ が従い、求める結論を得る。
\end{proof}













\section{Chow 群と K 群}
\label{chow-section-chow-and-K}

\noindent
この節では、連接層の圏の $K_0$ と Chow 群を比較する。

\medskip\noindent
$(S, \delta)$ を状況 \ref{chow-situation-setup} のとおりとし、
$X$ を $S$ 上局所有限型なスキームとする。$X$ 上の連接層の圏を
$\textit{Coh}(X) = \textit{Coh}(\mathcal{O}_X)$ と表す。
これはアーベル圏である。『スキームのコホモロジー』の補題 \ref{coherent-lemma-coherent-abelian-Noetherian} を参照せよ。
任意の $k \in \mathbf{Z}$ に対して、$\textit{Coh}_{\leq k}(X)$ を
$\dim_\delta(\text{Supp}(\mathcal{F})) \leq k$ を満たす連接層 $\mathcal{F}$ からなる
$\textit{Coh}(X)$ の充満部分圏とする。

\begin{lemma}
\label{chow-lemma-Serre-subcategories}
$(S, \delta)$ を状況 \ref{chow-situation-setup} のとおりとし、
$X$ を $S$ 上局所有限型なスキームとする。圏 $\textit{Coh}_{\leq k}(X)$ は、
アーベル圏 $\textit{Coh}(X)$ の Serre 部分圏である。
\end{lemma}

\begin{proof}
Serre 部分圏の定義は『ホモロジー代数』の定義 \ref{homology-definition-serre-subcategory} にある。
補題の証明は直ちに従うので省略する。
\end{proof}

\begin{lemma}
\label{chow-lemma-cycles-k-group}
$(S, \delta)$ を状況 \ref{chow-situation-setup} のとおりとし、
$X$ を $S$ 上局所有限型なスキームとする。写像
$$
Z_k(X)
\longrightarrow
K_0(\textit{Coh}_{\leq k}(X)/\textit{Coh}_{\leq k - 1}(X)),
\quad
\sum n_Z[Z] \mapsto
\left[\bigoplus\nolimits_{n_Z > 0} \mathcal{O}_Z^{\oplus n_Z}\right]
-
\left[\bigoplus\nolimits_{n_Z < 0} \mathcal{O}_Z^{\oplus -n_Z}\right]
$$
と
$$
K_0(\textit{Coh}_{\leq k}(X)/\textit{Coh}_{\leq k - 1}(X))
\longrightarrow
Z_k(X),\quad
\mathcal{F} \longmapsto [\mathcal{F}]_k
$$
は互いに逆な同型である。
\end{lemma}

\begin{proof}
$\sum n_Z[Z]$ が $Z_k(X)$ に属するならば、族 $\{Z \mid n_Z > 0\}$ は
$X$ 上局所有限なので、直和
$\bigoplus\nolimits_{n_Z > 0} \mathcal{O}_Z^{\oplus n_Z}$ と
$\bigoplus\nolimits_{n_Z < 0} \mathcal{O}_Z^{\oplus -n_Z}$ は $X$ 上の連接層である。
写像 $\mathcal{F} \to [\mathcal{F}]_k$ は $\textit{Coh}_{\leq k}(X)$ 上加法的である。
補題 \ref{chow-lemma-additivity-sheaf-cycle} を参照せよ。また、
$\mathcal{F} \in \textit{Coh}_{\leq k - 1}(X)$ ならば $[\mathcal{F}]_k = 0$ である。
『ホモロジー代数』の補題 \ref{homology-lemma-serre-subcategory-K-groups} の (1) により、
第二の写像もよく定義される。第一の写像と第二の写像の合成が恒等写像であることは明らかである。

\medskip\noindent
逆に、$X$ 上の連接層 $\mathcal{F}$ から始める。$n_i > 0$ とし、
$Z_i \subset X$, $i \in I$ を相異なる $\delta$-次元 $k$ の整閉部分スキームとして、
$[\mathcal{F}]_k = \sum_{i \in I} n_i[Z_i]$ と書く。示すべきことは
$$
[\mathcal{F}] = [\bigoplus\nolimits_{i \in I} \mathcal{O}_{Z_i}^{\oplus n_i}]
$$
が $K_0(\textit{Coh}_{\leq k}(X)/\textit{Coh}_{\leq k - 1}(X))$ において成り立つことである。
$\xi_i \in Z_i$ を生成点と表す。
$$
\mathcal{F}' = \Ker(\mathcal{F} \to \bigoplus \xi_{i, *}\mathcal{F}_{\xi_i})
$$
とおくと、$\mathcal{F}'$ は、台の次元が $\leq k - 1$ であるような
$\mathcal{F}$ の最大の連接部分加群である。特に $\mathcal{F}$ と
$\mathcal{F}/\mathcal{F}'$ は
$K_0(\textit{Coh}_{\leq k}(X)/\textit{Coh}_{\leq k - 1}(X))$ において同じ類をもつ。
したがって $\mathcal{F}$ を $\mathcal{F}/\mathcal{F}'$ で置き換え、上に表示した核
$\mathcal{F}'$ が零であると仮定してよいし、そう仮定する。

\medskip\noindent
各 $i \in I$ に対し、点 $\xi_i$ の剰余体上で各逐次商の次元が $1$ となるフィルトレーション
$$
\mathcal{F}_{\xi_i} = \mathcal{F}_i^0 \supset \mathcal{F}_i^1 \supset
\ldots \supset \mathcal{F}_i^{n_i} = 0
$$
を選ぶ。$\mathcal{O}_{X, \xi_i}$ 上の $\mathcal{F}_{\xi_i}$ の長さは $n_i$ なので、
これは可能である。$p > n_i$ に対して $\mathcal{F}_i^p = 0$ とおく。
$p \geq 0$ に対して
$$
\mathcal{F}^p =
\Ker\left(\mathcal{F} \longrightarrow \bigoplus
\xi_{i, *}(\mathcal{F}_{\xi_i}/\mathcal{F}_i^p)\right)
$$
と表す。このとき $\mathcal{F}^p$ は連接で、$\mathcal{F}^0 = \mathcal{F}$ であり、
$\mathcal{F}^p/\mathcal{F}^{p + 1}$ は $\xi_i$ のある開近傍上、
$n_i > p$ の場合には階数 $1$ の自由 $\mathcal{O}_{Z_i}$-加群に、
$n_i \leq p$ の場合には $0$ に同型である。さらに
$\mathcal{F}' = \bigcap \mathcal{F}^p = 0$ である。任意の準コンパクト開集合
$U \subset X$ は有限個の $\xi_i$ しか含まないので、$p \gg 0$ ならば
$\mathcal{F}^p|_U$ は零である。ゆえに $\bigoplus_{p \geq 0} \mathcal{F}^p$ は
連接 $\mathcal{O}_X$-加群である。連接 $\mathcal{O}_X$-加群の短完全列
$$
0 \to
\bigoplus\nolimits_{p > 0} \mathcal{F}^p \to
\bigoplus\nolimits_{p \geq 0} \mathcal{F}^p \to
\bigoplus\nolimits_{p > 0} \mathcal{F}^p/\mathcal{F}^{p + 1} \to 0
$$
および
$$
0 \to
\bigoplus\nolimits_{p > 0} \mathcal{F}^p \to
\bigoplus\nolimits_{p \geq 0} \mathcal{F}^p \to
\mathcal{F} \to 0
$$
を考える。これだけで
$$
[\mathcal{F}] = [\bigoplus \mathcal{F}^p/\mathcal{F}^{p + 1}]
$$
が $K_0(\textit{Coh}_{\leq k}(X)/\textit{Coh}_{\leq k - 1}(X))$ において成り立つ。
次に、$p \geq 0$, $i \in I$, $n_i > p$ を満たす各組に対して、
零でないイデアル層 $\mathcal{I}_{i, p} \subset \mathcal{O}_{Z_i}$ と、
上で述べた $\xi_i$ の開近傍上で同型となる $X$ 上の写像
$\mathcal{I}_{i, p} \to \mathcal{F}^p/\mathcal{F}^{p + 1}$ を選ぶ。
『スキームのコホモロジー』の補題 \ref{coherent-lemma-extend-coherent} により、これは可能である。
そこで短完全列
$$
0 \to
\bigoplus\nolimits_{p \geq 0, i \in I, n_i > p} \mathcal{I}_{i, p}
\to
\bigoplus \mathcal{F}^p/\mathcal{F}^{p + 1} \to
\mathcal{Q} \to 0
$$
と短完全列
$$
0 \to
\bigoplus\nolimits_{p \geq 0, i \in I, n_i > p} \mathcal{I}_{i, p}
\to
\bigoplus\nolimits_{p \geq 0, i \in I, n_i > p} \mathcal{O}_{Z_i}
\to
\mathcal{Q}' \to 0
$$
を考える。$\mathcal{Q}$ と $\mathcal{Q}'$ はともに点 $\xi_i$ の近傍で零であり、
$\bigcup Z_i$ 上に台をもつ。したがって $\mathcal{Q}$ と $\mathcal{Q}'$ は
$\textit{Coh}_{\leq k - 1}(X)$ に属する。
$$
\bigoplus\nolimits_{i \in I} \mathcal{O}_{Z_i}^{\oplus n_i} \cong
\bigoplus\nolimits_{p \geq 0, i \in I, n_i > p} \mathcal{O}_{Z_i}
$$
なので証明は完了する。
\end{proof}

\begin{lemma}
\label{chow-lemma-finite-cycles-k-group}
$\pi : X \to Y$ を、状況 \ref{chow-situation-setup} のとおりの $(S, \delta)$ 上
局所有限型なスキームの有限射とする。このとき
$\pi_* : \textit{Coh}(X) \to \textit{Coh}(Y)$ は完全関手であり、
$\textit{Coh}_{\leq k}(X)$ を $\textit{Coh}_{\leq k}(Y)$ に送り、
これらの圏とその商の $K_0$ 上に準同型を誘導する。補題
\ref{chow-lemma-cycles-k-group} の写像は可換図式
$$
\xymatrix{
Z_k(X) \ar[d]^{\pi_*} \ar[r] &
K_0(\textit{Coh}_{\leq k}(X)/\textit{Coh}_{\leq k - 1}(X))
\ar[d]^{\pi_*} \ar[r] &
Z_k(X) \ar[d]^{\pi_*} \\
Z_k(Y) \ar[r] &
K_0(\textit{Coh}_{\leq k}(Y)/\textit{Coh}_{\leq k - 1}(Y)) \ar[r] &
Z_k(Y)
}
$$
をなす。
\end{lemma}

\begin{proof}
有限射はアフィンなので、$\pi$ に沿う準連接加群の押し出しは完全関手である。
『スキームのコホモロジー』の補題 \ref{coherent-lemma-relative-affine-vanishing} を参照せよ。
有限射は固有なので、$\pi_*$ は連接層を連接層に送る。『スキームのコホモロジー』の命題 \ref{coherent-proposition-proper-pushforward-coherent} を参照せよ。
台の次元に関する主張は明らかである。右側の可換性は補題
\ref{chow-lemma-cycle-push-sheaf} から直ちに従う。横向きの矢印は全単射なので、
左側も可換である。
\end{proof}

\begin{lemma}
\label{chow-lemma-from-chow-to-K}
$X$ を、状況 \ref{chow-situation-setup} のとおりの $(S, \delta)$ 上
局所有限型なスキームとする。補題 \ref{chow-lemma-cycles-k-group} の写像
$Z_k(X) \to K_0(\textit{Coh}_{\leq k}(X)/\textit{Coh}_{\leq k - 1}(X))$
が誘導する標準写像
$$
\CH_k(X)
\longrightarrow
K_0(\textit{Coh}_{\leq k + 1}(X)/\textit{Coh}_{\leq k - 1}(X))
$$
が存在する。
\end{lemma}

\begin{proof}
零に有理同値な $Z_k(X)$ の元 $\alpha$ が
$K_0(\textit{Coh}_{\leq k + 1}(X)/\textit{Coh}_{\leq k - 1}(X))$ において
零に写ることを示せばよい。定義 \ref{chow-definition-rational-equivalence} のとおり
$\alpha = \sum (i_j)_*\text{div}(f_j)$ と書く。
$$
\pi = \coprod i_j : W = \coprod W_j \longrightarrow X
$$
は有限射であることに注意する。実際、各 $i_j : W_j \to X$ は閉埋め込みであり、
$W_j$ の族は $X$ 上局所有限である。したがって補題
\ref{chow-lemma-finite-cycles-k-group} を用いて $W$ の場合に帰着できる。
$W$ は整スキームの非交和なので、次の段落で扱う場合に帰着する。

\medskip\noindent
$X$ は $\delta$-次元 $k + 1$ の整スキームであると仮定する。
$f$ を $X$ 上の零でない有理関数とし、$\alpha = \text{div}(f)$ とおく。
$\alpha$ が
$K_0(\textit{Coh}_{\leq k + 1}(X)/\textit{Coh}_{\leq k - 1}(X))$ において
零に写ることを示す。$\mathcal{I} \subset \mathcal{O}_X$ を $f$ の分母イデアルとする。
『因子』の定義 \ref{divisors-definition-regular-meromorphic-ideal-denominators} を参照せよ。
このとき短完全列
$$
0 \to \mathcal{I} \to \mathcal{O}_X \to \mathcal{O}_X/\mathcal{I} \to 0
$$
と
$$
0 \to \mathcal{I} \xrightarrow{f} \mathcal{O}_X \to
\mathcal{O}_X/f\mathcal{I} \to 0
$$
がある。『因子』の補題 \ref{divisors-lemma-regular-meromorphic-ideal-denominators} を参照せよ。等式
$$
[\mathcal{O}_X/\mathcal{I}]_k - [\mathcal{O}_X/f\mathcal{I}]_k =
\text{div}(f)
$$
を主張する。この主張により、元 $\alpha = \text{div}(f)$ は
$K_0(\textit{Coh}_{\leq k}(X)/\textit{Coh}_{\leq k - 1}(X))$ において
$[\mathcal{O}_X/\mathcal{I}] - [\mathcal{O}_X/f\mathcal{I}]$ で表される。
すると上の短完全列から、この元は
$K_0(\textit{Coh}_{\leq k + 1}(X)/\textit{Coh}_{\leq k - 1}(X))$ において零に写る。

\medskip\noindent
主張を示すため、$Z \subset X$ を $\delta$-次元 $k$ の整閉部分スキームとし、
$\xi \in Z$ をその生成点とする。このとき
$I = \mathcal{I}_\xi \subset A = \mathcal{O}_{X, \xi}$ は $fI \subset A$ を満たす
イデアルである。$\text{div}(f)$ における $[Z]$ の係数は $\text{ord}_A(f)$ である。
（通常どおり、$X$ の関数体を $A$ の分数体と同一視する。）一方、
$[\mathcal{O}_X/\mathcal{I}] - [\mathcal{O}_X/f\mathcal{I}]$ における $[Z]$ の係数は
$$
\text{length}_A(A/I) - \text{length}_A(A/fI)
$$
である。『可換代数』の定義 \ref{algebra-definition-distance} の距離関数を用いると、
これは
$$
d(A, I) - d(A, fI) = d(I, fI) = \text{ord}_A(f)
$$
と書き換えられる。等式は『可換代数』の補題 \ref{algebra-lemma-properties-distance-function} および \ref{algebra-lemma-order-vanishing-determinant} から従う。
（これらの補題を用いる必要はないが、便利である。）
\end{proof}

\begin{remark}
\label{chow-remark-good-cases-K-A}
$(S, \delta)$ を状況 \ref{chow-situation-setup} のとおりとし、
$X$ を $S$ 上局所有限型なスキームとする。後で補題
\ref{chow-lemma-cycles-rational-equivalence-K-group} において、補題
\ref{chow-lemma-from-chow-to-K} の写像
$$
\CH_k(X)
\longrightarrow
K_0(\textit{Coh}_{k + 1}(X)/\textit{Coh}_{\leq k - 1}(X))
$$
が単射であることを見る。これを標準写像
$$
K_0(\textit{Coh}_{k + 1}(X)/\textit{Coh}_{\leq k - 1}(X))
\longrightarrow
K_0(\textit{Coh}(X)/\textit{Coh}_{\leq k - 1}(X))
$$
と合成すると、標準写像
$$
\CH_k(X)
\longrightarrow
K_0(\textit{Coh}(X)/\textit{Coh}_{\leq k - 1}(X)).
$$
を得る。この写像が単射であるべきか否かについて、文献中に主張または予想を
見いだすことはできなかった。この写像の核は捩れ群であると期待するのが妥当と思われる。
この問題には後で戻る（将来の参照を挿入する）。
\end{remark}

\begin{lemma}
\label{chow-lemma-K-coherent-supported-on-closed}
$X$ を局所 Noether スキームとし、$Z \subset X$ を閉部分スキームとする。
$\textit{Coh}_Z(X) \subset \textit{Coh}(X)$ を、集合論的な台が $Z$ に含まれる
連接 $\mathcal{O}_X$-加群からなる Serre 部分圏と表す。このとき完全な包含関手
$\textit{Coh}(Z) \to \textit{Coh}_Z(X)$ は同型
$$
K'_0(Z) = K_0(\textit{Coh}(Z)) \longrightarrow K_0(\textit{Coh}_Z(X))
$$
を誘導する。
\end{lemma}

\begin{proof}
$\mathcal{F}$ を $\textit{Coh}_Z(X)$ の対象とし、
$\mathcal{I} \subset \mathcal{O}_X$ を $Z$ の準連接イデアル層とする。降下フィルトレーション
$$
\ldots \subset
\mathcal{F}^p = \mathcal{I}^p \mathcal{F} \subset
\mathcal{F}^{p - 1} \subset \ldots \subset \mathcal{F}^0 = \mathcal{F}
$$
を考える。補題 \ref{chow-lemma-from-chow-to-K} の証明とまったく同様に、このフィルトレーションは
局所有限であり、したがって
$\bigoplus_{p \geq 0} \mathcal{F}^p$,
$\bigoplus_{p \geq 1} \mathcal{F}^p$, および
$\bigoplus_{p \geq 0} \mathcal{F}^p/\mathcal{F}^{p + 1}$
は $Z$ 上に台をもつ連接 $\mathcal{O}_X$-加群である。ゆえに補題
\ref{chow-lemma-from-chow-to-K} の証明とまったく同様に、$K_0(\textit{Coh}_Z(X))$ において
$$
[\mathcal{F}] =
[\bigoplus\nolimits_{p \geq 0} \mathcal{F}^p/\mathcal{F}^{p + 1}]
$$
を得る。連接加群 $\bigoplus_{p \geq 0} \mathcal{F}^p/\mathcal{F}^{p + 1}$ は
$\mathcal{I}$ で零化されるので、$[\mathcal{F}]$ は像に属する。
実際、写像
$$
\mathcal{F} \longmapsto 
c(\mathcal{F}) =
[\bigoplus\nolimits_{p \geq 0} \mathcal{F}^p/\mathcal{F}^{p + 1}]
$$
は $K_0(\textit{Coh}_Z(X))$ を経由し、補題の主張にある写像の逆写像になると主張する。
このためには、$\textit{Coh}_Z(X)$ における短完全列
$$
0 \to \mathcal{F} \to \mathcal{G} \to \mathcal{H} \to 0
$$
に対して $c(\mathcal{G}) = c(\mathcal{F}) + c(\mathcal{H})$ を得ることを示せばよい。
任意の $q \geq 0$ に対して短完全列
$$
0 \to
(\mathcal{F} \cap \mathcal{I}^q\mathcal{G})/
(\mathcal{F} \cap \mathcal{I}^{q + 1}\mathcal{G}) \to
\mathcal{G}^q/\mathcal{G}^{q + 1} \to
\mathcal{H}^q/\mathcal{H}^{q + 1} \to 0
$$
があることに注意する。$p, q \geq 0$ に対して連接部分加群
$$
\mathcal{F}^{p, q} = \mathcal{I}^p\mathcal{F} \cap \mathcal{I}^q\mathcal{G}
$$
を考える。上とまったく同様に論じ、フィルトレーション
$\mathcal{F}^p = \mathcal{I}^p\mathcal{F}$ と
$\mathcal{F} \cap \mathcal{I}^q\mathcal{G}$ が局所有限であることを用いると、
$K_0(\textit{Coh}(Z))$ において
$$
[\bigoplus\nolimits_{p \geq 0} \mathcal{F}^p/\mathcal{F}^{p + 1}] =
[\bigoplus\nolimits_{p, q \geq 0}
\mathcal{F}^{p, q}/(\mathcal{F}^{p + 1, q} + \mathcal{F}^{p, q + 1})] =
[\bigoplus\nolimits_{q \geq 0}
(\mathcal{F} \cap \mathcal{I}^q\mathcal{G})/
(\mathcal{F} \cap \mathcal{I}^{q + 1}\mathcal{G})]
$$
を得る。これを上の完全列と組み合わせれば所望の結果が従う。細部は省略する。
\end{proof}













\section{可逆層に付随する因子}
\label{chow-section-divisor-invertible-sheaf}

\noindent
次の定義は、現在の設定における『因子』の定義 \ref{divisors-definition-divisor-invertible-sheaf} の類似である。

\begin{definition}
\label{chow-definition-divisor-invertible-sheaf}
$(S, \delta)$ を状況 \ref{chow-situation-setup} のとおりとする。
$X$ を $S$ 上局所有限型とし、$X$ は整で $n = \dim_\delta(X)$ と仮定する。
$\mathcal{L}$ を可逆 $\mathcal{O}_X$-加群とする。
\begin{enumerate}
\item $\mathcal{L}$ の任意の零でない有理型切断 $s$ に対して、
{\it $s$ に付随する Weil 因子}を、『因子』の定義 \ref{divisors-definition-divisor-invertible-sheaf} で定義された $(n - 1)$-サイクル
$$
\text{div}_\mathcal{L}(s) =
\sum \text{ord}_{Z, \mathcal{L}}(s) [Z]
$$
と定義する。補題 \ref{chow-lemma-divisor-delta-dimension} により Weil 因子の
$\delta$-次元は $n - 1$ なので、これは意味をもつ。
\item {\it $\mathcal{L}$ に付随する Weil 因子}を
$$
c_1(\mathcal{L}) \cap [X] =
\text{class of }\text{div}_\mathcal{L}(s) \in \CH_{n - 1}(X)
$$
と定義する。ここで $s$ は $X$ 上の $\mathcal{L}$ の任意の零でない有理型切断である。
『因子』の補題 \ref{divisors-lemma-divisor-meromorphic-well-defined} により、
これはよく定義される。
\end{enumerate}
\end{definition}

\noindent
$X$ と $S$ を上の定義 \ref{chow-definition-divisor-invertible-sheaf} のとおりとし、
$n = \dim_\delta(X)$ とおく。定義から、$Cl(X)$ を『因子』の定義 \ref{divisors-definition-class-group} で定義される $X$ の Weil 因子類群とすると、
$Cl(X) = \CH_{n - 1}(X)$ であることは明らかである。写像
$$
\Pic(X) \longrightarrow \CH_{n - 1}(X), \quad
\mathcal{L} \longmapsto c_1(\mathcal{L}) \cap [X]
$$
は、任意の局所 Noether 整スキームについて Divisors, Equation
(\ref{divisors-equation-c1}) で構成された写像 $\Pic(X) \to Cl(X)$ と同じである。
特に、この写像はアーベル群の準同型であり、$X$ が正規スキームならば単射、
$X$ のすべての局所環が UFD ならば同型である。『因子』の補題 \ref{divisors-lemma-normal-c1-injective} および \ref{divisors-lemma-local-rings-UFD-c1-bijective} を参照せよ。
可逆層に付随する Weil 因子を容易に計算できる場合がいくつかある。

\begin{lemma}
\label{chow-lemma-compute-c1}
$(S, \delta)$ を状況 \ref{chow-situation-setup} のとおりとする。
$X$ を $S$ 上局所有限型とし、$X$ は整で $n = \dim_\delta(X)$ と仮定する。
$\mathcal{L}$ を可逆 $\mathcal{O}_X$-加群とし、
$s \in \Gamma(X, \mathcal{L})$ を零でない大域切断とする。このとき
$$
\text{div}_\mathcal{L}(s) = [Z(s)]_{n - 1}
$$
が $Z_{n - 1}(X)$ において成り立ち、
$$
c_1(\mathcal{L}) \cap [X] = [Z(s)]_{n - 1}
$$
が $\CH_{n - 1}(X)$ において成り立つ。
\end{lemma}

\begin{proof}
$Z \subset X$ を $\delta$-次元 $n - 1$ の整閉部分スキームとし、
$\xi \in Z$ をその生成点とする。生成元 $s_\xi \in \mathcal{L}_\xi$ を選び、
ある $f \in \mathcal{O}_{X, \xi}$ に対して $s = fs_\xi$ と書く。
$Z(s)$ の定義（『因子』の定義 \ref{divisors-definition-zero-scheme-s}）により、$Z(s)$ は
$\mathcal{I}_\xi = (f)$ を満たす準連接イデアル層
$\mathcal{I} \subset \mathcal{O}_X$ で切り出される。したがって所望のとおり
$\text{length}_{\mathcal{O}_{X, \xi}}(\mathcal{O}_{Z(s), \xi})
=
\text{length}_{\mathcal{O}_{X, \xi}}(\mathcal{O}_{X, \xi}/(f))
=
\text{ord}_{\mathcal{O}_{X, \xi}}(f)$ である。
\end{proof}

\noindent
次の補題は、より一般的な補題 \ref{chow-lemma-flat-pullback-cap-c1} によって後に置き換えられる。

\begin{lemma}
\label{chow-lemma-flat-pullback-divisor-invertible-sheaf}
$(S, \delta)$ を状況 \ref{chow-situation-setup} のとおりとする。
$X$, $Y$ を $S$ 上局所有限型とし、$X$, $Y$ は整で
$n = \dim_\delta(Y)$ と仮定する。$\mathcal{L}$ を可逆 $\mathcal{O}_Y$-加群とし、
$f : X \to Y$ を相対次元 $r$ の平坦射とする。このとき
$$
f^*(c_1(\mathcal{L}) \cap [Y]) = c_1(f^*\mathcal{L}) \cap [X]
$$
が $\CH_{n + r - 1}(X)$ において成り立つ。
\end{lemma}

\begin{proof}
$s$ を $\mathcal{L}$ の零でない有理型切断とする。実際に
$f^*\text{div}_\mathcal{L}(s) = \text{div}_{f^*\mathcal{L}}(f^*s)$
であることを示す。これから補題が従う。点 $\xi \in Y$ を取り、
生成元 $s_\xi \in \mathcal{L}_\xi$ を選ぶ。$g \in R(Y)^*$ に対して
$s = gs_\xi$ と書く。このとき $\xi$ の開近傍 $V \subset Y$ で、
$s_\xi \in \mathcal{L}(V)$ かつ $s_\xi$ が $\mathcal{L}|_V$ を生成するものが存在する。
したがって
$$
\text{div}_\mathcal{L}(s)|_V = \text{div}_Y(g)|_V.
$$
まったく同様に、$f^*s_\xi$ は $f^{-1}(V)$ 上で $f^*\mathcal{L}$ を生成し、
$f^*s = g f^*s_\xi$ なので、
$$
\text{div}_\mathcal{L}(f^*s)|_{f^{-1}(V)}
=
\text{div}_X(g)|_{f^{-1}(V)}.
$$
も成り立つ。ゆえに $f^{-1}(V)$ 上で所望のサイクルの等式は、主因子の引き戻しに
関する対応する結果、すなわち補題 \ref{chow-lemma-flat-pullback-principal-divisor} から従う。
\end{proof}



\section{可逆層とのキャップ積}
\label{chow-section-intersecting-with-divisors}

\noindent
この節では次の構成を調べる。

\begin{definition}
\label{chow-definition-cap-c1}
$(S, \delta)$ を状況 \ref{chow-situation-setup} のとおりとする。
$X$ を $S$ 上局所有限型とし、$\mathcal{L}$ を可逆 $\mathcal{O}_X$-加群とする。
各整数 $k$ に対し、{\it $\mathcal{L}$ の第一 Chern 類とのキャップ積}と呼ばれる作用
$$
c_1(\mathcal{L}) \cap - :
Z_{k + 1}(X) \to \CH_k(X)
$$
を次のように定義する。
\begin{enumerate}
\item $\dim_\delta(W) = k + 1$ を満たす整閉部分スキーム
$i : W \to X$ が与えられたとき、
$$
c_1(\mathcal{L}) \cap [W] = i_*(c_1({i^*\mathcal{L}}) \cap [W])
$$
と定義する。右辺は定義 \ref{chow-definition-divisor-invertible-sheaf} で定義されている。
\item 一般の $(k + 1)$-サイクル $\alpha = \sum n_i [W_i]$ に対して
$$
c_1(\mathcal{L}) \cap \alpha = \sum n_i c_1(\mathcal{L}) \cap [W_i]
$$
とおく。
\end{enumerate}
\end{definition}

\noindent
各 $c_1(\mathcal{L}) \cap W_i = \sum_j n_{i, j} [Z_{i, j}]$ を、
$\{Z_{i, j}\}_j$ が $W_i$ の整閉部分スキームの局所有限和となるように書く。
$\{W_i\}$ は $X$ 上の整閉部分スキームの局所有限族なので、
$\{Z_{i, j}\}_{i, j}$ が $X$ の閉部分スキームの局所有限族であることは容易に従う。
したがって $c_1(\mathcal{L}) \cap \alpha = \sum n_in_{i, j}[Z_{i, j}]$ はサイクルである。
別の、より便利な見方は、射 $\coprod W_i \to X$ が固有であることに注意することである。
ゆえに $c_1(\mathcal{L}) \cap \alpha$ は
$\CH_k(\coprod W_i) = \prod \CH_k(W_i)$ の類の押し出しとみなせる。
これは、結果が $X$ 上の有理同値を除いてよく定義される理由も説明する。

\medskip\noindent
次の数節の主目標は、$c_1(\mathcal{L})$ とのキャップ積が有理同値を経由することを
示すことである。これは自明ではない。

\begin{lemma}
\label{chow-lemma-c1-cap-additive}
$(S, \delta)$ を状況 \ref{chow-situation-setup} のとおりとする。
$X$ を $S$ 上局所有限型とし、$\mathcal{L}$, $\mathcal{N}$ を $X$ 上の可逆層とする。
このとき、任意の $\alpha \in Z_{k + 1}(X)$ に対し
$$
c_1(\mathcal{L}) \cap \alpha  + c_1(\mathcal{N}) \cap \alpha =
c_1(\mathcal{L} \otimes_{\mathcal{O}_X} \mathcal{N}) \cap \alpha
$$
が $\CH_k(X)$ において成り立つ。さらに、すべての $\alpha$ に対して
$c_1(\mathcal{O}_X) \cap \alpha = 0$ である。
\end{lemma}

\begin{proof}
加法性は『因子』の補題 \ref{divisors-lemma-c1-additive} と定義から直ちに従う。
$c_1(\mathcal{O}_X) \cap \alpha = 0$ を見るには、切断
$1 \in \Gamma(X, \mathcal{O}_X)$ を考える。これは任意の整閉部分スキーム
$W \subset X$ 上で至る所零でない切断に制限される。したがって所望のとおり
$c_1(\mathcal{O}_X) \cap [W] = 0$ である。
\end{proof}

\noindent
$Z(s) \subset X$ は、スキーム $X$ 上の可逆層の大域切断 $s$ の零スキームを表すことを
思い出そう。『因子』の定義 \ref{divisors-definition-zero-scheme-s} を参照せよ。

\begin{lemma}
\label{chow-lemma-prepare-geometric-cap}
$(S, \delta)$ を状況 \ref{chow-situation-setup} のとおりとする。
$Y$ を $S$ 上局所有限型とし、$\mathcal{L}$ を可逆 $\mathcal{O}_Y$-加群、
$s \in \Gamma(Y, \mathcal{L})$ とする。次を仮定する。
\begin{enumerate}
\item $\dim_\delta(Y) \leq k + 1$ である。
\item $\dim_\delta(Z(s)) \leq k$ である。
\item $Z(s)$ の $\delta$-次元 $k$ の既約成分の任意の生成点 $\xi$ において、
$s$ による乗法が単射 $\mathcal{O}_{Y, \xi} \to \mathcal{L}_\xi$ を誘導する。
\end{enumerate}
$Y_i \subset Y$ を $Y$ の $\delta$-次元 $k + 1$ の既約成分として
$[Y]_{k + 1} = \sum n_i[Y_i]$ と書く。
$s_i = s|_{Y_i} \in \Gamma(Y_i, \mathcal{L}|_{Y_i})$ とおく。このとき
\begin{equation}
\label{chow-equation-equal-as-cycles}
[Z(s)]_k =  \sum n_i[Z(s_i)]_k
\end{equation}
が $Y$ 上の $k$-サイクルとして成り立つ。
\end{lemma}

\begin{proof}
$Z \subset Y$ を $\delta$-次元 $k$ の整閉部分スキームとし、
$\xi \in Z$ をその生成点とする。式 $\sum n_i[Z(s_i)]_k$ における
$[Z]$ の係数 $n$ と、式 $[Z(s)]_k$ における $[Z]$ の係数 $m$ を比較する。
生成元 $s_\xi \in \mathcal{L}_\xi$ を選び、
$A = \mathcal{O}_{Y, \xi}$, $L = \mathcal{L}_\xi$ と書く。このとき $L = As_\xi$ である。
一意な $f \in A$ に対して $s = f s_\xi$ と書く。仮定 (3) は
$f : A \to A$ が単射であることを意味する。
$\dim_\delta(Y) \leq k + 1$ かつ $\dim_\delta(Z) = k$ なので、
$\dim(A) = 0$ または $1$ である。さらに
$$
m = \text{length}_A(A/(f))
$$
であり、これはいずれの場合にも有限である。

\medskip\noindent
$\dim(A) = 0$ の場合、$f : A \to A$ が単射であることから
$f \in A^*$ が従う。ゆえにこの場合 $m$ は零である。また条件
$\dim(A) = 0$ は、$\xi$ が $\delta$-次元 $k + 1$ のいかなる既約成分にも
属さないことを意味する。すなわち $n = 0$ でもある。

\medskip\noindent
次に $\dim(A) = 1$ とする。$A$ は Noether 局所環なので、有限個の極小素イデアル
$\mathfrak q_1, \ldots, \mathfrak q_t$ をもつ。これらは $\xi'$ を通る $Y_i$ と
一対一に対応する。さらに
$n_i = \text{length}_{A_{\mathfrak q_i}}(A_{\mathfrak q_i})$ である。
また $[Z(s_i)]_k$ における $[Z]$ の重複度は
$\text{length}_A(A/(f, \mathfrak q_i))$ である。したがって、この場合に示すべき等式は
$$
\text{length}_A(A/(f))
=
\sum \text{length}_{A_{\mathfrak q_i}}(A_{\mathfrak q_i})
\text{length}_A(A/(f, \mathfrak q_i))
$$
であり、これは補題 \ref{chow-lemma-additivity-divisors-restricted} から従う。
\end{proof}

\noindent
次の補題は、可逆層の $c_1$ と閉部分スキームに付随するサイクルとのキャップ積を
計算するために有用な結果である。$Z(s) \subset X$ は、スキーム $X$ 上の可逆層の
大域切断 $s$ の零スキームを表すことを思い出そう。『因子』の定義 \ref{divisors-definition-zero-scheme-s} を参照せよ。

\begin{lemma}
\label{chow-lemma-geometric-cap}
$(S, \delta)$ を状況 \ref{chow-situation-setup} のとおりとする。
$X$ を $S$ 上局所有限型とし、$\mathcal{L}$ を可逆 $\mathcal{O}_X$-加群とする。
$Y \subset X$ を閉部分スキームとし、$s \in \Gamma(Y, \mathcal{L}|_Y)$ とする。
次を仮定する。
\begin{enumerate}
\item $\dim_\delta(Y) \leq k + 1$ である。
\item $\dim_\delta(Z(s)) \leq k$ である。
\item $Z(s)$ の $\delta$-次元 $k$ の既約成分の任意の生成点 $\xi$ において、
$s$ による乗法が単射
$\mathcal{O}_{Y, \xi} \to (\mathcal{L}|_Y)_\xi$\footnote{たとえば、
$s$ が $\mathcal{L}|_Y$ の正則切断ならば、この条件は成り立つ。}を誘導する。
\end{enumerate}
このとき
$$
c_1(\mathcal{L}) \cap [Y]_{k + 1} = [Z(s)]_k
$$
が $\CH_k(X)$ において成り立つ。
\end{lemma}

\begin{proof}
$Y_i \subset Y$ を $Y$ の $\delta$-次元 $k + 1$ の既約成分とし、$n_i > 0$ として
$$
[Y]_{k + 1} = \sum n_i[Y_i]
$$
と書く。仮定により、制限
$s_i = s|_{Y_i} \in \Gamma(Y_i, \mathcal{L}|_{Y_i})$ は零でなく、
したがって正則切断である。補題 \ref{chow-lemma-compute-c1} により
$[Z(s_i)]_k$ は $c_1(\mathcal{L}|_{Y_i})$ を表す。ゆえに定義により
$$
c_1(\mathcal{L}) \cap [Y]_{k + 1} = \sum n_i[Z(s_i)]_k
$$
である。したがって結果は補題 \ref{chow-lemma-prepare-geometric-cap} から従う。
\end{proof}




\section{可逆層とのキャップ積と押し出し・引き戻し}
\label{chow-section-intersecting-with-divisors-push-pull}

\noindent
この節では、作用 $c_1(\mathcal{L}) \cap -$ が平坦引き戻しおよび固有押し出しと
可換であることを証明する。

\begin{lemma}
\label{chow-lemma-prepare-flat-pullback-cap-c1}
$(S, \delta)$ を状況 \ref{chow-situation-setup} のとおりとする。
$X$, $Y$ を $S$ 上局所有限型とし、$f : X \to Y$ を相対次元 $r$ の平坦射とする。
$\mathcal{L}$ を $Y$ 上の可逆層とする。$Y$ は整で
$n = \dim_\delta(Y)$ と仮定し、$s$ を $\mathcal{L}$ の零でない有理型切断とする。
このとき
$$
f^*\text{div}_\mathcal{L}(s) = \sum n_i\text{div}_{f^*\mathcal{L}|_{X_i}}(s_i)
$$
が $Z_{n + r - 1}(X)$ において成り立つ。ここで和は
$\delta$-次元 $n + r$ の既約成分 $X_i \subset X$ にわたり、切断
$s_i = f|_{X_i}^*(s)$ は $s$ の引き戻しであり、
$n_i = m_{X_i, X}$ は $X$ における $X_i$ の重複度である。
\end{lemma}

\begin{proof}
このサイクルの等式を示すには、$Y$ 上局所的に考えてよい。
したがって $Y$ はアフィンで、零でない切断
$p \in \Gamma(Y, \mathcal{L})$ と $q \in \Gamma(Y, \mathcal{O})$ に対して
$s = p/q$ と仮定してよい。明らかな記法の下で
$$
f^*\text{div}_\mathcal{L}(p) =
\sum n_i\text{div}_{f^*\mathcal{L}|_{X_i}}(p_i)
\quad\text{and}\quad
f^*\text{div}_\mathcal{O}(q) =
\sum n_i\text{div}_{\mathcal{O}_{X_i}}(q_i)
$$
をともに示せれば、『因子』の補題 \ref{divisors-lemma-c1-additive} の加法性から
結論を得る。ゆえに $s \in \Gamma(Y, \mathcal{L})$ と仮定してよい。
この場合、等式 (\ref{chow-equation-equal-as-cycles}) を適用して
$$
[Z(f^*(s))]_{k + r - 1} =
\sum n_i\text{div}_{f^*\mathcal{L}|_{X_i}}(s_i)
$$
を得る。ここで $f^*(s) \in f^*\mathcal{L}$ は $s$ の $X$ への引き戻しを表す。
一方、補題 \ref{chow-lemma-compute-c1} and \ref{chow-lemma-pullback-coherent} により
$$
f^*\text{div}_\mathcal{L}(s) = f^*[Z(s)]_{k - 1}
= [f^{-1}(Z(s))]_{k + r - 1},
$$
である。$Z(f^*(s)) = f^{-1}(Z(s))$ なので結論を得る。
\end{proof}

\begin{lemma}
\label{chow-lemma-flat-pullback-cap-c1}
$(S, \delta)$ を状況 \ref{chow-situation-setup} のとおりとする。
$X$, $Y$ を $S$ 上局所有限型とし、$f : X \to Y$ を相対次元 $r$ の平坦射とする。
$\mathcal{L}$ を $Y$ 上の可逆層とし、$\alpha$ を $Y$ 上の $k$-サイクルとする。
このとき
$$
f^*(c_1(\mathcal{L}) \cap \alpha) = c_1(f^*\mathcal{L}) \cap f^*\alpha
$$
が $\CH_{k + r - 1}(X)$ において成り立つ。
\end{lemma}

\begin{proof}
$\alpha = \sum n_i[W_i]$ と書く。$X$ の閉部分スキーム $f^{-1}(W_i)$ 上に
有理同値を構成することにより
$$
f^*(c_1(\mathcal{L}) \cap [W_i]) = c_1(f^*\mathcal{L}) \cap f^*[W_i]
$$
が $\CH_{k + r - 1}(X)$ において成り立つことを示す。
注意 \ref{chow-remark-infinite-sums-rational-equivalences} の議論により、
これで補題の等式が証明される。

\medskip\noindent
$W \subset Y$ を $\delta$-次元 $k$ の整閉部分スキームとする。
閉部分スキーム $W' = f^{-1}(W) = W \times_Y X$ を考えると、ファイバー積図式
$$
\xymatrix{
W' \ar[r] \ar[d]_h & X \ar[d]^f \\
W \ar[r] & Y
}
$$
がある。示すべきことは
$f^*(c_1(\mathcal{L}) \cap [W]) = c_1(f^*\mathcal{L}) \cap f^*[W]$ である。
$\mathcal{L}|_W$ の零でない有理型切断 $s$ を選ぶ。$W'_i \subset W'$ を
$\delta$-次元 $k + r$ の既約成分とし、定義のとおり $n_i$ を $W'$ における
$W'_i$ の重複度として $[W']_{k + r} = \sum n_i[W'_i]$ と書く。
すると $Z_{k + r}(X)$ において $f^*[W] = \sum n_i[W'_i]$ である。
各 $W'_i \to W$ は支配的なので、各 $i$ に対して
$s_i = s|_{W'_i}$ は零でない有理型切断である。補題
\ref{chow-lemma-prepare-flat-pullback-cap-c1} によりサイクルの等式
$$
h^*\text{div}_{\mathcal{L}|_W}(s) =
\sum n_i\text{div}_{f^*\mathcal{L}|_{W'_i}}(s_i)
$$
が $Z_{k + r - 1}(W')$ において成り立つ。左辺は $W'$ 上のサイクルであり、
$\CH_{k + r - 1}(X)$ において $f^*(c_1(\mathcal{L}) \cap [W])$ に押し出される。
右辺も $W'$ 上のサイクルであり、$\CH_{k + r - 1}(X)$ において
$c_1(f^*\mathcal{L}) \cap f^*[W]$ に押し出される。これで証明は完了する。
\end{proof}

\begin{lemma}
\label{chow-lemma-equal-c1-as-cycles}
$(S, \delta)$ を状況 \ref{chow-situation-setup} のとおりとする。
$X$, $Y$ を $S$ 上局所有限型とし、$f : X \to Y$ を固有射、
$\mathcal{L}$ を $Y$ 上の可逆層とする。$s$ を $Y$ 上の $\mathcal{L}$ の
零でない有理型切断 $s$ とする。$X$, $Y$ は整で、$f$ は支配的かつ
$\dim_\delta(X) = \dim_\delta(Y)$ と仮定する。このとき $Y$ 上のサイクルとして
$$
f_*\left(\text{div}_{f^*\mathcal{L}}(f^*s)\right) =
[R(X) : R(Y)]\text{div}_\mathcal{L}(s).
$$
が成り立つ。特に
$$
f_*(c_1(f^*\mathcal{L}) \cap [X]) =
[R(X) : R(Y)] c_1(\mathcal{L}) \cap [Y] =
c_1(\mathcal{L}) \cap f_*[X]
$$
である。
\end{lemma}

\begin{proof}
定義により $f_*[X] = [R(X) : R(Y)][Y]$ なので、最後の等式は最初の等式から従う。
この証明には補題 \ref{chow-lemma-proper-pushforward-alteration} を再利用できる。
実際、サイクルの等式を示すので $Y$ 上局所的に考えてよく、
$\mathcal{L} = \mathcal{O}_Y$ と仮定してよい。この場合 $s$ は有理関数
$g \in R(Y)$ に対応し、示すべきことは単に
$$
f_*\left(\text{div}_X(g)\right) =
[R(X) : R(Y)]\text{div}_Y(g).
$$
である。上記の補題 \ref{chow-lemma-proper-pushforward-alteration} の結果と比較すると、
$g \in R(Y)^*$ であることから
$\text{Nm}_{R(X)/R(Y)}(g) = g^{[R(X) : R(Y)]}$ なので、これは成り立つ。
\end{proof}

\begin{lemma}
\label{chow-lemma-pushforward-cap-c1}
$(S, \delta)$ を状況 \ref{chow-situation-setup} のとおりとする。
$X$, $Y$ を $S$ 上局所有限型とし、$p : X \to Y$ を固有射とする。
$\alpha \in Z_{k + 1}(X)$ とし、$\mathcal{L}$ を $Y$ 上の可逆層とする。このとき
$$
p_*(c_1(p^*\mathcal{L}) \cap \alpha) = c_1(\mathcal{L}) \cap p_*\alpha
$$
が $\CH_k(Y)$ において成り立つ。
\end{lemma}

\begin{proof}
任意の整閉部分スキーム $W \subset X$ に対して写像 $p|_W : W \to Y$ が
閉埋め込みとなる性質を $p$ がもつと仮定する。このとき $c_1(\mathcal{L})$ との
キャップ積の定義により補題は成り立つ。

\medskip\noindent
この注意を用いて特別な場合に帰着する。$n_i \not = 0$ とし、$W_i$ は相異なるとして
$\alpha = \sum n_i[W_i]$ と書く。$W'_i \subset Y$ を $W_i$ の像とする
（整閉部分スキームとして）。図式
$$
\xymatrix{
X' = \coprod W_i \ar[r]_-q \ar[d]_{p'} & X \ar[d]^p \\
Y' = \coprod W'_i \ar[r]^-{q'} & Y.
}
$$
を考える。$\{W_i\}$ は $X$ 上局所有限で、$p$ は固有なので、
$\{W'_i\}$ は $Y$ 上局所有限であり、$q, q', p'$ も固有射である。
$\sum n_i[W_i]$ を $k$-サイクル $\alpha' \in Z_k(X')$ とみなすこともできる。
明らかに $q_*\alpha' = \alpha$ である。証明冒頭の注意により
$q_*(c_1(q^*p^*\mathcal{L}) \cap \alpha')
= c_1(p^*\mathcal{L}) \cap q_*\alpha'$
および
$(q')_*(c_1((q')^*\mathcal{L}) \cap p'_*\alpha') =
c_1(\mathcal{L}) \cap q'_*p'_*\alpha'$ が成り立つ。したがって射 $p'$ と
サイクル $\sum n_i[W_i]$ に対して補題を示せば十分である。
明らかに、これは $X$, $Y$ が整で、$f : X \to Y$ が支配的かつ
$\alpha = [X]$ である場合を仮定してよいことを意味する。
この場合、結果は補題 \ref{chow-lemma-equal-c1-as-cycles} から従う。
\end{proof}






\section{鍵となる公式}
\label{chow-section-key}

\noindent
$(S, \delta)$ を状況 \ref{chow-situation-setup} のとおりとする。
$X$ を $S$ 上局所有限型とし、$X$ は整で $\dim_\delta(X) = n$ と仮定する。
$\mathcal{L}$ と $\mathcal{N}$ を $X$ 上の可逆層とする。
$s$ を $\mathcal{L}$ の零でない有理型切断、$t$ を $\mathcal{N}$ の
零でない有理型切断とする。$Z_i \subset X$, $i \in I$ を次の性質をもつ
余次元 $1$ の既約閉部分集合の局所有限集合とする。生成点が $\xi$ である
$Z \not \in \{Z_i\}$ に対して、$s$ は $\mathcal{L}_\xi$ の生成元であり、
$t$ は $\mathcal{N}_\xi$ の生成元である。このような集合は『因子』の補題 \ref{divisors-lemma-divisor-meromorphic-locally-finite} により存在する。このとき
$$
\text{div}_\mathcal{L}(s) = \sum \text{ord}_{Z_i, \mathcal{L}}(s) [Z_i]
$$
であり、同様に
$$
\text{div}_\mathcal{N}(t) = \sum \text{ord}_{Z_i, \mathcal{N}}(t) [Z_i]
$$
である。さらに定義を展開する。各 $i$ に対して生成元
$s_i \in \mathcal{L}_{\xi_i}$ と $t_i \in \mathcal{N}_{\xi_i}$ を選ぶ。
ここで $\xi_i$ は $Z_i$ の生成点である。このとき
$$
s = f_i s_i
\quad\text{and}\quad
t = g_i t_i
$$
と書ける。$B_i = \mathcal{O}_{X, \xi_i}$ とおく。定義により
$$
\text{ord}_{Z_i, \mathcal{L}}(s) = \text{ord}_{B_i}(f_i)
\quad\text{and}\quad
\text{ord}_{Z_i, \mathcal{N}}(t) = \text{ord}_{B_i}(g_i)
$$
である。$t_i$ は $\mathcal{N}_{\xi_i}$ の生成元なので、ファイバー
$\mathcal{N}_{\xi_i} \otimes \kappa(\xi_i)$ におけるその像は
$\mathcal{N}|_{Z_i}$ の零でない有理型切断である。この像を $t_i|_{Z_i}$ と表す。
定義から
$$
c_1(\mathcal{N}) \cap \text{div}_\mathcal{L}(s) =
\sum \text{ord}_{B_i}(f_i)
(Z_i \to X)_*\text{div}_{\mathcal{N}|_{Z_i}}(t_i|_{Z_i})
$$
であり、同様に
$$
c_1(\mathcal{L}) \cap \text{div}_\mathcal{N}(t) =
\sum \text{ord}_{B_i}(g_i)
(Z_i \to X)_*\text{div}_{\mathcal{L}|_{Z_i}}(s_i|_{Z_i})
$$
が $\CH_{n - 2}(X)$ において成り立つ。この二つのサイクルの間の有理同値を求める。
そのため、第 \ref{chow-section-tame-symbol} 節の tame 記号
$$
\partial_{B_i}(f_i, g_i) \in \kappa(\xi_i)^*
$$
を考える。

\begin{lemma}[鍵となる公式]
\label{chow-lemma-key-formula}
上の状況において、サイクル
$$
\sum
(Z_i \to X)_*\left(
\text{ord}_{B_i}(f_i) \text{div}_{\mathcal{N}|_{Z_i}}(t_i|_{Z_i}) -
\text{ord}_{B_i}(g_i) \text{div}_{\mathcal{L}|_{Z_i}}(s_i|_{Z_i}) \right)
$$
はサイクル
$$
\sum (Z_i \to X)_*\text{div}(\partial_{B_i}(f_i, g_i))
$$
に等しい。
\end{lemma}

\begin{proof}
まず、$B_i$ の単元 $u$ によって $s_i$ を $us_i$ に置き換えたときに何が起こるかを
調べる。このとき $f_i$ は $u^{-1} f_i$ に置き換わる。したがって補題の第一の式の
第一部分は変わらず、第二部分には『因子』の補題 \ref{divisors-lemma-divisor-meromorphic-well-defined} により
$$
-\text{ord}_{B_i}(g_i)\text{div}(u|_{Z_i})
$$
を加える。ここで $u|_{Z_i}$ は剰余体における $u$ の像である。一方、第二の式には
tame 記号の双線型性により
$$
\text{div}(\partial_{B_i}(u^{-1}, g_i))
$$
を加える。これらの項は tame 記号の性質 (\ref{chow-item-normalization}) により一致する。

\medskip\noindent
$Z \subset X$ を $\dim_\delta(Z) = n - 2$ を満たす既約閉部分集合とする。
補題の二つのサイクルにおける $Z$ の係数が同じであることを示す際には、前段落のように
$s_i \mapsto us_i$ と置き換えてよい。まったく同様に、$B_i$ の単元 $v$ によって
$t_i \mapsto vt_i$ と置き換えてよいことも分かる。

\medskip\noindent
サイクルの等式を示すので、係数ごとに論じてよい。そこで
$\dim_\delta(Z) = n - 2$ を満たす既約閉部分集合 $Z \subset X$ を選び、係数を比較する。
$\xi \in Z$ を生成点とし、$A = \mathcal{O}_{X, \xi}$ とおく。これは次元 $2$ の
Noether 局所整域である。$\mathcal{L}_\xi$ と $\mathcal{N}_\xi$ の生成元
$\sigma$ と $\tau$ を選ぶ。$X$ を縮小した後、$\sigma$ と $\tau$ が $X$ 全体上で
可逆層 $\mathcal{L}$ と $\mathcal{N}$ の自明化を定めると仮定してよいし、そう仮定する。
$X$ を縮小した後 $Z_i$ は局所有限なので、すべての $i \in I$ に対して
$Z \subset Z_i$ であり、$I$ は有限であると仮定してよい。このとき $\xi_i$ は
高さ $1$ の素イデアル $\mathfrak q_i \subset A$ に対応する。
$A_{\mathfrak q_i}$ の単元 $a_i$, $b_i$ によって
$s_i = a_i \sigma$ および $t_i = b_i \tau$ と書ける。
上の注意により、すべての $i$ に対して $a_i = b_i = 1$ である場合に補題を示せば十分である。

\medskip\noindent
すべての $i$ に対して $a_i = b_i = 1$ と仮定する。$\sigma$ と $\tau$ を
自明化切断として選んだので、補題の第一の式は零である。$A$ の分数体の元 $f$, $g$ によって
$s = f\sigma$ および $t = g \tau$ と書く。前段落により、すべての $i$ に対して
$f_i = f$ および $g_i = g$ である場合に帰着している。さらに
$\{\mathfrak q_i\}$ に属さない $A$ の高さ $1$ の素イデアル $\mathfrak q$ に対して、
$f$ と $g$ はともに $A_\mathfrak q$ の単元である。これは補題に先立つ議論での
族 $\{Z_i\}$ の選び方による。したがって補題の第二の式における $Z$ の係数は
$$
\sum\nolimits_i \text{ord}_{A/\mathfrak q_i}(\partial_{B_i}(f, g))
$$
であり、鍵補題 \ref{chow-lemma-milnor-gersten-low-degree} により零である。
\end{proof}

\begin{remark}
\label{chow-remark-higher-chow-pointwise}
$(S, \delta)$ を状況 \ref{chow-situation-setup} のとおりとする。
$X$ を $S$ 上局所有限型とし、$k \in \mathbf{Z}$ とする。複体
$$
\bigoplus\nolimits_{\delta(x) = k + 2}' K_2^M(\kappa(x))
\xrightarrow{\partial}
\bigoplus\nolimits_{\delta(x) = k + 1}' K_1^M(\kappa(x))
\xrightarrow{\partial}
\bigoplus\nolimits_{\delta(x) = k}' K_0^M(\kappa(x))
$$
が存在すると主張する。ここでは注意 \ref{chow-remark-chow-group-pointwise} で導入した
記法と規約を用い、さらに次を定める。
\begin{enumerate}
\item $K_2^M(\kappa(x))$ は、点 $x \in X$ の剰余体 $\kappa(x)$ の
Milnor K-理論の次数 $2$ の部分である（注意 \ref{chow-remark-gersten-complex-milnor} を参照）。これは
$\kappa(x)^* \otimes_\mathbf{Z} \kappa(x)^*$ を、
$\lambda \in \kappa(x) \setminus \{0, 1\}$ に対する
$\lambda \otimes (1 - \lambda)$ の形の元で生成される部分群で割った商である。
\item 第一の微分 $\partial$ を次のように定義する。第一項の元
$\xi = \sum_x \alpha_x$ が与えられたとき
$$
\partial(\xi) = \sum\nolimits_{x \leadsto x',\ \delta(x') = k + 1}
\partial_{\mathcal{O}_{W_x, x'}}(\alpha_x)
$$
とおく。ここで
$\partial_{\mathcal{O}_{W_x, x'}} : K_2^M(\kappa(x)) \to K_1^M(\kappa(x))$
は第 \ref{chow-section-tame-symbol} 節で構成された tame 記号である。
\end{enumerate}
これが複体、すなわち $\partial \circ \partial = 0$ であると主張する。
これを見るには、上のような元 $\xi$ と $\delta(x'') = k$ を満たす点
$x'' \in X$ を取り、元 $\partial(\partial(\xi))$ における $x''$ の係数が
零であることを調べれば十分である。$\xi = \sum \alpha_x$ は局所有限和なので、
加法性により、ある $\delta(x) = k + 2$ を満たす $x \in X$ と
$\alpha_x \in K_2^M(\kappa(x))$ に対して $\xi = \alpha_x$ と仮定してよい。
さらに線型性により、ある $f, g \in \kappa(x)^*$ に対して
$\alpha_x = f \otimes g$ と仮定してよい。$W \subset X$ を生成点 $x$ をもつ
整閉部分スキームと表す。$x'' \not \in W$ ならば、
$\partial(\partial(\xi))$ における $x$ の係数が零であることは直ちに明らかである。
$x'' \in W$ ならば、$\partial(\partial(x))$ における $x''$ の係数は
$$
\sum\nolimits_{x \leadsto x' \leadsto x'',\ \delta(x') = k + 1}
\text{ord}_{\mathcal{O}_{\overline{\{x'\}}, x''}}(
\partial_{\mathcal{O}_{W, x'}}(f, g))
$$
に等しい。鍵となる代数的補題 \ref{chow-lemma-milnor-gersten-low-degree} は、
まさにこれが零であると述べている。
\end{remark}

\begin{remark}
\label{chow-remark-higher-chow}
$(S, \delta)$ を状況 \ref{chow-situation-setup} のとおりとする。
$X$ を $S$ 上局所有限型とし、$k \in \mathbf{Z}$ とする。
注意 \ref{chow-remark-higher-chow-pointwise} の複体と注意 \ref{chow-remark-chow-group-pointwise} における $\CH_k(X)$ の表示は、第一高次 Chow 群を
$$
\CH^M_k(X, 1) =
H_1(\text{the complex of Remark \ref{chow-remark-higher-chow-pointwise}})
$$
と定義できることを示唆する。文献で定義されている高次 Chow 群、たとえば
Spencer Bloch の論文（\cite{Bloch} and \cite{Bloch-moving}）のものと区別するため、
上付き文字 ${}^M$ を用いる。$U \subset X$ を開部分集合とし、その補集合
$Y \subset X$ を被約閉部分スキームとみなす。このとき分裂短完全列
$$
0 \to
\bigoplus\nolimits_{y \in Y, \delta(y) = k + i}' K_i^M(\kappa(y)) \to
\bigoplus\nolimits_{x \in X, \delta(x) = k + i}' K_i^M(\kappa(x)) \to
\bigoplus\nolimits_{u \in U, \delta(u) = k + i}' K_i^M(\kappa(u)) \to 0
$$
を得る。これは $i = 2, 1, 0$ に対して注意 \ref{chow-remark-higher-chow-pointwise} の複体の境界写像と両立する。
蛇の補題（『ホモロジー代数』の補題 \ref{homology-lemma-long-exact-sequence-chain}）を適用すると、
補題 \ref{chow-lemma-restrict-to-open} の標準完全列を延長する六項完全列
$$
\CH^M_k(Y, 1) \to \CH^M_k(X, 1) \to \CH^M_k(U, 1) \to
\CH_k(Y) \to \CH_k(X) \to \CH_k(U) \to 0
$$
を得る。さらに作業を加えれば、第一高次 Chow 群 $\CH^M_k(X, 1)$ に対して
平坦引き戻しと固有押し出しも定義できる。この問題には後で戻る
（ここに将来の参照を挿入する）。
\end{remark}






\section{可逆層とのキャップ積と有理同値}
\label{chow-section-commutativity}

\noindent
鍵補題を適用して、可逆層とのキャップ積の基本的性質を得る。特に、
$c_1(\mathcal{L}) \cap -$ が有理同値を経由し、異なる可逆層に対するこれらの作用が
可換であることを見る。

\begin{lemma}
\label{chow-lemma-commutativity-on-integral}
$(S, \delta)$ を状況 \ref{chow-situation-setup} のとおりとする。
$X$ を $S$ 上局所有限型とし、$X$ は整で $\dim_\delta(X) = n$ と仮定する。
$\mathcal{L}$, $\mathcal{N}$ を $X$ 上の可逆層とする。
$\mathcal{L}$ の零でない有理型切断 $s$ と $\mathcal{N}$ の零でない有理型切断 $t$ を選び、
$\alpha = \text{div}_\mathcal{L}(s)$ および $\beta = \text{div}_\mathcal{N}(t)$ とおく。
このとき
$$
c_1(\mathcal{N}) \cap \alpha
=
c_1(\mathcal{L}) \cap \beta
$$
が $\CH_{n - 2}(X)$ において成り立つ。
\end{lemma}

\begin{proof}
鍵補題 \ref{chow-lemma-key-formula} とそれに先立つ議論から直ちに従う。
\end{proof}

\begin{lemma}
\label{chow-lemma-factors}
$(S, \delta)$ を状況 \ref{chow-situation-setup} のとおりとする。
$X$ を $S$ 上局所有限型とし、$\mathcal{L}$ を $X$ 上の可逆層とする。
作用 $\alpha \mapsto c_1(\mathcal{L}) \cap \alpha$ は有理同値を経由して作用
$$
c_1(\mathcal{L}) \cap - : \CH_{k + 1}(X) \to \CH_k(X)
$$
を与える。
\end{lemma}

\begin{proof}
$\alpha \in Z_{k + 1}(X)$ かつ $\alpha \sim_{rat} 0$ とする。
定義 \ref{chow-definition-cap-c1} で定義された $c_1(\mathcal{L}) \cap \alpha$ が零であることを
示せばよい。定義 \ref{chow-definition-rational-equivalence} により、
$\dim_\delta(W_j) = k + 2$ を満たす整閉部分スキームの局所有限族 $\{W_j\}$ と
有理関数 $f_j \in R(W_j)^*$ が存在して
$$
\alpha = \sum (i_j)_*\text{div}_{W_j}(f_j)
$$
となる。$p : \coprod W_j \to X$ は固有射であることに注意する。
したがって $\alpha' \in Z_{k + 1}(\coprod W_j)$ を主因子
$\text{div}_{W_j}(f_j)$ の和とすると、$\alpha = p_*\alpha'$ である。
補題 \ref{chow-lemma-pushforward-cap-c1} により
$c_1(\mathcal{L}) \cap \alpha = p_*(c_1(p^*\mathcal{L}) \cap \alpha')$ である。
ゆえに各 $c_1(\mathcal{L}|_{W_j}) \cap \text{div}_{W_j}(f_j)$ が零であることを
示せば十分である。言い換えると、$X$ は整で、ある $f \in R(X)^*$ に対して
$\alpha = \text{div}_X(f)$ であると仮定してよい。

\medskip\noindent
$X$ は整で、ある $f \in R(X)^*$ に対して $\alpha = \text{div}_X(f)$ と仮定する。
$f$ を可逆層 $\mathcal{N} = \mathcal{O}_X$ の正則有理型切断とみなせる。
$\mathcal{L}$ の有理型切断 $s$ を選び、$\beta = \text{div}_\mathcal{L}(s)$ と表す。
補題 \ref{chow-lemma-commutativity-on-integral} により
$$
c_1(\mathcal{L}) \cap \alpha = c_1(\mathcal{O}_X) \cap \beta.
$$
である。しかし補題 \ref{chow-lemma-c1-cap-additive} により、右辺は所望のとおり
$\CH_k(X)$ において零である。
\end{proof}

\noindent
$(S, \delta)$ を状況 \ref{chow-situation-setup} のとおりとする。
$X$ を $S$ 上局所有限型とし、$\mathcal{L}$ を $X$ 上の可逆層とする。
補題 \ref{chow-lemma-factors} により意味をもつ作用 $c_1(\mathcal{L}) \cap - $ を
$$
c_1(\mathcal{L}) \cap - : \CH_{k + 1}(X) \to \CH_k(X)
$$
と表す。また任意の $s \geq 0$ に対し、この作用の $s$ 回反復を
$c_1(\mathcal{L})^s \cap -$ と表す。

\begin{lemma}
\label{chow-lemma-cap-commutative}
$(S, \delta)$ を状況 \ref{chow-situation-setup} のとおりとする。
$X$ を $S$ 上局所有限型とし、$\mathcal{L}$, $\mathcal{N}$ を $X$ 上の可逆層とする。
任意の $\alpha \in \CH_{k + 2}(X)$ に対して
$$
c_1(\mathcal{L}) \cap c_1(\mathcal{N}) \cap \alpha
=
c_1(\mathcal{N}) \cap c_1(\mathcal{L}) \cap \alpha
$$
が $\CH_k(X)$ の元として成り立つ。
\end{lemma}

\begin{proof}
$\dim_\delta(Z_j) = k + 2$ を満たす整閉部分スキーム $Z_j \subset X$ の
ある局所有限族に対して $\alpha = \sum m_j[Z_j]$ と書く。
固有射 $p : \coprod Z_j \to X$ を考える。$\coprod Z_j$ 上の
$(k + 2)$-サイクルとして $\alpha' = \sum m_j[Z_j]$ とおく。
補題 \ref{chow-lemma-pushforward-cap-c1} を数回適用すると
$c_1(\mathcal{L}) \cap c_1(\mathcal{N}) \cap \alpha
= p_*(c_1(p^*\mathcal{L}) \cap c_1(p^*\mathcal{N}) \cap \alpha')$
および
$c_1(\mathcal{N}) \cap c_1(\mathcal{L}) \cap \alpha
= p_*(c_1(p^*\mathcal{N}) \cap c_1(p^*\mathcal{L}) \cap \alpha')$
を得る。ゆえに $X$ が整で $\alpha = [X]$ の場合に公式を示せば十分である。
この場合、結果は補題 \ref{chow-lemma-commutativity-on-integral} と定義から従う。
\end{proof}






\section{Gysin 準同型}
\label{chow-section-intersecting-effective-Cartier}

\noindent
本節では、可逆層の切断の零点スキーム $D$ に対する Gysin 写像を定義する。
$D$ が有効 Cartier 因子である場合が重要だが、任意の $D$ へ一般化しておくことで、
注意 \ref{chow-remark-pullback-pairs} および補題 \ref{chow-lemma-closed-in-X-gysin},
\ref{chow-lemma-gysin-flat-pullback}, \ref{chow-lemma-gysin-commutes-gysin} に見られる
種々の整合性を柔軟に定式化できる。これらの結果は、仮想正規束を備えた局所主閉部分スキーム
（注意 \ref{chow-remark-generalize-to-virtual}）や擬因子
（注意 \ref{chow-remark-generalize-to-pseudo-divisor}）にも一般化できる。

\medskip\noindent
有効 Cartier 因子は、$\mathcal{L}$ が可逆層で $s$ が正則な大域切断であるような対
$(\mathcal{L}, s)$ の同型類と $1$-to-$1$ に対応することを思い出そう。
「因子」補題 \ref{divisors-lemma-characterize-OD} を参照されたい。
$D$ が $(\mathcal{L}, s)$ に対応するならば、$\mathcal{L} = \mathcal{O}_X(D)$ である。
本節を読む際にはこの点を念頭に置く。

\begin{definition}
\label{chow-definition-gysin-homomorphism}
$(S, \delta)$ を状況 \ref{chow-situation-setup} のとおりとする。
$X$ を $S$ 上局所有限型とする。$(\mathcal{L}, s)$ を、可逆層と大域切断
$s \in \Gamma(X, \mathcal{L})$ からなる対とする。$D = Z(s)$ を $s$ の零点スキームとし、
$i : D \to X$ を閉埋め込みと表す。各整数 $k$ に対し、{\it Gysin 準同型}
$$
i^* : Z_{k + 1}(X) \to \CH_k(D).
$$
を次の規則によって定義する。
\begin{enumerate}
\item $\dim_\delta(W) = k + 1$ を満たす整閉部分スキーム $W \subset X$ に対しては、
\begin{enumerate}
\item $W \not \subset D$ ならば、$D$ 上の $k$-サイクルとして
$i^*[W] = [D \cap W]_k$ と定める。
\item $W \subset D$ ならば、誘導される閉埋め込み $i' : W \to D$ を用いて
$i^*[W] = i'_*(c_1(\mathcal{L}|_W) \cap [W])$ と定める。
\end{enumerate}
\item 一般の $(k + 1)$-サイクル $\alpha = \sum n_j[W_j]$ に対しては
$$
i^*\alpha = \sum n_j i^*[W_j]
$$
とおく。
\item $D$ が有効 Cartier 因子ならば、類 $i^*\alpha$ を $X$ 上の類へ押し出したものを
$D \cdot \alpha = i_*i^*\alpha$ と表す。
\end{enumerate}
\end{definition}

\noindent
実際、後で見るように、この Gysin 準同型 $i^*$ は非平坦な引き戻しの一例とみなせる。
したがって、類 $i^*\alpha$ を類 $\alpha$ の {\it 引き戻し} と非形式的に呼ぶことがある。

\begin{remark}
\label{chow-remark-generalize-to-virtual}
$X$ を状況 \ref{chow-situation-setup} のとおり $S$ 上局所有限型なスキームとする。
$(D, \mathcal{N}, \sigma)$ を、局所主な（「因子」定義
\ref{divisors-definition-effective-Cartier-divisor}）閉部分スキーム
$i : D \to X$、可逆 $\mathcal{O}_D$-加群 $\mathcal{N}$、および $\mathcal{O}_D$-加群の全射
$\sigma : \mathcal{N}^{\otimes -1} \to i^*\mathcal{I}_D$
からなる三つ組とする\footnote{この条件により、$D$ が有効 Cartier 因子ならば
$\mathcal{N} = \mathcal{O}_X(D)|_D$ となることが保証される。}。
ここで $\mathcal{N}$ は {\it $X$ における $D$ の仮想正規束} と考えるべきものである。
定義 \ref{chow-definition-gysin-homomorphism} の構成
$i^* : Z_{k + 1}(X) \to \CH_k(D)$ はこのような三つ組へ一般化される。
節 \ref{chow-section-gysin-higher-codimension} を参照されたい。
\end{remark}

\begin{remark}
\label{chow-remark-generalize-to-pseudo-divisor}
$X$ を状況 \ref{chow-situation-setup} のとおり $S$ 上局所有限型なスキームとする。
\cite{F} では $X$ 上の {\it 擬因子} を、$\mathcal{L}$ が可逆 $\mathcal{O}_X$-加群、
$Z \subset X$ が閉部分集合、$s \in \Gamma(X \setminus Z, \mathcal{L})$ が
至る所消えない切断であるような三つ組 $D = (\mathcal{L}, Z, s)$ と定義している。
上と同様に、$\CH_{k + 1}(X)$ の各 $\alpha$ に対し、
$|\alpha|$ を $\alpha$ の台として、$\CH_k(Z \cap |\alpha|)$ の積
$D \cdot \alpha$ を定義できる。
\end{remark}

\begin{lemma}
\label{chow-lemma-support-cap-effective-Cartier}
$(S, \delta)$ を状況 \ref{chow-situation-setup} のとおりとし、$X$ を $S$ 上局所有限型とする。
$(\mathcal{L}, s, i : D \to X)$ を定義 \ref{chow-definition-gysin-homomorphism} のとおりとし、
$\alpha$ を $X$ 上の $(k + 1)$-サイクルとする。このとき $\CH_k(X)$ において
$i_*i^*\alpha = c_1(\mathcal{L}) \cap \alpha$ である。特に $D$ が有効 Cartier 因子ならば、
$D \cdot \alpha = c_1(\mathcal{O}_X(D)) \cap \alpha$ である。
\end{lemma}

\begin{proof}
$i_j : W_j \to X$ が $\dim_\delta(W_j) = k$ を満たす整閉部分スキームであるように、
$\alpha = \sum n_j[W_j]$ と書く。$D$ は $s$ の零点スキームなので、
$D \cap W_j$ は制限 $s|_{W_j}$ の零点スキームである。したがって
$W_j \not \subset D$ を満たす各 $j$ に対し、補題 \ref{chow-lemma-geometric-cap} により
$c_1(\mathcal{L}) \cap [W_j] = [D \cap W_j]_k$ である。ゆえに
$$
c_1(\mathcal{L}) \cap \alpha
=
\sum\nolimits_{W_j \not \subset D} n_j[D \cap W_j]_k
+
\sum\nolimits_{W_j \subset D}
n_j i_{j, *}(c_1(\mathcal{L})|_{W_j}) \cap [W_j])
$$
を得る。定義 \ref{chow-definition-cap-c1} により、これは $\CH_k(X)$ における等式である。
右辺は、定義 \ref{chow-definition-gysin-homomorphism} による $D$ 上の類 $i^*\alpha$ の
押し出しと項ごとに一致する。したがって証明された。
\end{proof}

\begin{lemma}
\label{chow-lemma-easy-gysin}
$(S, \delta)$ を状況 \ref{chow-situation-setup} のとおりとし、$X$ を $S$ 上局所有限型とする。
$(\mathcal{L}, s, i : D \to X)$ を定義 \ref{chow-definition-gysin-homomorphism} のとおりとする。
\begin{enumerate}
\item $Z \subset X$ を、$\dim_\delta(Z) \leq k + 1$ かつ $D \cap Z$ が
$Z$ 上の有効 Cartier 因子であるような閉部分スキームとする。このとき
$i^*[Z]_{k + 1} = [D \cap Z]_k$ である。
\item $\mathcal{F}$ を、$\dim_\delta(\text{Supp}(\mathcal{F})) \leq k + 1$ かつ
$s : \mathcal{F} \to \mathcal{F} \otimes_{\mathcal{O}_X} \mathcal{L}$ が単射であるような
$X$ 上の連接層とする。このとき $\CH_k(D)$ において
$$
i^*[\mathcal{F}]_{k + 1} = [i^*\mathcal{F}]_k
$$
である。
\end{enumerate}
\end{lemma}

\begin{proof}
(1) のとおりの $Z \subset X$ を仮定し、$\mathcal{F} = \mathcal{O}_Z$ とおく。
$D \cap Z$ が有効 Cartier 因子であるという仮定は、
$s : \mathcal{F} \to \mathcal{F} \otimes_{\mathcal{O}_X} \mathcal{L}$ が単射であるという仮定と
同値である。さらに $[Z]_{k + 1} = [\mathcal{F}]_{k + 1}]$ かつ
$[D \cap Z]_k = [\mathcal{O}_{D \cap Z}]_k = [i^*\mathcal{F}]_k$ である。
補題 \ref{chow-lemma-cycle-closed-coherent} を参照されたい。したがって (1) は (2) から従う。

\medskip\noindent
$m_j > 0$ で、$W_j \subset X$ が互いに異なる $\delta$-次元 $k + 1$ の整閉部分スキームと
なるように $[\mathcal{F}]_{k + 1} = \sum m_j[W_j]$ と書く。
$s : \mathcal{F} \to \mathcal{F} \otimes_{\mathcal{O}_X} \mathcal{L}$ が単射であるという仮定から、
すべての $j$ について $W_j \not \subset D$ である。定義により
$$
i^*[\mathcal{F}]_{k + 1} = \sum m_j [D \cap W_j]_k.
$$
を得る。サイクルとして
$$
\sum [D \cap W_j]_k = [i^*\mathcal{F}]_k
$$
であると主張する。$Z \subset D$ を $\delta$-次元 $k$ の整閉部分スキームとし、
$\xi \in Z$ をその生成点とする。$A = \mathcal{O}_{X, \xi}$、$M = \mathcal{F}_\xi$ とおく。
$f \in A$ を $D$ のイデアルを生成する元、すなわち
$\mathcal{O}_{D, \xi} = A/fA$ となる元とする。仮定により
$\dim(\text{Supp}(M)) = 1$、写像 $f : M \to M$ は単射、かつ
$\text{length}_A(M/fM) < \infty$ である。さらに $\text{length}_A(M/fM)$ は
$[i^*\mathcal{F}]_k$ における $[Z]$ の係数である。他方、
$\mathfrak q_1, \ldots, \mathfrak q_t$ を $M$ の台における極小素イデアルとする。このとき
$$
\sum
\text{length}_{A_{\mathfrak q_i}}(M_{\mathfrak q_i})
\text{ord}_{A/\mathfrak q_i}(f)
$$
は $\sum [D \cap W_j]_k$ における $[Z]$ の係数である。
したがって補題 \ref{chow-lemma-additivity-divisors-restricted} により等式を得る。
\end{proof}

\begin{remark}
\label{chow-remark-gysin-on-cycles}
$X \to S$, $\mathcal{L}$, $s$, $i : D \to X$ を定義
\ref{chow-definition-gysin-homomorphism} のとおりとし、
$\mathcal{L}|_D \cong \mathcal{O}_D$ と仮定する。この場合、
$W \subset D$ が整閉部分スキームならば $i^*[W] = 0$ とすることにより、
サイクル上の標準的写像 $i^* : Z_{k + 1}(X) \to Z_k(D)$ を定義できる。
この事実は後で有用になる。
\end{remark}

\begin{remark}
\label{chow-remark-pullback-pairs}
$f : X' \to X$ を状況 \ref{chow-situation-setup} のとおり $S$ 上局所有限型なスキームの射とし、
$(\mathcal{L}, s, i : D \to X)$ を定義 \ref{chow-definition-gysin-homomorphism} のとおりの
三つ組とする。このとき $\mathcal{L}' = f^*\mathcal{L}$、$s' = f^*s$、
$D' = X' \times_X D = Z(s')$ とおける。これにより可換図式
$$
\xymatrix{
D' \ar[d]_g \ar[r]_{i'} & X' \ar[d]^f \\
D \ar[r]^i & X
}
$$
が得られ、$i^*$ と $(i')^*$ の間の種々の整合性を問うことができる。
\end{remark}

\begin{lemma}
\label{chow-lemma-closed-in-X-gysin}
$(S, \delta)$ を状況 \ref{chow-situation-setup} のとおりとする。
$f : X' \to X$ を $S$ 上局所有限型なスキームの固有射とし、
$(\mathcal{L}, s, i : D \to X)$ を定義 \ref{chow-definition-gysin-homomorphism} のとおりとする。
注意 \ref{chow-remark-pullback-pairs} のとおり図式
$$
\xymatrix{
D' \ar[d]_g \ar[r]_{i'} & X' \ar[d]^f \\
D \ar[r]^i & X
}
$$
を作る。$X'$ 上の任意の $(k + 1)$-サイクル $\alpha'$ に対し、$\CH_k(D)$ において
$i^*f_*\alpha' = g_*(i')^*\alpha'$ である
（$f_*$ はサイクルのレベルで定義されるので、この式には意味がある）。
\end{lemma}

\begin{proof}
ある整閉部分スキーム $W' \subset X'$ に対して $\alpha = [W']$ と仮定し、
$W = f(W') \subset X$ とおく。$W' \not \subset D'$ の場合は $W \not \subset D$ であり、
$$
[W' \cap D']_k = \text{div}_{\mathcal{L}'|_{W'}}({s'|_{W'}})
\quad\text{and}\quad
[W \cap D]_k = \text{div}_{\mathcal{L}|_W}(s|_W)
$$
である。したがって補題 \ref{chow-lemma-equal-c1-as-cycles} により、第一のサイクルの $f_*$ は
第二のサイクルに等しく、等式はサイクルとして成り立つ。$W' \subset D'$ の場合は
$W \subset D$ であり、補題 \ref{chow-lemma-equal-c1-as-cycles} の第二の主張により、
$f_*(c_1(\mathcal{L}|_{W'}) \cap [W'])$ は $\CH_k(W)$ において
$c_1(\mathcal{L}|_W) \cap [W]$ に等しい。注意
\ref{chow-remark-infinite-sums-rational-equivalences} により、一般の $\alpha'$ についても結論が従う。
\end{proof}

\begin{lemma}
\label{chow-lemma-gysin-flat-pullback}
$(S, \delta)$ を状況 \ref{chow-situation-setup} のとおりとする。
$f : X' \to X$ を、$S$ 上局所有限型なスキームの相対次元 $r$ の平坦射とする。
$(\mathcal{L}, s, i : D \to X)$ を定義 \ref{chow-definition-gysin-homomorphism} のとおりとし、
注意 \ref{chow-remark-pullback-pairs} のとおり図式
$$
\xymatrix{
D' \ar[d]_g \ar[r]_{i'} & X' \ar[d]^f \\
D \ar[r]^i & X
}
$$
を作る。$X$ 上の任意の $(k + 1)$-サイクル $\alpha$ に対し、
$\CH_{k + r}(D')$ において $(i')^*f^*\alpha = g^*i^*\alpha$ である
（$f^*$ はサイクルのレベルで定義されるので、この式には意味がある）。
\end{lemma}

\begin{proof}
ある整閉部分スキーム $W \subset X$ に対して $\alpha = [W]$ と仮定し、
$W' = f^{-1}(W) \subset X'$ とおく。$W \not \subset D$ の場合は
$W' \not \subset D'$ であり、$D'$ の閉部分スキームとして
$$
W' \cap D' = g^{-1}(W \cap D)
$$
である。したがって補題 \ref{chow-lemma-pullback-coherent} により、等式はサイクルとして成り立つ。
$W \subset D$ の場合は $W' \subset D'$ かつ $W' = g^{-1}(W)$ であり、
$[W']_{k + 1 + r} = g^*[W]$ である。補題 \ref{chow-lemma-flat-pullback-cap-c1} により、
等式は $\CH_{k + r}(D')$ において成り立つ。注意
\ref{chow-remark-infinite-sums-rational-equivalences} により、一般の $\alpha'$ についても結論が従う。
\end{proof}








\section{Gysin 準同型と有理同値}
\label{chow-section-gysin}

\noindent
本節では鍵となる公式を用い、Gysin 準同型が有理同値を通して定まることを示す。
さらに、重要な可換性も証明する。

\begin{lemma}
\label{chow-lemma-gysin-factors-general}
$(S, \delta)$ を状況 \ref{chow-situation-setup} のとおりとする。
$X$ を $S$ 上局所有限型とする。さらに $X$ は整であるとし、
$n = \dim_\delta(X)$ とおく。
$i : D \to X$ を有効 Cartier 因子とする。$\mathcal{N}$ を可逆 $\mathcal{O}_X$-加群、
$t$ を $\mathcal{N}$ の零でない有理型切断とする。このとき $\CH_{n - 2}(D)$ において
$i^*\text{div}_\mathcal{N}(t) = c_1(\mathcal{N}|_D) \cap [D]_{n - 1}$ である。
\end{lemma}

\begin{proof}
$\delta$-次元 $n - 1$ のある整閉部分スキーム $Z_i \subset X$ を用いて
$\text{div}_\mathcal{N}(t) = \sum \text{ord}_{Z_i, \mathcal{N}}(t)[Z_i]$ と書く。
族 $\{Z_i\}$ は局所有限であり、$U = X \setminus \bigcup Z_i$ 上で
$t \in \Gamma(U, \mathcal{N}|_U)$ は生成元であり、さらに $D$ の各既約成分がいずれかの
$Z_i$ であると仮定してよい。「因子」補題
\ref{divisors-lemma-components-locally-finite},
\ref{divisors-lemma-divisor-locally-finite},
\ref{divisors-lemma-divisor-meromorphic-locally-finite} を参照されたい。

\medskip\noindent
$\mathcal{L} = \mathcal{O}_X(D)$ とおき、標準切断を
$s \in \Gamma(X, \mathcal{O}_X(D)) = \Gamma(X, \mathcal{L})$ と表す。
現在の状況に節 \ref{chow-section-key} の議論を適用する。各 $i$ に対し、
$\xi_i \in Z_i$ をその生成点、$B_i = \mathcal{O}_{X, \xi_i}$ とする。
各 $i$ について $B_i$ 上の生成元 $s_i \in \mathcal{L}_{\xi_i}$ と
$t_i \in \mathcal{N}_{\xi_i}$ を選ぶが、$Z_i \not \subset D$ の場合には
$s_i = s$ を選ぶものとする。$f_i, g_i \in B_i$ によって
$s = f_i s_i$、$t = g_i t_i$ と書く。このとき
$\text{ord}_{Z_i, \mathcal{N}}(t) = \text{ord}_{B_i}(g_i)$ である。
他方、$f_i \in B_i$ であり、$s_i$ の選び方から
$$
[D]_{n - 1} = \sum \text{ord}_{B_i}(f_i)[Z_i]
$$
である。サイクルとして
$$
i^*\text{div}_\mathcal{N}(t) =
\sum \text{ord}_{B_i}(g_i) \text{div}_{\mathcal{L}|_{Z_i}}(s_i|_{Z_i})
$$
であると主張する。より正確には、右辺が左辺を表すサイクルである。
実際、これは $\text{div}_\mathcal{N}(t)$ の上の式と、$Z_i \not \subset D$ の場合には
$s_i|_{Z_i} = s|_{Z_i}$ が正則切断であるため
$\text{div}_{\mathcal{L}|_{Z_i}}(s_i|_{Z_i}) = [Z(s_i|_{Z_i})]_{n - 2} =
[Z_i \cap D]_{n - 2}$ となることから明らかである。補題 \ref{chow-lemma-compute-c1} を参照されたい。
同様に
$$
c_1(\mathcal{N}) \cap [D]_{n - 1} =
\sum \text{ord}_{B_i}(f_i) \text{div}_{\mathcal{N}|_{Z_i}}(t_i|_{Z_i})
$$
である。鍵となる公式（補題 \ref{chow-lemma-key-formula}）により、サイクルの等式
$$
\sum \left(
\text{ord}_{B_i}(f_i) \text{div}_{\mathcal{N}|_{Z_i}}(t_i|_{Z_i}) -
\text{ord}_{B_i}(g_i) \text{div}_{\mathcal{L}|_{Z_i}}(s_i|_{Z_i}) \right) =
\sum \text{div}_{Z_i}(\partial_{B_i}(f_i, g_i))
$$
を得る。$Z_i \not \subset D$ ならば $f_i = 1$ であり、したがって
$\text{div}_{Z_i}(\partial_{B_i}(f_i, g_i)) = 0$ である。ゆえに $D$ 上で、
$i^*\text{div}_\mathcal{N}(t)$ と $c_1(\mathcal{N}) \cap [D]_{n - 1}$ を表す
これらの具体的サイクルの間に有理同値が得られる。これで証明が完了する。
\end{proof}

\begin{lemma}
\label{chow-lemma-gysin-factors}
$(S, \delta)$ を状況 \ref{chow-situation-setup} のとおりとし、$X$ を $S$ 上局所有限型とする。
$(\mathcal{L}, s, i : D \to X)$ を定義 \ref{chow-definition-gysin-homomorphism} のとおりとする。
Gysin 準同型は有理同値を通して写像 $i^* : \CH_{k + 1}(X) \to \CH_k(D)$ を与える。
\end{lemma}

\begin{proof}
$\alpha \in Z_{k + 1}(X)$ とし、$\alpha \sim_{rat} 0$ と仮定する。
これは、$\delta$-次元 $k + 2$ の整閉部分スキーム $W_j \subset X$ の局所有限族と
$f_j \in R(W_j)^*$ が存在して
$\alpha = \sum i_{j, *}\text{div}_{W_j}(f_j)$ となることを意味する。
$X' = \coprod W_i$ とおき、注意 \ref{chow-remark-pullback-pairs} の図式
$$
\xymatrix{
D' \ar[d]_q \ar[r]_{i'} & X' \ar[d]^p \\
D \ar[r]^i & X
}
$$
を考える。$X' \to X$ は固有なので、補題 \ref{chow-lemma-closed-in-X-gysin} により
$i^*p_* = q_*(i')^*$ である。$q_*$ は有理同値を通して定まることが分かっているので
（補題 \ref{chow-lemma-proper-pushforward-rational-equivalence}）、$X'$ 上の
$\alpha' = \sum \text{div}_{W_j}(f_j)$ について示せば十分である。
明らかに、$X$ が整で、ある $f \in R(X)^*$ に対し $\alpha = \text{div}(f)$ である場合へ帰着する。

\medskip\noindent
$X$ は整で、ある $f \in R(X)^*$ に対して $\alpha = \text{div}(f)$ であると仮定する。
$X = D$ ならば、$i^*\alpha$ は $c_1(\mathcal{L}) \cap \alpha$ に等しい。
これは補題 \ref{chow-lemma-factors} により零と有理同値である。
$D \not = X$ ならば、補題 \ref{chow-lemma-gysin-factors-general} により
$i^*\text{div}_X(f)$ は $\CH_{n - 2}(D)$ において
$c_1(\mathcal{O}_D) \cap [D]_{n - 1}$ に等しい。当然、$c_1(\mathcal{O}_D)$ との
キャップ積は零写像である（補題 \ref{chow-lemma-c1-cap-additive}）。
\end{proof}

\begin{lemma}
\label{chow-lemma-gysin-back}
$(S, \delta)$ を状況 \ref{chow-situation-setup} のとおりとし、$X$ を $S$ 上局所有限型とする。
$(\mathcal{L}, s, i : D \to X)$ を定義 \ref{chow-definition-gysin-homomorphism} のとおりとする。
このとき $i^*i_* : \CH_k(D) \to \CH_{k - 1}(D)$ は $\alpha$ を
$c_1(\mathcal{L}|_D) \cap \alpha$ へ写す。
\end{lemma}

\begin{proof}
これはサイクル上の $i_*$ の定義と、定義 \ref{chow-definition-gysin-homomorphism} で与えた
$i^*$ の定義から直ちに従う。
\end{proof}

\begin{lemma}
\label{chow-lemma-gysin-commutes-cap-c1}
$(S, \delta)$ を状況 \ref{chow-situation-setup} のとおりとし、$X$ を $S$ 上局所有限型とする。
$(\mathcal{L}, s, i : D \to X)$ を定義 \ref{chow-definition-gysin-homomorphism} のとおりの
三つ組とし、$\mathcal{N}$ を可逆 $\mathcal{O}_X$-加群とする。このとき、任意の
$\alpha \in \CH_k(X)$ に対して $\CH_{k - 2}(D)$ において
$i^*(c_1(\mathcal{N}) \cap \alpha) = c_1(i^*\mathcal{N}) \cap i^*\alpha$ である。
\end{lemma}

\begin{proof}
補題 \ref{chow-lemma-gysin-factors} と全く同じ証明により、これは補題
\ref{chow-lemma-pushforward-cap-c1}, \ref{chow-lemma-cap-commutative},
\ref{chow-lemma-gysin-factors-general} から従う。
\end{proof}

\begin{lemma}
\label{chow-lemma-gysin-commutes-gysin}
$(S, \delta)$ を状況 \ref{chow-situation-setup} のとおりとし、$X$ を $S$ 上局所有限型とする。
$(\mathcal{L}, s, i : D \to X)$ と $(\mathcal{L}', s', i' : D' \to X)$ を
定義 \ref{chow-definition-gysin-homomorphism} のとおりの二つの三つ組とする。このとき図式
$$
\xymatrix{
\CH_k(X) \ar[r]_{i^*} \ar[d]_{(i')^*} & \CH_{k - 1}(D) \ar[d]^{j^*} \\
\CH_{k - 1}(D') \ar[r]^{(j')^*} & \CH_{k - 2}(D \cap D')
}
$$
は可換であり、ここに現れる写像はすべて Gysin 写像である。
\end{lemma}

\begin{proof}
$j : D \cap D' \to D$ および $j' : D \cap D' \to D'$ で、
$(\mathcal{L}|_{D'}, s|_{D'}$ と $(\mathcal{L}'_D, s|_D)$ に対応する閉埋め込みを表す。
すべての $\alpha \in \CH_k(X)$ に対して $(j')^*i^*\alpha = j^* (i')^*\alpha$ を示せばよい。
$W \subset X$ を次元 $k$ の整閉部分スキームとする。
$\alpha = [W]$ の場合に等式を証明し、それを鍵となる公式から導く。

\medskip\noindent
$\sigma$ を $\mathcal{L}|_W$ の零でない有理型切断とし、$W \not \subset D$ の場合には
$s|_W$ に等しく選ぶ。$\sigma'$ を $\mathcal{L}'|_W$ の零でない有理型切断とし、
$W \not \subset D'$ の場合には $s'|_W$ に等しく選ぶ。節 \ref{chow-section-key} の議論のとおり
$$
\text{div}_{\mathcal{L}|_W}(\sigma) =
\sum \text{ord}_{Z_i, \mathcal{L}|_W}(\sigma)[Z_i] = \sum n_i[Z_i]
$$
および同様に
$$
\text{div}_{\mathcal{L}'|_W}(\sigma') =
\sum \text{ord}_{Z_i, \mathcal{L}'|_W}(\sigma')[Z_i] = \sum n'_i[Z_i]
$$
と書く。このとき $n_i \not = 0$ ならば $Z_i \subset D$ であり、
$n'_i \not = 0$ ならば $Z'_i \subset D'$ である。各 $i$ について
$\xi_i \in Z_i$ を生成点とする。節 \ref{chow-section-key} と同様、各 $i$ に対して
$B_i = \mathcal{O}_{W, \xi_i}$ 上の生成元として、
$\sigma_i \in \mathcal{L}_{\xi_i}$ と、それぞれ\ $\sigma'_i \in \mathcal{L}'_{\xi_i}$ を選ぶ。
ただし\ 前者は $Z_i \not \subset D$ の場合には $s$ の像に、それぞれ\ 後者は
$Z_i \not \subset D'$ の場合には $s'$ の像に等しく選ぶ。
$\sigma = f_i \sigma_i$ および $\sigma' = f'_i\sigma'_i$ と書けば、
$n_i = \text{ord}_{B_i}(f_i)$ かつ $n'_i = \text{ord}_{B_i}(f'_i)$ である。
定義から、サイクルとして
$$
(j')^*i^*[W] =
\sum \text{ord}_{B_i}(f_i) \text{div}_{\mathcal{L}'|_{Z_i}}(\sigma'_i|_{Z_i})
$$
であり、また
$$
j^*(i')^*[W] =
\sum \text{ord}_{B_i}(f'_i) \text{div}_{\mathcal{L}|_{Z_i}}(\sigma_i|_{Z_i})
$$
である。ここで鍵となる公式（補題 \ref{chow-lemma-key-formula}）により、サイクルの等式
$$
\sum \left(
\text{ord}_{B_i}(f_i) \text{div}_{\mathcal{L}'|_{Z_i}}(\sigma'_i|_{Z_i}) -
\text{ord}_{B_i}(f'_i) \text{div}_{\mathcal{L}|_{Z_i}}(\sigma_i|_{Z_i})
\right) =
\sum \text{div}_{Z_i}(\partial_{B_i}(f_i, f'_i))
$$
を得る。$Z_i \not \subset D \cap D'$ ならば $f_i = 1$ または $f'_i = 1$ なので、
$\text{div}_{Z_i}(\partial_{B_i}(f_i, f'_i)) = 0$ である。
したがって $D \cap D' \cap W$ 上で、$(j')^*i^*[W]$ と $j^*(i')^*[W]$ を表す
これらの具体的サイクルの間に有理同値が得られる。注意
\ref{chow-remark-infinite-sums-rational-equivalences} により、一般の $\alpha$ についても結論が従う。
\end{proof}
















\section{相対有効 Cartier 因子}
\label{chow-section-relative-effective-cartier}

\noindent
相対有効 Cartier 因子は「因子」節 \ref{divisors-section-effective-Cartier-morphisms} で
定義され、研究されている。ベクトル束の Chern 類に関する基本結果を展開するためには、
周囲のスキームと有効 Cartier 因子の双方が基底上平坦である場合だけが必要である。

\begin{lemma}
\label{chow-lemma-relative-effective-cartier}
$(S, \delta)$ を状況 \ref{chow-situation-setup} のとおりとする。
$X$, $Y$ を $S$ 上局所有限型とし、$p : X \to Y$ を相対次元 $r$ の平坦射とする。
$i : D \to X$ を相対有効 Cartier 因子とする
（「因子」定義 \ref{divisors-definition-relative-effective-Cartier-divisor}）。
$\mathcal{L} = \mathcal{O}_X(D)$ とおく。任意の $\alpha \in \CH_{k + 1}(Y)$ に対し
$$
i^*p^*\alpha = (p|_D)^*\alpha
$$
が $\CH_{k + r}(D)$ で成り立ち、また
$$
c_1(\mathcal{L}) \cap p^*\alpha = i_* ((p|_D)^*\alpha)
$$
が $\CH_{k + r}(X)$ で成り立つ。
\end{lemma}

\begin{proof}
$W \subset Y$ を $\delta$-次元 $k + 1$ の整閉部分スキームとする。
「因子」補題 \ref{divisors-lemma-relative-Cartier} により、
$D \cap p^{-1}W$ は $p^{-1}W$ 上の有効 Cartier 因子である。
補題 \ref{chow-lemma-easy-gysin} により
$$
i^*[p^{-1}W]_{k + r + 1} =
[D \cap p^{-1}W]_{k + r} =
[(p|_D)^{-1}(W)]_{k + r}.
$$
の第一の等式を得る。第二の等式は、スキームとして
$D \cap p^{-1}(W) = (p|_D)^{-1}(W)$ であることから従う。
定義により $p^*[W] = [p^{-1}W]_{k + r + 1}$ なので、サイクルとして
$i^*p^*[W] = (p|_D)^*[W]$ である。$\alpha = \sum m_j[W_j]$ が一般の
$k + 1$-サイクルならば、サイクルとして
$i^*\alpha = \sum m_j i^*p^*[W_j] = \sum m_j(p|_D)^*[W_j]$ を得る。
これで第一の等式が示された。第二の等式は、第一の等式に補題
\ref{chow-lemma-support-cap-effective-Cartier} を適用すれば従う。
\end{proof}








\section{アフィン束}
\label{chow-section-affine-vector}

\noindent
アフィン束に対して、Chow 群上の引き戻し写像は全射である。

\begin{lemma}
\label{chow-lemma-pullback-affine-fibres-surjective}
$(S, \delta)$ を状況 \ref{chow-situation-setup} のとおりとする。
$X$, $Y$ を $S$ 上局所有限型とし、$f : X \to Y$ を相対次元 $r$ の平坦射とする。
各 $y \in Y$ に対し、開近傍 $U \subset Y$ が存在して
$f|_{f^{-1}(U)} : f^{-1}(U) \to U$ が射 $U \times \mathbf{A}^r \to U$ と
同一視されると仮定する。このとき、すべての $k \in \mathbf{Z}$ に対して
$f^* : \CH_k(Y) \to \CH_{k + r}(X)$ は全射である。
\end{lemma}

\begin{proof}
$\alpha \in \CH_{k + r}(X)$ とする。$m_j \not = 0$ で、$W_j$ が互いに異なる
$\delta$-次元 $k + r$ の整閉部分スキームとなるように
$\alpha = \sum m_j[W_j]$ と書く。このとき族 $\{W_j\}$ は $X$ で局所有限である。
任意の準コンパクト開集合 $V \subset Y$ に対し、$f^{-1}(V) \cap W_j$ が空でないのは
有限個の $j$ に限られる。したがって、像の閉包の族
$Z_j = \overline{f(W_j)}$ は $Y$ の整閉部分スキームの局所有限族である。

\medskip\noindent
ファイバー積図式
$$
\xymatrix{
f^{-1}(Z_j) \ar[r] \ar[d]_{f_j} & X \ar[d]^f \\
Z_j \ar[r] & Y
}
$$
を考える。ある $k$-サイクル $\beta_j \in \CH_k(Z_j)$ に対して
$[W_j] \in Z_{k + r}(f^{-1}(Z_j))$ が $f_j^*\beta_j$ と有理同値であると仮定する。
このとき $\beta = \sum m_j \beta_j$ は $Y$ 上の $k$-サイクルであり、
$f^*\beta = \sum m_j f_j^*\beta_j$ は $\alpha$ と有理同値になる
（注意 \ref{chow-remark-infinite-sums-rational-equivalences} を参照）。
これにより、$Y$ は整で、$\alpha = [W]$ が $Y$ を支配する $X$ のある整閉部分スキームに
よって与えられる場合へ帰着する。特に $d = \dim_\delta(Y) < \infty$ と仮定してよい。

\medskip\noindent
そこで $d = \dim_\delta(Y)$ に関する帰納法を用いる。
$d < k$ ならば $\CH_{k + r}(X) = 0$ なので補題は成り立つ。
仮定により、$V$ 上のスキームとして $f^{-1}(V) \cong V \times \mathbf{A}^r$ となる
稠密開集合 $V \subset Y$ が存在する。ある $\beta \in Z_k(V)$ に対して
$\alpha|_{f^{-1}(V)} = f^*\beta$ を示せたと仮定する。
補題 \ref{chow-lemma-exact-sequence-open} により、ある $\beta' \in Z_k(Y)$ に対して
$\beta = \beta'|_V$ である。補題 \ref{chow-lemma-restrict-to-open} の完全列
$\CH_k(f^{-1}(Y \setminus V)) \to \CH_k(X) \to \CH_k(f^{-1}(V))$ により、
$\alpha - f^*\beta'$ はあるサイクル
$\alpha' \in \CH_{k + r}(f^{-1}(Y \setminus V))$ から来る。
$\dim_\delta(Y \setminus V) < d$ なので、$d$ に関する帰納法により結論を得る。

\medskip\noindent
したがって $X = Y \times \mathbf{A}^r$ と仮定してよい。この場合 $f$ を
$$
X = Y \times \mathbf{A}^r \to
Y \times \mathbf{A}^{r - 1} \to \ldots \to
Y \times \mathbf{A}^1 \to Y.
$$
と分解できる。ゆえに $r = 1$ の場合を扱えば十分である。
証明の第二段落の議論により、$\alpha = [W]$、$Y$ は整、かつ $W \to Y$ は支配的である場合へ
帰着する。再び $d = \dim_\delta(Y)$ に関する帰納法を用いる。
$W = Y \times \mathbf{A}^1$ ならば $[W] = f^*[Y]$ である。
最後に $W \subset Y \times \mathbf{A}^1$ が真の包含ならば、$W \to Y$ は有限次体拡大
$R(W)/R(Y)$ を誘導する。$W \to Y$ の生成ファイバーが
$\mathbf{A}^1_{R(Y)}$ において $P$ で切り出されるようなモニック既約多項式
$P(T) \in R(Y)[T]$ をとる。$P \in \Gamma(V, \mathcal{O}_Y)[T]$ であり、かつ
$W \cap f^{-1}(V)$ がなお $P$ で切り出されるような空でない開集合 $V \subset Y$ をとる。
すると $\alpha|_{f^{-1}(V)} \sim_{rat} 0$ であり、したがって
$(Y \setminus V) \times \mathbf{A}^1$ 上のあるサイクル $\alpha'$ に対して
$\alpha \sim_{rat} \alpha'$ である。次元に関する帰納法によって結論を得る。
\end{proof}

\begin{lemma}
\label{chow-lemma-linebundle}
$(S, \delta)$ を状況 \ref{chow-situation-setup} のとおりとし、$X$ を $S$ 上局所有限型とする。
$\mathcal{L}$ を可逆 $\mathcal{O}_X$-加群とし、
$$
p :
L = \underline{\Spec}(\text{Sym}^*(\mathcal{L}))
\longrightarrow
X
$$
を対応する $X$ 上のベクトル束とする。このとき、すべての $k$ に対して
$p^* : \CH_k(X) \to \CH_{k + 1}(L)$ は同型である。
\end{lemma}

\begin{proof}
全射性については補題 \ref{chow-lemma-pullback-affine-fibres-surjective} を参照されたい。
$o : X \to L$ を $L \to X$ の零切断、すなわち、すべての $n > 0$ に対して
$\mathcal{L}^{\otimes n}$ を零へ写す全射
$\text{Sym}^*(\mathcal{L}) \to \mathcal{O}_X$ に対応する射とする。
このとき $p \circ o = \text{id}_X$ であり、$o(X)$ は $L$ 上の有効 Cartier 因子である。
したがって補題 \ref{chow-lemma-relative-effective-cartier} により
$o^* \circ p^* = \text{id}$ であり、$p^*$ は単射でもある。
\end{proof}

\begin{remark}
\label{chow-remark-when-isomorphism}
後で、$X$ が $Y$ 上階数 $r$ のベクトル束ならば、引き戻し写像
$\CH_k(Y) \to \CH_{k + r}(X)$ は同型であることを見る（補題 \ref{chow-lemma-vectorbundle}）。
これは $X \to Y$ が補題 \ref{chow-lemma-pullback-affine-fibres-surjective} の仮定を満たすときには
常に成り立つ。\cite[Lemma 2.2]{Totaro-group} を参照されたい。
注意 \ref{chow-remark-higher-chow-isomorphism} では、高次 Chow 群を用いた証明を概説する。
\end{remark}

\begin{lemma}
\label{chow-lemma-linebundle-formulae}
補題 \ref{chow-lemma-linebundle} の状況で、$o : X \to L$ を零切断と表す
（同補題の証明を参照）。このとき次が成り立つ。
\begin{enumerate}
\item $o(X)$ は $p^*\mathcal{L}^{\otimes -1}$ の正則な大域切断の零点スキームである。
\item $o$ は閉埋め込みなので $o_* : \CH_k(X) \to \CH_k(L)$ が定まる。
\item $o(X)$ は有効 Cartier 因子なので $o^* : \CH_{k + 1}(L) \to \CH_k(X)$ が定まる。
\item $o^* p^* : \CH_k(X) \to \CH_k(X)$ は恒等写像である。
\item 任意の $\alpha \in \CH_k(X)$ に対し
$o_*\alpha = - p^*(c_1(\mathcal{L}) \cap \alpha)$ である。
\item $o^* o_* : \CH_k(X) \to \CH_{k - 1}(X)$ は写像
$\alpha \mapsto - c_1(\mathcal{L}) \cap \alpha$ に等しい。
\end{enumerate}
\end{lemma}

\begin{proof}
射影公式（「コホモロジー」補題 \ref{cohomology-lemma-projection-formula}）により、
$p_*\mathcal{O}_L = \text{Sym}^*(\mathcal{L})$ から
$p_*(p^*\mathcal{L}^{\otimes -1}) =
\text{Sym}^*(\mathcal{L}) \otimes_{\mathcal{O}_X} \mathcal{L}^{\otimes -1}$
を得る。(1) の切断は標準的自明化
$\mathcal{O}_X \to
\mathcal{L} \otimes_{\mathcal{O}_X} \mathcal{L}^{\otimes -1}$ である。
この切断の零点集合がちょうど $o(X)$ であることの証明は省略する。これで (1) が示された。

\medskip\noindent
(2), (3), (4) は補題 \ref{chow-lemma-linebundle} の証明中に見た。
もちろん (4) は補題 \ref{chow-lemma-relative-effective-cartier} の第一の公式である。

\medskip\noindent
(5) は、補題 \ref{chow-lemma-relative-effective-cartier} の第二の公式、
$c_1$ とのキャップ積の加法性（補題 \ref{chow-lemma-c1-cap-additive}）、および
$c_1$ とのキャップ積が平坦引き戻しと可換であること
（補題 \ref{chow-lemma-flat-pullback-cap-c1}）から従う。

\medskip\noindent
(6) は、補題 \ref{chow-lemma-gysin-back} と
$o^*p^*\mathcal{L} = \mathcal{L}$ であることから従う。
\end{proof}

\begin{lemma}
\label{chow-lemma-decompose-section}
$Y$ をスキームとし、$\mathcal{L}_i$, $i = 1, 2$ を可逆 $\mathcal{O}_Y$-加群とする。
$s$ を $\mathcal{L}_1 \otimes_{\mathcal{O}_Y} \mathcal{L}_2$ の大域切断とし、
$i : D \to Y$ で $s$ の零点スキームを表す。このとき可換図式
$$
\xymatrix{
D_1 \ar[r]_{i_1} \ar[d]_{p_1} &
L \ar[d]^p &
D_2 \ar[l]^{i_2} \ar[d]^{p_2} \\
D \ar[r]^i &
Y &
D \ar[l]_i
}
$$
と $p^*\mathcal{L}_i$ の切断 $s_i$ が存在し、次を満たす。
\begin{enumerate}
\item $p^*s = s_1 \otimes s_2$ である。
\item $p$ は有限型で相対次元 $1$ の平坦射である。
\item $D_i$ は $s_i$ の零点スキームである。
\item $i = 1, 2$ に対し、$D$ 上で
$D_i \cong
\underline{\Spec}(\text{Sym}^*(\mathcal{L}_{3 - i}^{\otimes -1})|_D))$ である。
\item $p^{-1}D = D_1 \cup D_2$ である（スキーム論的和）。
\item $D_1 \cap D_2$（スキーム論的交わり）は $D$ へ同型に写る。
\item $i = 1, 2$ に対し、$D_1 \cap D_2 \to D_i$ は直線束 $D_i \to D$ の零切断である。
\end{enumerate}
さらに、この図式と切断 $s_i$ の構成は任意の基底変換と可換である。
\end{lemma}

\begin{proof}
$p : X \to Y$ を、準連接 $\mathcal{O}_Y$-代数層
$$
\mathcal{A} =
\left(\bigoplus\nolimits_{a_1, a_2 \geq 0}
\mathcal{L}_1^{\otimes -a_1} \otimes_{\mathcal{O}_Y}
\mathcal{L}_2^{\otimes -a_2}\right)/\mathcal{J}
$$
の相対スペクトルとする。ここで $\mathcal{J}$ は、$a_1, a_2 > 0$ を満たす任意の直和因子
$\mathcal{L}_1^{\otimes -a_1} \otimes \mathcal{L}_2^{\otimes -a_2}$ の局所切断 $t$ に対する
形 $st - t$ の局所切断で生成されるイデアルである。写像
$p^*\mathcal{L}_i^{\otimes -1} \to \mathcal{O}_X$ とみた切断 $s_i$ を、写像
$\mathcal{L}_i^{\otimes -1} \to \mathcal{A} = p_*\mathcal{O}_X$ の随伴として定義する。
任意の $y \in Y$ に対し、$y$ を含むアフィン開集合 $V \subset Y$、たとえば
$V = \Spec(A)$ と、自明化 $z_i : \mathcal{O}_V \to \mathcal{L}_i^{\otimes -1}|_V$ を選べる。
$f = s(z_1z_2) \in A$ が閉部分スキーム $D$ を切り出すことに注意する。このとき明らかに
$$
p^{-1}(V) = \Spec(B[z_1, z_2]/(z_1 z_2 - f))
$$
である。$D_i$ は $z_i$ で切り出されるので、すべて明らかである。
\end{proof}

\begin{lemma}
\label{chow-lemma-decompose-section-formulae}
補題 \ref{chow-lemma-decompose-section} の状況で、$Y$ は状況
\ref{chow-situation-setup} のとおり $(S, \delta)$ 上局所有限型であると仮定する。
このとき、すべての $\alpha \in \CH_k(Y)$ に対して $\CH_k(D_1)$ において
$i_1^*p^*\alpha = p_1^*i^*\alpha$ である。
\end{lemma}

\begin{proof}
$W \subset Y$ を $\delta$-次元 $k$ の整閉部分スキームとする。二つの場合に分ける。

\medskip\noindent
$W \subset D$ と仮定する。Gysin 準同型の定義と補題
\ref{chow-lemma-c1-cap-additive} の加法性により、$\CH_{k - 1}(D)$ において
$i^*[W] = c_1(\mathcal{L}_1) \cap [W] + c_1(\mathcal{L}_2) \cap [W]$ である。
したがって $p_1^*i^*[W] =
p_1^*(c_1(\mathcal{L}_1) \cap [W]) + p_1^*(c_1(\mathcal{L}_2) \cap [W])$ である。
他方、平坦引き戻しの構成により $p^*[W] = [p^{-1}(W)]_{k + 1}$ である。
また $p^{-1}(W) = W_1 \cup W_2$ であり（スキーム論的に）、補題と図式の構成が
基底変換と可換であることから、$W_i = p_i^{-1}(W)$ は $W$ 上の直線束である。
$W_i$ は $L$ の $\delta$-次元 $k + 1$ の整閉部分スキームなので、
$[p^{-1}(W)]_{k + 1} = [W_1] + [W_2]$ である。したがって
\begin{align*}
i_1^*p^*[W]
& =
i_1^*[p^{-1}(W)]_{k + 1} \\
& =
i_1^*([W_1] + [W_2]) \\
& =
c_1(p_1^*\mathcal{L}_1) \cap [W_1] + [W_1 \cap W_2]_k \\
& =
c_1(p_1^*\mathcal{L}_1) \cap p_1^*[W] + [W_1 \cap W_2]_k \\
& =
p_1^*(c_1(\mathcal{L}_1) \cap [W]) + [W_1 \cap W_2]_k
\end{align*}
である。ここでは Gysin 準同型の構成、平坦引き戻しの定義（第二の等式）、および
$c_1 \cap -$ と平坦引き戻しの整合性（補題 \ref{chow-lemma-flat-pullback-cap-c1}）を用いた。
$W_1 \cap W_2$ は直線束 $W_1 \to W$ の零切断なので、補題
\ref{chow-lemma-linebundle-formulae} により
$[W_1 \cap W_2]_k = p_1^*(c_1(\mathcal{L}_2) \cap [W])$ である。
ここでは、$D_1$ が $\mathcal{L}_2$ の逆の相対スペクトルである直線束だという事実を用いる。
したがって先と同じ式が得られる。

\medskip\noindent
$W \not \subset D$ と仮定する。この場合、$i_1^*p^*[W]$ と $p_1^*i^*[W]$ はともに、
$D_1$ における $W$ のスキーム論的逆像に付随する $k - 1$-サイクルによって表される。
\end{proof}

\begin{lemma}
\label{chow-lemma-normal-cone-effective-Cartier}
状況 \ref{chow-situation-setup} において、$X$ を $S$ 上局所有限型なスキームとし、
$(\mathcal{L}, s, i : D \to X)$ を定義 \ref{chow-definition-gysin-homomorphism} のとおりの
三つ組とする。このとき可換図式
$$
\xymatrix{
D' \ar[r]_{i'} \ar[d]_p & X' \ar[d]^g \\
D \ar[r]^i & X
}
$$
であって、次を満たすものが存在する。
\begin{enumerate}
\item $p$ と $g$ は有限型で相対次元 $1$ の平坦射である。
\item すべての $k$ に対して $p^* : \CH_k(D) \to \CH_{k + 1}(D')$ は単射である。
\item $D' \subset X'$ は大域切断 $s' \in \Gamma(X', \mathcal{O}_{X'})$ の零点スキームである。
\item 写像 $\CH_k(X) \to \CH_k(D')$ として $p^*i^* = (i')^*g^*$ である。
\end{enumerate}
さらに、これらの性質は局所有限型な射 $Y \to X$ による任意の基底変換の後にも成り立つ。
\end{lemma}

\begin{proof}
定義 \ref{chow-definition-gysin-homomorphism} のとおりの三つ組
$(\mathcal{O}_{X'}, s', i' : D' \to X')$ があるので $(i')^*$ は定義され、主張には意味がある。

\medskip\noindent
$\mathcal{L}_1 = \mathcal{O}_X$, $\mathcal{L}_2 = \mathcal{L}$ とおき、
$\mathcal{L} = \mathcal{L}_1 \otimes_{\mathcal{O}_X} \mathcal{L}_2$ の切断 $s$ に
補題 \ref{chow-lemma-decompose-section} を適用する。$D' = D_1$ ととる。
すると主張は同補題、補題 \ref{chow-lemma-decompose-section-formulae}、および補題
\ref{chow-lemma-linebundle} による単射性から従う。
\end{proof}

\begin{remark}
\label{chow-remark-higher-chow-isomorphism}
$(S, \delta)$ を状況 \ref{chow-situation-setup} のとおりとし、$Y$ を $S$ 上局所有限型とする。
$r \geq 0$ とし、$f : X \to Y$ をスキームの射とする。各 $y \in Y$ が、
$V$ 上のスキームとして $f^{-1}(V) \cong V \times \mathbf{A}^r$ となる開集合
$V \subset Y$ に含まれると仮定する。本注意では、
$f^* : \CH_k(Y) \to \CH_{k + r}(X)$ が同型であることの証明を概説する。
まず補題 \ref{chow-lemma-pullback-affine-fibres-surjective} によりこの写像は全射である。
$f^*\alpha = 0$ を満たす $\alpha \in \CH_k(Y)$ をとり、$\alpha = 0$ を示す。

\medskip\noindent
ステップ 1. $\dim_\delta(Y) < \infty$ と仮定してよい
（実際に現れるすべての場合には直ちに成り立つので、このステップは読み飛ばしてよい）。
実際、$X$ 上で $f^*\alpha = 0$ を証する任意の有理同値には、次元 $k + r + 1$ の
整閉部分スキームの局所有限族が使われる。それらの $Y$ における像の閉包の和をとると、
$\dim_\delta(Y') \leq k + r + 1$ を満たす閉部分スキーム $Y' \subset Y$ が得られ、
$\alpha$ はある $\alpha' \in \CH_k(Y')$ の像であり、さらに $f'$ を $f$ の $Y'$ への
基底変換として $(f')^*\alpha = 0$ となる。

\medskip\noindent
ステップ 2. $d = \dim_\delta(Y) < \infty$ と仮定し、$d$ に関する帰納法を用いる。
$d < k$ ならば $\alpha = 0$ であり、これが帰納法の基底である。
一般に $f$ に関する仮定から、$U = f^{-1}(V) = \mathbf{A}^r_V$ となる稠密開集合
$V \subset Y$ を選べる。$Y' \subset Y$ で $V$ の補集合を被約閉部分スキームとして表し、
$X' = f^{-1}(Y')$ とおく。図式
$$
\xymatrix{
\CH^M_{k + r}(U, 1) \ar[r] &
\CH_{k + r}(X') \ar[r] &
\CH_{k + r}(X) \ar[r] &
\CH_{k + r}(U) \ar[r] &
0 \\
\CH^M_k(V, 1) \ar[r] \ar[u] &
\CH_k(Y') \ar[r] \ar[u] &
\CH_k(Y) \ar[r] \ar[u] &
\CH_k(V) \ar[r] \ar[u] &
0
}
$$
を考える。ここでは $V$ と $U$ の第一高次 Chow 群、注意 \ref{chow-remark-higher-chow} で
構成した六項完全列、これらの高次 Chow 群の平坦引き戻し、および平坦引き戻しと
六項完全列との整合性を用いる。$U = \mathbf{A}^r_V$ なので右端の縦写像は同型である。
$\CH_k(Y') \to \CH_{k + r}(X')$ は $d$ に関する帰納法により全単射である。
したがって議論を完結させるには
$$
\CH^M_k(V, 1) \longrightarrow \CH^M_{k + r}(U, 1)
$$
が全射であることを示せば十分である。補題
\ref{chow-lemma-pullback-affine-fibres-surjective} の証明と同様に論じれば、
これは次のステップ 3 に帰着する。

\medskip\noindent
ステップ 3. $F$ を体とする。このとき $\CH^M_0(\mathbf{A}^1_F, 1) = 0$ である
（上で引用した補題の証明では、同様に $\CH_0(\mathbf{A}^1_F) = 0$ を示した）。実際
$$
\CH^M_0(\mathbf{A}^1_F, 1) = \Coker\left(
\partial : K^M_2(F(T)) \longrightarrow
\bigoplus\nolimits_{\mathfrak p \subset F[T]\text{ maximal}}
\kappa(\mathfrak p)^*\right)
$$
である。この余核が消えることの古典的議論は、$\kappa(\mathfrak p)/F$ の次数に関する
帰納法によって、$\mathfrak p$ に対応する直和因子が像に含まれることを示すものである。
$\mathfrak p$ が既約モニック多項式 $P(T) \in F[T]$ で生成され、
$u \in \kappa(x)^*$ が $\deg(Q) < \deg(P)$ を満たすある $Q(T) \in F[T]$ の剰余類ならば、
$\partial(Q, P)$ は $\mathfrak p$ において元 $u$ を生じ、さらに $Q$ を割る素イデアルにおいて
より低次数の単元を生じる可能性があることを示せる。これで証明の概略が完了する。
\end{remark}











\section{双変交点理論}
\label{chow-section-bivariant}

\noindent
ベクトル束の高次 Chern 類を適切に論じるため、\cite{F} に従って双変 Chow 類を導入する。
本書の定義は \cite{F} と二点で異なる。(1) 扱う設定が異なり、(2) 双変類に要求するのは、
可逆加群の切断の零点スキームに対する Gysin 準同型
（節 \ref{chow-section-intersecting-effective-Cartier}）との可換性だけである。
後に補題 \ref{chow-lemma-gysin-commutes} で、ここでの双変類が任意の高余次元 Gysin 準同型とも
可換であり、したがって \cite{F} で要求されるすべての性質を満たすことを見る。
\cite[Theorem 17.1]{F} も参照されたい。

\begin{definition}
\label{chow-definition-bivariant-class}
\begin{reference}
\cite[Definition 17.1]{F} と同様である。
\end{reference}
$(S, \delta)$ を状況 \ref{chow-situation-setup} のとおりとし、
$f : X \to Y$ を $S$ 上局所有限型なスキームの射とする。$p \in \mathbf{Z}$ とする。
{\it $f$ に対する次数 $p$ の双変類 $c$} とは、任意の局所有限型射 $Y' \to Y$ と
任意の $k$ に写像
$$
c \cap - : \CH_k(Y') \longrightarrow \CH_{k - p}(X')
$$
を対応させる規則であって、ここで $X' = Y' \times_Y X$ とし、次の条件を満たすものをいう。
\begin{enumerate}
\item $Y'' \to Y'$ が固有ならば、$X'' = Y'' \times_Y X$ として、$Y''$ 上の任意の
$\alpha''$ に対し
$c \cap (Y'' \to Y')_*\alpha'' = (X'' \to X')_*(c \cap \alpha'')$ である。
\item $Y'' \to Y'$ が固定された相対次元をもつ局所有限型平坦射ならば、$Y'$ 上の任意の
$\alpha'$ に対し
$c \cap (Y'' \to Y')^*\alpha' = (X'' \to X')^*(c \cap \alpha')$ である。
\item $(\mathcal{L}', s', i' : D' \to Y')$ が定義
\ref{chow-definition-gysin-homomorphism} のとおりで、その $X'$ への引き戻しが
$(\mathcal{N}', t', j' : E' \to X')$ ならば、$Y'$ 上の任意の $\alpha'$ に対し
$c \cap (i')^*\alpha' = (j')^*(c \cap \alpha')$ である。
\end{enumerate}
$f$ に対する次数 $p$ の双変類全体を $A^p(X \to Y)$ と表す。
\end{definition}

\noindent
$(S, \delta)$ を状況 \ref{chow-situation-setup} のとおりとする。
$X \to Y$ と $Y \to Z$ を $S$ 上局所有限型なスキームの射とし、
$p \in \mathbf{Z}$ とする。$A^p(X \to Y)$ がアーベル群であることは明らかである。
さらに、結合的な双線形合成
$$
A^p(X \to Y) \times A^q(Y \to Z) \to A^{p + q}(X \to Z)
$$
があることも明らかである。

\begin{lemma}
\label{chow-lemma-flat-pullback-bivariant}
$(S, \delta)$ を状況 \ref{chow-situation-setup} のとおりとする。
$f : X \to Y$ を、$S$ 上局所有限型なスキームの相対次元 $r$ の平坦射とする。
このとき、各 $Y' \to Y$ に $X' = X \times_Y Y'$ 上の写像
$(f')^* : \CH_k(Y') \to \CH_{k + r}(X')$ を対応させる規則は次数 $-r$ の双変類である。
\end{lemma}

\begin{proof}
これは補題 \ref{chow-lemma-flat-pullback-rational-equivalence},
\ref{chow-lemma-compose-flat-pullback}, \ref{chow-lemma-flat-pullback-proper-pushforward},
\ref{chow-lemma-gysin-flat-pullback} から従う。
\end{proof}

\begin{lemma}
\label{chow-lemma-gysin-bivariant}
$(S, \delta)$ を状況 \ref{chow-situation-setup} のとおりとし、$X$ を $S$ 上局所有限型とする。
$(\mathcal{L}, s, i : D \to X)$ を定義 \ref{chow-definition-gysin-homomorphism} のとおりの
三つ組とする。このとき、各 $f : X' \to X$ に、$D' = D \times_X X'$ 上の写像
$(i')^* : \CH_k(X') \to \CH_{k - 1}(D')$ を対応させる規則は次数 $1$ の双変類である。
\end{lemma}

\begin{proof}
これは補題 \ref{chow-lemma-gysin-factors}, \ref{chow-lemma-closed-in-X-gysin},
\ref{chow-lemma-gysin-flat-pullback}, \ref{chow-lemma-gysin-commutes-gysin} から従う。
\end{proof}

\begin{lemma}
\label{chow-lemma-push-proper-bivariant}
$(S, \delta)$ を状況 \ref{chow-situation-setup} のとおりとする。
$f : X \to Y$ と $g : Y \to Z$ を $S$ 上局所有限型なスキームの射とする。
$c \in A^p(X \to Z)$ とし、$f$ は固有であると仮定する。
このとき、各 $Z' \to Z$ に
$\alpha \longmapsto f'_*(c \cap \alpha)$ を対応させる規則は双変類であり、
$f_* \circ c \in A^p(Y \to Z)$ と表す。
\end{lemma}

\begin{proof}
これは補題 \ref{chow-lemma-compose-pushforward},
\ref{chow-lemma-flat-pullback-proper-pushforward}, \ref{chow-lemma-closed-in-X-gysin} から従う。
\end{proof}

\begin{remark}
\label{chow-remark-restriction-bivariant}
$(S, \delta)$ を状況 \ref{chow-situation-setup} のとおりとする。
$X \to Y$ と $Y' \to Y$ を $S$ 上局所有限型なスキームの射とし、
$X' = Y' \times_Y X$ とおく。このとき明らかな制限写像
$$
A^p(X \to Y) \longrightarrow A^p(X' \to Y'),\quad
c \longmapsto res(c)
$$
がある。これは、$Y'$ 上局所有限型なスキーム $Y''$ を $Y$ 上局所有限型なスキームとみなし、
任意の $\alpha'' \in \CH_k(Y'')$ に対して
$res(c) \cap \alpha'' = c \cap \alpha''$ とおくことで得られる。
この制限操作は明らかな仕方で合成と整合する。
\end{remark}

\begin{remark}
\label{chow-remark-bivariant-commute}
$(S, \delta)$ を状況 \ref{chow-situation-setup} のとおりとし、$X$ を $S$ 上局所有限型とする。
$i = 1, 2$ に対し $Z_i \to X$ を局所有限型なスキームの射とし、
$c_i \in A^{p_i}(Z_i \to X)$, $i = 1, 2$ を双変類とする。
任意の $\alpha \in \CH_k(X)$ に対し、$\CH_{k - p_1 - p_2}(Z_1 \times_X Z_2)$ において
$$
c_1 \cap c_2 \cap \alpha = c_2 \cap c_1 \cap \alpha
$$
かどうかを問える。この等式が成り立ち、さらに任意の局所有限型基底変換
$X' \to X$ の後にも成り立つとき、$c_1$ と $c_2$ は {\it 可換である} という。
これはもちろん、$A^{p_1 + p_2}(Z_1 \times_X Z_2 \to X)$ において
$$
res(c_1) \circ c_2 = res(c_2) \circ c_1
$$
であることと同じである。ここで
$res(c_1) \in A^{p_1}(Z_1 \times_X Z_2 \to Z_2)$ は注意
\ref{chow-remark-restriction-bivariant} の意味での $c_1$ の制限であり、$res(c_2)$ も同様である。
\end{remark}

\begin{example}
\label{chow-example-gysin-commute}
$(S, \delta)$ を状況 \ref{chow-situation-setup} のとおりとし、$X$ を $S$ 上局所有限型とする。
$(\mathcal{L}, s, i : D \to X)$ を定義 \ref{chow-definition-gysin-homomorphism} のとおりの
三つ組とする。$Z \to X$ を局所有限型なスキームの射とし、
$c \in A^p(Z \to X)$ を双変類とする。補題 \ref{chow-lemma-gysin-bivariant} の双変 Gysin 類
$c' \in A^1(D \to X)$ は、注意 \ref{chow-remark-bivariant-commute} の意味で $c$ と可換である。
これは定義 \ref{chow-definition-bivariant-class} の条件 (3) の言い換えにほかならない。
\end{example}

\begin{remark}
\label{chow-remark-more-general-bivariant}
文献では検討されていないと思われる、より一般的な型の双変類がある。
すなわち、状況 \ref{chow-situation-setup} のとおり $(S, \delta)$ 上局所有限型なスキームの図式
$$
X \longrightarrow Z \longleftarrow Y
$$
が与えられたとし、$p \in \mathbf{Z}$ とする。このとき、各局所有限型射
$Z' \to Z$ とすべての $k \in \mathbf{Z}$ に写像
$$
c \cap - : \CH_k(Y') \longrightarrow \CH_{k - p}(X')
$$
を対応させる規則 $c$ を考えられる。ここで $X' = X \times_Z Z'$ および
$Y' = Z' \times_Z Y$ とし、次と整合することを要求する。
\begin{enumerate}
\item $Z'' \to Z'$ が固有である場合の固有押し出し。
\item $Z'' \to Z'$ が固定された相対次元をもつ平坦射である場合の平坦引き戻し。
\item $D' \subset Z'$ が定義 \ref{chow-definition-gysin-homomorphism} のとおりである場合の Gysin 写像。
\end{enumerate}
詳細な定式化は省略する。このような操作全体を
$A^p(X \to Z \leftarrow Y)$ と表すことにする。この概念が有用な簡単な例として、
固有射 $f : X_2 \to X_1$ がある場合を考える。$f_*$ は定義
\ref{chow-definition-bivariant-class} の意味では双変操作ではないが、上の一般化された意味では
双変操作であり、実際 $f_* \in A^0(X_1 \to X_1 \leftarrow X_2)$ である。
\end{remark}





\section{Chow コホモロジーと第一 Chern 類}
\label{chow-section-chow-cohomology}

\noindent
主として関心をもつのは $A^p(X) = A^p(X \to X)$ であり、これは常に
$\text{id}_X$ に対する双変コホモロジー類を意味する。Chern 類はここに属する。

\begin{definition}
\label{chow-definition-chow-cohomology}
$(S, \delta)$ を状況 \ref{chow-situation-setup} のとおりとし、$X$ を $S$ 上局所有限型とする。
$X$ の {\it Chow コホモロジー} とは、次数 $p$ の成分が $A^p(X \to X)$ である
次数付き $\mathbf{Z}$-代数 $A^*(X)$ をいう。
\end{definition}

\noindent
注意：$A^*(X)$ の $\mathbf{Z}$-代数構造が可換かどうかは明らかでない。
しかし Chern 類がその中心に属することを示す。

\begin{remark}
\label{chow-remark-pullback-cohomology}
$(S, \delta)$ を状況 \ref{chow-situation-setup} のとおりとし、
$f : Y' \to Y$ を $S$ 上局所有限型なスキームの射とする。
注意 \ref{chow-remark-restriction-bivariant} の特別な場合として、標準的な
$\mathbf{Z}$-代数写像 $res : A^*(Y) \to A^*(Y')$ がある。
文献ではこの写像をしばしば $f^*$ と表す。
\end{remark}

\begin{lemma}
\label{chow-lemma-cap-c1-bivariant}
$(S, \delta)$ を状況 \ref{chow-situation-setup} のとおりとし、$X$ を $S$ 上局所有限型とする。
$\mathcal{L}$ を可逆 $\mathcal{O}_X$-加群とする。このとき、各 $f : X' \to X$ に
$c_1(f^*\mathcal{L}) \cap - : \CH_k(X') \to \CH_{k - 1}(X')$ を対応させる規則は
次数 $1$ の双変類である。
\end{lemma}

\begin{proof}
これは補題 \ref{chow-lemma-factors}, \ref{chow-lemma-pushforward-cap-c1},
\ref{chow-lemma-flat-pullback-cap-c1}, \ref{chow-lemma-gysin-commutes-cap-c1} から従う。
\end{proof}

\noindent
上の補題により、ようやく次の定義が可能になる。

\begin{definition}
\label{chow-definition-first-chern-class}
$(S, \delta)$ を状況 \ref{chow-situation-setup} のとおりとし、$X$ を $S$ 上局所有限型とする。
$\mathcal{L}$ を可逆 $\mathcal{O}_X$-加群とする。$\mathcal{L}$ の {\it 第一 Chern 類}
$c_1(\mathcal{L}) \in A^1(X)$ とは、補題 \ref{chow-lemma-cap-c1-bivariant} の双変類をいう。
\end{definition}

\noindent
有限局所自由加群の Chern 類は節 \ref{chow-section-intersecting-chern-classes} で構成する。
$c_1(\mathcal{L})$ が $A^*(X)$ の中心に属することを証明しよう。

\begin{lemma}
\label{chow-lemma-c1-center}
$(S, \delta)$ を状況 \ref{chow-situation-setup} のとおりとし、$X$ を $S$ 上局所有限型とする。
$\mathcal{L}$ を可逆 $\mathcal{O}_X$-加群とする。このとき次が成り立つ。
\begin{enumerate}
\item $c_1(\mathcal{L}) \in A^1(X)$ は $A^*(X)$ の中心に属する。
\item $f : X' \to X$ が局所有限型で $c \in A^*(X' \to X)$ ならば、
$c \circ c_1(\mathcal{L}) = c_1(f^*\mathcal{L}) \circ c$ である。
\end{enumerate}
\end{lemma}

\begin{proof}
もちろん (2) から (1) が従う。$p : L \to X$ を補題 \ref{chow-lemma-linebundle} のとおりとし、
$o : X \to L$ を零切断とする。それらの基底変換を
$p' : L' \to X'$ および $o' : X' \to L'$ と表す。
補題 \ref{chow-lemma-linebundle-formulae} により
$$
p^*(c_1(\mathcal{L}) \cap \alpha) = - o_* \alpha
\quad\text{and}\quad
(p')^*(c_1(f^*\mathcal{L}) \cap \alpha') = - o'_* \alpha'
$$
である。$c$ は双変類なので
\begin{align*}
(p')^*(c \cap c_1(\mathcal{L}) \cap \alpha)
& =
c \cap p^*(c_1(\mathcal{L}) \cap \alpha) \\
& =
- c \cap o_* \alpha \\
& =
- o'_*(c \cap \alpha) \\
& =
(p')^*(c_1(f^*\mathcal{L}) \cap c \cap \alpha)
\end{align*}
を得る。上で引用した補題の一つにより $(p')^*$ は単射なので
$c \cap c_1(\mathcal{L}) \cap \alpha =
c_1(f^*\mathcal{L}) \cap c \cap \alpha$ である。
同じことが任意の局所有限型基底変換 $Y \to X$ の後にも成り立つので、
(2) の双変類の等式を得る。
\end{proof}

\begin{lemma}
\label{chow-lemma-vanish-above-dimension}
$(S, \delta)$ を状況 \ref{chow-situation-setup} のとおりとする。$X$ を、豊富な可逆層をもつ
$S$ 上有限型なスキームとする。$d = \dim(X) < \infty$ と仮定する
（ここでいう次元は通常の次元であり、$\delta$-次元ではない）。
このとき、$X$ 上の任意の可逆層 $\mathcal{L}_1, \ldots, \mathcal{L}_{d + 1}$ に対して
$A^{d + 1}(X)$ において
$c_1(\mathcal{L}_1) \circ \ldots \circ c_1(\mathcal{L}_{d + 1}) = 0$
が成り立つ。
\end{lemma}

\begin{proof}
$d$ に関する帰納法で証明する。$d = 0$ の場合、$X$ は閉点の有限集合であり、
したがってすべての可逆加群は自明なので、主張は成り立つ。$d > 0$ と仮定する。
『因子』の補題 \ref{divisors-lemma-quasi-projective-Noetherian-pic-effective-Cartier}
により、ある有効 Cartier 因子 $D, D' \subset X$ を用いて
$\mathcal{L}_{d + 1} \cong \mathcal{O}_X(D) \otimes
\mathcal{O}_X(D')^{\otimes -1}$ と書ける。このとき $c_1(\mathcal{L}_{d + 1})$ は
$c_1(\mathcal{O}_X(D))$ と $c_1(\mathcal{O}_X(D'))$ の差であるから、
$\mathcal{L}_{d + 1} = \mathcal{O}_X(D)$（ここで D は有効 Cartier 因子）と仮定してよい。

\medskip\noindent
$i : D \to X$ を包含射とし、$i^* \in A^1(D \to X)$ を
補題 \ref{chow-lemma-gysin-bivariant} の Gysin 準同型が与える双変類とする。
補題 \ref{chow-lemma-support-cap-effective-Cartier} により
$A^1(X)$ において $i_* \circ i^* = c_1(\mathcal{L}_{d + 1})$ である
（左辺が意味をもつことには補題 \ref{chow-lemma-push-proper-bivariant} も用いる）。
$c_1(\mathcal{L}_i)$ は $i_*$ と $i^*$ の双方と可換なので
（双変類の定義による）、次を得る。
$$
c_1(\mathcal{L}_1) \circ \ldots \circ c_1(\mathcal{L}_{d + 1}) =
i_* \circ c_1(\mathcal{L}_1) \circ \ldots \circ c_1(\mathcal{L}_d) \circ i^* =
i_* \circ c_1(\mathcal{L}_1|_D) \circ \ldots \circ c_1(\mathcal{L}_d|_D)
\circ i^*
$$
$X$ の生成点はいずれも $D$ に属さないため $\dim(D) < d$ であり、
$d$ に関する帰納法から結論が従う。
\end{proof}

\begin{remark}
\label{chow-remark-ring-loc-classes}
$(S, \delta)$ を状況 \ref{chow-situation-setup} のとおりとする。
$Z \to X$ を $S$ 上局所有限型なスキームの閉埋め込みとし、$p \geq 0$ とする。
この設定で
$$
A^{(p)}(Z \to X) =
\prod\nolimits_{i \leq p - 1} A^i(X) \times
\prod\nolimits_{i \geq p} A^i(Z \to X).
$$
と定める。このとき $A^{(p)}(Z \to X)$ は標準的に次数付き代数の構造をもつ。
実際、より一般に積
$$
A^{(p)}(Z \to X) \times A^{(q)}(Z \to X)
\longrightarrow A^{(\max(p, q))}(Z \to X)
$$
がある。これらを定義するため、次の写像を定める。
\begin{align*}
A^i(Z \to X) \times A^j(X) & \to A^{i + j}(Z \to X) \\
A^i(X) \times A^j(Z \to X) & \to A^{i + j}(Z \to X) \\
A^i(Z \to X) \times A^j(Z \to X) & \to A^{i + j}(Z \to X)
\end{align*}
第一の写像には双変類の合成を用いる。第二の写像には制限
$A^i(X) \to A^i(Z)$（注意 \ref{chow-remark-restriction-bivariant}）と合成
$A^i(Z) \times A^j(Z \to X) \to A^{i + j}(Z \to X)$ を用いる。
第三の写像では $(c, c')$ を $res(c) \circ c'$ に送る。ここで
$res : A^i(Z \to X) \to A^i(Z)$ は制限写像である
（注意 \ref{chow-remark-restriction-bivariant} を参照）。
これらの積が適切な意味で結合的であることの検証は省略する。
\end{remark}

\begin{remark}
\label{chow-remark-res-push}
$(S, \delta)$ を状況 \ref{chow-situation-setup} のとおりとする。
$Z \to X$ を $S$ 上局所有限型なスキームの閉埋め込みとする。
注意 \ref{chow-remark-restriction-bivariant} の制限写像を
$res : A^p(Z \to X) \to A^p(Z)$ と表す。
$c \in A^p(Z \to X)$ に対し、$\alpha \in \CH_*(Z)$ ならば
$res(c) \cap \alpha = c \cap i_*\alpha$ である。
実際、$res(c) \cap \alpha = c \cap \alpha$ であり、$c$ と固有押し出しとの両立性から
$(Z \to Z)_*(c \cap \alpha) = c \cap (Z \to X)_*\alpha$ を得る。
\end{remark}






\section{双変類に関する補題}
\label{chow-section-bivariant-II}

\noindent
この節では双変類に関するいくつかの初等的な結果を証明する。
まず、演算が有理同値を通して定まるための判定法を与える。

\begin{lemma}
\label{chow-lemma-factors-through-rational-equivalence}
\begin{reference}
\cite[Theorem 17.1]{F} の非常に弱い形
\end{reference}
$(S, \delta)$ を状況 \ref{chow-situation-setup} のとおりとする。
$f : X \to Y$ を $S$ 上局所有限型なスキームの射とし、$p \in \mathbf{Z}$ とする。
任意の局所有限型射 $Y' \to Y$ と任意の $k$ に対して写像
$$
c \cap - : Z_k(Y') \longrightarrow \CH_{k - p}(X')
$$
を対応させる規則が与えられているとする。ここで $Y' = X' \times_X Y$ であり、
$\mathcal{L}'|_{D'} \cong \mathcal{O}_{D'}$ であるときには
定義 \ref{chow-definition-bivariant-class} の条件 (3) を満たすものとする。
このとき $c \cap -$ は有理同値を通して定まる。
\end{lemma}

\begin{proof}
この主張が意味をもつことをまず確認する。定義 \ref{chow-definition-gysin-homomorphism} の
三つ組 $(\mathcal{L}, s, i : D \to X)$ が
$\mathcal{L}|_D \cong \mathcal{O}_D$ を満たすとき、演算 $i^*$ はサイクルのレベルで
定義されるからである。注意 \ref{chow-remark-gysin-on-cycles} を参照。
$\alpha \in Z_k(X')$ を零と有理同値なサイクルとする。
$c \cap \alpha = 0$ を示せばよい。補題 \ref{chow-lemma-rational-equivalence-family} により、
$\alpha = i_0^*\beta - i_\infty^*\beta$ を満たすサイクル
$\beta \in Z_{k + 1}(X' \times \mathbf{P}^1)$ が存在する。ここで
$i_0, i_\infty : X' \to X' \times \mathbf{P}^1$ は $0, \infty$ 上の
$X'$ の閉埋め込みである。これらは自明な法束をもつ有効 Cartier 因子の例なので、
$c \cap i_0^*\beta = j_0^*(c \cap \beta)$ および
$c \cap i_\infty^*\beta = j_\infty^*(c \cap \beta)$ である。
ここで $j_0, j_\infty : Y' \to Y' \times \mathbf{P}^1$ は先と同様の閉埋め込みである。
補題 \ref{chow-lemma-rational-equivalence-family} から
$j_0^*(c \cap \beta) \sim_{rat} j_\infty^*(c \cap \beta)$ なので、結論を得る。
\end{proof}

\begin{lemma}
\label{chow-lemma-bivariant-weaker}
\begin{reference}
\cite[Theorem 17.1]{F} の弱い形
\end{reference}
$(S, \delta)$ を状況 \ref{chow-situation-setup} のとおりとする。
$f : X \to Y$ を $S$ 上局所有限型なスキームの射とし、$p \in \mathbf{Z}$ とする。
任意の局所有限型射 $Y' \to Y$ と任意の $k$ に対して写像
$$
c \cap - : \CH_k(Y') \longrightarrow \CH_{k - p}(X')
$$
を対応させる規則が与えられているとする。ここで $Y' = X' \times_X Y$ であり、
定義 \ref{chow-definition-bivariant-class} の条件 (1), (2) と、
$\mathcal{L}'|_{D'} \cong \mathcal{O}_{D'}$ である場合の条件 (3) を満たすものとする。
このとき $c \cap -$ は双変類である。
\end{lemma}

\begin{proof}
$Y' \to Y$ を局所有限型なスキームの射とする。
$(\mathcal{L}', s', i' : D' \to Y')$ を定義 \ref{chow-definition-gysin-homomorphism} の
とおりとし、その $X'$ への引き戻しを $(\mathcal{N}', t', j' : E' \to X')$ とする。
すべての $\alpha' \in \CH_k(Y')$ に対して
$c \cap (i')^*\alpha' = (j')^*(c \cap \alpha')$ を示す必要がある。

\medskip\noindent
補題 \ref{chow-lemma-normal-cone-effective-Cartier} で構成した、相対次元 $1$ の滑らかな射
$g : Y'' \to Y'$ と、$i'' : D'' \to Y''$ および $p : D'' \to D'$ を用いる。
（注意：$D''$ は $D'$ の完全な逆像ではない。）
$X' \to Y'$ によるそれらの基底変換を $f : X'' \to X'$ および
$E'' \subset X''$ と表す。図式は
$$
\xymatrix{
& X'' \ar[rr] \ar'[d][dd]_h & & Y'' \ar[dd]^g \\
E'' \ar[rr] \ar[dd]_q \ar[ru]^{j''} & & D'' \ar[dd]^p \ar[ru]^{i''} & \\
& X' \ar'[r][rr] & & Y' \\
E' \ar[rr] \ar[ru]^{j'} & & D' \ar[ru]^{i'}
}
$$
である。補題で与えられた性質により、$\beta' = (i')^*\alpha'$ は
$p^*\beta' = (i'')^*g^*\alpha'$ を満たす $\CH_{k - 1}(D')$ の一意な元である。
同様に、$\gamma' = (j')^*(c \cap \alpha')$ は
$q^*\gamma' = (j'')^*h^*(c \cap \alpha')$ を満たす
$\CH_{k - 1 - p}(E')$ の一意な元である。一方、$c$ に関する仮定から
$$
(j'')^*h^*(c \cap \alpha') =
(j'')^*(c \cap g^*\alpha') =
c \cap (i'')^*g^*\alpha'
$$
である。ここで、補題の主張で仮定した条件 (3) の修正版は、$i''$ とその基底変換
$j''$ に適用できることに注意する。同様に
$$
q^*(c \cap \beta') = c \cap p^*\beta'
$$
である。したがって、上で述べた一意性から $\gamma' = c \cap \beta'$ を得る。
\end{proof}

\noindent
次に、双変類が零となるための判定法を与える。

\begin{lemma}
\label{chow-lemma-bivariant-zero}
$(S, \delta)$ を状況 \ref{chow-situation-setup} のとおりとする。
$f : X \to Y$ を $S$ 上局所有限型なスキームの射とし、$c \in A^p(X \to Y)$ とする。
$Y'' \to Y' \to Y$ に対して $X'' = Y'' \times_Y X$ および
$X' = Y' \times_Y X$ とおく。次の条件は同値である。
\begin{enumerate}
\item $c$ は零である。
\item $Y$ 上局所有限型な任意の整スキーム $Y'$ に対し、
$\CH_*(X')$ において $c \cap [Y'] = 0$ である。
\item $Y$ 上局所有限型な任意の整スキーム $Y'$ に対し、
$\CH_*(X'')$ において $c \cap [Y''] = 0$ となるような固有双有理射
$Y'' \to Y'$ が存在する。
\end{enumerate}
\end{lemma}

\begin{proof}
含意 (1) $\Rightarrow$ (2) $\Rightarrow$ (3) は明らかである。
条件 (3) から (2) が従う。実際、$(Y'' \to Y')_*[Y''] = [Y']$ であり、
$c$ は双変類なので $c \cap [Y'] = (X'' \to X')_*(c \cap [Y''])$ である。
条件 (2) を仮定する。$Y' \to Y$ を局所有限型とし、$\alpha \in \CH_k(Y')$ とする。
$Y'_i \subset Y'$ を $\delta$-次元 $k$ の整閉部分スキームの局所有限族として
$\alpha = \sum n_i [Y'_i]$ と書く。このとき $\alpha$ は、固有射 $Y'' \to Y'$ による、
$Y'' = \coprod Y'_i$ 上のサイクル $\alpha' = \sum n_i[Y'_i]$ の押し出しである。
双変類の性質により、$\CH_{k - p}(X'')$ において
$c \cap \alpha' = 0$ であることを示せば十分である。
$X'_i = Y'_i \times_Y X$ とおけば
$\CH_{k - p}(X'') = \prod \CH_{k - p}(X'_i)$ である。これは定義から直ちに従う。
射影写像 $\CH_{k - p}(X'') \to \CH_{k - p}(X'_i)$ は平坦引き戻しで与えられる。
$c$ とのキャップ積は平坦引き戻しと可換なので、$\CH_{k - p}(X'_i)$ において
$c \cap [Y'_i]$ が零であることを示せば十分であり、これは仮定から成り立つ。
\end{proof}

\begin{lemma}
\label{chow-lemma-disjoint-decomposition-bivariant}
$(S, \delta)$ を状況 \ref{chow-situation-setup} のとおりとする。
$f : X \to Y$ を $S$ 上局所有限型なスキームの射とする。
開かつ閉な部分スキームによる非交和分解
$X = \coprod_{i \in I} X_i$ および $Y = \coprod_{j \in J} Y_j$ と、
$f(X_i) \subset Y_{a(i)}$ を満たす集合の写像 $a : I \to J$ があると仮定する。
このとき
$$
A^p(X \to Y) = \prod\nolimits_{i \in I} A^p(X_i \to Y_{a(i)})
$$
である。
\end{lemma}

\begin{proof}
元 $(c_i) \in \prod_i A^p(X_i \to Y_{a(i)})$ が与えられたとする。
$\beta \in \CH_k(Y)$ に対し、$X_i$ への制限が
$c_i \cap \beta|_{Y_{a(i)}}$ である $\CH_{k - p}(X)$ の元を対応させることができる。
これは $\CH_{k - p}(X) = \prod_i \CH_{k - p}(X_i)$ だからである。
同じ構成は任意の局所有限型基底変換 $Y' \to Y$ の後にも機能し、
$c \in A^p(X \to Y)$ を得る。したがって式の右辺から左辺への写像 $\Psi$ を得る。
逆に、$c \in A^p(X \to Y)$ と元 $\beta_i \in \CH_k(Y_{a(i)})$ が与えられたとき、
$\CH_{k - p}(X_i)$ の元 $(c \cap (Y_{a(i)} \to Y)_*\beta_i)|_{X_i}$ を考えられる。
同じ構成は任意の局所有限型基底変換 $Y' \to Y$ の後にも機能し、
$c_i \in A^p(X_i \to Y_{a(i)})$ を得る。したがって式の左辺から右辺への写像
$\Phi$ を得る。$\Phi \circ \Psi = \text{id}$ は直ちに分かる。
逆向きを示すため、$c \in A^p(X \to Y)$, $\beta \in \CH_k(Y)$ とし、
$\Phi(c) = (c_i)$ と書く。$j \in J$ とする。
$c$ は平坦引き戻しと可換なので
$$
(c \cap \beta)|_{\coprod_{a(i) = j} X_i} =
c \cap \beta|_{Y_j}
$$
を得る。また $c$ は固有押し出しと可換なので
$$
(\coprod\nolimits_{a(i) = j} X_i \to X)_*
((c \cap \beta)|_{\coprod_{a(i) = j} X_i})
=
c \cap (Y_j \to Y)_*\beta|_{Y_j}
$$
を得る。左辺は、$a(i) = j$ を満たす $i \in I$ について $X_i$ への制限が
$(c \cap \beta)|_{X_i}$ であり、それ以外では $0$ である $X$ 上のサイクルである。
右辺は、$a(i) = j$ を満たす $i \in I$ について $X_i$ への制限が
$c_i \cap \beta|_{Y_j}$ である $X$ 上のサイクルである。
したがって所望のとおり $c \cap \beta = \Psi((c_i))$ である。
\end{proof}

\begin{remark}
\label{chow-remark-completion-bivariant}
$(S, \delta)$ を状況 \ref{chow-situation-setup} のとおりとする。
$f : X \to Y$ を $S$ 上局所有限型なスキームの射とする。
$X = \coprod_{i \in I} X_i$ および $Y = \coprod_{j \in J} Y_j$ を、
$X$ と $Y$ の連結成分への分解とする（$X$ と $Y$ は局所 Noether なので連結成分は開である。
『位相』の補題 \ref{topology-lemma-locally-Noetherian-locally-connected} および
『スキームの性質』の補題 \ref{properties-lemma-Noetherian-topology} を参照）。
$f(X_i) \subset Y_{a(i)}$ を満たす添字を $a(i) \in J$ とする。
補題 \ref{chow-lemma-disjoint-decomposition-bivariant} により
$A^p(X \to Y) = \prod A^p(X_i \to Y_{a(i)})$ である。
この設定では
$$
A^*(X \to Y)^\wedge = \prod\nolimits_i A^*(X_i \to Y_{a(i)})
$$
とおくと便利である。この「完備化された」双変群は
$$
A^*(X \to Y)^\wedge \quad\subset\quad \prod\nolimits_{p \geq 0} A^p(X \to Y)
$$
のうち、各連結成分 $X_i$ について、ほとんどすべての $p$ に対して
$c_p$ の $A^p(X_i \to Y_{a(i)})$ における像が零となる元
$c = (c_0, c_1, c_2, \ldots)$ からなる部分集合である。
$Y \to Z$ をもう一つの射とすると、合成
$A^*(X \to Y) \times A^*(Y \to Z) \to A^*(X \to Z)$ は、完備化の合成
$A^*(X \to Y)^\wedge \times A^*(Y \to Z)^\wedge \to A^*(X \to Z)^\wedge$
へ拡張される。
$A^*(X)^\wedge = A^*(X \to X)^\wedge$ を $X$ の
{\it 完備化双変コホモロジー環} と呼ぶことがある。
\end{remark}

\begin{lemma}
\label{chow-lemma-envelope-bivariant}
$(S, \delta)$ を状況 \ref{chow-situation-setup} のとおりとする。
$f : X \to Y$ を $S$ 上局所有限型なスキームの射とする。
$g : Y' \to Y$ を包絡射（定義 \ref{chow-definition-envelope}）とし、
$X' = Y' \times_Y X$ と表す。$p \in \mathbf{Z}$ とし、
$c' \in A^p(X' \to Y')$ とする。二つの制限が等しく
$$
res_1(c') = res_2(c') \in A^p(X' \times_X X' \to Y' \times_Y Y')
$$
なるとする（証明を参照）。このとき、$A^p(X' \to Y')$ における制限が
$res(c) = c'$ となる一意な $c \in A^p(X \to Y)$ が存在する。
\end{lemma}

\begin{proof}
可換図式
$$
\xymatrix{
X' \times_X X' \ar[d]^{f''} \ar@<1ex>[r]^-a \ar@<-1ex>[r]_-b &
X' \ar[d]^{f'} \ar[r]_h &
X \ar[d]^f \\
Y' \times_Y Y' \ar@<1ex>[r]^-p \ar@<-1ex>[r]_-q &
Y' \ar[r]^g &
Y
}
$$
がある。元 $res_1(c')$ は、射 $a, f', p, f''$ からなる Cartesian 正方形に関する
$c'$ の制限（注意 \ref{chow-remark-restriction-bivariant} を参照）であり、元 $res_2(c')$ は
射 $b, f', q, f''$ からなる Cartesian 正方形に関する $c'$ の制限である。
$res_1(c') = res_2(c')$ と仮定し、$\beta \in \CH_k(Y)$ とする。
補題 \ref{chow-lemma-envelope} により $g_*\beta' = \beta$ を満たす
$\beta' \in \CH_k(Y')$ を選べる。そこで
$$
c \cap \beta = h_*(c' \cap \beta')
$$
とおく。これが $\beta'$ の選び方に依存しないことを示すには、
$\gamma \in \CH_k(Y' \times_Y Y')$ に対して
$h_*(c' \cap (p_*\gamma - q_*\gamma))$ が零であることを示せば十分である。
$c'$ は双変類なので
$$
h_*(c' \cap (p_*\gamma - q_*\gamma)) =
h_*(a_*(c' \cap \gamma) - b_*(c' \cap \gamma)) = 0
$$
である。最後の等式は $h \circ a = h \circ b$ から
$h_* \circ a_* = h_* \circ b_*$ となることによる。

\medskip\noindent
この $c \cap \beta$ の選び方は、$res(c) = c'$ という要請と、
双変類の固有押し出しとの両立性によって強制されることに注意する。

\medskip\noindent
もちろん、双変類 $c$ を定義するには、定義 \ref{chow-definition-bivariant-class} に列挙された
条件を満たす任意の局所有限型射 $Y_1 \to Y$ に対し、写像
$c \cap -: \CH_k(Y_1) \to \CH_{k + p}(Y_1 \times_Y X)$ を構成する必要がある。
$Y'_1 = Y' \times_Y Y_1$, $X_1 = X \times_Y Y_1$ と表す。
補題 \ref{chow-lemma-base-change-envelope} により射 $Y'_1 \to Y_1$ は包絡射である。
したがって、基底変換した図式
$$
\xymatrix{
X'_1 \times_{X_1} X'_1 \ar[d]^{f''_1} \ar@<1ex>[r]^-{a_1} \ar@<-1ex>[r]_-{b_1} &
X'_1 \ar[d]^{f'_1} \ar[r]_{h_1} &
X_1 \ar[d]^{f_1} \\
Y'_1 \times_{Y_1} Y'_1 \ar@<1ex>[r]^-{p_1} \ar@<-1ex>[r]_-{q_1} &
Y'_1 \ar[r]^{g_1} &
Y_1
}
$$
と同じ議論を用いて、先と同様に良定義な写像
$c \cap - : \CH_k(Y_1) \to \CH_{k + p}(X_1)$ を得る。

\medskip\noindent
次に、$c$ について定義 \ref{chow-definition-bivariant-class} の条件 (1), (2), (3) を
確認しなければならない。例えば、$t : Y_2 \to Y_1$ を
$Y$ 上局所有限型なスキームの固有射とする。上と同様に、第一の図式の
$Y_1$, resp.\ $Y_2$ への基底変換を添字 ${}_1$, resp.\ ${}_2$ で表す。
$t$ の $Y'$, $X$, $X'$ への基底変換をそれぞれ $t' : Y'_2 \to Y'_1$,
$s : X_2 \to X_1$, $s' : X'_2 \to X'_1$ と表す。示すべきことは
$$
s_*(c \cap \beta_2) = c \cap t_*\beta_2
$$
である。ここで $\beta_2 \in \CH_k(Y_2)$ である。
$g_{2, *}\beta'_2 = \beta_2$ を満たす $\beta'_2 \in \CH_k(Y'_2)$ を選ぶ。
$c'$ は双変類であり、図式
$$
\vcenter{
\xymatrix{
X'_2 \ar[d]_{s'} \ar[r]_{h_2} & X_2 \ar[d]^s \\
X'_1 \ar[r]^{h_1} & X_1
}
}
\quad\text{and}\quad
\vcenter{
\xymatrix{
X'_2 \ar[d]_{s'} \ar[r]_{f'_2} & Y'_2 \ar[d]^{t'} \\
X'_2 \ar[r]^{f'_1} & Y'_1
}
}
$$
は Cartesian なので
$$
s_*(c \cap \beta_2) =
s_*(h_{2, *}(c' \cap \beta'_2)) =
h_{1, *}s'_*(c' \cap \beta'_2) =
h_{1, *}(c' \cap (t'_*\beta'_2))
$$
を得る。最後の式は構成により $c \cap t_*\beta_2$ を計算している。
実際、$t'_*\beta'_2 \in \CH_k(Y'_1)$ の $g_{1, *}$ による像は
$t_*\beta_2$ である。これで条件 (1) が示された。
他の条件も同様に示されるので、詳しい議論は省略する。
\end{proof}















\section{射影空間束公式}
\label{chow-section-projective-space-bundle-formula}

\noindent
$(S, \delta)$ を状況 \ref{chow-situation-setup} のとおりとし、$X$ を $S$ 上局所有限型とする。
階数 $r$ の有限局所自由 $\mathcal{O}_X$-加群 $\mathcal{E}$ を考える。
本章では、$\mathcal{E}$ に {\it 付随する射影束} を射
$$
\xymatrix{
\mathbf{P}(\mathcal{E}) =
\underline{\text{Proj}}_X(\text{Sym}^*(\mathcal{E}))
\ar[r]^-\pi
& X
}
$$
とする。ただし $X$ 上で $\mathcal{O}_{\mathbf{P}(\mathcal{E})}(1)$ を
$\pi_*(\mathcal{O}_{\mathbf{P}(\mathcal{E})}(1)) = \mathcal{E}$ となるよう正規化する。
特に全射 $\pi^*\mathcal{E} \to \mathcal{O}_{\mathbf{P}(\mathcal{E})}(1)$ がある。
略式に「$(\pi : P \to X, \mathcal{O}_P(1))$ を $\mathcal{E}$ に付随する射影束とする」
と言うときは、$P = \mathbf{P}(\mathcal{E})$ かつ
$\mathcal{O}_P(1) = \mathcal{O}_{\mathbf{P}(\mathcal{E})}(1)$ である状況を意味する。

\begin{lemma}
\label{chow-lemma-cap-projective-bundle}
$(S, \delta)$ を状況 \ref{chow-situation-setup} のとおりとし、$X$ を $S$ 上局所有限型とする。
$\mathcal{E}$ を階数 $r$ の有限局所自由 $\mathcal{O}_X$-加群 $\mathcal{E}$ とする。
$(\pi : P \to X, \mathcal{O}_P(1))$ を $\mathcal{E}$ に付随する射影束とする。
任意の $\alpha \in \CH_k(X)$ に対し、元
$$
\pi_*\left(
c_1(\mathcal{O}_P(1))^s \cap \pi^*\alpha
\right)
\in
\CH_{k + r - 1 - s}(X)
$$
は $s < r - 1$ のとき $0$ であり、$s = r - 1$ のとき $\alpha$ に等しい。
\end{lemma}

\begin{proof}
$Z \subset X$ を $\delta$-次元 $k$ の整閉部分スキームとする。
$\pi^{-1}(Z)$ は $\delta$-次元 $r - 1$ の整スキームなので、
$\pi^*[Z] = [\pi^{-1}(Z)]$ であることに注意する。
$s < r - 1$ ならば、構成により
$c_1(\mathcal{O}_P(1))^s \cap \pi^*[Z]$ は $\pi^{-1}(Z)$ に台をもつ
$(k + r - 1 - s)$-サイクルで表される。したがって次元上の理由から、
このサイクルの押し出しは零である。

\medskip\noindent
$s = r - 1$ とする。上の議論により、ある $n \in \mathbf{Z}$ に対して
$\pi_*(c_1(\mathcal{O}_P(1))^s \cap \pi^*\alpha) = n [Z]$ である。
$n = 1$ を示したい。上と同じ次元上の理由により、$T \subset Z$ を真の閉部分集合として
$X$ を $X \setminus T$ に置き換えた後に主張を示せば十分である。
$\xi$ を $Z$ の生成点とする。$\mathcal{E}_\xi$ の基底の一部をなす元
$e_1, \ldots, e_{r - 1} \in \mathcal{E}_\xi$ を選べる。
これらは $\mathcal{O}_P(1)|_{\pi^{-1}(Z)}$ の有理切断
$s_1, \ldots, s_{r - 1}$ を与える。その共通零点集合は、
$\mathbf{P}(\mathcal{E}|_Z) \to Z$ のある有理切断の像の閉包と、台が $Z$ の真の閉部分集合
$T$ に写る閉部分集合との合併である。$X$ から $T$ を除き、それに応じて
$P$ から $\pi^{-1}(T)$ を除けば、$s_1, \ldots, s_n$ は大域切断の列
$s_i \in \Gamma(\pi^{-1}(Z), \mathcal{O}_{\pi^{-1}(Z)}(1))$ をなし、
その共通零点集合は切断 $Z \to \pi^{-1}(Z)$ の像である。したがって順次
\begin{eqnarray*}
\pi^*[Z] & = & [\pi^{-1}(Z)] \\
c_1(\mathcal{O}_P(1)) \cap \pi^*[Z] & = & [Z(s_1)] \\
c_1(\mathcal{O}_P(1))^2 \cap \pi^*[Z] & = & [Z(s_1) \cap Z(s_2)] \\
\ldots & = & \ldots \\
c_1(\mathcal{O}_P(1))^{r - 1} \cap \pi^*[Z] & = &
[Z(s_1) \cap \ldots \cap Z(s_{r - 1})]
\end{eqnarray*}
を得る。これは補題 \ref{chow-lemma-geometric-cap} を繰り返し適用したものである。
$Z$ 上の $\pi$ の切断の像を $\pi$ で押し出すと明らかに $[Z]$ になるので、
$\alpha = [Z]$ の場合の結論を得る。これらの議論から一般のサイクル
$\alpha = \sum n_j [Z_j]$ に対する結論が従うことの検証は省略する。
\end{proof}

\begin{lemma}[射影空間束公式]
\label{chow-lemma-chow-ring-projective-bundle}
$(S, \delta)$ を状況 \ref{chow-situation-setup} のとおりとし、$X$ を $S$ 上局所有限型とする。
$\mathcal{E}$ を階数 $r$ の有限局所自由 $\mathcal{O}_X$-加群
$\mathcal{E}$ とする。$(\pi : P \to X, \mathcal{O}_P(1))$ を
$\mathcal{E}$ に付随する射影束とする。写像
$$
\bigoplus\nolimits_{i = 0}^{r - 1}
\CH_{k + i}(X)
\longrightarrow
\CH_{k + r - 1}(P),
$$
$$
(\alpha_0, \ldots, \alpha_{r-1})
\longmapsto
\pi^*\alpha_0 +
c_1(\mathcal{O}_P(1)) \cap \pi^*\alpha_1
+ \ldots +
c_1(\mathcal{O}_P(1))^{r - 1} \cap \pi^*\alpha_{r-1}
$$
は同型である。
\end{lemma}

\begin{proof}
$k \in \mathbf{Z}$ を固定する。まず写像が単射であることを示す。
左辺の元 $(\alpha_0, \ldots, \alpha_{r - 1})$ の像が零であるとする。
補題 \ref{chow-lemma-cap-projective-bundle} により
$$
0 = \pi_*(\pi^*\alpha_0 +
c_1(\mathcal{O}_P(1)) \cap \pi^*\alpha_1
+ \ldots +
c_1(\mathcal{O}_P(1))^{r - 1} \cap \pi^*\alpha_{r-1})
= \alpha_{r - 1}
$$
である。次に
$$
0 = \pi_*(c_1(\mathcal{O}_P(1)) \cap (\pi^*\alpha_0 +
c_1(\mathcal{O}_P(1)) \cap \pi^*\alpha_1
+ \ldots +
c_1(\mathcal{O}_P(1))^{r - 2} \cap \pi^*\alpha_{r - 2}))
= \alpha_{r - 2}
$$
を得る。以下同様である。したがって写像は単射である。

\medskip\noindent
写像が全射であることを示せばよい。
$X_i$, $i \in I$ を $X$ の既約成分とする。このとき
$P_i = \mathbf{P}(\mathcal{E}|_{X_i})$, $i \in I$ は $P$ の既約成分である。
可換図式
$$
\xymatrix{
\coprod P_i \ar[d]_{\coprod \pi_i} \ar[r]_p & P \ar[d]^\pi \\
\coprod X_i \ar[r]^q & X
}
$$
を考える。$p_*$ は全射であることに注意する。
$\beta \in \CH_k(\coprod X_i)$ ならば
$\pi^* q_* \beta = p_*(\coprod \pi_i)^* \beta$ である。
補題 \ref{chow-lemma-flat-pullback-proper-pushforward} を参照。
補題 \ref{chow-lemma-pushforward-cap-c1} により
$c_1(\mathcal{O}(1))$ とのキャップ積についても同様である。
したがって、各射 $\pi_i : P_i \to X_i$ に対して補題の写像が全射ならば、
$\pi : P \to X$ に対しても全射である。ゆえに $X$ は既約と仮定してよい。
すると $\dim_\delta(X) < \infty$ であり、特に $\dim_\delta(X)$ に関する帰納法を使える。

\medskip\noindent
$\dim_\delta(X) < k$ ならば結論は明らかである。
$\alpha \in \CH_{k + r - 1}(P)$ とする。任意の局所閉部分スキーム $T \subset X$ に対し、
$\gamma_T : \bigoplus \CH_{k + i}(T) \to \CH_{k + r - 1}(\pi^{-1}(T))$ を写像
$$
\gamma_T(\alpha_0, \ldots, \alpha_{r - 1})
= \pi^*\alpha_0 + \ldots +
c_1(\mathcal{O}_{\pi^{-1}(T)}(1))^{r - 1} \cap \pi^*\alpha_{r - 1}.
$$
とする。ある空でない開集合 $U \subset X$ に対して
$\alpha|_{\pi^{-1}(U)} = \gamma_U(\alpha_0, \ldots, \alpha_{r - 1})$ であるとする。
このとき持ち上げ $\alpha'_i \in \CH_{k + i}(X)$ を選べる。
補題 \ref{chow-lemma-restrict-to-open} により
$\alpha - \gamma_X(\alpha'_0, \ldots, \alpha'_{r - 1})$ は、
$Y = X \setminus U$ を被約閉部分スキームとみたとき、
$P_Y = \mathbf{P}(\mathcal{E}|_Y)$ 上の $k$-サイクルと有理同値である。
$\dim_\delta(Y) < \dim_\delta(X)$ であることに注意する。
帰納法により $P_Y \to Y$ に対して結論が成り立つので、$\alpha$ に対しても成り立つ。
したがって $X$ を $X$ の任意の空でない開部分に置き換えてよい。

\medskip\noindent
特に $\mathcal{E} \cong \mathcal{O}_X^{\oplus r}$ と仮定してよい。
この場合 $\mathbf{P}(\mathcal{E}) = X \times \mathbf{P}^{r - 1}$ である。
成層
$$
\mathbf{P}^{r - 1} = \mathbf{A}^{r - 1}
\amalg \mathbf{A}^{r - 2}
\amalg \ldots
\amalg \mathbf{A}^0
$$
を用いる。各層の閉包は $\mathbf{P}^{r - 1 - i}$ であり、
$c_1(\mathcal{O}(1))^i \cap [\mathbf{P}^{r - 1}]$ の代表である。
したがって $P$ にも同様の成層
$$
P = U^{r - 1} \amalg U^{r - 2} \amalg \ldots \amalg U^0
$$
がある。$P^i$ を $U^i$ の閉包とし、$\pi^i : P^i \to X$ を
$\pi$ の $P^i$ への制限とする。
$\alpha \in \CH_{k + r - 1}(P)$ とする。補題
\ref{chow-lemma-pullback-affine-fibres-surjective} により、ある
$\alpha_0 \in \CH_k(X)$ に対して
$\alpha|_{U^{r - 1}} = \pi^*\alpha_0|_{U^{r - 1}}$ と書ける。
したがって差 $\alpha - \pi^*\alpha_0$ は、ある
$\alpha' \in \CH_{k + r - 1}(P^{r - 2})$ の像である。
再び補題 \ref{chow-lemma-pullback-affine-fibres-surjective} により、ある
$\alpha_1 \in \CH_{k + 1}(X)$ に対して
$\alpha'|_{U^{r - 2}} = (\pi^{r - 2})^*\alpha_1|_{U^{r - 2}}$ と書ける。
補題 \ref{chow-lemma-relative-effective-cartier} により
$(\pi^{r - 2})^*\alpha_1$ の像は
$c_1(\mathcal{O}_P(1)) \cap \pi^*\alpha_1$ を表す。
また
$\alpha - \pi^*\alpha_0 - c_1(\mathcal{O}_P(1)) \cap \pi^*\alpha_1$ は、
ある $\alpha'' \in \CH_{k + r - 1}(P^{r - 3})$ の像である。
以下同様である。
\end{proof}

\begin{lemma}
\label{chow-lemma-vectorbundle}
$(S, \delta)$ を状況 \ref{chow-situation-setup} のとおりとし、$X$ を $S$ 上局所有限型とする。
$\mathcal{E}$ を $X$ 上の階数 $r$ の有限局所自由層とする。
$$
p :
E = \underline{\Spec}(\text{Sym}^*(\mathcal{E}))
\longrightarrow
X
$$
を $X$ 上の付随するベクトル束とする。このときすべての $k$ に対して
$p^* : \CH_k(X) \to \CH_{k + r}(E)$ は同型である。
\end{lemma}

\begin{proof}
（直線束の場合については補題 \ref{chow-lemma-linebundle} を参照。）
全射性については補題 \ref{chow-lemma-pullback-affine-fibres-surjective} を参照する。
$(\pi  : P \to X, \mathcal{O}_P(1))$ を、有限局所自由層
$\mathcal{E} \oplus \mathcal{O}_X$ に付随する射影空間束とする。
大域切断 $(0, 1) \in \Gamma(X, \mathcal{E} \oplus \mathcal{O}_X)$ に対応する
$s \in \Gamma(P, \mathcal{O}_P(1))$ をとる。
$D = Z(s) \subset P$ とする。
$(\pi|_D : D \to X , \mathcal{O}_P(1)|_D)$ は $\mathcal{E}$ に付随する射影空間束である。
$\pi_D = \pi|_D$ および $\mathcal{O}_D(1) = \mathcal{O}_P(1)|_D$ と表す。
さらに $D$ は $P$ 上の有効 Cartier 因子である。したがって
$\mathcal{O}_P(D) = \mathcal{O}_P(1)$ である
（『因子』の補題 \ref{divisors-lemma-characterize-OD} を参照）。
また同型 $E \cong P \setminus D$ がある。対応する開埋め込みを
$j : E \to P$ と表す。
単射性の証明には、
$$
j^* :
\CH_{k + r}(P)
\longrightarrow
\CH_{k + r}(E)
$$
の核が有効 Cartier 因子 $D$ に台をもつサイクルであることを用いる。
補題 \ref{chow-lemma-restrict-to-open} を参照。
したがって $p^*\alpha = 0$ ならば、ある $\beta \in \CH_{k + r}(D)$ に対して
$\pi^*\alpha = i_*\beta$ である。
補題 \ref{chow-lemma-chow-ring-projective-bundle} により
$$
\beta = \pi_D^*\beta_0 +
\ldots + c_1(\mathcal{O}_D(1))^{r - 1} \cap \pi_D^* \beta_{r - 1}.
$$
と書ける。ここで $\beta_i \in \CH_{k + i}(X)$ である。
補題 \ref{chow-lemma-relative-effective-cartier} および
\ref{chow-lemma-pushforward-cap-c1} により、これは
$$
\pi^*\alpha = i_*\beta =
c_1(\mathcal{O}_P(1)) \cap \pi^*\beta_0 +
\ldots +
c_1(\mathcal{O}_D(1))^r \cap \pi^*\beta_{r - 1}.
$$
を含意する。$\mathcal{E} \oplus \mathcal{O}_X$ の階数は $r + 1$ なので、
$\alpha$ およびすべての $\beta_i$ が零でない限り、これは
補題 \ref{chow-lemma-pushforward-cap-c1} に矛盾する。
\end{proof}








\section{ベクトル束の Chern 類}
\label{chow-section-chern-classes-vector-bundles}

\noindent
射影空間束公式を用いると、階数 $r$ のベクトル束の Chern 類を、
$c_1(\mathcal{O}(1))^r$ の低い冪による展開を通じて定義できる。
式 (\ref{chow-equation-chern-classes}) を参照。符号の理由は後で説明する。

\begin{definition}
\label{chow-definition-chern-classes}
$(S, \delta)$ を状況 \ref{chow-situation-setup} のとおりとし、$X$ を $S$ 上局所有限型とする。
$X$ は整で $n = \dim_\delta(X)$ であると仮定する。
$\mathcal{E}$ を $X$ 上の階数 $r$ の有限局所自由層とし、
$(\pi : P \to X, \mathcal{O}_P(1))$ を $\mathcal{E}$ に付随する射影空間束とする。
\begin{enumerate}
\item 補題 \ref{chow-lemma-chow-ring-projective-bundle} により、
$c_0 = [X]$ かつ
\begin{equation}
\label{chow-equation-chern-classes}
\sum\nolimits_{i = 0}^r
(-1)^i c_1(\mathcal{O}_P(1))^i \cap \pi^*c_{r - i}
= 0.
\end{equation}
を満たす元 $c_i \in \CH_{n - i}(X)$, $i = 0, \ldots, r$ が存在する。
\item 上の記号のもとで、$\CH_{n - i}(X)$ の元として
$c_i(\mathcal{E}) \cap [X] = c_i$ とおく。
これらを {\it $X$ 上の $\mathcal{E}$ の Chern 類} と呼ぶ。
\item {\it $X$ 上の $\mathcal{E}$ の全 Chern 類} とは
$$
c({\mathcal E}) \cap [X] =
c_0({\mathcal E}) \cap [X]
+ c_1({\mathcal E}) \cap [X] + \ldots
+ c_r({\mathcal E}) \cap [X]
$$
という和をいう。これは
$\CH_*(X) = \bigoplus_{k \in \mathbf{Z}} \CH_k(X)$ の元である。
\end{enumerate}
\end{definition}

\noindent
ベクトル束の階数が $1$ の場合、この定義が新しい概念を与えないことを確認しよう。

\begin{lemma}
\label{chow-lemma-first-chern-class}
$(S, \delta)$ を状況 \ref{chow-situation-setup} のとおりとし、$X$ を $S$ 上局所有限型とする。
$X$ は整で $n = \dim_\delta(X)$ であると仮定する。
$\mathcal{L}$ を可逆 $\mathcal{O}_X$-加群とする。
定義 \ref{chow-definition-chern-classes} による $X$ 上の $\mathcal{L}$ の第一 Chern 類は、
定義 \ref{chow-definition-divisor-invertible-sheaf} により $\mathcal{L}$ に付随する
Weil 因子に等しい。
\end{lemma}

\begin{proof}
この証明では、$c_1(\mathcal{L}) \cap [X]$ で
定義 \ref{chow-definition-divisor-invertible-sheaf} の構成を表す。
$\mathcal{L}$ の階数は $1$ なので、本章の正規化により
$\mathbf{P}(\mathcal{L}) = X$ および
$\mathcal{O}_{\mathbf{P}(\mathcal{L})}(1) = \mathcal{L}$ である。
したがって (\ref{chow-equation-chern-classes}) は
$$
(-1)^1 c_1(\mathcal{L}) \cap c_0 + (-1)^0 c_1 = 0
$$
となる。$c_0 = [X]$ なので、所望の
$c_1 = c_1(\mathcal{L}) \cap [X]$ を得る。
\end{proof}

\begin{remark}
\label{chow-remark-equation-signs}
式 \ref{chow-equation-chern-classes} を
\begin{equation}
\label{chow-equation-signs}
\sum\nolimits_{i = 0}^r
c_1(\mathcal{O}_P(-1))^i \cap \pi^*c_{r - i}
= 0.
\end{equation}
と書き直すこともできる。しかし双対よりも、普遍的商層
$\mathcal{O}_P(1)$ を扱う方が容易である。
\end{remark}




\section{Chern 類との交叉}
\label{chow-section-intersecting-chern-classes}

\noindent
この節では、$X$ 上のベクトル束の Chern 類を $X$ 上の双変類として定義する。
補題 \ref{chow-lemma-cap-cp-bivariant} とそれに続く議論を参照。
構成は、まず素サイクル上で演算を定義し、次に和をとるという通常の手順に従う。
補題 \ref{chow-lemma-determine-intersections} では、付随する射影束上の通常の公式によって
結果が決定されることを示す。次に、Chern 類とのキャップ積が有理同値を通して定まり、
固有押し出し、平坦引き戻し、および有効 Cartier 因子の包含に対する Gysin 写像と
可換であることを示す。これらの補題は、補題 \ref{chow-lemma-determine-intersections} の
特徴づけと補題 \ref{chow-lemma-push-proper-bivariant} を直接用いれば避けることもできる。
詳細な展開を見たい読者は『代数空間の Chow 群』の補題 \ref{spaces-chow-lemma-segre-classes} を参照されたい。

\begin{definition}
\label{chow-definition-cap-chern-classes}
$(S, \delta)$ を状況 \ref{chow-situation-setup} のとおりとし、$X$ を $S$ 上局所有限型とする。
$\mathcal{E}$ を $X$ 上の階数 $r$ の有限局所自由層とする。
任意の整数 $k$ と任意の $0 \leq j \leq r$ に対して、演算
$$
c_j(\mathcal{E}) \cap - : Z_k(X) \to \CH_{k - j}(X)
$$
を次のように定義し、{\it $\mathcal{E}$ の第 $j$ Chern 類との交叉} と呼ぶ。
\begin{enumerate}
\item $\delta$-次元 $k$ の整閉部分スキーム $i : W \to X$ が与えられたとき
$$
c_j(\mathcal{E}) \cap [W] = i_*(c_j({i^*\mathcal{E}}) \cap [W])
\in
\CH_{k - j}(X)
$$
と定める。ここで $c_j({i^*\mathcal{E}}) \cap [W]$ は
定義 \ref{chow-definition-chern-classes} のとおりである。
\item 一般の $k$-サイクル $\alpha = \sum n_i [W_i]$ に対して
$$
c_j(\mathcal{E}) \cap \alpha = \sum n_i c_j(\mathcal{E}) \cap [W_i]
$$
とおく。
\end{enumerate}
\end{definition}

\noindent
$\mathcal{E}$ の階数が $1$ ならば、これは以前の定義
（定義 \ref{chow-definition-cap-c1}）と補題 \ref{chow-lemma-first-chern-class} により一致する。

\begin{lemma}
\label{chow-lemma-determine-intersections}
$(S, \delta)$ を状況 \ref{chow-situation-setup} のとおりとし、$X$ を $S$ 上局所有限型とする。
$\mathcal{E}$ を $X$ 上の階数 $r$ の有限局所自由層とする。
$(\pi : P \to X, \mathcal{O}_P(1))$ を $\mathcal{E}$ に付随する射影束とする。
$\alpha \in Z_k(X)$ に対し、元 $c_j(\mathcal{E}) \cap \alpha$ は、
$\alpha_0 = \alpha$ かつ
$$
\sum\nolimits_{i = 0}^r
(-1)^i c_1(\mathcal{O}_P(1))^i \cap
\pi^*(\alpha_{r - i}) = 0
$$
が $P$ の Chow 群で成り立つような $\CH_{k - j}(X)$ の一意な元
$\alpha_j$ である。
\end{lemma}

\begin{proof}
$\alpha_0 = \alpha$ であり、かつ表示された等式を満たす
$\alpha_0, \ldots, \alpha_r$ の一意性は、射影空間束公式
補題 \ref{chow-lemma-chow-ring-projective-bundle} から従う。
$W$ が $X$ の整閉部分スキームであるとき、等式は
$\alpha = [W]$ に対して定義により成り立つ。
一般の $X$ 上の $k$-サイクル $\alpha$ に対して、
$n_a \not = 0$ かつ $i_a : W_a \to X$ が相異なる整閉部分スキームとなるよう
$\alpha = \sum n_a[W_a]$ と書く。族 $\{W_a\}$ は $X$ 上局所有限である。
$P_a = \pi^{-1}(W_a) = \mathbf{P}(\mathcal{E}|_{W_a})$ とおく。
対応する閉埋め込みを $i'_a : P_a \to P$ と表す。ファイバー積図式
$$
\xymatrix{
P' \ar@{=}[r] \ar[d]_{\pi'} &
\coprod P_a \ar[d]_{\coprod \pi_a} \ar[r]_{\coprod i'_a} &
P \ar[d]^\pi \\
X' \ar@{=}[r] &
\coprod W_a \ar[r]^{\coprod i_a} &
X
}
$$
を考える。射 $p : X' \to X$ は固有である。さらに
$\pi' : P' \to X'$ は、可逆層
$\mathcal{O}_{P'}(1) = \coprod \mathcal{O}_{P_a}(1)$ と合わせて、
$\mathcal{E}' = p^*\mathcal{E}$ に付随する射影束である。
この可逆層は $\mathcal{O}_P(1)$ の引き戻しでもある。定義により
$$
c_j(\mathcal{E}) \cap [\alpha]
=
\sum i_{a, *}(c_j(\mathcal{E}|_{W_a}) \cap [W_a]).
$$
である。$\beta_{a, j} = c_j(\mathcal{E}|_{W_a}) \cap [W_a]$ と書く。
これは $\CH_{k - j}(W_a)$ の元である。定義により各 $a$ に対して
$$
\sum\nolimits_{i = 0}^r
(-1)^i c_1(\mathcal{O}_{P_a}(1))^i \cap \pi_a^*(\beta_{a, r - i}) = 0
$$
である。したがって明らかに
$$
\sum\nolimits_{i = 0}^r
(-1)^i c_1(\mathcal{O}_{P'}(1))^i \cap (\pi')^*(\beta_{r - i}) = 0
$$
である。ここで $\beta_j = \sum n_a\beta_{a, j} \in \CH_{k - j}(X')$ とする。
$p' : P' \to P$ を射 $\coprod i'_a$ とする。
補題 \ref{chow-lemma-flat-pullback-proper-pushforward} により
$\pi^*p_*\beta_j = p'_*(\pi')^*\beta_j$ である。
補題 \ref{chow-lemma-pushforward-cap-c1} の射影公式により
$$
\sum\nolimits_{i = 0}^r
(-1)^i c_1(\mathcal{O}_P(1))^i \cap \pi^*(p_*\beta_j) = 0
$$
を得る。$p_*\beta_j$ は $c_j(\mathcal{E}) \cap \alpha$ の代表なので、結論を得る。
\end{proof}

\noindent
以後、Chern 類のこの特徴づけを一貫して用い、さらに多くの性質を証明する。

\begin{lemma}
\label{chow-lemma-cap-chern-class-factors-rational-equivalence}
$(S, \delta)$ を状況 \ref{chow-situation-setup} のとおりとし、$X$ を $S$ 上局所有限型とする。
$\mathcal{E}$ を $X$ 上の階数 $r$ の有限局所自由層とする。
$\alpha \sim_{rat} \beta$ が $X$ 上の有理同値な $k$-サイクルならば、
$\CH_{k - j}(X)$ において
$c_j(\mathcal{E}) \cap \alpha = c_j(\mathcal{E}) \cap \beta$ である。
\end{lemma}

\begin{proof}
補題 \ref{chow-lemma-determine-intersections} により、元
$\alpha_j = c_j(\mathcal{E}) \cap \alpha$, $j \geq 1$ と
$\beta_j = c_j(\mathcal{E}) \cap \beta$, $j \geq 1$ は、
$\mathcal{E}$ に付随する射影束の Chow 群における {\it 同じ} 等式によって
一意に決定される。（これはもちろん、平坦引き戻しが有理同値と両立することを用いている。
補題 \ref{chow-lemma-flat-pullback-rational-equivalence} を参照。）したがって両者は等しい。
\end{proof}

\noindent
言い換えると、有限局所自由層の Chern 類とのキャップ積は有理同値を通して定まり、写像
$$
c_j(\mathcal{E}) \cap - : \CH_k(X) \to \CH_{k - j}(X).
$$
を与える。次の課題は、Chern 類が双変類であることを示すことである。
定義 \ref{chow-definition-bivariant-class} を参照。

\begin{lemma}
\label{chow-lemma-pushforward-cap-cj}
$(S, \delta)$ を状況 \ref{chow-situation-setup} のとおりとし、
$X$, $Y$ を $S$ 上局所有限型とする。
$\mathcal{E}$ を $X$ 上の階数 $r$ の有限局所自由層とする。
$p : X \to Y$ を固有射とし、$\alpha$ を $X$ 上の $k$-サイクルとする。
$\mathcal{E}$ を $Y$ 上の有限局所自由層とする。このとき
$$
p_*(c_j(p^*\mathcal{E}) \cap \alpha) = c_j(\mathcal{E}) \cap p_*\alpha
$$
である。
\end{lemma}

\begin{proof}
$(\pi : P \to Y, \mathcal{O}_P(1))$ を $\mathcal{E}$ に付随する射影束とする。
このとき $P_X = X \times_Y P$ は $p^*\mathcal{E}$ に付随する射影束であり、
$\mathcal{O}_{P_X}(1)$ は $\mathcal{O}_P(1)$ の引き戻しである。
$\alpha_j = c_j(p^*\mathcal{E}) \cap \alpha$ と書けば $\alpha_0 = \alpha$ である。
補題 \ref{chow-lemma-determine-intersections} により
$$
\sum\nolimits_{i = 0}^r
(-1)^i c_1(\mathcal{O}_P(1))^i \cap
\pi_X^*(\alpha_{r - i}) = 0
$$
が $P_X$ の Chow 群で成り立つ。ファイバー積図式
$$
\xymatrix{
P_X \ar[r]_-{p'} \ar[d]_{\pi_X} & P \ar[d]^\pi \\
X \ar[r]^p & Y
}
$$
を考える。上の等式に固有押し出し $p'_*$
（補題 \ref{chow-lemma-proper-pushforward-rational-equivalence}）を適用する。
補題 \ref{chow-lemma-pushforward-cap-c1} および
\ref{chow-lemma-flat-pullback-proper-pushforward} を用いると
$$
\sum\nolimits_{i = 0}^r
(-1)^i c_1(\mathcal{O}_P(1))^i \cap
\pi^*(p_*\alpha_{r - i}) = 0
$$
を $P$ の Chow 群で得る。補題
\ref{chow-lemma-determine-intersections} の特徴づけから結論が従う。
\end{proof}

\begin{lemma}
\label{chow-lemma-flat-pullback-cap-cj}
$(S, \delta)$ を状況 \ref{chow-situation-setup} のとおりとし、
$X$, $Y$ を $S$ 上局所有限型とする。
$\mathcal{E}$ を $Y$ 上の階数 $r$ の有限局所自由層とする。
$f : X \to Y$ を相対次元 $r$ の平坦射とし、$\alpha$ を $Y$ 上の $k$-サイクルとする。
このとき
$$
f^*(c_j(\mathcal{E}) \cap \alpha) = c_j(f^*\mathcal{E}) \cap f^*\alpha
$$
である。
\end{lemma}

\begin{proof}
$\alpha_j = c_j(\mathcal{E}) \cap \alpha$ と書けば $\alpha_0 = \alpha$ である。
補題 \ref{chow-lemma-determine-intersections} により
$$
\sum\nolimits_{i = 0}^r
(-1)^i c_1(\mathcal{O}_P(1))^i \cap
\pi^*(\alpha_{r - i}) = 0
$$
が $\mathcal{E}$ に付随する射影束
$(\pi : P \to Y, \mathcal{O}_P(1))$ の Chow 群で成り立つ。
ファイバー積図式
$$
\xymatrix{
P_X = \mathbf{P}(f^*\mathcal{E}) \ar[r]_-{f'} \ar[d]_{\pi_X} &
P \ar[d]^\pi \\
X \ar[r]^f & Y
}
$$
を考える。$\mathcal{O}_{P_X}(1)$ は $\mathcal{O}_P(1)$ の引き戻しである。
上の等式に平坦引き戻し $(f')^*$（補題
\ref{chow-lemma-flat-pullback-rational-equivalence}）を適用する。
補題 \ref{chow-lemma-flat-pullback-cap-c1} および
\ref{chow-lemma-compose-flat-pullback} により
$$
\sum\nolimits_{i = 0}^r
(-1)^i c_1(\mathcal{O}_{P_X}(1))^i \cap
\pi_X^*(f^*\alpha_{r - i}) = 0
$$
が $P_X$ の Chow 群で成り立つ。補題
\ref{chow-lemma-determine-intersections} の特徴づけから結論が従う。
\end{proof}

\begin{lemma}
\label{chow-lemma-cap-chern-class-commutes-with-gysin}
$(S, \delta)$ を状況 \ref{chow-situation-setup} のとおりとし、$X$ を $S$ 上局所有限型とする。
$\mathcal{E}$ を $X$ 上の階数 $r$ の有限局所自由層とする。
$(\mathcal{L}, s, i : D \to X)$ を定義
\ref{chow-definition-gysin-homomorphism} のとおりとする。このときすべての
$\alpha \in \CH_k(X)$ に対して
$c_j(\mathcal{E}|_D) \cap i^*\alpha = i^*(c_j(\mathcal{E}) \cap \alpha)$ である。
\end{lemma}

\begin{proof}
$\alpha_j = c_j(\mathcal{E}) \cap \alpha$ と書けば $\alpha_0 = \alpha$ である。
補題 \ref{chow-lemma-determine-intersections} により
$$
\sum\nolimits_{i = 0}^r
(-1)^i c_1(\mathcal{O}_P(1))^i \cap
\pi^*(\alpha_{r - i}) = 0
$$
が $\mathcal{E}$ に付随する射影束
$(\pi : P \to X, \mathcal{O}_P(1))$ の Chow 群で成り立つ。
ファイバー積図式
$$
\xymatrix{
P_D = \mathbf{P}(\mathcal{E}|_D) \ar[r]_-{i'} \ar[d]_{\pi_D} &
P \ar[d]^\pi \\
D \ar[r]^i & X
}
$$
を考える。$\mathcal{O}_{P_D}(1)$ は $\mathcal{O}_P(1)$ の引き戻しである。
表示された等式に Gysin 写像 $(i')^*$（補題 \ref{chow-lemma-gysin-factors}）を適用する。
補題 \ref{chow-lemma-gysin-commutes-cap-c1} および
\ref{chow-lemma-gysin-flat-pullback} を用いると
$$
\sum\nolimits_{i = 0}^r
(-1)^i c_1(\mathcal{O}_{P_D}(1))^i \cap
\pi_D^*(i^*\alpha_{r - i}) = 0
$$
を $P_D$ の Chow 群で得る。補題
\ref{chow-lemma-determine-intersections} の特徴づけから結論が従う。
\end{proof}

\noindent
以上で、Chern 類とのキャップ積が双変類を定めることを証明するための材料が揃った。

\begin{lemma}
\label{chow-lemma-cap-cp-bivariant}
$(S, \delta)$ を状況 \ref{chow-situation-setup} のとおりとし、$X$ を $S$ 上局所有限型とする。
$\mathcal{E}$ を階数 $r$ の局所自由 $\mathcal{O}_X$-加群とし、
$0 \leq p \leq r$ とする。このとき各 $f : X' \to X$ に
$c_p(f^*\mathcal{E}) \cap - : \CH_k(X') \to \CH_{k - p}(X')$
を対応させる規則は次数 $p$ の双変類である。
\end{lemma}

\begin{proof}
補題
\ref{chow-lemma-cap-chern-class-factors-rational-equivalence},
\ref{chow-lemma-pushforward-cap-cj},
\ref{chow-lemma-flat-pullback-cap-cj},
\ref{chow-lemma-cap-chern-class-commutes-with-gysin}
および定義 \ref{chow-definition-bivariant-class} から直ちに従う。
\end{proof}

\noindent
この補題により、有限局所自由加群の Chern 類を次のように定義できる。

\begin{definition}
\label{chow-definition-chern-classes-final}
$(S, \delta)$ を状況 \ref{chow-situation-setup} のとおりとし、$X$ を $S$ 上局所有限型とする。
$\mathcal{E}$ を階数 $r$ の局所自由 $\mathcal{O}_X$-加群とする。
$i = 0, \ldots, r$ に対し、$\mathcal{E}$ の {\it 第 $i$ Chern 類} とは、
補題 \ref{chow-lemma-cap-cp-bivariant} で構成した次数 $i$ の双変類
$c_i(\mathcal{E}) \in A^i(X)$ をいう。
$\mathcal{E}$ の {\it 全 Chern 類} とは形式和
$$
c(\mathcal{E}) = 
c_0(\mathcal{E}) + c_1(\mathcal{E}) + \ldots + c_r(\mathcal{E})
$$
をいう。これは $X$ 上の非斉次双変類とみなす。
\end{definition}

\noindent
定義 \ref{chow-definition-cap-chern-classes} に続く注意により、
$\mathcal{E}$ が可逆ならば、この定義は
定義 \ref{chow-definition-first-chern-class} と一致する。
次に、Chern 類が双変 Chow コホモロジー環 $A^*(X)$ の中心に属することを示す。

\begin{lemma}
\label{chow-lemma-cap-commutative-chern}
$(S, \delta)$ を状況 \ref{chow-situation-setup} のとおりとし、$X$ を $S$ 上局所有限型とする。
$\mathcal{E}$ を階数 $r$ の局所自由 $\mathcal{O}_X$-加群とする。このとき
\begin{enumerate}
\item $c_j(\mathcal{E}) \in A^j(X)$ は $A^*(X)$ の中心に属する。
\item $f : X' \to X$ が局所有限型で $c \in A^*(X' \to X)$ ならば、
$c \circ c_j(\mathcal{E}) = c_j(f^*\mathcal{E}) \circ c$ である。
\end{enumerate}
特に、$\mathcal{F}$ を $X$ 上の階数 $s$ のもう一つの局所自由
$\mathcal{O}_X$-加群とすると、すべての $\alpha \in \CH_k(X)$ に対して
$$
c_i(\mathcal{E}) \cap c_j(\mathcal{F}) \cap \alpha
=
c_j(\mathcal{F}) \cap c_i(\mathcal{E}) \cap \alpha
$$
が $\CH_{k - i - j}(X)$ の元として成り立つ。
\end{lemma}

\begin{proof}
(2) から (1) が従うことは直ちに分かる。
$\alpha \in \CH_k(X)$ とし、$\alpha_j = c_j(\mathcal{E}) \cap \alpha$ と書けば
$\alpha_0 = \alpha$ である。補題 \ref{chow-lemma-determine-intersections} により
$$
\sum\nolimits_{i = 0}^r
(-1)^i c_1(\mathcal{O}_P(1))^i \cap
\pi^*(\alpha_{r - i}) = 0
$$
が $\mathcal{E}$ に付随する射影束
$(\pi : P \to Y, \mathcal{O}_P(1))$ の Chow 群で成り立つ。
$\pi$ の $f$ による基底変換を $\pi' : P' \to X'$ と表す。
補題 \ref{chow-lemma-c1-center} と双変類の性質を用いると
\begin{align*}
0 & = c \cap \left(\sum\nolimits_{i = 0}^r
(-1)^i c_1(\mathcal{O}_P(1))^i \cap
\pi^*(\alpha_{r - i})\right) \\
& =
\sum\nolimits_{i = 0}^r
(-1)^i c_1(\mathcal{O}_{P'}(1))^i \cap
(\pi')^*(c \cap \alpha_{r - i})
\end{align*}
を $P'$ の Chow 群で得る（計算は省略する）。
したがって補題 \ref{chow-lemma-determine-intersections} の特徴づけにより、
$c \cap \alpha_j$ は $c_j(f^*\mathcal{E}) \cap (c \cap \alpha)$ に等しい。
これで補題が示された。
\end{proof}

\begin{remark}
\label{chow-remark-extend-to-finite-locally-free}
$(S, \delta)$ を状況 \ref{chow-situation-setup} のとおりとし、$X$ を $S$ 上局所有限型とする。
$\mathcal{E}$ を有限局所自由 $\mathcal{O}_X$-加群とする。
$\mathcal{E}$ の階数が一定でなくても、$\mathcal{E}$ の Chern 類を定義できる。
実際、この場合
$$
X = X_0 \amalg X_1 \amalg X_2 \amalg \ldots
$$
と書ける。ここで $X_r \subset X$ は $\mathcal{E}$ の階数が $r$ である
開かつ閉な部分空間である。補題
\ref{chow-lemma-disjoint-decomposition-bivariant} により
$A^p(X) = \prod A^p(X_r)$ である。
そこで $c_p(\mathcal{E})$ を、$A^p(X_r)$ における類
$c_p(\mathcal{E}|_{X_r})$ の積として定義できる。
具体的に、$X' \to X$ が局所有限型な射ならば、引き戻しにより
$X'$ の対応する分解を得て、定義から
$$
\CH_*(X') = \prod\nolimits_{r \geq 0} \CH_*(X'_r)
$$
となる。このとき $c_p(\mathcal{E}) \in A^p(X)$ は、これらの直積分解を保ち、
各因子上で既に定義した演算 $c_i(\mathcal{E}|_{X_r}) \cap -$ として作用する双変類である。
この設定では、無限個の $p$ に対して $c_p(\mathcal{E})$ が非零となりうることに注意する。
したがって全 Chern 類は、完備化双変コホモロジー環の元
$$
c(\mathcal{E}) =
c_0(\mathcal{E}) + c_1(\mathcal{E}) + c_2(\mathcal{E}) + \ldots
\in A^*(X)^\wedge
$$
である。注意 \ref{chow-remark-completion-bivariant} を参照。
この設定で $\mathcal{E}$ の「階数」を元 $r(\mathcal{E}) \in A^0(X)$ として定義する。
これは $(\alpha_r) \in \prod \CH_*(X'_r)$ を
$(r\alpha_r) \in \prod \CH_*(X'_r)$ に送る双変演算である。
$c_p(\mathcal{E})$ と $r(\mathcal{E})$ が依然として $A^*(X)$ の中心に属することに注意する。
\end{remark}

\begin{remark}
\label{chow-remark-top-chern-class}
$(S, \delta)$ を状況 \ref{chow-situation-setup} のとおりとし、$X$ を $S$ 上局所有限型とする。
$\mathcal{E}$ を有限局所自由 $\mathcal{O}_X$-加群とする。一般に、注意
\ref{chow-remark-extend-to-finite-locally-free} のとおり
$X = \coprod X_r$ と書く。
空でない $X_r$ が有限個しかなければ
$$
c_{top}(\mathcal{E}) = \sum\nolimits_r c_r(\mathcal{E}|_{X_r})
\in A^*(X) = \bigoplus A^*(X_r)
$$
とおける。ここで等式は補題 \ref{chow-lemma-disjoint-decomposition-bivariant} による。
空でない $X_r$ が無限個ある場合も、同じ記号で
$$
c_{top}(\mathcal{E}) = \prod c_r(\mathcal{E}|_{X_r})
\in \prod A^r(X_r) \subset A^*(X)^\wedge
$$
を表す。記号については注意 \ref{chow-remark-completion-bivariant} を参照。
\end{remark}











\section{Chern 類間の多項式関係}
\label{chow-section-relations-chern-classes}

\noindent
$(S, \delta)$ を状況 \ref{chow-situation-setup} のとおりとし、$X$ を $S$ 上局所有限型とする。
$\mathcal{E}_i$ を $X$ 上の有限局所自由層の有限族とする。
補題 \ref{chow-lemma-cap-commutative-chern} により、Chern 類
$$
c_j(\mathcal{E}_i) \in A^*(X)
$$
は Chow コホモロジー代数 $A^*(X)$ の可換な（実際には中心的な）
$\mathbf{Z}$-部分代数を生成する。したがって、これらの Chern 類の多項式が
零であること、または二つの多項式が等しいことの意味を定められる。
例えば
$c_1(\mathcal{E}_1)^5 + c_2(\mathcal{E}_2)c_3(\mathcal{E}_3) = 0$
とは、すべての局所有限型射 $f : Y \to X$ に対して演算
$$
\CH_k(Y) \longrightarrow \CH_{k - 5}(Y), \quad
\alpha \longmapsto
c_1(\mathcal{E}_1)^5 \cap \alpha +
c_2(\mathcal{E}_2) \cap c_3(\mathcal{E}_3) \cap \alpha
$$
が零であることを意味する。補題 \ref{chow-lemma-bivariant-zero} により、これは、
$Y$ が $S$ 上局所有限型な整スキームである任意の射
$f : Y \to X$ に対してサイクル
$$
c_1(\mathcal{E}_1)^5 \cap [Y] +
c_2(\mathcal{E}_2) \cap c_3(\mathcal{E}_3) \cap [Y]
$$
が $\CH_{\dim(Y) - 5}(Y)$ で零であることと同値である。

\medskip\noindent
具体例として、補題 \ref{chow-lemma-c1-cap-additive} で証明した関係
$$
c_1(\mathcal{L} \otimes_{\mathcal{O}_X} \mathcal{N})
=
c_1(\mathcal{L}) + c_1(\mathcal{N})
$$
がある。より一般に、任意の局所自由層を可逆層でテンソルすると次のようになる。

\begin{lemma}
\label{chow-lemma-chern-classes-E-tensor-L}
$(S, \delta)$ を状況 \ref{chow-situation-setup} のとおりとし、$X$ を $S$ 上局所有限型とする。
$\mathcal{E}$ を $X$ 上の階数 $r$ の有限局所自由層とし、
$\mathcal{L}$ を $X$ 上の可逆層とする。このとき
\begin{equation}
\label{chow-equation-twist}
c_i({\mathcal E} \otimes {\mathcal L})
=
\sum\nolimits_{j = 0}^i
\binom{r - i + j}{j} c_{i - j}({\mathcal E}) c_1({\mathcal L})^j
\end{equation}
が $A^*(X)$ で成り立つ。
\end{lemma}

\begin{proof}
これは任意の三つ組 $(X, \mathcal{E}, \mathcal{L})$ に対して成り立つべき関係である。
特に $X$ が整の場合に成り立つべきであり、補題
\ref{chow-lemma-bivariant-zero} により、そのような $X$ に対して
$[X]$ とのキャップ積をとったときに成り立つことを示せば十分である。
そこで $X$ は整と仮定する。$\mathcal{E}$, resp.\ $\mathcal{E} \otimes \mathcal{L}$ に付随する
射影空間束を $(\pi : P \to X, \mathcal{O}_P(1))$, resp.\ $(\pi' : P' \to X, \mathcal{O}_{P'}(1))$
とする。標準射
$$
\xymatrix{
P \ar[rd]_\pi \ar[rr]_g & & P' \ar[ld]^{\pi'} \\
& X &
}
$$
を考える。『スキームの構成』の補題 \ref{constructions-lemma-twisting-and-proj} を参照。この射は
$g^*\mathcal{O}_{P'}(1)
= \mathcal{O}_P(1) \otimes \pi^* {\mathcal L}$
を満たす。したがって
$$
\sum\nolimits_{i = 0}^r
(-1)^i
(\xi + x)^i \cap \pi^*(c_{r - i}(\mathcal{E} \otimes \mathcal{L}) \cap [X])
=
0
$$
が $\CH_*(P)$ で成り立つ。ここで $\xi$ は
$c_1(\mathcal{O}_P(1))$ を、$x$ は $c_1(\pi^*\mathcal{L})$ を表す。
初等的な代数計算により、これは
$$
\sum\nolimits_{i = 0}^r
(-1)^i \xi^i \left(
\sum\nolimits_{j = i}^r
(-1)^{j - i}
\binom{j}{i}
x^{j - i} \cap
\pi^*(c_{r - j}(\mathcal{E} \otimes \mathcal{L}) \cap [X])
\right)
=
0
$$
と同値である。式 (\ref{chow-equation-chern-classes}) と比較すると
$$
c_{r - i}(\mathcal{E}) \cap [X] =
\sum\nolimits_{j = i}^r
\binom{j}{i}
(-c_1(\mathcal{L}))^{j - i} \cap
c_{r - j}(\mathcal{E} \otimes \mathcal{L}) \cap [X]
$$
が従う。これを変形し、負号を除いて添字を付け直すと、所望の関係を得る。
\end{proof}

\noindent
(\ref{chow-equation-twist}) の具体例は
\begin{align*}
c_1(\mathcal{E} \otimes \mathcal{L})
& =
c_1(\mathcal{E}) +
r c_1(\mathcal{L}) \\
c_2(\mathcal{E} \otimes \mathcal{L})
& =
c_2(\mathcal{E}) +
(r - 1) c_1(\mathcal{E}) c_1(\mathcal{L}) +
\binom{r}{2} c_1(\mathcal{L})^2 \\
c_3(\mathcal{E} \otimes \mathcal{L})
& =
c_3(\mathcal{E}) +
(r - 2) c_2(\mathcal{E})c_1(\mathcal{L}) +
\binom{r - 1}{2} c_1(\mathcal{E})c_1(\mathcal{L})^2 +
\binom{r}{3} c_1(\mathcal{L})^3
\end{align*}
である。








\section{Chern 類の加法性}
\label{chow-section-additivity-chern-classes}

\noindent
以下の準備補題はいずれも、最後の結果から直ちに従う。

\begin{lemma}
\label{chow-lemma-get-rid-of-trivial-subbundle}
$(S, \delta)$ を状況 \ref{chow-situation-setup} のとおりとする。
$X$ を $S$ 上局所有限型とする。$\mathcal{E}$, $\mathcal{F}$ を $X$ 上の
それぞれ階数 $r$, $r - 1$ の有限局所自由層とし、短完全列
$$
0 \to \mathcal{O}_X \to \mathcal{E} \to \mathcal{F} \to 0
$$
に入ると仮定する。このとき $A^*(X)$ において
$$
c_r(\mathcal{E}) = 0, \quad
c_j(\mathcal{E}) = c_j(\mathcal{F}), \quad j = 0, \ldots, r - 1
$$
が成り立つ。
\end{lemma}

\begin{proof}
補題 \ref{chow-lemma-bivariant-zero} により、$X$ が整のとき
$c_j(\mathcal{E}) \cap [X] = c_j(\mathcal{F}) \cap [X]$ であることを示せば十分である。
$\mathcal{E}$, resp.\ $\mathcal{F}$ に付随する射影空間束を
$(\pi : P \to X, \mathcal{O}_P(1))$, resp.\ $(\pi' : P' \to X, \mathcal{O}_{P'}(1))$
と表す。全射 $\mathcal{E} \to \mathcal{F}$ から $X$ 上の閉埋め込み
$$
i : P' \longrightarrow P
$$
が得られる。さらに、元
$1 \in \Gamma(X, \mathcal{O}_X) \subset \Gamma(X, \mathcal{E})$ から大域切断
$s \in \Gamma(P, \mathcal{O}_P(1))$ が得られ、その零点集合はちょうど $P'$ である。
したがって $P'$ は $P$ 上の有効 Cartier 因子であり、
$\mathcal{O}_P(P') \cong \mathcal{O}_P(1)$ を満たす。ゆえに補題
\ref{chow-lemma-relative-effective-cartier} により、$X$ 上の任意のサイクル類 $\alpha$ に対して
$$
c_1(\mathcal{O}_P(1)) \cap \pi^*\alpha = i_*((\pi')^*\alpha)
$$
である。補題 \ref{chow-lemma-determine-intersections} により、
$\alpha_j = c_j(\mathcal{F}) \cap [X]$, $j = 0, \ldots, r - 1$ は
$$
\sum\nolimits_{j = 0}^{r - 1} (-1)^jc_1(\mathcal{O}_{P'}(1))^j
\cap (\pi')^*\alpha_j = 0
$$
を満たす。これを $P$ へ押し出し、上の注意と補題
\ref{chow-lemma-pushforward-cap-c1} を用いると
$$
\sum\nolimits_{j = 0}^{r - 1}
(-1)^j c_1(\mathcal{O}_P(1))^{j + 1}
\cap \pi^*\alpha_j = 0
$$
を得る。補題 \ref{chow-lemma-determine-intersections} の一意性により、
$c_r(\mathcal{E}) \cap [X] = 0$ および
$c_j(\mathcal{E}) \cap [X] = \alpha_j = c_j(\mathcal{F}) \cap [X]$
が $j = 0, \ldots, r - 1$ に対して成り立つ。したがって補題が従う。
\end{proof}

\begin{lemma}
\label{chow-lemma-additivity-invertible-subsheaf}
$(S, \delta)$ を状況 \ref{chow-situation-setup} のとおりとする。
$X$ を $S$ 上局所有限型とする。$\mathcal{E}$, $\mathcal{F}$ を $X$ 上の
それぞれ階数 $r$, $r - 1$ の有限局所自由層とし、短完全列
$$
0 \to \mathcal{L} \to \mathcal{E} \to \mathcal{F} \to 0
$$
に入ると仮定する。ただし $\mathcal{L}$ は可逆層とする。このとき $A^*(X)$ において
$$
c(\mathcal{E}) = c(\mathcal{L}) c(\mathcal{F})
$$
が成り立つ。
\end{lemma}

\begin{proof}
この関係は実質的に
$c_i(\mathcal{E}) = c_i(\mathcal{F}) + c_1(\mathcal{L})c_{i - 1}(\mathcal{F})$
というだけである。補題 \ref{chow-lemma-get-rid-of-trivial-subbundle} により、
$j = 0, \ldots, r$ に対して
$c_j(\mathcal{E} \otimes \mathcal{L}^{\otimes -1})
= c_j(\mathcal{F} \otimes \mathcal{L}^{\otimes -1})$ である。ここで慣例により
$c_r(\mathcal{F} \otimes \mathcal{L}^{-1}) = 0$ とおいた。
補題 \ref{chow-lemma-chern-classes-E-tensor-L} を適用すると
$$
\sum_{j = 0}^i
\binom{r - i + j}{j} (-1)^j c_{i - j}({\mathcal E}) c_1({\mathcal L})^j
=
\sum_{j = 0}^i
\binom{r - 1 - i + j}{j} (-1)^j c_{i - j}({\mathcal F}) c_1({\mathcal L})^j
$$
を得る。
$c_i(\mathcal{E}) = c_i(\mathcal{F}) + c_1(\mathcal{L})c_{i - 1}(\mathcal{F})$
とおけばこの方程式の「解」が得られる。これが唯一の可能な解であることを示せば
補題が従うが、その検証は省略する。
\end{proof}

\begin{lemma}
\label{chow-lemma-additivity-chern-classes}
$(S, \delta)$ を状況 \ref{chow-situation-setup} のとおりとする。
$X$ を $S$ 上局所有限型なスキームとする。${\mathcal E}$ が、
それぞれ階数 $r_i$ の有限局所自由層 $\mathcal{E}_i$ からなる完全列
$$
0
\to
{\mathcal E}_1
\to
{\mathcal E}
\to
{\mathcal E}_2
\to
0
$$
に入ると仮定する。全 Chern 類は $A^*(X)$ において
$$
c({\mathcal E}) = c({\mathcal E}_1) c({\mathcal E}_2)
$$
を満たす。
\end{lemma}

\begin{proof}
補題 \ref{chow-lemma-bivariant-zero} により $X$ は整と仮定でき、
$[X]$ とのキャップ積をとったときにこの等式を示せばよい。
$r_1$ に関する帰納法による。$r_1 = 1$ の場合は補題
\ref{chow-lemma-additivity-invertible-subsheaf} である。$r_1 > 1$ と仮定する。
$(\pi : P \to X, \mathcal{O}_P(1))$ を $\mathcal{E}_1$ に付随する射影空間束とする。
次に注意する。
\begin{enumerate}
\item $\pi^* : \CH_*(X) \to \CH_*(P)$ は単射である。
\item $\pi^*\mathcal{E}_1$ は短完全列
$0 \to \mathcal{F} \to \pi^*\mathcal{E}_1 \to \mathcal{L} \to 0$
に入る。ただし $\mathcal{L}$ は可逆である。
\end{enumerate}
第一の主張は射影空間束公式から、第二の主張は射影空間束の定義から従う。
(実際 $\mathcal{L} = \mathcal{O}_P(1)$ である。)
$Q = \pi^*\mathcal{E}/\mathcal{F}$ とおくと、これは完全列
$0 \to \mathcal{L} \to Q \to \pi^*\mathcal{E}_2 \to 0$ に入る。
帰納法により
\begin{eqnarray*}
c(\pi^*\mathcal{E}) \cap [P]
& = &
c(\mathcal{F}) \cap c(\pi^*\mathcal{E}/\mathcal{F}) \cap [P] \\
& = &
c(\mathcal{F}) \cap c(\mathcal{L}) \cap c(\pi^*\mathcal{E}_2) \cap [P] \\
& = &
c(\pi^*\mathcal{E}_1) \cap c(\pi^*\mathcal{E}_2) \cap [P]
\end{eqnarray*}
である。$[P] = \pi^*[X]$ なので、補題 \ref{chow-lemma-flat-pullback-cap-cj} から結論を得る。
\end{proof}

\begin{lemma}
\label{chow-lemma-chern-filter-by-linebundles}
$(S, \delta)$ を状況 \ref{chow-situation-setup} のとおりとする。
$X$ を $S$ 上局所有限型とする。${\mathcal L}_i$, $i = 1, \ldots, r$ を
$X$ 上の可逆 $\mathcal{O}_X$-加群とする。$\mathcal{E}$ を、フィルトレーション
$$
0 = \mathcal{E}_0 \subset \mathcal{E}_1 \subset \mathcal{E}_2
\subset \ldots \subset \mathcal{E}_r = \mathcal{E}
$$
を備えた局所自由 $\mathcal{O}_X$-加群とし、
$\mathcal{E}_i/\mathcal{E}_{i - 1} \cong \mathcal{L}_i$ を仮定する。
$c_1({\mathcal L}_i) = x_i$ とおく。このとき $A^*(X)$ において
$$
c(\mathcal{E})
=
\prod\nolimits_{i = 1}^r (1 + x_i)
$$
である。
\end{lemma}

\begin{proof}
補題 \ref{chow-lemma-additivity-invertible-subsheaf} と帰納法を用いる。
\end{proof}




\section{零サイクルの次数}
\label{chow-section-degree-zero-cycles}

\noindent
体上の固有スキーム上の零サイクルの次数を定義する。
一つの方法は、補題 \ref{chow-lemma-spell-out-degree-zero-cycle} のように直接定義し、
補題 \ref{chow-lemma-curve-principal-divisor} によって良定義性を示すことである。
ここでは代わりに次のように定義する。

\begin{definition}
\label{chow-definition-degree-zero-cycle}
$k$ を体とする（例 \ref{chow-example-field}）。$p : X \to \Spec(k)$ を固有とする。
$X$ 上の {\it 零サイクルの次数} は、固有押し出し
$$
p_* : \CH_0(X) \to \CH_0(\Spec(k))
$$
（補題 \ref{chow-lemma-proper-pushforward-rational-equivalence}）と、
$[\Spec(k)]$ を $1$ に写す自然な同型 $\CH_0(\Spec(k)) = \mathbf{Z}$
との合成によって与えられる。記法は $\deg(\alpha)$ とする。
\end{definition}

\noindent
この定義をより具体的に述べよう。

\begin{lemma}
\label{chow-lemma-spell-out-degree-zero-cycle}
$k$ を体、$X$ を $k$ 上固有とする。
$\alpha = \sum n_i[Z_i]$ を $Z_0(X)$ の元とする。このとき
$$
\deg(\alpha) = \sum n_i\deg(Z_i)
$$
である。ただし $\deg(Z_i)$ は $Z_i \to \Spec(k)$ の次数、すなわち
$\deg(Z_i) = \dim_k \Gamma(Z_i, \mathcal{O}_{Z_i})$ である。
\end{lemma}

\begin{proof}
これは固有押し出しの定義（定義 \ref{chow-definition-proper-pushforward}）である。
\end{proof}

\noindent
次に、Varieties の節 \ref{varieties-section-divisors-curves} で定義した、
体上の $1$ 次元固有スキーム上のベクトル束の次数との関係を与える。

\begin{lemma}
\label{chow-lemma-degree-vector-bundle}
$k$ を体とし、$X$ を次元 $\leq 1$ の $k$ 上固有なスキームとする。
$\mathcal{E}$ を一定階数の有限局所自由 $\mathcal{O}_X$-加群とする。このとき
$$
\deg(\mathcal{E}) = \deg(c_1(\mathcal{E}) \cap [X]_1)
$$
である。左辺は Varieties の定義
\ref{varieties-definition-degree-invertible-sheaf} で定義される。
\end{lemma}

\begin{proof}
$C_i \subset X$, $i = 1, \ldots, t$ を、誘導される既約スキーム構造を入れた
次元 $1$ の既約成分とし、$m_i$ を $X$ における $C_i$ の重複度とする。
このとき $[X]_1 = \sum m_i[C_i]$ であり、
$c_1(\mathcal{E}) \cap [X]_1$ はサイクル
$m_i c_1(\mathcal{E}|_{C_i}) \cap [C_i]$ の押し出しの和である。
Varieties の補題 \ref{varieties-lemma-degree-in-terms-of-components} により
$\mathcal{E}$ の次数にも同様の分解があるので、$X$ が $k$ 上の固有曲線の場合を
示せば十分である。

\medskip\noindent
$X$ を $k$ 上の固有曲線と仮定する。Divisors の補題
\ref{divisors-lemma-filter-after-modification} により、$f^*\mathcal{E}$ が逐次商を
可逆 $\mathcal{O}_{X'}$-加群とするフィルトレーションをもつような改変
$f : X' \to X$ が存在する。$f_*[X']_1 = [X]_1$ なので、補題
\ref{chow-lemma-pushforward-cap-cj} から
$$
\deg(c_1(\mathcal{E}) \cap [X]_1) = \deg(c_1(f^*\mathcal{E}) \cap [X']_1)
$$
を得る。Varieties の補題 \ref{varieties-lemma-degree-birational-pullback} により
次数にも同様の関係があるので、$\mathcal{E}$ が逐次商を可逆
$\mathcal{O}_X$-加群とするフィルトレーションをもつ場合に帰着する。
この場合、次数の加法性（Varieties の補題 \ref{varieties-lemma-degree-additive}）と
第一 Chern 類の加法性（補題 \ref{chow-lemma-additivity-chern-classes}）を用い、
次段落の場合に帰着できる。

\medskip\noindent
$X$ を $k$ 上の固有曲線とし、$\mathcal{E}$ を可逆 $\mathcal{O}_X$-加群と仮定する。
Divisors の補題
\ref{divisors-lemma-quasi-projective-Noetherian-pic-effective-Cartier} により、
$\mathcal{E}$ は $X$ 上のある有効 Cartier 因子 $D, D'$ を用いて
$\mathcal{O}_X(D) \otimes \mathcal{O}_X(D')^{\otimes -1}$ と同型になる（ここで $X$ が射影的であることも
用いた。例えば Varieties の補題 \ref{varieties-lemma-dim-1-proper-projective} を見よ）。
可逆層のテンソル積に関する次数の加法性
（Varieties の補題 \ref{varieties-lemma-degree-tensor-product}）と、
可逆層のテンソル積に関する $c_1$ の加法性
（補題 \ref{chow-lemma-c1-cap-additive} または \ref{chow-lemma-chern-classes-E-tensor-L}）により、
$\mathcal{E} = \mathcal{O}_X(D)$ の場合に帰着する。この場合、左辺は
$\deg(D)$（Varieties の補題 \ref{varieties-lemma-degree-effective-Cartier-divisor}）となり、
右辺は補題 \ref{chow-lemma-geometric-cap} により $\deg([D]_0)$ となる。
定義により
$$
[D]_0 = \sum\nolimits_{x \in D}
\text{length}_{\mathcal{O}_{X, x}}(\mathcal{O}_{D, x}) [x] =
\sum\nolimits_{x \in D}
\text{length}_{\mathcal{O}_{D, x}}(\mathcal{O}_{D, x}) [x]
$$
であるから、
$$
\deg([D]_0) = \sum\nolimits_{x \in D}
\text{length}_{\mathcal{O}_{D, x}}(\mathcal{O}_{D, x}) [\kappa(x) : k] =
\dim_k \Gamma(D, \mathcal{O}_D) = \deg(D)
$$
を得る。最後から二つ目の等式には、$D$ がアフィンであることと
Algebra の補題 \ref{algebra-lemma-pushdown-module} を用いた。
\end{proof}

\noindent
最後に、これまでの結果を Varieties の節 \ref{varieties-section-num} で定義した
数値的交点積と結びつける。

\begin{lemma}
\label{chow-lemma-degrees-and-numerical-intersections}
$k$ を体とし、$X$ を $k$ 上固有なスキームとする。
$Z \subset X$ を次元 $d$ の閉部分スキームとし、
$\mathcal{L}_1, \ldots, \mathcal{L}_d$ を可逆 $\mathcal{O}_X$-加群とする。このとき
$$
(\mathcal{L}_1 \cdots \mathcal{L}_d \cdot Z) =
\deg(
c_1(\mathcal{L}_1) \cap \ldots \cap c_1(\mathcal{L}_d) \cap [Z]_d)
$$
である。左辺は Varieties の定義
\ref{varieties-definition-intersection-number} で定義される。特に、
$\mathcal{L}$ が豊富な可逆 $\mathcal{O}_X$-加群ならば
$$
\deg_\mathcal{L}(Z) = \deg(c_1(\mathcal{L})^d \cap [Z]_d)
$$
である。
\end{lemma}

\begin{proof}
$d$ に関する帰納法で示す。$d = 0$ なら、結果は Varieties の補題
\ref{varieties-lemma-chi-tensor-finite} により成り立つ。$d > 0$ と仮定する。

\medskip\noindent
$Z_i \subset Z$, $i = 1, \ldots, t$ を、誘導される既約スキーム構造を入れた
次元 $d$ の既約成分とし、$m_i$ を $Z$ における $Z_i$ の重複度とする。
このとき $[Z]_d = \sum m_i[Z_i]$ であり、
$c_1(\mathcal{L}_1) \cap \ldots \cap c_1(\mathcal{L}_d) \cap [Z]_d$ はサイクル
$m_i c_1(\mathcal{L}_1) \cap \ldots \cap c_1(\mathcal{L}_d) \cap [Z_i]$ の和である。
Varieties の補題 \ref{varieties-lemma-numerical-polynomial-leading-term} により
$(\mathcal{L}_1 \cdots \mathcal{L}_d \cdot Z)$ にも同様の分解があるので、
$Z = X$ が $k$ 上の次元 $d$ の固有多様体の場合を示せば十分である。

\medskip\noindent
Chow の補題により、$Y$ が $k$ 上 H-射影的となる双有理固有射
$f : Y \to X$ が存在する。Cohomology of Schemes の補題
\ref{coherent-lemma-chow-Noetherian} および注意
\ref{coherent-remark-chow-Noetherian} を見よ。このとき Varieties の補題
\ref{varieties-lemma-intersection-number-and-pullback} により
$$
(f^*\mathcal{L}_1 \cdots f^*\mathcal{L}_d \cdot Y) =
(\mathcal{L}_1 \cdots \mathcal{L}_d \cdot X)
$$
であり、また補題 \ref{chow-lemma-pushforward-cap-c1} により
$$
f_*(c_1(f^*\mathcal{L}_1) \cap \ldots \cap c_1(f^*\mathcal{L}_d) \cap [Y]) =
c_1(\mathcal{L}_1) \cap \ldots \cap c_1(\mathcal{L}_d) \cap [X]
$$
である。したがって $X$ を $Y$ で置き換え、$X$ は $k$ 上射影的と仮定できる。

\medskip\noindent
$X$ が次元 $d$ の固有射影多様体ならば、Divisors の補題
\ref{divisors-lemma-quasi-projective-Noetherian-pic-effective-Cartier} により、
ある有効 Cartier 因子 $D, D' \subset X$ に対して
$\mathcal{L}_1 = \mathcal{O}_X(D) \otimes \mathcal{O}_X(D')^{\otimes -1}$ と書ける。
等式の両辺の加法性（Varieties の補題
\ref{varieties-lemma-intersection-number-additive} および補題 \ref{chow-lemma-c1-cap-additive}）により、
$\mathcal{L}_1 = \mathcal{O}_X(D)$ で $D$ が有効 Cartier 因子の場合に帰着する。
Varieties の補題
\ref{varieties-lemma-numerical-intersection-effective-Cartier-divisor} により
$$
(\mathcal{L}_1 \cdots \mathcal{L}_d \cdot X) =
(\mathcal{L}_2 \cdots \mathcal{L}_d \cdot D)
$$
であり、また補題 \ref{chow-lemma-geometric-cap} により
$$
c_1(\mathcal{L}_1) \cap \ldots \cap c_1(\mathcal{L}_d) \cap [X] =
c_1(\mathcal{L}_2) \cap \ldots \cap c_1(\mathcal{L}_d) \cap [D]_{d - 1}
$$
である。ゆえに帰納法の仮定から結論を得る。
\end{proof}








\section{指定された余次元のサイクル}
\label{chow-section-cycles-codimension}

\noindent
場合によっては、全サイクルのアーベル群に余次元による第二の次数付けがある。

\begin{lemma}
\label{chow-lemma-locally-equidimensional}
$(S, \delta)$ を状況 \ref{chow-situation-setup} のとおりとし、$X$ を $S$ 上局所有限型とする。
節 \ref{chow-section-setup} のとおり $\delta = \delta_{X/S}$ と書く。次の条件は同値である。
\begin{enumerate}
\item $X = \coprod_{n \in \mathbf{Z}} X_n$ という開かつ閉な部分スキームへの分解が存在し、
$\xi \in X_n$ が $X_n$ の既約成分の生成点なら常に $\delta(\xi) = n$ である。
\item 任意の $x \in X$ に対して、$x$ の開近傍 $U \subset X$ と整数 $n$ が存在し、
$\xi \in U$ が $U$ の既約成分の生成点なら常に $\delta(\xi) = n$ である。
\item 任意の $x \in X$ に対して整数 $n_x$ が存在し、$x$ を含む $X$ の既約成分の
任意の生成点 $\xi$ に対して $\delta(\xi) = n_x$ である。
\end{enumerate}
$X$ が正規または Cohen--Macaulay ならば、これらの条件は成り立つ
\footnote{実際には $X$ が $(S_2)$ であれば十分である。Local Cohomology の補題
\ref{local-cohomology-lemma-catenary-S2-equidimensional} と比較せよ。}。
\end{lemma}

\begin{proof}
(1) $\Rightarrow$ (2) $\Rightarrow$ (3) は明らかである。
逆に (3) が成り立つなら、$X_n = \{x \in X \mid n_x = n\}$ とおけば (1) のような分解を得る。
実際、与えられた $x$ に対して、$x$ を通る $X$ の既約成分の和から
$x$ を通らない $X$ の既約成分の和を除いたものが $x$ の開近傍になるので、$X_n$ は開である。
$X$ が正規なら、$X$ は整スキームの互いに交わらない和である
（Properties の補題 \ref{properties-lemma-normal-locally-Noetherian}）から、条件は成り立つ。
$X$ が Cohen--Macaulay なら、
$\delta' : X \to \mathbf{Z}$, $x \mapsto -\dim(\mathcal{O}_{X, x})$ は $X$ 上の次元関数である
（例 \ref{chow-example-CM-irreducible} を見よ）。$\delta - \delta'$ は局所定数であり
（Topology の補題 \ref{topology-lemma-dimension-function-unique}）、$X$ の任意の生成点 $\xi$ に対して
$\delta'(\xi) = 0$ なので、(2) が成り立つ。
\end{proof}

\noindent
$(S, \delta)$ を状況 \ref{chow-situation-setup} のとおりとする。$X$ を $S$ 上局所有限型で、
補題 \ref{chow-lemma-locally-equidimensional} の同値な条件を満たすものとする。
整閉部分スキーム $Z \subset X$ に対し、$X$ における $Z$ の余次元
$\text{codim}(Z, X)$ が定まる。Topology の定義 \ref{topology-definition-codimension} を見よ。
$X$ 上のサイクル $\alpha = \sum n_Z[Z]$ が
$n_Z \not = 0 \Rightarrow \text{codim}(Z, X) = p$ を満たすとき、これを
{\it 余次元 $p$ のサイクル} という。余次元 $p$ のサイクル全体のアーベル群を
$Z^p(X)$ と表す。$X = \coprod X_n$ を補題
\ref{chow-lemma-locally-equidimensional} の (1) で与えられる分解とする。
サイクルは局所有限和として定義したので、明らかに
$$
Z^p(X) = \prod\nolimits_n Z_{n - p}(X_n)
$$
である。さらに $\prod_p Z^p(X) = \prod_k Z_k(X)$ である。
ここで、$X$ 上の余次元 $p$ のサイクルの有理同値を前とまったく同じように定義し、
補題 \ref{chow-lemma-locally-equidimensional} の $X$ に対する指定された余次元のサイクルの理論全体を
初めから構成し直すこともできる。しかし代わりに、
{\it 余次元 $p$ のサイクルの Chow 群} を単に
$$
\CH^p(X) = \prod\nolimits_n \CH_{n - p}(X_n)
$$
と定義する。前と同様に $\prod_p \CH^p(X) = \prod_k \CH_k(X)$ である。
$X$ が準コンパクトなら、式中の積は有限積（したがって直和）であり、
$\bigoplus_p \CH^p(X) = \bigoplus_k \CH_k(X)$ を得る。$X$ が準コンパクトかつ有限次元なら、
これらの群のうち非零なものは有限個しかない。

\medskip\noindent
上で証明した Chow 群に関する多くの構成と結果は、Chow 群 $\CH^*(X)$ に対する
自然な対応物をもつ。それぞれは、関係するスキームを補題
\ref{chow-lemma-locally-equidimensional} のような「等次元」の部分に分解し、
上の積分解の各因子に既に証明した結果を適用すれば示される。
そのうちいくつかを列挙する。
\begin{enumerate}
\item $f : X \to Y$ が $S$ 上局所有限型なスキームの平坦射で、$X$, $Y$ が補題
\ref{chow-lemma-locally-equidimensional} の同値な条件を満たすなら、平坦引き戻しは写像
$$
f^* : \CH^p(Y) \to \CH^p(X)
$$
を定める。
\item $f : X \to Y$ が $S$ 上局所有限型なスキームの射で、$X$, $Y$ が補題
\ref{chow-lemma-locally-equidimensional} の同値な条件を満たすとする。既約成分の任意の組
$Z \subset X$, $W \subset Y$ が $f(Z) \subset W$ を満たすとき
$\dim_\delta(W) - \dim_\delta(Z) = r$ となるなら、$f$ は {\it 余次元}
$r \in \mathbf{Z}$ をもつという。
\item $f : X \to Y$ が $S$ 上局所有限型なスキームの固有射で、$X$, $Y$ が補題
\ref{chow-lemma-locally-equidimensional} の同値な条件を満たし、$f$ が余次元 $r$ をもつなら、
固有押し出しは写像
$$
f_* : \CH^p(X) \to \CH^{p + r}(Y)
$$
である。
\item $f : X \to Y$ が $S$ 上局所有限型なスキームの射で、$X$, $Y$ が補題
\ref{chow-lemma-locally-equidimensional} の同値な条件を満たし、$f$ が余次元 $r$ をもち、
$c \in A^q(X \to Y)$ なら、$c$ は写像
$$
c \cap -  : \CH^p(Y) \to \CH^{p + q - r}(X)
$$
を誘導する。
\item $X$ が $S$ 上局所有限型で補題
\ref{chow-lemma-locally-equidimensional} の同値な条件を満たすスキームであり、
$\mathcal{L}$ が可逆 $\mathcal{O}_X$-加群なら、
$$
c_1(\mathcal{L}) \cap - : \CH^p(X) \to \CH^{p + 1}(X)
$$
である。
\item $X$ が $S$ 上局所有限型で補題
\ref{chow-lemma-locally-equidimensional} の同値な条件を満たすスキームであり、
$\mathcal{E}$ が有限局所自由 $\mathcal{O}_X$-加群なら、
$$
c_i(\mathcal{E}) \cap - : \CH^p(X) \to \CH^{p + i}(X)
$$
である。
\end{enumerate}
注意：射が余次元 $r$ をもつという性質は基底変換で保たれない。

\begin{remark}
\label{chow-remark-fundamental-class}
$(S, \delta)$ を状況 \ref{chow-situation-setup} のとおりとする。$X$ を $S$ 上局所有限型で、
補題 \ref{chow-lemma-locally-equidimensional} の同値な条件を満たすものとする。
$X = \coprod X_n$ を開かつ閉な部分スキームへの分解で、$X_n$ の各既約成分が
$\delta$-次元 $n$ をもつものとする。この状況では、ときどき
$$
[X] = \sum\nolimits_n [X_n]_n \in \CH^0(X)
$$
とおく。この類は Chow 理論における $X$ の一種の「基本類」である。
\end{remark}




\section{分解原理}
\label{chow-section-splitting-principle}

\noindent
この設定で分解原理が正確に何を主張するかを述べるのはそれほど容易ではない。
一つの定式化は次のとおりである。

\begin{lemma}
\label{chow-lemma-splitting-principle}
$(S, \delta)$ を状況 \ref{chow-situation-setup} のとおりとし、$X$ を $S$ 上局所有限型とする。
$\mathcal{E}_i$ を階数 $r_i$ の局所自由 $\mathcal{O}_X$-加群の有限族とする。
次を満たす相対次元 $d$ の射影的平坦射 $\pi : P \to X$ が存在する。
\begin{enumerate}
\item 任意の射 $f : Y \to X$ に対して、写像
$\pi_Y^* : \CH_*(Y) \to \CH_{* + d}(Y \times_X P)$ は単射である。
\item 各 $\pi^*\mathcal{E}_i$ は、逐次商 $\mathcal{L}_{i, 1}, \ldots, \mathcal{L}_{i, r_i}$ が
可逆 ${\mathcal O}_P$-加群となるフィルトレーションをもつ。
\end{enumerate}
さらに、(1) が成り立つとき、制限写像 $A^*(X) \to A^*(P)$
（注意 \ref{chow-remark-pullback-cohomology}）は単射である。
\end{lemma}

\begin{proof}
すべての $i$ に対して $r_i \geq 1$ と仮定できる。$\sum (r_i - 1)$ に関する帰納法で示す。
この整数が $0$ なら、すべての $i$ に対して $\mathcal{E}_i$ は可逆であり、
$\pi = \text{id}_X$ ととればよい。そうでないなら、$r_i > 1$ となる $i$ を一つ選び、射
$\pi_i : P_i = \mathbf{P}(\mathcal{E}_i) \to X$ を考える。
階数がそれぞれ $r_i - 1$, $r_i$, $1$ である有限局所自由 $\mathcal{O}_{P_i}$-加群の短完全列
$$
0 \to \mathcal{F} \to \pi_i^*\mathcal{E}_i \to \mathcal{O}_{P_i}(1) \to 0
$$
がある。射影束公式（補題 \ref{chow-lemma-chow-ring-projective-bundle}）により、
任意の基底変換の後に $\pi_i^*$ は Chow 群上で単射である。
$\mathcal{F}$ および $i' \not = i$ に対する $\pi_{i'}^*\mathcal{E}_{i'}$ という
有限局所自由 $\mathcal{O}_{P_i}$-加群に帰納法の仮定を適用すると、
補題で述べた性質をもつ射 $\pi : P \to P_i$ が得られる。このとき合成
$\pi_i \circ \pi : P \to X$ が求める射である。いくつかの詳細は省略する。
\end{proof}

\begin{remark}
\label{chow-remark-the-proof-shows-more}
補題 \ref{chow-lemma-splitting-principle} の証明は、射 $\pi : P \to X$ が次の追加性質をもつことも示す。
\begin{enumerate}
\item $\pi$ は、有限一定階数の局所自由加群に付随する射影空間束の有限合成である。
\item 任意の $\alpha \in \CH_k(X)$ に対して
$\alpha = \pi_*(\xi_1 \cap \ldots \cap \xi_d \cap \pi^*\alpha)$ である。ただし $\xi_i$ はある可逆
$\mathcal{O}_P$-加群の第一 Chern 類である。
\end{enumerate}
第二の注意は第一の注意と補題 \ref{chow-lemma-cap-projective-bundle} から従う。
必要に応じてさらなる注意を追加する。
\end{remark}

\noindent
$(S, \delta)$, $X$, $\mathcal{E}_i$ を補題 \ref{chow-lemma-splitting-principle} のとおりとする。
{\it 分解原理} とは、形式的に
$$
c(\mathcal{E}_i) = \prod (1 + x_{i, j})
$$
と書く慣用をいう。記号 $x_{i, 1}, \ldots, x_{i, r_i}$ を $\mathcal{E}_i$ の
{\it Chern 根} という。言い換えれば、$\mathcal{E}_i$ の第 $p$ Chern 類は Chern 根の
第 $p$ 基本対称関数である。分解原理の有用性は、$\mathcal{E}_i$ の Chern 類間の
多項式関係を示すには、Chern 根間の対応する関係を示せば十分であるという主張に由来する。

\medskip\noindent
実際、$\pi : P \to X$ を補題 \ref{chow-lemma-splitting-principle} のとおりとする。
標準的な $\mathbf{Z}$-代数写像 $\pi^* : A^*(X) \to A^*(P)$ があることを思い出そう。
注意 \ref{chow-remark-pullback-cohomology} を見よ。$X$ 上の任意の $Y$ に対する Chow 群上の
$\pi_Y^*$ の単射性から、写像 $\pi^* : A^*(X) \to A^*(P)$ の単射性が従う
（詳細は省略する）。補題 \ref{chow-lemma-chern-filter-by-linebundles} により
$$
\pi^*c(\mathcal{E}_i) = \prod (1 + c_1(\mathcal{L}_{i, j}))
$$
である。したがって、Chern 根 $x_{i, j}$ を元
$c_1(\mathcal{L}_{i, j}) \in A^*(P)$ とみなし、表示式の左辺に単射
$\pi^* : A^*(X) \to A^*(P)$ を適用した後、その等式が $A^*(P)$ で成り立つとみなせる。

\medskip\noindent
この方法の働きを見るには、いくつかの例を与えるのがよい。

\begin{lemma}
\label{chow-lemma-chern-classes-dual}
状況 \ref{chow-situation-setup} において、$X$ を $S$ 上局所有限型とする。
$\mathcal{E}$ を双対 $\mathcal{E}^\vee$ をもつ有限局所自由 $\mathcal{O}_X$-加群とする。このとき
$$
c_i(\mathcal{E}^\vee) = (-1)^i c_i(\mathcal{E})
$$
が $A^i(X)$ で成り立つ。
\end{lemma}

\begin{proof}
$\pi : P \to X$ を補題 \ref{chow-lemma-splitting-principle} のような射とする。
任意の基底変換後の $\pi^*$ の単射性により、$P$ へ引き戻した後に
$\mathcal{E}$ と $\mathcal{E}^\vee$ の Chern 類間の関係を示せば十分である。
したがって、可逆 $\mathcal{O}_X$-加群 ${\mathcal L}_i$, $i = 1, \ldots, r$ と、
フィルトレーション
$$
0 = \mathcal{E}_0 \subset \mathcal{E}_1 \subset \mathcal{E}_2
\subset \ldots \subset \mathcal{E}_r = \mathcal{E}
$$
が存在し、$\mathcal{E}_i/\mathcal{E}_{i - 1} \cong \mathcal{L}_i$ であると仮定できる。
このとき双対フィルトレーション
$$
0 = \mathcal{E}_r^\perp \subset \mathcal{E}_1^\perp \subset \mathcal{E}_2^\perp
\subset \ldots \subset \mathcal{E}_0^\perp = \mathcal{E}^\vee
$$
が得られ、$\mathcal{E}_{i - 1}^\perp/\mathcal{E}_i^\perp \cong
\mathcal{L}_i^{\otimes -1}$ である。$x_i = c_1(\mathcal{L}_i)$ とおく。
補題 \ref{chow-lemma-c1-cap-additive} により
$c_1(\mathcal{L}_i^{\otimes -1}) = - x_i$ である。
補題 \ref{chow-lemma-chern-filter-by-linebundles} により、$A^*(X)$ において
$$
c(\mathcal{E}) = \prod\nolimits_{i = 1}^r (1 + x_i)
\quad\text{and}\quad
c(\mathcal{E}^\vee) = \prod\nolimits_{i = 1}^r (1 - x_i)
$$
である。したがって形式的な計算から結論が従うが、その計算は省略する。
\end{proof}

\begin{lemma}
\label{chow-lemma-chern-classes-tensor-product}
状況 \ref{chow-situation-setup} において、$X$ を $S$ 上局所有限型とする。
$\mathcal{E}$ および $\mathcal{F}$ を、それぞれ階数 $r$, $s$ の有限局所自由
$\mathcal{O}_X$-加群とする。このとき
$$
c_1(\mathcal{E} \otimes \mathcal{F})
=
r c_1(\mathcal{F}) + s c_1(\mathcal{E})
$$
および
$$
c_2(\mathcal{E} \otimes \mathcal{F})
=
r c_2(\mathcal{F}) + s c_2(\mathcal{E}) +
{r \choose 2} c_1(\mathcal{F})^2 +
(rs - 1) c_1(\mathcal{F})c_1(\mathcal{E}) +
{s \choose 2} c_1(\mathcal{E})^2
$$
が成り立ち、$A^*(X)$ において以下も同様である。
\end{lemma}

\begin{proof}
補題 \ref{chow-lemma-chern-classes-dual} の証明とまったく同様に論じると、
可逆 $\mathcal{O}_X$-加群 ${\mathcal L}_i$, $i = 1, \ldots, r$,
${\mathcal N}_i$, $i = 1, \ldots, s$ とフィルトレーション
$$
0 = \mathcal{E}_0 \subset \mathcal{E}_1 \subset \mathcal{E}_2
\subset \ldots \subset \mathcal{E}_r = \mathcal{E}
\quad\text{and}\quad
0 = \mathcal{F}_0 \subset \mathcal{F}_1 \subset \mathcal{F}_2
\subset \ldots \subset \mathcal{F}_s = \mathcal{F}
$$
があり、$\mathcal{E}_i/\mathcal{E}_{i - 1} \cong \mathcal{L}_i$ および
$\mathcal{F}_j/\mathcal{F}_{j - 1} \cong \mathcal{N}_j$ と仮定できる。
組 $(i, j)$ を辞書式に順序付けると、フィルトレーション
$$
0 \subset \ldots \subset
\mathcal{E}_i \otimes \mathcal{F}_j
+
\mathcal{E}_{i - 1} \otimes \mathcal{F}
\subset \ldots \subset \mathcal{E} \otimes \mathcal{F}
$$
が得られ、その逐次商は
$$
\mathcal{L}_1 \otimes \mathcal{N}_1,
\mathcal{L}_1 \otimes \mathcal{N}_2,
\ldots,
\mathcal{L}_1 \otimes \mathcal{N}_s,
\mathcal{L}_2 \otimes \mathcal{N}_1,
\ldots,
\mathcal{L}_r \otimes \mathcal{N}_s
$$
である。補題 \ref{chow-lemma-chern-filter-by-linebundles} により、$A^*(X)$ において
$$
c(\mathcal{E}) = \prod (1 + x_i),
\quad
c(\mathcal{F}) = \prod (1 + y_j),
\quad\text{and}\quad
c(\mathcal{E} \otimes \mathcal{F}) = \prod (1 + x_i + y_j),
$$
である。したがって形式的な計算から結論が従うが、その計算は省略する。
\end{proof}

\begin{remark}
\label{chow-remark-equalities-nonconstant-rank}
上で証明した等式は、階数が一定でないことを許した有限局所自由
$\mathcal{O}_X$-加群を扱う場合にも依然として正しい。
実際、階数と Chern 類の多項式を次のように扱える。次数付き多項式環
$\mathbf{Z}[r, c_1, c_2, c_3, \ldots]$ を考え、$r$ の次数を $0$、$c_i$ の次数を $i$ とする。
$$
P \in \mathbf{Z}[r, c_1, c_2, c_3, \ldots]
$$
を次数 $p$ の同次多項式とする。このとき、$X$ 上の任意の有限局所自由
$\mathcal{O}_X$-加群 $\mathcal{E}$ に対して
$$
P(\mathcal{E}) =
P(r(\mathcal{E}), c_1(\mathcal{E}), c_2(\mathcal{E}), c_3(\mathcal{E}), \ldots)
\in A^p(X)
$$
を考えられる。記法と約束は注意 \ref{chow-remark-extend-to-finite-locally-free} を見よ。
これらの多項式間の関係（複数の有限局所自由加群に対するもの）を証明するには、
$X$ 上局所的に作業し、上の分解原理を用いればよい。例えば
$$
c_2(\SheafHom_{\mathcal{O}_X}(\mathcal{E}, \mathcal{E})) =
P(\mathcal{E})
$$
と主張する。ただし $P = 2rc_2 - (r - 1)c_1^2$ である。
実際、$\SheafHom_{\mathcal{O}_X}(\mathcal{E}, \mathcal{E}) =
\mathcal{E} \otimes \mathcal{E}^\vee$ なので、上の補題
\ref{chow-lemma-chern-classes-dual} および \ref{chow-lemma-chern-classes-tensor-product} を適用できる。
ここで注意 \ref{chow-remark-extend-to-finite-locally-free} のように、$X$ を $\mathcal{E}$ の階数が
一定となる部分に分解すれば結論が容易に従う。
\end{remark}

\begin{example}
\label{chow-example-power-sum}
任意の $p \geq 1$ に対して、次数 $p$ の同次多項式
$P_p \in \mathbf{Z}[c_1, c_2, c_3, \ldots]$ で、任意の $n \geq p$ に対して
$$
P_p(s_1, s_2, \ldots, s_p) = \sum x_i^p
$$
が $\mathbf{Z}[x_1, \ldots, x_n]$ で成り立つものが唯一つ存在する。
ここで $s_1, \ldots, s_p$ は $x_1, \ldots, x_n$ の基本対称多項式であり、
$$
s_i = \sum\nolimits_{1 \leq j_1 < \ldots < j_i \leq n}
x_{j_1}x_{j_2} \ldots x_{j_i}
$$
である。$P_p$ の存在は、基本対称関数が整数上の対称関数の環全体を生成するという
よく知られた事実から従う。$P_p \in \mathbf{Z}[c_1, c_2, c_3, \ldots]$ を特徴付ける
別の方法は、形式的べき級数として
$$
\log(1 + c_1 + c_2 + c_3 + \ldots) =
\sum\nolimits_{p \geq 1} (-1)^{p - 1}\frac{P_p}{p}
$$
が成り立つことである。これは
$1 + c_1 + c_2 + \ldots = \prod (1 + x_i)$ と書き、対数関数のべき級数を適用すれば明らかである。
左辺を展開すると
\begin{align*}
& (c_1 + c_2 + \ldots) - (1/2)(c_1 + c_2 + \ldots)^2 +
(1/3)(c_1 + c_2 + \ldots)^3 - \ldots \\
& =
c_1 + (c_2 - (1/2)c_1^2) + (c_3 - c_1c_2 + (1/3)c_1^3) + \ldots
\end{align*}
を得る。これにより
\begin{align*}
P_1 & = c_1, \\
P_2 & = c_1^2 - 2c_2, \\
P_3 & = c_1^3 - 3c_1c_2 + 3c_3, \\
P_4 & = c_1^4 - 4c_1^2c_2 + 4c_1c_3 + 2c_2^2 - 4c_4,
\end{align*}
などが得られる。有限局所自由 $\mathcal{O}_X$-加群 $\mathcal{E}$ の Chern 類は
Chern 根 $x_i$ の基本対称多項式なので、
$$
P_p(\mathcal{E}) = \sum x_i^p
$$
である。便宜上、$\mathbf{Z}[r, c_1, c_2, c_3, \ldots]$ において $P_0 = r$ とおく。
すると双変類として $P_0(\mathcal{E}) = r(\mathcal{E})$ である
（注意 \ref{chow-remark-extend-to-finite-locally-free} および
\ref{chow-remark-equalities-nonconstant-rank} のとおり）。
\end{example}






\section{Chern 類と切断}
\label{chow-section-top-chern-class}

\noindent
この短い節の主結果は、有限局所自由加群の最高 Chern 類を
「正則切断」の零点軌跡を用いて計算できることである。

\medskip\noindent
$(S, \delta)$ を状況 \ref{chow-situation-setup} のとおりとする。$X$ を $S$ 上局所有限型なスキームとし、
$\mathcal{E}$ を有限局所自由 $\mathcal{O}_X$-加群とする。$f : X' \to X$ を局所有限型とし、
$$
s \in \Gamma(X', f^*\mathcal{E})
$$
を $X'$ への $\mathcal{E}$ の引き戻しの大域切断とする。
$Z(s) \subset X'$ を $s$ の零点スキームとする。より正確には、$Z(s)$ を、
イデアルの準連接層が写像
$s : f^*\mathcal{E}^\vee \to \mathcal{O}_{X'}$ の像である閉部分スキームと定義する。

\begin{lemma}
\label{chow-lemma-top-chern-class}
直前の状況で $\dim_\delta(X') = n$ と仮定する。$f^*\mathcal{E}$ は一定階数 $r$ をもち、
$\dim_\delta(Z(s)) \leq n - r$ であり、$\delta(\xi) = n - r$ を満たす任意の生成点
$\xi \in Z(s)$ に対して、$\mathcal{O}_{X', \xi}$ における $Z(s)$ のイデアルが
長さ $r$ の正則列で生成されるとする。このとき $\CH_*(X')$ において
$$
c_r(\mathcal{E}) \cap [X']_n = [Z(s)]_{n - r}
$$
である。
\end{lemma}

\begin{proof}
$c_r(\mathcal{E})$ は双変類なので（補題 \ref{chow-lemma-cap-cp-bivariant}）、
$X = X'$ と仮定し、$\CH_{n - r}(X)$ における等式
$c_r(\mathcal{E}) \cap [X]_n = [Z(s)]_{n - r}$ を示せばよい。
$r \geq 0$ に関する帰納法で示す。（$r = 0$ の場合は自明である。）
$r = 1$ の場合は補題 \ref{chow-lemma-geometric-cap} から従う。$r > 1$ と仮定する。

\medskip\noindent
$\pi : P \to X$ を $\mathcal{E}$ に付随する射影空間束とし、短完全列
$$
0 \to \mathcal{E}' \to \pi^*\mathcal{E} \to \mathcal{O}_P(1) \to 0
$$
を考える。射影空間束公式（補題 \ref{chow-lemma-chow-ring-projective-bundle}）により、
$\pi$ で引き戻した後に等式を示せば十分である。
$\pi^{-1}Z(s) = Z(\pi^*s)$ の $\delta$-次元は $\leq n - 1$ である。また、
$\delta$-次元 $n - 1$ の生成点における正則列の仮定は平坦引き戻しにより成り立つ。
Algebra の補題 \ref{algebra-lemma-flat-increases-depth} を見よ。
$t \in \Gamma(P, \mathcal{O}_P(1))$ を $\pi^*s$ の像とする。次を主張する。
$$
[Z(t)]_{n + r - 2} = c_1(\mathcal{O}_P(1)) \cap [P]_{n + r - 1}
$$
この主張を仮定し、次のように証明を完了する。
制限 $\pi^*s|_{Z(t)}$ は $\mathcal{O}_P(1)|_{Z(t)}$ で零に写るので、唯一の元
$s' \in \Gamma(Z(t), \mathcal{E}'|_{Z(t)})$ から来る。
$P$ の閉部分スキームとして $Z(s') = Z(\pi^*s)$ である。
$\xi \in Z(s')$ が $\delta(\xi) = n - 1$ を満たす生成点なら、
$\mathcal{O}_{Z(t), \xi}$ における $Z(s')$ のイデアルは長さ $r - 1$ の正則列で生成できる。
実際、このイデアルは $r - 1$ 個の元で生成され、それらは
$\mathcal{O}_{P, \xi}$ の $r - 1$ 個の元の像である。
これらの元と $\mathcal{O}_{P, \xi}$ における $Z(t)$ のイデアルの一つの生成元は、
$\mathcal{O}_{P, \xi}$ で長さ $r$ の正則列をなす。したがって $Z(t)$ 上の $s'$ に
帰納法の仮定を適用し、
$c_{r - 1}(\mathcal{E}') \cap [Z(t)]_{n + r - 2} = [Z(s')]_{n - 1}$ を得る。
以上を組み合わせると
\begin{align*}
c_r(\pi^*\mathcal{E}) \cap [P]_{n + r - 1}
& =
c_{r - 1}(\mathcal{E}') \cap c_1(\mathcal{O}_P(1)) \cap [P]_{n + r - 1} \\
& =
c_{r - 1}(\mathcal{E}') \cap [Z(t)]_{n + r - 2} \\
& =
[Z(s')]_{n - 1} \\
& = [Z(\pi^*s)]_{n - 1}
\end{align*}
となり、求める等式を得る。

\medskip\noindent
主張を証明する。これはすでに用いた補題
\ref{chow-lemma-geometric-cap} の適用から従う。
$\pi^{-1}(Z(s)) = Z(\pi^*s) \subset Z(t)$ である。一方、$x \in X$ に対して
$P_x \subset Z(t)$ なら、$t|_{P_x} = 0$ であり、$s$ はファイバー
$\mathcal{E} \otimes \kappa(x)$ で零となるから $x \in Z(s)$ である。
したがって $\dim_\delta(Z(t)) \leq n + (r - 1) - 1$ である。
最後に、$\xi \in Z(t)$ を $\delta(\xi) = n + r - 2$ を満たす生成点とする。
$\xi$ が $P \to X$ のファイバーの生成点でないなら、$Z(t)$ の局所方程式が
$\mathcal{O}_{P, \xi}$ の非零因子であることは直ちに従う（Algebra の補題
\ref{algebra-lemma-grothendieck} によりファイバー上で確認できるからである）。
$\xi$ がファイバーの生成点なら、$x = \pi(\xi) \in Z(s)$ かつ
$\delta(x) = n + r - 2 - (r - 1) = n - 1$ である。これは $r > 1$ のときの
$\dim_\delta(Z(s)) \leq n - r$ に矛盾する。よってこの場合は起こらない。
\end{proof}

\begin{lemma}
\label{chow-lemma-easy-virtual-class}
$(S, \delta)$ を状況 \ref{chow-situation-setup} のとおりとし、$X$ を $S$ 上局所有限型なスキームとする。
$$
0 \to \mathcal{N}' \to \mathcal{N} \to \mathcal{E} \to 0
$$
を有限局所自由 $\mathcal{O}_X$-加群の短完全列とし、閉埋め込み
$$
i :
N' = \underline{\Spec}_X(\text{Sym}((\mathcal{N}')^\vee))
\longrightarrow
N = \underline{\Spec}_X(\text{Sym}(\mathcal{N}^\vee))
$$
を考える。$\alpha \in \CH_k(X)$ に対して
$$
i_*(p')^*\alpha = p^*(c_{top}(\mathcal{E}) \cap \alpha)
$$
が成り立つ。ただし $p' : N' \to X$ および $p : N \to X$ は構造射である。
\end{lemma}

\begin{proof}
ここで $c_{top}(\mathcal{E})$ は注意 \ref{chow-remark-top-chern-class} で定義した双変類である。
その定義自体から、式を確かめるために $\mathcal{E}$ は一定階数をもつと仮定できる。
同様に $\mathcal{N}'$ と $\mathcal{N}$ の階数がそれぞれ一定で $r'$, $r$ と仮定できる。
したがって $\mathcal{E}$ の階数は $r - r'$ であり、
$c_{top}(\mathcal{E}) = c_{r - r'}(\mathcal{E})$ である。
$p^*\mathcal{E}$ は標準的な切断
$$
s \in \Gamma(N, p^*\mathcal{E}) = \Gamma(X, p_*p^*\mathcal{E}) =
\Gamma(X, \mathcal{E} \otimes_{\mathcal{O}_X} \text{Sym}(\mathcal{N}^\vee)
\supset \Gamma(X, \SheafHom(\mathcal{N}, \mathcal{E}))
$$
をもつ。これは補題の主張で与えられた全射 $\mathcal{N} \to \mathcal{E}$ に対応する。
この切断の零点スキームはちょうど $N' \subset N$ である。
$Y \subset X$ を $\delta$-次元 $n$ の整閉部分スキームとする。このとき次が成り立つ。
\begin{enumerate}
\item $p^{-1}(Y)$ は $\delta$-次元 $n + r$ の整スキームなので
$p^*[Y] = [p^{-1}(Y)]$ である。
\item $(p')^{-1}(Y)$ は $\delta$-次元 $n + r'$ の整スキームなので
$(p')^*[Y] = [(p')^{-1}(Y)]$ である。
\item $p^{-1}Y$ への $s$ の制限の零点スキームは $(p')^{-1}Y$ であり、
閉埋め込み $(p')^{-1}Y \to p^{-1}Y$ は正則埋め込み（局所的に正則列で切り出される）である。
\end{enumerate}
補題 \ref{chow-lemma-top-chern-class} により
$$
(p')^*[Y] = c_{r - r'}(p^*\mathcal{E}) \cap p^*[Y]
\quad\text{in}\quad \CH_*(N)
$$
を得る。これで補題が示された。
\end{proof}










\section{Chern 指標とテンソル積}
\label{chow-section-chern-classes-tensor}

\noindent
$(S, \delta)$ を状況 \ref{chow-situation-setup} のとおりとし、$X$ を $S$ 上局所有限型とする。
有限局所自由 $\mathcal{O}_X$-加群の {\it Chern 指標} を次の形式的表式で定義する。
$$
ch({\mathcal E}) = \sum\nolimits_{i=1}^r e^{x_i}
$$
ただし $x_i$ は ${\mathcal E}$ の Chern 根である。これを Chern 類の多項式として書くと
\begin{align*}
ch(\mathcal{E})
& =
r(\mathcal{E})
+
c_1(\mathcal{E}) +
\frac{1}{2}(c_1(\mathcal{E})^2 - 2c_2(\mathcal{E}))
+
\frac{1}{6}(c_1(\mathcal{E})^3 - 3c_1(\mathcal{E})c_2(\mathcal{E}) + 3c_3(\mathcal{E})) \\
& \quad\quad +
\frac{1}{24}(c_1(\mathcal{E})^4 - 4c_1(\mathcal{E})^2c_2(\mathcal{E}) + 4c_1(\mathcal{E})c_3(\mathcal{E}) + 2c_2(\mathcal{E})^2 - 4c_4(\mathcal{E}))
+
\ldots \\
& =
\sum\nolimits_{p = 0, 1, 2, \ldots} \frac{P_p(\mathcal{E})}{p!}
\end{align*}
を得る。ここで $P_p$ は例 \ref{chow-example-power-sum} の Chern 類の多項式である。
上の式の次数 $p$ の成分は
$$
ch_p(\mathcal{E}) = \frac{P_p(\mathcal{E})}{p!} \in A^p(X) \otimes \mathbf{Q}
$$
である。係数が有理数であるとは何を意味するのか。
これは単に、$ch_p(\mathcal{E})$ を $A^p(X) \otimes \mathbf{Q}$ の元とみなすという意味である。

\begin{remark}
\label{chow-remark-extend-chern-character-to-finite-locally-free}
上の議論では、$\mathcal{E}$ の階数が一定でない場合にも、$\mathcal{E}$ の Chern 指標の成分
$ch_p(\mathcal{E}) \in A^p(X) \otimes \mathbf{Q}$ を定義した。注意
\ref{chow-remark-extend-to-finite-locally-free} および \ref{chow-remark-equalities-nonconstant-rank} を見よ。
したがって $\mathcal{E}$ の Chern 指標全体は
$\prod_{p \geq 0} (A^p(X) \otimes \mathbf{Q})$ の元である。$X$ が準コンパクトで
$\dim(X) < \infty$（通常の次元）なら、補題 \ref{chow-lemma-vanish-above-dimension} と分解原理を用いて
$ch(\mathcal{E}) \in A^*(X) \otimes \mathbf{Q}$ となることを示せる。
\end{remark}

\begin{lemma}
\label{chow-lemma-chern-character-additive}
$(S, \delta)$ を状況 \ref{chow-situation-setup} のとおりとし、$X$ を $S$ 上局所有限型とする。
$
0 \to \mathcal{E}_1 \to \mathcal{E} \to \mathcal{E}_2 \to 0
$
を有限局所自由 $\mathcal{O}_X$-加群の短完全列とする。このとき
$$
ch(\mathcal{E}) = ch(\mathcal{E}_1) + ch(\mathcal{E}_2)
$$
である。より正確には、例 \ref{chow-example-power-sum} の $P_p$ に対して、$A^p(X)$ において
$P_p(\mathcal{E}) = P_p(\mathcal{E}_1) + P_p(\mathcal{E}_2)$ である。
\end{lemma}

\begin{proof}
より正確な主張を示せば十分である。節 \ref{chow-section-splitting-principle} によりこれは次のことから従う。
$x_{1, i}$, $i = 1, \ldots, r_1$ および $x_{2, i}$, $i = 1, \ldots, r_2$ がそれぞれ
$\mathcal{E}_1$ および $\mathcal{E}_2$ の Chern 根なら、
$x_{1, 1}, \ldots, x_{1, r_1}, x_{2, 1}, \ldots, x_{2, r_2}$ は $\mathcal{E}$ の Chern 根である。
したがって例 \ref{chow-example-power-sum} の $P_p$ の選び方から結論を得る。
\end{proof}

\begin{lemma}
\label{chow-lemma-chern-character-multiplicative}
$(S, \delta)$ を状況 \ref{chow-situation-setup} のとおりとし、$X$ を $S$ 上局所有限型とする。
$\mathcal{E}_1$ および $\mathcal{E}_2$ を有限局所自由 $\mathcal{O}_X$-加群とする。このとき
$$
ch(\mathcal{E}_1 \otimes_{\mathcal{O}_X} \mathcal{E}_2) =
ch(\mathcal{E}_1) ch(\mathcal{E}_2)
$$
である。より正確には、
$$
P_p(\mathcal{E}_1 \otimes_{\mathcal{O}_X} \mathcal{E}_2) =
\sum\nolimits_{p_1 + p_2 = p}
{p \choose p_1} P_{p_1}(\mathcal{E}_1) P_{p_2}(\mathcal{E}_2)
$$
が $A^p(X)$ で成り立つ。ただし $P_p$ は例 \ref{chow-example-power-sum} のとおりである。
\end{lemma}

\begin{proof}
より正確な主張を示せば十分である。節 \ref{chow-section-splitting-principle} によりこれは次のことから従う。
$x_{1, i}$, $i = 1, \ldots, r_1$ および $x_{2, i}$, $i = 1, \ldots, r_2$ がそれぞれ
$\mathcal{E}_1$ および $\mathcal{E}_2$ の Chern 根なら、
$x_{1, i} + x_{2, j}$, $1 \leq i \leq r_1$, $1 \leq j \leq r_2$ は
$\mathcal{E}_1 \otimes \mathcal{E}_2$ の Chern 根である。
したがって $(x_{1, i} + x_{2, j})^p$ の二項展開と、例
\ref{chow-example-power-sum} の多項式 $P_p$ の形から結論を得る。
\end{proof}

\begin{lemma}
\label{chow-lemma-chern-character-dual}
状況 \ref{chow-situation-setup} において、$X$ を $S$ 上局所有限型とする。
$\mathcal{E}$ を双対 $\mathcal{E}^\vee$ をもつ有限局所自由 $\mathcal{O}_X$-加群とする。このとき
$ch_i(\mathcal{E}^\vee) = (-1)^i ch_i(\mathcal{E})$ が $A^i(X) \otimes \mathbf{Q}$ で成り立つ。
\end{lemma}

\begin{proof}
Chern 類に対する対応する結果（補題 \ref{chow-lemma-chern-classes-dual}）から従う。
\end{proof}











\section{Chern 類と導来圏}
\label{chow-section-pre-derived}

\noindent
この節では、スキーム $X$ の導来圏の完全対象 $E$ が、$X$ の一つの包絡射上で
有限局所自由加群の有限複体で表されると仮定して、$E$ の全 Chern 類を定義する。

\medskip\noindent
$(S, \delta)$ を状況 \ref{chow-situation-setup} のとおりとし、$X$ を $S$ 上局所有限型とする。
$$
\mathcal{E}^a \to \mathcal{E}^{a + 1} \to \ldots \to \mathcal{E}^b
$$
を、一定階数の有限局所自由 $\mathcal{O}_X$-加群からなる有界複体とする。
このとき、{\it 複体の全 Chern 類} を
$$
c(\mathcal{E}^\bullet) = \prod\nolimits_{n = a, \ldots, b}
c(\mathcal{E}^n)^{(-1)^n} \in \prod\nolimits_{p \geq 0} A^p(X)
$$
で定義する。ここで逆元は形式的逆元であり、
$$
(1 + c_1 + c_2 + c_3 + \ldots)^{-1} =
1 - c_1 + c_1^2 - c_2 - c_1^3 + 2c_1 c_2 - c_3 + \ldots
$$
となる。$c(\mathcal{E}^\bullet)$ の次数 $p$ の部分を
$c_p(\mathcal{E}^\bullet) \in A^p(X)$ と表す。同様に、{\it 複体の Chern 指標} を
$$
ch(\mathcal{E}^\bullet) = \sum\nolimits_{n = a, \ldots, b}
(-1)^n ch(\mathcal{E}^n) \in
\prod\nolimits_{p \geq 0} (A^p(X) \otimes \mathbf{Q})
$$
で定義し、$ch(\mathcal{E}^\bullet)$ の次数 $p$ の部分を
$ch_p(\mathcal{E}^\bullet) \in A^p(X) \otimes \mathbf{Q}$ と表す。
最後に、例 \ref{chow-example-power-sum} の
$P_p \in \mathbf{Z}[r, c_1, c_2, c_3, \ldots]$ に対して
$$
P_p(\mathcal{E}^\bullet) = \sum\nolimits_{n = a, \ldots, b}
(-1)^n P_p(\mathcal{E}^n)
$$
を $A^p(X)$ で定義する。通常どおり
$ch_p(\mathcal{E}^\bullet) = (1/p!)P_p(\mathcal{E}^\bullet)$ である。
次の補題は、これらの構成が導来圏における複体の像のみに依存することを示す。

\begin{lemma}
\label{chow-lemma-pre-derived-chern-class}
$(S, \delta)$ を状況 \ref{chow-situation-setup} のとおりとし、$X$ を $S$ 上局所有限型とする。
$E \in D(\mathcal{O}_X)$ を、$E$ を表す有限局所自由 $\mathcal{O}_X$-加群の局所有界複体
$\mathcal{E}^\bullet$ が存在する対象とする。このとき、上の構成をわずかに一般化して得られる
$$
c(\mathcal{E}^\bullet) \in \prod\nolimits_{p \geq 0} A^p(X),\quad
ch(\mathcal{E}^\bullet) \in
\prod\nolimits_{p \geq 0} A^p(X) \otimes \mathbf{Q},\quad
P_p(\mathcal{E}^\bullet) \in A^p(X)
$$
は、複体 $\mathcal{E}^\bullet$ の選び方に依らない。
\end{lemma}

\begin{proof}
全 Chern 類について示す。他の二つの場合は、補題
\ref{chow-lemma-chern-character-additive} を補題 \ref{chow-lemma-additivity-chern-classes} の代わりに用いる
同じ議論から従う。

\medskip\noindent
注意 \ref{chow-remark-extend-to-finite-locally-free} と同様に、全 Chern 類 $c(\mathcal{E}^\bullet)$ を定義するため、
$X$ を開かつ閉な部分スキームに
$$
X = \coprod\nolimits_{i \in I} X_i
$$
と分解し、すべての $n$, $i$ に対して $X_i$ 上の $\mathcal{E}^n$ の階数が一定となるようにする。
（これらの階数は $X$ 上の局所定数関数なので、この分解は可能である。）
$\mathcal{E}^\bullet$ は局所有界なので、各固定した $i$ に対して非零な層
$\mathcal{E}^n|_{X_i}$ は有限個しかない。したがって上と同様に
$$
c(\mathcal{E}^\bullet|_{X_i}) =
\prod\nolimits_n c(\mathcal{E}^n|_{X_i})^{(-1)^n}
\in \prod\nolimits_{p \geq 0} A^p(X_i)
$$
と定義できる。補題 \ref{chow-lemma-disjoint-decomposition-bivariant} により
$A^p(X) = \prod_i A^p(X_i)$ である。ゆえに各 $p \in \mathbf{Z}$ に対して、すべての $i$ に対する
$X_i$ への制限が $c_p(\mathcal{E}^\bullet|_{X_i})$ に等しい唯一の元
$c_p(\mathcal{E}^\bullet) \in A^p(X)$ が存在する。

\medskip\noindent
$E$ を表す有限局所自由 $\mathcal{O}_X$-加群の第二の局所有界複体
$\mathcal{F}^\bullet$ があるとする。$g : Y \to X$ を局所有限型で $Y$ が整となる射とする。
補題 \ref{chow-lemma-bivariant-zero} により、
$c(g^*\mathcal{E}^\bullet) \cap [Y]$ と $c(g^*\mathcal{F}^\bullet) \cap [Y]$ が等しいことを示せば十分である。
さらに、$Y$ を $Y$ 上固有双有理な整スキームで置き換えた後にこれを示せばよい。
すると、まず $g^*\mathcal{E}^\bullet$ と $g^*\mathcal{F}^\bullet$ が、一定階数の有限局所自由
$\mathcal{O}_Y$-加群からなる有界複体であることが従う。次に More on Flatness の補題
\ref{flat-lemma-blowup-complex-integral} により、すべての $i \in \mathbf{Z}$ に対して
$H^i(Lg^*E)$ は Tor 次元 $\leq 1$ の完全対象であると仮定できる。
これで次段落の場合に帰着する。

\medskip\noindent
$X$ は整で、$\mathcal{E}^\bullet$, $\mathcal{F}^\bullet$ は一定階数の有限局所自由加群からなる有界複体であり、
すべての $i \in \mathbf{Z}$ に対して $H^i(E)$ は Tor 次元 $\leq 1$ の完全
$\mathcal{O}_X$-加群であると仮定する。
$c(\mathcal{E}^\bullet) \cap [X]$ と $c(\mathcal{F}^\bullet) \cap [X]$ が等しいことを示す必要がある。
複体の微分を $d_\mathcal{E}^i : \mathcal{E}^i \to \mathcal{E}^{i + 1}$ および
$d_\mathcal{F}^i : \mathcal{F}^i \to \mathcal{F}^{i + 1}$ と表す。More on Flatness の注意
\ref{flat-remark-when-you-have-a-complex} により、すべての $i$ に対して
$\Im(d_\mathcal{E}^i)$, $\Ker(d_\mathcal{E}^i)$, $\Im(d_\mathcal{F}^i)$, $\Ker(d_\mathcal{F}^i)$ は
有限局所自由 $\mathcal{O}_X$-加群である。加法性（補題 \ref{chow-lemma-additivity-chern-classes}）により
$$
c(\mathcal{E}^\bullet) = \prod\nolimits_i
c(\Ker(d_\mathcal{E}^i))^{(-1)^i} c(\Im(d_\mathcal{E}^i))^{(-1)^i}
$$
であり、$\mathcal{F}^\bullet$ についても同様である。短完全列
$$
0 \to \Im(d_\mathcal{E}^i) \to \Ker(d_\mathcal{E}^i) \to H^i(E) \to 0
\quad\text{and}\quad
0 \to \Im(d_\mathcal{F}^i) \to \Ker(d_\mathcal{F}^i) \to H^i(E) \to 0
$$
があるので、次段落で述べて解く問題に帰着する。

\medskip\noindent
$X$ は整であり、二つの短完全列
$$
0 \to \mathcal{E}' \to \mathcal{E} \to \mathcal{Q} \to 0
\quad\text{and}\quad
0 \to \mathcal{F}' \to \mathcal{F} \to \mathcal{Q} \to 0
$$
があって、$\mathcal{E}$, $\mathcal{E}'$, $\mathcal{F}$, $\mathcal{F}'$ は有限局所自由とする。
問題は
$c(\mathcal{E})c(\mathcal{E}')^{-1} \cap [X] =
c(\mathcal{F})c(\mathcal{F}')^{-1} \cap [X]$ を示すことである。
そのため、$\mathcal{G}$ を定義する短完全列
$$
0 \to \mathcal{G} \to \mathcal{E} \oplus \mathcal{F} \to \mathcal{Q} \to 0
$$
を考える。$\mathcal{Q}$ は Tor 次元 $\leq 1$ なので、$\mathcal{G}$ は有限局所自由である。
図式追跡により、全射 $\mathcal{G} \to \mathcal{F}$ の核は $\mathcal{E}$ 内の $\mathcal{E}'$ に同型に写り、
全射 $\mathcal{G} \to \mathcal{E}$ の核は $\mathcal{F}$ 内の $\mathcal{F}'$ に同型に写る。
（アフィン局所的には、これは Schanuel の補題から従うか、それと同値である。
Algebra の補題 \ref{algebra-lemma-Schanuel} を見よ。）したがって
$$
c(\mathcal{E})c(\mathcal{F}') = c(\mathcal{G}) =
c(\mathcal{F})c(\mathcal{E}')
$$
となり、望む等式を得る。
\end{proof}

\begin{lemma}
\label{chow-lemma-defined-by-envelope}
$(S, \delta)$ を状況 \ref{chow-situation-setup} のとおりとし、$X$ を $S$ 上局所有限型とする。
$E \in D(\mathcal{O}_X)$ を完全対象とする。$Lf^*E$ が $D(\mathcal{O}_Y)$ で、
有限局所自由 $\mathcal{O}_Y$-加群の局所有界複体 $\mathcal{E}^\bullet$ と同型になるような包絡射
$f : Y \to X$（定義 \ref{chow-definition-envelope}）が存在すると仮定する。
このとき、$f : Y \to X$ と $\mathcal{E}^\bullet$ の選び方に依らない唯一の双変類
$c(E) \in \prod_{p \geq 0} A^p(X)$,
$ch(E) \in \prod_{p \geq 0} A^p(X) \otimes \mathbf{Q}$, および
$P_p(E) \in A^p(X)$ が存在し、それらの $Y$ への制限はそれぞれ
$c(\mathcal{E}^\bullet)$, $ch(\mathcal{E}^\bullet)$, $P_p(\mathcal{E}^\bullet)$ に等しい。
\end{lemma}

\begin{proof}
$p \in \mathbf{Z}$ を固定する。Chern 類 $c_p(E) \in A^p(X)$ について補題を示し、
他の場合の議論は省略する。

\medskip\noindent
$g : T \to X$ を局所有限型な射で、$D(\mathcal{O}_T)$ において $Lg^*E$ を表す
有限局所自由 $\mathcal{O}_T$-加群の局所有界複体 $\mathcal{E}^\bullet$ が存在するものとする。
補題 \ref{chow-lemma-pre-derived-chern-class} により、双変類
$c_p(\mathcal{E}^\bullet) \in A^p(T)$ は $\mathcal{E}^\bullet$ の選び方に依らない。
この類を $c_p(Lg^*E) \in A^p(T)$ と書く。さらに $h : T' \to T$ を局所有限型な射とし、
$g' = g \circ h$ とおくと、
$L(g')^*E = L(g \circ h)^*E = Lh^*Lg^*E$ は $D(\mathcal{O}_{T'})$ で
$h^*\mathcal{E}^\bullet$ によって表される。したがって $c_p(L(g')^*E)$ は定義され、
$c_p(Lg^*E)$ の $T'$ への制限（注意 \ref{chow-remark-restriction-bivariant}）に等しい。
厳密にはここで補題 \ref{chow-lemma-cap-cp-bivariant} を適用する必要がある。

\medskip\noindent
$f : Y \to X$ と $\mathcal{E}^\bullet$ を補題の主張のとおりとする。
前段落で構成した類 $c' = c_p(Lf^*E)$ と $f : Y \to X$ に補題
\ref{chow-lemma-envelope-bivariant} を適用し、双変類 $c_p(E) \in A^p(X)$ を得る。
前段落の議論により、
$g = f \circ p = f \circ q : Y \times_X Y \to X$ に対して
$res_1(c') = c_p(Lg^*E) = res_2(c')$ であるから、補題の仮定は満たされる。

\medskip\noindent
最後に、$f' : Y' \to X$ を第二の包絡射で、$L(f')^*E$ が有限局所自由
$\mathcal{O}_{Y'}$-加群の有界複体で表されるものとする。このとき
$c_p(Lf^*E)$ と $c_p(L(f')^*E)$ の $Y \times_X Y'$ への制限は等しい。
$Y \times_X Y' \to X$ は包絡射なので（補題 \ref{chow-lemma-base-change-envelope} および
\ref{chow-lemma-composition-envelope}）、補題 \ref{chow-lemma-envelope-bivariant} の一意性により、
$c_p(E)$ の二つの候補は一致する。
\end{proof}

\begin{definition}
\label{chow-definition-defined-on-perfect}
$(S, \delta)$ を状況 \ref{chow-situation-setup} のとおりとし、$X$ を $S$ 上局所有限型とする。
$E \in D(\mathcal{O}_X)$ を完全対象とする。
\begin{enumerate}
\item $Lf^*E$ が $D(\mathcal{O}_Y)$ で有限局所自由 $\mathcal{O}_Y$-加群の局所有界複体と
同型になるような包絡射 $f : Y \to X$ が存在するとき、
{\it $E$ の Chern 類は定義される}という\footnote{いくつかの基準は補題
\ref{chow-lemma-chern-classes-defined} を見よ。}。
\item $E$ の Chern 類が定義されるとき、補題 \ref{chow-lemma-defined-by-envelope} を用いて
$$
c(E) \in \prod\nolimits_{p \geq 0} A^p(X),\quad
ch(E) \in
\prod\nolimits_{p \geq 0} A^p(X) \otimes \mathbf{Q},\quad
P_p(E) \in A^p(X)
$$
を定義する。
\end{enumerate}
\end{definition}

\noindent
この定義は、この章で想定する多くの状況に適用できるが、すべてに適用できるわけではない。
補題 \ref{chow-lemma-chern-classes-defined} を見よ。定義のような一般の
$E/X/(S,\delta)$ に対して、これらの双変類の初等的構成が存在するかもしれないが、われわれは知らない。

\begin{lemma}
\label{chow-lemma-chern-classes-defined}
$(S, \delta)$ を状況 \ref{chow-situation-setup} のとおりとし、$X$ を $S$ 上局所有限型とする。
$E \in D(\mathcal{O}_X)$ を完全対象とする。次の条件のいずれかが成り立つなら、
$E$ の Chern 類は定義される。
\begin{enumerate}
\item $Lf^*E$ が $D(\mathcal{O}_Y)$ で、有限局所自由 $\mathcal{O}_Y$-加群の局所有界複体と
同型になるような包絡射 $f : Y \to X$ が存在する。
\item $E$ は有限局所自由 $\mathcal{O}_X$-加群の有界複体で表される。
\item $X$ の既約成分は準コンパクトである。
\item $X$ は準コンパクトである。
\item $S$ 上局所有限型なスキームの射 $X \to X'$ が存在し、$E$ は $X'$ 上の Chern 類が
定義される完全対象 $E'$ の引き戻しである。
\item ここにさらに追加する。
\end{enumerate}
\end{lemma}

\begin{proof}
条件 (1) は定義 \ref{chow-definition-defined-on-perfect} の (1) そのものである。条件 (2) は (1) を含意する。

\medskip\noindent
(3) のように、$X$ の既約成分 $X_i$ が準コンパクトと仮定する。
$X_i$ には、$X$ 上の既約な整閉部分スキーム構造を入れる。
射 $\coprod X_i \to X$ は包絡射である。各 $i$ に対して、$X'_i$ が可逆加群の豊富族をもつような
包絡射 $X'_i \to X_i$ が存在する。More on Morphisms の命題
\ref{more-morphisms-proposition-envelope-with-resolution-property} を見よ。
$f : Y = \coprod X'_i \to X$ は包絡射である。小さな詳細は省略する。
Derived Categories of Schemes の補題
\ref{perfect-lemma-resolution-property-ample-family} により、各 $X'_i$ は分解性質をもつ。
したがって $D(\mathcal{O}_{X'_i})$ の完全対象 $L(f|_{X'_i})^*E$ は、有限局所自由
$\mathcal{O}_{X'_i}$-加群の有界複体で表される。Derived Categories of Schemes の補題
\ref{perfect-lemma-resolution-property-perfect-complex} を見よ。これで (3) が (1) を含意することが示された。

\medskip\noindent
(4) は (3) を含意する。

\medskip\noindent
$g : X \to X'$ と $E'$ を (5) のとおりとする。$L(f')^*E'$ が
$\mathcal{O}_{Y'}$-加群の局所有界複体 $(\mathcal{E}')^\bullet$ で表されるような包絡射
$f' : Y' \to X'$ が存在する。このとき基底変換 $f : Y \to X$ は補題
\ref{chow-lemma-base-change-envelope} により包絡射である。さらに引き戻し
$\mathcal{E}^\bullet = g^*(\mathcal{E}')^\bullet$ は $Lf^*E$ を表すので、$E$ の Chern 類は定義される。
\end{proof}

\begin{lemma}
\label{chow-lemma-chern-classes-computed}
$(S, \delta)$ を状況 \ref{chow-situation-setup} のとおりとし、$X$ を $S$ 上局所有限型とする。
$E \in D(\mathcal{O}_X)$ を完全対象とし、$E$ の Chern 類は定義されると仮定する。
$g : W \to X$ が局所有限型で $W$ が整なら、可換図式
$$
\xymatrix{
W' \ar[rd]_{g'} \ar[rr]_b & & W \ar[ld]^g \\
& X
}
$$
が存在する。ここで $W'$ は整、$b : W' \to W$ は固有双有理であり、
$L(g')^*E$ は一定階数の局所自由 $\mathcal{O}_{W'}$-加群の有界複体
$\mathcal{E}^\bullet$ で表され、$A^p(W')$ で
$res(c_p(E)) = c_p(\mathcal{E}^\bullet)$ である。
\end{lemma}

\begin{proof}
$Lf^*E$ が $D(\mathcal{O}_Y)$ で有限局所自由 $\mathcal{O}_Y$-加群の局所有界複体
$\mathcal{E}^\bullet$ と同型になる包絡射 $f : Y \to X$ を選ぶ。
$f$ の基底変換 $Y \times_X W \to W$ は補題 \ref{chow-lemma-base-change-envelope} により包絡射である。
$W$ の生成点に同じ剩余体をもって写る点 $\xi \in Y \times_X W$ を選ぶ。
$\xi$ を生成点とする整閉部分スキーム $W' \subset Y \times_X W$ を考える。
射影 $Y \times_X W \to W$ の $W'$ への制限は固有双有理射 $b : W' \to W$ である。
$g' = g \circ b$ とおき、引き戻し $(W' \to Y)^*\mathcal{E}^\bullet$ を考える。
これは $W'$ 上の有限局所自由加群からなる局所有界複体である。
$W'$ は整なので、この複体は有界で、各項の階数は一定である。
最後に、構成により $(W' \to Y)^*\mathcal{E}^\bullet$ は $L(g')^*E$ を表し、
その第 $p$ Chern 類は $W' \to X$ による $c_p(E)$ の制限を与える。これで証明が完了する。
\end{proof}

\begin{lemma}
\label{chow-lemma-commutative-chern-perfect}
$(S, \delta)$ を状況 \ref{chow-situation-setup} のとおりとし、$X$ を $S$ 上局所有限型とする。
$E \in D(\mathcal{O}_X)$ を完全対象とする。$E$ の Chern 類が定義されるなら、
\begin{enumerate}
\item $c_p(E)$ は代数 $A^*(X)$ の中心に属する。
\item $g : X' \to X$ が局所有限型で $c \in A^*(X' \to X)$ なら、
$c \circ c_p(E) = c_p(Lg^*E) \circ c$ である。
\end{enumerate}
\end{lemma}

\begin{proof}
(2) から (1) が直ちに従う。$g : X' \to X$ と $c \in A^*(X' \to X)$ を (2) のとおりとする。
$c \circ c_p(E) - c_p(Lg^*E) \circ c = 0$ を示すため、補題 \ref{chow-lemma-bivariant-zero} の基準を用いる。
したがって $X$ は整と仮定でき、さらに補題 \ref{chow-lemma-chern-classes-computed} により、
$E$ は一定階数の有限局所自由 $\mathcal{O}_X$-加群の有界複体
$\mathcal{E}^\bullet$ で表されると仮定できる。そこで $\CH_*(X')$ において
$$
c \cap c_p(\mathcal{E}^\bullet) \cap [X] =
c_p(\mathcal{E}^\bullet) \cap c \cap [X] 
$$
を示せばよい。これは補題 \ref{chow-lemma-cap-commutative-chern} と、
$c_p(\mathcal{E}^\bullet)$ が局所自由加群 $\mathcal{E}^n$ の Chern 類の多項式として
構成されることから直ちに従う。
\end{proof}

\begin{lemma}
\label{chow-lemma-additivity-on-perfect}
$(S, \delta)$ を状況 \ref{chow-situation-setup} のとおりとし、$X$ を $S$ 上局所有限型とする。
$$
E_1 \to E_2 \to E_3 \to E_1[1]
$$
を $D(\mathcal{O}_X)$ の完全対象からなる特別三角形とする。
次の条件のいずれかが成り立つと仮定する。
\begin{enumerate}
\item $Lf^*E_1 \to Lf^*E_2$ が、有限局所自由 $\mathcal{O}_Y$-加群の局所有界複体の写像で
表されるような包絡射 $f : Y \to X$ が存在する。
\item $E_1 \to E_2$ は、有限局所自由 $\mathcal{O}_X$-加群の局所有界複体の写像で表される。
\item $X$ の既約成分は準コンパクトである。
\item $X$ は準コンパクトである。
\item ここにさらに追加する。
\end{enumerate}
このとき $E_1$, $E_2$, $E_3$ の Chern 類は定義され、
$c(E_2) = c(E_1) c(E_3)$, $ch(E_2) = ch(E_1) + ch(E_3)$, および
$P_p(E_2) = P_p(E_1) + P_p(E_3)$ が成り立つ。
\end{lemma}

\begin{proof}
$f : Y \to X$ を包絡射とし、
$\alpha^\bullet : \mathcal{E}_1^\bullet \to \mathcal{E}_2^\bullet$ を、$Lf^*E_1 \to Lf^*E_2$ を表す
有限局所自由 $\mathcal{O}_Y$-加群の局所有界複体の写像とする。
このとき写像円錐 $C(\alpha)^\bullet$ は $Lf^*E_3$ を表す。
$C(\alpha)^n = \mathcal{E}_2^n \oplus \mathcal{E}_1^{n + 1}$ なので、$C(\alpha)^\bullet$ は
有限局所自由 $\mathcal{O}_Y$-加群の局所有界複体である。
したがって $E_1$, $E_2$, $E_3$ の Chern 類は定義される。
また $c_p(E_1)$ は、$A^p(Y)$ で $c_p(\mathcal{E}_1^\bullet)$ に制限される
$A^p(X)$ の唯一の元として定義されたことを思い出そう。$E_2$, $E_3$ についても同様である。
したがって $\prod_{p \geq 0} A^p(Y)$ において
$c(\mathcal{E}_2^\bullet) =
c(\mathcal{E}_1^\bullet) c(C(\alpha)^\bullet)$ を示せば十分である。
さらに $Y$ の連結成分へ制限した後に示せばよい。ゆえに複体
$\mathcal{E}_1^\bullet$ と $\mathcal{E}_2^\bullet$ は、一定階数の有限局所自由
$\mathcal{O}_Y$-加群の有界複体と仮定できる。この場合、求める等式は補題
\ref{chow-lemma-additivity-chern-classes} の乗法性から従う。$ch$ または $P_p$ の場合には補題
\ref{chow-lemma-chern-character-additive} を用いる。

\medskip\noindent
前段落で条件 (1) のときに補題が成り立つことを示した。(2) は (1) を含意するから、
これで第二の場合も扱われた。(3) を仮定する。補題
\ref{chow-lemma-chern-classes-defined} の証明とまったく同様に論じると、$Y$ が分解性質をもつ
準コンパクト（かつ準分離的）スキームの互いに交わらない和
$Y = \coprod Y_i$ となる包絡射 $f : Y \to X$ が得られる。
このとき $Y_i$ への $Lf^*E_1 \to Lf^*E_2$ の制限は、有限局所自由加群の有界複体の写像で
表される。Derived Categories of Schemes の命題
\ref{perfect-proposition-perfect-resolution-property} を見よ。
これで (3) が (1) を含意することがわかる。もちろん (4) は (3) を含意する。以上で証明が完了する。
\end{proof}

\begin{remark}
\label{chow-remark-splitting-principle-perfect}
完全複体の Chern 類は、定義されるときに一種の分解原理を満たす。
実際、$(S, \delta), X, E$ を定義 \ref{chow-definition-defined-on-perfect} のとおりとし、
$E$ の Chern 類が定義されると仮定する。双変類 $c_p(E)$, $P_p(E)$, $ch_p(E)$ の間の
関係を示したいとする。そのため、包絡射 $f : Y \to X$ と、$E$ を表す
有限局所自由 $\mathcal{O}_X$-加群の局所有界複体 $\mathcal{E}^\bullet$ を選ぶことができる。
補題 \ref{chow-lemma-defined-by-envelope} の一意性により、双変類
$c_p(\mathcal{E}^\bullet)$, $P_p(\mathcal{E}^\bullet)$, $ch_p(\mathcal{E}^\bullet)$ の間の求める関係を
示せば十分である。したがって $X$ を $Y$ の一つの連結成分で置き換え、
$E$ は一定階数の有限局所自由加群の有界複体 $\mathcal{E}^\bullet$ で表されると仮定できる。
分解原理（補題 \ref{chow-lemma-splitting-principle}）を用いると、各 $\mathcal{E}^i$ が、逐次商
$\mathcal{L}_{i, j}$ を可逆加群とするフィルトレーションをもつと仮定できる。
$x_{i, j} = c_1(\mathcal{L}_{i, j})$ とおくと
$$
c(E) =
\prod\nolimits_{i\text{ even}} (1 + x_{i, j})
\prod\nolimits_{i\text{ odd}} (1 + x_{i, j})^{-1}
$$
および
$$
P_p(E) =  \sum\nolimits_{i\text{ even}} (x_{i, j})^p -
\sum\nolimits_{i\text{ odd}} (x_{i, j})^p
$$
を得る。上の $c(E)$ の表式の対数を形式的にとると
$$
\log(c(E)) = \sum (-1)^{p - 1}\frac{P_p(E)}{p}
$$
である。例 \ref{chow-example-power-sum} の多項式 $P_p$ の構成を見ると、$P_p(E)$ は
ベクトル束の場合とまったく同じ $E$ の Chern 類の表式であることが従う。すなわち
\begin{align*}
P_1(E) & = c_1(E), \\
P_2(E) & = c_1(E)^2 - 2c_2(E), \\
P_3(E) & = c_1(E)^3 - 3c_1(E)c_2(E) + 3c_3(E), \\
P_4(E) & = c_1(E)^4 - 4c_1(E)^2c_2(E) + 4c_1(E)c_3(E) + 2c_2(E)^2 - 4c_4(E),
\end{align*}
などである。一方、双変類 $P_0(E) = r(E) = ch_0(E)$ は $E$ の Chern 類 $c(E)$ からは
復元できない。Chern 類は複体の階数の情報をもたない。
\end{remark}

\begin{lemma}
\label{chow-lemma-chern-classes-perfect-dual}
状況 \ref{chow-situation-setup} において、$X$ を $S$ 上局所有限型とする。
$E \in D(\mathcal{O}_X)$ を Chern 類が定義される完全対象とする。このとき
$c_i(E^\vee) = (-1)^i c_i(E)$, $P_i(E^\vee) = (-1)^iP_i(E)$, および
$ch_i(E^\vee) = (-1)^ich_i(E)$ が $A^i(X)$ で成り立つ。
\end{lemma}

\begin{proof}
第一の証明：補題 \ref{chow-lemma-commutative-chern-perfect} の証明と同様に論じ、
$E$ が一定階数の有限局所自由加群の有界複体で表される場合に帰着し、
補題 \ref{chow-lemma-chern-classes-dual} を適用する。第二の証明：注意
\ref{chow-remark-splitting-principle-perfect} で述べた分解原理を用い、$E^\vee$ の Chern 根が
$E$ の Chern 根の負であることを用いる。
\end{proof}

\begin{lemma}
\label{chow-lemma-chern-class-perfect-tensor-invertible}
状況 \ref{chow-situation-setup} において、$X$ を $S$ 上局所有限型とする。
$E$ を Chern 類が定義される $D(\mathcal{O}_X)$ の完全対象とし、
$\mathcal{L}$ を可逆 $\mathcal{O}_X$-加群とする。$E$ が一定階数 $r \in \mathbf{Z}$ をもつなら、
$$
c_i(E \otimes \mathcal{L}) =
\sum\nolimits_{j = 0}^i
\binom{r - i + j}{j} c_{i - j}(E) c_1(\mathcal{L})^j
$$
である。
\end{lemma}

\begin{proof}
$E$ が階数 $r$ の局所自由加群の場合、これは補題
\ref{chow-lemma-chern-classes-E-tensor-L} である。読者は形式的な計算により、この特別な場合から補題を導ける。
別の方法は注意 \ref{chow-remark-splitting-principle-perfect} の分解原理を用いることである。
この場合、次の代数的事実を示すことになる。形式的に
$$
\frac{\prod_{a = 1, \ldots, n} (1 + x_a)}{\prod_{n = 1, \ldots, m} (1 + y_b)}
= 1 + c_1 + c_2 + c_3 + \ldots
$$
と書き、$c_i$ が $\mathbf{Z}[x_i, y_j]$ における次数 $i$ の同次式なら、
$$
\frac{\prod_{a = 1, \ldots, n} (1 + x_a + t)}{\prod_{b = 1, \ldots, m} (1 + y_b + t)}
= \sum\nolimits_{i \geq 0} \sum\nolimits_{j = 0}^i
\binom{r - i + j}{j} c_{i - j} t^j
$$
である。ただし $r = n - m$ である。詳細は省略する。
\end{proof}

\begin{lemma}
\label{chow-lemma-chern-classes-perfect-tensor-product}
状況 \ref{chow-situation-setup} において、$X$ を $S$ 上局所有限型とする。
$E$ と $F$ を、Chern 類が定義される $D(\mathcal{O}_X)$ の完全対象とする。このとき
$$
c_1(E \otimes_{\mathcal{O}_X}^\mathbf{L} F) =
r(E) c_1(\mathcal{F}) + r(F) c_1(\mathcal{E})
$$
であり、$c_2(E \otimes_{\mathcal{O}_X}^\mathbf{L} F)$ は
$$
r(E) c_2(F) + r(F) c_2(E) + {r(E) \choose 2} c_1(F)^2 +
(r(E)r(F) - 1) c_1(F)c_1(E) + {r(F) \choose 2} c_1(E)^2
$$
で与えられる。$A^*(X)$ における高次の Chern 類についても同様の公式が成り立つ。
また $A^*(X) \otimes \mathbf{Q}$ において
$ch(E \otimes_{\mathcal{O}_X}^\mathbf{L} F) = ch(E) ch(F)$ である。
より正確には
$$
P_p(E \otimes_{\mathcal{O}_X}^\mathbf{L} F) = \sum\nolimits_{p_1 + p_2 = p}
{p \choose p_1} P_{p_1}(E) P_{p_2}(F)
$$
が $A^p(X)$ で成り立つ。
\end{lemma}

\begin{proof}
$Lf^*E$ と $Lf^*F$ が有限局所自由 $\mathcal{O}_X$-加群の局所有界複体で表されるような
包絡射 $f : Y \to X$ を選ぶと、ベクトル束に対する対応する結果、すなわち補題
\ref{chow-lemma-chern-classes-tensor-product} と
\ref{chow-lemma-chern-character-multiplicative} から計算により従う。
より適切な証明は、おそらく注意 \ref{chow-remark-splitting-principle-perfect} のように
分解原理を用い、補題を次段落で述べる多項式環の計算に帰着させるものである。

\medskip\noindent
$A$ を可換環とする（われわれの場合、これは $X$ の双変 Chow 環の、Chern 類で生成される
部分環である）。$S$ を有限集合とし、写像 $\epsilon : S \to \{\pm 1\}$ と
$f : S \to A$ が与えられているとする。$A$ において
$$
P_p(S, f , \epsilon) = \sum\nolimits_{s \in S} \epsilon(s) f(s)^p
$$
と定義する。第二の三つ組 $(S', \epsilon', f')$ が与えられたとき、$P_p$ について示すべき等式は
$$
P_p(S \times S', f + f' , \epsilon \epsilon') = 
\sum\nolimits_{p_1 + p_2 = p}
{p \choose p_1} P_{p_1}(S, f, \epsilon) P_{p_2}(S', f', \epsilon')
$$
である。これを示すには、$S \amalg S'$ を変数とする多項式環に帰着し、各項
$f(s)^if'(s')^j$ が左辺と右辺に同じ係数で現れることを確かめればよい。
$c_1(E \otimes_{\mathcal{O}_X}^\mathbf{L} F)$ と
$c_2(E \otimes_{\mathcal{O}_X}^\mathbf{L} F)$ の公式を検証するには、分解原理により、
これらの公式をねじれなし環で確かめることに帰着する。次に、注意
\ref{chow-remark-splitting-principle-perfect} で示した $P_j(E)$ と $c_i(E)$ の関係を用いる。例えば
$$
c_1(E \otimes F) = P_1(E \otimes F) = r(F)P_1(E) + r(E)P_1(F) =
r(F)c_1(E) + r(E)c_1(F)
$$
であり、中間の等式は定義により $r(E) = P_0(E)$ だからである。同様に
\begin{align*}
& 2c_2(E \otimes F) \\
& = c_1(E \otimes F)^2 - P_2(E \otimes F) \\
& =
(r(F)c_1(E) + r(E)c_1(F))^2 -
r(F)P_2(E) - P_1(E)P_1(F) - r(E)P_2(F) \\
& =
(r(F)c_1(E) + r(E)c_1(F))^2 -
r(F)(c_1(E)^2 - 2c_2(E)) - c_1(E)c_1(F) - \\
& \quad r(E)(c_1(F)^2 - 2c_2(F))
\end{align*}
を得る。これは因子 $2$ を除いて補題の主張の公式と一致することを、読者は確かめられる。
\end{proof}







\section{局所化 Chern 類の初歩的な場合}
\label{chow-section-preparation-localized-chern}

\noindent
この節では、次の補題で構成する双変類のいくつかの性質を論じる。その大部分は、
補題で与える特徴づけから直ちに従う。読者には、この節の残りを飛ばすことを勧める。

\begin{lemma}
\label{chow-lemma-silly}
$(S, \delta)$ を状況 \ref{chow-situation-setup} のとおりとし、$X$ を $S$ 上局所有限型とする。
$i_j : X_j \to X$, $j = 1, 2$ を、集合論的に $X = X_1 \cup X_2$ を満たす閉埋め込みとする。
$E_2 \in D(\mathcal{O}_{X_2})$ を完全対象とする。次を仮定する。
\begin{enumerate}
\item $E_2$ の Chern 類が定義される。
\item 制限 $E_2|_{X_1 \cap X_2}$ は零である。resp.\ コホモロジー次数 $0$ に置かれた、
階数 $< p$ の有限局所自由 $\mathcal{O}_{X_1 \cap X_2}$-加群と同型である。
\end{enumerate}
このとき標準的な双変類
$$
P'_p(E_2),\text{ resp. }c'_p(E_2) \in A^p(X_2 \to X)
$$
が存在し、これは
$$
P'_p(E_2) \cap i_{2, *} \alpha_2 = P_p(E_2) \cap \alpha_2
\quad\text{and}\quad
P'_p(E_2) \cap i_{1, *} \alpha_1 = 0,
$$
それぞれ
$$
c'_p(E_2) \cap i_{2, *} \alpha_2 = c_p(E_2) \cap \alpha_2
\quad\text{and}\quad
c'_p(E_2) \cap i_{1, *} \alpha_1 = 0
$$
を、$\alpha_i \in \CH_k(X_i)$ に対して、また任意の局所有限型基底変換
$X' \to X$ の後にも満たすという性質によって特徴づけられる。
\end{lemma}

\begin{proof}
以下では節 \ref{chow-section-pre-derived} の内容を断りなく用いる。

\medskip\noindent
$E_2|_{X_1 \cap X_2}$ が零であると仮定する。局所有限型なスキームの射
$X' \to X$ を考え、$i_j$ の基底変換を $i'_j : X'_j \to X'$ と表す。
補題 \ref{chow-lemma-exact-sequence-closed-chow} により、任意の元
$\alpha' \in \CH_k(X')$ は $i'_{1, *}\alpha'_1 + i'_{2, *}\alpha'_2$ と書ける。
ここで $\alpha'_2 \in \CH_k(X'_2)$ は、$X'_1 \cap X'_2 \to X'_2$ による前進像の像の元を
加える違いを除いて定まる。そこで
$P'_p(E_2) \cap \alpha' = P_p(E_2) \cap \alpha'_2 \in \CH_{k - p}(X'_2)$ とおける。
$E_2$ の $X_1 \cap X_2$ への制限が零であるという仮定により、これは良定義である。

\medskip\noindent
$E_2|_{X_1 \cap X_2}$ が、コホモロジー次数 $0$ に置かれた階数 $< p$ の有限局所自由
$\mathcal{O}_{X_1 \cap X_2}$-加群と同型なら、階数を考えることにより
$c_p(E_2|_{X_1 \cap X_2}) = 0$ であり、まったく同様に論じられる。
\end{proof}

\begin{lemma}
\label{chow-lemma-silly-independent}
補題 \ref{chow-lemma-silly} の双変類 $P'_p(E_2)$, resp.\ $c'_p(E_2)$ は
$A^p(X_2 \to X)$ において $X_1$ の選択に依存しない。
\end{lemma}

\begin{proof}
$X_1' \subset X$ を、集合論的に $X = X'_1 \cup X_2$ を満たす別の閉部分スキームとし、
制限 $E_2|_{X'_1 \cap X_2}$ は零である、resp.\ コホモロジー次数 $0$ に置かれた
階数 $< p$ の有限局所自由 $\mathcal{O}_{X'_1 \cap X_2}$-加群と同型であると仮定する。
このとき $X = (X_1 \cap X'_1) \cup X_2$ である。したがって任意の元
$\alpha \in \CH_k(X)$ は、$\alpha_2 \in \CH_k(X'_2)$ および
$\beta \in \CH_k(X_1 \cap X'_1)$ を用いて $i_*\beta + i_{2, *}\alpha_2$ と書ける。
よって
$P'_p(E_2) \cap \alpha = P_p(E_2) \cap \alpha_2 \in \CH_{k - p}(X_2)$,
resp.\  $c'_p(E_2) \cap \alpha = c_p(E_2) \cap \alpha_2 \in \CH_{k - p}(X_2)$
は、$X_1$ と $X'_1$ のどちらを用いるかに依存しないことが明らかである。
任意の基底変換の後も同様である。
\end{proof}

\begin{lemma}
\label{chow-lemma-base-change-silly}
補題 \ref{chow-lemma-silly} において、$X' \to X$ を局所有限型な射とする。
$X'$ への引き戻しを $X' = X'_1 \cup X'_2$ および
$E'_2 \in D(\mathcal{O}_{X'_2})$ と表す。このとき、$X' = X'_1 \cup X'_2$ と $E_2'$ を用いて
補題 \ref{chow-lemma-silly} で構成される $A^p(X_2' \to X')$ の類
$P'_p(E_2')$, resp.\ $c'_p(E_2')$ は、$A^p(X_2 \to X)$ の類
$P'_p(E_2)$, resp.\ $c'_p(E_2)$ の制限（注意 \ref{chow-remark-restriction-bivariant}）である。
\end{lemma}

\begin{proof}
補題 \ref{chow-lemma-silly} におけるこれらの類の特徴づけから直ちに従う。
\end{proof}

\begin{lemma}
\label{chow-lemma-silly-silly}
補題 \ref{chow-lemma-silly} において、$E_2$ は完全対象 $E \in D(\mathcal{O}_X)$ の制限であり、
$E|_{X_1}$ は零である、resp.\ コホモロジー次数 $0$ に置かれた階数 $< p$ の
有限局所自由 $\mathcal{O}_{X_1}$-加群と同型であるとする。
$E$ の Chern 類が定義されるなら
$i_{2, *} \circ P'_p(E_2) = P_p(E)$,
resp.\ $i_{2, *} \circ c'_p(E_2) = c_p(E)$
である（$\circ$ は補題 \ref{chow-lemma-push-proper-bivariant} のものとする）。
\end{lemma}

\begin{proof}
まず $E|_{X_1}$ が零であると仮定する。補題 \ref{chow-lemma-silly} の証明の記法を用いると、
この場合の補題は
\begin{align*}
P_p(E) \cap \alpha'
& =
i'_{1, *}(P_p(E) \cap \alpha'_1) +
i'_{2, *}(P_p(E) \cap \alpha'_2) \\
& =
i'_{1, *}(P_p(E|_{X_1}) \cap \alpha'_1) +
i'_{2, *}(P'_p(E_2) \cap \alpha') \\
& =
i'_{2, *}(P'_p(E_2) \cap \alpha')
\end{align*}
から従う。$E|_{X_1}$ がコホモロジー次数 $0$ に置かれた階数 $< p$ の有限局所自由
$\mathcal{O}_{X_1}$-加群と同型である場合も同様である。
\end{proof}

\begin{lemma}
\label{chow-lemma-silly-shrink}
補題 \ref{chow-lemma-silly} において、閉部分スキーム $X'_2 \subset X_2$ と
$X_1 \subset X'_1 \subset X$ があり、集合論的に $X = X'_1 \cup X'_2$ を満たすと仮定する。
$E_2|_{X'_1 \cap X_2}$ は零である、resp.\ 次数 $0$ に置かれた階数 $< p$ の
有限局所自由加群と同型であると仮定する。このとき
$(X'_2 \to X_2)_* \circ P'_p(E_2|_{X'_2}) = P'_p(E_2)$,
resp.\  $(X'_2 \to X_2)_* \circ c'_p(E_2|_{X'_2}) = c_p(E_2)$
である（$\circ$ は補題 \ref{chow-lemma-push-proper-bivariant} のものとする）。
\end{lemma}

\begin{proof}
補題 \ref{chow-lemma-silly} におけるこれらの類の特徴づけから直ちに従う。
\end{proof}

\begin{lemma}
\label{chow-lemma-silly-commutes}
補題 \ref{chow-lemma-silly} において、$f : Y \to X$ を局所有限型とし、
$c \in A^*(Y \to X)$ とする。このとき
$$
c \circ P'_p(E_2) = P'_p(Lf_2^*E_2) \circ c
\quad\text{resp.}\quad
c \circ c'_p(E_2) = c'_p(Lf_2^*E_2) \circ c
$$
が $A^*(Y_2 \to Y)$ で成り立つ。ここで $f_2 : Y_2 \to X_2$ は $f$ の基底変換である。
\end{lemma}

\begin{proof}
$\alpha \in \CH_k(X)$ とする。
$$
\alpha = \alpha_1 + \alpha_2
$$
と書ける。ここで $\alpha_i \in \CH_k(X_i)$ であり、閉埋め込み
$X_i \to X$ による前進像は記法から省略している。すると
$c'_p(E_2) \cap \alpha = c_p(E_2) \cap \alpha_2$,
$c \cap c'_p(E_2) \cap \alpha = c \cap c_p(E_2) \cap \alpha_2$,
$c \cap \alpha = c \cap \alpha_1 + c \cap \alpha_2$, および
$c'_p(Lf_2^*E_2) \cap c \cap \alpha = c_p(Lf_2^*E_2) \cap c \cap \alpha_2$
を確かめられる。補題 \ref{chow-lemma-commutative-chern-perfect} により結論を得る。
\end{proof}

\begin{lemma}
\label{chow-lemma-silly-compose}
補題 \ref{chow-lemma-silly} において、$E_2|_{X_1 \cap X_2}$ が零であると仮定する。このとき
\begin{align*}
P'_1(E_2) & = c'_1(E_2), \\
P'_2(E_2) & = c'_1(E_2)^2 - 2c'_2(E_2), \\
P'_3(E_2) & = c'_1(E_2)^3 - 3c'_1(E_2)c'_2(E_2) + 3c'_3(E_2), \\
P'_4(E_2) & = c'_1(E_2)^4 - 4c'_1(E_2)^2c'_2(E_2) +
4c'_1(E_2)c'_3(E_2) + 2c'_2(E_2)^2 - 4c'_4(E_2),
\end{align*}
などが成り立つ。積は注意 \ref{chow-remark-ring-loc-classes} のものとする。
\end{lemma}

\begin{proof}
零層は階数 $< 1$ をもつので、類 $c'_p(E_2)$ はすべての $p \geq 1$ に対して定義され、
主張は意味をもつ。等式は、補題 \ref{chow-lemma-silly} で得られた類の特徴づけと、
注意 \ref{chow-remark-splitting-principle-perfect} で与えた $E_2$ の Chern 類との交叉積に関する
対応する結果から直ちに従う。
\end{proof}

\begin{lemma}
\label{chow-lemma-silly-sum-c}
$(S, \delta)$ を状況 \ref{chow-situation-setup} のとおりとし、$X$ を $S$ 上局所有限型とする。
$i_j : X_j \to X$, $j = 1, 2$ を、集合論的に $X = X_1 \cup X_2$ を満たす閉埋め込みとする。
$E, F \in D(\mathcal{O}_X)$ を完全対象とする。次を仮定する。
\begin{enumerate}
\item $E$ と $F$ の Chern 類が定義される。
\item 制限 $E|_{X_1 \cap X_2}$ と $F|_{X_1 \cap X_2}$ は、コホモロジー次数 $0$ に置かれた、
それぞれ階数 $< p$ と $< q$ の有限局所自由 $\mathcal{O}_{X_1}$-加群と同型である。
\end{enumerate}
注意 \ref{chow-remark-ring-loc-classes} の記法により
$$
c^{(p)}(E) = 1 + c_1(E) + \ldots + c_{p - 1}(E) +
c'_p(E|_{X_2}) + c'_{p + 1}(E|_{X_2}) + \ldots \in A^{(p)}(X_2 \to X)
$$
とおく。ここで $c'_p(E|_{X_2})$ は補題 \ref{chow-lemma-silly} のものとする。
$c^{(q)}(F)$ と $c^{(p + q)}(E \oplus F)$ も同様に定める。このとき
$c^{(p + q)}(E \oplus F) = c^{(p)}(E)c^{(q)}(F)$ が
$A^{(p + q)}(X_2 \to X)$ で成り立つ。
\end{lemma}

\begin{proof}
補題 \ref{chow-lemma-silly} における類の特徴づけと、補題
\ref{chow-lemma-additivity-on-perfect} の加法性から直ちに従う。
\end{proof}

\begin{lemma}
\label{chow-lemma-silly-sum-P}
$(S, \delta)$ を状況 \ref{chow-situation-setup} のとおりとし、$X$ を $S$ 上局所有限型とする。
$i_j : X_j \to X$, $j = 1, 2$ を、集合論的に $X = X_1 \cup X_2$ を満たす閉埋め込みとする。
$E, F \in D(\mathcal{O}_{X_2})$ を完全対象とする。次を仮定する。
\begin{enumerate}
\item $E$ と $F$ の Chern 類が定義される。
\item 制限 $E|_{X_1 \cap X_2}$ と $F|_{X_1 \cap X_2}$ は零である。
\end{enumerate}
$P'_p(E), P'_p(F), P'_p(E \oplus F) \in A^p(X_2 \to X)$, $p \geq 0$ を、
補題 \ref{chow-lemma-silly} で構成される類とする。このとき
$P'_p(E \oplus F) = P'_p(E) + P'_p(F)$ である。
\end{lemma}

\begin{proof}
補題 \ref{chow-lemma-silly} における類の特徴づけと、補題
\ref{chow-lemma-additivity-on-perfect} の加法性から直ちに従う。
\end{proof}

\begin{lemma}
\label{chow-lemma-silly-tensor-invertible}
補題 \ref{chow-lemma-silly} において、$E_2$ が一定階数 $0$ をもつと仮定する。
$\mathcal{L}$ を可逆 $\mathcal{O}_X$-加群とする。このとき
$$
c'_i(E_2 \otimes \mathcal{L}) =
\sum\nolimits_{j = 0}^i
\binom{- i + j}{j} c'_{i - j}(E_2) c_1(\mathcal{L})^j
$$
である。
\end{lemma}

\begin{proof}
階数に関する仮定から、$E_2|_{X_1 \cap X_2}$ は零である。
したがって $c'_i(E_2)$ はすべての $i \geq 1$ に対して定義され、主張は意味をもつ。
実際の等式は、補題 \ref{chow-lemma-chern-class-perfect-tensor-invertible} と、
補題 \ref{chow-lemma-silly} における $c'_i$ の特徴づけから直ちに従う。
\end{proof}

\begin{lemma}
\label{chow-lemma-silly-tensor-product}
状況 \ref{chow-situation-setup} において、$X$ を $S$ 上局所有限型とする。
$$
X = X_1 \cup X_2 = X'_1 \cup X'_2
$$
を、$X$ を閉部分スキームの集合論的和として表す二通りの方法とする。
$E$, $E'$ を、Chern 類が定義される $D(\mathcal{O}_X)$ の完全対象とする。
$E|_{X_1}$ と $E'|_{X'_1}$ は零である\footnote{これらの制限が、それぞれ階数 $< p$ および
$< p'$ の有限局所自由加群と同型であるとのみ仮定する、この補題の変種もおそらく存在する。}
と仮定する（$i = 1, 2$）。補題 \ref{chow-lemma-silly} で構成される次の類を考える。
\begin{enumerate}
\item $r = P'_0(E) \in A^0(X_2 \to X)$ および
$r' = P'_0(E') \in A^0(X'_2 \to X)$。
\item $\gamma_p = c'_p(E|_{X_2}) \in A^p(X_2 \to X)$ および
$\gamma'_p = c'_p(E'|_{X'_2}) \in A^p(X'_2 \to X)$。
\item $\chi_p = P'_p(E|_{X_2}) \in A^p(X_2 \to X)$ および
$\chi'_p = P'_p(E'|_{X'_2}) \in A^p(X'_2 \to X)$。
\end{enumerate}
このとき
$$
c'_1((E \otimes_{\mathcal{O}_X}^\mathbf{L} E')|_{X_2 \cap X'_2}) =
r \gamma'_1 + r' \gamma_1
$$
が $A^1(X_2 \cap X'_2 \to X)$ で成り立ち、
$$
c'_2((E \otimes_{\mathcal{O}_X}^\mathbf{L} E')|_{X_2 \cap X'_2}) =
r \gamma'_2 + r' \gamma_2 + {r \choose 2} (\gamma'_1)^2 +
(rr' - 1) \gamma'_1\gamma_1 + {r' \choose 2} \gamma_1^2
$$
が $A^2(X_2 \cap X'_2 \to X)$ で成り立つ。高次の Chern 類についても同様である。
また
$$
P'_p((E \otimes_{\mathcal{O}_X}^\mathbf{L} E')|_{X_2 \cap X'_2}) =
\sum\nolimits_{p_1 + p_2 = p}
{p \choose p_1} \chi_{p_1} \chi'_{p_2}
$$
が $A^p(X_2 \cap X'_2 \to X)$ で成り立つ。
\end{lemma}

\begin{proof}
まず主張が意味をもつことを確認する。実際
$X = (X_2 \cap X'_2) \cup Y$ である。ただし
$Y = (X_1 \cap X'_1) \cup (X_1 \cap X'_2) \cup (X_2 \cap X'_1)$ とする。
対象 $E \otimes_{\mathcal{O}_X}^\mathbf{L} E'$ の $Y$ への制限は零である。
実際の等式は、補題 \ref{chow-lemma-silly} におけるこれらの類の特徴づけと、補題
\ref{chow-lemma-chern-classes-perfect-tensor-product} の等式から従う。詳細は省略する。
\end{proof}







\section{無限遠での Gysin}
\label{chow-section-gysin-at-infty}

\noindent
この節では、次の補題で構成する双変類を扱う。読者には、この節の残りを飛ばすことを勧める。

\begin{lemma}
\label{chow-lemma-gysin-at-infty}
$(S, \delta)$ を状況 \ref{chow-situation-setup} のとおりとし、$X$ を $S$ 上局所有限型とする。
$b : W \to \mathbf{P}^1_X$ を、$\mathbf{A}^1_X$ 上で同型となるスキームの固有射とする。
補集合が $\mathbf{A}^1_X$ である因子 $D_\infty \subset \mathbf{P}^1_X$ の逆像を
$i_\infty : W_\infty \to W$ と表す。このとき標準的な双変類
$$
C \in A^0(W_\infty \to X)
$$
が存在し、$\alpha \in \CH_k(X)$ に対して
$i_{\infty, *}(C \cap \alpha) = i_{0, *}\alpha$ を満たす。
任意の局所有限型基底変換 $X' \to X$ の後も同様である。
\end{lemma}

\begin{proof}
$\alpha \in \CH_k(X)$ が与えられたとする。補題 \ref{chow-lemma-exact-sequence-open} により、
$b^{-1}(\mathbf{A}^1_X)$ 上で $\alpha$ の平坦引き戻しに制限される
$\beta \in \CH_{k + 1}(W)$ が存在する。別の $\beta$ の選択と先の $\beta$ との差は、
補題 \ref{chow-lemma-restrict-to-open} により $W_\infty$ 上に台をもつサイクルである。
(\ref{chow-equation-zero-infty}) の有効 Cartier 因子 $D_\infty \subset \mathbf{P}^1_X$ の
法束は自明なので、Gysin 準同型 $i_\infty^*$ は $W_\infty$ 上に台をもつサイクル類を零にする。
注意 \ref{chow-remark-gysin-on-cycles} を見よ。したがって
$C \cap \alpha = i_\infty^*\beta$ とおくことは良定義である。

\medskip\noindent
$W_\infty$ と $W_0 = X \times \{0\}$ は、$\mathbf{P}^1_X$ において有理同値な有効 Cartier 因子
$D_0, D_\infty$ の引き戻しなので、$i_\infty^*\beta$ と $i_0^*\beta$ は $W$ 上の同じサイクル類に写る。
すなわち補題 \ref{chow-lemma-support-cap-effective-Cartier} により、両者はともに類
$c_1(\mathcal{O}_{\mathbf{P}^1_X}(1)) \cap \beta$ を表す。
$\beta$ の選び方から、$W_0 = X \times \{0\}$ 上のサイクルとして
$i_0^*\beta = \alpha$ である。例えば補題 \ref{chow-lemma-relative-effective-cartier} を見よ。
よって補題の主張どおり $i_{\infty, *}(C \cap \alpha) = i_{0, *}\alpha$ である。

\medskip\noindent
$b$ に関する仮定は、任意の局所有限型基底変換 $X' \to X$ で保たれる。
したがって上と同じ構成により演算 $C \cap - : \CH_k(X') \to \CH_k(W'_\infty)$ を得る。
この演算族が双変類を定めることを示すため、図式
$$
\xymatrix{
& & & \CH_*(X) \ar[d]^{\text{flat pullback}} \\
\CH_{* + 1}(W_\infty) \ar[r] \ar[rd]^0 &
\CH_{* + 1}(W) \ar[d]^{i_\infty^*} \ar[rr]^{\text{flat pullback}} & &
\CH_{* + 1}(\mathbf{A}^1_X) \ar[r] \ar@{..>}[lld]^{C \cap -} &
0 \\
& \CH_*(W_\infty)
}
$$
を、表示した $X$ に対して考え、さらに任意の $X' \to X$ によるこの図式の基底変換を考える。
平坦引き戻しと $i_\infty^*$ は双変演算である。補題
\ref{chow-lemma-flat-pullback-bivariant} と \ref{chow-lemma-gysin-bivariant} を見よ。
すると形式的な議論（スキームとその Chow 群の巨大な図式を用いる）により、
点線の矢印が双変演算であることがわかる。
\end{proof}

\begin{lemma}
\label{chow-lemma-base-change-gysin-at-infty}
補題 \ref{chow-lemma-gysin-at-infty} において、$X' \to X$ を局所有限型な射とする。
$b$ と $i_\infty$ の基底変換を、それぞれ $b' : W' \to \mathbf{P}^1_{X'}$ と
$i'_\infty : W'_\infty \to W'$ と表す。このとき、$b'$ を用いて補題
\ref{chow-lemma-gysin-at-infty} と同様に構成される類 $C' \in A^0(W'_\infty \to X')$ は、
$C$ の制限（注意 \ref{chow-remark-restriction-bivariant}）である。
\end{lemma}

\begin{proof}
構成と、平坦引き戻しおよび $i_\infty^*$ にも同様の主張が成り立つことから直ちに従う。
\end{proof}

\begin{lemma}
\label{chow-lemma-gysin-at-infty-independent}
補題 \ref{chow-lemma-gysin-at-infty} において、$g : W' \to W$ を
$\mathbf{A}^1_X$ 上で同型となる固有射とする。$C' \in A^0(W'_\infty \to X)$ と
$C \in A^0(W_\infty \to X)$ を、補題 \ref{chow-lemma-gysin-at-infty} で構成される類とする。
このとき $g_{\infty, *} \circ C' = C$ が $A^0(W_\infty \to X)$ で成り立つ。
\end{lemma}

\begin{proof}
$b' = b \circ g : W' \to \mathbf{P}^1_X$ とおく。包含射を
$i'_\infty : W'_\infty \to W'$ と表し、$g$ の制限を
$g_\infty : W'_\infty \to W_\infty$ と表す。
$\alpha \in \CH_k(X)$ が与えられたとき、$(b')^{-1}\mathbf{A}^1_X$ 上で $\alpha$ の
平坦引き戻しに制限される $\beta' \in \CH_{k + 1}(W')$ を選ぶ。
すると $\beta = g_*\beta' \in \CH_{k + 1}(W)$ は、$b^{-1}\mathbf{A}^1_X$ 上で
$\alpha$ の平坦引き戻しに制限される。補題 \ref{chow-lemma-closed-in-X-gysin} により
$i_\infty^*\beta = g_{\infty, *}(i'_\infty)^*\beta'$ である。
この事実と、任意の局所有限型な射 $X' \to X$ による基底変換後の対応する事実が、
補題の主張を与える。
\end{proof}

\begin{lemma}
\label{chow-lemma-homomorphism-pre}
補題 \ref{chow-lemma-gysin-at-infty} において
$C \circ (W_\infty \to X)_* \circ i_\infty^* = i_\infty^*$ である。
\end{lemma}

\begin{proof}
$\beta \in \CH_{k + 1}(W)$ とする。$\mathbf{P}^1$ の $0$ 上のファイバーの閉埋め込みを
$i_0 : X = X \times \{0\} \to W$ と表す。このとき
$(W_\infty \to X)_* i_\infty^* \beta = i_0^*\beta$ が $\CH_k(X)$ で成り立つ。
実際、$i_{\infty, *}i_\infty^*\beta$ と $i_{0, *}i_0^*\beta$ は $W$ 上で同じ類を表し
（例えば補題 \ref{chow-lemma-support-cap-effective-Cartier} による）、したがって $X$ 上の同じ類へ前進される。
$\beta$ の $b^{-1}(\mathbf{A}^1_X)$ への制限は
$i_0^*\beta = (W_\infty \to X)_* i_\infty^* \beta$ の平坦引き戻しに制限される。
これは $i_0$ による引き戻しの後で確かめられるからである。補題
\ref{chow-lemma-linebundle} と \ref{chow-lemma-linebundle-formulae} を見よ。
したがって $(W_\infty \to X)_*i_\infty^*\beta$ の $C$ による像を計算するときに
$\beta$ を用いることができ、求める結果を得る。
\end{proof}

\begin{lemma}
\label{chow-lemma-gysin-at-infty-commutes}
補題 \ref{chow-lemma-gysin-at-infty} において、$f : Y \to X$ を局所有限型な射とし、
$c \in A^*(Y \to X)$ とする。このとき $C \circ c = c \circ C$ が
$A^*(W_\infty \times_X Y \to X)$ で成り立つ。
\end{lemma}

\begin{proof}
Cartesian 正方形からなる可換図式
$$
\xymatrix{
W_\infty \times_X Y \ar@{=}[r] &
W_{Y, \infty} \ar[r]_{i_{Y, \infty}} \ar[d] &
W_Y \ar[r]_{b_Y} \ar[d] &
\mathbf{P}^1_Y \ar[r]_{p_Y} \ar[d] &
Y \ar[d]^f \\
& W_\infty \ar[r]^{i_\infty} &
W \ar[r]^b &
\mathbf{P}^1_X \ar[r]^p &
X
}
$$
を考える。元 $\alpha \in \CH_k(X)$ に対し、$b^{-1}(\mathbf{A}^1_X)$ への制限が
$\alpha$ の平坦引き戻しとなる $\beta \in \CH_{k + 1}(W)$ を選ぶ。
すると $c \cap \beta$ は $\CH_*(W_Y)$ の類であり、その $b_Y^{-1}(\mathbf{A}^1_Y)$ への制限は
$c \cap \alpha$ の平坦引き戻しである。さらに
$$
i_{Y, \infty}^*(c \cap \beta) = c \cap i_\infty^*\beta
$$
である。これは $c$ が双変類だからである。この等式はまさに
$C \cap c \cap \alpha = c \cap C \cap \alpha$ を意味する。
任意の局所有限型基底変換 $X' \to X$ の後にも同じ議論が成り立つ。これで補題が証明された。
\end{proof}





\section{局所化 Chern 類の準備}
\label{chow-section-preparation-localized-chern-II}

\noindent
この節では、次の補題で構成する双変類のいくつかの性質を論じる。
読者には、この節の残りを飛ばすことを勧める。

\begin{lemma}
\label{chow-lemma-localized-chern-pre}
$(S, \delta)$ を状況 \ref{chow-situation-setup} のとおりとし、$X$ を $S$ 上局所有限型とする。
$Z \subset X$ を閉部分スキームとし、
$$
b : W \longrightarrow \mathbf{P}^1_X
$$
をスキームの固有射とする。$Q \in D(\mathcal{O}_W)$ を完全対象とする。
補集合が $\mathbf{A}^1_X$ である因子 $D_\infty \subset \mathbf{P}^1_X$ の逆像を
$W_\infty \subset W$ と表す。次を仮定する。
\begin{enumerate}
\item[(A0)] $Q$ の Chern 類が定義される（節 \ref{chow-section-pre-derived}）。
\item[(A1)] $b$ は $\mathbf{A}^1_X$ 上で同型である。
\item[(A2)] $W_\infty$ の $X \setminus Z$ 上にあるすべての点を含む閉部分スキーム
$T \subset W_\infty$ が存在し、$Q|_T$ は零である、resp.\ コホモロジー次数 $0$ に置かれた
階数 $< p$ の有限局所自由 $\mathcal{O}_T$-加群と同型である。
\end{enumerate}
このとき標準的な双変類
$$
P'_p(Q),\text{ resp. }c'_p(Q) \in A^p(Z \to X)
$$
が存在し、
$(Z \to X)_* \circ P'_p(Q) = P_p(Q|_{X \times \{0\}})$,
resp.\ $(Z \to X)_* \circ c'_p(Q) = c_p(Q|_{X \times \{0\}})$ を満たす。
\end{lemma}

\begin{proof}
$Z$ の逆像を $E \subset W_\infty$ と表す。このとき
$W_\infty = T \cup E$ であり、$b$ は固有射 $E \to Z$ を誘導する。
補題 \ref{chow-lemma-gysin-at-infty} で構成される双変類を
$C \in A^0(W_\infty \to X)$ と表す。また、補題 \ref{chow-lemma-silly} で構成される
$A^p(E \to W_\infty)$ の双変類を $P'_p(Q|_E)$, resp.\ $c'_p(Q|_E)$ と表す。
これは意味をもつ。実際、仮定 (A2) により $(Q|_E)|_{T \cap E}$ は零である、resp.\ コホモロジー次数 $0$ に置かれた
階数 $< p$ の有限局所自由
$\mathcal{O}_{E \cap T}$-加群と同型である。そこで
$$
P'_p(Q) = (E \to Z)_* \circ P'_p(Q|_E) \circ C,\text{ resp. }
c'_p(Q) = (E \to Z)_* \circ c'_p(Q|_E) \circ C
$$
と定義する。補題 \ref{chow-lemma-push-proper-bivariant} により、これは双変類である。
$E \to Z \to X$ は $E \to W_\infty \to W \to X$ に等しいので
\begin{align*}
(Z \to X)_* \circ c'_p(Q)
& =
(W \to X)_* \circ i_{\infty, *} \circ (E \to W_\infty)_*
\circ c'_p(Q|_E) \circ C \\
& =
(W \to X)_* \circ i_{\infty, *} \circ c_p(Q|_{W_\infty}) \circ C \\
& =
(W \to X)_* \circ c_p(Q) \circ i_{\infty, *} \circ C \\
& =
(W \to X)_*\circ c_p(Q) \circ i_{0, *} \\
& =
(W \to X)_* \circ i_{0, *} \circ c_p(Q|_{X \times \{0\}}) \\
& =
c_p(Q|_{X \times \{0\}})
\end{align*}
を得る。第二の等式は補題 \ref{chow-lemma-silly-silly} による。
第三の等式は $c_p(Q)$ が双変類だからであり、第四の等式は補題
\ref{chow-lemma-gysin-at-infty} による。第五の等式は再び $c_p(Q)$ が双変類だからである。
最後の等式は、上のように $W_0$ を $X$ と同一視すると
$(W_0 \to W) \circ (W \to X)$ が $X$ 上の恒等射だからである。
まったく同じ等式列により $P'_p(Q)$ に対する性質も証明できる。
\end{proof}

\begin{lemma}
\label{chow-lemma-base-change-localized-chern-pre}
補題 \ref{chow-lemma-localized-chern-pre} において、$X' \to X$ を局所有限型な射とする。
$Z$, $b : W \to \mathbf{P}^1_X$, $T \subset W_\infty$ の基底変換を、それぞれ
$Z'$, $b' : W' \to \mathbf{P}^1_{X'}$, $T' \subset W'_\infty$ と表す。
$Q' = (W' \to W)^*Q$ とおく。このとき、$b'$, $Q'$, $T'$ を用いて補題
\ref{chow-lemma-localized-chern-pre} と同様に構成される $A^p(Z' \to X')$ の類
$P'_p(Q')$, resp.\ $c'_p(Q')$ は、$A^p(Z \to X)$ の類
$P'_p(Q)$, resp.\ $c'_p(Q)$ の制限（注意 \ref{chow-remark-restriction-bivariant}）である。
\end{lemma}

\begin{proof}
構成は
$$
P'_p(Q) = (E \to Z)_* \circ P'_p(Q|_E) \circ C,\text{ resp. }
c'_p(Q) = (E \to Z)_* \circ c'_p(Q|_E) \circ C
$$
であった。したがって補題は、$C$ に対する対応する基底変換性（補題
\ref{chow-lemma-base-change-gysin-at-infty}）と、補題 \ref{chow-lemma-silly} で構成される類にも
同じ基底変換性が成り立つことから従う（細部は省略する）。
\end{proof}

\begin{lemma}
\label{chow-lemma-localized-chern-pre-independent}
補題 \ref{chow-lemma-localized-chern-pre} の双変類 $P'_p(Q)$, resp.\ $c'_p(Q)$ は、
閉部分スキーム $T$ の選択に依存しない。さらに、$\mathbf{A}^1_X$ 上で同型となる固有射
$g : W' \to W$ が与えられ、$Q' = g^*Q$ とおくと、
$P'_p(Q) = P'_p(Q')$, resp.\ $c'_p(Q) = c'_p(Q')$ である。
\end{lemma}

\begin{proof}
$T$ からの独立性は補題 \ref{chow-lemma-silly-independent} から直ちに従う。

\medskip\noindent
$g : W' \to W$ を $\mathbf{A}^1_X$ 上で同型となる固有射とする。
$T' = g^{-1}(T) \subset W'_\infty$ とおけば、これは (A2) を満たす閉部分スキームである。
したがって $b' = b \circ g : W' \to \mathbf{P}^1_X$ と $Q'$ に対応する
$A^p(Z \to X)$ の演算 $P'_p(Q')$, resp.\ $c'_p(Q')$ は定義される。
$W'_\infty$ における $Z$ の逆像を $E' \subset W'_\infty$ と表す。構成を思い出すと
$$
c'_p(Q') = (E' \to Z)_* \circ c'_p(Q'|_{E'}) \circ C'
$$
である。ここで $C' \in A^0(W'_\infty \to X)$ かつ
$c'_p(Q'|_{E'}) \in A^p(E' \to W'_\infty)$ である。
補題 \ref{chow-lemma-gysin-at-infty-independent} により $g_{\infty, *} \circ C' = C$ である。
$E'$ はまた、$g_\infty$ による $W'_\infty$ における $E$ の逆像である。
さらに $Q' = g^*Q$ なので、$c'_p(Q'|_{E'})$ は単に、$W'_\infty$ 上のスキームへの
$c'_p(Q|_E)$ の制限である。注意 \ref{chow-remark-restriction-bivariant} を見よ。したがって
\begin{align*}
c'_p(Q')
& = 
(E' \to Z)_* \circ c'_p(Q'|_{E'}) \circ C' \\
& =
(E \to Z)_* \circ (E' \to E)_* \circ c'_p(Q|_E) \circ C' \\
& =
(E \to Z)_* \circ c'_p(Q|_E) \circ g_{\infty, *} \circ C' \\
& =
(E \to Z)_* \circ c'_p(Q|_E) \circ C \\
& =
c'_p(Q)
\end{align*}
を得る。第三の等式では、$c'_p(Q|_E)$ が双変類であるため固有前進像と可換であることを用いた。
等式 $P'_p(Q) = P'_p(Q')$ もまったく同様に証明される。
\end{proof}

\begin{lemma}
\label{chow-lemma-homomorphism}
補題 \ref{chow-lemma-localized-chern-pre} において、$Q|_T$ は階数 $< p$ の有限局所自由
$\mathcal{O}_T$-加群と同型であると仮定する。補題 \ref{chow-lemma-gysin-at-infty} の類を
$C \in A^0(W_\infty \to X)$ と表す。このとき
$$
C \circ c_p(Q|_{X \times \{0\}}) =
C \circ (Z \to X)_* \circ c'_p(Q) = c_p(Q|_{W_\infty}) \circ C
$$
である。
\end{lemma}

\begin{proof}
第一の等式は、補題 \ref{chow-lemma-localized-chern-pre} により
$c_p(Q|_{X \times \{0\}}) =
(Z \to X)_* \circ c'_p(Q)$ だからである。第二の等式はサイクル類ごとに証明してよい
（補題 \ref{chow-lemma-bivariant-zero} を見よ）。補題の双変類の構成は基底変換と両立するので、
ある $\alpha \in \CH_k(X)$ が与えられ、
$C \cap (Z \to X)_*(c'_p(Q) \cap \alpha) =
c_p(Q|_{W_\infty}) \cap C \cap \alpha$ を示せばよいと仮定できる。実際
\begin{align*}
C \cap (Z \to X)_*(c'_p(Q) \cap \alpha)
& =
C \cap (Z \to X)_* (E \to Z)_*(c'_p(Q|_E) \cap C \cap \alpha) \\
& =
C \cap (W_\infty \to X)_*(E \to W_\infty)_*(c'_p(Q|_E) \cap C \cap \alpha) \\
& =
C \cap (W_\infty \to X)_*(E \to W_\infty)_*(c'_p(Q|_E) \cap i_\infty^*\beta) \\
& =
C \cap (W_\infty \to X)_*(c_p(Q|_{W_\infty}) \cap i_\infty^*\beta) \\
& =
C \cap (W_\infty \to X)_*i_\infty^*(c_p(Q) \cap \beta) \\
& =
i_\infty^*(c_p(Q) \cap \beta) \\
& =
c_p(Q|_{W_\infty}) \cap i_\infty^*\beta \\
& =
c_p(Q|_{W_\infty}) \cap C \cap \alpha
\end{align*}
となり、求める等式を得る。第一の等式では、$E \subset W_\infty$ を $Z$ の逆像とし、
補題 \ref{chow-lemma-silly} で構成される類を $c'_p(Q|_E)$ としたときの
$c'_p(Q) = (E \to Z)_* \circ c'_p(Q|_E) \circ C$ を用いた。
第二の等式は、$E \to Z \to X$ が $E \to W_\infty \to X$ に等しいという主張である。
第三の等式では、$b^{-1}(\mathbf{A}^1_X)$ への制限が $\alpha$ の平坦引き戻しとなる
$\beta \in \CH_{k + 1}(W)$ を選ぶ。このとき構成により
$C \cap \alpha = i_\infty^*\beta$ である。第四の等式は補題
\ref{chow-lemma-silly-silly} による。第五の等式では、$c_p(Q)$ が双変類であり、したがって
$i_\infty^*$ と可換であることを用いた。第六の等式は補題
\ref{chow-lemma-homomorphism-pre} による。第七の等式では、再び $c_p(Q)$ が双変類であることを用いた。
最後の等式は $C \cap \alpha = i_\infty^*\beta$ による。
\end{proof}

\begin{lemma}
\label{chow-lemma-homomorphism-commute}
補題 \ref{chow-lemma-localized-chern-pre} において、$Y \to X$ を局所有限型な射とし、
$c \in A^*(Y \to X)$ を双変類とする。このとき
$$
P'_p(Q) \circ c = c \circ P'_p(Q)
\quad\text{resp.}\quad
c'_p(Q) \circ c = c \circ c'_p(Q)
$$
が $A^*(Y \times_X Z \to X)$ で成り立つ。
\end{lemma}

\begin{proof}
$Z$ の逆像を $E \subset W_\infty$ とする。構成を思い出すと
$P'_p(Q) = (E \to Z)_* \circ P'_p(Q|_E) \circ C$,
resp.\ $c'_p(Q) = (E \to Z)_* \circ c'_p(Q|_E) \circ C$ である。
ここで $C$ は補題 \ref{chow-lemma-gysin-at-infty} のもの、
$P'_p(Q|_E)$, resp.\ $c'_p(Q|_E)$ は補題 \ref{chow-lemma-silly} のものである。
補題 \ref{chow-lemma-gysin-at-infty-commutes} により $C$ は $c$ と可換であり、補題
\ref{chow-lemma-silly-commutes} により $P'_p(Q|_E)$, resp.\ $c'_p(Q|_E)$ は $c$ と可換である。
さらに $c$ は双変類なので、定義により $E \to Z$ による固有前進像と可換である。
これで証明が完了する。
\end{proof}

\begin{lemma}
\label{chow-lemma-localized-chern-pre-compose}
補題 \ref{chow-lemma-localized-chern-pre} において、$Q|_T$ が零であると仮定する。
$A^*(Z \to X)$ において
\begin{align*}
P'_1(Q) & = c'_1(Q), \\
P'_2(Q) & = c'_1(Q)^2 - 2c'_2(Q), \\
P'_3(Q) & = c'_1(Q)^3 - 3c'_1(Q)c'_2(Q) + 3c'_3(Q), \\
P'_4(Q) & = c'_1(Q)^4 - 4c'_1(Q)^2c'_2(Q) +
4c'_1(Q)c'_3(Q) + 2c'_2(Q)^2 - 4c'_4(Q),
\end{align*}
などが成り立つ。積は注意 \ref{chow-remark-ring-loc-classes} のものとする。
\end{lemma}

\begin{proof}
零層は階数 $< 1$ をもつので、類 $c'_p(Q)$ はすべての $p \geq 1$ に対して定義され、
主張は意味をもつ。補題 \ref{chow-lemma-localized-chern-pre} の証明では、補題
\ref{chow-lemma-gysin-at-infty} の双変類 $C \in A^0(W_\infty \to X)$ と、
$W_\infty$ における $Z$ の逆像 $E$ への $Q$ の制限 $Q|_E$ に対する、補題
\ref{chow-lemma-silly} の双変類 $P'_p(Q|_E)$ と $c'_p(Q|_E)$ を用いて、
類 $P'_p(Q)$ と $c'_p(Q)$ を構成した。補題 \ref{chow-lemma-silly-compose} により、
$P'_p(Q|_E)$ と $c'_p(Q|_E)$ の間には求める関係が成り立つ。構成を思い出すと
$$
P'_p(Q) = (E \to Z)_* \circ P'_p(Q|_E) \circ C
\quad\text{and}\quad
c'_p(Q) = (E \to Z)_* \circ c'_p(Q|_E) \circ C
$$
である。証明を終えるには、注意 \ref{chow-remark-ring-loc-classes} で定義される、
類 $a_p = c'_p(Q)$ 上の積と類 $b_p = c'_p(Q|_E)$ 上の積が両立すること、すなわち
$$
a_{p_1}a_{p_2} \ldots a_{p_r} =
(E \to Z)_* \circ b_{p_1}b_{p_2} \ldots b_{p_r} \circ C
$$
を示せば十分である。細部の一部は省略する。$r = 1$ なら主張は成り立つ。
$r > 1$ とする。注意 \ref{chow-remark-res-push} により、注意
\ref{chow-remark-ring-loc-classes} の積は、積 $a_{p_1}a_{p_2} \ldots a_{p_r}$, resp.\ $b_{p_1}b_{p_2} \ldots b_{p_r}$ の
因子の間に $(Z \to X)_*$, resp.\ $(E \to W_\infty)_*$ を挿入し、双変類として合成をとることで定められる。
補題 \ref{chow-lemma-silly} により
$$
(E \to W_\infty)_* \circ b_{p_i} = c_{p_i}(Q|_{W_\infty})
$$
であり、補題 \ref{chow-lemma-homomorphism} により
$$
C \circ (Z \to X)_* \circ a_{p_i} = c_{p_i}(Q|_{W_\infty}) \circ C
$$
が $i = 2, \ldots, r$ に対して成り立つ。計算すると、求める等式の左右両辺はともに
$$
(E \to Z)_* \circ c'_{p_1}(Q|_E) \circ
c_{p_2}(Q|_{W_\infty}) \circ \ldots \circ
c_{p_r}(Q|_{W_\infty}) \circ C
$$
に簡約される。これで証明が完了する。
\end{proof}

\begin{lemma}
\label{chow-lemma-localized-chern-pre-sum-c}
補題 \ref{chow-lemma-localized-chern-pre} において、$Q|_T$ が階数 $< p$ の有限局所自由
$\mathcal{O}_T$-加群と同型であると仮定する。さらに、Chern 類が定義される別の完全対象
$Q' \in D(\mathcal{O}_W)$ があり、$Q'|_T$ はコホモロジー次数 $0$ に置かれた階数 $< p'$ の
有限局所自由 $\mathcal{O}_T$-加群と同型であると仮定する。
注意 \ref{chow-remark-ring-loc-classes} の記法により
$$
c^{(p)}(Q) = 1 + c_1(Q|_{X \times \{0\}}) + \ldots +
c_{p - 1}(Q|_{X \times \{0\}}) +
c'_{p}(Q) + c'_{p + 1}(Q) + \ldots
$$
を $A^{(p)}(Z \to X)$ の元として定める。ここで $i \geq p$ に対する $c'_i(Q)$ は
補題 \ref{chow-lemma-localized-chern-pre} のものとする。$c^{(p')}(Q')$ と
$c^{(p + p')}(Q \oplus Q')$ も同様に定める。このとき
$c^{(p + p')}(Q \oplus Q') = c^{(p)}(Q)c^{(p')}(Q')$ が
$A^{(p + p')}(Z \to X)$ で成り立つ。
\end{lemma}

\begin{proof}
$A^p(X)$ における $c'_i(Q)$ の像は、$i \geq p$ に対して
$c_i(Q|_{X \times \{0\}})$ に等しいことを思い出そう。$Q'$ と $Q \oplus Q'$ についても同様である。
補題 \ref{chow-lemma-localized-chern-pre} を見よ。したがって次数 $< p + p'$ における等式は、
補題 \ref{chow-lemma-additivity-on-perfect} の加法性から従う。

\medskip\noindent
$n \geq p + p'$ とする。補題 \ref{chow-lemma-localized-chern-pre} の証明と同様に、
$Z$ の逆像を $E \subset W_\infty$ と表す。補題 \ref{chow-lemma-silly-sum-c} により、
$A^{(p + p')}(E \to W_\infty)$ において
$$
c^{(p + p')}(Q|_E \oplus Q'|_E) =
c^{(p)}(Q|_E)c^{(p')}(Q'|_E)
$$
である。一方、構成により
$$
c'_p(Q \oplus Q') = (E \to Z)_* \circ c'_p(Q|_E \oplus Q'|_E) \circ C
$$
である。したがって、すべての $i + j = n$ に対して
$$
(E \to Z)_* \circ c^{(p)}_i(Q|_E)c^{(p')}_j(Q'|_E) \circ C
=
c^{(p)}_i(Q)c^{(p')}_j(Q')
$$
が $A^n(Z \to X)$ で成り立つことを示せば十分である。両辺の積は注意
\ref{chow-remark-ring-loc-classes} のものとする。$i \geq p$, $j \geq p'$ のいずれが成り立つか、
または両方が成り立つかに応じて三つの場合がある。

\medskip\noindent
$i \geq p$ かつ $j \geq p'$ と仮定する。この場合の積は、二つの因子の間に
$(E \to W_\infty)_*$, resp.\ $(Z \to X)_*$ を挿入し、双変類として合成をとることで定義される。
注意 \ref{chow-remark-res-push} を見よ。すなわち示すべきことは
$$
(E \to Z)_* \circ c'_i(Q|_E) \circ
(E \to W_\infty)_* \circ c'_j(Q'|_E) \circ C =
c'_i(Q) \circ (Z \to X)_* \circ c'_j(Q')
$$
である。補題 \ref{chow-lemma-silly} により左辺は
$$
(E \to Z)_* \circ c'_i(Q|_E) \circ c_j(Q'|_{W_\infty}) \circ C
$$
に等しい。また $c'_i(Q) = (E \to Z)_* \circ c'_i(Q|_E) \circ C$ なので、右辺は
$$
(E \to Z)_* \circ c'_i(Q|_E) \circ C \circ (Z \to X)_* \circ c'_j(Q')
$$
に等しい。これは補題 \ref{chow-lemma-homomorphism} により直ちに上の式と等しい。

\medskip\noindent
$i \geq p$ かつ $j < p$ と仮定する。この場合に積の定義を展開すると、示すべきことは
$$
(E \to Z)_* \circ c'_i(Q|_E) \circ c_j(Q'|_{W_\infty}) \circ C =
c'_i(Q) \circ c_j(Q'|_{X \times \{0\}})
$$
である。再び $c'_i(Q) = (E \to Z)_* \circ c'_i(Q|_E) \circ C$ を用いると、
$c_j(Q'|_{W_\infty}) \circ C =
C \circ c_j(Q'|_{X \times \{0\}})$ を示せば十分である。これは補題
\ref{chow-lemma-homomorphism} の一部である。

\medskip\noindent
$i < p$ かつ $j \geq p'$ と仮定する。この場合に積の定義を展開すると、示すべきことは
$$
(E \to Z)_* \circ c_i(Q|_E) \circ c'_j(Q'|_E) \circ C =
c_i(Q|_{Z \times \{0\}}) \circ c'_j(Q')
$$
である。しかし $c'_j(Q|_E)$ と $c'_j(Q')$ は双変類なので、Chern 類との交叉積と可換である
（補題 \ref{chow-lemma-cap-commutative-chern}）。したがって
$$
(E \to Z)_* \circ c'_j(Q'|_E) \circ c_i(Q|_{W_\infty}) \circ C =
c'_j(Q') \circ c_i(Q|_{X \times \{0\}})
$$
を示せば十分であり、これは前段落で論じた場合に帰着する。
\end{proof}

\begin{lemma}
\label{chow-lemma-localized-chern-pre-sum-P}
補題 \ref{chow-lemma-localized-chern-pre} において、$Q|_T$ が零であると仮定する。
さらに、Chern 類が定義され、制限 $Q'|_T$ が零である別の完全対象
$Q' \in D(\mathcal{O}_W)$ があると仮定する。この場合、補題
\ref{chow-lemma-localized-chern-pre} で構成される類
$P'_p(Q), P'_p(Q'), P'_p(Q \oplus Q') \in A^p(Z \to X)$ は
$P'_p(Q \oplus Q') = P'_p(Q) + P'_p(Q')$ を満たす。
\end{lemma}

\begin{proof}
これらの類の構成と補題 \ref{chow-lemma-silly-sum-P} から直ちに従う。
\end{proof}









 
\section{局所化 Chern 類}
\label{chow-section-localized-chern}

\noindent
構成の概略。$F$ を体、$X$ を $F$ 上の多様体、$E$ を $D(\mathcal{O}_X)$ の完全対象とし、
$Z \subset X$ を $E|_{X \setminus Z} = 0$ を満たす閉部分スキームとする。
構成したいのは元
$$
c_p(Z \to X, E) \in A^p(Z \to X)
$$
である。そのために図式
$$
\xymatrix{
W \ar[d]_f \ar[r]_q & X \\
\mathbf{P}^1_F
}
$$
と $D(\mathcal{O}_W)$ の完全対象 $Q$ を、次を満たすように構成する。
\begin{enumerate}
\item $f$ は平坦で、$f$, $q$ は固有である。$t \in \mathbf{P}^1_F$ に対し、$f$ のファイバーを
$W_t$、$q$ の制限を $q_t : W_t \to X$、また $Q_t = Q|_{W_t}$ と表す。
\item $t \in \mathbf{A}^1_F$ に対し、$q_t : W_t \to X$ は同型で $Q_t = q_t^*E$ である。
\item $q_\infty : W_\infty \to X$ は $X \setminus Z$ 上で同型である。
\item $T \subset W_\infty$ を $q_\infty^{-1}(X \setminus Z)$ の閉包とすると、
$Q_\infty|_T$ は零である。
\end{enumerate}
これを $t \in \mathbf{P}^1$ でパラメータづけられた族 $\{(W_t, Q_t)\}$ と考えるのが要点である。
$t \not = \infty$ なら、$Q_t = X$ の Chow 群上で $c_p(Q_t)$ は単に $c_p(E)$ である。
一方 $t = \infty$ では、$Q_\infty|_T = 0$ なので、$c_p(Q_\infty)$ は $Q_\infty$ 上の類を
$E = q_\infty^{-1}(Z)$ 上に台をもつ類へ送る。$E$ を $q_\infty : W_\infty \to X$ の例外集合と考える。
任意の $\alpha \in \CH_*(X)$ はサイクルの「族」$\alpha_t \in \CH_*(W_t)$ を生じるので、
$c_p(Z \to X, E) \cap \alpha$ を前進像
$(E \to Z)_*(c_p(Q_\infty) \cap \alpha_\infty)$ と定義することが自然である。

\medskip\noindent
これを実現するには二つの主要な材料がある。(1) $W$ と $Q$ の構成は、一種の代数的な
Macpherson のグラフ構成であり、More on Flatness の節
\ref{flat-section-blowup-complexes-III} で行われる。(2) $W$ と $Q$ が与えられたときの
実際の類の構成は、節 \ref{chow-section-preparation-localized-chern-II} で、節
\ref{chow-section-gysin-at-infty} と \ref{chow-section-preparation-localized-chern} に基づいて行われる。

\begin{situation}
\label{chow-situation-loc-chern}
$(S, \delta)$ を状況 \ref{chow-situation-setup} のとおりとし、$X$ を $S$ 上局所有限型なスキームとする。
$i : Z \to X$ を閉埋め込み、$E \in D(\mathcal{O}_X)$ を対象とし、$p \geq 0$ とする。
次を仮定する。
\begin{enumerate}
\item $E$ は $D(\mathcal{O}_X)$ の完全対象である。
\item 制限 $E|_{X \setminus Z}$ は零である、resp.\ コホモロジー次数 $0$ に置かれた
階数 $< p$ の有限局所自由 $\mathcal{O}_{X \setminus Z}$-加群と同型である。
\item 次の少なくとも一つ\footnote{初読ではこの技術的条件を無視してよい。注意
\ref{chow-remark-loc-chern} の議論を見よ。}が成り立つ。
(a) $X$ は準コンパクトである。
(b) $X$ の既約成分は準コンパクトである。
(c) $E$ を表す有限局所自由 $\mathcal{O}_X$-加群の局所有界複体が存在する。
(d) $S$ 上局所有限型なスキームの射 $X \to X'$ が存在し、$E$ は $X'$ 上の完全対象の
引き戻しであり、$X'$ の既約成分は準コンパクトである。
\end{enumerate}
\end{situation}

\begin{lemma}
\label{chow-lemma-independent-loc-chern}
状況 \ref{chow-situation-loc-chern} において、標準的な双変類
$$
P_p(Z \to X, E) \in A^p(Z \to X),
\quad\text{resp.}\quad
c_p(Z \to X, E) \in A^p(Z \to X)
$$
が存在し、$X$ 上の双変類として
\begin{equation}
\label{chow-equation-defining-property-localized-classes}
i_* \circ P_p(Z \to X, E) = P_p(E),
\quad\text{resp.}\quad
i_* \circ c_p(Z \to X, E) = c_p(E)
\end{equation}
を満たす（$\circ$ は補題 \ref{chow-lemma-push-proper-bivariant} のものとする）。
\end{lemma}

\begin{proof}
これらの双変類を次のように構成する。
$$
b : W \longrightarrow \mathbf{P}^1_X
\quad\text{and}\quad
T \longrightarrow W_\infty
\quad\text{and}\quad
Q
$$
を、More on Flatness の節 \ref{flat-section-blowup-complexes-III} と補題
\ref{flat-lemma-graph-construction} で構成されるブローアップ、閉埋め込み、および
$D(\mathcal{O}_W)$ の完全対象 $Q$ とする。$T' \subset T$ を、$Q|_{T'}$ が零である、resp.\ コホモロジー次数 $0$ に置かれた
階数 $< p$ の有限局所自由 $\mathcal{O}_{T'}$-加群と
同型であるような開かつ閉な部分スキームとする。状況
\ref{chow-situation-loc-chern} の条件 (2) により、射
$$
T' \to T \to W_\infty \to X
$$
はすべて、$X$ の開部分スキーム $X \setminus Z$ 上でスキームの同型である。
以下で $Q$ の Chern 類が定義されることを確認する。
構成により $Q|_{X \times \{0\}} \cong E$ であることを思い出すと、$W, b, Q, T'$ を用いて
補題 \ref{chow-lemma-localized-chern-pre} で構成される双変類から
$$
P_p(Z \to X, E) = P'_p(Q) \in A^p(Z \to X)
$$
および
$$
c_p(Z \to X, E) = c'_p(Q) \in A^p(Z \to X)
$$
を得る。これらは (\ref{chow-equation-defining-property-localized-classes}) を満たす。

\medskip\noindent
この段落では $Q$ の Chern 類が定義されることを証明する（定義
\ref{chow-definition-defined-on-perfect}）。読者には飛ばすことを勧める。
状況 \ref{chow-situation-loc-chern} の仮定 (3)(a) または (3)(b) が成り立つ、すなわち $X$ の
既約成分が準コンパクトであるとする。このとき $W \to X$ は固有なので、$W$ についても同じである。
したがって補題 \ref{chow-lemma-chern-classes-defined} により、$D(\mathcal{O}_W)$ の任意の完全対象の
Chern 類は定義される。(3)(c) が成り立つ、すなわち $E$ が有限局所自由加群の局所有界複体で
表されるなら、More on Flatness の補題 \ref{flat-lemma-graph-construction} の (5) により、
対象 $Q$ は有限局所自由 $\mathcal{O}_W$-加群の局所有界複体で表される。したがって $Q$ の
Chern 類は定義される。最後に (3)(d) を仮定する。すなわち、$S$ 上局所有限型なスキームの射
$X \to X'$ があり、$E$ は $X'$ 上の完全対象 $E'$ の引き戻しで、$X'$ の既約成分は
準コンパクトであると仮定する。More on Flatness の節
\ref{flat-section-blowup-complexes-III} において三つ組
$(\mathbf{P}^1_{X'}, (\mathbf{P}^1_{X'})_\infty, L(p')^*E')$ から構成される射と完全対象を
$b' : W' \to \mathbf{P}^1_{X'}$ および $Q' \in D(\mathcal{O}_{W'})$ とする。
上の議論により $Q'$ の Chern 類は定義される。$b$ と $b'$ は More on Flatness の補題
\ref{flat-lemma-complex-and-divisor-blowup-pre} の適用によって構成されたので、More on Flatness の補題
\ref{flat-lemma-complex-and-divisor-blowup-base-change} により、$Q = L(W \to W')^*Q'$ を満たす射
$W \to W'$ が存在する。すると補題 \ref{chow-lemma-chern-classes-defined} により、$Q$ の Chern 類は定義される。
\end{proof}

\begin{definition}
\label{chow-definition-localized-chern}
$(S, \delta)$, $X$, $E \in D(\mathcal{O}_X)$, $i : Z \to X$ を状況
\ref{chow-situation-loc-chern} のとおりとする。
\begin{enumerate}
\item 制限 $E|_{X \setminus Z}$ が零なら、すべての $p \geq 0$ に対して、補題
\ref{chow-lemma-independent-loc-chern} の構成により
$$
P_p(Z \to X, E) \in A^p(Z \to X)
$$
を定義し、さらに {\it 局所化 Chern 指標}を
$$
ch(Z \to X, E) =
\sum\nolimits_{p = 0, 1, 2, \ldots} \frac{P_p(Z \to X, E)}{p!}
\quad\text{in}\quad \prod\nolimits_{p \geq 0} A^p(Z \to X) \otimes \mathbf{Q}
$$
で定義する。
\item 制限 $E|_{X \setminus Z}$ が、コホモロジー次数 $0$ に置かれた階数 $< p$ の
有限局所自由 $\mathcal{O}_{X \setminus Z}$-加群と同型なら、補題
\ref{chow-lemma-independent-loc-chern} の構成により {\it 局所化第 $p$ Chern 類}
$c_p(Z \to X, E)$ を定義する。
\end{enumerate}
\end{definition}

\noindent
定義の状況で $E|_{X \setminus Z}$ が零であると仮定する。上で等式
(\ref{chow-equation-defining-property-localized-classes}) を示したので、確かに
$$
i_* \circ ch(Z \to X, E) = ch(E)
$$
が $A^*(X) \otimes \mathbf{Q}$ で成り立つ。

\medskip\noindent
次は重要な整合性の確認である。

\begin{lemma}
\label{chow-lemma-base-change-loc-chern}
状況 \ref{chow-situation-loc-chern} において、$f : X' \to X$ を局所有限型なスキームの射とする。
$E' = f^*E$ および $Z' = f^{-1}(Z)$ と表す。このとき、定義
\ref{chow-definition-localized-chern} の双変類
$$
P_p(Z' \to X', E') \in A^p(Z' \to X'),
\quad\text{resp.}\quad
c_p(Z' \to X', E') \in A^p(Z' \to X')
$$
で、$X', Z', E'$ を用いて補題 \ref{chow-lemma-independent-loc-chern} と同様に構成されるものは、
双変類 $P_p(Z \to X, E) \in A^p(Z \to X)$, resp.\ $c_p(Z \to X, E) \in A^p(Z \to X)$ の制限（注意
\ref{chow-remark-restriction-bivariant}）である。
\end{lemma}

\begin{proof}
構造射を $p : \mathbf{P}^1_X \to X$ および $p' : \mathbf{P}^1_{X'} \to X'$ と表す。
$b : W \to \mathbf{P}^1_X$ と $b' : W' \to \mathbf{P}^1_{X'}$ は、More on Flatness の補題
\ref{flat-lemma-complex-and-divisor-blowup-pre} において三つ組
$(\mathbf{P}^1_X, (\mathbf{P}^1_X)\infty, p^*E)$ と
$(\mathbf{P}^1_{X'}, (\mathbf{P}^1_{X'})\infty, (p')^*E')$ から構成される射であった。
さらに $Q = L\eta_{\mathcal{I}_\infty}p^*E$ および
$Q = L\eta_{\mathcal{I}'_\infty}(p')^*E'$ である。ここで
$\mathcal{I}_\infty \subset \mathcal{O}_W$ は $W_\infty$ のイデアル層、
$\mathcal{I}'_\infty \subset \mathcal{O}_{W'}$ は $W'_\infty$ のイデアル層である。
次に、$h : \mathbf{P}^1_{X'} \to \mathbf{P}^1_X$ は、スキームの射であって、有効 Cartier 因子
$(\mathbf{P}^1_X)_\infty$ の引き戻しが有効 Cartier 因子
$(\mathbf{P}^1_{X'})_\infty$ となり、かつ $h^*p^*E = (p')^*E'$ を満たす。
More on Flatness の補題
\ref{flat-lemma-complex-and-divisor-blowup-base-change} により、可換図式
$$
\xymatrix{
W' \ar[rd]_{b'} \ar[r]_-g &
\mathbf{P}^1_{X'} \times_{\mathbf{P}^1_X} W \ar[d]_r \ar[r]_-q &
W \ar[d]^b \\
&
\mathbf{P}^1_{X'} \ar[r] &
\mathbf{P}^1_X
}
$$
を得る。ここで $W'$ は $b$ に関する $\mathbf{P}^1_{X'}$ の「狭義変換」であり、
$Q' = (q \circ g)^*Q$ である。さて $P_p(Z \to X, E) = P'_p(Q)$,
resp.\ $c_p(Z \to X, E) = c'_p(Q)$ であった。ここで $P'_p(Q)$, resp.\ $c'_p(Q)$ は、
補題 \ref{chow-lemma-localized-chern-pre} において $b, Q, T'$ を用いて構成される。
$T'$ は次の二性質をもつ閉部分スキーム $T' \subset W_\infty$ である。
(a) $T'$ は $W_\infty$ の $X \setminus Z$ 上にあるすべての点を含む。
(b) $Q|_{T'}$ は零である、resp.\ 次数 $0$ に置かれた階数 $< p$ の有限局所自由加群と同型である。
補題 \ref{chow-lemma-localized-chern-pre} の構成では、性質 (a), (b) をもつ特定の閉部分スキーム
$T'$ を選んだが、$T'$ の正確な選択は重要でない。補題
\ref{chow-lemma-localized-chern-pre-independent} を見よ。

\medskip\noindent
次に補題 \ref{chow-lemma-base-change-localized-chern-pre} により、双変類
$P_p(Z \to X, E) = P'_p(Q)$, resp.\ $c_p(Z \to X, E) = c_p(Q')$ の $X'$ への制限は、
補題 \ref{chow-lemma-localized-chern-pre} と同様に、
$r : \mathbf{P}^1_{X'} \times_{\mathbf{P}^1_X} W \to \mathbf{P}^1_{X'}$、複体 $q^*Q$、および逆像
$q^{-1}(T')$ を用いて構成される類 $P'_p(q^*Q)$, resp.\ $c'_p(q^*Q)$ に対応する。

\medskip\noindent
さらに補題 \ref{chow-lemma-localized-chern-pre-independent} の第二の主張により
$P'_p(Q') = P'_p(q^*Q)$, resp.\ $c'_p(q^*Q) = c'_p(Q')$ である。
$P_p(Z' \to X', E') = P'_p(Q')$, resp.\ $c_p(Z' \to X', E') = c'_p(Q')$ なので、補題が従う。
\end{proof}

\begin{remark}
\label{chow-remark-loc-chern}
状況 \ref{chow-situation-loc-chern} では、仮定 (3) を「$E$ の Chern 類が定義される」という仮定で
置き換える方が自然であった。実際、補題 \ref{chow-lemma-independent-loc-chern} と
\ref{chow-lemma-base-change-loc-chern} を補題 \ref{chow-lemma-envelope-bivariant} と組み合わせれば、
このわずかに一般的な場合へ定義を容易に拡張できる。必要になれば、ここでそうすることにする。
\end{remark}

\begin{lemma}
\label{chow-lemma-loc-chern-after-pushforward}
状況 \ref{chow-situation-loc-chern} において、任意の $\alpha \in \CH_*(Z)$ に対し
$$
P_p(Z \to X, E) \cap i_*\alpha = P_p(E|_Z) \cap \alpha,
\quad\text{resp.}\quad
c_p(Z \to X, E) \cap i_*\alpha = c_p(E|_Z) \cap \alpha
$$
が $\CH_*(Z)$ で成り立つ。
\end{lemma}

\begin{proof}
第二の等式だけを証明し、第一の証明は省略する。$c_p(Z \to X, E)$ は双変類であり、
$Z \to X$ を $Z \to X$ で基底変換すると $\text{id} : Z \to Z$ になるので
$c_p(Z \to X, E) \cap i_*\alpha = c_p(Z \to X, E) \cap \alpha$ である。
補題 \ref{chow-lemma-base-change-loc-chern} により、$c_p(Z \to X, E)$ の $Z$ への制限（！）は、
$\text{id} : Z \to Z$ と $E|_Z$ に対する局所化 Chern 類である。したがって
$X = Z$ とした (\ref{chow-equation-defining-property-localized-classes}) から結論が従う。
\end{proof}

\begin{lemma}
\label{chow-lemma-loc-chern-disjoint}
状況 \ref{chow-situation-loc-chern} において、$\alpha \in \CH_k(X)$ の台が $Z$ と交わらないなら
$P_p(Z \to X, E) \cap \alpha = 0$, resp.\ $c_p(Z \to X, E) \cap \alpha = 0$ である。
\end{lemma}

\begin{proof}
局所化 Chern 類の構成から直ちに従う。別の見方として、まず
$c_p(Z \to X, E) \cap \alpha$ を計算する際に $c_p(Z \to X, E)$ を $\alpha$ の台へ制限し、
次に補題 \ref{chow-lemma-base-change-loc-chern} によりこの制限が零であることを用いてもよい。
\end{proof}

\begin{lemma}
\label{chow-lemma-loc-chern-shrink-Z}
状況 \ref{chow-situation-loc-chern} において、$Z \subset Z' \subset X$ とし、$Z'$ を $X$ の
閉部分スキームと仮定する。このとき
$P_p(Z' \to X, E) = (Z \to Z')_* \circ P_p(Z \to X, E)$,
resp.\ $c_p(Z' \to X, E) = (Z \to Z')_* \circ c_p(Z \to X, E)$
である（$\circ$ は補題 \ref{chow-lemma-push-proper-bivariant} のものとする）。
\end{lemma}

\begin{proof}
補題 \ref{chow-lemma-independent-loc-chern} における $P_p(Z' \to X, E)$,
resp.\ $c_p(Z' \to X, E)$ の構成は、まったく同じ射
$b : W \to \mathbf{P}^1_X$ と $D(\mathcal{O}_W)$ の完全対象 $Q$ を用いる。
したがって補題 \ref{chow-lemma-silly-shrink} を用いて結論を得る。細部の一部は省略する。
\end{proof}

\begin{lemma}
\label{chow-lemma-loc-chern-agree}
補題 \ref{chow-lemma-silly} において、$E_2$ は完全対象 $E \in D(\mathcal{O}_X)$ の制限であり、
その $X_1$ への制限は零である、resp.\ コホモロジー次数 $0$ に置かれた階数 $< p$ の
有限局所自由 $\mathcal{O}_{X_1}$-加群と同型であるとする。このとき、補題
\ref{chow-lemma-silly} の類 $P'_p(E_2)$, resp.\ $c'_p(E_2)$ は、定義
\ref{chow-definition-localized-chern} の $P_p(X_2 \to X, E)$, resp.\ $c_p(X_2 \to X, E)$ と、
$E$ が状況 \ref{chow-situation-loc-chern} の仮定 (3) を満たすなら一致する。
\end{lemma}

\begin{proof}
$E$ に関する仮定から、$X_1$ を含む開部分スキーム $U \subset X$ が存在し、$E|_U$ は零である、
resp.\ 階数 $< p$ の有限局所自由 $\mathcal{O}_U$-加群と同型である。More on Algebra の補題
\ref{more-algebra-lemma-lift-perfect-from-residue-field} を見よ。
$Z \subset X$ を、$X$ における $U$ の補集合に被約誘導閉部分スキーム構造を入れたものとする。
補題 \ref{chow-lemma-loc-chern-shrink-Z} により
$P_p(X_2 \to X, E) = (Z \to X_2)_* \circ P_p(Z \to X, E)$,
resp.\ $c_p(X_2 \to X, E) = (Z \to X_2)_* \circ c_p(Z \to X, E)$ である。
そこで $P_p(X_2 \to X, E)$, resp.\ $c_p(X_2 \to X, E)$ が、補題
\ref{chow-lemma-silly} で与えた $P'_p(E_2)$, resp.\ $c'_p(E_2)$ の特徴づけを満たすことを示す。
実際、いま示した関係
$P_p(X_2 \to X, E) = (Z \to X_2)_* \circ P_p(Z \to X, E)$,
resp.\ $c_p(X_2 \to X, E) = (Z \to X_2)_* \circ c_p(Z \to X, E)$ と
$X_1 \cap Z = \emptyset$ から、合成 $P_p(X_2 \to X, E) \circ i_{1, *}$,
resp.\ $c_p(X_2 \to X, E) \circ i_{1, *}$ は補題
\ref{chow-lemma-loc-chern-disjoint} により零である。一方
$P_p(X_2 \to X, E) \circ i_{2, *} = P_p(E_2)$,
resp.\ $c_p(X_2 \to X, E) \circ i_{2, *} = c_p(E_2)$ が補題
\ref{chow-lemma-loc-chern-after-pushforward} により成り立つ。
\end{proof}





\section{二つの技術的補題}
\label{chow-section-tools-loc-chern}

\noindent
この節では、局所化 Chern 類をより扱いやすくするための補助的な道具を整備する。
次の補題は、補題 \ref{chow-lemma-localized-chern-pre-compose} と
\ref{chow-lemma-localized-chern-pre-sum-c} の証明ですでに現れた事実を、より精密に述べたものである。

\begin{lemma}
\label{chow-lemma-homomorphism-final}
$(S, \delta)$ を状況 \ref{chow-situation-setup} のとおりとし、$X$ を $S$ 上局所有限型とする。
$b : W \longrightarrow \mathbf{P}^1_X$ をスキームの固有射とし、$n \geq 1$ とする。
$i = 1, \ldots, n$ に対して、$Z_i \subset X$ を閉部分スキーム、
$Q_i \in D(\mathcal{O}_W)$ を完全対象、$p_i \geq 0$ を整数とし、
$T_i \subset W_\infty$, $i = 1, \ldots, n$ を閉とする。
$W_i = b^{-1}(\mathbf{P}^1_{Z_i})$ と表す。次を仮定する。
\begin{enumerate}
\item $i = 1, \ldots, n$ に対して、$b, Z_i, Q_i, T_i, p_i$ は補題
\ref{chow-lemma-localized-chern-pre} の仮定を満たす。
\item $Q_i|_{W \setminus W_i}$ は零である、resp.\ コホモロジー次数 $0$ に置かれた
階数 $< p_i$ の有限局所自由加群と同型である。
\item $W$ 上の $Q_i$ は状況 \ref{chow-situation-loc-chern} の仮定 (3) を満たす。
\end{enumerate}
このとき $P'_{p_n}(Q_n) \circ \ldots \circ P'_{p_1}(Q_1)$ は
$$
(W_{n, \infty} \cap \ldots \cap W_{1, \infty} \to
Z_n \cap \ldots \cap Z_1)_* \circ
P'_{p_n}(Q_n|_{W_{n, \infty}}) \circ \ldots \circ P'_{p_1}(Q_1|_{W_{1, \infty}})
\circ C
$$
に等しく、この等式は $A^{p_n + \ldots + p_1}(Z_n \cap \ldots \cap Z_1 \to X)$ で成り立つ。
resp.\ $c'_{p_n}(Q_n) \circ \ldots \circ c'_{p_1}(Q_1)$ は
$$
(W_{n, \infty} \cap \ldots \cap W_{1, \infty} \to
Z_n \cap \ldots \cap Z_1)_* \circ
c'_{p_n}(Q_n|_{W_{n, \infty}}) \circ \ldots \circ c'_{p_1}(Q_1|_{W_{1, \infty}})
\circ C
$$
に等しく、この等式は $A^{p_n + \ldots + p_1}(Z_n \cap \ldots \cap Z_1 \to X)$ で成り立つ。
\end{lemma}

\begin{proof}
Chern 類に関する主張を $n$ に関する帰納法で証明する。$P_p(-)$ に関する主張もまったく同様に証明される。
$n = 1$ の場合は $c'_{p_1}(Q_1)$ の構成そのものである。実際、$W_{1, \infty}$ は
$W_\infty$ における $Z_1$ の逆像である。$n > 1$ とする。帰納法により
$c'_{p_n}(Q_n) \circ \ldots \circ c'_{p_1}(Q_1)$ は
$$
c'_{p_n}(Q_n) \circ
(W_{n - 1, \infty} \cap \ldots \cap W_{1, \infty} \to
Z_{n - 1} \cap \ldots \cap Z_1)_* \circ
c'_{p_{n - 1}}(Q_{n - 1}|_{W_{n - 1}, \infty}) \circ \ldots \circ
c'_{p_1}(Q_1|_{W_{1, \infty}})
\circ C
$$
に等しい。補題 \ref{chow-lemma-base-change-localized-chern-pre} により、
$c'_{p_n}(Q_n)$ の $Z_{n - 1} \cap \ldots \cap Z_1$ への制限は、閉部分集合
$Z_n \cap \ldots \cap Z_1$、射
$b' : W_{n - 1} \cap \ldots \cap W_1 \to
\mathbf{P}^1_{Z_{n - 1} \cap \ldots \cap Z_1}$、および $Q_n$ の
$W_{n - 1} \cap \ldots \cap W_1$ への制限によって計算される。
ここで $(b')^{-1}(Z_n) = W_n \cap \ldots \cap W_1$ かつ
$(W_n \cap \ldots \cap W_1)_\infty =
W_{n, \infty} \cap \ldots \cap W_{1, \infty}$ である。
補題 \ref{chow-lemma-gysin-at-infty} の類を
$C_{n - 1} \in A^0(W_{n - 1, \infty} \cap \ldots \cap W_{1, \infty} \to
Z_{n - 1} \cap \ldots \cap Z_1)$ と表す。したがって
$c'_{p_n}(Q_n)$ の $Z_{n - 1} \cap \ldots \cap Z_1$ への制限は
\begin{align*}
&
(W_{n, \infty} \cap \ldots \cap W_{1, \infty} \to Z_n \cap \ldots \cap Z_1)_*
\circ
c'_{p_n}(Q_n|_{(W_n \cap \ldots \cap W_1)_\infty})
\circ
C_{n - 1} \\
& =
(W_{n, \infty} \cap \ldots \cap W_{1, \infty} \to Z_n \cap \ldots \cap Z_1)_*
\circ
c'_{p_n}(Q_n|_{W_{n, \infty}})
\circ
C_{n - 1}
\end{align*}
である。等式は補題 \ref{chow-lemma-base-change-silly} から従う（右辺では制限を記法から省略した）。
したがって上の式は
\begin{align*}
(W_{n, \infty} \cap \ldots \cap W_{1, \infty} \to
Z_n \cap \ldots \cap Z_1)_* \circ
c'_{p_n}(Q_n|_{W_n, \infty}) \circ \\
C_{n - 1} \circ
(W_{n - 1, \infty} \cap \ldots \cap W_{1, \infty} \to
Z_{n - 1} \cap \ldots \cap Z_1)_* \\
\circ
c'_{p_{n - 1}}(Q_{n - 1}|_{W_{n - 1}, \infty}) \circ \ldots \circ
c'_{p_1}(Q_1|_{W_{1, \infty}})
\circ C
\end{align*}
となる。補題 \ref{chow-lemma-homomorphism-pre} により、合成
$C_{n - 1} \circ (W_{n - 1, \infty} \cap \ldots \cap W_{1, \infty} \to
Z_{n - 1} \cap \ldots \cap Z_1)_*$ は Gysin 写像
$$
(W_{n - 1, \infty} \cap \ldots \cap W_{1, \infty} \to
W_{n - 1} \cap \ldots \cap W_1)^*
$$
の像の元上で恒等写像である。したがって
$c'_{p_{n - 1}}(Q_{n - 1}|_{W_{n - 1}, \infty}) \circ \ldots \circ
c'_{p_1}(Q_1|_{W_{1, \infty}}) \circ C$ の像の任意の元が Gysin 写像の像に属することを
示せば十分である。補題 \ref{chow-lemma-loc-chern-agree} と、補題の主張における $Q_i$ に関する
仮定 (2), (3) により
$$
c'_{p_i}(Q_i|_{W_{i, \infty}}) = \text{restriction of } c_{p_i}(W_i \to W, Q_i)
\text{ to } W_{i, \infty}
$$
と書ける。したがって、$\beta \in \CH_{k + 1}(W)$ が $b^{-1}(\mathbf{A}^1_X)$ 上で
$\alpha$ の平坦引き戻しに制限されるなら
\begin{align*}
& c'_{p_{n - 1}}(Q_{n - 1}|_{W_{n - 1}, \infty}) \cap \ldots \cap
c'_{p_1}(Q_1|_{W_{1, \infty}})
\cap C \cap \alpha \\
& =
c'_{p_{n - 1}}(Q_{n - 1}|_{W_{n - 1}, \infty}) \cap \ldots \cap
c'_{p_1}(Q_1|_{W_{1, \infty}})
\cap i_\infty^* \beta \\
& =
c_{p_{n - 1}}(W_{n - 1} \to W, Q_{n - 1}) \cap \ldots \cap
c_{p_{n - 1}}(W_1 \to W, Q_1) \cap i_\infty^* \beta \\
& =
(W_{n - 1, \infty} \cap \ldots \cap W_{1, \infty} \to
W_{n - 1} \cap \ldots \cap W_1)^*
\left(c_{p_{n - 1}}(W_{n - 1} \to W, Q_{n - 1}) \cap \ldots \cap
c_{p_1}(W_1 \to W, Q_1) \cap \beta\right)
\end{align*}
となり、求める結論を得る。最後の等式では、$c_{p_i}(W_i \to W, Q_i)$ が双変類であり、
したがって定義により $i_\infty^*$ と可換であることを用いた。
\end{proof}

\noindent
次の補題は、複体の局所化 Chern 類を計算する際に極めて大きな自由度を与える。

\begin{lemma}
\label{chow-lemma-independent-loc-chern-bQ}
$(S, \delta), X, Z, b : W \to \mathbf{P}^1_X, Q, T, p$ が
補題 \ref{chow-lemma-localized-chern-pre} の仮定を満たすとする。
$F \in D(\mathcal{O}_X)$ を、次の条件を満たす完全対象とする。
\begin{enumerate}
\item $Q$ の $b^{-1}(\mathbf{A}^1_X)$ への制限は $F$ の引き戻しと同型である。
\item $F|_{X \setminus Z}$ は零である、resp.\ コホモロジー次数 $0$ に置かれた
階数 $< p$ の有限局所自由 $\mathcal{O}_{X \setminus Z}$-加群と同型である。
\item $W$ 上の $Q$ と $X$ 上の $F$ は状況 \ref{chow-situation-loc-chern} の仮定 (3) を満たす。
\end{enumerate}
このとき、補題 \ref{chow-lemma-localized-chern-pre} で構成した
$A^p(Z \to X)$ の類 $P'_p(Q)$, resp.\ $c'_p(Q)$ は、定義
\ref{chow-definition-localized-chern} の $P_p(Z \to X, F)$, resp.\ $c_p(Z \to X, F)$ に等しい。
\end{lemma}

\begin{proof}
仮定は局所有限型な射 $X' \to X$ による基底変換で保たれる。したがって、任意の
$\alpha \in \CH_k(X)$ に対して
$P_p(Z \to X, F) \cap \alpha = P'_p(Q) \cap \alpha$,
resp.\ $c_p(Z \to X, F) \cap \alpha = c'_p(Q) \cap \alpha$
を示せば十分である。補題 \ref{chow-lemma-gysin-at-infty} における $C$ の構成と同様に、
$b^{-1}(\mathbf{A}^1_X)$ への制限が $\alpha$ の平坦引き戻しに等しい
$\beta \in \CH_{k + 1}(W)$ を選ぶ。$W' = b^{-1}(Z)$ とおき、
$W_\infty \to X$ による $Z$ の逆像を
$E = W'_\infty \subset W_\infty$ と表す。
補題は次の一連の等式から従う（$P_p$ の場合も同様である）。
\begin{align*}
c'_p(Q) \cap \alpha
& =
(E \to Z)_*(c'_p(Q|_E) \cap i_\infty^*\beta) \\
& =
(E \to Z)_*(c_p(E \to W_\infty, Q|_{W_\infty}) \cap i_\infty^*\beta) \\
& =
(W'_\infty \to Z)_*(c_p(W' \to W, Q) \cap i_\infty^*\beta) \\
& =
(W'_\infty \to Z)_*((i'_\infty)^*(c_p(W' \to W, Q) \cap \beta)) \\
& =
(W'_\infty \to Z)_*((i'_\infty)^*(c_p(Z' \to X, F) \cap \beta)) \\
& =
(W'_0 \to Z)_*((i'_0)^*(c_p(Z' \to X, F) \cap \beta)) \\
& =
(W'_0 \to Z)_*(c_p(Z' \to X, F) \cap i_0^*\beta)) \\
& =
c_p(Z \to X, F) \cap \alpha
\end{align*}
第一の等式は補題 \ref{chow-lemma-localized-chern-pre} における $c'_p(Q)$ の構成である。
第二の等式は補題 \ref{chow-lemma-loc-chern-agree} である。$W' \to W$ を
$W_\infty \to W$ で基底変換すると射 $E = W'_\infty \to W_\infty$ になるので、
第三の等式は補題 \ref{chow-lemma-base-change-loc-chern} から従う。
第四の等式では $i'_\infty : W'_\infty \to W'$ は包含射であり、この等式は
$c_p(W' \to W, Q)$ が双変類であることから従う。第五の等式については、
$c_p(W' \to W, Q)$ と $c_p(Z' \to X, F)$ が
$A^p((b')^{-1} \to b^{-1}(\mathbf{A}^1_X))$ において同じ双変類に制限されることに注意する。
実際、補題の仮定 (1) により $Q$ と $F$ は
$D(\mathcal{O}_{b^{-1}(\mathbf{A}^1_X)})$ の同じ対象に制限される。ここで補題
\ref{chow-lemma-base-change-loc-chern} を用いる。$(i'_\infty)^*$ は $W'_\infty$ 上に台をもつ
サイクルを零にするので（注意 \ref{chow-remark-gysin-on-cycles} を参照）、第五の等式が従う。
第六の等式は、$W'_\infty$ と $W'_0$ が $\mathbf{P}^1_Z$ における有理同値な有効 Cartier 因子
$D_0, D_\infty$ の引き戻しであり、したがって $i_\infty^*\beta$ と $i_0^*\beta$ が
$W'$ 上の同じサイクル類へ写ることから従う。すなわち、いずれも補題
\ref{chow-lemma-support-cap-effective-Cartier} により類
$c_1(\mathcal{O}_{\mathbf{P}^1_Z}(1)) \cap c_p(Z \to X, F_) \cap \beta$ を表す。
第七の等式は $c_p(Z \to X, F)$ が双変類であることから従う。構成により
$W'_0 = Z$ かつ $i_0^*\beta = \alpha$ であり、これが最後の等式を説明する。
\end{proof}







\section{局所化 Chern 類の性質}
\label{chow-section-properties-loc-chern}

\noindent
この節の主結果は、局所化 Chern 類の加法性と乗法性である。

\begin{lemma}
\label{chow-lemma-loc-chern-character}
状況 \ref{chow-situation-loc-chern} において $E|_{X \setminus Z}$ が零であると仮定する。このとき
\begin{align*}
P_1(Z \to X, E) & = c_1(Z \to X, E), \\
P_2(Z \to X, E) & = c_1(Z \to X, E)^2 - 2c_2(Z \to X, E), \\
P_3(Z \to X, E) & = c_1(Z \to X, E)^3 - 3c_1(Z \to X, E)c_2(Z \to X, E)
+ 3c_3(Z \to X, E),
\end{align*}
などが成り立つ。ここで積は注意 \ref{chow-remark-ring-loc-classes} の代数
$A^{(1)}(Z \to X)$ において取る。
\end{lemma}

\begin{proof}
零層の階数は $< 1$ なので、類 $c_p(Z \to X, E)$ はすべての $p \geq 1$ に対して
定義されており、主張には意味がある。結論は、より一般的な補題
\ref{chow-lemma-localized-chern-pre-compose} から直ちに従う。実際、局所化 Chern 類は
補題 \ref{chow-lemma-localized-chern-pre} の手順を用いて節
\ref{chow-section-localized-chern} で定義された。
\end{proof}

\begin{lemma}
\label{chow-lemma-loc-chern-classes-commute}
状況 \ref{chow-situation-loc-chern} において、$Y \to X$ を局所有限型とし、
$c \in A^*(Y \to X)$ とする。このとき
$$
P_p(Z \to X, E) \circ c = c \circ P_p(Z \to X, E),
$$
および
$$
c_p(Z \to X, E) \circ c = c \circ c_p(Z \to X, E)
$$
が $A^*(Y \times_X Z \to X)$ において成り立つ。
\end{lemma}

\begin{proof}
これは補題 \ref{chow-lemma-homomorphism-commute} から従う。より正確には、
$$
b : W \to \mathbf{P}^1_X
\quad\text{and}\quad
Q
\quad\text{and}\quad
T' \subset T \subset W_\infty
$$
を補題 \ref{chow-lemma-independent-loc-chern} の証明のとおりとする。定義により、
双変演算として $c_p(Z \to X, E) = c'_p(Q)$ である。ここで右辺は
$W, b, Q, T'$ を用いて補題 \ref{chow-lemma-localized-chern-pre} で構成した双変類である。
補題 \ref{chow-lemma-homomorphism-commute} により
$P'_p(Q) \circ c = c \circ P'_p(Q)$, resp.\ $c'_p(Q) \circ c = c \circ c'_p(Q)$
が $A^*(Y \times_X Z \to X)$ において成り立つので、結論を得る。
\end{proof}

\begin{remark}
\label{chow-remark-loc-chern-classes}
状況 \ref{chow-situation-loc-chern} においては
$$
c^{(p)}(Z \to X, E) = 1 + c_1(E) + \ldots + c_{p - 1}(E) +
c_p(Z \to X, E) + c_{p + 1}(Z \to X, E) + \ldots
$$
を、注意 \ref{chow-remark-ring-loc-classes} で考えた代数 $A^{(p)}(Z \to X)$ の元として
定義しておくと便利である。
\end{remark}

\begin{lemma}
\label{chow-lemma-additivity-loc-chern-c}
$(S, \delta)$ を状況 \ref{chow-situation-setup} のとおりとする。
$X$ を $S$ 上局所有限型とし、$Z \to X$ を閉埋め込みとする。
$D(\mathcal{O}_X)$ の完全対象の識別三角
$$
E_1 \to E_2 \to E_3 \to E_1[1]
$$
を考える。次を仮定する。
\begin{enumerate}
\item 制限 $E_1|_{X \setminus Z}$ と $E_3|_{X \setminus Z}$ は、次数 $0$ に置かれた、
それぞれ階数 $< p_1$ と $< p_3$ の有限局所自由
$\mathcal{O}_{X \setminus Z}$-加群と同型である。
\item 次の少なくとも一つが成り立つ。
(a) $X$ は準コンパクトである。
(b) $X$ の既約成分は準コンパクトである。
(c) $E_3 \to E_1[1]$ は、有限局所自由 $\mathcal{O}_X$-加群からなる
局所有界複体の写像で表現できる。
(d) 包絡 $f : Y \to X$ が存在し、$Lf^*E_3 \to Lf^*E_1[1]$ は
有限局所自由 $\mathcal{O}_Y$-加群からなる局所有界複体の写像で表現できる。
\end{enumerate}
注意 \ref{chow-remark-loc-chern-classes} の記法を用いると
$$
c^{(p_1 + p_3)}(Z \to X, E_2) = c^{(p_1)}(Z \to X, E_1)c^{(p_3)}(Z \to X, E_3)
$$
が $A^{(p_1 + p_3)}(Z \to X)$ において成り立つ。
\end{lemma}

\begin{proof}
仮定から、$E_2|_{X \setminus Z}$ は零である、resp.\ 次数零に置かれた階数
$< p_1 + p_3$ の有限局所自由 $\mathcal{O}_{X \setminus Z}$-加群と同型である。
したがって主張には意味がある。

\medskip\noindent
$f : Y \to X$ を包絡とする。補題の公式の両辺を展開すると、$A^*(X)$ と
$A^*(Z \to X)$ における類のいくつかの等式を証明すべきことが分かる。
補題 \ref{chow-lemma-envelope-bivariant} の一意性により、対応する関係を
$A^*(Y)$ と $A^*(Z \to Y)$ において証明すれば十分である。さらに、関係する類の構成は
基底変換と両立するので（補題 \ref{chow-lemma-base-change-loc-chern}）、$X$ を $Y$ で、
識別三角をその引き戻しで置き換えてよい。

\medskip\noindent
補題 \ref{chow-lemma-additivity-on-perfect} の証明において、条件 (2)(a), (2)(b), (2)(c) が
条件 (2)(d) を含意することを見た。これと前段落の議論を合わせると、
次の段落で扱う場合に帰着する。

\medskip\noindent
$\varphi^\bullet : \mathcal{E}_3^\bullet[-1] \to \mathcal{E}_1^\bullet$ を、導来圏の写像
$E_3[-1] \to E_1$ を表す、有限局所自由 $\mathcal{O}_X$-加群からなる局所有界複体の写像とする。
射影 $g : X' \to X$ を備えたスキーム $X' = \mathbf{A}^1 \times X$ を考え、
$Z' = g^{-1}(Z) = \mathbf{A}^1 \times Z$ とおく。$\mathbf{A}^1$ の座標を $t$ と表す。
$X'$ 上の複体の写像
$$
t g^*\varphi^\bullet :
g^*\mathcal{E}_3^\bullet[-1]
\longrightarrow
g^*\mathcal{E}_1^\bullet
$$
の錐 $\mathcal{C}^\bullet$ を考える。識別三角
$$
g^*\mathcal{E}_1^\bullet \to \mathcal{C}^\bullet \to
g^*\mathcal{E}_3^\bullet \to g^*\mathcal{E}_1^\bullet[1]
$$
を得る。ここで最初の三項は、各次数で分裂する複体の短完全列をなす。
明らかに $\mathcal{C}^\bullet$ は有限局所自由 $\mathcal{O}_{X'}$-加群からなる有界複体であり、
その $X' \setminus Z'$ への制限は、次数 $0$ に置かれた階数 $< p_1 + p_3$ の
有限局所自由 $\mathcal{O}_{X' \setminus Z'}$-加群と同型である。したがって $p \geq p_1 + p_3$ に対して
局所化 Chern 類
$$
c_p(Z' \to X', \mathcal{C}^\bullet) \in A^p(Z' \to X')
$$
をもつ。任意の $\alpha \in \CH_k(X)$ に対して
$$
c_p(Z' \to X', \mathcal{C}^\bullet) \cap g^*\alpha
\in \CH_{k + 1 - p}(\mathbf{A}^1 \times X)
$$
を考える。$t = 0$ に制限すると、写像 $t g^*\varphi^\bullet$ は零に制限され、
$\mathcal{C}^\bullet|_{t = 0}$ は $\mathcal{E}_1^\bullet$ と $\mathcal{E}_3^\bullet$ の直和になる。
局所化 Chern 類の基底変換との両立性（補題 \ref{chow-lemma-base-change-loc-chern}）により
$$
i_0^* \circ c^{(p_1 + p_3)}(Z' \to X', \mathcal{C}^\bullet) \circ g^* =
c^{(p_1 + p_2)}(Z \to X, E_1 \oplus E_3)
$$
が $A^{(p_1 + p_3)}(Z \to X)$ において成り立つ。一方、$t = 1$ に制限すると、
写像 $t g^*\varphi^\bullet$ は $\varphi$ に制限され、
$\mathcal{C}^\bullet|_{t = 1}$ は $E_2$ を表す有限局所自由加群からなる有界複体になる。
したがって
$$
i_1^* \circ c^{(p_1 + p_3)}(Z' \to X', \mathcal{C}^\bullet) \circ g^* =
c^{(p_1 + p_2)}(Z \to X, E_2)
$$
が $A^{(p_1 + p_3)}(Z \to X)$ において成り立つ。有理同値の定義により
$i_0^* = i_1^*$ である（より正確には補題 \ref{chow-lemma-linebundle-formulae} の公式から従う）ので、
$$
c^{(p_1 + p_2)}(Z \to X, E_2) = c^{(p_1 + p_2)}(Z \to X, E_1 \oplus E_3)
$$
を得る。これで次の段落で扱う場合に帰着した。

\medskip\noindent
$E_2 = E_1 \oplus E_3$ であり、三つ組 $(X, Z, E_i)$ が状況
\ref{chow-situation-loc-chern} のとおりであると仮定する。$i = 1, 3$ に対して
$$
b_i : W_i \to \mathbf{P}^1_X
\quad\text{and}\quad
Q_i
\quad\text{and}\quad
T'_i \subset T_i \subset W_{i, \infty}
$$
を補題 \ref{chow-lemma-independent-loc-chern} の証明のとおりとする。定義により
$$
c_p(Z \to X, E_i) = c'_p(Q_i)
$$
である。右辺は $W_i, b_i, Q_i, T'_i$ を用いて補題
\ref{chow-lemma-localized-chern-pre} で構成した双変類である。
$W = W_1 \times_{b_1, \mathbf{P}^1_X, b_2} W_2$ とおき、Cartesian 図式
$$
\xymatrix{
W \ar[d]_{g_1} \ar[rd]^b \ar[r]_{g_3} & W_3 \ar[d]^{b_3} \\
W_1 \ar[r]^{b_1} & \mathbf{P}^1_X
}
$$
を考える。もちろん $b^{-1}(\mathbf{A}^1)$ は $\mathbf{A}^1_X$ へ同型に写る。
$T' = g_1^{-1}(T'_1) \cap g_2^{-1}(T'_2)$ は依然として、$X \setminus Z$ 上にある
$W_\infty$ のすべての点を含む。補題 \ref{chow-lemma-localized-chern-pre-independent} により、
$i = 1, 3$ に対する $c_p(Z \to X, E_i)$ の構成に
$W$, $b$, $g_i^*\mathcal{Q}_i$, $T'$ を用いてよい。さらに補題
\ref{chow-lemma-independent-loc-chern-bQ} のより強い独立性により、類
$c_p(Z \to X, E_2)$ の計算に $W$, $b$, $g_1^*Q_1 \oplus g_3^*Q_3$, $T'$ を用いてよい。
したがって、求める等式は補題 \ref{chow-lemma-localized-chern-pre-sum-c} から従う。
\end{proof}

\begin{lemma}
\label{chow-lemma-additivity-loc-chern-P}
$(S, \delta)$ を状況 \ref{chow-situation-setup} のとおりとする。
$X$ を $S$ 上局所有限型とし、$Z \to X$ を閉埋め込みとする。
$D(\mathcal{O}_X)$ の完全対象の識別三角
$$
E_1 \to E_2 \to E_3 \to E_1[1]
$$
を考える。次を仮定する。
\begin{enumerate}
\item 制限 $E_1|_{X \setminus Z}$ と $E_3|_{X \setminus Z}$ は零である。
\item 次の少なくとも一つが成り立つ。
(a) $X$ は準コンパクトである。
(b) $X$ の既約成分は準コンパクトである。
(c) $E_3 \to E_1[1]$ は、有限局所自由 $\mathcal{O}_X$-加群からなる
局所有界複体の写像で表現できる。
(d) 包絡 $f : Y \to X$ が存在し、$Lf^*E_3 \to Lf^*E_1[1]$ は
有限局所自由 $\mathcal{O}_Y$-加群からなる局所有界複体の写像で表現できる。
\end{enumerate}
このとき
$$
P_p(Z \to X, E_2) = P_p(Z \to X, E_1) + P_p(Z \to X, E_3)
$$
がすべての $p \in \mathbf{Z}$ に対して成り立ち、したがって
$ch(Z \to X, E_2) = ch(Z \to X, E_1) + ch(Z \to X, E_3)$ である。
\end{lemma}

\begin{proof}
証明は補題 \ref{chow-lemma-additivity-loc-chern-c} の証明とまったく同じであるが、
最後に補題 \ref{chow-lemma-localized-chern-pre-sum-P} を用いる点だけが異なる。
$p > 0$ に対しては、$p_1 = p_3 = 1$ とした補題
\ref{chow-lemma-additivity-loc-chern-c} と、補題 \ref{chow-lemma-loc-chern-character} における
$P_p(Z \to X, E)$ と $c_p(Z \to X, E)$ の関係からもこの補題を導ける。
$p = 0$ の場合は直接示せる（$X$ が $Z$ に完全に含まれる連結成分をもつ場合にのみ問題になる）。
\end{proof}

\begin{lemma}
\label{chow-lemma-loc-chern-tensor-product}
状況 \ref{chow-situation-setup} において、$X$ を $S$ 上局所有限型とする。
$Z_i \subset X$, $i = 1, 2$ を閉部分スキームとし、$F_i$, $i = 1, 2$ を
$D(\mathcal{O}_X)$ の完全対象とする。$i = 1, 2$ に対して
$F_i|_{X \setminus Z_i}$ が零であり\footnote{おそらく、この補題には
$F_i|_{X \setminus Z_i}$ が階数 $< p_i$ の有限局所自由
$\mathcal{O}_{X \setminus Z_i}$-加群と同型であることだけを仮定する変種がある。}、
$X$ 上の $F_i$ が状況 \ref{chow-situation-loc-chern} の仮定 (3) を満たすと仮定する。
$r_i = P_0(Z_i \to X, F_i) \in A^0(Z_i \to X)$ と表す。このとき
$$
c_1(Z_1 \cap Z_2 \to X, F_1 \otimes_{\mathcal{O}_X}^\mathbf{L} F_2) =
r_1 c_1(Z_2 \to X, F_2) + r_2 c_1(Z_1 \to X, F_1)
$$
が $A^1(Z_1 \cap Z_2 \to X)$ において成り立ち、
\begin{align*}
c_2(Z_1 \cap Z_2 \to X, F_1 \otimes_{\mathcal{O}_X}^\mathbf{L} F_2)
& =
r_1 c_2(Z_2 \to X, F_2) +
r_2 c_2(Z_1 \to X, F_1) + \\
& {r_1 \choose 2} c_1(Z_2 \to X, F_2)^2 + \\
& (r_1r_2 - 1) c_1(Z_2 \to X, F_2)c_1(Z_1 \to X, F_1) + \\
& {r_2 \choose 2} c_1(Z_1 \to X, F_1)^2
\end{align*}
が $A^2(Z_1 \cap Z_2 \to X)$ において成り立つ。高次の Chern 類についても同様である。
同様に
$$
ch(Z_1 \cap Z_2 \to X, F_1 \otimes_{\mathcal{O}_X}^\mathbf{L} F_2) =
ch(Z_1 \to X, F_1) ch(Z_2 \to X, F_2)
$$
が $\prod_{p \geq 0} A^p(Z_1 \cap Z_2 \to X) \otimes \mathbf{Q}$ において成り立つ。
より正確には
$$
P_p(Z_1 \cap Z_2 \to X, F_1 \otimes_{\mathcal{O}_X}^\mathbf{L} F_2) =
\sum\nolimits_{p_1 + p_2 = p}
{p \choose p_1} P_{p_1}(Z_1 \to X, F_1) P_{p_2}(Z_2 \to X, F_2)
$$
が $A^p(Z_1 \cap Z_2 \to X)$ において成り立つ。
\end{lemma}

\begin{proof}
$F_i$ の局所化 Chern 類の構成におけるとおり、より一般には補題
\ref{chow-lemma-independent-loc-chern-bQ} のとおり、固有射
$b_i : W_i \to \mathbf{P}^1_X$、$Q_i \in D(\mathcal{O}_{W_i})$、および閉部分スキーム
$T_i \subset W_{i, \infty}$ を選ぶ。可換図式
$$
\xymatrix{
W \ar[d]^{g_1} \ar[rd]^b \ar[r]_{g_2} & W_2 \ar[d]^{b_2} \\
W_1 \ar[r]^{b_1} & \mathbf{P}^1_X
}
$$
を選び、すべての射が固有であり、$\mathbf{A}^1_X$ 上で同型であるようにする。
例えば $W$ として、$b_1^{-1}(\mathbf{A}^1_X)$ と $b_2^{-1}(\mathbf{A}^1_X)$ の間の
同型のグラフの閉包を取れる。補題 \ref{chow-lemma-independent-loc-chern-bQ} により、
局所化 Chern 類 $c_p(Z_i \to X, F_i)$ の構成には
$W$, $b = b_i \circ g_i$, $Lg_i^*Q_i$, $g_i^{-1}(T_i)$ を用いてよい。
したがって次の段落で述べる状況に帰着する。

\medskip\noindent
次をもつと仮定する。
\begin{enumerate}
\item $\mathbf{A}^1_X$ 上で同型である固有射 $b : W \to \mathbf{P}^1_X$。
\item $E_i \subset W_\infty$ は $Z_i$ の逆像である。
\item Chern 類が定義されている完全対象 $Q_i \in D(\mathcal{O}_W)$ であって、次を満たすもの。
\begin{enumerate}
\item $Q_i$ の $b^{-1}(\mathbf{A}^1_X)$ への制限は $F_i$ の引き戻しである。
\item $X \setminus Z_i$ 上にある $W_\infty$ のすべての点を含む閉部分スキーム
$T_i \subset W_\infty$ が存在し、$Q_i|_{T_i}$ は零である。
\end{enumerate}
\end{enumerate}
補題 \ref{chow-lemma-independent-loc-chern-bQ} により
$$
c_p(Z_i \to X, F_i) = c'_p(Q_i) =
(E_i \to Z_i)_* \circ c'_p(Q_i|_{E_i}) \circ C
$$
および
$$
P_p(Z_i \to X, F_i) = P'_p(Q_i) =
(E_i \to Z_i)_* \circ P'_p(Q_i|_{E_i}) \circ C
$$
が $i = 1, 2$ に対して成り立つ。次に
$Q = Q_1 \otimes_{\mathcal{O}_W}^\mathbf{L} Q_2$ は
$F_1 \otimes_{\mathcal{O}_X}^\mathbf{L} F_2$ と $T_1 \cup T_2$ に対して
(3)(a) と (3)(b) を満たす。したがって
$$
c_p(Z_1 \cap Z_2 \to X, F_1 \otimes_{\mathcal{O}_X}^\mathbf{L} F_2) =
(E_1 \cap E_2 \to Z_1 \cap Z_2)_* \circ
c'_p(Q|_{E_1 \cap E_2}) \circ C
$$
および
$$
P_p(Z_1 \cap Z_2 \to X, F_1 \otimes_{\mathcal{O}_X}^\mathbf{L} F_2) =
(E_1 \cap E_2 \to Z_1 \cap Z_2)_* \circ
P'_p(Q|_{E_1 \cap E_2}) \circ C
$$
が同じ補題により成り立つ。補題 \ref{chow-lemma-silly-tensor-product} により、類
$c'_p(Q|_{E_1 \cap E_2})$ と $P'_p(Q|_{E_1 \cap E_2})$ は、類
$c'_p(Q_i|_{E_i})$ と $P'_p(Q_i|_{E_i})$ を用いて正しい形に展開できる。
最後に補題 \ref{chow-lemma-homomorphism-final} により、$c'_p(Q_i|_{E_i})$ と
$P'_p(Q_i|_{E_i})$ の多項式は、求めるとおり $c'_p(Q_i)$ と $P'_p(Q_i)$ の
対応する多項式に一致する。
\end{proof}






\section{無限遠におけるブローアップ}
\label{chow-section-blowup-Z-first}

\noindent
$X$ をスキームとし、$Z \subset X$ を有限型準連接イデアル層で切り出される閉部分スキームとする。
$X' \to X$ を中心 $Z$ のブローアップと表す。$b : W \to \mathbf{P}^1_X$ を
中心 $\infty(Z)$ のブローアップとし、$E \subset W$ を例外因子と表す。可換図式
$$
\xymatrix{
X' \ar[r] \ar[d] & W \ar[d]^b \\
X \ar[r]^\infty & \mathbf{P}^1_X
}
$$
があり、その水平射は閉埋め込みである（「因子」補題
\ref{divisors-lemma-strict-transform}）。$E \subset W$ を例外因子、
$W_\infty \subset W$ を $(\mathbf{P}^1_X)_\infty$ の逆像と表す。このとき次が成り立つ。
\begin{enumerate}
\item $b$ は $\mathbf{A}^1_X \cup \mathbf{P}^1_{X \setminus Z}$ 上で同型である。
\item $X'$ は $W$ 上の有効 Cartier 因子である。
\item $X' \cap E$ は $X' \to X$ の例外因子である。
\item $W$ 上の有効 Cartier 因子として $W_\infty = X' + E$ である。
\item $E = \underline{\text{Proj}}_Z(\mathcal{C}_{Z/X, *}[S])$ である。ここで $S$ は
次数 $1$ に置かれた変数である。
\item $X' \cap E = \underline{\text{Proj}}_Z(\mathcal{C}_{Z/X, *})$ である。
\item
\label{chow-item-cone-is-open}
$E \setminus X' = E \setminus (X' \cap E) =
\underline{\Spec}_Z(\mathcal{C}_{Z/X, *}) = C_ZX$ である。
\item
\label{chow-item-find-Z-in-blowup}
閉埋め込み $\mathbf{P}^1_Z \to W$ が存在し、$b$ との合成は包含射
$\mathbf{P}^1_Z \to \mathbf{P}^1_X$ であり、$\infty$ によるその基底変換は合成
$Z \to C_ZX \to E \to W_\infty$ である。ここで最初の射は錐の頂点である。
\end{enumerate}
$\mathcal{C}_{Z/X, *}$ は $X$ における $Z$ の余法代数であることを想起する
（「因子」定義 \ref{divisors-definition-conormal-sheaf}）。また $C_ZX$ は
$X$ における $Z$ の正規錐である（「因子」定義 \ref{divisors-definition-normal-cone}）。

\medskip\noindent
ここで上の番号付き主張を証明する。読者には、以下の証明を読む代わりに、いくつかの例を
自ら計算してみることを強く勧める。

\medskip\noindent
(1) は「因子」補題
\ref{divisors-lemma-blowing-up-gives-effective-Cartier-divisor} の対応する主張から従う。
同じ補題により $E \subset W$ は有効 Cartier 因子である。

\medskip\noindent
「因子」補題 \ref{divisors-lemma-blow-up-pullback-effective-Cartier} により
$W_\infty$ は有効 Cartier 因子である。$E \subset W_\infty$ なので、ある有効 Cartier 因子
$D$ に対して $W_\infty = D + E$ と書ける（「因子」補題
\ref{divisors-lemma-difference-effective-Cartier-divisors}）。後で $D = X'$ を示す。
これにより (2) と (4) が証明される。

\medskip\noindent
$X'$ は閉埋め込み $\infty : X \to \mathbf{P}^1_X$ の狭義変換なので（上を参照）、
$X' \to X$ の例外因子は交わり $X' \cap E$ に等しい。例えば、どちらも $Z$ の
イデアル層の $X'$ への引き戻しで切り出されるからである。これで (3) が従う。

\medskip\noindent
$\infty(Z)$ と $\mathbf{P}^1_Z$ の交わりは有効 Cartier 因子
$(\mathbf{P}^1_Z)_\infty$ である。したがって、ブローアップ $b$ による
$\mathbf{P}^1_Z$ の狭義変換は $\mathbf{P}^1_Z$ へ同型に写る
（「因子」補題 \ref{divisors-lemma-strict-transform} および
\ref{divisors-lemma-blow-up-effective-Cartier-divisor}）。これにより (8) で述べた射
$\mathbf{P}^1_Z \to W$ を得る。$b$ は分離的なので、この射は閉埋め込みである
（「スキーム」補題 \ref{schemes-lemma-section-immersion}）。

\medskip\noindent
$\Spec(A) \subset X$ をアフィン開集合とし、$Z \cap \Spec(A)$ が有限生成イデアル
$I \subset A$ に対応すると仮定する。$\infty(Z \cap \Spec(A))$ のアフィン近傍は、
座標 $s = T_0/T_1$ をもつ $A$ 上のアフィン空間である。
$I$ と $s$ で生成されるイデアルを $J = (I, s) \subset A[s]$ と表す。
$B = A[s] \oplus J \oplus J^2 \oplus \ldots$ を $(A[s], J)$ の Rees 代数とする。
すべての $n \geq 0$ に対し、$A[s]$ の $A$-部分加群として
$$
J^n =
I^n \oplus sI^{n - 1} \oplus s^2I^{n - 2} \ldots \oplus s^nA
\oplus s^{n + 1}A \oplus \ldots
$$
である。開部分スキーム
$$
\text{Proj}(B) = \text{Proj}(A[s] \oplus J \oplus J^2 \oplus \ldots)
\subset W
$$
を考える。最後に、$s \in J$ を $B$ の次数 $1$ の元とみなしたものを $S$ と表す。

\medskip\noindent
$\text{Proj}$ の形成は基底変換と可換なので（「構成」補題
\ref{constructions-lemma-base-change-map-proj}）、
$$
E = \text{Proj}(B \otimes_{A[s]} A/I) =
\text{Proj}((A/I \oplus I/I^2 \oplus I^2/I^3 \oplus \ldots)[S])
$$
である。$B \otimes_{A[s]} A/I = \bigoplus J^n/J^{n + 1}$ が表示された形になることは、
上の $J^n$ の冪の記述から直ちに分かる。「因子」補題
\ref{divisors-lemma-affine-conormal-sheaf} により、$\Spec(A)$ における
$Z \cap \Spec(A)$ の余法代数は次数付き $A$-代数
$A/I \oplus I/I^2 \oplus I^2/I^3 \oplus \ldots$ に対応するので、これで (5) が従う。

\medskip\noindent
$\text{Proj}(B)$ は、$f \in I$ に対するアフィン開集合 $D_+(S)$ と
$D_+(f^{(1)})$ で被覆されることを想起する。これらはそれぞれアフィン・ブローアップ代数
$A[s][\frac{J}{s}]$ と $A[s][\frac{J}{f}]$ のスペクトルである
（「因子」補題 \ref{divisors-lemma-blowing-up-affine} および「代数」定義
\ref{algebra-definition-blow-up}）。これらのアフィン開集合をそれぞれ記述すれば、証明が完了する。

\medskip\noindent
まず開集合 $D_+(S)$、すなわち $A[s][\frac{J}{s}]$ のスペクトルを考える。
上の $J$ の冪の記述から
$$
A[s][\textstyle{\frac{J}{s}}] = \sum s^{-n}I^n[s] \subset A[s, s^{-1}]
$$
を得る。元 $s$ はこの環の非零因子であり、例外因子 $E$ と $W_\infty$ の双方を定める。
したがって $D \cap D_+(S) = \emptyset$ である。最後に、
$A[s][\frac{J}{s}]$ の $s$ による商は余法代数
$$
A/I \oplus I/I^2 \oplus I^2/I^3 \oplus \ldots
$$
である。これで (7) が従う。

\medskip\noindent
次に開集合 $D_+(f^{(1)})$、すなわち $A[s][\frac{J}{f}]$ のスペクトルを考える。
上の $J$ の冪の記述から
$$
A[s][\textstyle{\frac{J}{f}}] =
A[\textstyle{\frac{I}{f}}][\textstyle{\frac{s}{f}}]
$$
を得る。ここで $\frac{s}{f}$ は変数である。元 $f$ はこの環の非零因子であり、
その零スキームが例外因子 $E$ を定める。$s$ は $W_\infty$ を定め、
$s = f \cdot \frac{s}{f}$ なので、$\frac{s}{f}$ は上で構成した因子 $D$ を定める。
したがって
$$
D \cap D_+(f^{(1)}) = \Spec(A[\textstyle{\frac{I}{f}}])
$$
であり、これは $\Spec(A)$ 上のブローアップ $X'$ の対応する開集合である。
実際、$(A, I)$ の Rees 代数への全射な次数付き $A[s]$-代数写像
$B \to A \oplus I \oplus I^2 \oplus \ldots$ は、$\Spec(A[s])$ 上の閉埋め込み
$X' \to W$ に対応する。これで求める $D = X'$ が従う。

\medskip\noindent
(6) を証明しよう。前段落の $\frac{s}{f}$ の零スキームは、アフィン開集合
$D_+(f^{(1)})$ 上における $S$ の零スキームの制限である。したがって $E$ 上で
$S = 0$ が $X' \cap E$ を定める。ゆえに (6) は (5) から従う。

\medskip\noindent
最後に (8) の最後の部分を証明する。これは明らかである。実際、写像
$\mathbf{P}^1_Z \to W$ はアフィン局所的に全射
$$
B \to B \otimes_{A[s]} A/I =
(A/I \oplus I/I^2 \oplus I^2/I^3 \oplus \ldots)[S] \to
A/I[S]
$$
と同一視 $\text{Proj}(A/I[S]) = \Spec(A/I)$ によって与えられる。
いくつかの詳細は省略する。






\section{高余次元の Gysin 準同型}
\label{chow-section-gysin-higher-codimension}

\noindent
$(S, \delta)$ を状況 \ref{chow-situation-setup} のとおりとし、$X$ を $S$ 上局所有限型な
スキームとする。この節では三つ組
$$
(Z \to X, \mathcal{N}, \sigma : \mathcal{N}^\vee \to \mathcal{C}_{Z/X})
$$
を考える。これは閉埋め込み $Z \to X$、局所自由 $\mathcal{O}_Z$-加群
$\mathcal{N}$、および $\mathcal{N}$ の双対から $X$ における $Z$ の余法層への全射
$\sigma : \mathcal{N}^\vee \to \mathcal{C}_{Z/X}$ からなる
（「射」節 \ref{morphisms-section-conormal-sheaf}）。
$\mathcal{N}$ を {\it $X$ における $Z$ の仮想法層} と呼ぶ。

\begin{lemma}
\label{chow-lemma-pullback-virtual-normal-sheaf}
$(S, \delta)$ を状況 \ref{chow-situation-setup} のとおりとする。
$$
\xymatrix{
Z' \ar[r] \ar[d]_g & X' \ar[d]^f \\
Z \ar[r] & X
}
$$
を、水平射が閉埋め込みである、$S$ 上局所有限型なスキームの Cartesian 図式とする。
$\mathcal{N}$ が $X$ における $Z$ の仮想法層ならば、
$\mathcal{N}' = g^*\mathcal{N}$ は $X'$ における $Z'$ の仮想法層である。
\end{lemma}

\begin{proof}
これは「射」補題 \ref{morphisms-lemma-conormal-functorial-flat} で証明された写像
$g^*\mathcal{C}_{Z/X} \to \mathcal{C}_{Z'/X'}$ の全射性から従う。
\end{proof}

\noindent
$(S, \delta)$ を状況 \ref{chow-situation-setup} のとおりとし、$X$ を $S$ 上局所有限型な
スキームとする。$\mathcal{N}$ を閉埋め込み $Z \to X$ の仮想法束とする。この状況で
$$
p : N = \underline{\Spec}_Z(\text{Sym}(\mathcal{N}^\vee)) \longrightarrow Z
$$
を、切断が $\mathcal{N}$ の切断に対応する $Z$ 上のベクトル束とする。
この状況では標準的な閉埋め込み
$$
C_ZX \longrightarrow N_ZX \longrightarrow N
$$
がある。第一の閉埋め込みは「因子」の式
(\ref{divisors-equation-normal-cone-in-normal-bundle}) であり、第二の閉埋め込みは
$\sigma$ が誘導する全射
$\text{Sym}(\mathcal{N}^\vee) \to \text{Sym}(\mathcal{C}_{Z/X})$ に対応する。
$$
b : W \longrightarrow \mathbf{P}^1_X
$$
を節 \ref{chow-section-blowup-Z-first} で構成した $\infty(Z)$ におけるブローアップとする。
補題 \ref{chow-lemma-gysin-at-infty} により、標準的な双変類
$$
C \in A^0(W_\infty \to X)
$$
がある。項目 (\ref{chow-item-cone-is-open}) の開埋め込み $j : C_ZX \to W_\infty$ と、
上で構成した閉埋め込み $i : C_ZX \to N$ を考える。補題 \ref{chow-lemma-vectorbundle} により、
すべての $\alpha \in \CH_k(X)$ に対して
$$
i_*j^*(C \cap \alpha) = p^*\beta
$$
を満たす唯一の $\beta \in \CH_*(Z)$ が存在する。
$c(Z \to X, \mathcal{N}) \cap \alpha = \beta$ と定める。

\begin{lemma}
\label{chow-lemma-construction-gysin}
上の構成は双変類\footnote{記法 $A^*(Z \to X)^\wedge$ は注意
\ref{chow-remark-completion-bivariant} で論じる。$X$ が準コンパクトならば
$A^*(Z \to X)^\wedge = A^*(Z \to X)$ である。}
$$
c(Z \to X, \mathcal{N}) \in A^*(Z \to X)^\wedge
$$
を定め、さらにその構成は補題 \ref{chow-lemma-pullback-virtual-normal-sheaf} の基底変換と両立する。
$\mathcal{N}$ の階数が定数 $r$ ならば、
$c(Z \to X, \mathcal{N}) \in A^r(Z \to X)$ である。
\end{lemma}

\begin{proof}
$i_* \circ j^* \circ C$ と $p^*$ はともに双変類なので（補題
\ref{chow-lemma-flat-pullback-bivariant} および \ref{chow-lemma-push-proper-bivariant}）、
等式
$$
i_* \circ j^* \circ C  = p^* \circ c(Z \to X, \mathcal{N})
$$
を適切に解釈して用い、$c(Z \to X, \mathcal{N})$ を双変類として定義できる。
補題 \ref{chow-lemma-vectorbundle} により $p^*$ は Chow 群上で常に全単射なので、
この定義は有効である。

\medskip\noindent
$X' \to X$, $Z' \to X'$, $\mathcal{N}'$ を補題
\ref{chow-lemma-pullback-virtual-normal-sheaf} のとおりとする。
$c = c(Z \to X, \mathcal{N})$ および $c' = c(Z' \to X', \mathcal{N}')$ と書く。
補題の第二の主張は、注意 \ref{chow-remark-restriction-bivariant} の意味で $c'$ が $c$ の制限であることをいう。
これがすべての局所有限型な $X'/X$ に対して成り立つと主張しているので、形式的な議論により、
$\alpha' \in \CH_k(X')$ に対して $c' \cap \alpha' = c \cap \alpha'$ を確認すれば十分である。
このため、可換図式
$$
\xymatrix{
C_{Z'}X' \ar[d] \ar[r] &
W'_\infty \ar[d] \ar[r] &
W' \ar[d] \ar[r] &
\mathbf{P}^1_{X'} \ar[d] \\
C_ZX \ar[r] &
W_\infty \ar[r] &
W \ar[r] &
\mathbf{P}^1_X
}
$$
があることに注意する。これは閉埋め込み
$$
W' \to W \times_{\mathbf{P}^1_X} \mathbf{P}^1_{X'},\quad
W'_\infty \to W_\infty \times_X X',\quad
C_{Z'}X' \to C_ZX \times_Z Z'
$$
を誘導する。$c \cap \alpha'$ を得るには、補題 \ref{chow-lemma-gysin-at-infty} において射
$W \times_{\mathbf{P}^1_X} \mathbf{P}^1_{X'} \to \mathbf{P}^1_{X'}$ を用いて定義した類
$C \cap \alpha'$ を用いる。一方、$c' \cap \alpha'$ を得るには、射
$W' \to \mathbf{P}^1_{X'}$ を用いて定義した類 $C' \cap \alpha'$ を用いる。
補題 \ref{chow-lemma-gysin-at-infty-independent} により、閉埋め込み
$W'_\infty \to (W \times_{\mathbf{P}^1_X} \mathbf{P}^1_{X'})_\infty$ による
$C' \cap \alpha'$ の押し出しは $C \cap \alpha'$ に等しい。したがって補題
\ref{chow-lemma-flat-pullback-proper-pushforward} により、開集合への引き戻し
$$
C_{Z'}X' \subset W'_\infty,\quad
C_ZX \times_Z Z' \subset (W \times_{\mathbf{P}^1_X} \mathbf{P}^1_{X'})_\infty
$$
についても同じことが成り立つ。可換図式
$$
\xymatrix{
C_{Z'} X' \ar[d] \ar[r] & N' \ar@{=}[d] \\
C_ZX \times_Z Z' \ar[r] & N \times_Z Z'
}
$$
があるので、これらの類は $N'$ 上の同じ類へ押し出される。よって
$\CH_*(Z')$ において同じ元 $c \cap \alpha' = c' \cap \alpha'$ を得る。
\end{proof}

\begin{lemma}
\label{chow-lemma-gysin-decompose}
$(S, \delta)$ を状況 \ref{chow-situation-setup} のとおりとし、$X$ を $S$ 上局所有限型な
スキームとする。$\mathcal{N}$ を $X$ の閉部分スキーム $Z$ の仮想法層とする。
有限局所自由 $\mathcal{O}_Z$-加群の短完全列
$0 \to \mathcal{N}' \to \mathcal{N} \to \mathcal{E} \to 0$ があり、与えられた全射
$\sigma : \mathcal{N}^\vee \to \mathcal{C}_{Z/X}$ が写像
$\sigma' : (\mathcal{N}')^\vee \to \mathcal{C}_{Z/X}$ を経由すると仮定する。
このとき双変類として
$$
c(Z \to X, \mathcal{N}) = c_{top}(\mathcal{E}) \circ c(Z \to X, \mathcal{N}')
$$
である。
\end{lemma}

\begin{proof}
全射 $\mathcal{N}^\vee \to (\mathcal{N}')^\vee$ に対応するベクトル束の閉埋め込みを
$N' \to N$ と表す。このとき閉埋め込み
$$
C_ZX \to N' \to N
$$
がある。したがって、求める双変類の関係は補題 \ref{chow-lemma-easy-virtual-class} から直ちに従う。
\end{proof}

\begin{lemma}
\label{chow-lemma-gysin-excess}
$(S, \delta)$ を状況 \ref{chow-situation-setup} のとおりとする。
水平射が閉埋め込みである、$S$ 上局所有限型なスキームの Cartesian 図式
$$
\xymatrix{
Z' \ar[r] \ar[d]_g & X' \ar[d]^f \\
Z \ar[r] & X
}
$$
を考える。$\mathcal{N}$, resp.\ $\mathcal{N}'$ を
$Z \subset X$, resp.\ $Z' \to X'$ の仮想法層とする。
$Z'$ 上の有限局所自由加群の短完全列
$0 \to \mathcal{N}' \to g^*\mathcal{N} \to \mathcal{E} \to 0$ が与えられ、図式
$$
\xymatrix{
g^*\mathcal{N}^\vee \ar[r] \ar[d] &
(\mathcal{N}')^\vee \ar[d] \\
g^*\mathcal{C}_{Z/X} \ar[r] &
\mathcal{C}_{Z'/X'}
}
$$
が可換であると仮定する。このとき
$$
res(c(Z \to X, \mathcal{N})) =
c_{top}(\mathcal{E}) \circ c(Z' \to X', \mathcal{N}')
$$
が $A^*(Z' \to X')^\wedge$ において成り立つ。
\end{lemma}

\begin{proof}
補題 \ref{chow-lemma-construction-gysin} により
$res(c(Z \to X, \mathcal{N})) = c(Z' \to X', g^*\mathcal{N})$ であり、
等式は補題 \ref{chow-lemma-gysin-decompose} から従う。
\end{proof}

\noindent
$(S, \delta)$ を状況 \ref{chow-situation-setup} のとおりとし、$X$ を $S$ 上局所有限型な
スキームとする。$\mathcal{N}$ を $X$ の閉部分スキーム $Z$ の仮想法層とする。
$Y \to X$ を局所有限型な射とする。$Z \times_X Y \to Y$ が正則閉埋め込みであると仮定する
（「因子」節 \ref{divisors-section-regular-immersions}）。この場合、余法層
$\mathcal{C}_{Z \times_X Y/Y}$ は有限局所自由 $\mathcal{O}_{Z \times_X Y}$-加群であり、短完全列
$$
0 \to \mathcal{E}^\vee \to
\mathcal{N}^\vee|_{Z \times_X Y} \to \mathcal{C}_{Z \times_X Y/Y} \to 0
$$
を得る。商 $\mathcal{N}|_{Y \times_X Z} \to \mathcal{E}$ を、この状況の
{\it 過剰法層} と呼ぶ。

\begin{lemma}
\label{chow-lemma-gysin-fundamental}
直前に述べた状況で $\dim_\delta(Y) = n$ とし、
$\mathcal{C}_{Y \times_X Z/Z}$ の階数が定数 $r$ であると仮定する。このとき
$$
c(Z \to X, \mathcal{N}) \cap [Y]_n =
c_{top}(\mathcal{E}) \cap [Z \times_X Y]_{n - r}
$$
が $\CH_*(Z \times_X Y)$ において成り立つ。
\end{lemma}

\begin{proof}
双変類 $c_{top}(\mathcal{E}) \in A^*(Z \times_X Y)$ は注意
\ref{chow-remark-top-chern-class} で定義した。補題 \ref{chow-lemma-construction-gysin} により
$X$ を $Y$ で置き換えてよい。したがって、$Z \to X$ は余次元 $r$ の正則閉埋め込みであり、
$\dim_\delta(X) = n$ であると仮定して、$\CH_*(Z)$ における等式
$c(Z \to X, \mathcal{N}) \cap [X]_n =
c_{top}(\mathcal{E}) \cap [Z]_{n - r}$ を示せばよい。
さらに補題 \ref{chow-lemma-gysin-decompose} により、
$\mathcal{N}^\vee \to \mathcal{C}_{Z/X}$ が同型であると仮定してよい。すなわち
$\CH_*(Z)$ において
$c(Z \to X, \mathcal{C}_{Z/X}^\vee) \cap [X]_n = [Z]_{n - r}$ を示せばよい。

\medskip\noindent
$c(Z \to X, \mathcal{C}_{Z/X}^\vee) \cap [X]_n$ の定義の各段階を追おう。
$b : W \to \mathbf{P}^1_X$ を $\infty(Z)$ のブローアップとする。まず、
$C \in A^0(W_\infty \to X)$ が補題 \ref{chow-lemma-gysin-at-infty} の類であるとき、
$C \cap [X]_n$ を計算する。$[W]_{n + 1}$ は $W$ 上のサイクルであり、その
$\mathbf{A}^1_X$ への制限は $[X]_n$ の平坦引き戻しに等しい。したがって
$C \cap [X]_n$ は $i_\infty^*[W]_{n + 1}$ に等しい。
$W_\infty$ は $W$ 上の有効 Cartier 因子なので、補題 \ref{chow-lemma-easy-gysin} により
$i_\infty^*[W]_{n + 1} = [W_\infty]_n$ である。この類の開集合
$C_ZX \subset W_\infty$ への制限は、もちろん $[C_ZX]_n$ である。
$Z \subset X$ は正則埋め込みなので、次数付き $\mathcal{O}_Z$-代数として
$$
\mathcal{C}_{Z/X, *} = \text{Sym}(\mathcal{C}_{Z/X})
$$
である（「因子」補題 \ref{divisors-lemma-quasi-regular-immersion}）。したがって
$p : N = C_ZX \to Z$ は、階数 $r$ の有限局所自由加群 $\mathcal{C}_{Z/X}$ に対応する
ベクトル束の構造射である。このとき $p^*[Z]_{n - r} = [C_ZX]_n$ は明らかであり、証明が完了する。
\end{proof}

\begin{lemma}
\label{chow-lemma-gysin-easy}
$(S, \delta)$ を状況 \ref{chow-situation-setup} のとおりとし、$X$ を $S$ 上局所有限型な
スキームとする。$\mathcal{N}$ を $X$ の閉部分スキーム $Z$ の仮想法層とし、
$Y \to X$ を局所有限型な射とする。整数 $r$, $n$ に対して次を仮定する。
\begin{enumerate}
\item $\mathcal{N}$ は階数 $r$ の局所自由層である。
\item $Y$ のすべての既約成分の $\delta$-次元は $n$ である。
\item $\dim_\delta(Z \times_X Y) \leq n - r$ である。
\item $\delta(\xi) = n - r$ を満たす $\xi \in Z \times_X Y$ に対して、局所環
$\mathcal{O}_{Y, \xi}$ は Cohen--Macaulay である。
\end{enumerate}
このとき
$c(Z \to X, \mathcal{N}) \cap [Y]_n = [Z \times_X Y]_{n - r}$ が
$\CH_{n - r}(Z \times_X Y)$ において成り立つ。
\end{lemma}

\begin{proof}
$Z \times_X Y$ は $Y$ の閉部分スキームなので、主張には意味がある。
$\mathcal{N}$ の階数が $r$ なので、$c(Z \to X, \mathcal{N}) \cap [Y]_n$ は
$\CH_{n - r}(Z \times_X Y)$ に属する。$\dim_\delta(Z \cap Y) \leq n - r$ なので、
Chow 群 $\CH_{n - r}(Z \times_X Y)$ は、$\delta$-次元 $n - r$ の既約成分
$W \subset Z \times_X Y$ のサイクル類で自由に生成される。$\xi \in W$ を生成点とする。
仮定 (2) により $\dim(\mathcal{O}_{Y, \xi}) = r$ である。一方、$\mathcal{N}$ の階数は
$r$ であり、$\mathcal{N}^\vee \to \mathcal{C}_{Z/X}$ は全射なので、$Z$ のイデアル層は
局所的に $r$ 個の方程式で切り出される。したがって、$Y$ における $Z \times_X Y$ の
準連接イデアル層 $\mathcal{I} \subset \mathcal{O}_Y$ は局所的に $r$ 個の元で生成される。
$\mathcal{O}_{Y, \xi}$ は次元 $r$ の Cohen--Macaulay 環であり、$\mathcal{I}_\xi$ は
定義イデアルである（$\xi$ は $Z \times_X Y$ の生成点である）から、$\mathcal{I}_\xi$ は
正則列で生成される（「代数」補題 \ref{algebra-lemma-reformulate-CM}）。
「因子」補題 \ref{divisors-lemma-Noetherian-scheme-regular-ideal} により、
$\xi$ のある開近傍 $V \subset Y$ 上で $\mathcal{I}$ は正則列で生成される。
$\CH_{n - r}(Z \times_X Y)$ の上の記述により、
$c(Z \to X, \mathcal{N}) \cap [V]_n = [Z \times_X V]_{n - r}$ を
$\CH_{n - r}(Z \times_X V)$ において示せば十分である。これは、$V$ 上の過剰法層が
$0$ なので、補題 \ref{chow-lemma-gysin-fundamental} から従う。
\end{proof}

\begin{lemma}
\label{chow-lemma-gysin-agrees}
$(S, \delta)$ を状況 \ref{chow-situation-setup} のとおりとし、$X$ を $S$ 上局所有限型な
スキームとする。$(\mathcal{L}, s, i : D \to X)$ を定義
\ref{chow-definition-gysin-homomorphism} のとおりの三つ組とする。
$A^1(D \to X)$ の元とみなした Gysin 準同型 $i^*$（補題
\ref{chow-lemma-gysin-bivariant}）は、$\mathcal{N} = i^*\mathcal{L}$ を
$X$ における $D$ の仮想法層とみなして構成した双変類
$c(D \to X, \mathcal{N}) \in A^1(D \to X)$ と同じである。
\end{lemma}

\begin{proof}
補題 \ref{chow-lemma-bivariant-zero} の判定法を用いる。したがって $X$ は整スキームであると仮定してよく、
$i^*[X]$ が $c \cap [X]$ に等しいことを示せばよい。$n = \dim_\delta(X)$ とおく。
通常どおり、二つの場合がある。

\medskip\noindent
$X = D$ ならば、両方の類は $c_1(\mathcal{N}) \cap [X]_n$ で表される。
補題 \ref{chow-lemma-gysin-fundamental} および定義 \ref{chow-definition-gysin-homomorphism} を参照せよ。

\medskip\noindent
$D \not = X$ ならば、$D \to X$ は有効 Cartier 因子であり、特に余次元 $1$ の正則閉埋め込みである。
再び補題 \ref{chow-lemma-gysin-fundamental} により
$c(D \to X, \mathcal{N}) \cap [X]_n = [D]_{n - 1}$ である。
Gysin 準同型についても定義により同じことが成り立つので、再び結論を得る。
\end{proof}

\begin{lemma}
\label{chow-lemma-gysin-commutes}
$(S, \delta)$ を状況 \ref{chow-situation-setup} のとおりとし、$X$ を $S$ 上局所有限型な
スキームとする。$Z \subset X$ を仮想法層 $\mathcal{N}$ を備えた閉部分スキームとする。
$Y \to X$ を局所有限型とし、$c \in A^*(Y \to X)$ とする。このとき $c$ と
$c(Z \to X, \mathcal{N})$ は可換である（注意 \ref{chow-remark-bivariant-commute}）。
\end{lemma}

\begin{proof}
補題 \ref{chow-lemma-bivariant-zero} を用いて確認できる。したがって $X$ は整スキームであると仮定してよく、
$c \cap c(Z \to X, \mathcal{N}) \cap [X] =
c(Z \to X, \mathcal{N}) \cap c \cap [X]$ を $\CH_*(Z \times_X Y)$ において示せばよい。

\medskip\noindent
$Z = X$ ならば、補題 \ref{chow-lemma-gysin-fundamental} により
$c(Z \to X, \mathcal{N}) = c_{top}(\mathcal{N})$ である。これは双変類 $c$ と可換である
（補題 \ref{chow-lemma-cap-commutative-chern}）。

\medskip\noindent
$Z$ は $X$ に等しくないと仮定する。補題 \ref{chow-lemma-bivariant-zero} により、
$X$ を非零イデアルでブローアップした後に結論を証明すれば十分である。
$X$ を $Z$ のイデアル層でブローアップする。「因子」補題
\ref{divisors-lemma-blowing-up-gives-effective-Cartier-divisor} により、
$Z$ が有効 Cartier 因子である場合に帰着する。

\medskip\noindent
$Z$ が有効 Cartier 因子ならば
$$
c(Z \to X, \mathcal{N}) =
c_{top}(\mathcal{E}) \circ i^*
$$
である。ここで $i^* \in A^1(Z \to X)$ は $i : Z \to X$ に付随する Gysin 準同型
（補題 \ref{chow-lemma-gysin-bivariant}）であり、$\mathcal{E}$ は
$\mathcal{N}^\vee \to \mathcal{C}_{Z/X}$ の核の双対である
（補題 \ref{chow-lemma-gysin-decompose} および \ref{chow-lemma-gysin-agrees}）。
Chern 類は双変環の中心に属し（補題 \ref{chow-lemma-cap-commutative-chern} の強い意味で）、
$c$ は双変類の定義により Gysin 準同型 $i^*$ と可換なので、結論を得る。
\end{proof}

\noindent
$(S, \delta)$ を状況 \ref{chow-situation-setup} のとおりとし、$X$ を $S$ 上局所有限型で
$\delta$-次元 $n$ の整スキームとする。$Z \subset Y \subset X$ を、ともに $X$ の
有効 Cartier 因子である閉部分スキームとする。$X$ における $Y$ の法線錐の零切断を
$o : Y \to C_Y X$ と表す。$C_YX$ は $Y$ 上の直線束なので、補題
\ref{chow-lemma-gysin-bivariant} により双変類 $o^* \in A^1(Y \to C_YX)$ を得る。

\begin{lemma}
\label{chow-lemma-relation-normal-cones}
上の記法のもとで
$$
o^*[C_ZX]_n = [C_Z Y]_{n - 1}
$$
が $\CH_{n - 1}(Y \times_{o, C_Y X} C_ZX)$ において成り立つ。
\end{lemma}

\begin{proof}
$W \to \mathbf{P}^1_X$ を、節 \ref{chow-section-blowup-Z-first} のとおり $\infty(Z)$ の
ブローアップと表す。同様に $W' \to \mathbf{P}^1_X$ を $\infty(Y)$ のブローアップと表す。
$\infty(Z) \subset \infty(Y)$ なので、イデアル層の逆向きの包含、したがってこれらの
ブローアップを定める次数付き代数の写像を得る。これは $W$ から $W'$ への有理射を生じ、
実際に標準的な代表
$$
W \supset U \longrightarrow W'
$$
をもつ（「構成」補題 \ref{constructions-lemma-morphism-relative-proj}）。
局所計算（省略）により、$U$ は少なくとも $\infty$ 上にない $W$ のすべての点と、
特殊ファイバーの開部分スキーム $C_Z X$ を含む。$U$ を縮小して
$U_\infty = C_Z X$ かつ $\mathbf{A}^1_X \subset U$ と仮定してよい。
別の局所計算（省略）により、射 $U_\infty \to W'_\infty$ は、$Z \subset Y$ から来る
イデアル層の包含が誘導する法線錐の標準射
$C_Z X \to C_Y X \subset W'_\infty$ を誘導する。
$W'' \subset W$ を、$W$ における $\mathbf{P}^1_Y \subset \mathbf{P}^1_X$ の狭義変換と表す。
「因子」補題 \ref{divisors-lemma-strict-transform} により、$W''$ は
$\infty(Z)$ における $\mathbf{P}^1_Y$ のブローアップであり、したがって
$(W'' \cap U)_\infty = C_ZY$ である。

\medskip\noindent
項目 (\ref{chow-item-find-Z-in-blowup}) の有効 Cartier 因子
$i : \mathbf{P}^1_Y \to W'$ と、補題 \ref{chow-lemma-gysin-bivariant} のそれに付随する双変類
$i^* \in A^1(\mathbf{P}^1_Y \to W')$ を考える。同様に、無限遠における Gysin 写像を
$(i'_\infty)^* \in A^1(W'_\infty \to W')$ と表す。$i'_\infty$ の $U$ への制限
（注意 \ref{chow-remark-restriction-bivariant}）は、$i_\infty^* \in A^1(W_\infty \to W)$ の
$U$ への制限である。一方では
$$
(i'_\infty)^* i^* [U]_{n + 1} =
i_\infty^* i^* [U]_{n + 1} =
i_\infty^* [(W'' \cap U)_\infty]_{n + 1} =
[C_ZY]_n
$$
である。実際、$i_\infty^*$ は $\infty$ 上に台をもつすべての類を零にし、
$i^*[U]$ と $[W'']$ は $\mathbf{A}^1$ 上のサイクルとして一致し、$C_ZY$ は
$\infty$ 上の $W'' \cap U$ のファイバーである。他方では
$$
(i'_\infty)^* i^* [U]_{n + 1} =
i^* i_\infty^*[U]_{n + 1} =
i^* [U_\infty] =
o^*[C_YX]_n
$$
である。$(i'_\infty)^*$ と $i^*$ は可換であり（補題
\ref{chow-lemma-gysin-commutes-gysin}）、$\infty$ 上の
$i : \mathbf{P}^1_Y \to W'$ のファイバーは $o : Y \to C_YX$ と開埋め込み
$C_YX \to W'_\infty$ の合成として分解するからである。補題が従う。
\end{proof}

\begin{lemma}
\label{chow-lemma-gysin-composition}
$(S, \delta)$ を状況 \ref{chow-situation-setup} のとおりとする。
$Z \subset Y \subset X$ を、$S$ 上局所有限型なスキームの閉部分スキームとする。
$\mathcal{N}$ を $Z \subset X$ の仮想法層、$\mathcal{N}'$ を $Z \subset Y$ の仮想法層、
$\mathcal{N}''$ を $Y \subset X$ の仮想法層とする。可換図式
$$
\xymatrix{
(\mathcal{N}'')^\vee|_Z \ar[r] \ar[d] &
\mathcal{N}^\vee \ar[r] \ar[d] &
(\mathcal{N}')^\vee \ar[d] \\
\mathcal{C}_{Y/X}|_Z \ar[r] &
\mathcal{C}_{Z/X} \ar[r] &
\mathcal{C}_{Z/Y}
}
$$
があると仮定する。ここで下の列は「射の続論」補題
\ref{more-morphisms-lemma-transitivity-conormal} の列であり、上の列は短完全列である。
このとき
$$
c(Z \to X, \mathcal{N}) =
c(Z \to Y, \mathcal{N}') \circ c(Y \to X, \mathcal{N}'')
$$
が $A^*(Z \to X)^\wedge$ において成り立つ。
\end{lemma}

\begin{proof}
仮定は任意の局所有限型な射 $X' \to X$ による基底変換後も満たされることに注意する。
実際、仮想法層の短完全列は局所的に分裂するので、任意の基底変換後も完全である。
したがって双変類の等式を確認するには補題 \ref{chow-lemma-bivariant-zero} を用いてよい。
ゆえに $X$ は整スキームであると仮定して、
$c(Z \to X, \mathcal{N}) \cap [X] =
c(Z \to Y, \mathcal{N}') \cap c(Y \to X, \mathcal{N}'') \cap [X]$ を示せばよい。

\medskip\noindent
$Y = X$ ならば
\begin{align*}
c(Z \to Y, \mathcal{N}') \cap c(Y \to X, \mathcal{N}'') \cap [X]
& =
c(Z \to Y, \mathcal{N}') \cap c_{top}(\mathcal{N}'') \cap [Y] \\
& =
c_{top}(\mathcal{N}''|_Z) \cap c(Z \to Y, \mathcal{N}') \cap [Y] \\
& =
c(Z \to X, \mathcal{N}) \cap [X] 
\end{align*}
である。第一の等式は補題 \ref{chow-lemma-gysin-decompose} による。
第二の等式は Chern 類が双変類と可換であることによる
（補題 \ref{chow-lemma-cap-commutative-chern}）。第三の等式は補題
\ref{chow-lemma-gysin-decompose} による。

\medskip\noindent
$Y \not = X$ と仮定する。補題 \ref{chow-lemma-bivariant-zero} により、
$X$ を非零イデアルでブローアップした後に結論を証明すれば十分である。
$X$ を $Y$ のイデアル層と $Z$ のイデアル層の積でブローアップする。
「因子」補題 \ref{divisors-lemma-blowing-up-gives-effective-Cartier-divisor} および
\ref{divisors-lemma-blowing-up-two-ideals} により、$Y$ と $Z$ がともに
$X$ 上の有効 Cartier 因子である場合に帰着する。

\medskip\noindent
$0 \to \mathcal{E}^\vee \to (\mathcal{N}'')^\vee \to \mathcal{C}_{Y/X} \to 0$ が
短完全列になるような有限局所自由 $\mathcal{O}_Z$-加群の全射を
$\mathcal{N}'' \to \mathcal{E}$ と表す。このとき $\mathcal{N} \to \mathcal{E}|_Z$ も全射である。
この写像の有限局所自由な核を $\mathcal{N}_1$ と表すと、
$\mathcal{N}^\vee \to \mathcal{C}_{Z/X}$ は $\mathcal{N}_1$ を経由する。
補題 \ref{chow-lemma-gysin-decompose} により
$$
c(Y \to X, \mathcal{N}'') = c_{top}(\mathcal{E}) \circ
c(Y \to X, \mathcal{C}_{Y/X}^\vee)
$$
および
$$
c(Z \to X, \mathcal{N}) = c_{top}(\mathcal{E}|_Z) \circ
c(Z \to X, \mathcal{N}_1)
$$
である。束の Chern 類は双変類と可換なので（補題
\ref{chow-lemma-cap-commutative-chern}）、
$$
c(Z \to X, \mathcal{N}_1) =
c(Z \to Y, \mathcal{N}') \circ c(Y \to X, \mathcal{C}_{Y/X}^\vee)
$$
を $A^*(Z \to X)$ において証明すれば十分である。ここで
$\mathcal{N}'' = \mathcal{C}_{Y/X}$ と仮定してよい。これで次の段落で扱う場合に帰着する。

\medskip\noindent
この段落では、$Z$ と $Y$ は次元 $n$ の整スキーム $X$ 上の有効 Cartier 因子であり、
$\mathcal{N}'' = \mathcal{C}_{Y/X}$ である。この場合、補題
\ref{chow-lemma-gysin-fundamental} により
$c(Y \to X, \mathcal{C}_{Y/X}^\vee) \cap [X] = [Y]_{n - 1}$ である。
したがって
$c(Z \to X, \mathcal{N}) \cap [X] = c(Z \to Y, \mathcal{N}') \cap [Y]_{n - 1}$
を示せばよい。$\mathcal{N}$ と $\mathcal{N}'$ に付随する $Z$ 上のベクトル束を
$N$ と $N'$ と表す。$Z$ 上の錐とベクトル束の可換図式
$$
\xymatrix{
N' \ar[r]_i &
N \ar[r] &
(C_Y X) \times_Y Z \\
C_Z Y \ar[r] \ar[u] &
C_Z X \ar[u]
}
$$
を考える。$N'$ は $Z$ 上 $N$ の相対有効 Cartier 因子であり、図式
$$
\xymatrix{
N' \ar[d] \ar[r]_i & N \ar[d] \\
Z \ar[r]^-o & (C_Y X) \times_Y Z
}
$$
は Cartesian である。ここで $o$ は $Y$ 上の直線束 $C_Y X$ の零切断である。
補題 \ref{chow-lemma-relation-normal-cones} により
$o^*[C_ZX]_n = [C_Z Y]_{n - 1}$ が
$$
\CH_{n - 1}(Y \times_{o, C_Y X} C_ZX) =
\CH_{n - 1}(Z \times_{o, (C_Y X) \times_Y Z} C_ZX)
$$
において成り立つ。上の正方形が Cartesian であることから
$$
i^*[C_ZX]_n = [C_Z Y]_{n - 1}
$$
が $\CH_{n - 1}(N')$ において成り立つ。ここで
$\gamma = c(Z \to X, \mathcal{N}) \cap [X]$ と
$\gamma' = c(Z \to Y, \mathcal{N}') \cap [Y]_{n - 1}$ は、それぞれ
$\CH_n(N)$ における $p^*\gamma = [C_Z X]_n$ と、$\CH_{n - 1}(N')$ における
$(p')^*\gamma' = [C_Z Y]_{n - 1}$ によって特徴づけられる。
補題 \ref{chow-lemma-relative-effective-cartier} により $i^* \circ p^* = (p')^*$ なので、証明が完了する。
\end{proof}

\begin{remark}[埋め込みに対する変種]
\label{chow-remark-gysin-for-immersion}
$(S, \delta)$ を状況 \ref{chow-situation-setup} のとおりとする。
$X$ を $S$ 上局所有限型なスキームとし、$i : Z \to X$ をスキームの埋め込みとする。
この状況では次のようにする。
\begin{enumerate}
\item $X$ における $Z$ の余法層 $\mathcal{C}_{Z/X}$ は定義されている
（「射」定義 \ref{morphisms-definition-conormal-sheaf}）。
\item 有限局所自由 $\mathcal{O}_Z$-加群 $\mathcal{N}$ と全射
$\sigma : \mathcal{N}^\vee \to \mathcal{C}_{Z/X}$ からなる対を、埋め込み
$Z \to X$ の仮想法束と呼ぶ。
\item $Z \to X$ が閉埋め込み $Z \to U$ を経由するような開部分スキーム
$U \subset X$ を選び、
$c(Z \to X, \mathcal{N}) = c(Z \to U, \mathcal{N}) \circ (U \to X)^*$ とおく。
\end{enumerate}
双変類 $c(Z \to X, \mathcal{N})$ は開部分スキーム $U$ の選択に依存しない。
すべての補題には、このわずかに一般的な構成に対する直ちに得られる対応物がある。
詳細は省略する。
\end{remark}







\section{いくつかの類の計算}
\label{chow-section-calculate}

\noindent
先へ進むため、上で構成したいくつかの類の値を計算する必要がある。

\begin{lemma}
\label{chow-lemma-compute-koszul}
$(S, \delta)$ を状況 \ref{chow-situation-setup} のとおりとし、$X$ を $S$ 上局所有限型な
スキームとする。$\mathcal{E}$ を階数 $r$ の局所自由 $\mathcal{O}_X$-加群とする。このとき
$$
\prod\nolimits_{n = 0, \ldots, r} c(\wedge^n \mathcal{E})^{(-1)^n} =
1 - (r - 1)! c_r(\mathcal{E}) + \ldots
$$
である。
\end{lemma}

\begin{proof}
分裂原理により、これは $\mathcal{E}$ の Chern 根 $x_1, \ldots, x_r$ 上の多項式環における
計算へ帰着する。節 \ref{chow-section-splitting-principle} を参照せよ。ここで
$$
c(\wedge^n \mathcal{E}) =
\prod\nolimits_{1 \leq i_1 < \ldots < i_n \leq r}
(1 + x_{i_1} + \ldots + x_{i_n})
$$
である。したがって、補題の等式の左辺の対数は
$$
-
\sum\nolimits_{p \geq 1}
\sum\nolimits_{n = 0}^r
\sum\nolimits_{1 \leq i_1 < \ldots < i_n \leq r}
\frac{(-1)^{p + n}}{p}(x_{i_1} + \ldots + x_{i_n})^p
$$
である。先頭の負号に注意せよ。一方
$$
\sum\nolimits_{p \geq 0}
\sum\nolimits_{n = 0}^r
\sum\nolimits_{1 \leq i_1 < \ldots < i_n \leq r}
\frac{(-1)^{p + n}}{p!}(x_{i_1} + \ldots + x_{i_n})^p
=
\prod (1 - e^{-x_i})
$$
である。したがって、この Chern 類の最初の非零項は次数 $r$ にあり、予想された値に等しい。
\end{proof}

\begin{lemma}
\label{chow-lemma-compute-section}
$(S, \delta)$ を状況 \ref{chow-situation-setup} のとおりとし、$X$ を $S$ 上局所有限型な
スキームとする。$\mathcal{C}$ を階数 $r$ の局所自由 $\mathcal{O}_X$-加群とする。射
$$
X = \underline{\text{Proj}}_X(\mathcal{O}_X[T])
\xrightarrow{i}
E = \underline{\text{Proj}}_X(\text{Sym}^*(\mathcal{C})[T])
\xrightarrow{\pi}
X
$$
を考える。このとき $t = 1, \ldots, r - 1$ に対して
$c_t(i_*\mathcal{O}_X) = 0$ であり、$A^0(C \to E)$ において
$$
p^* \circ \pi_* \circ c_r(i_*\mathcal{O}_X) = (-1)^{r - 1}(r - 1)! j^*
$$
である。ここで $j : C \to E$ と $p : C \to X$ は、ベクトル束
$C = \underline{\Spec}(\text{Sym}^*(\mathcal{C}))$ の包含射と構造射である。
\end{lemma}

\begin{proof}
標準写像 $\pi^*\mathcal{C} \to \mathcal{O}_E(1)$ は $i(X)$ に沿ってちょうど消える。
したがって、写像
$$
\pi^*\mathcal{C} \otimes \mathcal{O}_E(-1) \to \mathcal{O}_E
$$
の Koszul 複体は $i_*\mathcal{O}_X$ の分解である。特に
$i_*\mathcal{O}_X$ は $D(\mathcal{O}_E)$ の完全対象であり、その Chern 類は定義されている。
$t = 1, \ldots, t - 1$ に対する $c_t(i_*\mathcal{O}_X)$ の消滅は補題
\ref{chow-lemma-compute-koszul} から従う。同じ補題から
$$
c_r(i_*\mathcal{O}_X) = - (r - 1)!
c_r(\pi^*\mathcal{C} \otimes \mathcal{O}_E(-1))
$$
も得る。一方、補題 \ref{chow-lemma-chern-classes-dual} により
$$
c_r(\pi^*\mathcal{C} \otimes \mathcal{O}_E(-1)) =
(-1)^r c_r(\pi^*\mathcal{C}^\vee \otimes \mathcal{O}_E(1))
$$
であり、$\pi^*\mathcal{C}^\vee \otimes \mathcal{O}_E(1)$ は $i(X)$ に沿って
ちょうど消える切断 $s$ をもつ。

\medskip\noindent
$X$ を $X$ 上局所有限型なスキームで置き換えた後、等式の両辺が
$\alpha \in \CH_*(E)$ に同じ作用をすることを証明すれば十分である。
$C \to X$ はベクトル束なので、$C$ 上のすべてのサイクル類は、ある
$\beta \in \CH_*(X)$ に対する $p^*\beta$ の形である（補題 \ref{chow-lemma-vectorbundle}）。
したがって補題 \ref{chow-lemma-restrict-to-open} により、$\gamma$ が $E \setminus C$ 上に
台をもつように $\alpha = \pi^*\beta + \gamma$ と書ける。上の等式を用いると、
$W \subset E$ が、(a) $C$ と交わらないか、または (b) ある整閉部分スキーム
$Y \subset X$ に対して $W = \pi^{-1}Y$ の形である整閉部分スキームのときに
$$
p^*(\pi_*(c_r(\pi^*\mathcal{C}^\vee \otimes \mathcal{O}_E(1)) \cap [W])) =
j^*[W]
$$
を示せば十分である。切断 $s$ と補題 \ref{chow-lemma-top-chern-class} を用いると、場合 (a) では
$c_r(\pi^*\mathcal{C}^\vee \otimes \mathcal{O}_E(1)) \cap [W] = 0$ であり、場合 (b) では
$c_r(\pi^*\mathcal{C}^\vee \otimes \mathcal{O}_E(1)) \cap [W] = [i(Y)]$ である。
これから結論は容易に従う。詳細は省略する。
\end{proof}

\begin{lemma}
\label{chow-lemma-agreement-with-loc-chern}
$(S, \delta)$ を状況 \ref{chow-situation-setup} のとおりとする。
$i : Z \to X$ を、$S$ 上局所有限型なスキームの間の余次元 $r$ の正則閉埋め込みとする。
$\mathcal{N} = \mathcal{C}_{Z/X}^\vee$ を法層とする。$X$ が準コンパクトである
（または準コンパクトな既約成分をもつ）ならば、$t = 1, \ldots, r - 1$ に対して
$c_t(Z \to X, i_*\mathcal{O}_Z) = 0$ であり、
$$
c_r(Z \to X, i_*\mathcal{O}_Z) = (-1)^{r - 1} (r - 1)! c(Z \to X, \mathcal{N})
\quad\text{in}\quad
A^r(Z \to X)
$$
である。ここで $c_t(Z \to X, i_*\mathcal{O}_Z)$ は定義
\ref{chow-definition-localized-chern} の局所化 Chern 類である。
\end{lemma}

\begin{proof}
任意の $x \in Z$ に対し、$Z \cap \Spec(A) = V(f_1, \ldots, f_r)$ となる
アフィン開近傍 $\Spec(A) \subset X$ を選べる。ここで
$f_1, \ldots, f_r \in A$ は正則列である。「因子」定義
\ref{divisors-definition-regular-immersion} および補題
\ref{divisors-lemma-Noetherian-scheme-regular-ideal} を参照せよ。
このとき、例えば「代数の続論」補題
\ref{more-algebra-lemma-regular-koszul-regular} により、$f_1, \ldots, f_r$ の Koszul 複体は
$A/(f_1, \ldots, f_r)$ の分解である。したがって $A/(f_1, \ldots, f_r)$ は
$A$-加群として完全である。ゆえに $F = i_*\mathcal{O}_Z$ は $D(\mathcal{O}_X)$ の完全対象であり、
その $X \setminus Z$ への制限は零である。$X$ が準コンパクトである
（または準コンパクトな既約成分をもつ）という仮定により、局所化 Chern 類
$c_t(Z \to X, i_*\mathcal{O}_Z)$ は定義されている。状況
\ref{chow-situation-loc-chern} および定義 \ref{chow-definition-localized-chern} を参照せよ。
以上により主張には意味がある。

\medskip\noindent
$b : W \to \mathbf{P}^1_X$ を節 \ref{chow-section-blowup-Z-first} のとおり
$\infty(Z)$ におけるブローアップと表す。項目 (\ref{chow-item-find-Z-in-blowup}) により閉埋め込み
$$
i' : \mathbf{P}^1_Z \longrightarrow W
$$
がある。$Q = i'_*\mathcal{O}_{\mathbf{P}^1_Z}$ は $D(\mathcal{O}_W)$ の完全対象であり、
$F$ と $Q$ は補題 \ref{chow-lemma-independent-loc-chern-bQ} の仮定を満たすと主張する。

\medskip\noindent
この主張を仮定する。補題 \ref{chow-lemma-independent-loc-chern-bQ} の結論により、
すべての $p \geq 1$ に対して
$$
c_p(Z \to X, F) = c'_p(Q) = (E \to Z)_* \circ c'_p(Q|_E) \circ C
$$
である。$Q|_E$ は、$i'$ を $\infty$ で基底変換して得られる射 $Z \to E$ による
$Z$ の構造層の押し出しに等しい。したがって $E \to Z$ に補題
\ref{chow-lemma-compute-section} を適用すると、$1 \leq t \leq r - 1$ に対する
$c_t(Z \to X, F)$ の消滅が従う。$\mathcal{C}_{Z/X} = \mathcal{N}^\vee$ は局所自由なので、
双変類 $c(Z \to X, \mathcal{N})$ は関係
$$
j^* \circ C = p^* \circ c(Z \to X, \mathcal{N})
$$
によって特徴づけられる。ここで $j : C_ZX \to W_\infty$ と $p : C_ZX \to Z$ は
与えられた写像である（$C \in A^0(W_\infty \to X)$ は補題
\ref{chow-lemma-gysin-at-infty} の類であることを想起せよ）。したがって、補題の主張の表示式は
補題 \ref{chow-lemma-compute-section} の対応する等式から従う。

\medskip\noindent
主張を証明する。$A$ と $f_1, \ldots, f_r$ を上のとおりとする。節
\ref{chow-section-blowup-Z-first} のとおり、アフィン開集合
$\Spec(A[s]) \subset \mathbf{P}^1_X$ を考える。この開集合上で $s = 0$ が
$(\mathbf{P}^1_X)_\infty$ を定めることを想起せよ。したがって $\Spec(A[s])$ 上では、
正則列 $s, f_1, \ldots, f_r$ で生成されるイデアルをブローアップしている。
「代数の続論」補題 \ref{more-algebra-lemma-blowup-regular-sequence} により、$r + 1$ 個の
アフィン・チャートは $A[s]$ 上の大域的完全交差である。アフィン・ブローアップ代数
$$
A[s][f_1/s, \ldots, f_r/s] = A[s, y_1, \ldots, y_r]/(sy_i - f_i)
$$
に対応するチャートは、$y_1, \ldots, y_r$ で切り出される閉部分スキームとして
$i'(Z \cap \Spec(A))$ を含む。$y_1, \ldots, y_r, sy_1 - f_1, \ldots, sy_r - f_r$ は
多項式環 $A[s, y_1, \ldots, y_r]$ の正則列なので、$i'$ は正則埋め込みである。
いくつかの詳細は省略する。上と同様に、$Q = i'_*\mathcal{O}_{\mathbf{P}^1_Z}$ は
$D(\mathcal{O}_W)$ の完全対象である。補題 \ref{chow-lemma-independent-loc-chern-bQ} における
$F$ と $Q$ に関するその他の仮定（および補題 \ref{chow-lemma-localized-chern-pre} の仮定）は
直ちに確認できる。
\end{proof}

\begin{lemma}
\label{chow-lemma-actual-computation}
補題 \ref{chow-lemma-agreement-with-loc-chern} の状況で $\dim_\delta(X) = n$ とする。
このとき次が成り立つ。
\begin{enumerate}
\item $t = 1, \ldots, r - 1$ に対して
$c_t(Z \to X, i_*\mathcal{O}_Z) \cap [X]_n = 0$ である。
\item $c_r(Z \to X, i_*\mathcal{O}_Z) \cap [X]_n =
(-1)^{r - 1}(r - 1)![Z]_{n - r}$ である。
\item $t = 0, \ldots, r - 1$ に対して
$ch_t(Z \to X, i_*\mathcal{O}_Z) \cap [X]_n = 0$ である。
\item $ch_r(Z \to X, i_*\mathcal{O}_Z) \cap [X]_n = [Z]_{n - r}$ である。
\end{enumerate}
\end{lemma}

\begin{proof}
(1) と (2) は補題 \ref{chow-lemma-agreement-with-loc-chern} と補題
\ref{chow-lemma-gysin-fundamental} を合わせれば直ちに従う。
次に、補題 \ref{chow-lemma-loc-chern-character} で与えられた
$ch_p = (1/p!)P_p$ と $c_p$ の関係を用いて (3) と (4) を導く。
実際、$c_1 = c_2 = \ldots = c_{r - 1} = 0$ ならば
$(-1)^{r - 1}(r - 1)!ch_r = c_r$ である。
\end{proof}







\section{Adams 作用素}
\label{chow-section-adams}

\noindent
第二 Adams 作用素を定義するために必要な最小限の作業を行う。
$X$ をスキームとする。$\textit{Vect}(X)$ は有限局所自由 $\mathcal{O}_X$-加群の圏を
表すことを想起する。さらに「スキームの導来圏」節
\ref{perfect-section-K-groups} で、この圏に付随する第零 $K$-群
$K_0(\textit{Vect}(X))$ を構成した。また $K_0(\textit{Vect}(X))$ は環である
（「スキームの導来圏」注意 \ref{perfect-remark-K-ring}）。

\begin{lemma}
\label{chow-lemma-second-adams-operator}
$X$ をスキームとする。環準同型
$$
\psi^2 :
K_0(\textit{Vect}(X))
\longrightarrow
K_0(\textit{Vect}(X))
$$
が存在する。これは $\mathcal{L}$ が可逆ならば $[\mathcal{L}]$ を
$[\mathcal{L}^{\otimes 2}]$ へ送り、引き戻しと両立する。
\end{lemma}

\begin{proof}
$X$ をスキームとし、$\mathcal{E}$ を有限局所自由 $\mathcal{O}_X$-加群とする。
$K_0(\textit{Vect}(X))$ の元
$$
\psi^2(\mathcal{E}) = [\text{Sym}^2(\mathcal{E})] - [\wedge^2(\mathcal{E})]
$$
を考える。

\medskip\noindent
$X$ をスキームとし、有限局所自由 $\mathcal{O}_X$-加群の短完全列
$$
0 \to \mathcal{E} \to \mathcal{F} \to \mathcal{G} \to 0
$$
を考える。これを $\mathcal{F}$ 上の $2$ 段階のフィルトレーションとみなす。
$\text{Sym}^2(\mathcal{F})$ 上に誘導されるフィルトレーションは $3$ 段階をもち、
次数付き部分は $\text{Sym}^2(\mathcal{E})$, $\mathcal{E} \otimes \mathcal{F}$,
$\text{Sym}^2(\mathcal{G})$ である。したがって
$$
[\text{Sym}^2(\mathcal{F})] =
[\text{Sym}^2(\mathcal{E})] +
[\mathcal{E} \otimes \mathcal{F}] +
[\text{Sym}^2(\mathcal{G})]
$$
である。まったく同様に
$$
[\wedge^2(\mathcal{F})] =
[\wedge^2(\mathcal{E})] +
[\mathcal{E} \otimes \mathcal{F}] +
[\wedge^2(\mathcal{G})]
$$
を得る。したがって
$\psi^2(\mathcal{F}) = \psi^2(\mathcal{E}) + \psi^2(\mathcal{G})$ である。
よって、よく定義された加法写像
$\psi^2 : K_0(\textit{Vect}(X)) \to K_0(\textit{Vect}(X))$ を得る。

\medskip\noindent
この写像が引き戻しと可換であることは明らかである。

\medskip\noindent
残るのは $\psi^2$ が環準同型であることの証明である。$X$ をスキームとし、
$\mathcal{E}$ と $\mathcal{F}$ を有限局所自由 $\mathcal{O}_X$-加群とする。
短完全列
$$
0 \to \wedge^2(\mathcal{E}) \otimes \wedge^2(\mathcal{F}) \to
\text{Sym}^2(\mathcal{E} \otimes \mathcal{F}) \to
\text{Sym}^2(\mathcal{E}) \otimes \text{Sym}^2(\mathcal{F}) \to 0
$$
がある。第一の写像は $(e \wedge e') \otimes (f \wedge f')$ を
$(e \otimes f)(e' \otimes f') - (e' \otimes f)(e \otimes f')$ へ送り、
第二の写像は $(e \otimes f) (e' \otimes f')$ を $ee' \otimes ff'$ へ送る。
同様に短完全列
$$
0 \to \text{Sym}^2(\mathcal{E}) \otimes \wedge^2(\mathcal{F}) \to
\wedge^2(\mathcal{E} \otimes \mathcal{F}) \to
\wedge^2(\mathcal{E}) \otimes \text{Sym}^2(\mathcal{F}) \to 0
$$
がある。第一の写像は $e e' \otimes f \wedge f'$ を
$(e \otimes f) \wedge (e' \otimes f') + (e' \otimes f) \wedge (e \otimes f')$ へ送り、
第二の写像は $(e \otimes f) \wedge (e' \otimes f')$ を
$(e \wedge e') \otimes (f f')$ へ送る。上と同様に、これで $\psi^2$ が乗法的であることが分かる。
$\psi^2(1) = 1$ は明らかなので、証明が完了する。
\end{proof}

\begin{remark}
\label{chow-remark-adams-derived}
$X$ をスキームとし、$X$ 上で $2$ が可逆であるとする。このとき Adams 作用素 $\psi^2$ は
$K$-群 $K_0(X) = K_0(D_{perf}(\mathcal{O}_X))$ 上に直接的な方法で定義できる
（「スキームの導来圏」定義 \ref{perfect-definition-K-group}）。
実際、$X$ 上の完全複体 $L$ が与えられると、因子の交換により群 $\{\pm 1\}$ が
$L \otimes^\mathbf{L} L$ に作用する。そこで
$$
\psi^2(L) = [(L \otimes^\mathbf{L} L)^+] -
[(L \otimes^\mathbf{L} L)^-]
$$
とおける。ここで $(-)^+$ は不変部分を取ること、$(-)^-$ は反不変部分を取ることを
適切に定義したものである。不変部分と反不変部分を取る操作の完全性を用いれば、
補題 \ref{chow-lemma-second-adams-operator} の証明と同様に、これがよく定義されることを示せる。
$2$ が $X$ 上で可逆でない場合は状況がはるかに複雑になり、別の方法が必要である。
\end{remark}

\begin{lemma}
\label{chow-lemma-minus-adams-operator}
$X$ をスキームとする。環準同型
$\psi^{-1} : K_0(\textit{Vect}(X)) \to K_0(\textit{Vect}(X))$ が存在する。
これは $\mathcal{E}$ が有限局所自由ならば $[\mathcal{E}]$ を $[\mathcal{E}^\vee]$ へ送り、
引き戻しと両立する。
\end{lemma}

\begin{proof}
確認すべきことは、双対を取る操作が短完全列および引き戻しと両立することだけである。
これは明らかである。
\end{proof}

\begin{remark}
\label{chow-remark-chern-classes-K}
$(S, \delta)$ を状況 \ref{chow-situation-setup} のとおりとし、$X$ を $S$ 上局所有限型とする。
Chern 類写像は、左辺の生成元 $[\mathcal{E}]$ を
$c(\mathcal{E}) = 1 + c_1(\mathcal{E}) + c_2(\mathcal{E}) + \ldots$ へ送り、
乗法的に拡張することにより、標準写像
$$
c : K_0(\textit{Vect}(X)) \longrightarrow \prod\nolimits_{i \geq 0} A^i(X)
$$
を定める。したがって $-[\mathcal{E}]$ は形式的逆元 $c(\mathcal{E})^{-1}$ へ写る。
右辺に無限積が現れるのはこのためである。これは補題
\ref{chow-lemma-additivity-chern-classes} によりよく定義される。
\end{remark}

\begin{remark}
\label{chow-remark-chern-character-K}
$(S, \delta)$ を状況 \ref{chow-situation-setup} のとおりとし、$X$ を $S$ 上局所有限型とする。
Chern 指標写像は、左辺の生成元 $[\mathcal{E}]$ を $ch(\mathcal{E})$ へ送り、
加法的に拡張することにより、標準的な環準同型
$$
ch : K_0(\textit{Vect}(X)) \longrightarrow
\prod\nolimits_{i \geq 0} A^i(X) \otimes \mathbf{Q}
$$
を定める。これは補題 \ref{chow-lemma-chern-character-additive} によりよく定義され、
補題 \ref{chow-lemma-chern-character-multiplicative} により環準同型である。
\end{remark}

\begin{lemma}
\label{chow-lemma-adams-and-chern}
$(S, \delta)$ を状況 \ref{chow-situation-setup} のとおりとし、$X$ を $S$ 上局所有限型とする。
$\psi^2$ を補題 \ref{chow-lemma-second-adams-operator} のとおりとし、$c$ と $ch$ を注意
\ref{chow-remark-chern-classes-K} および \ref{chow-remark-chern-character-K} のとおりとする。
このとき、すべての $\alpha \in K_0(\textit{Vect}(X))$ に対して
$c_i(\psi^2(\alpha)) = 2^i c_i(\alpha)$ および
$ch_i(\psi^2(\alpha)) = 2^i ch_i(\alpha)$ である。
\end{lemma}

\begin{proof}
$A^i(X)$ 上で $2^i$ を掛ける写像
$\prod_{i \geq 0} A^i(X) \to \prod_{i \geq 0} A^i(X)$ は環準同型である。
$\psi^2$ も環準同型なので、$K_0(\textit{Vect}(X))$ の加法的生成元に対して公式を
証明すれば十分である。したがって、ある有限局所自由 $\mathcal{O}_X$-加群
$\mathcal{E}$ に対して $\alpha = [\mathcal{E}]$ と仮定してよい。
$\mathcal{E}$ の Chern 類の構成により、$\mathcal{E}$ の階数が定数 $r$ である場合に
直ちに帰着する（注意 \ref{chow-remark-extend-to-finite-locally-free}）。
この場合、制限 $A^*(X) \to A^*(P)$ が単射であり、$p^*\mathcal{E}$ が
可逆 $\mathcal{O}_P$-加群 $\mathcal{L}_j$ を次数付き部分とする有限フィルトレーションを
もつような射影的平滑射 $p : P \to X$ を選べる
（補題 \ref{chow-lemma-splitting-principle}）。このとき
$[p^*\mathcal{E}] = \sum [\mathcal{L}_j]$ であり、したがって $\psi^2$ の定義により
$\psi^2([p^\mathcal{E}]) = \sum [\mathcal{L}_j^{\otimes 2}]$ である。
$x_j  = c_1(\mathcal{L}_j)$ とおくと
$$
c(\alpha) = \prod (1 + x_j)
\quad\text{and}\quad
c(\psi^2(\alpha)) = \prod (1 + 2 x_j)
$$
が $\prod A^i(P)$ において成り立ち、
$$
ch(\alpha) = \sum \exp(x_j)
\quad\text{and}\quad
ch(\psi^2(\alpha)) = \sum \exp(2 x_j)
$$
も $\prod A^i(P)$ において成り立つ。これらの公式から求める結論が従う。
\end{proof}

\begin{remark}
\label{chow-remark-perf-Z-cohomology-K}
$X$ を局所 Noether スキームとし、$Z \subset X$ を閉部分スキームとする。
コホモロジー層が $Z$ 上に集合論的に台をもつ $\mathcal{O}_X$-加群の完全複体からなる、
厳密充満で飽和な三角部分圏
$$
D_{Z, perf}(\mathcal{O}_X) \subset D(\mathcal{O}_X)
$$
を考える。集合論的台が $Z$ に含まれる連接 $\mathcal{O}_X$-加群の Serre 部分圏を
$\textit{Coh}_Z(X) \subset \textit{Coh}(X)$ と表す。
$E \in D_{Z, perf}(\mathcal{O}_X)$ が与えられると、$X$ 上 Zariski 局所的には有限個の
コホモロジー層 $H^i(E)$ だけが非零であり、それらはすべて $Z$ 上に集合論的に台をもつ。
したがって、規則
$$
E \longmapsto
[\bigoplus\nolimits_{i \in \mathbf{Z}} H^{2i}(E)] -
[\bigoplus\nolimits_{i \in \mathbf{Z}} H^{2i + 1}(E)]
$$
により写像
$$
K_0(D_{Z, perf}(\mathcal{O}_X))
\longrightarrow
K_0(\textit{Coh}_Z(X)) = K'_0(Z)
$$
を定義できる。等号は補題 \ref{chow-lemma-K-coherent-supported-on-closed} による。
これは、$D_{Z, perf}(\mathcal{O}_X)$ の識別三角が与えられるとコホモロジー層の
長完全列を得るため有効である。
\end{remark}

\begin{remark}
\label{chow-remark-perf-Z-regular}
$X$, $Z$, $D_{Z, perf}(\mathcal{O}_X)$ を注意
\ref{chow-remark-perf-Z-cohomology-K} のとおりとし、$X$ が正則であると仮定する。
このとき標準写像
$$
K_0(\textit{Coh}(Z)) \longrightarrow K_0(D_{Z, perf}(\mathcal{O}_X))
$$
があり、次のように定義される。任意の連接 $\mathcal{O}_Z$-加群 $\mathcal{F}$ に対し、
$D(\mathcal{O}_X)$ のうち次数 $0$ に $\mathcal{F}$ をもち、その他の次数では零である対象を
$\mathcal{F}[0]$ と表す。「スキームの導来圏」補題
\ref{perfect-lemma-perfect-on-regular} により $\mathcal{F}[0]$ は $X$ 上の完全複体である。
したがって $\mathcal{F}[0]$ は $D_{Z, perf}(\mathcal{O}_X)$ の対象である。一方、
連接 $\mathcal{O}_Z$-加群の短完全列
$0 \to \mathcal{F} \to \mathcal{F}' \to \mathcal{F}'' \to 0$ が与えられると、識別三角
$\mathcal{F}[0] \to \mathcal{F}'[0] \to \mathcal{F}''[0] \to \mathcal{F}[1]$ を得る
（「導来圏」節 \ref{derived-section-canonical-delta-functor}）。したがって
$[\mathcal{F}]$ を $[\mathcal{F}[0]]$ へ送ることで写像
$K_0(\textit{Coh}(Z)) \to K_0(D_{Z, perf}(\mathcal{O}_X))$ を得る。
この恐ろしい記法については容赦されたい。
\end{remark}

\begin{lemma}
\label{chow-lemma-perf-Z-regular}
$X$ を Noether 正則スキームとし、$Z \subset X$ を閉部分スキームとする。
注意 \ref{chow-remark-perf-Z-cohomology-K} および \ref{chow-remark-perf-Z-regular} で構成した写像は
互いに逆であり、$K'_0(Z) = K_0(D_{Z, perf}(\mathcal{O}_X))$ を得る。
\end{lemma}

\begin{proof}
合成
$$
K_0(\textit{Coh}(Z)) \longrightarrow
K_0(D_{Z, perf}(\mathcal{O}_X)) \longrightarrow
K_0(\textit{Coh}(Z))
$$
が恒等写像であることは明らかである。したがって、第一の射が全射であることを示せば十分である。
$E$ を $D_{Z, perf}(\mathcal{O}_X)$ の対象とする。「スキームの導来圏」補題
\ref{perfect-lemma-perfect-on-regular} により
$D_{perf}(\mathcal{O}_X) = D^b_{\textit{Coh}}(\mathcal{O}_X)$ であることを想起する。
したがって、コホモロジー $H^i(E)$ は連接であり、$D_{Z, perf}(\mathcal{O}_X)$ の対象と
みなせ、有限個だけが非零である。標準切断の識別三角を用いると
$$
[E] = \sum (-1)^i[H^i(E)[0]]
$$
が $K_0(D_{Z, perf}(\mathcal{O}_X))$ において成り立つことが分かる。
したがって、$Z$ 上に集合論的に台をもつ任意の連接 $\mathcal{O}_X$-加群に対して
$[\mathcal{F}[0]]$ が写像の像に属することを示せば十分である。
$\mathcal{F}$ には、商が $\mathcal{O}_Z$-加群である有限フィルトレーションを取れるので、
証明が完了する。
\end{proof}

\begin{remark}
\label{chow-remark-localized-chern-classes-K}
$(S, \delta)$ を状況 \ref{chow-situation-setup} のとおりとし、$X$ を $S$ 上局所有限型とする。
$Z \subset X$ を閉部分スキームとし、$D_{Z, perf}(\mathcal{O}_X)$ を注意
\ref{chow-remark-perf-Z-cohomology-K} のとおりとする。$X$ が準コンパクトである
（または、より一般に $X$ の既約成分が準コンパクトである）ならば、局所化 Chern 類は、
左辺の生成元 $[E]$ を
$$
c(Z \to X, E) = 1 + c_1(Z \to X, E) + c_2(Z \to X, E) + \ldots
$$
へ送り、乗法的に拡張することにより、標準写像
$$
c(Z \to X, -) : K_0(D_{Z, perf}(\mathcal{O}_X)) \longrightarrow
A^0(X) \times \prod\nolimits_{i \geq 1} A^i(Z \to X)
$$
を定める。右辺の積は注意 \ref{chow-remark-ring-loc-classes} のとおりである。
$X$ の準コンパクト性条件により、局所化 Chern 類が定義され
（状況 \ref{chow-situation-loc-chern} および定義 \ref{chow-definition-localized-chern}）、
これらの局所化 Chern 類が識別三角を双変 Chow 環における対応する積へ変換することが
保証される（補題 \ref{chow-lemma-additivity-loc-chern-c}）。
\end{remark}

\begin{remark}
\label{chow-remark-localized-chern-character-K}
$(S, \delta)$ を状況 \ref{chow-situation-setup} のとおりとし、$X$ を $S$ 上局所有限型とする。
$Z \subset X$ を閉部分スキームとし、$D_{Z, perf}(\mathcal{O}_X)$ を注意
\ref{chow-remark-perf-Z-cohomology-K} のとおりとする。$X$ の既約成分が準コンパクトならば、
局所化 Chern 指標は、左辺の生成元 $[E]$ を $ch(Z \to X, E)$ へ送り、加法的に拡張することで、
標準的な加法的かつ乗法的な写像
$$
ch(Z \to X, -) : K_0(D_{Z, perf}(\mathcal{O}_X)) \longrightarrow
\prod\nolimits_{i \geq 0} A^i(Z \to X) \otimes \mathbf{Q}
$$
を定める。実際、$X$ の既約成分に関する条件により、局所化 Chern 指標が定義され
（状況 \ref{chow-situation-loc-chern} および定義 \ref{chow-definition-localized-chern}）、
これらの局所化 Chern 指標が識別三角を双変 Chow 環における対応する和へ変換することが
保証される（補題 \ref{chow-lemma-additivity-loc-chern-P}）。
$K_0(D_{Z, perf}(X))$ 上の乗法は導来テンソル積を用いて定義される
（「スキームの導来圏」注意 \ref{perfect-remark-perf-Z}）。したがって補題
\ref{chow-lemma-loc-chern-tensor-product} により
$ch(Z \to X, \alpha \beta) = ch(Z \to X, \alpha) ch(Z \to X, \beta)$ である。
$X$ が準コンパクトならば、写像 $ch(Z \to X, -)$ の像は
$A^*(Z \to X) \otimes \mathbf{Q}$ に含まれる。詳細は省略する。
\end{remark}

\begin{remark}
\label{chow-remark-chern-classes-agree}
$(S, \delta)$ を状況 \ref{chow-situation-setup} のとおりとする。$X$ を $S$ 上局所有限型とし、
$X$ が準コンパクトである（または、より一般に $X$ の既約成分が準コンパクトである）と仮定する。
$Z = X$ とし、注意 \ref{chow-remark-localized-chern-classes-K} および
\ref{chow-remark-localized-chern-character-K} の記法を用いると、
$D_{Z, perf}(\mathcal{O}_X) = D_{perf}(\mathcal{O}_X)$ であり、
$$
K_0(D_{Z, perf}(\mathcal{O}_X)) = K_0(D_{perf}(\mathcal{O}_X)) = K_0(X)
$$
である。「スキームの導来圏」定義 \ref{perfect-definition-K-group} を参照せよ。
したがって、注意 \ref{chow-remark-localized-chern-classes-K} および
\ref{chow-remark-localized-chern-character-K} の特別な場合として
$$
c : K_0(X) \to \prod A^i(X)
\quad\text{and}\quad
ch : K_0(X) \to \prod A^i(X) \otimes \mathbf{Q}
$$
を得る。もちろん代わりに、定義 \ref{chow-definition-defined-on-perfect}、補題
\ref{chow-lemma-additivity-on-perfect}、補題 \ref{chow-lemma-chern-classes-perfect-tensor-product} を
直接用いてこれらの写像を構成することもできる。これが同じ類を与えることは直ちに分かる。
有限局所自由加群を、次数 $0$ に置かれた完全複体とみなしてそれ自身へ送る標準写像
$K_0(\textit{Vect}(X)) \to K_0(X)$ があることを想起する
（「スキームの導来圏」節 \ref{perfect-section-K-groups}）。このとき図式
$$
\xymatrix{
K_0((\textit{Vect}(X)) \ar[rd]_c \ar[rr] & &
K_0(D_{perf}(\mathcal{O}_X)) = K_0(X) \ar[ld]^c \\
& \prod A^i(X)
}
$$
は可換である。ここで南東向きの射は注意 \ref{chow-remark-chern-classes-K} で構成したものである。
同様に図式
$$
\xymatrix{
K_0((\textit{Vect}(X)) \ar[rd]_{ch} \ar[rr] & &
K_0(D_{perf}(\mathcal{O}_X)) = K_0(X) \ar[ld]^{ch} \\
& \prod A^i(X) \otimes \mathbf{Q}
}
$$
は可換である。ここで南東向きの射は注意 \ref{chow-remark-chern-character-K} で構成したものである。
\end{remark}











\section{Chow 群と K 群再訪}
\label{chow-section-chow-and-K-II}

\noindent
本節は節 \ref{chow-section-chow-and-K} の続きである。
$(S, \delta)$ を状況 \ref{chow-situation-setup} のとおりとし、
$X$ を $S$ 上局所有限型とする。$X$ 上の連接層の K 群
$K'_0(X) = K_0(\textit{Coh}(X))$ には標準的な増大フィルトレーション
$$
F_kK'_0(X) =
\Im\Big(K_0(\textit{Coh}_{\leq k}(X)) \to K_0(\textit{Coh}(X)\Big)
$$
がある。これを台の次元によるフィルトレーションという。さらに
$$
\text{gr}_k K'_0(X) \subset K'_0(X)/F_{k - 1}K'_0(X) =
K_0(\textit{Coh}(X)/\textit{Coh}_{\leq k - 1}(X))
$$
であり、ここで等号は『ホモロジー代数』の補題 \ref{homology-lemma-serre-subcategory-K-groups} による。
注意 \ref{chow-remark-good-cases-K-A} の議論から、標準的な写像
$$
\CH_k(X) \longrightarrow \text{gr}_k K'_0(X)
$$
が得られる。これは、$\delta$-次元 $k$ の整閉部分スキーム
$Z \subset X$ の類を右辺の $[\mathcal{O}_Z]$ の類へ送ることで定義される。

\begin{proposition}
\label{chow-proposition-K-tensor-Q}
$(S, \delta)$ を状況 \ref{chow-situation-setup} のとおりとする。$S$ 上局所有限型な
スキームの閉埋め込み $X \to Y$ が与えられ、$Y$ は正則かつ準コンパクトであると仮定する。
このとき、注意 \ref{chow-remark-perf-Z-regular} の写像
$\mathcal{F} \mapsto \mathcal{F}[0]$、注意
\ref{chow-remark-localized-chern-character-K} の写像 $ch(X \to Y, -)$、および
写像 $c \mapsto c \cap [Y]$ の合成
$$
K'_0(X) \to
K_0(D_{X, perf}(\mathcal{O}_Y)) \to
A^*(X \to Y) \otimes \mathbf{Q} \to
\CH_*(X) \otimes \mathbf{Q}
$$
は、同型
$$
K'_0(X) \otimes \mathbf{Q}
\longrightarrow
\CH_*(X) \otimes \mathbf{Q}
$$
を誘導する。この同型は $Y$ の選択に依存する。さらに、上で述べた標準的な写像
$$
\CH_k(X) \otimes \mathbf{Q}
\longrightarrow
\text{gr}_k K'_0(X) \otimes \mathbf{Q}
$$
は、すべての $k \in \mathbf{Z}$ に対して $\mathbf{Q}$-ベクトル空間の同型である。
\end{proposition}

\begin{proof}
$Y$ は正則なので、注意 \ref{chow-remark-perf-Z-regular} の構成を適用できる。
$Y$ は準コンパクトなので、注意 \ref{chow-remark-localized-chern-character-K} の構成も適用できる。
さらに $Y$ は局所純次元であり（補題 \ref{chow-lemma-locally-equidimensional}）、したがって
「基本サイクル」$[Y]$ は $\CH_*(Y)$ の元として定義される。注意
\ref{chow-remark-fundamental-class} を参照されたい。これと上で構成した写像
$\CH_k(X) \to \text{gr}_kK'_0(X)$ を合わせると、次を示せば十分である。
\begin{enumerate}
\item $\mathcal{F}$ が連接 $\mathcal{O}_X$-加群で、その台の $\delta$-次元が
$\leq k$ ならば、上の合成は $[\mathcal{F}]$ を
$\bigoplus_{k' \leq k} \CH_{k'}(X) \otimes \mathbf{Q}$ へ送る。
\item $Z \subset X$ が $\delta$-次元 $k$ の整閉部分スキームならば、上の合成による
$[\mathcal{O}_Z]$ の像の次数 $k$ 成分は、$\CH_k(X) \otimes \mathbf{Q}$ における
$[Z]$ の類である。
\end{enumerate}
実際、これらが成り立てば、我々の写像は
$\text{gr}_kK'_0(X) \otimes \mathbf{Q} \to CH_k(X) \otimes \mathbf{Q}$
を誘導し、これは命題の上で与えた標準的な写像 
$\CH_k(X) \otimes \mathbf{Q} \to \text{gr}_k K'_0(X) \otimes \mathbf{Q}$
の逆写像となる。

\medskip\noindent
連接 $\mathcal{O}_X$-加群 $\mathcal{F}$ に対し、上の合成は $[\mathcal{F}]$ を
$$
ch(X \to Y, \mathcal{F}[0]) \cap [Y] \in \CH_*(X) \otimes \mathbf{Q}
$$
へ送る。$\mathcal{F}$ が集合論的に閉部分スキーム $Z \subset X$ 上に台をもつならば、
補題 \ref{chow-lemma-loc-chern-shrink-Z} により
$$
ch(X \to Y, \mathcal{F}[0]) = (Z \to X)_* \circ ch(Z \to Y, \mathcal{F}[0])
$$
である。この場合、得られる元は $\CH_*(Z) \to \CH_*(X)$ の像に属する。
したがって条件 (1) を得る。

\medskip\noindent
$Z \subset X$ を $\delta$-次元 $k$ の整閉部分スキームとする。同じ議論により、
上の合成は $[\mathcal{O}_Z]$ を
$$
ch(X \to Y, \mathcal{O}_Z[0]) \cap [Y] =
(Z \to X)_* ch(Z \to Y, \mathcal{O}_Z[0]) \cap [Y]
$$
へ送る。ゆえに
$ch(Z \to Y, \mathcal{O}_Z[0]) \cap [Y] \in
\CH_*(Z) \otimes \mathbf{Q}$ の次数 $k$ 成分が $[Z]$ であることを示せば十分である。
$\CH_k(Z) = \mathbf{Z}$ なので、これを示すにあたり $Y$ を $Z$ の生成点 $\xi$ の
開近傍で置き換えてよい。正則局所環 $\mathcal{O}_{X, \xi}$ の極大イデアルは
正則列で生成されるので（『可換代数』の補題 \ref{algebra-lemma-regular-ring-CM}）、
$Z$ のイデアルは正則列で生成されると仮定してよい。『因子』の補題 \ref{divisors-lemma-Noetherian-scheme-regular-ideal} を参照されたい。
したがって補題 \ref{chow-lemma-actual-computation} から結論が従う。
\end{proof}







\section{正則スキーム上の有理交叉積}
\label{chow-section-intersection-regular}

\noindent
$X$ が Noether、正則、有限次元で、対角射がアフィンならば、
$\CH_*(X) \otimes \mathbf{Q}$ に交叉積が存在することを示す。
構成の基礎は次の結果であり、これは前節の命題の系である。

\begin{lemma}
\label{chow-lemma-K-tensor-Q}
$(S, \delta)$ を状況 \ref{chow-situation-setup} のとおりとする。
$X$ を $S$ 上有限型な準コンパクト正則スキームとし、その対角射はアフィンで、
$\delta_{X/S} : X \to \mathbf{Z}$ は有界であるとする。このとき、注意
\ref{chow-remark-chern-character-K} の写像 $ch$ と写像 $c \mapsto c \cap [X]$ の合成
$$
K_0(\textit{Vect}(X)) \otimes \mathbf{Q}
\longrightarrow
A^*(X) \otimes \mathbf{Q}
\longrightarrow
\CH_*(X) \otimes \mathbf{Q}
$$
は同型である。
\end{lemma}

\begin{proof}
『スキームの導来圏』の補題 \ref{perfect-lemma-Kprime-K}、\ref{perfect-lemma-regular-resolution-property} および \ref{perfect-lemma-K-is-old-K} により
$K'_0(X) = K_0(X) = K_0(\textit{Vect}(X))$ である。
注意 \ref{chow-remark-chern-classes-agree} により、与えられた合成は $X = Y$ の場合の
命題 \ref{chow-proposition-K-tensor-Q} の写像と一致する。したがって結論は同命題から従う。
\end{proof}

\noindent
$X, S, \delta$ を補題 \ref{chow-lemma-K-tensor-Q} のとおりとする。
簡単のため、節 \ref{chow-section-cycles-codimension} のように一定の余次元のサイクルを扱う。
$[X]$ を $X$ の基本サイクルとする。注意 \ref{chow-remark-fundamental-class} を参照されたい。
$\alpha \in CH^i(X)$ および $\beta \in CH^j(X)$ を取る。
補題により
$ch(\alpha') \cap [X]  = \alpha$ を満たす唯一の
$\alpha' \in K_0(\textit{Vect}(X)) \otimes \mathbf{Q}$ が存在する。
これはもちろん、$i' \not = i$ ならば $ch_{i'}(\alpha') \cap [X] = 0$ であり、
$ch_i(\alpha') \cap [X] = \alpha$ であることを意味する。
補題 \ref{chow-lemma-adams-and-chern} により
$\alpha'' = 2^{-i}\psi^2(\alpha')$ も別の解である。一意性から
$\alpha'' = \alpha'$ を得るので、$i' \not = i$ に対して
$A^{i'}(X) \otimes \mathbf{Q}$ 内で $ch_{i'}(\alpha) = 0$ である。
そこで、上で見た $\alpha'$ の性質により、
$\CH^{i + j}(X) \otimes \mathbf{Q}$ において
$$
\alpha \cdot \beta = ch(\alpha') \cap \beta =
ch_i(\alpha') \cap \beta
$$
と定義できる。 これは対称な対である。実際、$\beta$ を持ち上げる
$\beta' \in K_0(\textit{Vect}(X)) \otimes \mathbf{Q}$ を取れば、
$$
\alpha \cdot \beta = ch(\alpha') \cap \beta =
ch(\alpha') \cap ch(\beta') \cap [X]
$$
を得るが、Chern 類どうしは可換である。同じ理由により交叉積は結合的である。すなわち
$$
(\alpha \cdot \beta) \cdot \gamma =
ch(\alpha') \cap ch(\beta') \cap ch(\gamma') \cap [X]
$$
であり、双変類の合成が結合的であることを用いた。
よりよい定式化は次のとおりであろう。次数と両立する可換かつ結合的な交叉積で、同型
$K_0(\textit{Vect}(X)) \otimes \mathbf{Q} \to \CH^*(X) \otimes \mathbf{Q}$
が環の同型となるものが、$\CH^*(X) \otimes \mathbf{Q}$ 上に一意に存在する。






\section{局所完全交叉射に対する Gysin 写像}
\label{chow-section-koszul}

\noindent
本節を読む前に、正則埋め込み
（『因子』の第 \ref{divisors-section-regular-immersions} 節）および
局所完全交叉射
（『射の続論』の第 \ref{more-morphisms-section-lci} 節）を復習することを勧める。

\medskip\noindent
$(S, \delta)$ を状況 \ref{chow-situation-setup} のとおりとする。
$i : X \to Y$ を $S$ 上局所有限型なスキームの正則埋め込み\footnote{『因子』の定義 \ref{divisors-definition-regular-immersion} を参照されたい。
局所 Noether スキームに対しては、正則埋め込み、Koszul 正則埋め込み、準正則埋め込みは
同じ概念である。『因子』の補題 \ref{divisors-lemma-regular-immersion-noetherian} を参照されたい。
本節では以後これを断らずに用いる。}とする。
とくに余法層 $\mathcal{C}_{X/Y}$ は有限局所自由である
（『因子』の補題 \ref{divisors-lemma-quasi-regular-immersion} を参照）。したがって法層
$$
\mathcal{N}_{X/Y} = \SheafHom_{\mathcal{O}_X}(\mathcal{C}_{X/Y}, \mathcal{O}_X)
$$
も有限局所自由であり、全射 $\mathcal{N}_{X/Y}^\vee \to \mathcal{C}_{X/Y}$ がある
（同型はとくに全射だからである）。節 \ref{chow-section-gysin-higher-codimension} の構成により、
標準的な双変類
$$
i^! = c(X \to Y, \mathcal{N}_{X/Y}) \in A^*(X \to Y)^\wedge
$$
を得る。この概念についていくつかの補題が必要である。

\begin{lemma}
\label{chow-lemma-composition-regular-immersion}
$(S, \delta)$ を状況 \ref{chow-situation-setup} のとおりとする。
$i : X \to Y$ および $j : Y \to Z$ を $S$ 上局所有限型なスキームの正則埋め込みとする。
このとき $j \circ i$ は正則埋め込みであり、
$(j \circ i)^! = i^! \circ j^!$ である。
\end{lemma}

\begin{proof}
第一の主張は『因子』の補題 \ref{divisors-lemma-composition-regular-immersion} である。
『因子』の補題 \ref{divisors-lemma-transitivity-conormal-quasi-regular} により短完全列
$$
0 \to
i^*(\mathcal{C}_{Y/Z}) \to
\mathcal{C}_{X/Z} \to
\mathcal{C}_{X/Y} \to 0
$$
がある。したがって、より一般的な補題 \ref{chow-lemma-gysin-composition} から結論が従う。
\end{proof}

\begin{lemma}
\label{chow-lemma-section-smooth}
$(S, \delta)$ を状況 \ref{chow-situation-setup} のとおりとする。
$p : P \to X$ を $S$ 上局所有限型なスキームの滑らかな射とし、
$s : X \to P$ をその切断とする。このとき $s$ は正則埋め込みであり、
$A^*(X)^\wedge$ において $1 = s^! \circ p^*$ である。ここで
$p^* \in A^*(P \to X)^\wedge$ は補題 \ref{chow-lemma-flat-pullback-bivariant} の双変類である。
\end{lemma}

\begin{proof}
第一の主張は『因子』の補題 \ref{divisors-lemma-section-smooth-regular-immersion} である。
仮定は任意の局所有限型基底変換 $X' \to X$ で保たれるので、任意の整閉部分スキーム
$Z \subset X$ に対して $\CH_*(X)$ 内で
$s^! \cap p^*[Z] = [Z]$ であることを示せば十分である。
$P$ を $s(Z)$ の開近傍で置き換えることにより、$P \to X$ は一定相対次元 $r$ の
滑らかな射であると仮定してよい。$\dim_\delta(Z) = n$ とする。
このとき $p^{-1}(Z)$ の各既約成分は次元 $r + n$ をもち、
$p^*[Z]$ は $[p^{-1}(Z)]_{n + r}$ で与えられる。
スキーム論的に $s(X) \cap p^{-1}(Z) = s(Z)$ である。
したがって上と同じ参照結果により、$s(X) \cap p^{-1}(Z)$ は
$\overline{p}^{-1}(Z)$ に余次元 $r$ で正則に埋め込まれた閉部分スキームである。
補題 \ref{chow-lemma-gysin-fundamental} から結論が従う。
\end{proof}

\noindent
$(S, \delta)$ を状況 \ref{chow-situation-setup} のとおりとする。$S$ 上局所有限型なスキームの
可換図式
$$
\xymatrix{
X \ar[rd]_f \ar[rr]_i & & P \ar[ld]^g \\
& Y
}
$$
で、$g$ が滑らか、$i$ が正則埋め込みであるものを考える。上の双変類 $i^!$ と、
補題 \ref{chow-lemma-flat-pullback-bivariant} の双変類
$g^* \in A^*(P \to Y)^\wedge$ を合成して
$$
f^! = i^! \circ g^* \in A^*(X \to Y)
$$
を得る。射 $f$ は局所完全交叉射である。『射の続論』の定義 \ref{more-morphisms-definition-lci} を参照されたい。逆に、局所 Noether スキームの
局所完全交叉射 $f : X \to Y$ が $g$ を滑らかな射として $f = g \circ i$ と分解されるならば、
$i$ は正則埋め込みである。$f^!$ のこの構成は射 $f$ のみに依存し、分解
$f = g \circ i$ の選択には依存しないことを示す。

\begin{lemma}
\label{chow-lemma-lci-gysin-well-defined}
$(S, \delta)$ を状況 \ref{chow-situation-setup} のとおりとする。
$f : X \to Y$ を $S$ 上局所有限型なスキームの局所完全交叉射とする。
双変類 $f^!$ は、$g$ が滑らかな射である分解 $f = g \circ i$ の選択に依存しない
（そのような分解が存在すると仮定する）。
\end{lemma}

\begin{proof}
もう一つの分解 $f = g' \circ i'$ が与えられたとする。滑らかな射
$g'' : P \times_Y P' \to Y$、埋め込み $i'' : X \to P \times_Y P'$、および分解
$f = g'' \circ i''$ を考えられる。したがって、$p$ が滑らかな射である図式
$$
\xymatrix{
& P' \ar[d]^p \ar[rd]^{g'} \\
X \ar[r]^i \ar[ru]^{i'} & P \ar[r]^g & Y
}
$$
がある場合に帰着してよい。このとき補題 \ref{chow-lemma-compose-flat-pullback} により
$(g')^* = p^* \circ g^*$ なので、$A^*(X \to P)$ において
$i^! = (i')^! \circ p^*$ を示せば十分である。可換図式
$$
\xymatrix{
& X \times_P P' \ar[d]^{\overline{p}} \ar[r]_j & P' \ar[d]^p \\
X \ar[ru]^s \ar[r]^1 & X \ar[r]^i & P
}
$$
を考える。ここで $s =(1, i')$ である。『因子』の補題 \ref{divisors-lemma-section-smooth-regular-immersion} および \ref{divisors-lemma-flat-base-change-regular-immersion} により、$s$ と $j$ は正則埋め込みであり、
$i' = j \circ s$ である。補題 \ref{chow-lemma-composition-regular-immersion} により
$(i')^! = s^! \circ j^!$ を得る。正方形は Cartesian なので、双変類 $j^!$ は
$i^!$ の $P'$ への制限（注意 \ref{chow-remark-restriction-bivariant}）である。
補題 \ref{chow-lemma-construction-gysin} を参照されたい。双変類は平坦引き戻しと可換なので
$j^! \circ p^* = \overline{p}^* \circ i^!$ である。したがって
$s^! \circ \overline{p}^* = \text{id}$ を示せば十分であり、これは補題
\ref{chow-lemma-section-smooth} で証明されている。
\end{proof}

\begin{definition}
\label{chow-definition-lci-gysin}
$(S, \delta)$ を状況 \ref{chow-situation-setup} のとおりとする。
$f : X \to Y$ を $S$ 上局所有限型なスキームの局所完全交叉射とする。
$g$ が滑らかで $i$ が埋め込みであるように $f = g \circ i$ と書けるとき、
{\it $f$ に対する Gysin 写像が存在する}という。この場合、上のように
{\it Gysin 写像} $f^! = i^! \circ g^* \in A^*(X \to Y)$ を定義する。
\end{definition}

\noindent
定義から、正則埋め込みに対しては先の構成と一致し、滑らかな射に対しては
平坦引き戻しと一致する。実際、この一致はすべての syntomic 射に対して成り立つ。

\begin{lemma}
\label{chow-lemma-lci-gysin-flat}
$(S, \delta)$ を状況 \ref{chow-situation-setup} のとおりとする。
$f : X \to Y$ を $S$ 上局所有限型なスキームの局所完全交叉射とする。
$f$ に対する Gysin 写像が存在し、$f$ が平坦ならば、$f^!$ は補題
\ref{chow-lemma-flat-pullback-bivariant} の双変類に等しい。
\end{lemma}

\begin{proof}
$i : X \to P$ が埋め込みで $g : P \to Y$ が滑らかである分解 $f = g \circ i$ を選ぶ。
局所有限型な任意の射 $Y' \to Y$ に対し、$f'$, $g'$, $i'$ の基底変換は同じ仮定を満たす。
『スキームの射』の補題 \ref{morphisms-lemma-base-change-smooth} および \ref{morphisms-lemma-base-change-syntomic} および『射の続論』の補題 \ref{more-morphisms-lemma-flat-lci} を参照されたい。したがって補題
\ref{chow-lemma-bivariant-zero} により、$Y$ が整である場合に
$f^*[Y] = i^!(g^*[Y])$ を示せばよい。$n = \dim_\delta(Y)$ とおく。
$X$ と $P$ を連結成分に分解することにより、$f$ は相対次元 $r$ の平坦射、
$g$ は相対次元 $t$ の滑らかな射であると仮定してよい。このとき
$f^*[Y] = [X]_{n + s}$ および $g^*[Y] = [P]_{n + t}$ である。
一方、$i$ は余次元 $t - s$ の正則埋め込みである。ゆえに補題
\ref{chow-lemma-gysin-fundamental} により $i^![P]_{n + t} = [X]_{n + s}$ であり、証明は完了する。
\end{proof}

\begin{lemma}
\label{chow-lemma-lci-gysin-composition}
$(S, \delta)$ を状況 \ref{chow-situation-setup} のとおりとする。
$f : X \to Y$ および $g : Y \to Z$ を $S$ 上局所有限型なスキームの局所完全交叉射とする。
$g \circ f$ と $g$ に対する Gysin 写像が存在すると仮定する。このとき $f$ に対する
Gysin 写像が存在し、$(g \circ f)^! = f^! \circ g^!$ である。
\end{lemma}

\begin{proof}
『射の続論』の補題 \ref{more-morphisms-lemma-composition-lci} により
$g \circ f$ は局所完全交叉射なので、補題の主張は意味をもつ。
$X \to P$ が $X$ から $Z$ 上滑らかなスキーム $P$ への埋め込みならば、
$X \to P \times_Z Y$ は $X$ から $Y$ 上滑らかなスキームへの埋め込みである。
これで補題の第一の主張が示される。$Y \to P'$ を $Y$ から $Z$ 上滑らかなスキーム
$P'$ への埋め込みとする。可換図式
$$
\xymatrix{
X \ar[r] \ar[d] &
P \times_Z Y \ar[r]_a \ar[ld]^p &
P \times_Z P' \ar[ld]^q \\
Y \ar[r]_b \ar[d] &
P' \ar[ld] \\
Z
}
$$
を考える。水平射は正則埋め込み、南西向きの射は滑らかであり、正方形は Cartesian である。
双変類は平坦引き戻しと可換なので $a^! \circ q^* = p^* \circ b^!$ である。
この事実と補題 \ref{chow-lemma-composition-regular-immersion},
\ref{chow-lemma-compose-flat-pullback} を合わせれば、補題の主張が従う。
細部は省略する。
\end{proof}

\begin{lemma}
\label{chow-lemma-lci-gysin-commutes}
$(S, \delta)$ を状況 \ref{chow-situation-setup} のとおりとする。両正方形が Cartesian である
$S$ 上局所有限型なスキームの可換図式
$$
\xymatrix{
X'' \ar[d] \ar[r] &
X' \ar[d] \ar[r] &
X \ar[d]^f \\
Y'' \ar[r] &
Y' \ar[r] &
Y
}
$$
を考える。$f : X \to Y$ は局所完全交叉射で、$f$ に対する Gysin 写像が存在すると仮定する。
$c \in A^*(Y'' \to Y')$ とする。$f^!$ の $Y'$ への制限を
$res(f^!) \in A^*(X' \to Y')$ と表す（注意 \ref{chow-remark-restriction-bivariant}）。
このとき $c$ と $res(f^!)$ は可換である（注意 \ref{chow-remark-bivariant-commute}）。
\end{lemma}

\begin{proof}
$g$ が滑らかで $i$ が埋め込みである分解 $f = g \circ i$ を選ぶ。
$f^! = i^! \circ g^!$ なので、$g^!$（平坦引き戻しで与えられる）と $i^!$ のそれぞれについて
補題を示せば十分である。平坦引き戻しの場合は双変類の定義の一部であり、
$i^!$ の場合は補題 \ref{chow-lemma-gysin-commutes} から直ちに従う。
\end{proof}

\begin{lemma}
\label{chow-lemma-lci-gysin-easy}
$(S, \delta)$ を状況 \ref{chow-situation-setup} のとおりとする。
$S$ 上局所有限型なスキームの Cartesian 図式
$$
\xymatrix{
X' \ar[d]_{f'} \ar[r] &
X \ar[d]^f \\
Y' \ar[r] &
Y
}
$$
を考える。次を仮定する。
\begin{enumerate}
\item $f$ は局所完全交叉射であり、$f$ に対する Gysin 写像が存在する。
\item $X$, $X'$, $Y$, $Y'$ は補題 \ref{chow-lemma-locally-equidimensional} の同値な条件を満たす。
\item $x' \in X'$ の $X$, $Y'$, $Y$ における像をそれぞれ $x$, $y'$, $y$ とすると、
$n_{x'} - n_{y'} = n_x - n_y$ である。ここで $n_{x'}$, $n_x$, $n_{y'}$, $n_y$ は同補題のとおりである。
\item 任意の生成点 $\xi \in X'$ に対して局所環
$\mathcal{O}_{Y', f'(\xi)}$ は Cohen--Macaulay である。
\end{enumerate}
このとき $f^![Y'] = [X']$ である。ここで $[Y']$ と $[X']$ は注意
\ref{chow-remark-fundamental-class} のとおりである。
\end{lemma}

\begin{proof}
$x'$ を通る既約成分の生成点を $\xi$ としたとき、$n_{x'}$ は $\delta(\xi)$ の共通値である。
さらに関数 $x' \mapsto n_{x'}$, $x \mapsto n_x$, $y' \mapsto n_{y'}$,
$y \mapsto n_y$ は局所定数である。$X'$, $X$, $Y'$, $Y$ においてこの関数が値 $n$ を取る
開閉部分スキームをそれぞれ $X'_n$, $X_n$, $Y'_n$, $Y_n$ とする。
$[X'] = \sum [X'_n]_n$ および $[Y'] = \sum [Y'_n]_n$ である。
したがって補題の証明では、$X'$ をその一つの連結成分で置き換え、$X$, $Y'$, $Y$ を
その像を含む連結成分で置き換えてよい。このとき $X'$, $X$, $Y'$, $Y$ は
$\delta$-純次元、すなわち各既約成分が同じ $\delta$-次元をもつ。
$X'$, $X$, $Y'$, $Y$ に対するこの共通値をそれぞれ $n'$, $n$, $m'$, $m$ とする。
最後の仮定は $n' - m' = n - m$ を意味する。

\medskip\noindent
$i : X \to P$ が埋め込みで $g : P \to Y$ が滑らかである分解 $f = g \circ i$ を選ぶ。
$X$ は連結なので、$i(X)$ の点における $P \to Y$ の相対次元は一定である。
したがって $P$ を $i(X)$ の開近傍で置き換えることにより、$P \to Y$ は一定相対次元をもち、
$i : X \to P$ は閉埋め込みであると仮定してよい。$g$ と $i$ の基底変換をそれぞれ
$g' : Y' \times_Y P \to Y'$ および $i' : X' \to Y' \times_Y P$ と表す。
$g^*[Y] = [P]$ および $(g')^*[Y'] = [Y' \times_Y P]$ である。
最後に、$\xi' \in X'$ が生成点ならば
$\mathcal{O}_{Y' \times_Y P, i'(\xi)}$ は Cohen--Macaulay である。
実際、局所環の写像
$\mathcal{O}_{Y', f'(\xi)} \to \mathcal{O}_{Y' \times_Y P, i'(\xi)}$
は正則ファイバーをもつ平坦写像である
（『可換代数』の第 \ref{algebra-section-smooth-overview} 節）。正則局所環は Cohen--Macaulay であり
（『可換代数』の補題 \ref{algebra-lemma-regular-ring-CM}）、仮定 (4) により
$\mathcal{O}_{Y', f'(\xi)}$ は Cohen--Macaulay なので、『可換代数』の補題 \ref{algebra-lemma-CM-goes-up} から主張が従う。よって次段落の場合に帰着する。

\medskip\noindent
$f$ が正則閉埋め込みで、$X'$, $X$, $Y'$, $Y$ がそれぞれ $\delta$-次元
$n'$, $n$, $m'$, $m$ の $\delta$-純次元であり、$m' - n' = m - n$ と仮定する。
この場合、補題 \ref{chow-lemma-gysin-easy} から直ちに結論を得る。
\end{proof}

\begin{remark}
\label{chow-remark-gysin-chern-classes}
$(S, \delta)$ を状況 \ref{chow-situation-setup} のとおりとする。
$f : X \to Y$ を $S$ 上局所有限型なスキームの局所完全交叉射とし、
$f$ に対する Gysin 写像が存在すると仮定する。このとき
$f^! \circ c_i(\mathcal{E}) = c_i(f^*\mathcal{E}) \circ f^!$
であり、Chern 指標についても同様である。補題 \ref{chow-lemma-lci-gysin-commutes} を参照されたい。
さらに $X$ と $Y$ が補題 \ref{chow-lemma-locally-equidimensional} の同値な条件を満たし、
たとえば $Y$ が Cohen--Macaulay ならば、補題 \ref{chow-lemma-lci-gysin-easy} により
$f^![Y] = [X]$ である。この場合
$f^!(c_i(\mathcal{E}) \cap [Y]) = c_i(f^*\mathcal{E}) \cap [X]$
も得られ、Chern 指標についても同様である。
\end{remark}

\begin{lemma}
\label{chow-lemma-compare-gysin-base-change}
$(S, \delta)$ を状況 \ref{chow-situation-setup} のとおりとする。
$S$ 上局所有限型なスキームの Cartesian 正方形
$$
\xymatrix{
X' \ar[d]_{f'} \ar[r]_{g'} &
X \ar[d]^f \\
Y' \ar[r]^g &
Y
}
$$
を考える。次を仮定する。
\begin{enumerate}
\item $f$ と $f'$ はともに局所完全交叉射である。
\item $f$ に対する Gysin 写像が存在する。
\end{enumerate}
このとき $\mathcal{C} = \Ker(H^{-1}((g')^*\NL_{X/Y}) \to H^{-1}(\NL_{X'/Y'}))$
は有限局所自由 $\mathcal{O}_{X'}$-加群であり、$f'$ に対する Gysin 写像が存在し、
$A^*(X' \to Y')$ において
$$
res(f^!) = c_{top}(\mathcal{C}^\vee) \circ (f')^!
$$
である。
\end{lemma}

\begin{proof}
$\mathcal{C}$ が有限局所自由であることは『代数の続論』の補題 \ref{more-algebra-lemma-base-change-lci-bis} から直ちに従う。
$h : P \to Y$ が滑らかで $i$ が埋め込みである分解 $f = h \circ i$ を選ぶ。
このとき、$h' : P' \to Y'$ および $i' : X' \to P'$ をそれぞれ基底変換として、
$f' = h' \circ i'$ と分解できる。図式は
$$
\xymatrix{
X' \ar[r] \ar[d] &
P' \ar[r] \ar[d] &
Y' \ar[d] \\
X \ar[r]^i &
P \ar[r]^h &
Y
}
$$
である。とくに $f'$ に対する Gysin 写像が存在する。『射の続論』の補題 \ref{more-morphisms-lemma-get-NL} により
$$
\NL_{X/Y} = \left( \mathcal{C}_{X/P} \to i^*\Omega_{P/Y} \right)
$$
である。ここで $\mathcal{C}_{X/P}$ は埋め込み $i$ の余法層である。
プライム付きの場合も同様である。『スキームの射』の補題 \ref{morphisms-lemma-base-change-differentials} により $\Omega_{P/Y}$ は
$\Omega_{P'/Y'}$ に引き戻されるので、
$(g')^*i^*\Omega_{P/Y} = (i')^*\Omega_{P'/Y'}$ である。また
$(g')^*\mathcal{C}_{X/P} \to \mathcal{C}_{X'/P'}$ は全射である。
『スキームの射』の補題 \ref{morphisms-lemma-conormal-functorial-flat} を参照されたい。
したがって層 $\mathcal{C}$ は有限局所自由加群の写像
$(g')^*\mathcal{C}_{X/P} \to \mathcal{C}_{X'/P'}$ の核と標準的に同型である。
$i^!$ は $\mathcal{N} = \mathcal{C}_{Z/X}^\vee$ を用いて定義され、$(i')^!$ についても同様である。
ゆえに補題 \ref{chow-lemma-gysin-excess} を適用すると、$A^*(X' \to P')$ において
$$
res(i^!) = c_{top}(\mathcal{C}^\vee) \circ (i')^!
$$
を得る。最後に $f^! = i^! \circ h^*$、$(f')^! = (i')^! \circ (h')^*$、
$(h')^* = res(h^*)$ なので結論が従う。
\end{proof}

\begin{lemma}[ブローアップ公式]
\label{chow-lemma-blow-up-formula}
$(S, \delta)$ を状況 \ref{chow-situation-setup} のとおりとする。
$i : Z \to X$ を $S$ 上局所有限型なスキームの正則閉埋め込みとし、
$b : X' \to X$ を中心 $Z$ のブローアップとする。図式は
$$
\xymatrix{
E \ar[r]_j \ar[d]_\pi & X' \ar[d]^b \\
Z \ar[r]^i & X
}
$$
である。$b$ に対する Gysin 写像が存在すると仮定する。このとき $A^*(E \to Z)$ において
$$
res(b^!) = c_{top}(\mathcal{F}^\vee) \circ \pi^*
$$
である。ここで $\mathcal{F}$ は標準的な写像
$\pi^*\mathcal{C}_{Z/X} \to \mathcal{C}_{E/X'}$ の核である。
\end{lemma}

\begin{proof}
『代数の続論』の補題 \ref{more-algebra-lemma-blowup-regular-sequence} により
射 $b$ は局所完全交叉射なので、主張は意味をもつ。$Z \to X$ は正則埋め込みであり、
したがってとくに準正則なので、$\mathcal{C}_{Z/X}$ は有限局所自由で、写像
$\text{Sym}^*(\mathcal{C}_{Z/X}) \to \mathcal{C}_{Z/X, *}$ は同型である。
『因子』の補題 \ref{divisors-lemma-quasi-regular-immersion} を参照されたい。
$E = \text{Proj}(\mathcal{C}_{Z/X, *})$ なので
$E = \mathbf{P}(\mathcal{C}_{Z/X})$ は $Z$ 上の射影空間束である。
したがって $E \to Z$ は滑らかで、とくに局所完全交叉射である。
補題 \ref{chow-lemma-compare-gysin-base-change} を適用すると
$$
res(b^!) = c_{top}(\mathcal{C}^\vee) \circ \pi^!
$$
を得る。ここで $\mathcal{C}$ は同補題のとおりである。
補題 \ref{chow-lemma-lci-gysin-flat} によりもちろん $\pi^* = \pi^!$ である。
残るのは、$\mathcal{F}$ が写像
$H^{-1}(j^*\NL_{X'/X}) \to H^{-1}(\NL_{E/Z})$ の核 $\mathcal{C}$ に等しいことを示すことである。

\medskip\noindent
$E \to Z$ は滑らかなので $H^{-1}(\NL_{E/Z}) = 0$ である。『射の続論』の補題 \ref{more-morphisms-lemma-NL-smooth} を参照されたい。したがって $\mathcal{F}$ を
$H^{-1}(j^*\NL_{X'/X})$ と同一視できることを示せば十分である。
『射の続論』の補題 \ref{more-morphisms-lemma-get-triangle-NL} および \ref{more-morphisms-lemma-NL-immersion} により完全列
$$
0 \to H^{-1}(j^*\NL_{X'/X}) \to H^{-1}(\NL_{E/X}) \to
\mathcal{C}_{E/X'} \to \ldots
$$
がある。同じ補題を $E \to Z \to X$ に適用すると同型
$\pi^*\mathcal{C}_{Z/X} = H^{-1}(\pi^*\NL_{Z/X}) \to H^{-1}(\NL_{E/X})$ を得る。
これで結論が従う。
\end{proof}

\begin{lemma}
\label{chow-lemma-lci-gysin-product-regular}
$(S, \delta)$ を状況 \ref{chow-situation-setup} のとおりとする。
$f : X \to Y$ を $S$ 上局所有限型なスキームの射とし、$X$ と $Y$ はともに
準コンパクトかつ正則で、対角射がアフィンであり、有限次元であるとする。
このとき $f$ は局所完全交叉射である。さらに $f$ に対する Gysin 写像が存在すると仮定する。
このとき、節 \ref{chow-section-intersection-regular} の交叉積について
$\CH^*(X) \otimes \mathbf{Q}$ 内で
$$
f^!(\alpha \cdot \beta) = f^!\alpha \cdot f^!\beta
$$
である。
\end{lemma}

\begin{proof}
第一の主張は『射の続論』の補題 \ref{more-morphisms-lemma-morphism-regular-schemes-is-lci} から従う。
補題 \ref{chow-lemma-lci-gysin-easy} により $f^![Y] = [X]$ である。
節 \ref{chow-section-intersection-regular} のとおり、
$\alpha', \beta' \in K_0(\textit{Vect}(X)) \otimes \mathbf{Q}$ を用いて
$\alpha = ch(\alpha') \cap [Y]$ および $\beta = ch(\beta') \cap [Y]$ と書く。
$c = ch(\alpha')$ および $c' = ch(\beta')$ とおくと、構成により
$\alpha \cdot \beta = c \cap c' \cap [Y]$ である。補題
\ref{chow-lemma-lci-gysin-commutes} により $f^!$ は $c$ と $c'$ の双方と可換である。したがって
\begin{align*}
f^!(\alpha \cdot \beta)
& =
f^!(c \cap c' \cap [Y]) \\
& =
c \cap c' \cap f^![Y] \\
& =
c \cap c' \cap [X] \\
& =
(c \cap [X]) \cdot (c' \cap [X]) \\
& =
(c \cap f^![Y]) \cdot (c' \cap f^![Y]) \\
& =
f^!(\alpha) \cdot f^!(\beta)
\end{align*}
となり、主張を得る。
\end{proof}

\begin{lemma}
\label{chow-lemma-projection-formula-regular}
$(S, \delta)$ を状況 \ref{chow-situation-setup} のとおりとする。
$f : X \to Y$ を $S$ 上局所有限型なスキームの射とし、$X$ と $Y$ はともに
準コンパクトかつ正則で、対角射がアフィンであり、有限次元であるとする。
このとき $f$ は局所完全交叉射である。さらに $f$ に対する Gysin 写像が存在し、
$f$ は固有であると仮定する。このとき、節 \ref{chow-section-intersection-regular} の交叉積について
$\CH^*(Y) \otimes \mathbf{Q}$ 内で
$$
f_*(\alpha \cdot f^!\beta) = f_*\alpha \cdot \beta
$$
である。
\end{lemma}

\begin{proof}
第一の主張は『射の続論』の補題 \ref{more-morphisms-lemma-morphism-regular-schemes-is-lci} から従う。
補題 \ref{chow-lemma-lci-gysin-easy} により $f^![Y] = [X]$ である。
節 \ref{chow-section-intersection-regular} のとおり、
$\alpha' \in K_0(\textit{Vect}(X)) \otimes \mathbf{Q}$ および
$\beta' \in K_0(\textit{Vect}(Y)) \otimes \mathbf{Q}$ を用いて
$\alpha = ch(\alpha') \cap [X]$ および $\beta = ch(\beta') \cap [Y]$ と書く。
$c = ch(\alpha')$ および $c' = ch(\beta')$ とおく。このとき
\begin{align*}
f_*(\alpha \cdot f^!\beta)
& =
f_*(c \cap f^!(c' \cap [Y]_e)) \\
& =
f_*(c \cap c' \cap f^![Y]_e) \\
& =
f_*(c \cap c' \cap [X]_d) \\
& =
f_*(c' \cap c \cap [X]_d) \\
& =
c' \cap f_*(c \cap [X]_d) \\
& =
\beta \cdot f_*(\alpha)
\end{align*}
である。第一の等号は交叉積の構成による。補題 \ref{chow-lemma-lci-gysin-commutes} により
$f^!$ は $c'$ と可換である。Chern 類が双変環の中心に属することから、
$[X]$ と $c$, $c'$ の交叉順序を入れ替えられる。$c'$ は双変類なので、$c'$ と $f_*$ を可換にできる。
最後の等号も交叉積の構成による。
\end{proof}













\section{対角射に対する Gysin 写像}
\label{chow-section-gysin-for-diagonal}

\noindent
$(S, \delta)$ を状況 \ref{chow-situation-setup} のとおりとし、$f : X \to Y$ を
$S$ 上局所有限型なスキームの滑らかな射とする。このとき対角射
$\Delta : X \longrightarrow X \times_Y X$ は正則埋め込みである。『射の続論』の補題 \ref{more-morphisms-lemma-smooth-diagonal-perfect} を参照されたい。
したがって節 \ref{chow-section-koszul} で構成した Gysin 写像
$$
\Delta^! \in A^*(X \to X \times_Y X)^\wedge
$$
がある。$X \to Y$ の相対次元が一定値 $d$ ならば
$\Delta^! \in A^d(X \to X \times_Y X)$ である。

\begin{lemma}
\label{chow-lemma-diagonal-identity}
上の状況において、$A^0(X)$ 内で $\Delta^! \circ \text{pr}_i^! = 1$ である。
\end{lemma}

\begin{proof}
射影 $\text{pr}_i : X \times_Y X \to X$ は滑らかなので、これらの射影に対しても
Gysin 写像がある。したがって補題は意味をもち、補題
\ref{chow-lemma-lci-gysin-composition} の特別な場合である。
\end{proof}

\begin{proposition}
\label{chow-proposition-compute-bivariant}
\begin{reference}
\cite[Proposition 17.4.2]{F}
\end{reference}
$(S, \delta)$ を状況 \ref{chow-situation-setup} のとおりとする。
$f : X \to Y$ および $g : Y \to Z$ を $S$ 上局所有限型なスキームの射とする。
$g$ が相対次元 $d$ の滑らかな射ならば
$A^p(X \to Y) = A^{p - d}(X \to Z)$ である。
\end{proposition}

\begin{proof}
滑らかな射は Gysin 写像が存在する局所完全交叉射であることを用いる
（節 \ref{chow-section-koszul} を参照）。とくに $g^! \in A^{-d}(Y \to Z)$ がある。
そこで $c \in A^p(X \to Y)$ を $c \circ g^! \in A^{p - d}(X \to Z)$ へ送れる。

\medskip\noindent
逆に $c' \in A^{p - d}(X \to Z)$ とする。射 $Y \to Z$ による $c'$ の制限を
$res(c')$ と表す（注意 \ref{chow-remark-restriction-bivariant}）。図式
$$
\xymatrix{
X \times_Z Y \ar[r]_{\text{pr}_2} \ar[d]_{\text{pr}_1} & Y \ar[d]^g \\
X \ar[r]^f & Z
}
$$
は Cartesian なので $res(c') \in A^{p - d}(X \times_Z Y \to Y)$ である。
$\Delta : Y \to Y \times_Z Y$ を対角射とし、射
$X \times_Z Y \to Y \times_Z Y$ による $\Delta^!$ の $X \times_Z Y$ への制限を
$res(\Delta^!)$ と表す。図式
$$
\xymatrix{
X \ar[r] \ar[d] & X \times_Z Y \ar[d] \\
Y \ar[r]^-\Delta & Y \times_Z Y
}
$$
は Cartesian なので $res(\Delta^!) \in A^d(X \to X \times_Z Y)$ である。
これら二つの制限を合成して
$$
res(\Delta^!) \circ res(c') \in A^p(X \to Y)
$$
を得る。したがって写像 $A^p(X \to Y) \to A^{p - d}(X \to Z)$ と
$A^{p - d}(X \to Z) \to A^p(X \to Y)$ を構成した。これらが互いに逆であることを示す。

\medskip\noindent
まず $c \in A^p(X \to Y)$ から始める。各正方形が Cartesian である図式
$$
\xymatrix{
X \ar[d] \ar[r] & Y \ar[d] \\
X \times_Z Y \ar[r] \ar[d]^{\text{pr}_1} &
Y \times_Z Y \ar[r]_{p_2} \ar[d]^{p_1} &
Y \ar[d]^g \\
X \ar[r]^f &
Y \ar[r]^g &
Z
}
$$
を考える。この図式の下二つの正方形から
$res(c \circ g^!) = res(c) \cap p_2^!$ を得る。ここで $res(c)$ は $p_1$ による
$c$ の制限を表す。上の正方形を見て補題 \ref{chow-lemma-lci-gysin-commutes} を用いると
$c \circ \Delta^! = res(\Delta^!) \circ res(c)$ である。したがって
\begin{align*}
res(\Delta^!) \circ res(c \circ g^!)
& =
res(\Delta^!) \circ res(c) \circ p_2^! \\
& =
c \circ \Delta^! \circ p_2^! \\
& =
c
\end{align*}
を得る。最後の等号は補題 \ref{chow-lemma-diagonal-identity} による。

\medskip\noindent
逆に $c' \in A^{p - d}(X \to Z)$ から始める。上の図式の下の長方形を見ると
$res(c') \circ g^! = \text{pr}_1^! \circ c'$ である。したがって
\begin{align*}
res(\Delta^!) \circ res(c') \circ g^!
& =
res(\Delta^!) \circ \text{pr}_1^! \circ c' \\
& =
c'
\end{align*}
を得る。最後の等号は、図式の左二つの正方形から
$\text{id} = res(\Delta^! \circ p_1^!) = res(\Delta^!) \circ \text{pr}_1^!$
が従うためである。これで証明が完了する。
\end{proof}









\section{外積}
\label{chow-section-exterior-product}

\noindent
$k$ を体とする。本節では $S = \Spec(k)$ 上で考え、
$\delta : S \to \mathbf{Z}$ は唯一の点を $0$ へ送る写像とする。
例 \ref{chow-example-field} を参照されたい。

\medskip\noindent
$k$ 上局所有限型なスキームの Cartesian 正方形
$$
\xymatrix{
X \times_k Y \ar[r] \ar[d] & Y \ar[d] \\
X \ar[r] & \Spec(k) = S
}
$$
を考える。このとき標準的な写像
$$
\times :
\CH_n(X) \otimes_{\mathbf{Z}} \CH_m(Y)
\longrightarrow
\CH_{n + m}(X \times_k Y)
$$
があり、次の規則によって一意に定まる。次元 $n$, $m$ の整閉部分スキーム
$X' \subset X$, $Y' \subset Y$ に対し、$\CH_{n + m}(X \times_k Y)$ 内で
$$
[X'] \times [Y'] = [X' \times_k Y']_{n + m}
$$
である。

\begin{lemma}
\label{chow-lemma-exterior-product-well-defined}
写像
$\times : \CH_n(X) \otimes_{\mathbf{Z}} \CH_m(Y) \to \CH_{n + m}(X \times_k Y)$
は well-defined である。
\end{lemma}

\begin{proof}
$\alpha = \sum n_i[X_i]$ および $\beta = \sum m_j[Y_j]$ とし、
$X_i \subset X$ と $Y_j \subset Y$ がそれぞれ次元 $n$, $m$ の整閉部分スキームの
局所有限族であるとする。このとき $X_i \times_k Y_j$ は $X \times_k Y$ の
次元 $n + m$ の閉部分スキームの局所有限族なので、
$$
\alpha \times \beta = \sum n_i m_j [X_i \times_k Y_j]_{n + m}
$$
を $X \times_k Y$ 上の $(n + m)$-サイクルとして実際に考えられる。
このようにして加法的写像
$\times : Z_n(X) \otimes_{\mathbf{Z}} Z_m(Y) \to Z_{n + m}(X \times_k Y)$
を得る。問題は、この手続きが有理同値と両立することを示すことである。

\medskip\noindent
$i : X' \to X$ を次元 $n$ の整閉部分スキームの包含射とする。
補題 \ref{chow-lemma-flat-pullback-bivariant} により、射 $p' : X' \to \Spec(k)$ に沿う
平坦引き戻しは $(p')^* \in A^{-n}(X' \to \Spec(k))$ である。
したがって補題 \ref{chow-lemma-push-proper-bivariant} により
$c' = i_* \circ (p')^* \in A^{-n}(X \to \Spec(k))$ である。これから写像
$$
c' \cap - : \CH_m(Y) \longrightarrow \CH_{m + n}(X \times_k Y)
$$
が得られ、次元 $m$ の任意の整閉部分スキーム $Y' \subset Y$ に対して
$[Y']$ を $[X' \times_k Y']_{n + m}$ へ送ることが直ちに分かる。
したがって構成 $([X'], [Y']) \mapsto [X' \times_k Y']_{n + m}$ は第二変数について
有理同値を経由し、well-defined な写像
$Z_n(X) \otimes_{\mathbf{Z}} \CH_m(Y) \to \CH_{n + m}(X \times_k Y)$ を与える。
対称性により他方の変数についても同様であり、結論が従う。
\end{proof}

\begin{lemma}
\label{chow-lemma-chow-cohomology-towards-point}
$k$ を体とし、$X$ を $k$ 上局所有限型なスキームとする。このとき、すべての
$p \in \mathbf{Z}$ に対して標準的な同一視
$$
A^p(X \to \Spec(k)) = \CH_{-p}(X)
$$
がある。
\end{lemma}

\begin{proof}
元 $[\Spec(k)] \in \CH_0(\Spec(k))$ を考える。$c$ を $c \cap [\Spec(k)]$ へ送ることで
写像 $A^p(X \to \Spec(k)) \to \CH_{-p}(X)$ を得る。

\medskip\noindent
逆に $\alpha \in \CH_{-p}(X)$ とする。$X' \to \Spec(k)$ および
$\alpha' \in \CH_n(X')$ が与えられたとき、$\CH_{n - p}(X \times_k X')$ において
$$
c_\alpha \cap \alpha' = \alpha \times \alpha'
$$
とおくことで $c_\alpha \in A^p(X \to \Spec(k))$ を定義する。
これが双変類であることを示すため、定義 \ref{chow-definition-cycles} のように
$\alpha = \sum n_i[X_i]$ と書く。合成
$$
\coprod X_i \xrightarrow{g} X \to \Spec(k)
$$
を考え、合成を $f : \coprod X_i \to \Spec(k)$ と表す。
$g$ は固有で、$f$ は相対次元 $-p$ の平坦射である。補題
\ref{chow-lemma-flat-pullback-bivariant} により、$f$ に沿う引き戻しは双変類
$f^* \in A^p(\coprod X_i \to \Spec(k))$ である。
$i$ 番目の成分上でサイクルを $n_i$ 倍する双変類を
$\nu \in A^0(\coprod X_i)$ と表す。このとき
$\nu \circ f^* \in A^p(\coprod X_i \to X)$ である。最後に補題
\ref{chow-lemma-push-proper-bivariant} により双変類
$$
g_* \circ \nu \circ f^*
$$
を得る。$c_\alpha$ はこの類に等しく、したがってそれ自身も双変類であることが直ちに分かる。

\medskip\noindent
証明を終えるには、二つの構成が互いに逆であることを示さなければならない。
$c_\alpha \cap [\Spec(k)] = \alpha$ なので一方は明らかである。
他方について、$c \in A^p(X \to \Spec(k))$ とし
$\alpha = c \cap [\Spec(k)]$ とおく。補題 \ref{chow-lemma-bivariant-zero} により、
$X'$ が $\Spec(k)$ 上局所有限型な整スキームである場合に
$$
c \cap [X'] = c_\alpha \cap [X']
$$
を示せば十分である。このとき $p' : X' \to \Spec(k)$ は相対次元 $\dim(X')$ の平坦射であり、
$[X'] = (p')^*[\Spec(k)]$ である。したがって双変類 $c$ と $c_\alpha$ が
$[\Spec(k)]$ 上で一致することから、$[X']$ との交叉においても一致し、証明が完了する。
\end{proof}

\begin{lemma}
\label{chow-lemma-chow-cohomology-towards-point-commutes}
$k$ を体とし、$X$ を $k$ 上局所有限型なスキームとする。
$c \in A^p(X \to \Spec(k))$ とする。$Y \to Z$ を $k$ 上局所有限型なスキームの射とし、
$c' \in A^q(Y \to Z)$ とする。このとき $A^{p + q}(X \times_k Y \to Z)$ において
$c \circ c' = c' \circ c$ である。
\end{lemma}

\begin{proof}
補題 \ref{chow-lemma-chow-cohomology-towards-point} の証明で、$c$ は固有前送り、
連結成分ごとの整数倍、および平坦引き戻しの組合せで与えられることを見た。
双変類の定義により $c'$ はこれらの各演算と可換なので、結論が従う。
いくつかの細部は省略する。
\end{proof}

\begin{remark}
\label{chow-remark-commuting-exterior}
補題 \ref{chow-lemma-chow-cohomology-towards-point} と
\ref{chow-lemma-chow-cohomology-towards-point-commutes} の帰結は次のとおりである。
$k$ を体とし、$X$ を $k$ 上局所有限型なスキーム、$\alpha \in \CH_*(X)$ とする。
$Y \to Z$ を $k$ 上局所有限型なスキームの射とし、$c' \in A^q(Y \to Z)$ とする。
このとき任意の $\beta \in \CH_*(Z)$ に対し、$\CH_*(X \times_k Y)$ 内で
$$
\alpha \times (c' \cap \beta) = c' \cap (\alpha \times \beta)
$$
である。実際、$\alpha$ に対応する双変類
$c = c_\alpha \in A^*(X \to \Spec(k))$ を取ればよい。補題
\ref{chow-lemma-chow-cohomology-towards-point} の証明を参照されたい。
\end{remark}

\begin{lemma}
\label{chow-lemma-exterior-product-associative}
外積は結合的である。より正確には、$k$ を体、$X, Y, Z$ を $k$ 上局所有限型なスキームとし、
$\alpha \in \CH_*(X)$, $\beta \in \CH_*(Y)$, $\gamma \in \CH_*(Z)$ とする。
このとき $\CH_*(X \times_k Y \times_k Z)$ 内で
$(\alpha \times \beta) \times \gamma =
\alpha \times (\beta \times \gamma)$ である。
\end{lemma}

\begin{proof}
省略する。ヒント：スキームのファイバー積の結合性を用いる。
\end{proof}







\section{交叉積}
\label{chow-section-intersection-product}

\noindent
$k$ を体とする。本節では $S = \Spec(k)$ 上で考え、
$\delta : S \to \mathbf{Z}$ は唯一の点を $0$ へ送る写像とする。
例 \ref{chow-example-field} を参照されたい。

\medskip\noindent
$X$ を $k$ 上滑らかなスキームとする。節 \ref{chow-section-gysin-for-diagonal} の双変類
$\Delta^!$ により、$X$ 上局所有限型なスキームの Chow 群に一種の交叉積を定義できる。
実際、$Y \to X$ と $Z \to X$ を局所有限型なスキームの射とする。このとき
$$
Y \times_X Z =  (Y \times_k Z) \times_{X \times_k X, \Delta} X
$$
である。したがって写像列
$$
\CH_n(Y) \otimes_\mathbf{Z} \CH_m(Z)
\xrightarrow{\times}
\CH_{n + m}(Y \times_k Z)
\xrightarrow{\Delta^!}
\CH_{n + m - *}(Y \times_X Z)
$$
を考えられる。第一の矢印は節 \ref{chow-section-exterior-product} で構成した外積であり、
第二の矢印は節 \ref{chow-section-gysin-for-diagonal} で調べた対角射の Gysin 写像である。
$X$ が次元 $d$ の純次元ならば像は $\CH_{n + m - d}(Y \times_X Z)$ に入り、
一般には、$X$ が所与の次元をもつような $X$ の開閉部分スキーム上の部分へ分解できる。
$\alpha \in \CH_*(Y)$ と $\beta \in \CH_*(Z)$ に対して
$$
\alpha \cdot \beta = \Delta^!(\alpha \times \beta)
\in \CH_*(Y \times_X Z)
$$
と表す。特別な場合 $X = Y = Z$ には乗法
$$
\CH_*(X) \times \CH_*(X) \to \CH_*(X),\quad
(\alpha, \beta) \mapsto \alpha \cdot \beta
$$
を得る。これを {\it 交叉積} という。この積は明らかに対称であり、結合性は次の補題から従う。

\begin{lemma}
\label{chow-lemma-associative}
上で定義した積は結合的である。より正確には、$k$ を体、$X$ を $k$ 上滑らかなスキーム、
$Y, Z, W$ を $X$ 上局所有限型なスキームとし、
$\alpha \in \CH_*(Y)$, $\beta \in \CH_*(Z)$, $\gamma \in \CH_*(W)$ とする。
このとき $\CH_*(Y \times_X Z \times_X W)$ 内で
$(\alpha \cdot \beta) \cdot \gamma =
\alpha \cdot (\beta \cdot \gamma)$ である。
\end{lemma}

\begin{proof}
補題 \ref{chow-lemma-exterior-product-associative} により
$\CH_*(Y \times_k Z \times_k W)$ 内で
$(\alpha \times \beta) \times \gamma =
\alpha \times (\beta \times \gamma)$ である。閉埋め込み
$$
\Delta_{12} : X \times_k X \longrightarrow X \times_k X \times_k X,
\quad (x, x') \mapsto (x, x, x')
$$
および
$$
\Delta_{23} : X \times_k X \longrightarrow X \times_k X \times_k X,
\quad (x, x') \mapsto (x, x', x')
$$
を考え、対応する双変類を $\Delta_{12}^!$ と $\Delta_{23}^!$ と表す。
$\Delta_{12}^!$ は写像 $\text{pr}_{12}$ による $\Delta^!$ の
$X \times_k X \times_k X$ への制限（注意 \ref{chow-remark-restriction-bivariant}）であり、
$\Delta_{23}^!$ は写像 $\text{pr}_{23}$ による $\Delta^!$ の
$X \times_k X \times_k X$ への制限である。
したがって $\Delta_{23}$ による $\Delta_{12}^!$ の制限は $\Delta^!$ であり、
$\Delta_{12}$ による $\Delta_{23}^!$ の制限も $\Delta^!$ である。
補題 \ref{chow-lemma-gysin-commutes} により
$$
\Delta^! \circ \Delta_{12}^! =
\Delta^! \circ \Delta_{23}^! 
$$
を得る。次の等式列から補題が従う。
\begin{align*}
(\alpha \cdot \beta) \cdot \gamma
& =
\Delta^!(\Delta^!(\alpha \times \beta) \times \gamma) \\
& =
\Delta^!(\Delta_{12}^!((\alpha \times \beta) \times \gamma)) \\
& =
\Delta^!(\Delta_{23}^!((\alpha \times \beta) \times \gamma)) \\
& =
\Delta^!(\Delta_{23}^!(\alpha \times (\beta \times \gamma)) \\
& =
\Delta^!(\alpha \times \Delta^!(\beta \times \gamma)) \\
& =
\alpha \cdot (\beta \cdot \gamma)
\end{align*}
上の議論から、第二および最後から二番目以外の等号は明らかである。等式
$\Delta_{23}^!(\alpha \times (\beta \times \gamma)) =
\alpha \times \Delta^!(\beta \times \gamma)$ は注意
\ref{chow-remark-commuting-exterior} から従う。第二の等号も同様である。
\end{proof}

\begin{lemma}
\label{chow-lemma-identify-chow-for-smooth}
$k$ を体とし、$X$ を $k$ 上滑らかで次元 $d$ の純次元スキームとする。写像
$$
A^p(X) \longrightarrow \CH_{d - p}(X),\quad
c \longmapsto c \cap [X]_d
$$
は同型である。この同型により、双変類の合成は上で定義した交叉積に移る。
\end{lemma}

\begin{proof}
$g : X \to \Spec(k)$ を構造射と表す。問題の写像は同型の合成
$$
A^p(X) \to A^{p - d}(X \to \Spec(k)) \to \CH_{d - p}(X)
$$
である。第一は命題 \ref{chow-proposition-compute-bivariant} の同型
$c \mapsto c \circ g^*$ であり、第二は補題
\ref{chow-lemma-chow-cohomology-towards-point} の同型
$c \mapsto c \cap [\Spec(k)]$ である。
補題 \ref{chow-lemma-chow-cohomology-towards-point} の証明から、第二の矢印の逆は
$\alpha \in \CH_{d - p}(X)$ を双変類 $c_\alpha$ へ送る。この類は、$k$ 上局所有限型な
$Y$ に対する $\beta \in \CH_*(Y)$ を $\CH_*(X \times_k Y)$ 内の
$\alpha \times \beta$ へ送る。命題 \ref{chow-proposition-compute-bivariant} の証明から、
第一の矢印の逆はさらに $c_\alpha$ を、局所有限型な $Y \to X$ に対する
$\beta \in \CH_*(Y)$ を
$\Delta^!(\alpha \times \beta) = \alpha \cdot \beta$ へ送る双変類に移す。
これから補題の最後の主張が従う。
\end{proof}

\begin{lemma}
\label{chow-lemma-lci-gysin-product}
$k$ を体とし、$f : X \to Y$ を $k$ 上滑らかなスキームの射とする。
このとき $f$ に対する Gysin 写像が存在し、
$f^!(\alpha \cdot \beta) = f^!\alpha \cdot f^!\beta$ である。
\end{lemma}

\begin{proof}
$X \to X \times_k Y$ は $X$ から $Y$ 上滑らかなスキームへの埋め込みである。
したがって $f$ に対する Gysin 写像が存在する（定義 \ref{chow-definition-lci-gysin}）。
公式の証明では $X$ と $Y$ を連結成分に分解してよい。したがって $X$ は
$k$ 上滑らかで次元 $d$ の純次元、$Y$ は $k$ 上滑らかで次元 $e$ の純次元と仮定してよい。
たとえば補題 \ref{chow-lemma-lci-gysin-easy} により $f^![Y]_e = [X]_d$ である。
$\alpha = c \cap [Y]_e$ および $\beta = c' \cap [Y]_e$ と書く。このとき補題
\ref{chow-lemma-identify-chow-for-smooth} により
$\alpha \cdot \beta = c \cap c' \cap [Y]_e$ である。補題
\ref{chow-lemma-lci-gysin-commutes} により $f^!$ は $c$ と $c'$ の双方と可換である。したがって
\begin{align*}
f^!(\alpha \cdot \beta)
& =
f^!(c \cap c' \cap [Y]_e) \\
& =
c \cap c' \cap f^![Y]_e \\
& =
c \cap c' \cap [X]_d \\
& =
(c \cap [X]_d) \cdot (c' \cap [X]_d) \\
& =
(c \cap f^![Y]_e) \cdot (c' \cap f^![Y]_e) \\
& =
f^!(\alpha) \cdot f^!(\beta)
\end{align*}
を得る。ここでは $X$ に対しても補題 \ref{chow-lemma-identify-chow-for-smooth} を用いた。

\medskip\noindent
別証明として、$(f \times f)^!(\alpha \times \beta) = f^!\alpha \times f^!\beta$
を示し、補題 \ref{chow-lemma-lci-gysin-composition} を用いてもよい。
\end{proof}

\begin{lemma}
\label{chow-lemma-projection-formula}
$k$ を体とし、$f : X \to Y$ を $k$ 上滑らかなスキームの固有射とする。
このとき $f$ に対する Gysin 写像が存在し、
$f_*(\alpha \cdot f^!\beta) = f_*\alpha \cdot \beta$ である。
\end{lemma}

\begin{proof}
$X \to X \times_k Y$ は $X$ から $Y$ 上滑らかなスキームへの埋め込みである。
したがって $f$ に対する Gysin 写像が存在する（定義 \ref{chow-definition-lci-gysin}）。
公式の証明では $X$ と $Y$ を連結成分に分解してよい。したがって $X$ は
$k$ 上滑らかで次元 $d$ の純次元、$Y$ は $k$ 上滑らかで次元 $e$ の純次元と仮定してよい。
たとえば補題 \ref{chow-lemma-lci-gysin-easy} により $f^![Y]_e = [X]_d$ である。
補題 \ref{chow-lemma-identify-chow-for-smooth} のように
$\alpha = c \cap [X]_d$ および $\beta = c' \cap [Y]_e$ と書く。このとき
\begin{align*}
f_*(\alpha \cdot f^!\beta)
& =
f_*(c \cap f^!(c' \cap [Y]_e)) \\
& =
f_*(c \cap c' \cap f^![Y]_e) \\
& =
f_*(c \cap c' \cap [X]_d) \\
& =
f_*(c' \cap c \cap [X]_d) \\
& =
c' \cap f_*(c \cap [X]_d) \\
& =
\beta \cdot f_*(\alpha)
\end{align*}
である。第一の等号は $X$ に対する補題 \ref{chow-lemma-identify-chow-for-smooth} による。
補題 \ref{chow-lemma-lci-gysin-commutes} により $f^!$ は $c'$ と可換である。
交叉積の可換性により、$[X]_d$ と $c$, $c'$ の交叉順序を入れ替えられる。
$c'$ は双変類なので、$c'$ と $f_*$ を可換にできる。最後の等号も同補題による。
\end{proof}

\begin{lemma}
\label{chow-lemma-intersect-properly}
$k$ を体とし、$X$ を $k$ 上滑らかな整スキームとする。
$Y, Z \subset X$ を整閉部分スキームとし、
$d = \dim(Y) + \dim(Z) - \dim(X)$ とおく。次を仮定する。
\begin{enumerate}
\item $\dim(Y \cap Z) \leq d$ である。
\item $\delta(\xi) = d$ を満たすすべての $\xi \in Y \cap Z$ に対し、
$\mathcal{O}_{Y, \xi}$ と $\mathcal{O}_{Z, \xi}$ は Cohen--Macaulay である。
\end{enumerate}
このとき $\CH_d(X)$ 内で $[Y] \cdot [Z] = [Y \cap Z]_d$ である。
\end{lemma}

\begin{proof}
$[Y] \cdot [Z] = \Delta^!([Y \times Z])$ である。ここで
$\Delta^! = c(\Delta : X \to X \times X, \mathcal{T}_{X/k})$ は
節 \ref{chow-section-gysin-higher-codimension} の高余次元 Gysin 写像であり、
$\mathcal{T}_{X/k} = \SheafHom(\Omega_{X/k}, \mathcal{O}_X)$ は階数
$\dim(X)$ の局所自由層である。スキームの等式
$$
Y \cap Z = X \times_{\Delta, (X \times X)} (Y \times Z)
$$
と $\dim(Y \times Z) = \dim(Y) + \dim(Z)$ が成り立つので、補題
\ref{chow-lemma-gysin-easy} の条件 (1), (2), (3) が満たされる。
最後に $\xi \in Y \cap Z$ ならば、ファイバーが $\mathcal{O}_{Z, \xi}$ である平坦局所準同型
$$
\mathcal{O}_{Y, \xi} \longrightarrow
\mathcal{O}_{Y \times Z, \xi}
$$
がある。したがって $\mathcal{O}_{Y, \xi}$ と $\mathcal{O}_{Z, \xi}$ がともに
Cohen--Macaulay ならば、『可換代数』の補題 \ref{algebra-lemma-CM-goes-up} により
$\mathcal{O}_{Y \times Z, \xi}$ も Cohen--Macaulay である。このようにして補題
\ref{chow-lemma-gysin-easy} のすべての仮定が満たされ、結論が従う。
\end{proof}

\begin{lemma}
\label{chow-lemma-intersect-regularly-embedded}
$k$ を体、$X$ を $k$ 上滑らかなスキームとし、$i : Y \to X$ を正則閉埋め込みとする。
$\alpha \in \CH_*(X)$ とする。$Y$ が次元 $e$ の純次元ならば、
$\CH_*(X)$ 内で $\alpha \cdot [Y]_e = i_*(i^!(\alpha))$ である。
\end{lemma}

\begin{proof}
$X$ を連結成分に分解することにより、$X$ は次元 $d$ の純次元と仮定してよい。
補題 \ref{chow-lemma-identify-chow-for-smooth} のように、$x \in A^*(X)$ を用いて
$\alpha = c \cap [X]_n$ と書く。このとき
$$
i_*(i^!(\alpha)) = i_*(i^!(c \cap [X]_n)) =
i_*(c \cap i^![X]_n) = i_*(c \cap [Y]_e) =
c \cap i_*[Y]_e = \alpha \cdot [Y]_e
$$
である。第一の等号は $c$ の選び方による。第二の等号は補題
\ref{chow-lemma-lci-gysin-commutes} による。第三の等号は
$\CH_*(Y)$ 内で $i^![X]_d = [Y]_e$ であることによる
（補題 \ref{chow-lemma-lci-gysin-easy}）。第四の等号は双変類が固有前送りと可換であるためである。
最後の等号は補題 \ref{chow-lemma-identify-chow-for-smooth} による。
\end{proof}

\begin{lemma}
\label{chow-lemma-intersection-regular-smooth}
$k$ を体とし、$X$ を $k$ 上滑らかで準コンパクトかつ対角射がアフィンなスキームとする。
このとき、本節で構成した $\CH^*(X)$ 上の交叉積は $\mathbf{Q}$ をテンソルした後、
節 \ref{chow-section-intersection-regular} で構成した交叉積と一致する。
\end{lemma}

\begin{proof}
$\alpha \in \CH^i(X)$ および $\beta \in \CH^j(X)$ とする。
節 \ref{chow-section-intersection-regular} のように、
$\alpha', \beta' \in K_0(\textit{Vect}(X)) \otimes \mathbf{Q}$ を用いて
$\alpha = ch(\alpha') \cap [X]$ および $\beta = ch(\beta') \cap [X]$ と書く。
$c = ch(\alpha')$ および $c' = ch(\beta')$ とおく。
節 \ref{chow-section-intersection-regular} の交叉積は $c \cap c' \cap [X]$ を与える。
補題 \ref{chow-lemma-identify-chow-for-smooth} により、これは $\alpha \cdot \beta$ に等しい
（より正確には、任意の $k$ 上滑らかなスキーム $X$ に対して
$A^i(X) \to \CH^i(X)$, $c \mapsto c \cap [X]$ が同型であるという一般化を用いる）。
\end{proof}








\section{Dedekind 整域上の外積}
\label{chow-section-exterior-product-dim-1}

\noindent
$S$ を Dedekind 整域のスペクトルによる開被覆をもつ局所 Noether スキームとする。
閉点 $s \in S$ に対して $\delta(s) = 0$、それ以外では $\delta(s) = 1$ とおく。
このとき $(S, \delta)$ は一般的な状況 \ref{chow-situation-setup} の特別な場合である。
例 \ref{chow-example-domain-dimension-1} を参照されたい。
$S$ は正規であり（『可換代数』の補題 \ref{algebra-lemma-characterize-Dedekind}）、
したがって正規整スキームの非交和である
（『スキームの性質』の補題 \ref{properties-lemma-normal-locally-Noetherian}）。
よって以下の議論はすべて $S$ が既約な場合に帰着する。一方、$S$ が非分離であることは許す
（たとえば $S$ は点 $0$ を二重化したアフィン直線でもよい）。

\medskip\noindent
$S$ 上局所有限型なスキームの Cartesian 正方形
$$
\xymatrix{
X \times_S Y \ar[r] \ar[d] & Y \ar[d] \\
X \ar[r] & S
}
$$
を考える。標準的な写像
$$
\times :
\CH_n(X) \otimes_{\mathbf{Z}} \CH_m(Y)
\longrightarrow
\CH_{n + m - 1}(X \times_S Y)
$$
があり、次の規則によって一意に定まると主張する。$\delta$-次元がそれぞれ $n$, $m$ の
整閉部分スキーム $X' \subset X$, $Y' \subset Y$ に対して、次のようにおく。
\begin{enumerate}
\item $X'$ または $Y'$ が $S$ の既約成分を支配するならば
$[X'] \times [Y'] = [X' \times_S Y']_{n + m - 1}$ とする。
\item $X'$ も $Y'$ も $S$ の既約成分を支配しないならば
$[X'] \times [Y'] = 0$ とする。
\end{enumerate}

\begin{lemma}
\label{chow-lemma-exterior-product-well-defined-dim-1}
写像 $\times : \CH_n(X) \otimes_{\mathbf{Z}} \CH_m(Y) \to
\CH_{n + m - 1}(X \times_S Y)$ は well-defined である。
\end{lemma}

\begin{proof}
$n$-サイクル $\alpha = \sum_{i \in I} n_i[X_i]$ と $m$-サイクル
$\beta = \sum_{j \in J} m_j[Y_j]$ を考える。ここで $X_i \subset X$ と $Y_j \subset Y$ は
それぞれ $\delta$-次元 $n$, $m$ の整閉部分スキームの局所有限族である。
$K \subset I \times J$ を、$X_i$ または $Y_j$ が $S$ の既約成分を支配するような
組 $(i, j) \in I \times J$ 全体とする。このとき
$\{X_i \times_S Y_j\}_{(i, j) \in K}$ は $X \times_S Y$ の
$\delta$-次元 $n + m - 1$ の閉部分スキームの局所有限族である。したがって
$$
\alpha \times \beta =
\sum\nolimits_{(i, j) \in K} n_i m_j [X_i \times_S Y_j]_{n + m - 1}
$$
を $X \times_S Y$ 上の $(n + m - 1)$-サイクルとして考えられる。
このようにして加法的写像
$\times : Z_n(X) \otimes_{\mathbf{Z}} Z_m(Y) \to Z_{n + m}(X \times_S Y)$
を得る。問題は、この手続きが有理同値と両立することを示すことである。

\medskip\noindent
$i : X' \to X$ を、$S$ の既約成分を支配する $\delta$-次元 $n$ の整閉部分スキームの
包含射とする。『代数の続論』の補題 \ref{more-algebra-lemma-dedekind-torsion-free-flat} により
$p' : X' \to S$ は相対次元 $n - 1$ の平坦射である。したがって補題
\ref{chow-lemma-flat-pullback-bivariant} により $p'$ に沿う平坦引き戻しは
$(p')^* \in A^{-n + 1}(X' \to S)$ であり、補題
\ref{chow-lemma-push-proper-bivariant} により
$c' = i_* \circ (p')^* \in A^{-n + 1}(X \to S)$ である。これから写像
$$
c' \cap - : \CH_m(Y) \longrightarrow \CH_{m + n - 1}(X \times_S Y)
$$
が得られ、$\delta$-次元 $m$ の任意の整閉部分スキーム $Y' \subset Y$ に対して
$[Y']$ を $[X' \times_S Y']_{n + m - 1}$ へ送る。

\medskip\noindent
$i : X' \to X$ を、合成 $X' \to X \to S$ が閉点 $s \in S$ を経由する
$\delta$-次元 $n$ の整閉部分スキームの包含射とする。$s$ は Dedekind 整域のスペクトルの
閉点なので、$s$ は法束が自明な $S$ 上の有効 Cartier 因子である。Gysin 準同型を
$c \in A^1(s \to S)$ と表す。補題 \ref{chow-lemma-gysin-bivariant} を参照されたい。
射 $p' : X' \to s$ は相対次元 $n$ の平坦射である。したがって補題
\ref{chow-lemma-flat-pullback-bivariant} により、$p'$ に沿う平坦引き戻しは
$(p')^* \in A^{-n}(X' \to S)$ である。ゆえに補題
\ref{chow-lemma-push-proper-bivariant} により
$$
c' = i_* \circ (p')^* \circ c \in A^{-n}(X \to S)
$$
である。これから写像
$$
c' \cap - : \CH_m(Y) \longrightarrow \CH_{m + n - 1}(X \times_S Y)
$$
が得られる。$\delta$-次元 $m$ の任意の整閉部分スキーム $Y' \subset Y$ に対して、
$Y'$ が $S$ の既約成分を支配するならば $[Y']$ を
$[X' \times_S Y']_{n + m - 1}$ へ送り、そうでなければ $0$ へ送る。

\medskip\noindent
以上の二段落から、構成 $([X'], [Y']) \mapsto [X' \times_S Y']_{n + m - 1}$ は
第二変数について有理同値を経由し、well-defined な写像
$Z_n(X) \otimes_{\mathbf{Z}} \CH_m(Y) \to \CH_{n + m - 1}(X \times_S Y)$ を与える。
対称性により他方の変数についても同様であり、結論が従う。
\end{proof}

\begin{lemma}
\label{chow-lemma-chow-cohomology-towards-base-dim-1}
$(S, \delta)$ を上のとおりとし、$X$ を $S$ 上局所有限型なスキームとする。
このとき、すべての $p \in \mathbf{Z}$ に対して標準的な同一視
$$
A^p(X \to S) = \CH_{1 - p}(X)
$$
がある。
\end{lemma}

\begin{proof}
元 $[S]_1 \in \CH_1(S)$ を考える。$c$ を $c \cap [S]_1$ へ送ることで写像
$A^p(X \to S) \to \CH_{1 - p}(X)$ を得る。

\medskip\noindent
逆に $\alpha \in \CH_{1 - p}(X)$ とする。$X' \to S$ と
$\alpha' \in \CH_n(X')$ が与えられたとき、$\CH_{n - p}(X \times_S X')$ において
$$
c_\alpha \cap \alpha' = \alpha \times \alpha'
$$
とおくことで $c_\alpha \in A^p(X \to S)$ を定義する。
これが双変類であることを示すため、定義 \ref{chow-definition-cycles} のように
$\alpha = \sum_{i \in I} n_i[X_i]$ と書く。とくに射
$$
g : \coprod\nolimits_{i \in I} X_i \longrightarrow X
$$
は固有である。$i \in I$ を取る。$X_i$ が $S$ の既約成分を支配するならば、
構造射 $p_i : X_i \to S$ は平坦であり
$\xi_i = p_i^* \in A^p(X_i \to S)$ である。一方、$p_i$ が閉点の包含 $s_i \to S$ の前に
$p'_i : X_i \to s_i$ を経由するならば
$\xi_i = (p'_i)^* \circ c_i \in A^p(X_i \to S)$ である。
ここで $c_i \in A^1(s_i \to S)$ は Gysin 準同型、$(p'_i)^*$ は平坦引き戻しである。
$$
A^p(\coprod\nolimits_{i \in I} X_i \to S) =
\prod\nolimits_{i \in I} A^p(X_i \to S)
$$
なので
$$
\xi = \sum n_i \xi_i \in A^p(\coprod\nolimits_{i \in I} X_i \to S)
$$
を得る。最後に $g$ は固有なので、補題 \ref{chow-lemma-push-proper-bivariant} により双変類
$$
g_* \circ \xi \in A^p(X \to S)
$$
がある。$c_\alpha$ はこの類に等しいことが直ちに確かめられ、したがってそれ自身も双変類である
（補題 \ref{chow-lemma-exterior-product-well-defined-dim-1} の証明と比較されたい）。

\medskip\noindent
証明を終えるには、二つの構成が互いに逆であることを示さなければならない。
$c_\alpha \cap [S]_1 = \alpha$ なので一方は明らかである。他方について、
$c \in A^p(X \to S)$ とし $\alpha = c \cap [S]_1$ とおく。
補題 \ref{chow-lemma-bivariant-zero} により、$X'$ が $S$ 上局所有限型な整スキームである場合に
$$
c \cap [X'] = c_\alpha \cap [X']
$$
を示せば十分である。ところが $p' : X' \to S$ が相対次元
$\dim_\delta(X') - 1$ の平坦射で $[X'] = (p')^*[S]_1$ であるか、または
$X' \to S$ が $X' \to s \to S$ と分解され
$[X'] = (p')^*(s \to S)^*[S]_1$ である。したがって双変類 $c$ と $c_\alpha$ が
$[S]_1$ 上で一致することから、平坦引き戻しおよび Gysin 写像と可換することにより
$[X']$ との交叉においても一致し、証明が完了する。
\end{proof}

\begin{lemma}
\label{chow-lemma-chow-cohomology-towards-base-dim-1-commutes}
$(S, \delta)$ を上のとおりとし、$X$ を $S$ 上局所有限型なスキームとする。
$c \in A^p(X \to S)$ とする。$Y \to Z$ を $S$ 上局所有限型なスキームの射とし、
$c' \in A^q(Y \to Z)$ とする。このとき $A^{p + q}(X \times_S Y \to X \times_S Z)$ 内で
$c \circ c' = c' \circ c$ である。
\end{lemma}

\begin{proof}
補題 \ref{chow-lemma-chow-cohomology-towards-base-dim-1} の証明で、$c$ は固有前送り、
連結成分ごとの整数倍、平坦引き戻し、および Gysin 写像の組合せで与えられることを見た。
双変類の定義により $c'$ はこれらの各演算と可換なので、結論が従う。
いくつかの細部は省略する。
\end{proof}

\begin{remark}
\label{chow-remark-commuting-exterior-dim-1}
補題 \ref{chow-lemma-chow-cohomology-towards-base-dim-1} と
\ref{chow-lemma-chow-cohomology-towards-base-dim-1-commutes} の帰結は次のとおりである。
$(S, \delta)$ を上のとおりとし、$X$ を $S$ 上局所有限型なスキーム、
$\alpha \in \CH_*(X)$ とする。$Y \to Z$ を $S$ 上局所有限型なスキームの射とし、
$c' \in A^q(Y \to Z)$ とする。このとき任意の $\beta \in \CH_*(Z)$ に対し、
$\CH_*(X \times_S Y)$ 内で
$$
\alpha \times (c' \cap \beta) = c' \cap (\alpha \times \beta)
$$
である。実際、$\alpha$ に対応する双変類
$c = c_\alpha \in A^*(X \to S)$ を取ればよい。補題
\ref{chow-lemma-chow-cohomology-towards-base-dim-1} の証明を参照されたい。
\end{remark}

\begin{lemma}
\label{chow-lemma-exterior-product-associative-dim-1}
外積は結合的である。より正確には、$(S, \delta)$ を上のとおりとし、
$X, Y, Z$ を $S$ 上局所有限型なスキーム、
$\alpha \in \CH_*(X)$, $\beta \in \CH_*(Y)$, $\gamma \in \CH_*(Z)$ とする。
このとき $\CH_*(X \times_S Y \times_S Z)$ 内で
$(\alpha \times \beta) \times \gamma =
\alpha \times (\beta \times \gamma)$ である。
\end{lemma}

\begin{proof}
省略する。ヒント：スキームのファイバー積の結合性を用いる。
\end{proof}















\section{Dedekind 整域上の交叉積}
\label{chow-section-intersection-product-dim-1}

\noindent
$S$ を Dedekind 整域のスペクトルによる開被覆をもつ局所 Noether スキームとする。
閉点 $s \in S$ に対して $\delta(s) = 0$、それ以外では $\delta(s) = 1$ とおく。
このとき $(S, \delta)$ は一般的な状況 \ref{chow-situation-setup} の特別な場合である。
例 \ref{chow-example-domain-dimension-1} および節
\ref{chow-section-exterior-product-dim-1} の議論を参照されたい。

\medskip\noindent
$X$ を $S$ 上滑らかなスキームとする。節 \ref{chow-section-gysin-for-diagonal} の双変類
$\Delta^!$ により、$X$ 上局所有限型なスキームの Chow 群に一種の交叉積を定義できる。
実際、$Y \to X$ と $Z \to X$ を局所有限型なスキームの射とする。このとき
$$
Y \times_X Z =  (Y \times_S Z) \times_{X \times_S X, \Delta} X
$$
である。したがって写像列
$$
\CH_n(Y) \otimes_\mathbf{Z} \CH_m(Y)
\xrightarrow{\times}
\CH_{n + m - 1}(Y \times_S Z)
\xrightarrow{\Delta^!}
\CH_{n + m - *}(Y \times_X Z)
$$
を考えられる。第一の矢印は節 \ref{chow-section-exterior-product-dim-1} で構成した外積であり、
第二の矢印は節 \ref{chow-section-gysin-for-diagonal} で調べた対角射の Gysin 写像である。
$X$ が次元 $d$ の純次元ならば $X \to S$ は相対次元 $d - 1$ の滑らかな射なので、
像は $\CH_{n + m - d}(Y \times_X Z)$ に入る。一般には、$X$ が所与の次元をもつような
$X$ の開閉部分スキーム上の部分へ分解できる。
$\alpha \in \CH_*(Y)$ と $\beta \in \CH_*(Z)$ に対して
$$
\alpha \cdot \beta = \Delta^!(\alpha \times \beta)
\in \CH_*(Y \times_X Z)
$$
と表す。特別な場合 $X = Y = Z$ には乗法
$$
\CH_*(X) \times \CH_*(X) \to \CH_*(X),\quad
(\alpha, \beta) \mapsto \alpha \cdot \beta
$$
を得る。これを {\it 交叉積} という。この積は明らかに対称であり、結合性は次の補題から従う。

\begin{lemma}
\label{chow-lemma-associative-dim-1}
上で定義した積は結合的である。より正確には、$(S, \delta)$ を上のとおりとし、
$X$ を $S$ 上滑らかなスキーム、$Y, Z, W$ を $X$ 上局所有限型なスキームとする。
$\alpha \in \CH_*(Y)$, $\beta \in \CH_*(Z)$, $\gamma \in \CH_*(W)$ とすれば、
$\CH_*(Y \times_X Z \times_X W)$ 内で
$(\alpha \cdot \beta) \cdot \gamma =
\alpha \cdot (\beta \cdot \gamma)$ である。
\end{lemma}

\begin{proof}
補題 \ref{chow-lemma-exterior-product-associative-dim-1} により
$\CH_*(Y \times_S Z \times_S W)$ 内で
$(\alpha \times \beta) \times \gamma =
\alpha \times (\beta \times \gamma)$ である。閉埋め込み
$$
\Delta_{12} : X \times_S X \longrightarrow X \times_S X \times_S X,
\quad (x, x') \mapsto (x, x, x')
$$
および
$$
\Delta_{23} : X \times_S X \longrightarrow X \times_S X \times_S X,
\quad (x, x') \mapsto (x, x', x')
$$
を考え、対応する双変類を $\Delta_{12}^!$ と $\Delta_{23}^!$ と表す。
$\Delta_{12}^!$ は写像 $\text{pr}_{12}$ による $\Delta^!$ の
$X \times_S X \times_S X$ への制限（注意 \ref{chow-remark-restriction-bivariant}）であり、
$\Delta_{23}^!$ は写像 $\text{pr}_{23}$ による $\Delta^!$ の
$X \times_S X \times_S X$ への制限である。
したがって $\Delta_{23}$ による $\Delta_{12}^!$ の制限は $\Delta^!$ であり、
$\Delta_{12}$ による $\Delta_{23}^!$ の制限も $\Delta^!$ である。
補題 \ref{chow-lemma-gysin-commutes} により
$$
\Delta^! \circ \Delta_{12}^! =
\Delta^! \circ \Delta_{23}^! 
$$
を得る。次の等式列から補題が従う。
\begin{align*}
(\alpha \cdot \beta) \cdot \gamma
& =
\Delta^!(\Delta^!(\alpha \times \beta) \times \gamma) \\
& =
\Delta^!(\Delta_{12}^!((\alpha \times \beta) \times \gamma)) \\
& =
\Delta^!(\Delta_{23}^!((\alpha \times \beta) \times \gamma)) \\
& =
\Delta^!(\Delta_{23}^!(\alpha \times (\beta \times \gamma)) \\
& =
\Delta^!(\alpha \times \Delta^!(\beta \times \gamma)) \\
& =
\alpha \cdot (\beta \cdot \gamma)
\end{align*}
上の議論から、第二および最後から二番目以外の等号は明らかである。等式
$\Delta_{23}^!(\alpha \times (\beta \times \gamma)) =
\alpha \times \Delta^!(\beta \times \gamma)$ は注意
\ref{chow-remark-commuting-exterior} から従う。第二の等号も同様である。
\end{proof}

\begin{lemma}
\label{chow-lemma-identify-chow-for-smooth-dim-1}
$(S, \delta)$ を上のとおりとし、$X$ を $S$ 上滑らかで次元 $d$ の純次元スキームとする。
写像
$$
A^p(X) \longrightarrow \CH_{d - p}(X),\quad
c \longmapsto c \cap [X]_d
$$
は同型である。この同型により、双変類の合成は上で定義した交叉積に移る。
\end{lemma}

\begin{proof}
$g : X \to S$ を構造射と表す。問題の写像は同型の合成
$$
A^p(X) \to A^{p - d + 1}(X \to S) \to \CH_{d - p}(X)
$$
である。第一は命題 \ref{chow-proposition-compute-bivariant} の同型
$c \mapsto c \circ g^*$ であり、第二は補題
\ref{chow-lemma-chow-cohomology-towards-base-dim-1} の同型
$c \mapsto c \cap [S]_1$ である。
補題 \ref{chow-lemma-chow-cohomology-towards-base-dim-1} の証明から、第二の矢印の逆は
$\alpha \in \CH_{d - p}(X)$ を双変類 $c_\alpha$ へ送る。この類は、$k$ 上局所有限型な
$Y$ に対する $\beta \in \CH_*(Y)$ を $\CH_*(X \times_k Y)$ 内の
$\alpha \times \beta$ へ送る。命題 \ref{chow-proposition-compute-bivariant} の証明から、
第一の矢印の逆はさらに $c_\alpha$ を、局所有限型な $Y \to X$ に対する
$\beta \in \CH_*(Y)$ を
$\Delta^!(\alpha \times \beta) = \alpha \cdot \beta$ へ送る双変類に移す。
これから補題の最後の主張が従う。
\end{proof}










\section{Todd 類}
\label{chow-section-todd-classes}

\noindent
階数 $r$ のベクトル束 $\mathcal{E}$ に付随する最後の類は、
その {\it Todd 類} $Todd(\mathcal{E})$ である。
Chern 根 $x_1, \ldots, x_r$ を用いると、これは
$$
Todd(\mathcal{E})
=
\prod\nolimits_{i = 1}^r
\frac{x_i}{1 - e^{-x_i}}
$$
と定義される。Chern 類 $c_i = c_i(\mathcal{E})$ を用いると
$$
Todd(\mathcal{E})
=
1
+
\frac{1}{2}c_1
+
\frac{1}{12}(c_1^2 + c_2)
+
\frac{1}{24}c_1c_2
+
\frac{1}{720}(-c_1^4 + 4c_1^2c_2 + 3c_2^2 + c_1c_3 - c_4)
+
\ldots
$$
を得る。分母について必要な注意は前節で述べた。完全列
$$
0
\to
{\mathcal E}_1
\to
{\mathcal E}
\to
{\mathcal E}_2
\to
0
$$
が与えられると
$$
Todd({\mathcal E}) = Todd({\mathcal E}_1) Todd({\mathcal E}_2).
$$
が成り立つ。







\section{Grothendieck--Riemann--Roch}
\label{chow-section-grr}

\noindent
$(S, \delta)$ を状況 \ref{chow-situation-setup} のとおりとする。
$X, Y$ を $S$ 上局所有限型とする。
$\mathcal{E}$ を $X$ 上の階数 $r$ の有限局所自由層
${\mathcal E}$ とする。$f : X \to Y$ を固有かつ滑らかな射とする。
$R^if_*\mathcal{E}$ が $Y$ 上の有限階数局所自由層であると仮定する。
この場合、Grothendieck--Riemann--Roch の定理は
$$
f_*(Todd(T_{X/Y}) ch(\mathcal{E}))
=
\sum (-1)^i ch(R^if_*\mathcal{E})
$$
を述べる。ここで
$$
T_{X/Y} = \SheafHom_{\mathcal{O}_X}(\Omega_{X/Y}, \mathcal{O}_X)
$$
は $Y$ 上の $X$ の相対接束である。$k$ が体で $Y = \Spec(k)$ ならば、これは
$$
\chi(X, \mathcal{E}) = \deg(Todd(T_{X/k}) ch(\mathcal{E}))
$$
と言い換えられる。この定理はより一般的であり、正しい一般性で定式化すると
証明も容易になる。これについては別の箇所で改めて扱う
（将来の参照をここに挿入する）。







\section{基底スキームの変更}
\label{chow-section-change-base}

\noindent
本節では、異なる基底スキームに対する理論を比較する方法を説明する。

\begin{situation}
\label{chow-situation-setup-base-change}
$(S, \delta)$ および $(S', \delta')$ を状況 \ref{chow-situation-setup} のとおりとする。
さらに $g : S' \to S$ をスキームの平坦射とし、$c \in \mathbf{Z}$ を次の条件を
満たす整数とする。任意の $s \in S$ と、$g^{-1}(\{s\})$ の既約成分の生成点
$s' \in S'$ に対し
$\delta'(s') = \delta(s) + c$
が成り立つ。
\end{situation}

\noindent
$S$ 上局所有限型なスキーム $X$ に対し、Chow 群の整定義された写像
$\CH_k(X) \to \CH_{k + c}(X \times_S S')$ が存在し、これは概ね本章で定義した
演算と可換であることを示す。

\begin{lemma}
\label{chow-lemma-dimension-base-change}
状況 \ref{chow-situation-setup-base-change} において $X \to S$ を局所有限型とし、
$S' \to S$ による基底変換を $X' \to S'$ と表す。
$X$ が整で $\dim_\delta(X) = k$ ならば、$X'$ の任意の既約成分 $Z'$ は
$\dim_{\delta'}(Z') = k + c$ を満たす。
\end{lemma}

\begin{proof}
射影 $X' \to X$ は平坦射 $S' \to S$ の基底変換なので平坦である
（『スキームの射』の補題 \ref{morphisms-lemma-base-change-flat}）。
したがって $X'$ の既約成分の任意の生成点 $x'$ は生成点 $x \in X$ に写る
（『スキームの射』の補題 \ref{morphisms-lemma-generalizations-lift-flat} により、
一般化は $X' \to X$ に沿って持ち上がる）。$x$ の像を $s \in S$ とする。
スキーム $S'_s = S' \times_S s$ の台位相空間は $g^{-1}(\{s\})$ と同じである
（『スキーム』の補題 \ref{schemes-lemma-fibre-topological}）。
$x'$ は、単射 $S'_s \times_s x \to S' \times_S X$ を備えたスキーム
$S'_s \times_s x$ の点とみなせる。もちろん $x'$ は
$S'_s \times_s x$ の既約成分の生成点でもある。
$\Spec(\kappa(x)) \to \Spec(\kappa(s)) = s$ の平坦性を用いて上と同様に論じると、
$x'$ は $g^{-1}(\{s\})$ の既約成分の生成点 $s'$ に写る。よって仮定から
$\delta'(s') = \delta(s) + c$ である。
『スキームの射』の補題 \ref{morphisms-lemma-dimension-fibre-after-base-change} により
$\dim_x(X_s) = \dim_{x'}(X_{s'})$ である。
$x$ は既約成分 $X_s$ の生成点であり（実際これは既約スキームだが、そのことは不要である）、
$x'$ は $X'_{s'}$ の既約成分の生成点なので、
『スキームの射』の補題 \ref{morphisms-lemma-dimension-fibre-at-a-point} により
$\text{trdeg}_{\kappa(s)}(\kappa(x)) = 
\text{trdeg}_{\kappa(s')}(\kappa(x'))$
を得る。したがって
$$
\delta_{X'/S'}(x') = \delta(s') + \text{trdeg}_{\kappa(s')}(\kappa(x')) =
\delta(s) + c + \text{trdeg}_{\kappa(s)}(\kappa(x)) = \delta_{X/S}(x) + c
$$
である。定義 \ref{chow-definition-delta-dimension} により、これが求める主張を示す。
\end{proof}

\noindent
状況 \ref{chow-situation-setup-base-change} において $X \to S$ を局所有限型とし、
$g : S' \to S$ による基底変換を $X' \to S'$ と表す。このとき一意な準同型
$$
g^* : Z_k(X) \longrightarrow Z_{k + c}(X')
$$
があり、$\delta$-次元 $k$ の整閉部分スキーム $Z \subset X$ に対して
$[Z]$ を $[Z \times_S S']_{k + c}$ に送る。これは補題
\ref{chow-lemma-dimension-base-change} により意味をもつ。

\begin{lemma}
\label{chow-lemma-pullback-coherent-base-change}
状況 \ref{chow-situation-setup-base-change} において $X \to S$ を局所有限型とし、
$S' \to S$ による基底変換を $X' \to S$ とする。
\begin{enumerate}
\item $Z \subset X$ を $\dim_\delta(Z) \leq k$ を満たす閉部分スキームとし、
その基底変換を $Z' \subset X'$ とする。このとき
$\dim_{\delta'}(Z')) \leq k + c$ であり、$Z_{k + c}(X')$ において
$[Z']_{k + c} = g^*[Z]_k$ である。
\item $\mathcal{F}$ を $\dim_\delta(\text{Supp}(\mathcal{F})) \leq k$ を満たす
$X$ 上の連接層とし、$X'$ 上の基底変換を $\mathcal{F}'$ とする。このとき
$\dim_\delta(\text{Supp}(\mathcal{F}')) \leq k + c$ であり、
$Z_{k + c}(X')$ において
$g^*[\mathcal{F}]_k = [\mathcal{F}']_{k + c}$ である。
\end{enumerate}
\end{lemma}

\begin{proof}
証明は補題 \ref{chow-lemma-pullback-coherent} の証明とまったく同じであり、
読者には読み飛ばすことを勧める。

\medskip\noindent
次元に関する主張は補題 \ref{chow-lemma-dimension-base-change} から従う。
(1) は (2)、補題 \ref{chow-lemma-cycle-closed-coherent}、および連接加群
$\mathcal{O}_Z$ の基底変換が $\mathcal{O}_{Z'}$ であることから従う。

\medskip\noindent
(2) を証明する。$X$, $X'$ は局所 Noether なので、『スキームのコホモロジー』の補題 \ref{coherent-lemma-coherent-Noetherian} により $\mathcal{F}$ は有限型である。
したがって $\mathcal{F}'$ も有限型であり
（『加群の層』の補題 \ref{modules-lemma-pullback-finite-type}）、再び
『スキームのコホモロジー』の補題 \ref{coherent-lemma-coherent-Noetherian} により
$\mathcal{F}'$ は連接である。よって補題の主張は意味をもつ。
$W \subset X$ を $\delta$-次元 $k$ の整閉部分スキームとし、
$W' \subset X'$ を $\delta'$-次元 $k + c$ の整閉部分スキームで、
$X' \to X$ により $W$ に写るものとする。$g^*[\mathcal{F}]_k$ における
$[W']$ の係数 $n$ が、$[\mathcal{F}']_{k + c}$ における $[W']$ の係数 $m$ と
一致することを示せばよい。$\xi \in W$ および $\xi' \in W'$ を生成点とする。
$A = \mathcal{O}_{X, \xi}$、$B = \mathcal{O}_{X', \xi'}$ とし、
$M = \mathcal{F}_\xi$ を $A$-加群とみなす（次元の仮定により $M$ は有限長だが、
実際にこれを確認する必要はない。補題 \ref{chow-lemma-length-finite} を参照）。
$\mathcal{F}'_{\xi'} = B \otimes_A M$ である。したがって
$$
n = \text{length}_B(B \otimes_A M)
\quad
\text{and}
\quad
m = \text{length}_A(M) \text{length}_B(B/\mathfrak m_AB)
$$
を得る。ゆえに等式は『可換代数』の補題 \ref{algebra-lemma-pullback-module} から従う。
\end{proof}

\begin{lemma}
\label{chow-lemma-pullback-base-change}
状況 \ref{chow-situation-setup-base-change} において $X \to S$ を局所有限型とし、
$S' \to S$ による基底変換を $X' \to S'$ とする。上の写像
$g^* : Z_k(X) \to Z_{k + c}(X')$ は有理同値を経由して、Chow 群の写像
$$
g^* : \CH_k(X) \longrightarrow \CH_{k + c}(X')
$$
を与える。
\end{lemma}

\begin{proof}
$\alpha \in Z_k(X)$ を零と有理同値な $k$-サイクルとする。
補題 \ref{chow-lemma-rational-equivalence-family} により、$\delta$-次元 $k$ の整閉部分スキーム
$W_i \subset X \times \mathbf{P}^1$ からなる局所有限族が存在し、その各成分は
$X \times \mathbf{P}^1$ の因子 $(X \times \mathbf{P}^1)_0$ および
$(X \times \mathbf{P}^1)_\infty$ のいずれにも含まれず、
$\alpha = \sum ([(W_i)_0]_k - [(W_i)_\infty]_k)$ を満たす。
したがって、$X \times \mathbf{P}^1$ の因子
$(X \times \mathbf{P}^1)_0$, $(X \times \mathbf{P}^1)_\infty$ のいずれにも
含まれない $\delta$-次元 $k$ の整閉部分スキーム
$W \subset X \times \mathbf{P}^1$ に対し、次を示せば十分である。
\begin{enumerate}
\item 基底変換 $W' \subset X' \times \mathbf{P}^1$ は、$k$ を $k + c$ に
置き換えた補題 \ref{chow-lemma-closed-subscheme-cross-p1} の仮定を満たす。
\item $g^*[W_0]_k = [(W')_0]_{k + c}$ および
$g^*[W_\infty]_k = [(W')_\infty]_{k + c}$ である。
\end{enumerate}
(2) は補題 \ref{chow-lemma-pullback-coherent-base-change} と、ファイバー積の結合律により
$(W')_0$ が $W_0$ の基底変換であることから直ちに従う。
(1) について、まず次元に関する主張は補題 \ref{chow-lemma-dimension-base-change} から従う。
$w' \in (W')_0$ とし、その像を $w \in W_0$、さらに $z \in \mathbf{P}^1_S$ とする。
$t \in \mathcal{O}_{\mathbf{P}^1_S, z}$ を $0 : S \to \mathbf{P}^1_S$ の通常の方程式とする。
$\mathcal{O}_{W, w} \to \mathcal{O}_{W', w'}$ は平坦であり、$W$ は整かつ
$W \not = W_0$ なので $t$ は $\mathcal{O}_{W, w}$ 上の非零因子である。
したがって $t$ は $\mathcal{O}_{W', w'}$ 上でも非零因子である。
ゆえに $W'$ は $W'_0$ 上に随伴点をもたない。
\end{proof}

\begin{lemma}
\label{chow-lemma-pullback-base-change-pullback}
状況 \ref{chow-situation-setup-base-change} において $Y \to X \to S$ を局所有限型とし、
$S' \to S$ による基底変換を $Y' \to X' \to S'$ とする。
$f : Y \to X$ が相対次元 $r$ の平坦射であると仮定する。このとき
$f' : Y' \to X'$ は相対次元 $r$ の平坦射であり、サイクル群および Chow 群の図式
$$
\vcenter{
\xymatrix{
Z_{k + r}(Y) \ar[r]_{g^*} & Z_{k + c + r}(Y') \\
Z_k(X) \ar[r]^{g^*} \ar[u]^{f^*} & Z_{k + c}(X') \ar[u]_{(f')^*}
}
}
\quad\text{and}\quad
\vcenter{
\xymatrix{
\CH_{k + r}(Y) \ar[r]_{g^*} & \CH_{k + c + r}(Y') \\
\CH_k(X) \ar[r]^{g^*} \ar[u]^{f^*} & \CH_{k + c}(X') \ar[u]_{(f')^*}
}
}
$$
は可換である。
\end{lemma}

\begin{proof}
第一の図式が可換であることを示せば十分である。
$Z \subset X$ を $\delta$-次元 $k$ の整閉部分スキームとし、
その基底変換を $Z' \subset X'$ と表す。構成により
$g^*[Z] = [Z']_{k + c}$ である。補題 \ref{chow-lemma-pullback-coherent} により
$(f')^*g^*[Z] = [Z' \times_{X'} Y']_{k + c + r}$ である。
逆に、定義 \ref{chow-definition-flat-pullback} により
$f^*[Z] = [Z \times_X Y]_{k + r}$ であり、補題
\ref{chow-lemma-pullback-coherent-base-change} により
$g^*f^*[Z] = [(Z \times_X Y)']_{k + r + c}$ である。
ファイバー積の結合律により $(Z \times_X Y)' = Z' \times_{X'} Y'$ なので結論を得る。
\end{proof}

\begin{lemma}
\label{chow-lemma-pullback-base-change-pushforward}
状況 \ref{chow-situation-setup-base-change} において $Y \to X \to S$ を局所有限型とし、
$S' \to S$ による基底変換を $Y' \to X' \to S'$ とする。
$f : Y \to X$ が固有であると仮定する。このとき $f' : Y' \to X'$ は固有であり、
サイクル群および Chow 群の図式
$$
\vcenter{
\xymatrix{
Z_k(Y) \ar[r]_{g^*} \ar[d]_{f_*} & Z_{k + c}(Y') \ar[d]^{f'_*} \\
Z_k(X) \ar[r]^{g^*} & Z_{k + c}(X')
}
}
\quad\text{and}\quad
\vcenter{
\xymatrix{
\CH_k(Y) \ar[r]_{g^*} \ar[d]_{f_*} & \CH_{k + c}(Y') \ar[d]^{f'_*} \\
\CH_k(X) \ar[r]^{g^*} & \CH_{k + c}(X')
}
}
$$
は可換である。
\end{lemma}

\begin{proof}
第一の図式が可換であることを示せば十分である。
$Z \subset Y$ を $\delta$-次元 $k$ の整閉部分スキームとし、その基底変換を
$Z' \subset X'$ と表す。構成により $g^*[Z] = [Z']_{k + c}$ である。
補題 \ref{chow-lemma-cycle-push-sheaf} により
$(f')_*g^*[Z] = [f'_*\mathcal{O}_{Z'}]_{k + c}$ である。
同じ補題により $f_*[Z] = [f_*\mathcal{O}_Z]_k$ であり、補題
\ref{chow-lemma-pullback-coherent-base-change} により
$g^*f_*[Z] = [(X' \to X)^*f_*\mathcal{O}_Z]_{k + r}$ である。
したがって
$$
(X' \to X)^*f_*\mathcal{O}_Z \cong f'_*\mathcal{O}_{Z'}
$$
が $X'$ 上の連接加群として成り立つことを示せば十分である。
$X' \to X$ は平坦であり、$\mathcal{O}_{Z'} = (Y' \to Y)^*\mathcal{O}_Z$ なので、
これは平坦基底変換、すなわち『スキームのコホモロジー』の補題 \ref{coherent-lemma-flat-base-change-cohomology} から従う。
\end{proof}

\begin{lemma}
\label{chow-lemma-pullback-base-change-c1}
状況 \ref{chow-situation-setup-base-change} において $X \to S$ を局所有限型とし、
$S' \to S$ による基底変換を $X' \to S'$ とする。
$\mathcal{L}$ を可逆 $\mathcal{O}_X$-加群とし、$X'$ 上の基底変換を
$\mathcal{L}'$ とする。このとき Chow 群の図式
$$
\xymatrix{
\CH_k(X) \ar[r]_{g^*} \ar[d]_{c_1(\mathcal{L}) \cap -} &
\CH_{k + c}(X') \ar[d]^{c_1(\mathcal{L}') \cap -} \\
\CH_{k - 1}(X) \ar[r]^{g^*} & \CH_{k + c - 1}(X')
}
$$
は可換である。
\end{lemma}

\begin{proof}
$p : L \to X$ を $\mathcal{L}$ に付随する直線束とし、零切断を
$o : X \to L$ とする。$\alpha \in CH_k(X)$ に対し、
$\beta = c_1(\mathcal{L}) \cap \alpha$ は
$o_*\alpha = - p^*\beta$ を満たす $\CH_{k - 1}(X)$ の一意な元である
（補題 \ref{chow-lemma-linebundle} および \ref{chow-lemma-linebundle-formulae}）。
同じ特徴づけは引き戻し後にも成り立つ。したがって補題は
補題 \ref{chow-lemma-pullback-base-change-pullback} および
\ref{chow-lemma-pullback-base-change-pushforward} から従う。
\end{proof}

\begin{lemma}
\label{chow-lemma-pullback-base-change-chern-classes}
状況 \ref{chow-situation-setup-base-change} において $X \to S$ を局所有限型とし、
$S' \to S$ による基底変換を $X' \to S'$ とする。
$\mathcal{E}$ を階数 $r$ の有限局所自由 $\mathcal{O}_X$-加群とし、
$X'$ 上の基底変換を $\mathcal{E}'$ とする。このとき任意の $i$ に対し、
Chow 群の図式
$$
\xymatrix{
\CH_k(X) \ar[r]_{g^*} \ar[d]_{c_i(\mathcal{E}) \cap -} &
\CH_{k + c}(X') \ar[d]^{c_i(\mathcal{E}') \cap -} \\
\CH_{k - i}(X) \ar[r]^{g^*} & \CH_{k + c - i}(X')
}
$$
は可換である。
\end{lemma}

\begin{proof}
$P = \mathbf{P}(\mathcal{E})$ とおく。$P$ の基底変換 $P'$ は
$\mathbf{P}(\mathcal{E}')$ に等しい。平坦引き戻しおよび可逆加群の $c_1$ との積が
基底変換と可換であることは既に分かっている
（補題 \ref{chow-lemma-pullback-base-change-pullback} および
\ref{chow-lemma-pullback-base-change-c1}）。したがって補題は、補題
\ref{chow-lemma-determine-intersections} で与えた $c_i(\mathcal{E})$ とのキャップ積の
特徴づけから従う。
\end{proof}

\begin{lemma}
\label{chow-lemma-compose-base-change}
$(S, \delta)$、$(S', \delta')$、$(S'', \delta'')$ を状況
\ref{chow-situation-setup} のとおりとする。$g : S' \to S$ および
$g' : S'' \to S'$ をスキームの平坦射とし、$c, c' \in \mathbf{Z}$ を、
$S, \delta, S', \delta', g, c$ および $S', \delta', S'', g', c'$ が
状況 \ref{chow-situation-setup-base-change} のとおりとなる整数とする。
$X \to S$ を局所有限型とし、$S' \to S$ および $S'' \to S$ による基底変換を
$X' \to S'$ および $X'' \to S''$ と表す。このとき次が成り立つ。
\begin{enumerate}
\item $S, \delta, S'', \delta'', g \circ g', c + c'$ は
状況 \ref{chow-situation-setup-base-change} のとおりである。
\item 写像 $g^* : Z_k(X) \to Z_{k + c}(X')$ および
$(g')^* : Z_{k + c}(X') \to Z_{k + c + c'}(X'')$ の合成は
写像 $(g \circ g')^* : Z_k(X) \to Z_{k + c + c'}(X'')$ である。
\item 補題 \ref{chow-lemma-pullback-base-change} の写像
$g^* : \CH_k(X) \to \CH_{k + c}(X')$ および
$(g')^* : \CH_{k + c}(X') \to \CH_{k + c + c'}(X'')$ の合成は、
補題 \ref{chow-lemma-pullback-base-change} の写像
$(g \circ g')^* : \CH_k(X) \to \CH_{k + c + c'}(X'')$ である。
\end{enumerate}
\end{lemma}

\begin{proof}
$s \in S$ とし、$s'' \in S''$ を $(g \circ g')^{-1}(\{s\})$ の
ある既約成分の生成点とする。$s' = g'(s'')$ とおく。
明らかに $s''$ は $(g')^{-1}(\{s'\})$ の既約成分の生成点である。
さらに $g'$ は平坦であり、したがって一般化は $g'$ に沿って持ち上がる
（『スキームの射』の補題 \ref{morphisms-lemma-base-change-flat}）ので、
$s'$ も $g^{-1}(\{s\})$ の既約成分の生成点である。
よって仮定から $\delta'(s') = \delta(s) + c$ および
$\delta''(s'') = \delta'(s') + c'$ である。したがって
$\delta''(s'') = \delta(s) + c + c'$ であり、第一の主張が成り立つ。

\medskip\noindent
第二の主張について、$Z \subset X$ を $\delta$-次元 $k$ の整閉部分スキームとし、
その基底変換を $Z' \subset X'$ および $Z'' \subset X''$ と表す。
定義により $g^*[Z] = [Z']_{k + c}$ である。補題
\ref{chow-lemma-pullback-coherent-base-change} により
$(g')^*[Z']_{k + c} = [Z'']_{k + c + c'}$ である。これで最後の主張も示された。
\end{proof}

\begin{lemma}
\label{chow-lemma-chow-limit}
状況 \ref{chow-situation-setup-base-change} において $c = 0$ と仮定する。
さらに $S' = \lim_{i \in I} S_i$ が $S$ 上アフィンなスキーム $S_i$ の
フィルター付き極限であり、次を満たすと仮定する。
\begin{enumerate}
\item $\delta_i$ を $S_i \to S \xrightarrow{\delta} \mathbf{Z}$ に等しいものとすると、
組 $(S_i, \delta_i)$ は状況 \ref{chow-situation-setup} のとおりである。
\item $S_i, \delta_i, S, \delta, S \to S_i, c = 0$ は
状況 \ref{chow-situation-setup-base-change} のとおりである。
\item $i \geq i'$ に対し
$S_i, \delta_i, S_{i'}, \delta_{i'}, S_i \to S_{i'}, c = 0$ は
状況 \ref{chow-situation-setup-base-change} のとおりである。
\end{enumerate}
このとき $X$ を $S$ 上有限型な準コンパクトスキームとし、$S' \to S$ および
$S_i \to S$ による基底変換を $X'$ および $X_i$ とすれば、
$Z_k(X') = \colim Z_k(X_i)$ および $\CH_k(X') = \colim \CH_k(X_i)$ である。
\end{lemma}

\begin{proof}
補題 \ref{chow-lemma-compose-base-change} により、サイクル群 $Z_k(X_i)$ の系、
Chow 群 $\CH_k(X_i)$ の系、および写像
$\colim Z_k(X_i) \to Z_k(X')$ と $\colim \CH_i(X_i) \to \CH_k(X')$ を得る。
$X \to S$ が経由する準コンパクト開で $S$ を置き換えられるので、
この証明に現れるすべてのスキームが Noether であると仮定してよいし、そう仮定する
（したがってそれらは準コンパクトかつ準分離でもある）。

\medskip\noindent
写像 $\colim Z_k(X_i) \to Z_k(X')$ が全射であることを示す。
$Z' \subset X'$ を $\delta'$-次元 $k$ の整閉部分スキームとする。
『スキームの極限』の補題 \ref{limits-lemma-descend-finite-presentation} により、ある $i$ と、
基底変換が $Z'$ となる有限表示の射 $Z_i \to X_i$ を取れる。
$i$ を大きくすれば、『スキームの極限』の補題 \ref{limits-lemma-descend-closed-immersion-finite-presentation} により
$Z_i$ が $X_i$ の閉部分スキームであると仮定できる。
すると $Z' \to X_i$ は $Z_i$ を経由し、$Z_i$ を $Z' \to X_i$ の
スキーム論的像で置き換えられる。このようにして $Z_i$ は $X_i$ の
整閉部分スキームであると仮定できる。補題 \ref{chow-lemma-dimension-base-change} により
$\dim_{\delta_i}(Z_i) = \dim_{\delta'}(Z') = k$ である。
したがって $Z_k(X_i) \to Z_k(X')$ は $[Z_i]$ を $[Z']$ に送り、全射性を得る。

\medskip\noindent
写像 $\colim Z_k(X_i) \to Z_k(X')$ が単射であることを示す。
$\alpha_i = \sum n_j[Z_j] \in Z_k(X_i)$ を、$Z_k(X')$ における像が零である
サイクルとする。$j \not = j'$ ならば $Z_j \not = Z_{j'}$ であり、すべての $j$ に対し
$n_j \not = 0$ であると仮定してよいし、そう仮定する。
$Z_j$ の基底変換を $Z'_j \subset X'$ と表す。
補題 \ref{chow-lemma-dimension-base-change} により $Z'_j$ の各既約成分は
$\delta'$-次元 $k$ をもつ。さらに $Z_j$ は既約で、$Z'_j \to Z_j$ は
$S' \to S$ の基底変換として平坦なので、$Z'_j \to Z_j$ は支配的である。
したがって $Z'_j$ が空でなければ、$Z'_j$ のある既約成分 $Z'$ は $Z_j$ を支配する。
すると $Z'$ は $j' \not = j$ に対して $Z'_{j'}$ の既約成分にはなりえない。
ゆえに $Z'_j$ が空でなければ、$Z'$ の係数が非零となるため
$(S' \to S_i)^*\alpha_i
= \sum [Z'_j]_r$ は非零である。
したがってすべての $j$ に対し $Z'_j = \emptyset$ である。
しかしこれは、『スキームの極限』の補題 \ref{limits-lemma-limit-nonempty} により、ある遷移写像
$S_{i'} \to S_i$ による $Z_j$ の基底変換が空であることを意味する。
よって望みどおり $\alpha_i$ は余極限で零になる。

\medskip\noindent
$\colim Z_k(X_i) \to Z_k(X')$ の全射性から
$\colim \CH_k(X_i) \to \CH_k(X')$ の全射性が従う。
証明を終えるため、この写像が単射であることを示す。
$\alpha_i \in \CH_k(X_i)$ を、その像 $\alpha' \in \CH_k(X')$ が零であるサイクルとする。
このとき $\delta"$-次元 $k + 1$ の整閉部分スキーム
$W_l' \subset X'$、$l = 1, \ldots, r$ と、$W'_l$ 上の非零有理関数 $f'_l$ が存在して
$\alpha' = \sum_{l = 1, \ldots, r} \text{div}_{W'_l}(f'_l)$ を満たす。
上と同様に論じると、ある $i$ と、基底変換が $W'_l$ となる
$\delta_i$-次元 $k + 1$ の整閉部分スキーム $W_{i, l} \subset X_i$ を取れる。
$i$ を大きくすれば、$W_{i, l}$ 上の有理関数 $f_{i, l}$ があると仮定できる。
実際、$f'_l$ を空でない開集合 $U'_l \subset W'_l$ 上の構造層の切断とみなし、
『スキームの極限』の補題 \ref{limits-lemma-descend-opens} によりこれらの開集合を降下でき、
$i$ を大きくした後、『スキームの極限』の補題 \ref{limits-lemma-descend-section} により
$f'_l$ を降下できる。場合によってはさらに $i$ を大きくした後に
$$
\alpha_i = \sum\nolimits_{l = 1, \ldots, r} \text{div}_{W_{i, l}}(f_{i, l})
$$
が成り立つと主張する。

\medskip\noindent
この主張を示すため、$Z'_{l, j} \subset W'_l$ を、$f'_l$ が
$\bigcup_j Y'_{l, j}$ の外で可逆正則関数となるような、$\delta'$-次元 $k$ の
整閉部分スキームの有限族とする。上の議論により $i$ を大きくすれば、
$\delta_i$-次元 $k$ の整閉部分スキーム $Z_{i, l, j} \subset W_i$ が存在し、
$f_{i, l}$ が $\bigcup_j Z_{i, l, j}$ の外で可逆正則関数となると仮定できる。
このとき
$$
\text{div}_{W'_l}(f'_l) = \sum n_{l, j} [Z'_{l, j}]
$$
および
$$
\text{div}_{W_{i, l}}(f_{i, l}) = \sum n_{i, l, j} [Z_{i, l, j}]
$$
と書ける。主張を示すには $n_{l, i} = n_{i, l, j}$ を示せば十分である。
実際、これは
$\beta_i =
\alpha_i - \sum\nolimits_{l = 1, \ldots, r} \text{div}_{W_{i, l}}(f_{i, l})$
が $X_i$ 上のサイクルであり、その $X'$ への引き戻しがサイクルとして零であることを含意する。
したがって、ここでは省略する容易な議論により、ある $i' \geq i$ に対して
$\beta_i$ の $X_{i'}$ への引き戻しはサイクルとして零になる。

\medskip\noindent
等式 $n_{l, i} = n_{i, l, j}$ を示すため、生成点 $\xi' \in Z'_{l, j}$ を選び、
その像を $\xi \in Z_{i, l, j}$ と表す。これも生成点である。このとき局所環の写像
$$
\mathcal{O}_{W_{i, l}, \xi}
\longrightarrow
\mathcal{O}_{W'_l, \xi'}
$$
は平坦である。実際、$W'_l \to W_{i, l}$ は平坦射 $S' \to S_i$ の基底変換である。
また、$Z_{i, l, j}$ は $Z'_{l, j}$ に引き戻されるので
$\mathfrak m_\xi \mathcal{O}_{W'_l, \xi'} = \mathfrak m_{\xi'}$ である。
したがって
$$
n_{l, j} = \text{ord}_{Z'_{l, j}}(f'_l) =
\text{ord}_{\mathcal{O}_{W'_l, \xi'}}(f'_l)
\quad\text{and}\quad
n_{i, l, j} = \text{ord}_{Z_{i, l, j}}(f_{i, l}) =
\text{ord}_{\mathcal{O}_{W_{i, l}, \xi}}(f_{i, l})
$$
の等しさは『可換代数』の補題 \ref{algebra-lemma-pullback-module} と、
『可換代数』の第 \ref{algebra-section-orders-of-vanishing} 節における
$\text{ord}$ の構成から従う。
\end{proof}












\section{付録 A：鍵となる補題への別のアプローチ}
\label{chow-section-appendix-A}

\noindent
この付録では、まず局所環 $(R, \mathfrak m, \kappa)$ 上の有限長加群 $M$ の行列式
$\det_\kappa(M)$ を定義する。小節 \ref{chow-subsection-determinants-finite-length} を参照せよ。
行列式 $\det_\kappa(M)$ は $1$-次元 $\kappa$-ベクトル空間である。
小節 \ref{chow-subsection-periodic-complexes-determinants} ではこれを用いて、
$M$ が有限長である完全な $(2, 1)$-周期複体 $(M, \varphi, \psi)$ の行列式
$\det_\kappa(M, \varphi, \psi) \in \kappa^*$ を定義する。
小節 \ref{chow-subsection-symbols} では、これらの行列式を用いて、$R$ が次元 $1$ の
Noether 環であるときの非零因子の組 $a, b \in R$ に対する tame 記号
$d_R(a, b) = \det_\kappa(R/ab, a, b)$ を構成する。
$$
d_R(a, b) = \partial_R(a, b)
$$
であることに疑いはない。ここで $\partial_R$ は節 \ref{chow-section-tame-symbol} のとおりであるが、
その確認はまだ追加していない。この付録で構成する tame 記号の利点は、例えば、
次元 $1$ の有限 $R$-加群 $M$ の単射自己準同型の組 $\varphi, \psi$ で
$\varphi(\psi(M)) = \psi(\varphi(M))$ を満たすものへ拡張できることである。
小節 \ref{chow-subsection-length-determinant} では Herbrand 商と行列式を関係づける。
主結果（命題 \ref{chow-proposition-length-determinant-periodic-complex}）の述べやすい形は、
分数体 $K$ をもつ $1$-次元 Noether 局所整域 $R$ 上で、Herbrand 商 $e_R$
（定義 \ref{chow-definition-periodic-length}）が定義される $(2, 1)$-周期複体
$(M, \varphi, \psi)$ に対する公式
$$
-e_R(M, \varphi, \psi) =
\text{ord}_R(\det\nolimits_K(M_K, \varphi, \psi))
$$
である。この命題を用いて、この付録で構成する tame 記号に対する鍵となる補題
（補題 \ref{chow-lemma-milnor-gersten-low-degree}）の別証明を与える。
補題 \ref{chow-lemma-secondary-ramification} を参照せよ。


\subsection{有限長加群の行列式}
\label{chow-subsection-determinants-finite-length}

\noindent
本節の内容は論文 \cite{determinant} および学位論文 \cite{Joe} の内容に関係する。

\medskip\noindent
$(R, \mathfrak m, \kappa)$ を局所環とし、$\varphi : M \to M$ を
有限長 $R$-加群 $M$ の $R$-線形自己準同型とする。『代数の続論』の第 \ref{more-algebra-section-determinants-finite-length} 節では、$\kappa$ に関する
$\varphi$ の行列式 $\det_\kappa(\varphi)$（およびトレースと特性多項式）を既に定義した。
本節では、$\det_\kappa(\varphi : M \to M) : \det_\kappa(M) \to \det_\kappa(M)$ が
$\det_\kappa(\varphi)$ 倍に等しくなるような標準的な $1$-次元 $\kappa$-ベクトル空間
$\det_\kappa(M)$ を構成する。$M$ が $\mathfrak m$ で零化されるならば、$M$ は
有限次元 $\kappa$-ベクトル空間とみなせ、このとき
$n = \dim_\kappa(M)$ として $\det_\kappa(M) = \wedge^n_\kappa(M)$ である。
我々の構成はこれを $R$ 上のすべての有限長加群へ一般化する。また $R$ がその剰余体を
含むならば、$\det_\kappa(M)$ は適切な意味で通常の行列式により与えられる。
注意 \ref{chow-remark-explain-determinant} を参照せよ。

\begin{definition}
\label{chow-definition-determinant}
$R$ を極大イデアル $\mathfrak m$ と剰余体 $\kappa$ をもつ局所環とし、
$M$ を有限長 $R$-加群とする。$l = \text{length}_R(M)$ とおく。
\begin{enumerate}
\item 元 $x_1, \ldots, x_r \in M$ に対し、$x_1, \ldots, x_r$ が生成する
$M$ の $R$-部分加群を
$\langle x_1, \ldots, x_r \rangle = Rx_1 + \ldots + Rx_r$ と表す。
\item $M$ の元の $l$-組 $(e_1, \ldots, e_l)$ が、
$i = 1, \ldots, l$ に対して
$\mathfrak m e_i \subset \langle e_1, \ldots, e_{i - 1} \rangle$ を満たすとき、
これを {\it admissible} であるという。
\item {\it 記号} $[e_1, \ldots, e_l]$ とは、
$(e_1, \ldots, e_l)$ が admissible な $l$-組であることを意味する。
\item 記号間の {\it admissible 関係式}とは、次のいずれかをいう。
\begin{enumerate}
\item $(e_1, \ldots, e_l)$ が admissible 列で、ある $1 \leq a \leq l$ に対し
$e_a \in \langle e_1, \ldots, e_{a - 1}\rangle$ ならば
$[e_1, \ldots, e_l] = 0$ とする。
\item $(e_1, \ldots, e_l)$ が admissible 列で、ある $1 \leq a \leq l$ に対し、
$\lambda \in R^*$ および $x \in \langle e_1, \ldots, e_{a - 1}\rangle$ を用いて
$e_a = \lambda e'_a + x$ と書けるならば
$$
[e_1, \ldots, e_l] =
\overline{\lambda} [e_1, \ldots, e_{a - 1}, e'_a, e_{a + 1}, \ldots, e_l]
$$
とする。ここで $\overline{\lambda} \in \kappa^*$ は $\lambda$ の剰余体における像である。
\item $(e_1, \ldots, e_l)$ が admissible 列で
$\mathfrak m e_a \subset \langle e_1, \ldots, e_{a - 2}\rangle$ ならば
$$
[e_1, \ldots, e_l] =
- [e_1, \ldots, e_{a - 2}, e_a, e_{a - 1}, e_{a + 1}, \ldots, e_l].
$$
とする。
\end{enumerate}
\item {\it 有限長 $R$-加群 $M$ の行列式}を
$$
\det\nolimits_\kappa(M) =
\left\{
\frac{\kappa\text{-vector space generated by symbols}}
{\kappa\text{-linear combinations of admissible relations}}
\right\}
$$
と定義する。
\end{enumerate}
\end{definition}

\noindent
常に $l = \text{length}_R(M)$ であることを強調しておく。また、記号
$[e_1, \ldots, e_l]$ が各成分について加法的であるとは限らない
（通常、実際に加法的ではない）。行列式 $\det_\kappa(M)$ が実際に次元 $1$ を
もつことを示す前に、その次元が高々 $1$ であることを示す必要がある。

\begin{lemma}
\label{chow-lemma-dimension-at-most-one}
上の記法のもとで $\dim_\kappa(\det_\kappa(M)) \leq 1$ である。
\end{lemma}

\begin{proof}
$M$ の admissible 列 $(f_1, \ldots, f_l)$ で、$i = 1, \ldots, l$ に対し
$$
\text{length}_R(\langle f_1, \ldots, f_i\rangle) = i
$$
を満たすものを一つ固定する。$M$ の長さが $l$ だからこそ、そのような
admissible 列が存在する。$\det_\kappa(M)$ の任意の元が記号
$[f_1, \ldots, f_l]$ の $\kappa$-倍であることを示す。これで補題が示される。

\medskip\noindent
$(e_1, \ldots, e_l)$ を $M$ の admissible 列とする。
$[e_1, \ldots, e_l]$ が $[f_1, \ldots, f_l]$ の倍数であることを示せば十分である。
まず $\langle e_1, \ldots, e_l\rangle \not = M$ と仮定する。
このとき、$e_i \in \langle e_1, \ldots, e_{i - 1}\rangle$ を満たす
$i \in [1, \ldots, l]$ が存在する。第一の admissible 関係式から直ちに
$\det_\kappa(M)$ において $[e_1, \ldots, e_n] = 0$ が従う。
したがって $\langle e_1, \ldots, e_l\rangle = M$ と仮定してよい。
特に $f_1 \in \langle e_1, \ldots, e_i\rangle$ を満たす最小の添字
$i \in \{1, \ldots, l\}$ が存在する。これは
$x \in \langle e_1, \ldots, e_{i - 1}\rangle$ および $\lambda \in R^*$ を用いて
$e_i = \lambda f_1 + x$ と書けることを意味する。第二の admissible 関係式により
$[e_1, \ldots, e_l] =
\overline{\lambda}[e_1, \ldots, e_{i - 1}, f_1, e_{i + 1}, \ldots, e_l]$
である。$\mathfrak m f_1 = 0$ に注意せよ。したがって第三の admissible 関係式を
$i - 1$ 回適用すると
$$
[e_1, \ldots, e_l] =
(-1)^{i - 1}\overline{\lambda}
[f_1, e_1, \ldots, e_{i - 1}, e_{i + 1}, \ldots, e_l].
$$
を得る。また
$ \langle f_1, e_1, \ldots, e_{i - 1}, e_{i + 1}, \ldots, e_l\rangle = M$
でもある。帰納的に、もとの記号が
$$
[f_1, \ldots, f_j, e_{j + 1}, \ldots, e_l]
$$
のスカラー倍に等しいと示せたと仮定する。ここで
$(f_1, \ldots, f_j, e_{j + 1}, \ldots, e_l)$ は、その元が $M$ を生成する
admissible 列、すなわち \ with
$\langle f_1, \ldots, f_j, e_{j + 1}, \ldots, e_l\rangle = M$
を満たす列である。このとき
$f_{j + 1} \in \langle f_1, \ldots, f_j, e_{j + 1}, \ldots, e_i\rangle$
を満たす最小の $i$ を取り、上と同じ手順を行えば
$$
[f_1, \ldots, f_j, e_{j + 1}, \ldots, e_l]
=
(\text{scalar}) [f_1, \ldots, f_j, f_{j + 1}, e_{j + 1},
\ldots, \hat{e_i}, \ldots, e_l]
$$
を得る。この操作を続けると所望の結果を得る。
\end{proof}

\noindent
$\det_\kappa(M)$ が常に次元 $1$ をもつことを示す前に、加群が $\kappa$ 上の
ベクトル空間である場合には通常の最高外冪と一致することを示そう。

\begin{lemma}
\label{chow-lemma-compare-det}
$R$ を極大イデアル $\mathfrak m$ と剰余体 $\kappa$ をもつ局所環とする。
$M$ を $\mathfrak m$ で零化される有限長 $R$-加群とし、
$l = \dim_\kappa(M)$ とおく。このとき写像
$$
\det\nolimits_\kappa(M) \longrightarrow \wedge^l_\kappa(M),
\quad
[e_1, \ldots, e_l] \longmapsto e_1 \wedge \ldots \wedge e_l
$$
は同型である。
\end{lemma}

\begin{proof}
通常の記号 $e_1 \wedge \ldots \wedge e_l$ はすべての admissible 関係式を満たすので、
補題に記した規則が $\kappa$-線形写像を与えることは明らかである。
また、この写像が全射であることも明らかである。補題
\ref{chow-lemma-dimension-at-most-one} により左辺は高々1次元なので、
この写像は同型である。
\end{proof}

\begin{lemma}
\label{chow-lemma-determinant-dimension-one}
$R$ を極大イデアル $\mathfrak m$ と剰余体 $\kappa$ をもつ局所環とし、
$M$ を有限長 $R$-加群とする。上で定義した行列式 $\det_\kappa(M)$ は
次元 $1$ の $\kappa$-ベクトル空間である。これは
$\langle f_1, \ldots f_l \rangle = M$ を満たす任意の admissible 列に対する記号
$[f_1, \ldots, f_l]$ によって生成される。
\end{lemma}

\begin{proof}
補題 \ref{chow-lemma-dimension-at-most-one} とその証明により、$\det_\kappa(M)$ の次元は
高々 $1$ であり、実際 $[f_1, \ldots, f_l]$ によって生成されることが分かっている。
$l = \text{length}(M)$ に関する帰納法で、これが非零であることを示す。
$l = 1$ の場合は補題 \ref{chow-lemma-compare-det} から従う。
$\mathfrak m f = 0$ を満たす非零元 $f \in M$ を選ぶ。
$\overline{M} = M /\langle f \rangle$ とおき、商写像を
$x \mapsto \overline{x}$ と表す。全射
$$
\psi : \det\nolimits_k(M) \to \det\nolimits_\kappa(\overline{M})
$$
を定義する。帰納法の仮定により $\overline{M}$ の行列式は非零なので、
これにより補題が示される。

\medskip\noindent
記号上の $\psi$ を次のように定義する。
$(e_1, \ldots, e_l)$ を admissible 列とする。
$f \not \in \langle e_1, \ldots, e_l \rangle$ ならば、単に
$\psi([e_1, \ldots, e_l]) = 0$ とおく。
$f \in \langle e_1, \ldots, e_l \rangle$ ならば、
$f \in \langle e_1, \ldots, e_i \rangle$ を満たす最小の $i$ を選ぶ。
ある単元 $\lambda \in R$ と $x \in \langle e_1, \ldots, e_{i - 1} \rangle$ を用いて
$e_i = \lambda f + x$ と書ける。この場合
$$
\psi([e_1, \ldots, e_l]) =
(-1)^i
\overline{\lambda}[\overline{e}_1, \ldots,
\overline{e}_{i - 1},
\overline{e}_{i + 1}, \ldots, \overline{e}_l].
$$
とおく。実際
$(\overline{e}_1, \ldots,
\overline{e}_{i - 1},
\overline{e}_{i + 1}, \ldots, \overline{e}_l)$
は $\overline{M}$ の admissible 列なので、これは意味をもつ。
この規則を記号の線形結合へ $\kappa$-線形に拡張すると、実際に行列式上の写像が
得られることを示そう。そのためには admissible 関係式が零に写ることを示す必要がある。

\medskip\noindent
型 (a) の関係式。$(e_1, \ldots, e_l)$ を admissible 列とし、ある
$1 \leq a \leq l$ に対し
$e_a \in \langle e_1, \ldots, e_{a - 1}\rangle$ と仮定する。
$f \in \langle e_1, \ldots, e_i\rangle$ を満たす $i$ が最小であると仮定する。
このとき $i \not = a$ であり、$i < a$ ならば
$\overline{e}_a \in \langle \overline{e}_1, \ldots,
\hat{\overline{e}_i}, \ldots, \overline{e}_{a - 1}\rangle$、
$i > a$ ならば
$\overline{e}_a \in \langle \overline{e}_1, \ldots,
\overline{e}_{a - 1}\rangle$ である。
したがって $\det_\kappa(\overline{M})$ に対する同じ admissible 関係式により、記号
$[\overline{e}_1, \ldots,
\overline{e}_{i - 1},
\overline{e}_{i + 1}, \ldots, \overline{e}_l]$
は望みどおり零になる。

\medskip\noindent
型 (b) の関係式。$(e_1, \ldots, e_l)$ を admissible 列とし、ある
$1 \leq a \leq l$ に対し、$\lambda \in R^*$ および
$x \in \langle e_1, \ldots, e_{a - 1}\rangle$ を用いて
$e_a = \lambda e'_a + x$ と書けると仮定する。
$f \in \langle e_1, \ldots, e_i\rangle$ を満たす $i$ が最小であるとし、
$y \in \langle e_1, \ldots, e_{i - 1}\rangle$ を用いて
$e_i = \mu f + y$ と書く。$i < a$ ならば、示すべき等式は
$$
(-1)^i
\overline{\lambda}
[\overline{e}_1,
\ldots,
\overline{e}_{i - 1},
\overline{e}_{i + 1},
\ldots,
\overline{e}_l]
=
(-1)^i
\overline{\lambda}
[\overline{e}_1,
\ldots,
\overline{e}_{i - 1},
\overline{e}_{i + 1},
\ldots,
\overline{e}_{a - 1},
\overline{e}'_a,
\overline{e}_{a + 1},
\ldots,
\overline{e}_l]
$$
であり、これは $\overline{e}_a = \lambda \overline{e}'_a + \overline{x}$ と
$\det_\kappa(\overline{M})$ に対する対応する admissible 関係式から従う。
$i > a$ ならば、示すべき等式は
$$
(-1)^i
\overline{\lambda}
[\overline{e}_1,
\ldots,
\overline{e}_{i - 1},
\overline{e}_{i + 1},
\ldots,
\overline{e}_l]
=
(-1)^i
\overline{\lambda}
[\overline{e}_1,
\ldots,
\overline{e}_{a - 1},
\overline{e}'_a,
\overline{e}_{a + 1},
\ldots,
\overline{e}_{i - 1},
\overline{e}_{i + 1},
\ldots,
\overline{e}_l]
$$
であり、これも $\overline{e}_a = \lambda \overline{e}'_a + \overline{x}$ と
$\det_\kappa(\overline{M})$ に対する対応する admissible 関係式から従う。
興味深いのは $i = a$ の場合である。このとき
$e_a = \lambda e'_a + x = \mu f + y$ であり、したがって
$e'_a = \lambda^{-1}(\mu f + y - x)$ でもある。ゆえに
$$
\psi([e_1, \ldots, e_l])
= (-1)^i \overline{\mu}
[\overline{e}_1, \ldots,
\overline{e}_{i - 1},
\overline{e}_{i + 1}, \ldots, \overline{e}_l]
=
\psi(
\overline{\lambda}
[e_1, \ldots, e_{a - 1}, e'_a, e_{a + 1}, \ldots, e_l]
)
$$
を得て、望みどおりである。

\medskip\noindent
型 (c) の関係式。$(e_1, \ldots, e_l)$ を admissible 列とし、
$\mathfrak m e_a \subset \langle e_1, \ldots, e_{a - 2}\rangle$ と仮定する。
$f \in \langle e_1, \ldots, e_i\rangle$ を満たす $i$ が最小であると仮定し、
$x \in \langle e_1, \ldots, e_{i - 1}\rangle$ を用いて
$e_i = \lambda f + x$ と書く。$4$ つの場合に分ける。

\medskip\noindent
場合 1：$i < a - 1$。示すべき等式は
\begin{align*}
& (-1)^i
\overline{\lambda}
[\overline{e}_1,
\ldots,
\overline{e}_{i - 1},
\overline{e}_{i + 1},
\ldots,
\overline{e}_l] \\
& =
(-1)^{i + 1}
\overline{\lambda}
[\overline{e}_1,
\ldots,
\overline{e}_{i - 1},
\overline{e}_{i + 1},
\ldots,
\overline{e}_{a - 2},
\overline{e}_a,
\overline{e}_{a - 1},
\overline{e}_{a + 1},
\ldots,
\overline{e}_l]
\end{align*}
であり、$\det_\kappa(\overline{M})$ に対する型 (c) の admissible 関係式から従う。

\medskip\noindent
場合 2：$i > a$。示すべき等式は
\begin{align*}
& (-1)^i
\overline{\lambda}
[\overline{e}_1,
\ldots,
\overline{e}_{i - 1},
\overline{e}_{i + 1},
\ldots,
\overline{e}_l] \\
& =
(-1)^{i + 1}
\overline{\lambda}
[\overline{e}_1,
\ldots,
\overline{e}_{a - 2},
\overline{e}_a,
\overline{e}_{a - 1},
\overline{e}_{a + 1},
\ldots,
\overline{e}_{i - 1},
\overline{e}_{i + 1},
\ldots,
\overline{e}_l]
\end{align*}
であり、$\det_\kappa(\overline{M})$ に対する型 (c) の admissible 関係式から従う。

\medskip\noindent
場合 3：$i = a$。$y \in \langle e_1, \ldots, e_{a - 2}\rangle$ を用いて
$e_a = \lambda f + \mu e_{a - 1} + y$ と書く。このとき定義により
$$
\psi([e_1, \ldots, e_l]) =
(-1)^a
\overline{\lambda}
[\overline{e}_1,
\ldots,
\overline{e}_{a - 1},
\overline{e}_{a + 1},
\ldots,
\overline{e}_l]
$$
である。$\overline{\mu}$ が非零ならば
$e_{a - 1} = - \mu^{-1} \lambda f + \mu^{-1}e_a - \mu^{-1} y$ であり、定義により
$$
\psi(-[e_1, \ldots, e_{a - 2}, e_a, e_{a - 1}, e_{a + 1}, \ldots, e_l]) =
(-1)^a
\overline{\mu^{-1}\lambda}
[\overline{e}_1,
\ldots,
\overline{e}_{a - 2},
\overline{e}_a,
\overline{e}_{a + 1},
\ldots,
\overline{e}_l]
$$
を得る。$\overline{M}$ において
$\overline{e}_a = \mu \overline{e}_{a - 1} + \overline{y}$ なので、この二つの結果は
$\det_\kappa(\overline{M})$ に対する関係式 (a) により等しい。
一方、$\overline{\mu}$ が零ならば、$y \in \langle e_1, \ldots, e_{a - 2}\rangle$ を用いて
$e_a = \lambda f + y$ と書け、このとき
$$
\psi(-[e_1, \ldots, e_{a - 2}, e_a, e_{a - 1}, e_{a + 1}, \ldots, e_l]) =
(-1)^a
\overline{\lambda}
[\overline{e}_1,
\ldots,
\overline{e}_{a - 1},
\overline{e}_{a + 1},
\ldots,
\overline{e}_l]
$$
であり、これは $\psi([e_1, \ldots, e_l])$ に等しい。

\medskip\noindent
場合 4：$i = a - 1$。ここでは定義により
$$
\psi([e_1, \ldots, e_l]) =
(-1)^{a - 1}
\overline{\lambda}
[\overline{e}_1,
\ldots,
\overline{e}_{a - 2},
\overline{e}_a,
\ldots,
\overline{e}_l]
$$
である。$f \not \in \langle e_1, \ldots, e_{a - 2}, e_a \rangle$ ならば
$$
\psi(-[e_1, \ldots, e_{a - 2}, e_a, e_{a - 1}, e_{a + 1}, \ldots, e_l]) =
(-1)^{a + 1}\overline{\lambda}
[\overline{e}_1,
\ldots,
\overline{e}_{a - 2},
\overline{e}_a,
\ldots,
\overline{e}_l]
$$
である。$(-1)^{a - 1} = (-1)^{a + 1}$ なので二つの式は同じである。
最後に $f \in \langle e_1, \ldots, e_{a - 2}, e_a \rangle$ と仮定する。
この場合、$x \in \langle e_1, \ldots, e_{a - 2}\rangle$ および
$y \in \langle e_1, \ldots, e_{a - 2}\rangle$ と単元 $\lambda, \mu \in R$ を用いて
$e_{a - 1} = \lambda f + x$ および $e_a = \mu f + y$ と書ける。
したがって
$e_a \in \langle e_1, \ldots, e_{a - 1} \rangle$ および
$e_{a - 1} \in \langle e_1, \ldots, e_{a - 2}, e_a\rangle$
の両方が成り立つ。この場合、型 (a) の関係式は
$[e_1, \ldots, e_l]$ および
$[e_1, \ldots, e_{a - 2}, e_a, e_{a - 1}, e_{a + 1}, \ldots, e_l]$
の両方に適用できる。上で示したこれらと $\psi$ の両立性により
$$
\psi([e_1, \ldots, e_l])
\quad\text{and}\quad
\psi([e_1, \ldots, e_{a - 2}, e_a, e_{a - 1}, e_{a + 1}, \ldots, e_l])
$$
はいずれも望みどおり零である。

\medskip\noindent
以上で $\psi$ が整定義されることを示した。残るのは全射性だけである。
$(\overline{f}_2, \ldots, \overline{f}_l)$ を $\overline{M}$ の admissible 列とする。
持ち上げ $f_2, \ldots, f_l \in M$ を選べば、
$(f, f_2, \ldots, f_l)$ は $M$ の admissible 列である。
$\psi([f, f_2, \ldots, f_l]) = [f_2, \ldots, f_l]$ なので結論を得る。
\end{proof}

\noindent
$R$ を極大イデアル $\mathfrak m$ と剰余体 $\kappa$ をもつ局所環とする。
$\varphi : M \to N$ が有限長 $R$-加群の同型ならば、任意の $M$ の記号
$[e_1, \ldots, e_l]$ に対する規則
$$
\det\nolimits_\kappa(\varphi) :
\det\nolimits_\kappa(M)
\to
\det\nolimits_\kappa(N)
$$
および
$$
\det\nolimits_\kappa(\varphi)([e_1, \ldots, e_l])
=
[\varphi(e_1), \ldots, \varphi(e_l)]
$$
により同型を得ることに注意せよ。したがって $\det\nolimits_\kappa$ は関手
\begin{equation}
\label{chow-equation-functor}
\left\{
\begin{matrix}
\text{finite length }R\text{-modules}\\
\text{with isomorphisms}
\end{matrix}
\right\}
\longrightarrow
\left\{
\begin{matrix}
1\text{-dimensional }\kappa\text{-vector spaces}\\
\text{with isomorphisms}
\end{matrix}
\right\}
\end{equation}
である。これは「行列式関手」に典型的な性質であり（\cite{Knudsen} を参照）、
次の加法性も同様である。

\begin{lemma}
\label{chow-lemma-det-exact-sequences}
$(R, \mathfrak m, \kappa)$ を局所環とする。有限長 $R$-加群の任意の短完全列
$$
0 \to K \to L \to M \to 0
$$
に対し、非零記号上の規則
$$
\gamma_{K \to L \to M} :
\det\nolimits_\kappa(K) \otimes_\kappa \det\nolimits_\kappa(M)
\longrightarrow
\det\nolimits_\kappa(L)
$$
および
$$
[e_1, \ldots, e_k]
\otimes
[\overline{f}_1, \ldots, \overline{f}_m]
\longrightarrow
[e_1, \ldots, e_k, f_1, \ldots, f_m]
$$
で定義される標準的同型が存在し、次の性質をもつ。
\begin{enumerate}
\item 短完全列の任意の同型、すなわち短完全な行と同型 $u, v, w$ をもつ任意の可換図式
$$
\xymatrix{
0 \ar[r] &
K \ar[r] \ar[d]^u &
L \ar[r] \ar[d]^v &
M \ar[r] \ar[d]^w &
0 \\
0 \ar[r] &
K' \ar[r] &
L' \ar[r] &
M' \ar[r] &
0
}
$$
に対し
$$
\gamma_{K' \to L' \to M'} \circ
(\det\nolimits_\kappa(u) \otimes \det\nolimits_\kappa(w))
=
\det\nolimits_\kappa(v) \circ
\gamma_{K \to L \to M},
$$
である。
\item 行と列が完全な有限長 $R$-加群の任意の可換正方形
$$
\xymatrix{
& 0 \ar[d] & 0 \ar[d] & 0 \ar[d] & \\
0 \ar[r] & A \ar[r] \ar[d] & B \ar[r] \ar[d] & C \ar[r] \ar[d] & 0 \\
0 \ar[r] & D \ar[r] \ar[d] & E \ar[r] \ar[d] & F \ar[r] \ar[d] & 0 \\
0 \ar[r] & G \ar[r] \ar[d] & H \ar[r] \ar[d] & I \ar[r] \ar[d] & 0 \\
& 0  & 0  & 0  &
}
$$
に対し、次の図式は可換である。
$$
\xymatrix{
\det\nolimits_\kappa(A) \otimes
\det\nolimits_\kappa(C) \otimes
\det\nolimits_\kappa(G) \otimes
\det\nolimits_\kappa(I)
\ar[dd]_{\epsilon}
\ar[rrr]_-{\gamma_{A \to B \to C} \otimes \gamma_{G \to H \to I}}
& & &
\det\nolimits_\kappa(B) \otimes
\det\nolimits_\kappa(H)
\ar[d]^{\gamma_{B \to E \to H}}
\\
& & & \det\nolimits_\kappa(E)
\\
\det\nolimits_\kappa(A) \otimes
\det\nolimits_\kappa(G) \otimes
\det\nolimits_\kappa(C) \otimes
\det\nolimits_\kappa(I)
\ar[rrr]^-{\gamma_{A \to D \to G} \otimes \gamma_{C \to F \to I}}
& & &
\det\nolimits_\kappa(D) \otimes
\det\nolimits_\kappa(F)
\ar[u]_{\gamma_{D \to E \to F}}
}
$$
ここで $\epsilon$ はテンソル積の因子を入れ替える写像に
$(-1)^{cg}$ を乗じたものであり、$c = \text{length}_R(C)$ および
$g = \text{length}_R(G)$ である。
\item $0 \to K \to L \to M \to 0$ が実際に $\kappa$-ベクトル空間の短完全列ならば、
写像 $\gamma_{K \to L \to M}$ は通常の同型と一致する。
\end{enumerate}
\end{lemma}

\begin{proof}
写像 $\gamma_{K \to L \to M}$ の明示的記述で非零記号を取る意味は次の点にある。
$(e_1, \ldots, e_l)$ が $K$ の admissible 列で、
$(\overline{f}_1, \ldots, \overline{f}_m)$ が $M$ の admissible 列であっても、
$(e_1, \ldots, e_l, f_1, \ldots, f_m)$ が $L$ の admissible 列であるとは限らない
（もちろん $f_i \in L$ は $\overline{f}_i$ の持ち上げを表す）。しかし記号
$[e_1, \ldots, e_l]$ が $\det_\kappa(K)$ において非零ならば、必ず
$K = \langle e_1, \ldots, e_k\rangle$ であり
（補題 \ref{chow-lemma-dimension-at-most-one} の証明を参照）、この場合
$(e_1, \ldots, e_k, f_1, \ldots, f_m)$ は実際に admissible 列である。
さらに $\det_\kappa(L)$ に対する型 (b) の admissible 関係式により、
$\det_\kappa(L)$ における $[e_1, \ldots, e_k, f_1, \ldots, f_m]$ の値は、
この場合も持ち上げ $f_i$ の選択に依存しない。この注意から、$K$ における
$e_1, \ldots, e_k$ の admissible 関係式は、$L$ における
$e_1, \ldots, e_k, f_1, \ldots, f_m$ の admissible 関係式に移り、
$\overline{f}_1, \ldots, \overline{f}_m$ の admissible 関係式についても同様である。
したがって $\gamma$ は補題で述べたベクトル空間の線形写像を定義する。

\medskip\noindent
補題 \ref{chow-lemma-determinant-dimension-one} により、$\det_\kappa(L)$ は
$L = \langle x_1, \ldots, x_{k + m}\rangle$ を満たす任意の admissible 列
$(x_1, \ldots, x_{k + m})$ に対する一つの記号
$[x_1, \ldots, x_{k + m}]$ により生成される。したがって写像
$\gamma_{K \to L \to M}$ は全射であり、ゆえに同型である。

\medskip\noindent
性質 (1) は次の計算から従う。
\begin{eqnarray*}
& & \det\nolimits_\kappa(v)([e_1, \ldots, e_k, f_1, \ldots, f_m]) \\
& = &
[v(e_1), \ldots, v(e_k), v(f_1), \ldots, v(f_m)] \\
& = &
\gamma_{K' \to L' \to M'}([u(e_1), \ldots, u(e_k)]
\otimes [w(f_1), \ldots, w(f_m)]).
\end{eqnarray*}
性質 (2) は、$\det_\kappa(A)$ を生成する記号
$[\alpha_1, \ldots, \alpha_a]$、$\det_\kappa(C)$ を生成する記号
$[\gamma_1, \ldots, \gamma_c]$、$\det_\kappa(G)$ を生成する記号
$[\zeta_1, \ldots, \zeta_g]$、および $\det_\kappa(I)$ を生成する記号
$[\iota_1, \ldots, \iota_i]$ が与えられたとき
\begin{eqnarray*}
& & [\alpha_1, \ldots, \alpha_a, \tilde\gamma_1, \ldots, \tilde\gamma_c,
\tilde\zeta_1, \ldots, \tilde\zeta_g, \tilde\iota_1, \ldots, \tilde\iota_i] \\
& = &
(-1)^{cg} [\alpha_1, \ldots, \alpha_a, \tilde\zeta_1, \ldots, \tilde\zeta_g,
\tilde\gamma_1, \ldots, \tilde\gamma_c, \tilde\iota_1, \ldots, \tilde\iota_i]
\end{eqnarray*}
が $\det_\kappa(E)$ において成り立つことを意味する
（$E$ における適切な持ち上げ $\tilde{x}$ を取る）。これは型 (c) の admissible 関係式を
$cg$ 回、次の順で用いれば従う。まず $\tilde\zeta_1$ を
$\tilde\gamma_c, \ldots, \tilde\gamma_1$ の向こうへ移す
（$\mathfrak m\tilde\zeta_1 \subset A$ なので許される）。次に $\tilde\zeta_2$ を
$\tilde\gamma_c, \ldots, \tilde\gamma_1$ の向こうへ移す
（$\mathfrak m\tilde\zeta_2 \subset A + R\tilde\zeta_1$ なので許される）。以下同様である。

\medskip\noindent
補題の (3) は明らかである。これで証明が完了する。
\end{proof}

\noindent
補題の写像 $\gamma$ を用いて、より一般の写像 $\gamma$ を次のように定義できる。
$(R, \mathfrak m, \kappa)$ を局所環とし、$M$ を有限長 $R$-加群とする。
有限フィルトレーション（『ホモロジー代数』の定義 \ref{homology-definition-filtered} を参照）
$$
0 = F^m \subset F^{m - 1} \subset \ldots \subset F^{n + 1} \subset F^n = M
$$
が与えられたと仮定する。このとき整定義された標準的同型
$$
\gamma_{(M, F)} :
\det\nolimits_\kappa(F^{m - 1}/F^m) \otimes_\kappa \ldots \otimes_k 
\det\nolimits_\kappa(F^n/F^{n + 1})
\longrightarrow
\det\nolimits_\kappa(M)
$$
が存在する。これを構成するには、短完全列
$0 \to F^{i - 1}/F^i \to M/F^i \to M/F^{i - 1} \to 0$ から得られる
補題 \ref{chow-lemma-det-exact-sequences} の同型を用いる。
$G = 0$ とした同補題 \ref{chow-lemma-det-exact-sequences} の (2) により、短完全列
$0 \to F^i \to F^{i - 1} \to F^{i - 1}/F^i \to 0$ を用いても同じ同型を得る。

\medskip\noindent
行列式関手に典型的なもう一つの結果を挙げる。証明は難しくない。
通常、難しいのは行列式関手の存在を示すことである。

\begin{lemma}
\label{chow-lemma-uniqueness-det}
$(R, \mathfrak m, \kappa)$ を任意の局所環とする。
写像 $\gamma_{K \to L \to M}$ を備えた関手
$$
\det\nolimits_\kappa :
\left\{
\begin{matrix}
\text{finite length }R\text{-modules} \\
\text{with isomorphisms}
\end{matrix}
\right\}
\longrightarrow
\left\{
\begin{matrix}
1\text{-dimensional }\kappa\text{-vector spaces} \\
\text{with isomorphisms}
\end{matrix}
\right\}
$$
は次の性質によって特徴づけられる。
\begin{enumerate}
\item $\mathfrak m$ で零化される加群の部分圏への制限は、通常の行列式関手と同型である
（補題 \ref{chow-lemma-compare-det} を参照）。
\item 補題 \ref{chow-lemma-det-exact-sequences} の (1)、(2)、(3) が成り立つ。
\end{enumerate}
\end{lemma}

\begin{proof}
省略する。
\end{proof}

\begin{lemma}
\label{chow-lemma-determinant-quotient-ring}
$(R', \mathfrak m') \to (R, \mathfrak m)$ を、剰余体 $\kappa$ 上に同型を誘導する
局所環準同型とする。このとき任意の有限長 $R$-加群に対し、制限 $M_{R'}$ は
有限長 $R'$-加群であり、標準的同型
$$
\det\nolimits_{R, \kappa}(M)
\longrightarrow
\det\nolimits_{R', \kappa}(M_{R'})
$$
が存在する。この同型は $M$ に関して関手的であり、補題
\ref{chow-lemma-det-exact-sequences} の $\det_{R, \kappa}$ および
$\det_{R', \kappa}$ に対して定義された同型 $\gamma_{K \to L \to M}$ と両立する。
\end{lemma}

\begin{proof}
$R$-加群としての $M$ の長さが $l$ ならば、$R'$-加群としての $M$
（すなわち $M_{R'}$）の長さも $l$ である。『可換代数』の補題 \ref{algebra-lemma-pushdown-module} を参照せよ。
$\mathfrak m'$ は $\mathfrak m$ に写るので、$R$ 上の $M$ の admissible 列
$x_1, \ldots, x_l$ は $R'$ 上の $M$ の admissible 列でもある。
求める同型は、記号
$[x_1, \ldots, x_l] \in \det\nolimits_{R, \kappa}(M)$ を対応する記号
$[x_1, \ldots, x_l] \in \det\nolimits_{R', \kappa}(M)$ に送ることで得られる。
これが同型について関手的であり、補題 \ref{chow-lemma-det-exact-sequences} の同型
$\gamma$ と両立することは直ちに確認できる。
\end{proof}

\begin{remark}
\label{chow-remark-explain-determinant}
$(R, \mathfrak m, \kappa)$ を局所環とし、$\kappa$ の標数が零であるか、
または標数が $p$ で $p R = 0$ であると仮定する。
$M_1, \ldots, M_n$ を有限長 $R$-加群とする。以下で、$i = 1, \ldots, n$ に対して
$M_i$ を零化するイデアル $I \subset \mathfrak m$ と、標準的全射
$R/I \to \kappa$ の切断 $\sigma : \kappa \to R/I$ が存在することを示す。
$\sigma$ による $M_i$ の制限 $M_{i, \kappa}$ は、次元
$l_i = \text{length}_R(M_i)$ の $\kappa$-ベクトル空間であり、補題
\ref{chow-lemma-determinant-quotient-ring} を用いると
$$
\det\nolimits_\kappa(M_i) = \wedge_\kappa^{l_i}(M_{i, \kappa})
$$
を得る。これらの同型は、$I$ で零化される有限長 $R$-加群の短完全列に対する
補題 \ref{chow-lemma-det-exact-sequences} の同型 $\gamma_{K \to M \to L}$ と両立する。
したがって $\det_\kappa$ の性質の確認は、しばしば有限次元ベクトル空間の圏上の
通常の行列式の対応する性質の確認に帰着する。

\medskip\noindent
$I$ として、加群 $M = \bigoplus M_i$ の零化イデアル
（『可換代数』の定義 \ref{algebra-definition-annihilator}）を取れる。
この場合 $R/I \subset \text{End}_R(M)$ なので有限長である。
したがって $R/I$ は剰余体 $\kappa$ をもつ Artin 局所環である。
Artin 局所環は完備なので、Cohen 構造定理
（『可換代数』の定理 \ref{algebra-theorem-cohen-structure-theorem}）により
$R/I$ は係数環をもち、$R$ に関する仮定によりそれは体である。
\end{remark}

\noindent
線形写像の行列式を計算できる場合を一つ示す。実際にはどの場合にも不可解な点はない。
任意に選んだ例については例 \ref{chow-example-determinant-map} を参照せよ。

\begin{lemma}
\label{chow-lemma-times-u-determinant}
$R$ を剰余体 $\kappa$ をもつ局所環とし、$u \in R^*$ を単元とする。
$M$ を $R$ 上の有限長加群とし、$u$ 倍写像を $u_M : M \to M$ と表す。
このとき
$$
\det\nolimits_\kappa(u_M) :
\det\nolimits_\kappa(M)
\longrightarrow
\det\nolimits_\kappa(M)
$$
は $\overline{u}^l$ 倍である。ここで $l = \text{length}_R(M)$ であり、
$\overline{u} \in \kappa^*$ は $u$ の像である。
\end{lemma}

\begin{proof}
$\det\nolimits_\kappa(u_M) = f_M \text{id}_{\det\nolimits_\kappa(M)}$ を満たす元を
$f_M \in \kappa^*$ と表す。$0 \to K \to L \to M \to 0$ を有限 $R$-加群の
短完全列とする。このとき $u_k$, $u_L$, $u_M$ は短完全列の同型を与える。
したがって補題 \ref{chow-lemma-det-exact-sequences} (1) により
$f_K f_M = f_L$ である。ゆえに長さに関する帰納法により、長さ $1$ の場合を
証明すれば十分であり、その場合は自明である。
\end{proof}

\begin{example}
\label{chow-example-determinant-map}
局所環 $R = \mathbf{Z}_p$ を考え、
$M = \mathbf{Z}_p/(p^2) \oplus \mathbf{Z}_p/(p^3)$ とおく。
$u : M \to M$ を行列
$$
u =
\left(
\begin{matrix}
a & b \\
pc & d
\end{matrix}
\right)
$$
で与えられる写像とする。ここで $a, b, c, d \in \mathbf{Z}_p$ かつ
$a, d \in \mathbf{Z}_p^*$ である。この場合 $\det_\kappa(u)$ は
$a^2d^3 \bmod p \in \mathbf{F}_p^*$ 倍に等しい。これは、$u$ が記号
$[p^2e, pe, pf, e, f]$ に及ぼす作用を考えれば容易に分かる。ここで
$e = (0 , 1) \in M$ および $f = (1, 0) \in M$ である。
\end{example}





\subsection{周期複体と行列式}
\label{chow-subsection-periodic-complexes-determinants}

\noindent
$R$ を剰余体 $\kappa$ をもつ局所環とし、$(M, \varphi, \psi)$ を
$R$ 上の $(2, 1)$-周期複体とする。$M$ が有限長で、
$(M, \varphi, \psi)$ が完全であると仮定する。行列式の構成を用いて、
この状況の不変量を定義する。小節 \ref{chow-subsection-determinants-finite-length} を参照せよ。
$K_\varphi = \Ker(\varphi)$、$I_\varphi = \Im(\varphi)$、
$K_\psi = \Ker(\psi)$、$I_\psi = \Im(\psi)$ と略記する。
短完全列
$$
0 \to K_\varphi \to M \to I_\varphi \to 0, \quad
0 \to K_\psi \to M \to I_\psi \to 0
$$
は同型
$$
\gamma_\varphi :
\det\nolimits_\kappa(K_\varphi)
\otimes
\det\nolimits_\kappa(I_\varphi)
\longrightarrow
\det\nolimits_\kappa(M), \quad
\gamma_\psi :
\det\nolimits_\kappa(K_\psi)
\otimes
\det\nolimits_\kappa(I_\psi)
\longrightarrow
\det\nolimits_\kappa(M),
$$
を与える。補題 \ref{chow-lemma-det-exact-sequences} を参照せよ。
一方、複体の完全性は等式 $K_\varphi = I_\psi$ および
$K_\psi = I_\varphi$ を与え、したがって因子を入れ替えることで同型
$$
\sigma :
\det\nolimits_\kappa(K_\varphi)
\otimes
\det\nolimits_\kappa(I_\varphi)
\longrightarrow
\det\nolimits_\kappa(K_\psi)
\otimes
\det\nolimits_\kappa(I_\psi)
$$
を得る。この記法を用いて不変量を定義する。

\begin{definition}
\label{chow-definition-periodic-determinant}
$R$ を剰余体 $\kappa$ をもつ局所環とし、$(M, \varphi, \psi)$ を
$R$ 上の $(2, 1)$-周期複体とする。$M$ が有限長で、
$(M, \varphi, \psi)$ が完全であると仮定する。
{\it $(M, \varphi, \psi)$ の行列式}とは、合成
$$
\det\nolimits_\kappa(M)
\xrightarrow{\gamma_\psi \circ \sigma \circ \gamma_\varphi^{-1}}
\det\nolimits_\kappa(M)
$$
が
$(-1)^{\text{length}_R(I_\varphi)\text{length}_R(I_\psi)}
\det\nolimits_\kappa(M, \varphi, \psi)$ 倍となるような元
$$
\det\nolimits_\kappa(M, \varphi, \psi) \in \kappa^*
$$
をいう。
\end{definition}

\begin{remark}
\label{chow-remark-more-elementary}
上で導入した行列式を、より具体的に記述する。
$R$ を剰余体 $\kappa$ をもつ局所環とし、$(M, \varphi, \psi)$ を
$R$ 上の $(2, 1)$-周期複体とする。$M$ が有限長で、
$(M, \varphi, \psi)$ が完全であると仮定する。上と同様に
$I_\varphi = \Im(\varphi)$、$I_\psi = \Im(\psi)$ と略記する。
$\text{length}_R(I_\varphi) = a$ および $\text{length}_R(I_\psi) = b$ と仮定する。
完全性により $a + b = \text{length}_R(M)$ である。
admissible 列 $x_1, \ldots, x_a \in I_\varphi$ および
$y_1, \ldots, y_b \in I_\psi$ を、記号 $[x_1, \ldots, x_a]$ が
$\det_\kappa(I_\varphi)$ を生成し、記号 $[x_1, \ldots, x_b]$ が
$\det_\kappa(I_\psi)$ を生成するように選ぶ。
$\varphi(\tilde x_i) = x_i$ を満たす $\tilde x_i \in M$ と、
$\psi(\tilde y_j) = y_j$ を満たす $\tilde y_j \in M$ を選ぶ。
このとき $\det_\kappa(M, \varphi, \psi)$ は、$\det_\kappa(M)$ における等式
$$
[x_1, \ldots, x_a, \tilde y_1, \ldots, \tilde y_b]
=
(-1)^{ab} \det\nolimits_\kappa(M, \varphi, \psi)
[y_1, \ldots, y_b, \tilde x_1, \ldots, \tilde x_a]
$$
によって特徴づけられる。これは符号の理由も説明する。
\end{remark}

\begin{lemma}
\label{chow-lemma-periodic-determinant-shift}
$R$ を剰余体 $\kappa$ をもつ局所環とし、$(M, \varphi, \psi)$ を
$R$ 上の $(2, 1)$-周期複体とする。$M$ が有限長で、
$(M, \varphi, \psi)$ が完全であると仮定する。このとき
$$
\det\nolimits_\kappa(M, \varphi, \psi)
\det\nolimits_\kappa(M, \psi, \varphi)
= 1.
$$
\end{lemma}

\begin{proof}
省略する。
\end{proof}

\begin{lemma}
\label{chow-lemma-periodic-determinant-sign}
$R$ を剰余体 $\kappa$ をもつ局所環とし、$(M, \varphi, \varphi)$ を
$R$ 上の $(2, 1)$-周期複体とする。$M$ が有限長で、
$(M, \varphi, \varphi)$ が完全であると仮定する。このとき
$\text{length}_R(M) = 2 \text{length}_R(\Im(\varphi))$ であり
$$
\det\nolimits_\kappa(M, \varphi, \varphi)
=
(-1)^{\text{length}_R(\Im(\varphi))}
=
(-1)^{\frac{1}{2}\text{length}_R(M)}
$$
である。
\end{lemma}

\begin{proof}
定義の符号規則から直ちに従う。
\end{proof}

\begin{lemma}
\label{chow-lemma-periodic-determinant-easy-case}
$R$ を剰余体 $\kappa$ をもつ局所環とし、$M$ を有限長 $R$-加群とする。
\begin{enumerate}
\item $\varphi : M \to M$ が同型ならば
$\det_\kappa(M, \varphi, 0) = \det_\kappa(\varphi)$ である。
\item $\psi : M \to M$ が同型ならば
$\det_\kappa(M, 0, \psi) = \det_\kappa(\psi)^{-1}$ である。
\end{enumerate}
\end{lemma}

\begin{proof}
(1) を証明する。$\psi = 0$ とおく。定義
\ref{chow-definition-periodic-determinant} の前の記法により
$K_\varphi = I_\psi = 0$、$I_\varphi = K_\psi = M$ と同一視できる。
この同一視のもとで写像
$$
\gamma_\varphi :
\kappa \otimes \det\nolimits_\kappa(M)
=
\det\nolimits_\kappa(K_\varphi)
\otimes
\det\nolimits_\kappa(I_\varphi)
\longrightarrow
\det\nolimits_\kappa(M)
$$
は $\det_\kappa(\varphi^{-1})$ と同一視される。一方、写像
$\gamma_\psi$ は恒等写像と同一視される。したがってこの場合
$\gamma_\psi \circ \sigma \circ \gamma_\varphi^{-1}$ は
$\det_\kappa(\varphi)$ に等しい。これで結論を得る。(2) の証明は省略する。
\end{proof}

\begin{lemma}
\label{chow-lemma-periodic-determinant}
$R$ を剰余体 $\kappa$ をもつ局所環とする。
$(2, 1)$-周期複体の短完全列
$$
0 \to (M_1, \varphi_1, \psi_1)
\to (M_2, \varphi_2, \psi_2)
\to (M_3, \varphi_3, \psi_3)
\to 0
$$
があり、すべての $M_i$ が有限長で、各 $(M_1, \varphi_1, \psi_1)$ が
完全であると仮定する。このとき $\kappa^*$ において
$$
\det\nolimits_\kappa(M_2, \varphi_2, \psi_2) =
\det\nolimits_\kappa(M_1, \varphi_1, \psi_1)
\det\nolimits_\kappa(M_3, \varphi_3, \psi_3).
$$
である。
\end{lemma}

\begin{proof}
$I_{\varphi, i} = \Im(\varphi_i)$、$K_{\varphi, i} = \Ker(\varphi_i)$、
$I_{\psi, i} = \Im(\psi_i)$、$K_{\psi, i} = \Ker(\psi_i)$ と略記する。
行と列が完全な有限長 $R$-加群の可換正方形
$$
\xymatrix{
& 0 \ar[d] & 0 \ar[d] & 0 \ar[d] & \\
0 \ar[r] &
K_{\varphi, 1} \ar[r] \ar[d] &
K_{\varphi, 2} \ar[r] \ar[d] &
K_{\varphi, 3} \ar[r] \ar[d] &
0 \\
0 \ar[r] &
M_1 \ar[r] \ar[d] &
M_2 \ar[r] \ar[d] &
M_3 \ar[r] \ar[d] &
0 \\
0 \ar[r] &
I_{\varphi, 1} \ar[r] \ar[d] &
I_{\varphi, 2} \ar[r] \ar[d] &
I_{\varphi, 3} \ar[r] \ar[d] &
0 \\
& 0  & 0  & 0  &
}
$$
があることに注意する。上段は像の列
$I_{\psi, 1} \to I_{\psi, 2} \to I_{\psi, 3} \to 0$ と同一視できるので完全であり、
下段についても同様である。加群 $I_{\psi, i}$ と $K_{\psi, i}$ を含む同様の図式もある。
定義により $\det_\kappa(M_2, \varphi_2, \psi_2)$ は、符号を除いて、次の図式の
左側の縦写像の合成に対応する。
$$
\xymatrix{
\det_\kappa(M_1) \otimes
\det_\kappa(M_3) \ar[r]^\gamma
\ar[d]^{\gamma^{-1} \otimes \gamma^{-1}} &
\det_\kappa(M_2)
\ar[d]^{\gamma^{-1}} \\
\det\nolimits_\kappa(K_{\varphi, 1})
\otimes
\det\nolimits_\kappa(I_{\varphi, 1})
\otimes
\det\nolimits_\kappa(K_{\varphi, 3})
\otimes
\det\nolimits_\kappa(I_{\varphi, 3})
\ar[d]^{\sigma \otimes \sigma}
\ar[r]^-{\gamma \otimes \gamma} &
\det\nolimits_\kappa(K_{\varphi, 2})
\otimes
\det\nolimits_\kappa(I_{\varphi, 2})
\ar[d]^\sigma
\\
\det\nolimits_\kappa(K_{\psi, 1})
\otimes
\det\nolimits_\kappa(I_{\psi, 1})
\otimes
\det\nolimits_\kappa(K_{\psi, 3})
\otimes
\det\nolimits_\kappa(I_{\psi, 3})
\ar[d]^{\gamma \otimes \gamma}
\ar[r]^-{\gamma \otimes \gamma}
&
\det\nolimits_\kappa(K_{\psi, 2})
\otimes
\det\nolimits_\kappa(I_{\psi, 2})
\ar[d]^\gamma \\
\det_\kappa(M_1)
\otimes
\det_\kappa(M_3) \ar[r]^\gamma
&
\det_\kappa(M_2)
}
$$
補題 \ref{chow-lemma-det-exact-sequences} (2) を適用すると、上と下の正方形は符号を除いて
可換である。中央の正方形は自明に可換である（単に因子を入れ替えている）。したがって
$\det\nolimits_\kappa(M_2, \varphi_2, \psi_2) =
\epsilon \det\nolimits_\kappa(M_1, \varphi_1, \psi_1)
\det\nolimits_\kappa(M_3, \varphi_3, \psi_3)
$
がある符号 $\epsilon$ に対して成り立つ。符号は実際に計算でき、上の図式の外側の長方形は
\begin{eqnarray*}
\epsilon & = &
(-1)^{\text{length}(I_{\varphi, 1})\text{length}(K_{\varphi, 3})
+ \text{length}(I_{\psi, 1})\text{length}(K_{\psi, 3})} \\
& = &
(-1)^{\text{length}(I_{\varphi, 1})\text{length}(I_{\psi, 3})
+ \text{length}(I_{\psi, 1})\text{length}(I_{\varphi, 3})}
\end{eqnarray*}
を法として可換である（証明は省略する）。ここから符号も整合することが容易に従う。
\end{proof}

\begin{example}
\label{chow-example-dual-numbers}
$k$ を体とし、$k$ 上の双対数の環 $R = k[T]/(T^2)$ を考える。
$R$ における $T$ の類を $t$ と表す。$M = R$ とし、$u, v \in k^*$ を用いて
$\varphi = ut$、$\psi = vt$ とする。この場合 $\det_k(M)$ は生成元
$e = [t, 1]$ をもつ。$I_\varphi = K_\varphi = I_\psi = K_\psi = (t)$ と同一視する。
すると $\gamma_\varphi(t \otimes t) = u^{-1}[t, 1]$ である
（$u^{-1} \in M$ は $t \in I_\varphi$ の持ち上げだからである）。また同じ理由で
$\gamma_\psi(t \otimes t) = v^{-1}[t, 1]$ である。したがって
$\det_k(M, \varphi, \psi) = -u/v \in k^*$ である。
\end{example}

\begin{example}
\label{chow-example-Zp}
$R = \mathbf{Z}_p$、$M = \mathbf{Z}_p/(p^l)$ とする。
$a, b \geq 0$、$a + b = l$、$u, v \in \mathbf{Z}_p^*$ とし、
$\varphi = p^b u$ および $\varphi = p^a v$ とする。
例 \ref{chow-example-dual-numbers} と同様の計算により
\begin{eqnarray*}
\det\nolimits_{\mathbf{F}_p}(\mathbf{Z}_p/(p^l), p^bu, p^av) & = &
(-1)^{ab}u^a/v^b \bmod p \\
& = &
(-1)^{\text{ord}_p(\alpha)\text{ord}_p(\beta)}
\frac{\alpha^{\text{ord}_p(\beta)}}{\beta^{\text{ord}_p(\alpha)}} \bmod p
\end{eqnarray*}
を得る。ここで $\alpha = p^bu, \beta = p^av \in \mathbf{Z}_p$ である。
より一般の場合とその証明については補題 \ref{chow-lemma-symbol-is-usual-tame-symbol} を参照せよ。
\end{example}

\begin{example}
\label{chow-example-generic-vector-space}
$R = k$ を体とする。$M = k^{\oplus a} \oplus k^{\oplus b}$ を
$l = a + b$ 次元とする。$\varphi$ と $\psi$ を対角行列
$$
\varphi = \text{diag}(u_1, \ldots, u_a, 0, \ldots, 0),
\quad
\psi = \text{diag}(0, \ldots, 0, v_1, \ldots, v_b)
$$
とする。ここで $u_i, v_j \in k^*$ である。この場合
$$
\det\nolimits_k(M, \varphi, \psi)
=
\frac{u_1 \ldots u_a}{v_1 \ldots v_b}.
$$
である。これは直接計算するか、$l = 1$ の場合を計算して補題
\ref{chow-lemma-periodic-determinant} の加法性を用いれば分かる。
\end{example}

\begin{example}
\label{chow-example-special-vector-space}
$R = k$ を体とする。$M = k^{\oplus a} \oplus k^{\oplus a}$ を
$l = 2a$ 次元とする。$\varphi$ と $\psi$ をブロック行列
$$
\varphi =
\left(
\begin{matrix}
0 & U \\
0 & 0
\end{matrix}
\right),
\quad
\psi =
\left(
\begin{matrix}
0 & V \\
0 & 0
\end{matrix}
\right),
$$
とする。ここで $U, V \in \text{Mat}(a \times a, k)$ は可逆である。この場合
$$
\det\nolimits_k(M, \varphi, \psi)
=
(-1)^a\frac{\det(U)}{\det(V)}.
$$
であり、直接計算で分かる。$a = 1$ の場合は
例 \ref{chow-example-dual-numbers} の計算と同様である。
\end{example}

\begin{example}
\label{chow-example-a-la-oort}
$R = k$ を体とし、$M = k^{\oplus 4}$ とする。
$$
\varphi =
\left(
\begin{matrix}
  0 &   0 &   0 &   0 \\
u_1 &   0 &   0 &   0 \\
  0 &   0 &   0 &   0 \\
  0 &   0 & u_2 &   0
\end{matrix}
\right)
\quad
\varphi =
\left(
\begin{matrix}
  0 &   0 &   0 &   0 \\
  0 &   0 & v_2 &   0 \\
  0 &   0 &   0 &   0 \\
v_1 &   0 &   0 &   0
\end{matrix}
\right)
\quad
$$
とし、$u_1, u_2, v_1, v_2 \in k^*$ とする。このとき
$$
\det\nolimits_k(M, \varphi, \psi) = -\frac{u_1u_2}{v_1v_2}.
$$
である。
\end{example}

\noindent
次に、線形自己準同型の合成の行列式が行列式の積であるという事実の類似を扱う。
式が非常に長くなるのを避けるため、任意の $R$-加群写像
$\varphi : M \to M$ に対して $I_\varphi = \Im(\varphi)$ および
$K_\varphi = \Ker(\varphi)$ と書く。また射の組
$\varphi, \psi : M \to M$ に対して $\varphi\psi = \varphi \circ \psi$ と表す。

\begin{lemma}
\label{chow-lemma-multiplicativity-determinant}
$R$ を剰余体 $\kappa$ をもつ局所環とし、$M$ を有限長 $R$-加群とする。
$\alpha, \beta, \gamma$ を $M$ の自己準同型とし、次を仮定する。
\begin{enumerate}
\item $I_\alpha = K_{\beta\gamma}$ であり、$\alpha, \beta, \gamma$ の任意の置換についても同様である。
\item $K_\alpha = I_{\beta\gamma}$ であり、$\alpha, \beta, \gamma$ の任意の置換についても同様である。
\end{enumerate}
このとき次が成り立つ。
\begin{enumerate}
\item 組 $(M, \alpha, \beta\gamma)$ は完全な $(2, 1)$-周期複体である。
\item 組 $(I_\gamma, \alpha, \beta)$ は完全な $(2, 1)$-周期複体である。
\item 組 $(M/K_\beta, \alpha, \gamma)$ は完全な $(2, 1)$-周期複体である。
\item
$$
\det\nolimits_\kappa(M, \alpha, \beta\gamma)
=
\det\nolimits_\kappa(I_\gamma, \alpha, \beta)
\det\nolimits_\kappa(M/K_\beta, \alpha, \gamma).
$$
である。
\end{enumerate}
\end{lemma}

\begin{proof}
仮定から補題の (1) が従うことは明らかである。

\medskip\noindent
(1) を見るため、仮定から $I_{\gamma\alpha} = I_{\alpha\gamma}$ が従い、
核および他の任意の射の組についても同様であることに注意する。
さらに $I_{\gamma\beta} =I_{\beta\gamma} = K_\alpha \subset I_\gamma$ であり、
他の任意の組についても同様である。特に短完全列
$$
0 \to I_{\beta\gamma} \to I_\gamma \xrightarrow{\alpha} I_{\alpha\gamma} \to 0
$$
を得る。同様に短完全列
$$
0 \to I_{\alpha\gamma} \to I_\gamma \xrightarrow{\beta} I_{\beta\gamma} \to 0.
$$
を得る。これにより $(I_\gamma, \alpha, \beta)$ は完全な $(2, 1)$-周期複体である。
したがって補題の (2) が成り立つ。

\medskip\noindent
$\alpha$, $\gamma$ が $M/K_\beta$ の整定義された自己準同型を与えることを見るには、
$\alpha(K_\beta) \subset K_\beta$ および $\gamma(K_\beta) \subset K_\beta$ を確認すればよい。
これは $\alpha(K_\beta) = \alpha(I_{\gamma\alpha}) = I_{\alpha\gamma\alpha}
\subset I_{\alpha\gamma} = K_\beta$ から従い、他方も同様である。
写像 $\alpha : M/K_\beta \to M/K_\beta$ の核は
$K_{\beta\alpha}/K_\beta = I_\gamma/K_\beta$ である。同様に、
$\gamma : M/K_\beta \to M/K_\beta$ の核は $I_\alpha/K_\beta$ に等しい。
したがって (3) が成り立つ。

\medskip\noindent
$r = \text{length}_R(K_\alpha)$、$s = \text{length}_R(K_\beta)$、
$t = \text{length}_R(K_\gamma)$ とおく。上の完全列と仮定により
$\text{length}_R(I_\alpha) = s + t$、$\text{length}_R(I_\beta) = r + t$、
$\text{length}_R(I_\gamma) = r + s$、$\text{length}(M) = r + s + t$ である。
次を選ぶ。
\begin{enumerate}
\item $K_\alpha$ を生成する admissible 列 $x_1, \ldots, x_r \in K_\alpha$
\item $K_\beta$ を生成する admissible 列 $y_1, \ldots, y_s \in K_\beta$、
\item $K_\gamma$ を生成する admissible 列 $z_1, \ldots, z_t \in K_\gamma$、
\item $\beta\gamma\tilde x_i = x_i$ を満たす元 $\tilde x_i \in M$、
\item $\alpha\gamma\tilde y_i = y_i$ を満たす元 $\tilde y_i \in M$、
\item $\beta\alpha\tilde z_i = z_i$ を満たす元 $\tilde z_i \in M$。
\end{enumerate}
この選択のもとで列
$y_1, \ldots, y_s, \alpha\tilde z_1, \ldots, \alpha\tilde z_t$ は
$I_\alpha$ を生成する admissible 列である。したがって注意
\ref{chow-remark-more-elementary} により、行列式
$D = \det_\kappa(M, \alpha, \beta\gamma)$ は、次を満たす $\kappa^*$ の一意な元である。
\begin{align*}
[y_1, \ldots, y_s,
\alpha\tilde z_1, \ldots, \alpha\tilde z_s,
\tilde x_1, \ldots, \tilde x_r] \\
= (-1)^{r(s + t)} D
[x_1, \ldots, x_r,
\gamma\tilde y_1, \ldots, \gamma\tilde y_s,
\tilde z_1, \ldots, \tilde z_t]
\end{align*}
同じ注意により、$D_1 = \det_\kappa(M/K_\beta, \alpha, \gamma)$ は
$$
[y_1, \ldots, y_s,
\alpha\tilde z_1, \ldots, \alpha\tilde z_t,
\tilde x_1, \ldots, \tilde x_r]
=
(-1)^{rt} D_1
[y_1, \ldots, y_s,
\gamma\tilde x_1, \ldots, \gamma\tilde x_r,
\tilde z_1, \ldots, \tilde z_t]
$$
によって特徴づけられる。同じ注意により
$D_2 = \det_\kappa(I_\gamma, \alpha, \beta)$ は
$$
[y_1, \ldots, y_s,
\gamma\tilde x_1, \ldots, \gamma\tilde x_r,
\tilde z_1, \ldots, \tilde z_t]
=
(-1)^{rs} D_2
[x_1, \ldots, x_r,
\gamma\tilde y_1, \ldots, \gamma\tilde y_s,
\tilde z_1, \ldots, \tilde z_t]
$$
によって特徴づけられる。以上の公式を組み合わせると、望みどおり
$D = D_1 D_2$ を得る。
\end{proof}


\begin{lemma}
\label{chow-lemma-tricky}
$R$ を剰余体 $\kappa$ をもつ局所環とする。
$\alpha : (M, \varphi, \psi) \to (M', \varphi', \psi')$ を
$R$ 上の $(2, 1)$-周期複体の射とし、次を仮定する。
\begin{enumerate}
\item $M$, $M'$ は有限長である。
\item $(M, \varphi, \psi)$, $(M', \varphi', \psi')$ は完全である。
\item 写像 $\varphi$, $\psi$ が $K = \Ker(\alpha)$ 上に誘導する写像は零である。
\item 写像 $\varphi$, $\psi$ が $Q = \Coker(\alpha)$ 上に誘導する写像は零である。
\end{enumerate}
$N = \alpha(M) \subset M'$ と表す。$(2, 1)$-周期複体の二つの短完全列
$$
\begin{matrix}
0 \to (N, \varphi', \psi') \to (M', \varphi', \psi') \to (Q, 0, 0) \to 0 \\
0 \to (K, 0, 0) \to (M, \varphi, \psi) \to (N, \varphi', \psi') \to 0
\end{matrix}
$$
を得て、これらは二つの同型 $\alpha_i : Q \to K$、$i = 0, 1$ を誘導する。このとき
$$
\det\nolimits_\kappa(M, \varphi, \psi)
=
\det\nolimits_\kappa(\alpha_0^{-1} \circ \alpha_1)
\det\nolimits_\kappa(M', \varphi', \psi')
$$
である。特に $\alpha_0 = \alpha_1$ ならば
$\det\nolimits_\kappa(M, \varphi, \psi) =
\det\nolimits_\kappa(M', \varphi', \psi')$ である。
\end{lemma}

\begin{proof}
この補題の証明法は（少なくとも）二つある。一つは行列式の性質を用いて巨大な可換図式を
作る方法であり、もう一つは元の admissible 列による行列式の特徴づけを用いる方法である。
ここでは後者を用いる。

\medskip\noindent
まず写像 $\alpha_i$ を正確に説明する。$\alpha_0$ は、補題に掲げた複体の短完全列の
境界写像から得られる合成
$$
\alpha_0 : Q = H^0(Q, 0, 0) \to H^1(N, \varphi', \psi') \to H^2(K, 0, 0) = K
$$
であり、$\alpha_1$ は合成
$$
\alpha_1 : Q = H^1(Q, 0, 0) \to H^2(N, \varphi', \psi') \to H^3(K, 0, 0) = K
$$
である。複体 $(M, \varphi, \psi)$, $(M', \varphi', \psi')$ の完全性により、
これらの写像は同型である。

\medskip\noindent
$I_\varphi = \Im(\varphi)$、$K_\varphi = \Ker(\varphi)$ と書き、他の写像にも
同様の記法を用いる。$M$ および $M'$ の完全性は
$K_\varphi = I_\psi$ と、それと同様の三つの等式を意味する。
$k = \text{length}_R(K)$、$a = \text{length}_R(I_\varphi)$、
$b = \text{length}_R(I_\psi)$ とおく。このとき
$\text{length}_R(M) = a + b$、$\text{length}_R(N) = a + b - k$、
$\text{length}_R(Q) = k$、$\text{length}_R(M') = a + b$ である。
以下の完全列から $\text{length}_R(I_{\varphi'}) = a$ および
$\text{length}_R(I_{\psi'}) = b$ も従う。

\medskip\noindent
仮定 $K \subset K_\varphi = I_\psi$ は、$\varphi$ が $N$ を経由して完全列
$$
0 \to \alpha(I_\psi) \to N \xrightarrow{\varphi\alpha^{-1}} I_\psi \to 0.
$$
を与えることを意味する。ここで $\varphi\alpha^{-1}(x') = y$ とは、
$x' = \alpha(x)$ かつ $y = \varphi(x)$ であることを意味する。同様に
$$
0 \to \alpha(I_\varphi) \to N \xrightarrow{\psi\alpha^{-1}} I_\varphi \to 0.
$$
を得る。$\psi'$ が $Q$ 上に誘導する写像が零であるという仮定は
$I_{\psi'} = K_{\varphi'} \subset N$ を意味する。これは商
$\varphi'(N) \subset I_{\varphi'}$ が $Q$ と同一視されることを意味する。
$\varphi'(N) = \alpha(I_\varphi)$ に注意せよ。したがって同型
$$
\varphi' : Q \to I_{\varphi'}/\alpha(I_\varphi)
$$
があり、これは単に
$\varphi'(x' \bmod N) = \varphi'(x') \bmod \alpha(I_\varphi)$ と記述される。
まったく同様に
$$
\psi' : Q \to I_{\psi'}/\alpha(I_\psi)
$$
を得る。最後に $\alpha_0$ は合成
$$
\xymatrix{
Q \ar[r]^-{\varphi'} &
I_{\varphi'}/\alpha(I_\varphi)
\ar[rrr]^-{\psi\alpha^{-1}|_{I_{\varphi'}/\alpha(I_\varphi)}} & & &
K
}
$$
であり、同様に
$\alpha_1 = \varphi\alpha^{-1}|_{I_{\psi'}/\alpha(I_\psi)} \circ \psi'$ である。

\medskip\noindent
以下の式を短くするため、$\alpha(x)$ の代わりに $\alpha x$ と書く。
すべての写像はギリシャ文字、元はローマ字で表すので混同は生じない。
次を選ぶ。
\begin{enumerate}
\item $K$ を生成する admissible 列 $z_1, \ldots, z_k \in K$、
\item $\varphi z'_i = z_i$ を満たす元 $z'_i \in M$、
\item $\psi z''_i = z_i$ を満たす元 $z''_i \in M$、
\item $z_1, \ldots, z_k, x_{k + 1}, \ldots, x_a$ が $I_\varphi$ を生成する
admissible 列となるような元 $x_{k + 1}, \ldots, x_a \in I_\varphi$、
\item $\varphi \tilde x_i = x_i$ を満たす元 $\tilde x_i \in M$、
\item $z_1, \ldots, z_k, y_{k + 1}, \ldots, y_b$ が $I_\psi$ を生成する
admissible 列となるような元 $y_{k + 1}, \ldots, y_b \in I_\psi$、
\item $\psi \tilde y_i = y_i$ を満たす元 $\tilde y_i \in M$、
\item $w_1 \bmod N, \ldots, w_k \bmod N$ が $Q$ の admissible 列で $Q$ を生成するような
元 $w_1, \ldots, w_k \in M'$。
\end{enumerate}
注意 \ref{chow-remark-more-elementary} により、元
$D = \det_\kappa(M, \varphi, \psi) \in \kappa^*$ は次の等式によって特徴づけられる。
\begin{eqnarray*}
& &
[z_1, \ldots, z_k,
x_{k + 1}, \ldots, x_a,
z''_1, \ldots, z''_k,
\tilde y_{k + 1}, \ldots, \tilde y_b] \\
& = &
(-1)^{ab} D
[z_1, \ldots, z_k,
y_{k + 1}, \ldots, y_b,
z'_1, \ldots, z'_k,
\tilde x_{k + 1}, \ldots, \tilde x_a]
\end{eqnarray*}
上の議論により
$\alpha x_{k + 1}, \ldots, \alpha x_a, \varphi w_1, \ldots, \varphi w_k$ は
$I_{\varphi'}$ を生成する admissible 列であり、
$\alpha y_{k + 1}, \ldots, \alpha y_b, \psi w_1, \ldots, \psi w_k$ は
$I_{\psi'}$ を生成する admissible 列である。したがって注意
\ref{chow-remark-more-elementary} により、元
$D' = \det_\kappa(M', \varphi', \psi') \in \kappa^*$ は次の等式によって特徴づけられる。
\begin{eqnarray*}
& &
[\alpha x_{k + 1}, \ldots, \alpha x_a,
\varphi' w_1, \ldots, \varphi' w_k,
\alpha \tilde y_{k + 1}, \ldots, \alpha \tilde y_b,
w_1, \ldots, w_k]
\\
& = &
(-1)^{ab} D'
[\alpha y_{k + 1}, \ldots, \alpha y_b,
\psi' w_1, \ldots, \psi' w_k,
\alpha \tilde x_{k + 1}, \ldots, \alpha \tilde x_a,
w_1, \ldots, w_k]
\end{eqnarray*}
第一の表示式、resp.\ 第二の表示式では、両辺の記号の最初の、resp.\ 最後の
$k$ 個の成分が同じであることに注意せよ。したがってこれらの式は実際には等式
\begin{eqnarray*}
& &
[\alpha x_{k + 1}, \ldots, \alpha x_a,
\alpha z''_1, \ldots, \alpha z''_k,
\alpha \tilde y_{k + 1}, \ldots, \alpha \tilde y_b] \\
& = &
(-1)^{ab} D
[\alpha y_{k + 1}, \ldots, \alpha y_b,
\alpha z'_1, \ldots, \alpha z'_k,
\alpha \tilde x_{k + 1}, \ldots, \alpha \tilde x_a]
\end{eqnarray*}
および
\begin{eqnarray*}
& &
[\alpha x_{k + 1}, \ldots, \alpha x_a,
\varphi' w_1, \ldots, \varphi' w_k,
\alpha \tilde y_{k + 1}, \ldots, \alpha \tilde y_b]
\\
& = &
(-1)^{ab} D'
[\alpha y_{k + 1}, \ldots, \alpha y_b,
\psi' w_1, \ldots, \psi' w_k,
\alpha \tilde x_{k + 1}, \ldots, \alpha \tilde x_a]
\end{eqnarray*}
と同値であり、これらは $\det_\kappa(N)$ における等式である。
$\varphi' w_1, \ldots, \varphi' w_k$ および
$\alpha z''_1, \ldots, z''_k$ は加群 $I_{\varphi'}/\alpha(I_\varphi)$ を生成する
admissible 列であることに注意する。ある $\lambda_0 \in \kappa^*$ に対して
$\det_\kappa(I_{\varphi'}/\alpha(I_\varphi))$ において
$$
[\varphi' w_1, \ldots, \varphi' w_k]
= \lambda_0 [\alpha z''_1, \ldots, \alpha z''_k]
$$
と書く。同様に、ある $\lambda_1 \in \kappa^*$ に対して
$\det_\kappa(I_{\psi'}/\alpha(I_\psi))$ において
$$
[\psi' w_1, \ldots, \psi' w_k]
= \lambda_1 [\alpha z'_1, \ldots, \alpha z'_k]
$$
と書く。一方、上の $\alpha_i$ の記述により、$i = 0, 1$ に対し明らかに
$$
\alpha_i([w_1, \ldots, w_k]) = \lambda_i[z_1, \ldots, z_k]
$$
であり、これは
$$
\det\nolimits_\kappa(\alpha_0^{-1} \circ \alpha_1)
=
\lambda_1/\lambda_0
$$
を意味する。他方、明らかに
\begin{eqnarray*}
& &
\lambda_0 [\alpha x_{k + 1}, \ldots, \alpha x_a,
\alpha z''_1, \ldots, \alpha z''_k,
\alpha \tilde y_{k + 1}, \ldots, \alpha \tilde y_b] \\
& = &
[\alpha x_{k + 1}, \ldots, \alpha x_a,
\varphi' w_1, \ldots, \varphi' w_k,
\alpha \tilde y_{k + 1}, \ldots, \alpha \tilde y_b]
\end{eqnarray*}
および
\begin{eqnarray*}
& &
\lambda_1[\alpha y_{k + 1}, \ldots, \alpha y_b,
\alpha z'_1, \ldots, \alpha z'_k,
\alpha \tilde x_{k + 1}, \ldots, \alpha \tilde x_a] \\
& = &
[\alpha y_{k + 1}, \ldots, \alpha y_b,
\psi' w_1, \ldots, \psi' w_k,
\alpha \tilde x_{k + 1}, \ldots, \alpha \tilde x_a]
\end{eqnarray*}
である。これらは $\lambda_0 D = \lambda_1 D'$ を含意し、補題が従う。
\end{proof}








\subsection{記号}
\label{chow-subsection-symbols}

\noindent
この構成に適切な一般性は、おそらく次の補題の状況である。

\begin{lemma}
\label{chow-lemma-pre-symbol}
$A$ を Noether 局所環とし、$M$ を次元 $1$ の有限 $A$-加群とする。
$\varphi, \psi : M \to M$ を二つの単射 $A$-加群写像とし、
$\varphi(\psi(M)) = \psi(\varphi(M))$ と仮定する。例えば $\varphi$ と $\psi$ が
可換ならばこの仮定は成り立つ。このとき
$\text{length}_A(M/\varphi\psi M) < \infty$ であり、
$(M/\varphi\psi M, \varphi, \psi)$ は完全な $(2, 1)$-周期複体である。
\end{lemma}

\begin{proof}
$\mathfrak q$ を $M$ の台の極小素イデアルとする。
『可換代数』の補題 \ref{algebra-lemma-support-point} により、
$M_{\mathfrak q}$ は有限長 $A_{\mathfrak q}$-加群である。
したがって $\varphi$ と $\psi$ はいずれも同型
$M_{\mathfrak q} \to M_{\mathfrak q}$ を誘導する。ゆえに
$M/\varphi\psi M$ の台は $\{\mathfrak m_A\}$ であり、有限長である
（上で引用した補題を参照）。最後に、$M/\varphi\psi M$ 上の $\varphi$ の核は明らかに
$\psi M/\varphi\psi M$ なので、$\varphi$ の核は $M/\varphi\psi M$ 上の
$\psi$ の像である。仮定により $M/\varphi\psi M = M/\psi\varphi M$ なので、
逆方向も同様である。
\end{proof}

\begin{lemma}
\label{chow-lemma-symbol-defined}
$A$ を Noether 局所環とし、$a, b \in A$ とする。
\begin{enumerate}
\item $M$ が次元 $1$ の有限 $A$-加群で、$a, b$ が $M$ 上の非零因子ならば、
$\text{length}_A(M/abM) < \infty$ であり、$(M/abM, a, b)$ は
完全な $(2, 1)$-周期複体である。
\item $a, b$ が非零因子で $\dim(A) = 1$ ならば、
$\text{length}_A(A/(ab)) < \infty$ であり、$(A/(ab), a, b)$ は
完全な $(2, 1)$-周期複体である。
\end{enumerate}
特にこれらの場合、$\det_\kappa(M/abM, a, b) \in \kappa^*$、
resp.\ $\det_\kappa(A/(ab), a, b) \in \kappa^*$ が定義される。
\end{lemma}

\begin{proof}
補題 \ref{chow-lemma-pre-symbol} から従う。
\end{proof}

\begin{definition}
\label{chow-definition-symbol-M}
$A$ を剰余体 $\kappa$ をもつ Noether 局所環とし、$a, b \in A$ とする。
$M$ を次元 $1$ の有限 $A$-加群で、$a, b$ が $M$ 上の非零因子であるものとする。
{\it $M, a, b$ に付随する記号}を元
$$
d_M(a, b) =
\det\nolimits_\kappa(M/abM, a, b) \in \kappa^*
$$
として定義する。
\end{definition}

\begin{lemma}
\label{chow-lemma-multiplicativity-symbol}
$A$ を Noether 局所環とする。$a, b, c \in A$ とし、$M$ を
$\dim(\text{Supp}(M)) = 1$ を満たす有限 $A$-加群とする。
$a, b, c$ が $M$ 上の非零因子であると仮定する。このとき
$$
d_M(a, bc) = d_M(a, b) d_M(a, c)
$$
および $d_M(a, b)d_M(b, a) = 1$ が成り立つ。
\end{lemma}

\begin{proof}
第一の主張は、$M/abcM$ と $a, b, c$ 倍で与えられる自己準同型
$\alpha, \beta, \gamma$ に補題 \ref{chow-lemma-multiplicativity-determinant} を適用すれば従う。
第二の主張は補題 \ref{chow-lemma-periodic-determinant-shift} から従う。
\end{proof}

\begin{definition}
\label{chow-definition-tame-symbol}
$A$ を剰余体 $\kappa$ をもつ次元 $1$ の Noether 局所整域とする。
$K$ を $A$ の分数体とする。$A$ の {\it tame 記号}を写像
$$
K^* \times K^* \longrightarrow \kappa^*,
\quad
(x, y) \longmapsto d_A(x, y)
$$
と定義する。ここで $d_A(x, y)$ は補題 \ref{chow-lemma-multiplicativity-symbol} の乗法性により
$K^* \times K^*$ へ拡張したものである。
\end{definition}

\noindent
より一般に $d_M(-, -)$ を $A$ のある種の分数環へ拡張できることは明らかである
（$A$ が整域でない場合も含む）。

\begin{lemma}
\label{chow-lemma-symbol-when-equal}
$A$ を Noether 局所環、$M$ を次元 $1$ の有限 $A$-加群とする。
$a \in A$ を $M$ 上の非零因子とする。このとき
$d_M(a, a) = (-1)^{\text{length}_A(M/aM)}$ である。
\end{lemma}

\begin{proof}
補題 \ref{chow-lemma-periodic-determinant-sign} から直ちに従う。
\end{proof}

\begin{lemma}
\label{chow-lemma-symbol-when-one-is-a-unit}
$A$ を Noether 局所環とし、$M$ を次元 $1$ の有限 $A$-加群とする。
$b \in A$ を $M$ 上の非零因子とし、$u \in A^*$ とする。このとき
$$
d_M(u, b) = u^{\text{length}_A(M/bM)} \bmod \mathfrak m_A.
$$
特に $M = A$ ならば
$d_A(u, b) = u^{\text{ord}_A(b)} \bmod \mathfrak m_A$ である。
\end{lemma}

\begin{proof}
この場合 $M/ubM = M/bM$ であり、その上で $b$ 倍は零であることに注意する。
したがって補題 \ref{chow-lemma-periodic-determinant-easy-case} により
$d_M(u, b) = \det_\kappa(u|_{M/bM})$ である。補題
\ref{chow-lemma-times-u-determinant} から結論が従う。
\end{proof}

\begin{lemma}
\label{chow-lemma-symbol-short-exact-sequence}
$A$ を Noether 局所環とし、$a, b \in A$ とする。
$a, b$ が三つすべての $A$-加群上の非零因子であるような、
次元 $1$ の $A$-加群の短完全列
$$
0 \to M \to M' \to M'' \to 0
$$
を考える。このとき $\kappa^*$ において
$$
d_{M'}(a, b) = d_M(a, b) d_{M''}(a, b)
$$
である。
\end{lemma}

\begin{proof}
これが完全な $(2, 1)$-周期複体の短完全列
$$
0 \to
(M/abM, a, b) \to
(M'/abM', a, b) \to
(M''/abM'', a, b) \to 0
$$
を導くことは容易に分かる。したがって補題
\ref{chow-lemma-periodic-determinant} から結論が従う。
\end{proof}

\begin{lemma}
\label{chow-lemma-symbol-compare-modules}
$A$ を Noether 局所環とし、$\alpha : M \to M'$ を次元 $1$ の有限 $A$-加群の
準同型とする。$a, b \in A$ とし、次を仮定する。
\begin{enumerate}
\item $a$, $b$ は $M$ と $M'$ の両方の上で非零因子である。
\item $\dim(\Ker(\alpha)), \dim(\Coker(\alpha)) \leq 0$ である。
\end{enumerate}
このとき $d_M(a, b) = d_{M'}(a, b)$ である。
\end{lemma}

\begin{proof}
$a \in A^*$ ならば、等式は
$\text{length}(M/bM) = \text{length}(M'/bM')$ と補題
\ref{chow-lemma-symbol-when-one-is-a-unit} から従う。同様に $b$ が単元ならば、
補題 \ref{chow-lemma-multiplicativity-symbol} の対称性により補題は成り立つ。
したがって $a, b \in \mathfrak m_A$ と仮定してよい。
特にこれは $\mathfrak m$ が $M$ の随伴素イデアルではないことを含意するので、
$\alpha : M \to M'$ は単射である。よって $M$ を $M'$ の部分加群とみなせる。
仮定により $M'/M$ は台 $\{\mathfrak m_A\}$ をもつ有限 $A$-加群なので有限長である。
$M \subset M'' \subset M'$ を満たす任意の第三の加群 $M''$ に対し、写像
$M \to M''$ および $M'' \to M'$ も補題の仮定を満たすことに注意する。
したがって $M'/M$ の長さに関する帰納法により
$\text{length}_A(M'/M) = 1$ の場合に帰着できる。最後にこの場合、写像
$$
\overline{\alpha} : M/abM \longrightarrow M'/abM'.
$$
を考える。構成により $\overline{\alpha}$ の余核 $Q$ の長さは $1$ である。
$a, b \in \mathfrak m_A$ なので、これらは $Q$ 上で自明に作用する。
また $\overline{\alpha}$ の核 $K$ の長さも $1$ となるので、$a$, $b$ は $K$ 上でも
自明に作用する。したがって補題 \ref{chow-lemma-tricky} を適用できる。
よって二つの写像 $\alpha_i : Q \to K$ が等しいことを示せば十分である。
実際、いずれの写像も
$q = x' \bmod \Im(\overline{\alpha}) \mapsto abx' \in K$ に等しい。
確認は省略する。
\end{proof}

\begin{lemma}
\label{chow-lemma-compute-symbol-M}
$A$ を Noether 局所環とし、$M$ を $\dim(\text{Supp}(M)) = 1$ を満たす
有限 $A$-加群とする。$a, b \in A$ を $M$ 上の非零因子とする。
$\mathfrak q_1, \ldots, \mathfrak q_t$ を $M$ の台の極小素イデアルとする。このとき
$$
d_M(a, b)
=
\prod\nolimits_{i = 1, \ldots, t}
d_{A/\mathfrak q_i}(a, b)^{
\text{length}_{A_{\mathfrak q_i}}(M_{\mathfrak q_i})}
$$
が $\kappa^*$ の元として成り立つ。
\end{lemma}

\begin{proof}
$A$-部分加群によるフィルトレーション
$$
0 = M_0 \subset M_1 \subset \ldots \subset M_n = M
$$
で、各商 $M_j/M_{j - 1}$ が $A$ のある素イデアル $\mathfrak p_j$ に対する
$A/\mathfrak p_j$ と同型になるものを選ぶ。『可換代数』の補題 \ref{algebra-lemma-filter-Noetherian-module} を参照せよ。
各 $j$ に対し、ある $i$ について $\mathfrak p_j = \mathfrak q_i$ であるか、
$\mathfrak p_j = \mathfrak m_A$ である。さらに固定した $i$ に対し、
$\mathfrak p_j = \mathfrak q_i$ を満たす $j$ の個数は、『可換代数』の補題 \ref{algebra-lemma-filter-minimal-primes-in-support} により
$\text{length}_{A_{\mathfrak q_i}}(M_{\mathfrak q_i})$ に等しい。
したがって各 $j$ に対して $d_{M_j}(a, b)$ が定義され、
$$
d_{M_j}(a, b)
=
\left\{
\begin{matrix}
d_{M_{j - 1}}(a, b) d_{A/\mathfrak q_i}(a, b) &
\text{if} & \mathfrak p_j = \mathfrak q_i \\
d_{M_{j - 1}}(a, b) & \text{if} & \mathfrak p_j = \mathfrak m_A
\end{matrix}
\right.
$$
である。第一の場合には補題 \ref{chow-lemma-symbol-short-exact-sequence}、第二の場合には
補題 \ref{chow-lemma-symbol-compare-modules} を用いた。よって補題が従う。
\end{proof}

\begin{lemma}
\label{chow-lemma-symbol-is-usual-tame-symbol}
$A$ を分数体 $K$ をもつ離散付値環とする。非零元 $x, y \in K$ に対し
$$
d_A(x, y)
=
(-1)^{\text{ord}_A(x)\text{ord}_A(y)}
\frac{x^{\text{ord}_A(y)}}{y^{\text{ord}_A(x)}} \bmod \mathfrak m_A,
$$
である。すなわち、この記号は通常の tame 記号に等しい。
\end{lemma}

\begin{proof}
乗法性により $x, y \in A$ の場合を示せば十分である。
$t \in A$ を素元とする。ある $a, b \geq 0$ および $u, v \in A^*$ を用いて
$x = t^bu$、$y = t^bv$ と書き、$l = a + b$ とおく。このとき
$t^{l - 1}, \ldots, t^b$ は $(x)/(xy)$ の admissible 列であり、
$t^{l - 1}, \ldots, t^a$ は $(y)/(xy)$ の admissible 列である。
したがって注意 \ref{chow-remark-more-elementary} により、$d_A(x, y)$ は等式
$$
[t^{l - 1}, \ldots, t^b, v^{-1}t^{b - 1}, \ldots, v^{-1}]
=
(-1)^{ab} d_A(x, y)
[t^{l - 1}, \ldots, t^a, u^{-1}t^{a - 1}, \ldots, u^{-1}].
$$
によって特徴づけられる。したがって記号 $[x_1, \ldots, x_l]$ の admissible 関係式により
$$
d_A(x, y) = (-1)^{ab} u^a/v^b \bmod \mathfrak m_A
$$
を得て、望みどおりである。
\end{proof}

\begin{lemma}
\label{chow-lemma-symbol-is-steinberg-prepare}
$A$ を Noether 局所環とし、$a, b \in A$ とする。
$M$ を次元 $1$ の有限 $A$-加群で、$a$, $b$, $b - a$ の各々がその上の
非零因子であるものとする。このとき $\kappa^*$ において
$$
d_M(a, b - a)d_M(b, b) = d_M(b, b - a)d_M(a, b)
$$
である。
\end{lemma}

\begin{proof}
補題 \ref{chow-lemma-compute-symbol-M} により、ある素イデアル
$\mathfrak q \subset A$ について $M = A/\mathfrak q$ かつ
$\dim(A/\mathfrak q) = 1$ の場合に関係式を示せば十分である。

\medskip\noindent
$M = A/\mathfrak q$ の場合、$A$ を $A/\mathfrak q$ で、$a, b$ を
$A/\mathfrak q$ における像で置き換えられる。したがって $A = M$ で、
$A$ が次元 $1$ の Noether 局所整域であると仮定してよい。
これは、$A$ と $A/\mathfrak q$ の剰余体 $\kappa$ が同じであり、任意の
$A/\mathfrak q$-加群 $M$ に対して $A$ 上で取った行列式と
$A/\mathfrak q$ 上で取った行列式が標準的に同一視されるからである。
補題 \ref{chow-lemma-determinant-quotient-ring} を参照せよ。

\medskip\noindent
$a, b$ がともに極大イデアルに属する場合に関係式を示せば十分である。
実際、一方または両方が単元の場合は補題
\ref{chow-lemma-symbol-when-one-is-a-unit} および \ref{chow-lemma-symbol-when-equal} から従う。

\medskip\noindent
補題 \ref{chow-lemma-Noetherian-domain-dim-1-two-elements} のとおり、拡大
$A \subset A'$ と分解 $a = ta'$、$b = tb'$ を選ぶ。
$b - a = t(b' - a')$ であり、
$A' = (a', b') = (a', b' - a') = (b' - a', b')$ でもあることに注意する。
以下、$A'$ を $A$-加群とみなし、$a, b, a', b', t$ を $A'$ の
$A$-加群自己準同型とみなす。記法 $d^A_{A'}(a', b')$ などを
$$
d^A_{A'}(a', b')
=
\det\nolimits_\kappa(A'/a'b'A', a', b')
$$
の意味で用いる。これは補題 \ref{chow-lemma-pre-symbol} により定義される。
上付き添字 ${}^A$ は、既に定義した記号 $d_{A'}(a', b')$ と区別するために用いる。
後者は異なる記号である（例えば値を取る $A'$ の剰余体は $\kappa$ と異なることがある）。
補題 \ref{chow-lemma-symbol-compare-modules} により
$d_A(a, b) = d^A_{A'}(a, b)$ であり、他の組合せについても同様である。
これと乗法性を用いると
$$
d^A_{A'}(a', b' - a') d^A_{A'}(b', b')
=
d^A_{A'}(b', b' - a') d^A_{A'}(a', b')
$$
を示せば十分である。ここで $(a', b') = A'$ などであるから
$$
\begin{matrix}
A'/(a'(b' - a')) & \cong & A'/(a') \oplus A'/(b' - a') \\
A'/(b'(b' - a')) & \cong & A'/(b') \oplus A'/(b' - a') \\
A'/(a'b') & \cong & A'/(a') \oplus A'/(b')
\end{matrix}
$$
を得る。さらに $A/(a')$ 上の $b' - a'$ 倍は $b'$ 倍に等しく、
$A/(b')$ 上の $b' - a'$ 倍は $-a'$ 倍に等しいことに注意する。
補題 \ref{chow-lemma-periodic-determinant-easy-case} および
\ref{chow-lemma-periodic-determinant} を用いると
$$
\begin{matrix}
d^A_{A'}(a', b' - a') & = &
\det\nolimits_\kappa(b'|_{A'/(a')})^{-1}
\det\nolimits_\kappa(a'|_{A'/(b' - a')}) \\
d^A_{A'}(b', b' - a') & = &
\det\nolimits_\kappa(-a'|_{A'/(b')})^{-1}
\det\nolimits_\kappa(b'|_{A'/(b' - a')}) \\
d^A_{A'}(a', b') & = &
\det\nolimits_\kappa(b'|_{A'/(a')})^{-1}
\det\nolimits_\kappa(a'|_{A'/(b')})
\end{matrix}
$$
を得る。したがって
$$
(-1)^{\text{length}_A(A'/(b'))}
d^A_{A'}(a', b' - a')
=
d^A_{A'}(b', b' - a') d^A_{A'}(a', b')
$$
である。符号は上の第二の等式の $-a'$ から生じる。一方、補題
\ref{chow-lemma-periodic-determinant-sign} により
$d^A_{A'}(b', b') = (-1)^{\text{length}_A(A'/(b'))}$ であり、補題が示された。
\end{proof}

\noindent
tame 記号は Steinberg 記号である。

\begin{lemma}
\label{chow-lemma-symbol-is-steinberg}
$A$ を分数体 $K$ をもつ次元 $1$ の Noether 局所整域とする。
$x \in K \setminus \{0, 1\}$ に対し
$$
d_A(x, 1 -x) = 1
$$
である。
\end{lemma}

\begin{proof}
$a, b \in A$ を用いて $x = a/b$ と書く。
$1 - x = (b - a)/b$ なので、仮定から $b - a \not = 0$ も従う。したがって
$$
d_A(x, 1 - x)
=
d_A(a, b - a)d_A(a, b)^{-1}d_A(b, b - a)^{-1}d_A(b, b)
$$
を得る。ゆえに
$d_A(a, b - a) d_A(b, b) = d_A(b, b - a) d_A(a, b)$ を示せばよい。
これは補題 \ref{chow-lemma-symbol-is-steinberg-prepare} である。
\end{proof}










\subsection{長さと行列式}
\label{chow-subsection-length-determinant}

\noindent
この節では、格子を比較するために行列式を用いる。
鍵となる補題は次のものである。

\begin{lemma}
\label{chow-lemma-key-lemma}
$R$ を Noether 局所環とする。
$\mathfrak q \subset R$ を $\dim(R/\mathfrak q) = 1$ を満たす素イデアルとする。
$\varphi : M \to N$ を有限 $R$-加群の準同型とする。
次の性質をもつ $x_1, \ldots, x_l \in M$ および $y_1, \ldots, y_l \in M$
が存在すると仮定する。
\begin{enumerate}
\item $M = \langle x_1, \ldots, x_l\rangle$,
\item $i = 1, \ldots, l$ に対して
$\langle x_1, \ldots, x_i\rangle / \langle x_1, \ldots, x_{i - 1}\rangle
\cong R/\mathfrak q$,
\item $N = \langle y_1, \ldots, y_l\rangle$, および
\item $i = 1, \ldots, l$ に対して
$\langle y_1, \ldots, y_i\rangle / \langle y_1, \ldots, y_{i - 1}\rangle
\cong R/\mathfrak q$.
\end{enumerate}
このとき、$\varphi$ が単射であることと $\varphi_{\mathfrak q}$ が同型であることは同値であり、
この場合
$$
\text{length}_R(\Coker(\varphi)) = \text{ord}_{R/\mathfrak q}(f)
$$
が成り立つ。ここで $f \in \kappa(\mathfrak q)$ は、
$\det_{\kappa(\mathfrak q)}(N_{\mathfrak q})$ において
$$
[\varphi(x_1), \ldots, \varphi(x_l)] = f [y_1, \ldots, y_l]
$$
を満たす元である。
\end{lemma}

\begin{proof}
まず、$l = 1$ の場合に補題が成り立つことに注意する。
実際、この場合 $x_1$ は $R/\mathfrak q$ 上の $M$ の基底であり、
$y_1$ は $R/\mathfrak q$ 上の $N$ の基底であって、ある $f \in R$ に対し
$\varphi(x_1) = fy_1$ である。したがって、$\varphi$ が単射であることと
$f \not \in \mathfrak q$ であることは同値である。さらに
$\Coker(\varphi) = R/(f, \mathfrak q)$ であるから、
$\text{ord}_{R/q}(f)$ の定義
（『可換代数』の定義 \ref{algebra-definition-ord} 参照）により補題が成り立つ。

\medskip\noindent
実際、より一般に、ある $f_i \in R$, $f_i \not \in \mathfrak q$ に対して
$\varphi(x_i) = f_iy_i$ であると仮定する。このとき誘導される写像
$$
\langle x_1, \ldots, x_i\rangle / \langle x_1, \ldots, x_{i - 1}\rangle
\longrightarrow
\langle y_1, \ldots, y_i\rangle / \langle y_1, \ldots, y_{i - 1}\rangle
$$
はすべて単射であり、その余核は $R/(f_i, \mathfrak q)$ と同型である。
したがって
$$
\text{length}_R(\Coker(\varphi)) = \sum \text{ord}_{R/\mathfrak q}(f_i).
$$
を得る。一方、この場合、記号に対する許容関係式 (b) から
$$
[\varphi(x_1), \ldots, \varphi(x_l)] = f_1 \ldots f_l [y_1, \ldots, y_l]
$$
であることは明らかである。よって、この場合にも結論が成り立つ。

\medskip\noindent
一般の場合を $l$ に関する帰納法で証明する。$l > 1$ と仮定する。
$\varphi(x_1) \in \langle y_1, \ldots, y_i\rangle$ を満たす最小の
$i \in \{1, \ldots, l\}$ をとる。$i$ に関する帰納法で論じる。
$i = 1$ ならば、可換図式
$$
\xymatrix{
0 \ar[r] &
\langle x_1 \rangle \ar[r] \ar[d] &
\langle x_1, \ldots, x_l \rangle \ar[r] \ar[d] &
\langle x_1, \ldots, x_l \rangle / \langle x_1 \rangle \ar[r] \ar[d] &
0 \\
0 \ar[r] &
\langle y_1 \rangle \ar[r] &
\langle y_1, \ldots, y_l \rangle \ar[r] &
\langle y_1, \ldots, y_l \rangle / \langle y_1 \rangle \ar[r] &
0
}
$$
を得るので、蛇の補題と $l$ に関する帰納法から補題が従う。
以下 $i > 1$ と仮定する。
$a_j, a \in R$ および $a \not \in \mathfrak q$ を用いて
$\varphi(x_1) = a_1 y_1 + \ldots + a_{i - 1} y_{i - 1} + a y_i$
と書く（そうでなければ $i$ は最小でない）。次のようにおく。
$$
x'_j =
\left\{
\begin{matrix}
x_j & \text{if} & j = 1 \\
ax_j & \text{if} & j \geq 2
\end{matrix}
\right.
\quad\text{and}\quad
y'_j =
\left\{
\begin{matrix}
y_j & \text{if} & j < i \\
ay_j & \text{if} & j \geq i
\end{matrix}
\right.
$$
$M' = \langle x'_1, \ldots, x'_l \rangle$ および
$N' = \langle y'_1, \ldots, y'_l \rangle$ とおく。
構成により
$\varphi(x'_1) = a_1 y'_1 + \ldots + a_{i - 1} y'_{i - 1} + y'_i$ であり、
$j > 1$ に対して
$\varphi(x'_j) = a\varphi(x_i) \in \langle y'_1, \ldots, y'_l\rangle$
であるから、$R$-加群と写像の可換図式
$$
\xymatrix{
M' \ar[d] \ar[r]_{\varphi'} & N' \ar[d] \\
M \ar[r]^\varphi & N
}
$$
を得る。証明の第 2 段落の結果により
$\text{length}_R(M/M') = (l - 1)\text{ord}_{R/\mathfrak q}(a)$
であり、同様に
$\text{length}_R(M/M') = (l - i + 1)\text{ord}_{R/\mathfrak q}(a)$
であることが分かる。図式追跡により、これは
$$
\text{length}_R(\Coker(\varphi')) =
\text{length}_R(\Coker(\varphi)) + i\ \text{ord}_{R/\mathfrak q}(a).
$$
を含意する。一方、
$$
[\varphi(x_1), \ldots, \varphi(x_l)] = f [y_1, \ldots, y_l],
\quad
[\varphi'(x'_1), \ldots, \varphi(x'_l)] = f' [y'_1, \ldots, y'_l]
$$
と書けば、$f' = a^if$ であることは明らかである。したがって、次の場合に補題を証明すれば十分である。
$\varphi(x_1) = a_1y_1 + \ldots a_{i - 1}y_{i - 1} + y_i$ の場合、
すなわち $a = 1$ の場合である。次に、
$$
[y_1, \ldots, y_l] = [y_1, \ldots, y_{i - 1},
a_1y_1 + \ldots a_{i - 1}y_{i - 1} + y_i, y_{i + 1}, \ldots, y_l]
$$
であることを思い出そう。これは記号に対する許容関係式から従う。列
$y_1, \ldots, y_{i - 1},
a_1y_1 + \ldots + a_{i - 1}y_{i - 1} + y_i, y_{i + 1}, \ldots, y_l$
も補題の条件 (3), (4) を満たす。したがって実際には
$\varphi(x_1) = y_i$ と仮定してよい。この場合、
$\mathfrak q x_1 = 0$ であり、したがって $\mathfrak q y_i = 0$ でもあることに注意する。
また
$$
[y_1, \ldots, y_l] =
- [y_1, \ldots, y_{i - 2}, y_i, y_{i - 1}, y_{i + 1}, \ldots, y_l]
$$
である。これは $\det_{\kappa(\mathfrak q)}(N_{\mathfrak q})$ を定義する
許容関係式の第 3 のものによる。したがって $y_1, \ldots, y_l$ を、列
$y'_1, \ldots, y'_l =
y_1, \ldots, y_{i - 2}, y_i, y_{i - 1}, y_{i + 1}, \ldots, y_l$
で置き換えてよい（この列も補題の条件 (3), (4) を満たす）。
これにより不変量 $i$ は明らかに $1$ 減少するから、
$i$ に関する帰納法により結論を得る。
\end{proof}

\noindent
前の補題を用いるため、必要な性質をもつ元の列がしばしば存在することを示す。

\begin{lemma}
\label{chow-lemma-good-sequence-exists}
$R$ を Noether 局所環とする。
$\mathfrak q \subset R$ を素イデアルとする。
$M$ を有限 $R$-加群とし、$\mathfrak q$ が $M$ の台の極小素イデアルの
一つであるとする。このとき、次を満たす $x_1, \ldots, x_l \in M$ が存在する。
\begin{enumerate}
\item $M / \langle x_1, \ldots, x_l\rangle$ の台は $\mathfrak q$ を含まない。
\item $i = 1, \ldots, l$ に対して
$\langle x_1, \ldots, x_i\rangle / \langle x_1, \ldots, x_{i - 1}\rangle
\cong R/\mathfrak q$ である。
\end{enumerate}
さらに、この場合 $l = \text{length}_{R_\mathfrak q}(M_\mathfrak q)$ である。
\end{lemma}

\begin{proof}
$\mathfrak q$ が $M$ の台の極小素イデアルであるという条件から、
$l = \text{length}_{R_\mathfrak q}(M_\mathfrak q)$ は有限である
（『可換代数』の補題 \ref{algebra-lemma-support-point} 参照）。
したがって、$i = 1, \ldots, l$ に対して
$\langle y_1, \ldots, y_i\rangle / \langle y_1, \ldots, y_{i - 1}\rangle
\cong \kappa(\mathfrak q)$
となる $y_1, \ldots, y_l \in M_{\mathfrak q}$ をとれる。
ある $z_i \in M$ の像が $f_i y_i$ となるような
$f_i \in R$, $f_i \not \in \mathfrak q$ をとれる。
さらに $R$ は Noether 環なので、ある $g_j \in R$ により
$\mathfrak q = (g_1, \ldots, g_t)$ と書ける。
仮定により、加群 $M_{\mathfrak q}$ の中で
$g_j y_i \in \langle y_1, \ldots, y_{i - 1} \rangle$ である。
$z_i$ の選び方により、さらに
$f_{ji} \in R$, $f_{ij} \not \in \mathfrak q$ であって
$f_{ij} g_j z_i \in \langle z_1, \ldots, z_{i - 1} \rangle$
となるものをとれる（等式は加群 $M$ の中で考える）。
補題は
$$
x_1 = f_{11}f_{12}\ldots f_{1t}z_1,
\quad
x_2 = f_{11}f_{12}\ldots f_{1t}f_{21}f_{22}\ldots f_{2t}z_2,
$$
などととることから従う。実際、すべての元 $f_i, f_{ij}$ は
$R_{\mathfrak q}$ において可逆なので、なお
$R_{\mathfrak q}x_1 + \ldots + R_{\mathfrak q}x_i /
R_{\mathfrak q}x_1 + \ldots + R_{\mathfrak q}x_{i - 1}
\cong \kappa(\mathfrak q)$ が $i = 1, \ldots, l$ に対して成り立つ。
構成により
$\mathfrak q x_i \in \langle x_1, \ldots, x_{i - 1}\rangle$ である。
したがって
$\langle x_1, \ldots, x_i\rangle / \langle x_1, \ldots, x_{i - 1}\rangle$
は一つの元で生成され、$\mathfrak q$ により零化される $R$-加群であって、
$\mathfrak q$ で局所化すると $\kappa(\mathfrak q)$ 上の $q$-次元ベクトル空間となる。
ゆえにこれは $R/\mathfrak q$ と同型である。
\end{proof}

\noindent
この節の主結果は次のものである。
以下では、この命題から得られるさまざまな帰結を見る。
まず、より容易な補題 \ref{chow-lemma-application-herbrand-quotient} を
読者自身で証明することを勧める。

\begin{proposition}
\label{chow-proposition-length-determinant-periodic-complex}
$R$ を剰余体 $\kappa$ をもつ Noether 局所環とする。
$(M, \varphi, \psi)$ を $R$ 上の $(2, 1)$-周期複体とする。
次を仮定する。
\begin{enumerate}
\item $M$ は有限 $R$-加群である。
\item $(M, \varphi, \psi)$ のコホモロジー加群は有限長である。
\item $\dim(\text{Supp}(M)) = 1$ である。
\end{enumerate}
$\mathfrak q_i$, $i = 1, \ldots, t$ を $M$ の台の極小素イデアルとする。
このとき\footnote{
もちろん、$\det_\kappa(M, \varphi, \psi)$ を現在の値の逆元として
再定義すれば負号を除ける。定義 \ref{chow-definition-periodic-determinant} を参照。}
$$
- e_R(M, \varphi, \psi) =
\sum\nolimits_{i = 1, \ldots, t}
\text{ord}_{R/\mathfrak q_i}\left(
\det\nolimits_{\kappa(\mathfrak q_i)}
(M_{\mathfrak q_i}, \varphi_{\mathfrak q_i}, \psi_{\mathfrak q_i})
\right)
$$
が成り立つ。
\end{proposition}

\begin{proof}
まず、次のようにして $t = 1$ の場合に帰着する。
$\text{Supp}(M) = \{\mathfrak m, \mathfrak q_1, \ldots, \mathfrak q_t\}$,
であることに注意する。ここで $\mathfrak m \subset R$ は極大イデアルである。
$M_i$ を $M \to M_{\mathfrak q_i}$ の像と表すと、
$\text{Supp}(M_i) = \{\mathfrak m, \mathfrak q_i\}$ である。
写像 $\varphi$（resp.\ $\psi$）は $R$-加群の写像
$\varphi_i : M_i \to M_i$（resp.\ $\psi_i : M_i \to M_i$）を誘導する。
したがって、$(2, 1)$-周期複体の射
$$
(M, \varphi, \psi) \longrightarrow
\bigoplus\nolimits_{i = 1, \ldots, t} (M_i, \varphi_i, \psi_i).
$$
を得る。この写像の核と余核の台は $\{\mathfrak m\}$ に含まれる。
したがって、補題 \ref{chow-lemma-compare-periodic-lengths} により
$$
e_R(M, \varphi, \psi) =
\sum\nolimits_{i = 1, \ldots, t}
e_R(M_i, \varphi_i, \psi_i)
$$
である。一方、明らかに $M_{\mathfrak q_i} = M_{i, \mathfrak q_i}$ であるから、
補題の公式の右辺の各項は
$$
\text{ord}_{R/\mathfrak q_i}\left(
\det\nolimits_{\kappa(\mathfrak q_i)}
(M_{i, \mathfrak q_i}, \varphi_{i, \mathfrak q_i}, \psi_{i, \mathfrak q_i})
\right)
$$
に等しい。言い換えると、各加群 $M_i$ について補題を証明できれば、
補題が成り立つ。これで $t = 1$ の場合に帰着された。

\medskip\noindent
$M$ が有限 $R$-加群であり、
$\text{Supp}(M) = \{\mathfrak m, \mathfrak q\}$ で、コホモロジー加群が有限長であるような
Noether 局所環上の $(2, 1)$-周期複体 $(M, \varphi, \psi)$ を考える。
この場合の証明は、補題 \ref{chow-lemma-key-lemma} と注意深い帳尻合わせから従う。
$K_\varphi = \Ker(\varphi)$,
$I_\varphi = \Im(\varphi)$,
$K_\psi = \Ker(\psi)$, および
$I_\psi = \Im(\psi)$
と表す。$R$ は Noether 環なので、これらはすべて有限 $R$-加群である。次のようにおく。
$$
a = \text{length}_{R_{\mathfrak q}}(I_{\varphi, \mathfrak q})
= \text{length}_{R_{\mathfrak q}}(K_{\psi, \mathfrak q}),
\quad
b = \text{length}_{R_{\mathfrak q}}(I_{\psi, \mathfrak q})
= \text{length}_{R_{\mathfrak q}}(K_{\varphi, \mathfrak q}).
$$
これらの等号は、$\mathfrak q$ で局所化すると複体が完全になることによる。
$l = \text{length}_{R_{\mathfrak q}}(M_{\mathfrak q})$ は
$l = a + b$ に等しいことに注意する。

\medskip\noindent
台が $\{\mathfrak m, \mathfrak q\}$ に含まれる有限 $R$-加群 $N$ において
元の列を選ぶため、補題 \ref{chow-lemma-good-sequence-exists} を用いる。
この場合 $N_{\mathfrak q}$ は有限長であり、その長さを $n \in \mathbf{N}$ とする。
補題 \ref{chow-lemma-good-sequence-exists} の性質 (1), (2) をもつ列
$w_1, \ldots, w_n \in N$ を「良い列」と呼ぶことにする。
良い列が生成する部分加群による $N$ の商
$N/\langle w_1, \ldots, w_n \rangle$ の台は $\{\mathfrak m\}$ に含まれるので、
有限長である（『可換代数』の補題 \ref{algebra-lemma-support-point}）。
さらに、記号
$[w_1, \ldots, w_n] \in \det_{\kappa(\mathfrak q)}(N_{\mathfrak q})$
は生成元である。補題 \ref{chow-lemma-determinant-dimension-one} を参照。

\medskip\noindent
以上を踏まえ、次の良い列を選ぶ。
$$
\begin{matrix}
x_1, \ldots, x_b & \text{in} & K_\varphi, &
t_1, \ldots, t_a & \text{in} & K_\psi, \\
y_1, \ldots, y_a & \text{in} & I_\varphi \cap \langle t_1, \ldots t_a\rangle, &
s_1, \ldots, s_b & \text{in} & I_\psi \cap \langle x_1, \ldots, x_b\rangle.
\end{matrix}
$$
この選択を次のように少し調整する。
$y_i \in I_\varphi$ の持ち上げ $\tilde y_i \in M$ と、
$s_i \in I_\psi$ の持ち上げ $\tilde s_i \in M$ を選ぶ。
$\mathfrak q \tilde y_1 \subset \langle x_1, \ldots, x_b\rangle$
であるとは限らず、
$\mathfrak q \tilde s_1 \subset \langle t_1, \ldots, t_a\rangle$
であるとも限らない。しかし、$\mathfrak q$ が有限生成であることを用いると
（補題 \ref{chow-lemma-good-sequence-exists} の証明と同様に）、
$d \in R$, $d \not \in \mathfrak q$ であって
$\mathfrak q d\tilde y_1 \subset \langle x_1, \ldots, x_b\rangle$
かつ
$\mathfrak q d\tilde s_1 \subset \langle t_1, \ldots, t_a\rangle$
となるものをとれる。したがって $y_i$ を $dy_i$ で、
$\tilde y_i$ を $d\tilde y_i$ で、$s_i$ を $ds_i$ で、$\tilde s_i$ を
$d\tilde s_i$ で置き換えることにより、
$x_1, \ldots, x_b, \tilde y_1, \ldots, \tilde y_b$
および $t_1, \ldots, t_a, \tilde s_1, \ldots, \tilde s_b$
も $M$ の良い列であると仮定してよい。

\medskip\noindent
最後に、有限 $R$-加群
$$
\langle
x_1, \ldots, x_b, \tilde y_1, \ldots, \tilde y_a
\rangle
\cap
\langle
t_1, \ldots, t_a, \tilde s_1, \ldots, \tilde s_b
\rangle.
$$
の良い列 $z_1, \ldots, z_l$ を選ぶ。
これは $M$ においても良い列であることに注意する。

\medskip\noindent
$I_{\varphi, \mathfrak q} = K_{\psi, \mathfrak q}$ なので、
$\det_{\kappa(\mathfrak q)}(K_{\psi, \mathfrak q})$ の中で
$[y_1, \ldots, y_a] = h [t_1, \ldots, t_a]$
となる一意な元 $h \in \kappa(\mathfrak q)$ が存在する。
同様に $I_{\psi, \mathfrak q} = K_{\varphi, \mathfrak q}$ なので、
$\det_{\kappa(\mathfrak q)}(K_{\varphi, \mathfrak q})$ の中で
$[s_1, \ldots, s_b] = g [x_1, \ldots, x_b]$
となる一意な元 $h \in \kappa(\mathfrak q)$ が存在する。
$M$ にある三つの良い列についても同様にできる。合わせて、ある
$g, h, f_\varphi, f_\psi \in \kappa(\mathfrak q)$ に対して
次の恒等式を得る。
\begin{align*}
[y_1, \ldots, y_a]
& =
h [t_1, \ldots, t_a] \\
[s_1, \ldots, s_b]
& =
g [x_1, \ldots, x_b] \\
[z_1, \ldots, z_l]
& =
f_\varphi [x_1, \ldots, x_b, \tilde y_1, \ldots, \tilde y_a] \\
[z_1, \ldots, z_l]
& =
f_\psi [t_1, \ldots, t_a, \tilde s_1, \ldots, \tilde s_b].
\end{align*}

\medskip\noindent
以上の記法を準備したので、
$\det_{\kappa(\mathfrak q)}(M, \varphi, \psi)$ を計算しよう。
すなわち、元 $[z_1, \ldots, z_l]$ を考える。
定義 \ref{chow-definition-periodic-determinant} の写像
$\gamma_\psi \circ \sigma \circ \gamma_\varphi^{-1}$ のもとで
\begin{eqnarray*}
[z_1, \ldots, z_l] & = &
f_\varphi [x_1, \ldots, x_b, \tilde y_1, \ldots, \tilde y_a] \\
& \mapsto & f_\varphi [x_1, \ldots, x_b] \otimes [y_1, \ldots, y_a] \\
& \mapsto &
f_\varphi h/g [t_1, \ldots, t_a] \otimes [s_1, \ldots, s_b] \\
& \mapsto &
f_\varphi h/g [t_1, \ldots, t_a, \tilde s_1, \ldots, \tilde s_b] \\
& = &
f_\varphi h/f_\psi g [z_1, \ldots, z_l]
\end{eqnarray*}
となる。これは
$\det_{\kappa(\mathfrak q)}
(M_{\mathfrak q}, \varphi_{\mathfrak q}, \psi_{\mathfrak q})$
が符号を除いて $f_\varphi h/f_\psi g$ に等しいことを意味する。

\medskip\noindent
次の量を略記する。
\begin{eqnarray*}
k_\varphi & = & \text{length}_R(K_\varphi/\langle x_1, \ldots, x_b\rangle) \\
k_\psi    & = & \text{length}_R(K_\psi/\langle t_1, \ldots, t_a\rangle) \\
i_\varphi & = & \text{length}_R(I_\varphi/\langle y_1, \ldots, y_a\rangle) \\
i_\psi    & = & \text{length}_R(I_\psi/\langle s_1, \ldots, s_a\rangle) \\
m_\varphi & = & \text{length}_R(M/
\langle x_1, \ldots, x_b, \tilde y_1, \ldots, \tilde y_a\rangle) \\
m_\psi    & = & \text{length}_R(M/
\langle t_1, \ldots, t_a, \tilde s_1, \ldots, \tilde s_b\rangle) \\
\delta_\varphi & = & \text{length}_R(
\langle x_1, \ldots, x_b, \tilde y_1, \ldots, \tilde y_a\rangle
\langle z_1, \ldots, z_l\rangle) \\
\delta_\psi & = & \text{length}_R(
\langle t_1, \ldots, t_a, \tilde s_1, \ldots, \tilde s_b\rangle
\langle z_1, \ldots, z_l\rangle)
\end{eqnarray*}
完全列 $0 \to K_\varphi \to M \to I_\varphi \to 0$ を用いると
$m_\varphi = k_\varphi + i_\varphi$ を得る。同様に
$m_\psi = k_\psi + i_\psi$ である。
$\delta_\varphi + m_\varphi = \delta_\psi + m_\psi$ も成り立つ。
実際、これは $M$ における $\langle z_1, \ldots, z_l \rangle$ の余長に等しい。
最後に、鍵補題 \ref{chow-lemma-key-lemma} の最初の適用により
$$
\delta_\varphi = \text{ord}_{R/\mathfrak q}(f_\varphi),
\quad
\delta_\psi = \text{ord}_{R/\mathfrak q}(f_\psi)
$$
である。

\medskip\noindent
次に、周期複体の重複度を計算する。
\begin{eqnarray*}
e_R(M, \varphi, \psi) & = &
\text{length}_R(K_\varphi/I_\psi) - \text{length}_R(K_\psi/I_\varphi) \\
& = &
\text{length}_R(
\langle x_1, \ldots, x_b\rangle/
\langle s_1, \ldots, s_b\rangle)
+ k_\varphi - i_\psi \\
& & -
\text{length}_R(
\langle t_1, \ldots, t_a\rangle/
\langle y_1, \ldots, y_a\rangle)
- k_\psi + i_\varphi \\
& = &
\text{ord}_{R/\mathfrak q}(g/h) + k_\varphi - i_\psi - k_\psi + i_\varphi \\
& = &
\text{ord}_{R/\mathfrak q}(g/h) + m_\varphi - m_\psi \\
& = &
\text{ord}_{R/\mathfrak q}(g/h) + \delta_\psi - \delta_\varphi \\
& = &
\text{ord}_{R/\mathfrak q}(f_\psi g/f_\varphi h)
\end{eqnarray*}
ここで第 3 の等式では鍵補題 \ref{chow-lemma-key-lemma} を 2 回用いた。
$\det_{\kappa(\mathfrak q)}
(M_{\mathfrak q}, \varphi_{\mathfrak q}, \psi_{\mathfrak q})$
についての上の計算により、命題が証明された。
\end{proof}

\noindent
ほとんどの応用では、次の補題で十分である。

\begin{lemma}
\label{chow-lemma-application-herbrand-quotient}
$R$ を極大イデアル $\mathfrak m$ をもつ Noether 局所環とする。
$M$ を有限 $R$-加群とし、$\psi : M \to M$ を $R$-加群の写像とする。
次を仮定する。
\begin{enumerate}
\item $\Ker(\psi)$ と $\Coker(\psi)$ は有限長である。
\item $\dim(\text{Supp}(M)) \leq 1$ である。
\end{enumerate}
$\text{Supp}(M) = \{\mathfrak m, \mathfrak q_1, \ldots, \mathfrak q_t\}$
と書き、
$\det_{\kappa(\mathfrak q_i)}(\psi_{\mathfrak q_i}) :
\det_{\kappa(\mathfrak q_i)}(M_{\mathfrak q_i})
\to \det_{\kappa(\mathfrak q_i)}(M_{\mathfrak q_i})$
が $f_i$ 倍であるような元を $f_i \in \kappa(\mathfrak q_i)^*$ と表す。
このとき
$$
\text{length}_R(\Coker(\psi))
-
\text{length}_R(\Ker(\psi))
=
\sum\nolimits_{i = 1, \ldots, t}
\text{ord}_{R/\mathfrak q_i}(f_i).
$$
が成り立つ。
\end{lemma}

\begin{proof}
$H^0(M, 0, \psi) = \Coker(\psi)$ および
$H^1(M, 0, \psi) = \Ker(\psi)$ であることを思い出そう。
定義 \ref{chow-definition-periodic-length} の前の注記を参照。
命題 \ref{chow-proposition-length-determinant-periodic-complex} と
補題 \ref{chow-lemma-periodic-determinant-easy-case} を組み合わせれば補題が従う。

\medskip\noindent
別証明。命題 \ref{chow-proposition-length-determinant-periodic-complex} の証明と同様に、
$\text{Supp}(M) = \{\mathfrak m, \mathfrak q\}$ の場合に帰着する。
次に、補題 \ref{chow-lemma-key-lemma} と
\ref{chow-lemma-good-sequence-exists} を直接組み合わせて、
命題 \ref{chow-proposition-length-determinant-periodic-complex} のこの特別な場合を証明する。
この場合は帳尻合わせがはるかに少なくて済むので、読者が詳細を確かめることを勧める。
詳細は省略する。
\end{proof}




\subsection{鍵補題への応用}
\label{chow-subsection-application-tame-symbol}

\noindent
この節では上の結果を適用し、上で構成した tame 記号 $d_A$ に対する
鍵補題（補題 \ref{chow-lemma-milnor-gersten-low-degree}）の類似を示す。
Milnor $K$-理論との関係については
注意 \ref{chow-remark-gersten-complex-milnor} を参照されたい。

\begin{lemma}[鍵補題]
\label{chow-lemma-secondary-ramification}
\begin{reference}
$A$ が excellent 環であるとき、これは \cite[Proposition 1]{Kato-Milnor-K} である。
\end{reference}
$A$ を分数体 $K$ をもつ $2$-次元 Noether 局所整域とする。
$f, g \in K^*$ とする。
$\mathfrak q_1, \ldots, \mathfrak q_t$ を、$f$ または $g$ の少なくとも一方が
$A^*_{\mathfrak q}$ の元でないような $A$ の高さ $1$ の素イデアル
$\mathfrak q$ とする。このとき
$$
\sum\nolimits_{i = 1, \ldots, t}
\text{ord}_{A/\mathfrak q_i}(d_{A_{\mathfrak q_i}}(f, g))
=
0
$$
が成り立つ。これは
$$
\sum\nolimits_{\text{height}(\mathfrak q) = 1}
\text{ord}_{A/\mathfrak q}(d_{A_{\mathfrak q}}(f, g))
=
0
$$
とも書ける。実際、$f, g \in A^*_{\mathfrak q}$ となる $A$ の任意の
高さ 1 の素イデアル $\mathfrak q$ に対して、補題
\ref{chow-lemma-symbol-when-one-is-a-unit} により
$d_{A_{\mathfrak q}}(f, g) = 1$ である。
\end{lemma}

\begin{proof}
tame 記号 $d_{A_{\mathfrak q}}(f, g)$ は加法的であり
（補題 \ref{chow-lemma-multiplicativity-symbol}）、
位数関数 $\text{ord}_{A/\mathfrak q}$ も加法的なので
（『可換代数』の補題 \ref{algebra-lemma-ord-additive}）、
$f = a \in A$ および $g = b \in A$ の場合に公式を証明すれば十分である。
この場合、示すべきことは
$$
\sum\nolimits_{\text{height}(\mathfrak q) = 1}
\text{ord}_{A/\mathfrak q}(\det\nolimits_\kappa(A_{\mathfrak q}/(ab), a, b))
= 0
$$
である。命題 \ref{chow-proposition-length-determinant-periodic-complex} により、
これは
$$
e_A(A/(ab), a, b) = 0.
$$
を示すことと同値である。複体
$A/(ab) \xrightarrow{a} A/(ab) \xrightarrow{b} A/(ab) \xrightarrow{a} A/(ab)$
は完全なので、結論を得る。
\end{proof}





\section{付録 B：別のアプローチ}
\label{chow-section-appendix-chow}

\noindent
この付録ではまず、本文の内容を連接層の $K$-理論と結び付けることを簡潔に試みる。
特に、可逆加群の $c_1$ とのカップ積が、この可逆加群とのテンソル積に
どのように関係するかを述べる。補題 \ref{chow-lemma-coherent-sheaf-cap-c1} を参照。
この内容は明らかに非常に興味深く、より詳細かつ広範な解説に値する。


\subsection{有理同値と K-群}
\label{chow-subsection-rational-equivalence-K-groups}

\noindent
この節は第 \ref{chow-section-chow-and-K} 節の続きである。
次の補題の動機は
『ホモロジー代数』の補題 \ref{homology-lemma-serre-subcategory-K-groups} である。

\begin{lemma}
\label{chow-lemma-maps-between-coherent-sheaves}
$(S, \delta)$ を状況 \ref{chow-situation-setup} のとおりとする。
$X$ を $S$ 上局所有限型なスキームとする。
$\mathcal{F}$ を $X$ 上の連接層とする。
$$
\xymatrix{
\ldots \ar[r] &
\mathcal{F} \ar[r]^\varphi &
\mathcal{F} \ar[r]^\psi &
\mathcal{F} \ar[r]^\varphi &
\mathcal{F} \ar[r] & \ldots
}
$$
を Homology, Equation (\ref{homology-equation-cyclic-complex}) のような複体とする。
次を仮定する。
\begin{enumerate}
\item $\dim_\delta(\text{Supp}(\mathcal{F})) \leq k + 1$.
\item $i = 0, 1$ に対して
$\dim_\delta(\text{Supp}(H^i(\mathcal{F}, \varphi, \psi))) \leq k$.
\end{enumerate}
このとき、$X$ 上の $k$-サイクルとして
$$
[H^0(\mathcal{F}, \varphi, \psi)]_k
\sim_{rat}
[H^1(\mathcal{F}, \varphi, \psi)]_k
$$
が成り立つ。
\end{lemma}

\begin{proof}
$\text{Supp}(\mathcal{F})$ の $\delta$-次元 $k + 1$ の既約成分全体を
$\{W_j\}_{j \in J}$ とする。補題 \ref{chow-lemma-length-finite} により、
$\{W_j\}$ は $X$ の閉部分集合からなる局所有限族であることに注意する。
各 $j$ に対し、$\xi_j \in W_j$ を生成点とする。次のようにおく。
$$
f_j = \det\nolimits_{\kappa(\xi_j)}
(\mathcal{F}_{\xi_j}, \varphi_{\xi_j}, \psi_{\xi_j})
\in
R(W_j)^*.
$$
記法については定義 \ref{chow-definition-periodic-determinant} を参照。
次を主張する。
$$
- [H^0(\mathcal{F}, \varphi, \psi)]_k + [H^1(\mathcal{F}, \varphi, \psi)]_k
=
\sum (W_j \to X)_*\text{div}(f_j)
$$
これを証明すれば補題が従う。

\medskip\noindent
$Z \subset X$ を $\delta$-次元 $k$ の整閉部分スキームとする。
上の等式を証明するには、
$
[H^0(\mathcal{F}, \varphi, \psi)]_k - [H^1(\mathcal{F}, \varphi, \psi)]_k
$
における $[Z]$ の係数 $n$ が、
$
\sum (W_j \to X)_*\text{div}(f_j)
$
における $[Z]$ の係数 $m$ と等しいことを示せば十分である。
$\xi \in Z$ を生成点とする。
局所環 $A = \mathcal{O}_{X, \xi}$ を考え、
$M = \mathcal{F}_\xi$ を $A$-加群とみなす。
茎上の $\varphi, \psi$ の作用を
$\varphi, \psi : M \to M$ と表す。
$\xi \in Z$ の選び方により $\delta(\xi) = k$ であり、
したがって $\dim(\text{Supp}(M)) = 1$ である。
最後に、$\xi$ を通る整閉部分スキーム $W_j$ は、
$\text{Supp}(M)$ の極小素イデアル $\mathfrak q_i$ に対応する。
各場合に、元 $f_j \in R(W_j)^*$ は $\kappa(\mathfrak q_i)^*$ の元
$\det_{\kappa(\mathfrak q_i)}(M_{\mathfrak q_i}, \varphi, \psi)$ に対応する。
したがって
$$
n = - e_A(M, \varphi, \psi)
$$
および
$$
m =
\sum
\text{ord}_{A/\mathfrak q_i}
(\det\nolimits_{\kappa(\mathfrak q_i)}(M_{\mathfrak q_i}, \varphi, \psi))
$$
を得る。ゆえに命題 \ref{chow-proposition-length-determinant-periodic-complex}
から結論が従う。
\end{proof}

\begin{lemma}
\label{chow-lemma-cycles-rational-equivalence-K-group}
$(S, \delta)$ を状況 \ref{chow-situation-setup} のとおりとする。
$X$ を $S$ 上局所有限型なスキームとする。
補題 \ref{chow-lemma-from-chow-to-K} の写像
$$
\CH_k(X) \longrightarrow
K_0(\textit{Coh}_{\leq k + 1}(X)/\textit{Coh}_{\leq k - 1}(X))
$$
は、$\CH_k(X)$ から写像
$$
K_0(\textit{Coh}_{\leq k}(X)/\textit{Coh}_{\leq k - 1}(X))
\longrightarrow
K_0(\textit{Coh}_{\leq k + 1}(X)/\textit{Coh}_{\leq k - 1}(X)).
$$
の像 $B_k(X)$ への全単射を誘導する。
\end{lemma}

\begin{proof}
補題 \ref{chow-lemma-cycles-k-group} により、
補題 \ref{chow-lemma-from-chow-to-K} の写像と両立する等式
$Z_k(X) = K_0(\textit{Coh}_{\leq k}(X)/\textit{Coh}_{\leq k - 1}(X))$
がある。そこで、
$K_0(\textit{Coh}_{\leq k}(X)/\textit{Coh}_{\leq k - 1}(X))$ の元
$[A] - [B]$ が $B_k(X)$ において零に写る、すなわち
$K_0(\textit{Coh}_{\leq k + 1}(X)/\textit{Coh}_{\leq k - 1}(X))$
において零に写ると仮定する。
$[A] - [B]$ が $X$ 上で零に有理同値なサイクルに対応することを示さなければならない。
$X$ 上の、台の $\delta$-次元が $\leq k$ である連接層
$\mathcal{A}, \mathcal{B}$ に対して
$[A] = [\mathcal{A}]$ および $[B] = [\mathcal{B}]$ であると仮定する。
$[A] - [B]$ が群
$K_0(\textit{Coh}_{\leq k + 1}(X)/\textit{Coh}_{\leq k - 1}(X))$
において零に写るという仮定は、台の $\delta$-次元が $\leq k - 1$ である
$X$ 上の連接層 $\mathcal{A}', \mathcal{B}'$ であって、
$[\mathcal{A} \oplus \mathcal{A}'] - [\mathcal{B} \oplus \mathcal{B}']$
が $K_0(\textit{Coh}_{k + 1}(X))$ において零となるものが存在することを意味する
（『ホモロジー代数』の補題 \ref{homology-lemma-serre-subcategory-K-groups} の (1) を用いる）。
『ホモロジー代数』の補題 \ref{homology-lemma-serre-subcategory-K-groups} の (2) により、
これは圏 $\textit{Coh}_{\leq k + 1}(X)$ の $(2, 1)$-周期複体
$(\mathcal{F}, \varphi, \psi)$ であって
$\mathcal{A} \oplus \mathcal{A}' = H^0(\mathcal{F}, \varphi, \psi)$
および $\mathcal{B} \oplus \mathcal{B}' = H^1(\mathcal{F}, \varphi, \psi)$
を満たすものが存在することを意味する。
補題 \ref{chow-lemma-maps-between-coherent-sheaves} により、これは
$$
[\mathcal{A} \oplus \mathcal{A}']_k
\sim_{rat}
[\mathcal{B} \oplus \mathcal{B}']_k
$$
を含意する。したがって、補題 \ref{chow-lemma-cycles-k-group} の写像
$$
K_0(\textit{Coh}_{\leq k}(X)/\textit{Coh}_{\leq k - 1}(X))
\longrightarrow Z_k(X)
$$
により、$[A] - [B]$ は零に有理同値なサイクルに写る。
これが示すべきことであり、証明は完了する。
\end{proof}














\subsection{Cartier 因子と K-群}
\label{chow-subsection-cartier-coherent}

\noindent
この節では、可逆層 $\mathcal{L}$ の第 1 Chern 類との交叉が、
$K$-群における $\mathcal{L} - \mathcal{O}$ とのテンソル積に
どのように対応するかを述べる。

\begin{lemma}
\label{chow-lemma-no-embedded-points-modules}
$A$ を Noether 局所環とする。
$M$ を有限 $A$-加群とする。
$a, b \in A$ とする。
次を仮定する。
\begin{enumerate}
\item $\dim(A) = 1$,
\item $a$ と $b$ はともに $A$ の非零因子である。
\item $A$ は埋め込まれた素イデアルをもたない。
\item $M$ は埋め込まれた随伴素イデアルをもたない。
\item $\text{Supp}(M) = \Spec(A)$.
\end{enumerate}
$I = \{x \in A \mid x(a/b) \in A\}$ とする。
$\mathfrak q_1, \ldots, \mathfrak q_t$ を $A$ の極小素イデアルとする。
このとき $(a/b)IM \subset M$ であり、
$$
\text{length}_A(M/(a/b)IM)
-
\text{length}_A(M/IM)
=
\sum\nolimits_i
\text{length}_{A_{\mathfrak q_i}}(M_{\mathfrak q_i})
\text{ord}_{A/\mathfrak q_i}(a/b)
$$
が成り立つ。
\end{lemma}

\begin{proof}
$M$ は埋め込まれた随伴素イデアルをもたず、$M$ の台は $\Spec(A)$ なので、
$\text{Ass}(M) = \{\mathfrak q_1, \ldots, \mathfrak q_t\}$ である。
したがって $a$, $b$ は $M$ 上の非零因子である。次に注意する。
\begin{align*}
& \text{length}_A(M/(a/b)IM) \\
& = \text{length}_A(bM/aIM) \\
& = \text{length}_A(M/aIM)
-
\text{length}_A(M/bM) \\
& = \text{length}_A(M/aM) + \text{length}_A(aM/aIM) - \text{length}_A(M/bM) \\
& = \text{length}_A(M/aM) + \text{length}_A(M/IM) - \text{length}_A(M/bM)
\end{align*}
実際、単射 $b : M \to bM$ は $(a/b)IM$ を $aIM$ に写し、
単射 $a : M \to aM$ は $IM$ を $aIM$ に写す。
したがって補題の等式の左辺は
$$
\text{length}_A(M/aM) - \text{length}_A(M/bM).
$$
に等しい。補題 \ref{chow-lemma-additivity-divisors-restricted} の第 2 の公式を
それぞれ $x = a, b$ に適用し、『可換代数』の定義 \ref{algebra-definition-ord} の $\text{ord}$-関数の定義を用いれば結論を得る。
\end{proof}

\begin{lemma}
\label{chow-lemma-no-embedded-points}
$(S, \delta)$ を状況 \ref{chow-situation-setup} のとおりとする。
$X$ を $S$ 上局所有限型とする。
$\mathcal{L}$ を可逆 $\mathcal{O}_X$-加群とする。
$\mathcal{F}$ を連接 $\mathcal{O}_X$-加群とする。
$s \in \Gamma(X, \mathcal{K}_X(\mathcal{L}))$ を
$\mathcal{L}$ の有理型切断とする。
次を仮定する。
\begin{enumerate}
\item $\dim_\delta(X) \leq k + 1$,
\item $X$ は埋め込まれた点をもたない。
\item $\mathcal{F}$ は埋め込まれた随伴点をもたない。
\item $\mathcal{F}$ の台は $X$ である。
\item 切断 $s$ は正則有理型である。
\end{enumerate}
この状況で $\mathcal{I} \subset \mathcal{O}_X$ を $s$ の分母のイデアルとする。
『因子』の定義 \ref{divisors-definition-regular-meromorphic-ideal-denominators} を参照。
このとき次が成り立つ。
\begin{enumerate}
\item 次の短完全列が存在する。
$$
\begin{matrix}
0 &
\to &
\mathcal{I}\mathcal{F} &
\xrightarrow{1} &
\mathcal{F} &
\to &
\mathcal{Q}_1 &
\to &
0 \\
0 &
\to &
\mathcal{I}\mathcal{F} &
\xrightarrow{s} &
\mathcal{F} \otimes_{\mathcal{O}_X} \mathcal{L} &
\to &
\mathcal{Q}_2 &
\to &
0
\end{matrix}
$$
\item 連接層 $\mathcal{Q}_1$, $\mathcal{Q}_2$ の台の
$\delta$-次元は $\leq k$ である。
\item 切断 $s$ は、$X$ の $\delta$-次元 $k + 1$ の各既約成分 $X_i$ 上で
正則有理型切断 $s_i$ に制限される。
\item $[\mathcal{F}]_{k + 1} = \sum m_i[X_i]$ と書くと、
$Z_k(X)$ において
$$
[\mathcal{Q}_2]_k - [\mathcal{Q}_1]_k
=
\sum m_i(X_i \to X)_*\text{div}_{\mathcal{L}|_{X_i}}(s_i)
$$
であり、特に $\CH_k(X)$ において
$$
[\mathcal{Q}_2]_k - [\mathcal{Q}_1]_k
=
c_1(\mathcal{L}) \cap [\mathcal{F}]_{k + 1}
$$
である。
\end{enumerate}
\end{lemma}

\begin{proof}
『因子』の補題 \ref{divisors-lemma-make-maps-regular-section} から、
余核の台が閉で至る所稠密でない部分集合 $T$ に含まれる単射
$1 : \mathcal{I}\mathcal{F} \to \mathcal{F}$ および
$s : \mathcal{I}\mathcal{F} \to \mathcal{F} \otimes_{\mathcal{O}_X}\mathcal{L}$
が存在することを思い出そう。その余核を補題のとおり $\mathcal{Q}_i$ と表す。
すると $\dim_\delta(\text{Supp}(\mathcal{Q}_i)) \leq k$ が従う。
『因子』の補題 \ref{divisors-lemma-pullback-meromorphic-sections-defined}
および \ref{divisors-lemma-meromorphic-sections-pullback} により、
引き戻し $s_i$ は定義され、$\mathcal{L}|_{X_i}$ の正則有理型切断である。
(4) のサイクルの等式は、(4) のサイクル類の等式を含意する。
したがって、残るのは
$$
[\mathcal{Q}_2]_k - [\mathcal{Q}_1]_k
=
\sum m_i(X_i \to X)_*\text{div}_{\mathcal{L}|_{X_i}}(s_i)
$$
が $Z_k(X)$ において成り立つことを示すことだけである。
これを見るため、$Z \subset X$ を $\delta$-次元 $k$ の整閉部分スキームとし、
$\xi \in Z$ を生成点とする。
$A = \mathcal{O}_{X, \xi}$ および $M = \mathcal{F}_\xi$ とおく。
さらに、生成元 $s_\xi \in \mathcal{L}_\xi$ を選ぶ。
このとき $a, b \in A$ を非零因子として
$s = (a/b) s_\xi$ と書ける。この場合
$I = \mathcal{I}_\xi = \{x \in A \mid x(a/b) \in A\}$ である。
左辺における $[Z]$ の係数は
$$
\text{length}_A(M/(a/b)IM) - \text{length}_A(M/IM)
$$
であり、右辺における $[Z]$ の係数は
$$
\sum
\text{length}_{A_{\mathfrak q_i}}(M_{\mathfrak q_i})
\text{ord}_{A/\mathfrak q_i}(a/b)
$$
である。ここで $\mathfrak q_1, \ldots, \mathfrak q_t$ は
$1$-次元局所環 $A$ の極小素イデアルである。
したがって補題 \ref{chow-lemma-no-embedded-points-modules} から結論が従う。
\end{proof}

\begin{lemma}
\label{chow-lemma-coherent-sheaf-cap-c1}
$(S, \delta)$ を状況 \ref{chow-situation-setup} のとおりとする。
$X$ を $S$ 上局所有限型とする。
$\mathcal{L}$ を可逆 $\mathcal{O}_X$-加群とする。
$\mathcal{F}$ を連接 $\mathcal{O}_X$-加群とする。
$\dim_\delta(\text{Supp}(\mathcal{F})) \leq k + 1$ と仮定する。
このとき、元
$$
[\mathcal{F} \otimes_{\mathcal{O}_X} \mathcal{L}]
-
[\mathcal{F}]
\in
K_0(\textit{Coh}_{\leq k + 1}(X)/\textit{Coh}_{\leq k - 1}(X))
$$
は補題 \ref{chow-lemma-cycles-rational-equivalence-K-group} の部分群 $B_k(X)$ に属し、
写像 $B_k(X) \to \CH_k(X)$ により
元 $c_1(\mathcal{L}) \cap [\mathcal{F}]_{k + 1}$ に写る。
\end{lemma}

\begin{proof}
『因子』の補題 \ref{divisors-lemma-remove-embedded-points} で構成された短完全列
$$
0 \to \mathcal{K} \to \mathcal{F} \to \mathcal{F}' \to 0
$$
をとる。これは特に、$\mathcal{F}'$ が埋め込まれた随伴点をもたないことを意味する。
$\mathcal{K}$ の台は $\mathcal{F}$ の台の中で至る所稠密でないので、
$\dim_\delta(\text{Supp}(\mathcal{K})) \leq k$ である。
$\mathcal{K}$ から始めて
『因子』の補題 \ref{divisors-lemma-remove-embedded-points} を再び適用すると、
短完全列
$$
0 \to \mathcal{K}'' \to \mathcal{K} \to \mathcal{K}' \to 0
$$
を得る。ここで $\dim_\delta(\text{Supp}(\mathcal{K}'')) < k$ であり、
$\mathcal{K}'$ は埋め込まれた随伴点をもたない。
連接層 $\mathcal{F}'$ と $\mathcal{K}'$ に対して補題を証明できたと仮定する。
すると、等式
$$
[\mathcal{F}]_{k + 1}
=
[\mathcal{F}']_{k + 1}
+ [\mathcal{K}']_{k + 1}
+ [\mathcal{K}'']_{k + 1}
$$
（補題 \ref{chow-lemma-additivity-sheaf-cycle} を用いる）、
$$
[\mathcal{F} \otimes_{\mathcal{O}_X} \mathcal{L}]
-
[\mathcal{F}]
=
[\mathcal{F}' \otimes_{\mathcal{O}_X} \mathcal{L}]
-
[\mathcal{F}']
+
[\mathcal{K}' \otimes_{\mathcal{O}_X} \mathcal{L}]
-
[\mathcal{K}']
+
[\mathcal{K}'' \otimes_{\mathcal{O}_X} \mathcal{L}]
-
[\mathcal{K}'']
$$
（$\otimes \mathcal{L}$ が完全であることを用いる）、
および $[\mathcal{K}'']_{k + 1}$ と
$[\mathcal{K}'' \otimes_{\mathcal{O}_X} \mathcal{L}]
- [\mathcal{K}'']$ が
$K_0(\textit{Coh}_{\leq k + 1}(X)/\textit{Coh}_{\leq k - 1}(X))$
において自明に零であることから、$\mathcal{F}$ に対して結論が成り立つ。
したがって、層 $\mathcal{F}$ は埋め込まれた随伴点をもたないと仮定してよい。

\medskip\noindent
$X$, $\mathcal{F}$ を補題のとおりとし、さらに $\mathcal{F}$ は
埋め込まれた随伴点をもたないと仮定する。
『因子』の補題 \ref{divisors-lemma-no-embedded-points-endos} で構成された
イデアル層 $\mathcal{I} \subset \mathcal{O}_X$、対応する閉部分スキーム
$i : Z \to X$、および連接 $\mathcal{O}_Z$-加群 $\mathcal{G}$ を考える。
$Z$ は埋め込まれた点をもたない局所 Noether スキームであり、
$\mathcal{G}$ は埋め込まれた随伴点をもたない連接層で、
$\text{Supp}(\mathcal{G}) = Z$ かつ $i_*\mathcal{G} = \mathcal{F}$ であることを思い出そう。
さらに $\mathcal{N} = \mathcal{L}|_Z$ とおく。

\medskip\noindent
『因子』の補題 \ref{divisors-lemma-regular-meromorphic-section-exists} により、
可逆層 $\mathcal{N}$ は $Z$ 上に正則有理型切断 $s$ をもつ。
$\mathcal{J} \subset \mathcal{O}_Z$ を $s$ の分母の層と表す。
補題 \ref{chow-lemma-no-embedded-points} により、短完全列
$$
\begin{matrix}
0 &
\to &
\mathcal{J}\mathcal{G} &
\xrightarrow{1} &
\mathcal{G} &
\to &
\mathcal{Q}_1 &
\to &
0 \\
0 &
\to &
\mathcal{J}\mathcal{G} &
\xrightarrow{s} &
\mathcal{G} \otimes_{\mathcal{O}_Z} \mathcal{N} &
\to &
\mathcal{Q}_2 &
\to &
0
\end{matrix}
$$
が存在し、$\dim_\delta(\text{Supp}(\mathcal{Q}_i)) \leq k$ であり、
サイクル
$
[\mathcal{Q}_2]_k - [\mathcal{Q}_1]_k
$
は $c_1(\mathcal{N}) \cap [\mathcal{G}]_{k + 1}$ の代表元である。
射影公式により
$i_*(\mathcal{G} \otimes \mathcal{N}) = \mathcal{F} \otimes \mathcal{L}$
であることを用いると
（『層のコホモロジー』の補題 \ref{cohomology-lemma-projection-formula} 参照）、
$$
[\mathcal{F} \otimes_{\mathcal{O}_X} \mathcal{L}]
-
[\mathcal{F}]
=
[i_*\mathcal{Q}_2] - [i_*\mathcal{Q}_1]
$$
が
$K_0(\textit{Coh}_{\leq k + 1}(X)/\textit{Coh}_{\leq k - 1}(X))$
において成り立つ。これだけで
$[\mathcal{F} \otimes_{\mathcal{O}_X} \mathcal{L}] - [\mathcal{F}]$
が $B_k(X)$ の元であることが分かる。さらに
\begin{eqnarray*}
[i_*\mathcal{Q}_2]_k - [i_*\mathcal{Q}_1]_k
& = &
i_*\left( [\mathcal{Q}_2]_k - [\mathcal{Q}_1]_k \right) \\
& = &
i_*\left(c_1(\mathcal{N}) \cap [\mathcal{G}]_{k + 1} \right) \\
& = &
c_1(\mathcal{L}) \cap i_*[\mathcal{G}]_{k + 1} \\
& = &
c_1(\mathcal{L}) \cap [\mathcal{F}]_{k + 1}
\end{eqnarray*}
である。これは上の結果と補題 \ref{chow-lemma-pushforward-cap-c1} および
\ref{chow-lemma-cycle-push-sheaf} による。また定義により、
写像 $B_k(X) \to \CH_k(X)$ のもとでのこの元の像と一致する。
したがって補題が証明された。
\end{proof}
